% ============================================================
%  四合一多版本条件编译引擎（\ifteacher × \ifbasic）
% ------------------------------------------------------------
%  用法：在 main.tex 导言区最后 % ============================================================
%  四合一多版本条件编译引擎（\ifteacher × \ifbasic）
% ------------------------------------------------------------
%  用法：在 main.tex 导言区最后 % ============================================================
%  四合一多版本条件编译引擎（\ifteacher × \ifbasic）
% ------------------------------------------------------------
%  用法：在 main.tex 导言区最后 % ============================================================
%  四合一多版本条件编译引擎（\ifteacher × \ifbasic）
% ------------------------------------------------------------
%  用法：在 main.tex 导言区最后 \input{versions}，
%        再设置两行开关，然后 \input{main_body}。
%
%  四种组合（改两行即可，无需复制任何文件）：
%    \teachertrue  \basicfalse   完整教师版（默认母版）
%    \teacherfalse \basicfalse   完整学生版（隐藏解析与填空答案）
%    \teachertrue  \basictrue    薄弱生教师版（显化脚手架 + 完整答案）
%    \teacherfalse \basictrue    薄弱生学生版（保留 \tk 留白横线）
%
%  本文件只依赖已加载的宏包（xcolor / tcolorbox / environ），
%  不修改 common/，故对 ../hv-bank 等其他分册零影响。
% ============================================================

% ---------- 布尔开关 ----------
\newif\ifteacher          % 教师解析版 / 学生留白版
\newif\ifbasic            % 基础通关版 / 完整高观点版
\newif\ifdemomode         % 是否挂载示范章节 chapters/demo_multiversion

% 安全默认（main.tex 的两行开关覆盖此处）
\teachertrue
\basicfalse
\demomodefalse

% ---------- 零残留隐藏所需 ----------
\usepackage{environ}      % \NewEnviron（隐藏时不留空白）

% ---------- 版本标签（封面 / 页眉自动显示） ----------
\newcommand{\versionlabel}{%
  \ifteacher 教师解析版\else 学生留白版\fi
  \,\textperiodcentered\,
  \ifbasic 基础通关版\else 完整高观点版\fi}

% ============================================================
%  一、版本内容分流环境
% ============================================================

% 仅完整版渲染；薄弱生版彻底忽略（不占位）
% 用途：莫比乌斯矩阵、射影同构、压轴第 3–4 阶梯难题等
\NewEnviron{advancedonly}{%
  \ifbasic\else
    \par\medskip
    \BODY
    \par\medskip
  \fi
}

% 仅薄弱生版渲染；完整版彻底忽略
% 用途：课前 15 分钟计算过关、大白话口诀、半成品步步填空
\NewEnviron{basiconly}{%
  \ifbasic
    \par\medskip
    \BODY
    \par\medskip
  \fi
}

% ============================================================
%  二、双模填空宏 \tk{...} 与内联版本选择 \vchoice{}{}
% ============================================================

% 教师版：红色下划线答案；学生版：等宽下划线留白（版面稳定）
% 内容一律写成裸数学式（不带 $）。\underline 继承当前模式，
% 而 \sbox 的盒子内容始终是文本模式、必须自带数学定界符。
\newcommand{\tk}[1]{%
  \ifteacher
    \ifmmode
      {\color{second}\underline{#1}}%
    \else
      {\color{second}\underline{$\displaystyle#1$}}%
    \fi
  \else
    \sbox0{$\displaystyle#1$}%
    \underline{\hspace{\wd0}}%
  \fi
}

% 同一句子里的措辞切换：完整版文本 / 薄弱生大白话文本
\newcommand{\vchoice}[2]{%
  \ifbasic #2\else #1\fi}

% ============================================================
%  三、解答环境重定义（教师版带边框【剖析与标准解答】）
% ============================================================
\RenewDocumentEnvironment{solution}{+b}{%
  \ifteacher
    \par\medskip
    \begin{tcolorbox}[
      breakable, enhanced, sharp corners,
      colback=main!3, colframe=main!60!black, boxrule=0.5pt,
      title={\bfseries 剖析与标准解答}, coltitle=black, fonttitle=\small,
      top=5pt,
    ]
    #1
    \end{tcolorbox}%
  \else
    % 学生版：不留内容，只保留少量书写留白
    \par\noteblank[6ex]%
  \fi
}{}

% ============================================================
%  四、语义化盒子（corebox / methodbox / trapbox / warmupbox）
% ============================================================

% 核心概念盒（绿）
\newtcolorbox{corebox}[1]{
  breakable, enhanced, sharp corners,
  colback=main!4, colframe=main!60!black, boxrule=0.6pt,
  title={\bfseries #1}, coltitle=black, fonttitle=\small,
  attach boxed title to top left={xshift=2mm, yshift=-2mm},
  boxed title style={sharp corners, colback=white,
    colframe=main!60!black, boxrule=0.4pt},
  top=7pt,
}

% 方法归纳盒（青）
\newtcolorbox{methodbox}[1]{
  breakable, enhanced, sharp corners,
  colback=third!5, colframe=third!60!black, boxrule=0.5pt,
  title={\bfseries #1}, coltitle=black, fonttitle=\small,
  attach boxed title to top left={xshift=2mm, yshift=-2mm},
  boxed title style={sharp corners, colback=white,
    colframe=third!60!black, boxrule=0.4pt},
  top=7pt,
}

% 易错陷阱盒（红）
\newtcolorbox{trapbox}[1]{
  breakable, enhanced, sharp corners,
  colback=second!4, colframe=second!65!black, boxrule=0.5pt,
  title={\bfseries 陷阱：#1}, coltitle=black, fonttitle=\small,
  attach boxed title to top left={xshift=2mm, yshift=-2mm},
  boxed title style={sharp corners, colback=white,
    colframe=second!65!black, boxrule=0.4pt},
  top=7pt,
}

% 课前热身盒（蓝）
\newtcolorbox{warmupbox}[1]{
  breakable, enhanced, sharp corners,
  colback=structurecolor!4, colframe=structurecolor!70!black, boxrule=0.5pt,
  title={\bfseries #1}, coltitle=black, fonttitle=\small,
  top=6pt,
}

% ============================================================
%  五、按章（节）自动编号的 example / exercise 环境
%  注：elegantnote 自带 example 计数器（\c@example），故另立
%      hvexample / hvexercise 计数，避免与模板冲突；编号仍按节累计。
% ============================================================
\newcounter{hvexample}[section]
\renewcommand{\thehvexample}{\thesection.\arabic{hvexample}}

\RenewDocumentEnvironment{example}{O{}}{%
  \refstepcounter{hvexample}%
  \par\medskip\noindent
  {\bfseries\color{main}例题\thehvexample\IfValueT{#1}{：#1}\quad}%
}{\par\medskip}

\newcounter{hvexercise}[section]
\renewcommand{\thehvexercise}{\thesection.\arabic{hvexercise}}

\NewDocumentEnvironment{exercise}{O{}}{%
  \refstepcounter{hvexercise}%
  \par\medskip\noindent
  {\bfseries\color{third}练习\thehvexercise\IfValueT{#1}{：#1}\quad}%
}{\par\medskip}

% ============================================================
%  六、七问连环母题（Scaffolded Master Problem）载体环境
% ------------------------------------------------------------
%  定位：一道母题挂载 (1)–(7) 个递进小问，学生做题留白处写下的
%        推理本身就是这一章的学习笔记。写作规范见 ARCHITECTURE.md。
% ============================================================

\newcounter{hvmasterq}[section]
\renewcommand{\thehvmasterq}{\thesection.\arabic{hvmasterq}}

% 母题盒：标题自动编号；盒内用 enumerate[label=(\arabic*)] 写七问
% 注意：计数器由调用方在 \begin{masterq} 前手动 \refstepcounter{hvmasterq}，
%       不要把 \refstepcounter 写进本定义的选项里（tcolorbox 会反复展开选项，
%       导致编译卡在分页计算）。
\newtcolorbox{masterq}[1]{
  breakable, enhanced, sharp corners,
  colback=structurecolor!3, colframe=structurecolor!70!black, boxrule=0.6pt,
  title={\bfseries 七问连环母题\thehvmasterq：#1},
  coltitle=black, fonttitle=\small,
  attach boxed title to top left={xshift=2mm, yshift=-2mm},
  boxed title style={sharp corners, colback=white,
    colframe=structurecolor!70!black, boxrule=0.4pt},
  top=7pt,
}

% 每问末尾的笔记归档标签：\logotag{存档类别}（如「核心公理」「大招公式」）
\newcommand{\logotag}[1]{%
  \hfill\textcolor{third}{\small\ttfamily【笔记归档·#1】}}
，
%        再设置两行开关，然后 % ============================================================
%  main_body.tex —— 正文主体（由 main.tex 在设置好版本开关后 \input）
%  拆分原因：四合一多版本编译时，包装文件只需复写「两行开关」，
%  再 \input{main_body}，避免导言区内容四处重复、彼此漂移。
% ============================================================

\title{数学导学案·高观点系列}
\author{Shawn Shang}
\date{2026 年 11 月}

\graphicspath{{../figures/}{../map/figures/}}

\begin{document}

\begin{titlepage}
  \thispagestyle{empty}
  \centering
  \makebox[\textwidth]{%
    \includegraphics[width=\paperwidth]{cover_cyan.png}}\par
  \vfill
  {\Huge\bfseries\color{ecolor!70!black} 数学导学案·高观点系列\par}
  \vskip 1.2em
  {\large 结构主义视角下的高中数学\par}
  \vskip 0.8em
  {\large 第 0 章 — 第 9 章\par}
  \vskip 0.8em
  {\large\versionlabel\par}
  \vskip 2.5em
  {\large Shawn Shang\par}
  \vfill
  {\normalsize 2026 年 11 月\par}
\end{titlepage}
\tableofcontents

% 让第一个 \section 编号为 0（第 0 章是总纲）
\setcounter{section}{-1}

\input{chapters/00}
\input{chapters/01}
\input{chapters/02}
\input{chapters/03}
\input{chapters/04}
\input{chapters/05}
\input{chapters/06}
\input{chapters/07}
\input{chapters/08_homogenization}
\input{chapters/09_fixed_point}

% 示范章节（仅当 main.tex 中 \demomotrue 时挂载）
\ifdemomode
  \input{chapters/demo_multiversion}
\fi

\appendix
\input{chapters/appA_legalization}
\input{chapters/appB_honors_bridge}

\end{document}
。
%
%  四种组合（改两行即可，无需复制任何文件）：
%    \teachertrue  \basicfalse   完整教师版（默认母版）
%    \teacherfalse \basicfalse   完整学生版（隐藏解析与填空答案）
%    \teachertrue  \basictrue    薄弱生教师版（显化脚手架 + 完整答案）
%    \teacherfalse \basictrue    薄弱生学生版（保留 \tk 留白横线）
%
%  本文件只依赖已加载的宏包（xcolor / tcolorbox / environ），
%  不修改 common/，故对 ../hv-bank 等其他分册零影响。
% ============================================================

% ---------- 布尔开关 ----------
\newif\ifteacher          % 教师解析版 / 学生留白版
\newif\ifbasic            % 基础通关版 / 完整高观点版
\newif\ifdemomode         % 是否挂载示范章节 chapters/demo_multiversion

% 安全默认（main.tex 的两行开关覆盖此处）
\teachertrue
\basicfalse
\demomodefalse

% ---------- 零残留隐藏所需 ----------
\usepackage{environ}      % \NewEnviron（隐藏时不留空白）

% ---------- 版本标签（封面 / 页眉自动显示） ----------
\newcommand{\versionlabel}{%
  \ifteacher 教师解析版\else 学生留白版\fi
  \,\textperiodcentered\,
  \ifbasic 基础通关版\else 完整高观点版\fi}

% ============================================================
%  一、版本内容分流环境
% ============================================================

% 仅完整版渲染；薄弱生版彻底忽略（不占位）
% 用途：莫比乌斯矩阵、射影同构、压轴第 3–4 阶梯难题等
\NewEnviron{advancedonly}{%
  \ifbasic\else
    \par\medskip
    \BODY
    \par\medskip
  \fi
}

% 仅薄弱生版渲染；完整版彻底忽略
% 用途：课前 15 分钟计算过关、大白话口诀、半成品步步填空
\NewEnviron{basiconly}{%
  \ifbasic
    \par\medskip
    \BODY
    \par\medskip
  \fi
}

% ============================================================
%  二、双模填空宏 \tk{...} 与内联版本选择 \vchoice{}{}
% ============================================================

% 教师版：红色下划线答案；学生版：等宽下划线留白（版面稳定）
% 内容一律写成裸数学式（不带 $）。\underline 继承当前模式，
% 而 \sbox 的盒子内容始终是文本模式、必须自带数学定界符。
\newcommand{\tk}[1]{%
  \ifteacher
    \ifmmode
      {\color{second}\underline{#1}}%
    \else
      {\color{second}\underline{$\displaystyle#1$}}%
    \fi
  \else
    \sbox0{$\displaystyle#1$}%
    \underline{\hspace{\wd0}}%
  \fi
}

% 同一句子里的措辞切换：完整版文本 / 薄弱生大白话文本
\newcommand{\vchoice}[2]{%
  \ifbasic #2\else #1\fi}

% ============================================================
%  三、解答环境重定义（教师版带边框【剖析与标准解答】）
% ============================================================
\RenewDocumentEnvironment{solution}{+b}{%
  \ifteacher
    \par\medskip
    \begin{tcolorbox}[
      breakable, enhanced, sharp corners,
      colback=main!3, colframe=main!60!black, boxrule=0.5pt,
      title={\bfseries 剖析与标准解答}, coltitle=black, fonttitle=\small,
      top=5pt,
    ]
    #1
    \end{tcolorbox}%
  \else
    % 学生版：不留内容，只保留少量书写留白
    \par\noteblank[6ex]%
  \fi
}{}

% ============================================================
%  四、语义化盒子（corebox / methodbox / trapbox / warmupbox）
% ============================================================

% 核心概念盒（绿）
\newtcolorbox{corebox}[1]{
  breakable, enhanced, sharp corners,
  colback=main!4, colframe=main!60!black, boxrule=0.6pt,
  title={\bfseries #1}, coltitle=black, fonttitle=\small,
  attach boxed title to top left={xshift=2mm, yshift=-2mm},
  boxed title style={sharp corners, colback=white,
    colframe=main!60!black, boxrule=0.4pt},
  top=7pt,
}

% 方法归纳盒（青）
\newtcolorbox{methodbox}[1]{
  breakable, enhanced, sharp corners,
  colback=third!5, colframe=third!60!black, boxrule=0.5pt,
  title={\bfseries #1}, coltitle=black, fonttitle=\small,
  attach boxed title to top left={xshift=2mm, yshift=-2mm},
  boxed title style={sharp corners, colback=white,
    colframe=third!60!black, boxrule=0.4pt},
  top=7pt,
}

% 易错陷阱盒（红）
\newtcolorbox{trapbox}[1]{
  breakable, enhanced, sharp corners,
  colback=second!4, colframe=second!65!black, boxrule=0.5pt,
  title={\bfseries 陷阱：#1}, coltitle=black, fonttitle=\small,
  attach boxed title to top left={xshift=2mm, yshift=-2mm},
  boxed title style={sharp corners, colback=white,
    colframe=second!65!black, boxrule=0.4pt},
  top=7pt,
}

% 课前热身盒（蓝）
\newtcolorbox{warmupbox}[1]{
  breakable, enhanced, sharp corners,
  colback=structurecolor!4, colframe=structurecolor!70!black, boxrule=0.5pt,
  title={\bfseries #1}, coltitle=black, fonttitle=\small,
  top=6pt,
}

% ============================================================
%  五、按章（节）自动编号的 example / exercise 环境
%  注：elegantnote 自带 example 计数器（\c@example），故另立
%      hvexample / hvexercise 计数，避免与模板冲突；编号仍按节累计。
% ============================================================
\newcounter{hvexample}[section]
\renewcommand{\thehvexample}{\thesection.\arabic{hvexample}}

\RenewDocumentEnvironment{example}{O{}}{%
  \refstepcounter{hvexample}%
  \par\medskip\noindent
  {\bfseries\color{main}例题\thehvexample\IfValueT{#1}{：#1}\quad}%
}{\par\medskip}

\newcounter{hvexercise}[section]
\renewcommand{\thehvexercise}{\thesection.\arabic{hvexercise}}

\NewDocumentEnvironment{exercise}{O{}}{%
  \refstepcounter{hvexercise}%
  \par\medskip\noindent
  {\bfseries\color{third}练习\thehvexercise\IfValueT{#1}{：#1}\quad}%
}{\par\medskip}

% ============================================================
%  六、七问连环母题（Scaffolded Master Problem）载体环境
% ------------------------------------------------------------
%  定位：一道母题挂载 (1)–(7) 个递进小问，学生做题留白处写下的
%        推理本身就是这一章的学习笔记。写作规范见 ARCHITECTURE.md。
% ============================================================

\newcounter{hvmasterq}[section]
\renewcommand{\thehvmasterq}{\thesection.\arabic{hvmasterq}}

% 母题盒：标题自动编号；盒内用 enumerate[label=(\arabic*)] 写七问
% 注意：计数器由调用方在 \begin{masterq} 前手动 \refstepcounter{hvmasterq}，
%       不要把 \refstepcounter 写进本定义的选项里（tcolorbox 会反复展开选项，
%       导致编译卡在分页计算）。
\newtcolorbox{masterq}[1]{
  breakable, enhanced, sharp corners,
  colback=structurecolor!3, colframe=structurecolor!70!black, boxrule=0.6pt,
  title={\bfseries 七问连环母题\thehvmasterq：#1},
  coltitle=black, fonttitle=\small,
  attach boxed title to top left={xshift=2mm, yshift=-2mm},
  boxed title style={sharp corners, colback=white,
    colframe=structurecolor!70!black, boxrule=0.4pt},
  top=7pt,
}

% 每问末尾的笔记归档标签：\logotag{存档类别}（如「核心公理」「大招公式」）
\newcommand{\logotag}[1]{%
  \hfill\textcolor{third}{\small\ttfamily【笔记归档·#1】}}
，
%        再设置两行开关，然后 % ============================================================
%  main_body.tex —— 正文主体（由 main.tex 在设置好版本开关后 \input）
%  拆分原因：四合一多版本编译时，包装文件只需复写「两行开关」，
%  再 % ============================================================
%  main_body.tex —— 正文主体（由 main.tex 在设置好版本开关后 \input）
%  拆分原因：四合一多版本编译时，包装文件只需复写「两行开关」，
%  再 \input{main_body}，避免导言区内容四处重复、彼此漂移。
% ============================================================

\title{数学导学案·高观点系列}
\author{Shawn Shang}
\date{2026 年 11 月}

\graphicspath{{../figures/}{../map/figures/}}

\begin{document}

\begin{titlepage}
  \thispagestyle{empty}
  \centering
  \makebox[\textwidth]{%
    \includegraphics[width=\paperwidth]{cover_cyan.png}}\par
  \vfill
  {\Huge\bfseries\color{ecolor!70!black} 数学导学案·高观点系列\par}
  \vskip 1.2em
  {\large 结构主义视角下的高中数学\par}
  \vskip 0.8em
  {\large 第 0 章 — 第 9 章\par}
  \vskip 0.8em
  {\large\versionlabel\par}
  \vskip 2.5em
  {\large Shawn Shang\par}
  \vfill
  {\normalsize 2026 年 11 月\par}
\end{titlepage}
\tableofcontents

% 让第一个 \section 编号为 0（第 0 章是总纲）
\setcounter{section}{-1}

\input{chapters/00}
\input{chapters/01}
\input{chapters/02}
\input{chapters/03}
\input{chapters/04}
\input{chapters/05}
\input{chapters/06}
\input{chapters/07}
\input{chapters/08_homogenization}
\input{chapters/09_fixed_point}

% 示范章节（仅当 main.tex 中 \demomotrue 时挂载）
\ifdemomode
  \input{chapters/demo_multiversion}
\fi

\appendix
\input{chapters/appA_legalization}
\input{chapters/appB_honors_bridge}

\end{document}
，避免导言区内容四处重复、彼此漂移。
% ============================================================

\title{数学导学案·高观点系列}
\author{Shawn Shang}
\date{2026 年 11 月}

\graphicspath{{../figures/}{../map/figures/}}

\begin{document}

\begin{titlepage}
  \thispagestyle{empty}
  \centering
  \makebox[\textwidth]{%
    \includegraphics[width=\paperwidth]{cover_cyan.png}}\par
  \vfill
  {\Huge\bfseries\color{ecolor!70!black} 数学导学案·高观点系列\par}
  \vskip 1.2em
  {\large 结构主义视角下的高中数学\par}
  \vskip 0.8em
  {\large 第 0 章 — 第 9 章\par}
  \vskip 0.8em
  {\large\versionlabel\par}
  \vskip 2.5em
  {\large Shawn Shang\par}
  \vfill
  {\normalsize 2026 年 11 月\par}
\end{titlepage}
\tableofcontents

% 让第一个 \section 编号为 0（第 0 章是总纲）
\setcounter{section}{-1}

% 高观点导学案：集合、逻辑演算与证明武器库（第 0 章·总纲）
% 主线：把“集合与逻辑”从概念背诵升级为贯穿高中解题的“题干同解变形与严格证明通用武器库”
% 四模块：集合语言与同解变形公理 → 命题演算三大变形引擎 → 两大间接证明支柱 → 武器库索引表与跨模块实战
% 全程用“集合包含 ↔ 逻辑等价 ↔ 量词否命题”三种语言互译，并用韦恩图与多米诺图把抽象逻辑可视化。

\section{集合、逻辑演算与证明武器库（高观点总纲）}

本章不引入新公式，而是把教材第 1 章的集合、常用逻辑用语重新锻造为一套“解题武器库”：
面对题干，先问它是哪一种\key{对象集合}（解集、值域、参数范围），再问它需要哪一种\key{变形引擎}
（同解变形、德摩根律、逆否等价、补集转化、反证法、数学归纳法）。四个模块环环相扣：
\key{集合语言与同解变形公理}（$A=B\iff A\subseteq B\land B\subseteq A$）$\to$
\key{命题演算三大变形引擎}（德摩根律、逆否等价、量词否定）$\to$
\key{两大间接证明支柱}（反证法、数学归纳法）$\to$
\key{武器库索引表与跨模块实战}。
核心思想：\key{题干中的每一个关键词，都对应一次集合运算或一次逻辑等价变形}；
解题失误往往不是计算错误，而是\key{没审清变形方向}（增根或漏根）或\key{没找对证明支柱}。

\begin{learnobjective}
\begin{itemize}
  \item 掌握集合相等判定公理 $A=B\iff(A\subseteq B)\land(B\subseteq A)$，能用它审查方程、不等式变形的同解性；
  \item 熟记同解变形防增漏根清单：平方（范围扩大易增根）、同除含参因式（范围缩小易漏根）、对数真数（定义域子集）；
  \item 掌握德摩根律 $\neg(p\land q)\iff\neg p\lor\neg q$ 与逆否等价 $(p\Rightarrow q)\iff(\neg q\Rightarrow\neg p)$，能用补集转化处理“至多、至少、不都、无解”；
  \item 掌握量词否定 $\neg\forall=\exists\neg$、$\neg\exists=\forall\neg$，会把“恒成立／能成立”翻译为值域的子集包含／交集非空；
  \item 理解反证法即证 $A\cap\complement_U B=\varnothing$，识别“唯一性、不存在、无理数、无限性”四类触发信号词；
  \item 理解数学归纳法的皮亚诺公理与良序原理本质，会用最小反例原理分析“归纳基础＋归纳递推”两步的缺一不可。
\end{itemize}
\end{learnobjective}

\begin{preview}
\subsection*{知识回顾}
\begin{itemize}
  \item $A\subseteq B$ 的定义是：对任意 $a\in A$，都有 \blank；
  \item 若 $A=\{x\mid x^2=1\}$，$B=\{-1,1\}$，则 $A$ 与 $B$ 的关系是 \blank（相等／互不包含）；
  \item 命题“若 $p$ 则 $q$”的逆否命题是 \blank；
  \item 德摩根律：$\neg(p\land q)\iff$ \blank。
\end{itemize}

\subsection*{教材速读}
阅读教材第 1 章“集合与常用逻辑用语”及本章导学案，完成下列问题：
\begin{enumerate}
  \item 补集 $\complement_U A$ 的定义是 $\{x\in U\mid \blank\}$；
  \item “$p$ 是 $q$ 的充分条件”用集合语言说是 \blank；
  \item “方程 $f(x)=0$ 无解”用集合语言说是 \blank；
  \item 数学归纳法证明的两步是 \blank 与 \blank。
\end{enumerate}
\end{preview}

\begin{explore}

\subsection{集合语言与同解变形公理}

\key{探究一：集合相等——互相包含才是相等}
集合相等不是“看起来一样”，而是\key{互相包含}：
\[
  A=B\iff(A\subseteq B)\land(B\subseteq A).
\]
证明 $A=B$ 必须证两个方向；只证 $A\subseteq B$ 只能得到包含，不能得到相等。
这条公理还是审查方程（不等式）变形的“合法性标准”：
把原方程的解集记为 $A$，变形后方程的解集记为 $B$，则
\begin{itemize}
  \item $A=B$：变形\key{同解}，解集不变，安全；
  \item $A\subsetneqq B$：变形后解集\key{变大}，多出的解是\key{增根}；
  \item $B\subsetneqq A$：变形后解集\key{变小}，丢掉的解是\key{漏根}。
\end{itemize}
因此解方程、解不等式的每一步都要问一句：\key{这一步把解集往哪个方向变了？}

\begin{center}
\footnotesize
\setlength{\tabcolsep}{3pt}
\renewcommand{\arraystretch}{1.4}
\begin{tabular}{p{2.6cm}p{2.8cm}p{3.2cm}p{3.2cm}}
\toprule
变形操作 & 解集关系 & 方向与风险 & 应对 \\
\midrule
两边加/减同一整式 & $A=B$ & 同解，安全 & 直接使用 \\
两边乘/除以非零常数 & $A=B$ & 同解，安全 & 先确认非零 \\
两边平方（或去绝对值） & $A\subsetneqq B$ & 范围扩大，易增根 & 必须验根 \\
两边同除含参/未知因式 & $B\subsetneqq A$ & 范围缩小，易漏根 & 移项分解，讨论零 \\
取对数/解对数方程 & $A\subsetneqq B$ & 真数 $>0$ 限制，易增根 & 先限定定义域 \\
\bottomrule
\end{tabular}
\end{center}

\key{风险一：平方变形范围扩大，易增根。}
由 $f(x)=g(x)$ 平方得 $f^2(x)=g^2(x)$，而
\[
  f^2=g^2\iff(f-g)(f+g)=0\iff f=g\ \text{或}\ f=-g,
\]
凭空多出 $f=-g$ 的解。必须验根。

\begin{example}
解方程 $\sqrt{x+1}=x-1$，并指出直接平方会多出哪个根。
\end{example}
\begin{solution}
两边平方（须先注意右边非负：$x-1\geq0$，即 $x\geq1$）：
\[
  x+1=(x-1)^2=x^2-2x+1\implies x^2-3x=0\implies x=0\ \text{或}\ x=3.
\]
检验：$x=0$ 时左边 $=1$，右边 $=-1$，不相等，是平方后产生的\key{增根}，舍去；
$x=3$ 时左边 $=2$，右边 $=2$，成立。故解集为 $\{3\}$。
\end{solution}

\key{风险二：同除含参因式范围缩小，易漏根。}
方程两边同除以含未知数（或参数）的因式，若该因式可能为 $0$，就会丢掉使因式为 $0$ 的解。
正确做法是移项因式分解，并单独讨论因式为 $0$ 的情形。

\begin{example}
解关于 $x$ 的方程 $x^2=ax$（$a$ 为参数），并说明“两边同除以 $x$”会漏掉什么。
\end{example}
\begin{solution}
移项分解：$x(x-a)=0$，故 $x=0$ 或 $x=a$。
当 $a\neq0$ 时解集为 $\{0,a\}$；当 $a=0$ 时方程为 $x^2=0$，解集为 $\{0\}$（两个相等实根）。
若直接两边同除以 $x$，得到 $x=a$，会漏掉 $x=0$（$a\neq0$ 时），
因为 $x$ 可能为 $0$，同除使解集由 $\{0,a\}$ 缩小为 $\{a\}$，属于“范围缩小易漏根”。
\end{solution}

\key{风险三：对数真数限制定义域子集，易增根。}
\[
  \log_a f(x)=\log_a g(x)\iff f(x)=g(x)>0,
\]
其中“$>0$”来自真数的定义域限制。定义域
$D=\{x\mid f(x)>0,\ g(x)>0\}$ 是一个\key{子集约束}，解必须落在 $D$ 内；忽略它就增根。
（同理，$a^{\log_a x}=x$ 要求 $x>0$，而 $\log_a a^x=x$ 对一切实数 $x$ 成立。）

\begin{example}
解方程 $\log_2(x^2-4)=\log_2(3x)$。
\end{example}
\begin{solution}
由对数定义，真数必须为正：$x^2-4>0$ 且 $3x>0$，即 $x>2$。
在此定义域内 $\log_2$ 是单射，故 $x^2-4=3x$，即 $x^2-3x-4=0$，解得 $x=4$ 或 $x=-1$。
$x=-1$ 不在定义域 $x>2$ 内，是增根，舍去。故解集为 $\{4\}$。
\end{solution}

\begin{conccard}[一级]{集合相等判定公理：$A=B\iff(A\subseteq B)\land(B\subseteq A)$}{证明集合相等、审查方程（不等式）同解变形}{只证 $A\subseteq B$ 一个方向就宣布相等；只算“怎么变形”，不算解集往哪变}
集合相等的本质是“互相包含”。证 $A=B$ 必须证两个方向：$A\subseteq B$ 且 $B\subseteq A$；
只证一个方向只能得到包含。把方程（不等式）的解集记为 $A$、变形后解集记为 $B$，
只有 $A=B$ 时变形才同解；$A\subsetneqq B$ 则增根，$B\subsetneqq A$ 则漏根。
这条公理是全部同解变形的“合法性审查标准”。
\end{conccard}

\begin{commonerror}{同解变形不审方向}
常见错误：平方后不验根、同除含参因式不讨论零、解对数方程不写定义域。
它们都源于同一个疏漏——只算“怎么变形”，不算“解集往哪个方向变”。
例如由 $\sqrt{x+1}=x-1$ 平方得 $x^2-3x=0$，解集由 $\{3\}$ 扩大为 $\{0,3\}$，多出的 $0$ 就是增根；
由 $x^2=ax$ 两边除以 $x$ 得 $x=a$，解集由 $\{0,a\}$ 缩小为 $\{a\}$，漏掉了 $0$。
\end{commonerror}

\subsection{命题演算三大变形引擎}

\key{探究二：德摩根律——“且”的否定是“或”}
\[
  \neg(p\land q)\iff\neg p\lor\neg q,\qquad
  \neg(p\lor q)\iff\neg p\land\neg q.
\]
对应的集合版本（补集的分配律）：
\[
  \complement_U(A\cap B)=\complement_U A\cup\complement_U B,\qquad
  \complement_U(A\cup B)=\complement_U A\cap\complement_U B.
\]
记忆：否定“且”要变“或”，否定“或”要变“且”，并且每个命题都要否定。

\key{引擎二：逆否等价——正难则反}
\[
  (p\Rightarrow q)\iff(\neg q\Rightarrow\neg p).
\]
原命题与逆否命题同真同假，这是“正难则反”的合法性来源。它还有集合版本：

\begin{derivation}
要证 $A\subseteq B\iff\complement_U B\subseteq\complement_U A$。
\begin{align*}
  A\subseteq B &\iff \forall x\,(x\in A\Rightarrow x\in B)\\
  &\iff \forall x\,(x\notin B\Rightarrow x\notin A) &&(\text{逆否等价})\\
  &\iff \forall x\,(x\in\complement_U B\Rightarrow x\in\complement_U A)\\
  &\iff \complement_U B\subseteq\complement_U A.
\end{align*}
\end{derivation}

可见：要证“小集合含于大集合”（$A\subseteq B$），可改证“大集合的补集含于小集合的补集”
（$\complement_U B\subseteq\complement_U A$），往往更容易。

\key{引擎三：量词否定与补集转化流程}
\[
  \neg(\forall x,P(x))\iff\exists x,\neg P(x),\qquad
  \neg(\exists x,P(x))\iff\forall x,\neg P(x).
\]
口诀：\key{量词互换，结论否定}。与德摩根律联用：
\[
  \neg\bigl(\forall x,\,p(x)\land q(x)\bigr)\iff\exists x,\,\neg p(x)\lor\neg q(x).
\]
题干关键词的补集转化流程：
\begin{itemize}
  \item “至少有一个”（$\exists$）的否定是“一个都没有”（$\forall,\neg$）；
  \item “至多有一个”= 否定“至少有两个”；
  \item “不都”= $\neg\forall$ = “至少有一个不”；
  \item “无解”= 解集 $\varnothing$ = 补集为全集 $U$。
\end{itemize}
流程：把“至多／至少／不都／无解”改写成“$\exists/\forall$ + 否定” $\to$ 用德摩根与量词否定化简
$\to$ 求补集（\key{补集法}）$\to$ 取补并检查端点。

\key{恒成立／能成立：值域的子集包含与交集非空。}
设 $f:D\to\mathbb{R}$ 的值域为 $R_f=\{f(x)\mid x\in D\}$，目标集合为 $M$，则
\[
  \forall x\in D,\ f(x)\in M \iff R_f\subseteq M,\qquad
  \exists x\in D,\ f(x)\in M \iff R_f\cap M\neq\varnothing.
\]
特例 $M=(a,+\infty)$：恒成立 $\iff\min R_f>a$（看最小值）；能成立 $\iff\max R_f>a$（看最大值）。
例如 $f(x)=x^2-2x$ 在 $[0,3]$ 上的值域为 $[-1,3]$，则 $f(x)>a$ 恒成立 $\iff a<-1$，
能成立 $\iff a<3$。

\begin{example}
写出命题“$\exists x\in\mathbb{R},\ x^2-2x+a\leq0$”的否定；
并说明“方程 $x^2-2x+a=0$ 无实根”用集合语言如何表述。
\end{example}
\begin{solution}
否定为“$\forall x\in\mathbb{R},\ x^2-2x+a>0$”（量词互换，结论否定）。
“无实根”即解集为空集：$\{x\mid x^2-2x+a=0\}=\varnothing$，
等价于“对任意 $x\in\mathbb{R}$，$x^2-2x+a\neq0$”，也等价于 $\Delta=4-4a<0$，即 $a>1$。
\end{solution}

\begin{example}
已知“存在 $x\in[1,2]$，使 $x^2-ax+1<0$”为真，求 $a$ 的取值范围（用补集法）。
\end{example}
\begin{solution}
因为 $x\in[1,2]$ 为正，原命题等价于
\[
  \exists x\in[1,2],\ a>x+\frac1x
  \iff a>\min_{x\in[1,2]}\Bigl(x+\frac1x\Bigr).
\]
由均值不等式 $x+\frac1x\geq2$，当 $x=1$ 时取等号，故最小值为 $2$，于是 $a>2$。

补集法验证：其否定为“$\forall x\in[1,2],\ x^2-ax+1\geq0$”，即 $a\leq x+\frac1x$ 恒成立，
得 $a\leq2$；取补集得 $a>2$，与上面一致。
\end{solution}

\begin{problem}
设 $f(x)=x^2-2x$ 在 $[0,3]$ 上的值域为 $R_f$。分别求“$f(x)>a$ 恒成立”与“$f(x)>a$ 能成立”时 $a$ 的范围。
\end{problem}
\begin{solution}
$R_f=[-1,3]$。恒成立 $\iff R_f\subseteq(a,+\infty)\iff a<-1$；
能成立 $\iff R_f\cap(a,+\infty)\neq\varnothing\iff a<3$。
\end{solution}

\begin{conccard}[二级]{逆否等价与补集反向包含：$p\Rightarrow q\iff\neg q\Rightarrow\neg p$；$A\subseteq B\iff\complement_U B\subseteq\complement_U A$}{证明充分条件、正难则反的题目}{把“逆命题”与“逆否命题”混同；把“否命题”当“命题的否定”}
原命题与逆否命题同真同假，这是“正难则反”的合法性来源。集合语言下，
要证 $A\subseteq B$ 可改证 $\complement_U B\subseteq\complement_U A$——
把“小推大”换成“大补推小补”，往往更容易。
\end{conccard}

\begin{conccard}[二级]{恒成立／能成立：值域的子集包含与交集非空}{含参恒成立、能成立求参数范围}{把“恒成立”与“能成立”方向用反；忽略值域端点是否取到}
设 $f$ 在 $D$ 上的值域为 $R_f$。$\forall x\in D,\ f(x)\in M\iff R_f\subseteq M$（看最小值）；
$\exists x\in D,\ f(x)\in M\iff R_f\cap M\neq\varnothing$（看最大值）。
特例 $M=(a,+\infty)$：恒成立 $\iff\min R_f>a$，能成立 $\iff\max R_f>a$。
\end{conccard}

\begin{commonerror}{把“否命题”当“命题的否定”}
命题“若 $p$ 则 $q$”的\key{否定}是“$p$ 且非 $q$”（$p\land\neg q$），它的\key{否命题}是“若非 $p$ 则非 $q$”（$\neg p\Rightarrow\neg q$）。
两者不同：否定只否定结论，否命题条件与结论都否定。
例如“若 $x>2$ 则 $x^2>4$”的否定是“存在 $x>2$ 使 $x^2\leq4$”；
其否命题“若 $x\leq2$ 则 $x^2\leq4$”是假命题（取 $x=-3$）。
\end{commonerror}

\begin{method}
补集转化法（正难则反）：遇到“至多、至少、不都、无解、并非”等关键词，先写出它的否定
（量词互换、联结词用德摩根），把“存在”转“任意”或把“任意”转“存在”，
求出否定命题成立的参数范围，再取补集。步骤：
\begin{enumerate}
  \item 定位量词（$\forall/\exists$）与联结词（$\land/\lor$）；
  \item 否定量词、否定结论，用德摩根律化简；
  \item 求否定命题对应的集合（补集法求参数）；
  \item 取补集，并检查端点。
\end{enumerate}
\end{method}

\subsection{两大间接证明支柱}

\key{探究三：反证法——证明反例集为空}
要证命题 $p\Rightarrow q$。设 $U$ 为论域，$A=\{x\mid p(x)\}$（条件成立），
$B=\{x\mid q(x)\}$（结论成立），则要证 $A\subseteq B$。
反证法不是直接证 $A\subseteq B$，而是证明\key{反例集为空}：
\[
  A\cap\complement_U B=\varnothing.
\]
做法：假设存在反例 $x_0\in A\cap\complement_U B$（满足条件却不满足结论），推出矛盾，
从而交集为空，即 $A\subseteq B$。

\begin{figure}[htbp]
  \centering
  \begin{tikzpicture}[>=Stealth, semithick, font=\footnotesize]
    \fill[gray!12] (0,0) rectangle (8,5);
    \draw[thick, main!70!black] (0,0) rectangle (8,5);
    \node[main!70!black] at (7.5,4.5) {$U$};
    \fill[main!10, main!70!black] (4.2,2.5) circle (1.8);
    \fill[second!20, second] (3.7,2.5) circle (0.9);
    \node[main!70!black] at (4.9,3.7) {$B$};
    \node[second] at (3.7,2.5) {$A$};
    \node[gray!60!black] at (1.1,2.6) {$\complement_U B$};
    \node[red!70!black] at (1.1,1.6) {$A\cap\complement_U B=\varnothing$};
  \end{tikzpicture}
  \caption{反证法的集合图景：条件集合 $A$ 完全落在结论集合 $B$ 内，反例集 $A\cap\complement_U B$ 为空}
  \label{fig:contradiction}
\end{figure}

反证法的第一步是明确写出“反设”：$A\not\subseteq B$，即存在 $x\in A$ 且 $x\notin B$。
图 \ref{fig:contradiction-assumption} 画出这一反设态；随后由条件与 $B$ 的定义推出矛盾，红色月牙收缩为空，回到图 \ref{fig:contradiction} 的结论态。

\begin{figure}[htbp]
  \centering
  \begin{tikzpicture}[>=Stealth, semithick, font=\footnotesize]
    \fill[gray!12] (0,0) rectangle (8,5);
    \draw[thick, main!70!black] (0,0) rectangle (8,5);
    \node[main!70!black] at (7.5,4.5) {$U$};

    \coordinate (B) at (2.7,2.6);
    \coordinate (A) at (4.9,2.6);

    % B（灰底）与 A；再把 A\B 涂成红色月牙
    \fill[main!14] (B) circle (1.6);
    \begin{scope}
      \clip (A) circle (1.4);
      \fill[red!50] (A) circle (1.4);
      \fill[main!14] (B) circle (1.6);
    \end{scope}

    \draw[main!60!black, thick] (B) circle (1.6);
    \draw[second!70!black, thick] (A) circle (1.4);

    \node[main!70!black] at (1.85,3.9) {$B$};
    \node[second!80!black] at (5.5,3.55) {$A$};
    \node[gray!60!black] at (0.75,4.4) {$\complement_U B$};

    \fill[red!80!black] (5.8,2.6) circle (2.2pt);
    \node[red!80!black, right] at (5.87,2.6) {$x$};
    \node[red!70!black, align=center] at (6.05,1.0) {$A\cap\complement_U B$\\ $\neq\varnothing$};
    \draw[red!60!black, ->, shorten >=1pt] (6.25,1.3) -- (5.9,2.1);

    \node[red!70!black] at (4.0,0.55) {反设：$\exists x\in A$ 且 $x\notin B$};
  \end{tikzpicture}
  \caption{反证法的反设态（正难则反）：假设 $A\not\subseteq B$，集合 $A$ 有一部分越出 $B$，红色月牙区域 $A\cap\complement_U B$ 非空，反例点 $x$ 落在其中}
  \label{fig:contradiction-assumption}
\end{figure}

\key{触发信号词：}
\begin{itemize}
  \item \key{唯一性}（“只有一个”“恰有一个”“唯一”）——假设存在两个不同的对象，推出矛盾（或证明任意两个相等）；
  \item \key{不存在}（“不存在”“没有”“无”）——假设存在，推出矛盾；
  \item \key{无理数}（“$\sqrt2$ 是无理数”）——假设可写成既约分数，推出分子分母同为偶数，与既约矛盾；
  \item \key{无限性}（“有无穷多个”）——假设只有有限个，推出还能构造出下一个，矛盾。
\end{itemize}

\begin{example}
用反证法证明：$\sqrt2$ 不是有理数。
\end{example}
\begin{solution}
假设 $\sqrt2$ 是有理数，则可写成既约分数 $\sqrt2=\dfrac pq$
（$p,q\in\mathbb{Z}^*$，$q>0$，$\gcd(p,q)=1$）。
两边平方得 $p^2=2q^2$，故 $p^2$ 为偶数，从而 $p$ 为偶数（奇数的平方仍为奇数）。
设 $p=2k$，代入得 $4k^2=2q^2$，即 $q^2=2k^2$，故 $q$ 也为偶数。
于是 $p,q$ 均为偶数，与 $\gcd(p,q)=1$ 矛盾。故假设不成立，$\sqrt2$ 是无理数。
\end{solution}

\key{探究四：数学归纳法——皮亚诺公理与最小反例原理}
数学归纳法不是“递推直觉”，而是自然数结构的公理推论。皮亚诺公理（$\mathbb{N}^*$）要点：
\begin{enumerate}
  \item $1\in\mathbb{N}^*$；
  \item 每个 $n\in\mathbb{N}^*$ 有唯一后继 $n+1$；
  \item $1$ 不是任何数的后继；
  \item 后继函数是单射：$m+1=n+1\Rightarrow m=n$；
  \item \key{归纳公理}：若 $S\subseteq\mathbb{N}^*$ 满足 $1\in S$ 且 $n\in S\Rightarrow n+1\in S$，则 $S=\mathbb{N}^*$。
\end{enumerate}
与归纳公理等价的是\key{良序原理}：$\mathbb{N}^*$ 的任一非空子集必有最小元。
由此得\key{最小反例原理}：若命题 $P(n)$ 不恒成立，则存在最小反例 $m$。
由 $m$ 最小知 $P(m-1)$ 成立，再用归纳递推得 $P(m)$ 成立，矛盾；故 $P(n)$ 恒成立。

因此数学归纳法两步缺一不可：
\begin{itemize}
  \item \key{归纳基础}：验证 $P(n_0)$ 成立（第一块骨牌倒下）；
  \item \key{归纳递推}：假设 $P(k)$ 成立，证 $P(k+1)$ 成立（前一块倒下必带倒后一块）。
\end{itemize}
只有基础是“只倒第一块”；只有递推是“没有启动”。归纳起点 $n_0$ 可以平移（不一定从 $1$ 开始）。

\begin{figure}[htbp]
  \centering
  \begin{tikzpicture}[>=Stealth, semithick, font=\footnotesize]
    \draw[->, gray] (0,0)--(10.8,0) node[right]{$n\in\mathbb{N}^*$};
    \foreach \i/\lab in {0/1,1/2,2/3,3/4,4/5,5/6} {
      \pgfmathsetmacro{\xx}{0.8+1.5*\i}
      \draw[main!70!black, fill=main!12, rounded corners=0.8pt] (\xx,0) rectangle ++(0.5,2.0);
      \node at (\xx+0.25,1.0) {$\lab$};
    }
    \foreach \i in {0,1,2,3,4} {
      \pgfmathsetmacro{\xx}{1.3+1.5*\i}
      \draw[->, second, thick] (\xx,1.0)--++(0.55,0);
    }
    \node[below, main] at (0.8,-0.15) {推倒第 $1$ 块};
    \node[above, second] at (5.0,2.3) {若第 $k$ 块倒下，则第 $k+1$ 块必倒下};
    \node[below, gray!70!black] at (8.2,-0.15) {$\cdots$};
    \node[below, gray!70!black] at (9.8,-0.15) {全部倒下};
  \end{tikzpicture}
  \caption{数学归纳法的多米诺模型：归纳基础推倒第 1 块，归纳递推保证“倒一块必带倒下一块”，于是全部倒下}
  \label{fig:induction}
\end{figure}

\begin{example}
用数学归纳法证明：$7^n-1$ 能被 $6$ 整除（$n\in\mathbb{N}^*$）。
\end{example}
\begin{solution}
(1) 基础：$n=1$ 时 $7^1-1=6$，能被 $6$ 整除。
(2) 递推：假设 $n=k$ 时 $7^k-1$ 能被 $6$ 整除，即 $7^k-1=6m$。则 $n=k+1$ 时
\[
  7^{k+1}-1=7\cdot7^k-1=7(7^k-1)+6=7\cdot6m+6=6(7m+1),
\]
能被 $6$ 整除。由(1)(2)，命题对一切 $n\in\mathbb{N}^*$ 成立。
\end{solution}

\subsection{武器库索引表与跨模块实战}

\key{《高中数学逻辑变形与证明触发信号对照表》}
把题干信号、集合／逻辑翻译、变形引擎与操作风险放在同一张表，解题时“对号入座”。

\begin{center}
\footnotesize
\setlength{\tabcolsep}{3pt}
\renewcommand{\arraystretch}{1.4}
\begin{tabular}{p{2.0cm}p{3.2cm}p{2.6cm}p{3.6cm}}
\toprule
题干信号 & 逻辑／集合翻译 & 变形引擎 & 操作与风险 \\
\midrule
集合相等 $A=B$ & 双向包含 & 集合相等公理 & 分证 $A\subseteq B$ 与 $B\subseteq A$ \\
\hline
平方、去绝对值 & $f^2=g^2\iff f=\pm g$ & 同解变形 & 解集扩大，必须验根 \\
\hline
同除含参因式 & 可能丢根 & 移项分解＋讨论 & 解集缩小，防漏根 \\
\hline
对数方程 & 真数 $>0$ & 同解变形＋定义域 & 忽略定义域则增根 \\
\hline
至少有一个 & $\exists x$ & 量词否定 & 否定为 $\forall x,\ \neg$ \\
\hline
至多有一个 & 否定“至少两个” & 量词否定＋德摩根 & 命题的否定改写 \\
\hline
不都 & $\neg\forall$ & 量词否定 & 否定为“至少有一个不” \\
\hline
无解、为空集 & $S=\varnothing$ & 补集 & 补集为全集 $U$ \\
\hline
恒成立 & $\forall x,\ f(x)\in M$ & 值域子集 & $R_f\subseteq M$，看最小值 \\
\hline
能成立 & $\exists x,\ f(x)\in M$ & 值域交集 & $R_f\cap M\neq\varnothing$，看最大值 \\
\hline
唯一性 & 存在且唯一 & 反证法 & 假设第二个，推出矛盾 \\
\hline
不存在、没有 & $\neg\exists$ & 反证法 & 假设存在，推出矛盾 \\
\hline
无理数 & 非有理数 & 反证法 & 设 $\sqrt2=\frac pq$ 既约 \\
\hline
无限性 & 非有限 & 反证法 & 假设有限，构造下一个 \\
\hline
对所有 $n\in\mathbb{N}^*$ & 自然数命题 & 数学归纳法 & 基础＋递推，缺一不可 \\
\hline
整除、离散不等式 & 自然数命题 & 数学归纳法／最小反例 & 归纳起点可平移至 $n_0$ \\
\bottomrule
\end{tabular}
\end{center}

\begin{example}
解不等式 $\dfrac{x^2-4}{x-2}\geq0$，并说明直接约去 $(x-2)$ 会丢掉什么。
\end{example}
\begin{solution}
定义域要求 $x-2\neq0$。在定义域内
\[
  \frac{x^2-4}{x-2}=\frac{(x-2)(x+2)}{x-2}=x+2,
\]
故原不等式等价于 $x+2\geq0$ 且 $x\neq2$，即 $x\geq-2$ 且 $x\neq2$，
解集为 $[-2,2)\cup(2,+\infty)$。
若直接约去 $(x-2)$ 得 $x+2\geq0$，解集为 $[-2,+\infty)$，把分母不能为 $0$ 的限制丢了（漏掉限制，不是漏根，而是扩大了解集），
所以必须保留 $x\neq2$。
\end{solution}

\begin{example}
若“存在 $x\in\mathbb{R}$，使 $x^2+2ax+4<0$”是假命题，求 $a$ 的取值范围。
\end{example}
\begin{solution}
该命题为假 $\iff$ 其否定“$\forall x\in\mathbb{R},\ x^2+2ax+4\geq0$”为真。
这是恒成立问题：$\Delta=(2a)^2-16\leq0\iff4a^2-16\leq0\iff-2\leq a\leq2$。
故 $a\in[-2,2]$。
\end{solution}

\begin{example}
证明：函数 $f(x)=x^3+x+1$ 有且只有一个零点。
\end{example}
\begin{solution}
\key{存在性}：$f(-1)=-1-1+1=-1<0$，$f(0)=1>0$，且 $f$ 是多项式（连续），
由零点存在定理，$(-1,0)$ 内至少有一个零点。
\key{唯一性（反证）}：假设有两个不相等的零点 $x_1,x_2$（设 $x_1<x_2$），则
\[
  x_1^3+x_1+1=0,\qquad x_2^3+x_2+1=0.
\]
两式相减得 $(x_1^3-x_2^3)+(x_1-x_2)=0$，即
\[
  (x_1-x_2)(x_1^2+x_1x_2+x_2^2+1)=0.
\]
因 $x_1\neq x_2$，需 $x_1^2+x_1x_2+x_2^2+1=0$。
但 $x_1^2+x_1x_2+x_2^2=\left(x_1+\frac{x_2}{2}\right)^2+\frac34x_2^2\geq0$，
故 $x_1^2+x_1x_2+x_2^2+1\geq1>0$，矛盾。故零点唯一。
\end{solution}

\begin{example}
用数学归纳法证明：$n^3+5n$ 能被 $6$ 整除（$n\in\mathbb{N}^*$）。
\end{example}
\begin{solution}
(1) 基础：$n=1$ 时 $1+5=6$，能被 $6$ 整除。
(2) 递推：假设 $n=k$ 时 $k^3+5k=6m$。则 $n=k+1$ 时
\begin{align*}
  (k+1)^3+5(k+1)&=k^3+3k^2+3k+1+5k+5\\
  &=(k^3+5k)+3k^2+3k+6\\
  &=6m+3k(k+1)+6.
\end{align*}
因为 $k(k+1)$ 是两个连续整数之积，必为偶数，故 $3k(k+1)$ 能被 $6$ 整除；又 $6$ 能被 $6$ 整除。
所以 $(k+1)^3+5(k+1)$ 能被 $6$ 整除。由(1)(2)，结论对一切 $n\in\mathbb{N}^*$ 成立。
\end{solution}

\begin{example}
用数学归纳法证明：对一切 $n\geq5$，$2^n>n^2$。
\end{example}
\begin{solution}
(1) 基础：$n=5$ 时 $2^5=32>25=5^2$，成立。
(2) 递推：假设 $n=k$（$k\geq5$）时 $2^k>k^2$。则 $n=k+1$ 时
\[
  2^{k+1}=2\cdot2^k>2k^2.
\]
要证 $2k^2>(k+1)^2=k^2+2k+1$，即 $k^2-2k-1>0$，即 $(k-1)^2>2$。
因 $k\geq5$，$(k-1)^2\geq16>2$，成立。故 $2^{k+1}>(k+1)^2$。
由数学归纳法，对一切 $n\geq5$，$2^n>n^2$（注意归纳起点平移到了 $5$）。
\end{solution}

\begin{example}
设命题 $p$：$\forall x\in\mathbb{R},\ x^2+ax+1>0$；命题 $q$：$\exists x\in\mathbb{R},\ x^2-2x+a=0$。
若 $p\lor q$ 为假，求 $a$ 的取值范围。
\end{example}
\begin{solution}
$p\lor q$ 为假 $\iff$ $p$ 假且 $q$ 假（德摩根律）。
$p$ 真 $\iff\Delta_1=a^2-4<0\iff-2<a<2$，故 $p$ 假 $\iff a\leq-2$ 或 $a\geq2$。
$q$ 真 $\iff\Delta_2=4-4a\geq0\iff a\leq1$，故 $q$ 假 $\iff a>1$。
取交：$(a\leq-2\ \text{或}\ a\geq2)\cap(a>1)$，得 $a\geq2$。
故 $a\in[2,+\infty)$。
\end{solution}

\begin{conccard}[三级]{反证法触发信号：唯一性、不存在、无理数、无限性}{结论含“只有一个／不存在／无理数／无穷多”的证明题}{把反证法当成“随便假设”；假设后没有推出矛盾就下结论}
反证法即证反例集 $A\cap\complement_U B=\varnothing$：假设反例存在，推出与公理、定理或临时假设矛盾的结果。
唯一性假设“至少两个”；不存在假设“存在”；无理数设既约分数 $\frac pq$；无限性设“只有有限个”。
矛盾必须明确写出，否则证明不完整。
\end{conccard}

\begin{conccard}[三级]{数学归纳法：归纳基础＋归纳递推}{证明对全体自然数（或 $n\geq n_0$）成立的命题}{只做递推不做基础；归纳假设没有真正用上；归纳起点误写为 $1$}
归纳公理：若 $1\in S$ 且 $n\in S\Rightarrow n+1\in S$，则 $S=\mathbb{N}^*$。
其根基是良序原理与最小反例原理：若存在反例，则存在最小反例，与递推矛盾。
使用时要检查两点：归纳基础是否验证；递推中是否用到了归纳假设。起点可按命题需要平移到 $n_0$。
\end{conccard}

\begin{summary}
本章把集合与逻辑锻造成一套“武器库”：集合相等公理 $A=B\iff(A\subseteq B)\land(B\subseteq A)$ 是审查一切同解变形的标准（$A\subsetneqq B$ 增根、$B\subsetneqq A$ 漏根）；德摩根律、逆否等价、量词否定与补集转化处理“至多、至少、不都、无解、恒成立、能成立”；反证法即证反例集 $A\cap\complement_U B=\varnothing$，由“唯一性、不存在、无理数、无限性”四类信号词触发；数学归纳法以皮亚诺公理与良序原理为根基，用最小反例原理保证“归纳基础＋归纳递推”两步缺一不可。解题时先翻译题干（集合／逻辑语言），再选引擎，最后审方向（防增漏根）。
\end{summary}

\end{explore}

\subsection{四视角等价地图：同一对象，多副面孔}

前四个模块其实都在做同一件事：给同一个数学对象换一副面孔。
代数、几何、集合、映射（对象本身）不是四门课，而是同一结构的四种\key{表示}；
只要两种表示之间存在\emph{保持结构的对应}（同构），命题就能互相翻译。
但「等价」不等于「一样难」：\textbf{同一问题在不同框架下的可解释性与计算量差别极大}，
所以解题的关键动作是\key{找一个同构映射，把问题搬到更顺的那一侧}。

\begin{center}
\footnotesize
\begin{tabular}{c|p{2.6cm}|p{3.2cm}|p{3.0cm}|p{3.4cm}}
\toprule
视角 & 看什么 & 擅长 & 短板 & 典型「搬运」 \\
\midrule
代数 & 符号、方程、恒等变形 & 机械化、可执行、可检验 & 缺直观，易迷失方向 & 几何 $\to$ 方程（解析几何）\\
几何 & 图形、变换、位置关系 & 直观、好找思路、便于猜答案 & 精度与严格性弱 & 方程 $\to$ 图形（数形结合）\\
集合 & 元素、包含、交并补、量词 & 严谨分类、盯住边界与反例 & 不给算法 & 命题 $\to$ 集合关系（充分必要）\\
映射／对象 & 对应关系、迭代、结构本身 & 统一语言、跨领域搬运结论 & 太抽象，需落地实例 & 矩阵 $=$ 线性映射；坐标 $=$ 同构 \\
\bottomrule
\end{tabular}
\end{center}

\begin{figure}[htbp]
\centering
\begin{tikzpicture}[>=Stealth,semithick,font=\footnotesize,
  box/.style={draw,rounded corners,inner sep=4pt,align=center,minimum width=2.1cm}]
  \node[box,draw=main] (A) at (-2.2,1.9) {代数\\方程};
  \node[box,draw=second] (G) at (2.2,1.9) {几何\\图形};
  \node[box,draw=third] (S) at (-2.2,-1.9) {集合\\逻辑};
  \node[box,draw=gray!60!black] (M) at (2.2,-1.9) {映射\\结构};
  \node[box,draw=main!70!black,minimum width=2.4cm] (C) at (0,0) {\textbf{同一个}\\数学对象};
  \draw[<->,main] (C)--(A) node[midway,above left,font=\scriptsize] {列方程};
  \draw[<->,second] (C)--(G) node[midway,above right,font=\scriptsize] {画出来};
  \draw[<->,third] (C)--(S) node[midway,below left,font=\scriptsize] {下定义};
  \draw[<->,gray!60!black] (C)--(M) node[midway,below right,font=\scriptsize] {建模型};
  \draw[->,gray!55,dashed] (A)--(G) node[midway,above,font=\scriptsize] {数形结合};
  \draw[->,gray!55,dashed] (S)--(M) node[midway,below,font=\scriptsize] {关系即映射};
\end{tikzpicture}
\caption{四视角等价地图：中心是对象本身，四个顶点是它的四种表示；
箭头都是「保持结构的映射」。解题时沿箭头换到顺手的一侧，算完再换回来}
\label{fig:four-view-map}
\end{figure}

\begin{figure}[htbp]
\centering
\begin{tikzpicture}[>=Stealth,semithick,font=\footnotesize,
  ch/.style={draw,rounded corners,inner sep=5pt,align=center,minimum width=2.0cm},
  lanename/.style={rotate=90,anchor=south,font=\scriptsize\bfseries,text=main!70!black}]
  % 三条主线 + 一条独立线
  \node[lanename] at (-1.1,6) {几何·结构线};
  \node[lanename] at (-1.1,3) {离散动力线};
  \node[lanename] at (-1.1,0.6) {连续分析线};
  % 第 0 层：语言底座
  \node[ch,draw=main,fill=main!8] (c0) at (1.6,7.6) {第 0 章\\集合·逻辑·证明};
  \node[ch,draw=gray!60!black,fill=gray!8] (c7) at (11.4,7.6) {第 7 章\\概率·统计};
  \draw[->,main] (c0)--(c7) node[midway,above,font=\scriptsize] {测度＝归一化计数};
  % 几何·结构线
  \node[ch,draw=second] (c1) at (1.6,6) {第 1 章\\函数};
  \node[ch,draw=second] (c2) at (5.0,6) {第 2 章\\三角·复数};
  \node[ch,draw=second] (c3) at (8.2,6) {第 3 章\\向量};
  \node[ch,draw=second] (c4) at (11.4,6) {第 4 章\\立几·圆锥曲线};
  \draw[->,second] (c1)--(c2);
  \draw[->,second] (c2)--(c3);
  \draw[->,second] (c3)--(c4);
  % 离散动力线
  \node[ch,draw=third] (c5) at (1.6,3) {第 5 章\\数列·递推};
  \node[ch,draw=third] (c9) at (11.4,3) {第 9 章\\不动点};
  \draw[->,third] (c5)--(c9) node[midway,above,font=\scriptsize] {递推求通项·收敛};
  % 连续分析线
  \node[ch,draw=main!60!black] (c6) at (5.0,0.6) {第 6 章\\导数·局部线性化};
  \node[ch,draw=main!60!black] (c8) at (8.2,0.6) {第 8 章\\齐次化};
  \draw[->,main!60!black] (c6)--(c9) node[midway,right,font=\scriptsize] {稳定性};
  \draw[->,main!60!black] (c8)--(c9) node[midway,right,font=\scriptsize] {射影};
  % 跨线接口
  \draw[->,gray!55,dashed] (c1)--(c5) node[midway,left,font=\scriptsize] {离散函数};
  \draw[->,gray!55,dashed] (c1)--(c6) node[midway,left,font=\scriptsize] {变化率};
  \draw[<->,gray!55,dashed] (c5)--(c6) node[midway,below,font=\scriptsize] {差分 $\leftrightarrow$ 导数};
  \draw[->,gray!55,dashed] (c4)--(c8) node[midway,right,font=\scriptsize] {条件最值};
  \draw[->,main!35] (c0)--(c1);
  \draw[->,gray!55,dashed] (c3)--(c9) node[midway,above right,font=\scriptsize] {特征方向};
\end{tikzpicture}
\caption{全书知识脉络图：实线是各自主线的推进，虚线是跨主线接口。
三条主线（几何·结构／离散动力／连续分析）最终在第 9 章不动点汇合；
第 8 章齐次化是射影观点的入口，第 0 章语言底座供全局调用}
\label{fig:book-map}
\end{figure}

\textbf{脉络怎么读}：几何·结构线从函数到圆锥曲线，解决「形的问题」；
离散动力线与连续分析线在 \key{差分 $\leftrightarrow$ 导数}处对接（同一个
「以直代曲／以比代差」思想），分别解决「列的问题」与「极限的问题」，
并共同汇入第 9 章的不动点理论；第 8 章齐次化把二元问题压到一元比例，
是通往射影／映射观点的接口。每一章章首的「多框架对读」卡，就是在这张
图上标明「你现在在哪条线、哪一侧」。

\textbf{怎么用它找思路}：拿到题先问「我现在站在哪个视角」——
卡在代数就画图（第 2、4 章的向量与圆锥曲线）；卡在几何就设坐标列方程
（第 4 章的仿射、第 8 章的齐次化）；被式子的次数与比值绕晕就升到射影／映射
观点（第 9 章的不动点与 Möbius）；讨论「有没有、有几个」就回到集合与量词
（本章的同解变形与反证法）。\textbf{同一结论往往能在四处各证一遍}，选最省力的那处书写。
此后各章章首统一附「多框架对读」卡，把这条暗线显性化。

\begin{inclasstest}
\begin{enumerate}
  \item 写出“$\exists x\in\mathbb{R},\ x^2-2x+1\leq0$”的否定，并判断其真假；\level{易}
  \item 解方程 $\sqrt{x+3}=x+1$，并说明验根的必要性；\level{中}
  \item 已知“存在 $x\in[1,3]$，使 $x^2-ax+2\leq0$”为真，求 $a$ 的取值范围；\level{中}
  \item 用反证法证明：若 $a,b\in\mathbb{Z}$ 且 $ab$ 能被 $3$ 整除，则 $a,b$ 中至少有一个能被 $3$ 整除。\level{中}
\end{enumerate}
\end{inclasstest}

\begin{afterclass}
\subsection*{基础巩固}

\begin{enumerate}
  \item 写出德摩根律与量词否定的四条公式：$\neg(p\land q)$、$\neg(p\lor q)$、$\neg(\forall x,P(x))$、$\neg(\exists x,P(x))$；
  \item 解方程 $x^2=2x$，并说明若两边同除以 $x$ 会漏掉哪个根。
\end{enumerate}

\subsection*{进阶提升}
\begin{enumerate}
  \item 已知“$\forall x\in\mathbb{R},\ x^2+2x+a>0$”为真，求 $a$；再写出它的否定命题为真时 $a$ 的范围；
  \item 用数学归纳法证明：$5^n-1$ 能被 $4$ 整除。
\end{enumerate}

\subsection*{拓展探究}
\begin{enumerate}
  \item 用反证法证明：$\sqrt3$ 是无理数。（提示：用到上面第 4 题的结论“若 $3\mid ab$，则 $3\mid a$ 或 $3\mid b$”。）
\end{enumerate}

\subsection*{教师版答案}
\begin{solution}
\textbf{基础 1：}$\neg(p\land q)\iff\neg p\lor\neg q$；$\neg(p\lor q)\iff\neg p\land\neg q$；
$\neg(\forall x,P(x))\iff\exists x,\neg P(x)$；$\neg(\exists x,P(x))\iff\forall x,\neg P(x)$。
\end{solution}
\begin{solution}
\textbf{基础 2：}移项分解 $x(x-2)=0$，得 $x=0$ 或 $x=2$，解集 $\{0,2\}$。
若两边同除以 $x$ 得 $x=2$，漏掉 $x=0$（同除使解集由 $\{0,2\}$ 缩小为 $\{2\}$）。
\end{solution}
\begin{solution}
\textbf{进阶 1：}$x^2+2x+a>0$ 恒成立 $\iff\Delta=4-4a<0\iff a>1$。
其否定“$\exists x\in\mathbb{R},\ x^2+2x+a\leq0$”为真 $\iff\Delta\geq0\iff a\leq1$。
\end{solution}
\begin{solution}
\textbf{进阶 2：}(1) $n=1$ 时 $5-1=4$，能被 $4$ 整除。
(2) 假设 $5^k-1=4m$，则 $5^{k+1}-1=5(5^k-1)+4=20m+4=4(5m+1)$，能被 $4$ 整除。
由数学归纳法，结论成立。
\end{solution}
\begin{solution}
\textbf{拓展：}假设 $\sqrt3$ 是有理数，写成既约分数 $\sqrt3=\frac pq$（$\gcd(p,q)=1$）。
平方得 $p^2=3q^2$，故 $3\mid p^2$；由提示（取 $a=b=p$）得 $3\mid p$，设 $p=3k$。
代入得 $9k^2=3q^2$，即 $q^2=3k^2$，故 $3\mid q$。于是 $p,q$ 均为 $3$ 的倍数，与既约矛盾。
故 $\sqrt3$ 是无理数。
\end{solution}
\reading{拓展阅读 00}{逻辑与证明的武器库}
\end{afterclass}

% 高观点导学案：函数的概念与基本初等函数
% 主线：关系 R⊆A×B（三档舞台）→ 方程/不等式＝关系的条件表达 → 两条限制轴（单值性→映射/函数；连续性与结构保持→基本初等函数）
% 全程三视角互译：代数有序数对集合 ↔ 几何坐标系点集 ↔ 映射箭头图

\section{函数的概念与基本初等函数（高观点结构版）}

\begin{methodbox}{多框架对读}
\small
\begin{center}
\begin{tabular}{c|p{2.9cm}|p{5.0cm}|p{3.1cm}}
\toprule
视角 & 本章看点 & 别的视角怎么说 & 哪侧最省力 \\
\midrule
代数 & 解析式、复合与解方程 & 有序数对集合的坐标运算 & 变形与恒等式 \\
几何 & 图象的平移／对称变换 & 坐标系中点集的移动 & 猜性质、定范围 \\
集合 & 定义域、值域、$f:A\to B$ & 取值范围即约束集合 & 审边界、找反例 \\
映射 & $y=f(x)$ 即对应法则 & 变换＝作用在图象上的映射 & 统一语言、整体搬迁 \\
\bottomrule
\end{tabular}
\end{center}
卡的使用：卡住时先判断自己站在哪一栏，再沿箭头换栏（总图见第 0 章图~\ref{fig:four-view-map}）。
\end{methodbox}

% ---------- 前置引入：舞台的诞生 ----------
\subsection{前置引入：舞台的诞生——集合自身的四级进化}

研究任何数学对象，现代数学靠两根支柱：\key{舞台}（集合自身的结构）与 \key{规则}（集合之间的映射）。本导学案主线一直在讲\key{规则}：$A\times B$ 中怎样挑点、怎样过滤出单值对应。但高中函数之所以拥有单调性、奇偶性、零点存在定理与导数，原因不在规则里，而在\key{舞台}上——定义域与值域选中的是哪种集合。

\begin{conccard}[一级]{集合装了多少「装备」，函数就有多少性质}{适用于所有「新定义」与概念辨析题}{只背性质条款，不问性质从何而来}
一般集合（例如 $\{\text{张三},\text{李四}\}$）只是一堆孤立元素：不能比大小，不能做加减，也没有「挨得很近」的概念。一旦把定义域与值域锁定在实数集 $\mathbb{R}$ 上，$\mathbb{R}$ 自带的三层结构便一次性注入函数。
\end{conccard}

\begin{itemize}
  \item \level{Lv.1} \key{原始集合}——一袋散装弹珠。只有「属于」与「元素个数」，无法比大小、不能运算：既谈不上单调性，也画不出坐标系中的图象。
  \item \level{Lv.2} \key{有序与代数集合}（$\mathbb{Z},\mathbb{Q}$）——排好队且能算账的刻度尺。\key{全序结构}（$x_1<x_2$ 可比大小）解锁\key{单调性与最值}；\key{代数运算与相反数}（$+,-,\times,\div$ 与 $-x$）解锁\key{奇偶性与周期性}。复数域 $\mathbb{C}$ 不能比大小，所以复变函数根本没有「单调递增」可言。
  \item \level{Lv.3} \key{稠密但有孔的集合}（$\mathbb{Q}$）——布满微观针眼的筛子。任意两点之间仍有无穷多个点，却漏掉了 $\sqrt{2},\pi,e$ 等无数无理数。致命后果见图~\ref{fig:qline}：$f(x)=x^2-2$ 从负值升到正值时，正好从 $\sqrt{2}$ 的孔洞中漏过去，\key{根本没有零点}。
  \item \level{Lv.4} \key{完备连续集合}（$\mathbb{R}$）——无缝焊接的实心金条。无理数填满全部孔洞，曲线无处可漏，\key{零点存在定理}（介值性）、极限以及「以直代曲」的导数才真正成为可能。
\end{itemize}

Lv.1 与 Lv.2 的差别不在元素本身，而在元素之间被装上了什么关系：前者只有「属于」与计数（图 \ref{fig:bag2axis} 左），后者多出了「谁大谁小」的全序与「怎么算账」的代数运算（图 \ref{fig:bag2axis} 右）。

\begin{figure}[htbp]
  \centering
  \begin{tikzpicture}[>=stealth,font=\small]
    % ---------- 左：Lv.1 一袋散装弹珠（无序点集） ----------
    \begin{scope}
      \draw[gray!60,very thick,fill=main!8] (0,0) ellipse (2.25cm and 1.6cm);
      \foreach \px/\py/\lab in {
        -1.45/0.62/-2, -0.95/-0.88/-1, -0.42/0.28/-\tfrac12,
        0.02/-0.55/0, 0.42/0.78/\tfrac12, 0.88/-0.18/1,
        1.35/0.5/\tfrac32, 1.12/-1.0/2, -1.78/-0.32/\tfrac52}
      {
        \fill[second] (\px,\py) circle (0.075);
        \node[gray!80,font=\scriptsize] at (\px+0.16,\py+0.16) {$\lab$};
      }
      \node at (0,-1.95) {\small 只有「属于」与元素个数};
      \node[second] at (0,-2.42) {\small 不能比大小，也不能加减};
    \end{scope}
    % ---------- 右：Lv.2 带刻度的数轴（有序集合 Q） ----------
    \begin{scope}[xshift=7.2cm]
      \draw[->,gray,very thick] (-2.7,0)--(3.1,0) node[right,gray] {$x$};
      \foreach \t in {-2,-1,0,1,2}
        \draw[gray,thick] (\t,0.09)--(\t,-0.09);
      \foreach \t in {-2,-1,0,1,2}
        \node[gray,font=\scriptsize] at (\t,-0.34) {$\t$};
      \foreach \t in {-2,-1,-0.5,0,0.5,1,1.5,2,2.5}
        \fill[main] (\t,0) circle (0.05);
      \node[third] at (0,1.15) {$(<)\quad +,\,-,\,\times,\,\div$};
      \node at (0,-1.35) {\small 排好队且能算账的刻度尺};
      \node[main] at (0,-1.82) {\small 全序解锁单调最值，运算解锁奇偶周期};
    \end{scope}
  \end{tikzpicture}
  \caption{从「属于」到「可比大小」：左为 Lv.1 原始集合——一袋散装弹珠，只有属于关系与元素个数，无法比大小、不能运算；右为 Lv.2 有序与代数集合——同一批元素排上带刻度的数轴 $\mathbb{Q}$，全序 $(<)$ 解锁单调性与最值，$+,-,\times,\div$ 与相反数解锁奇偶性与周期性}
  \label{fig:bag2axis}
\end{figure}

\begin{figure}[htbp]
  \centering
  \begin{tikzpicture}[x=0.92cm,y=0.92cm,>=stealth,font=\small]
    % ---------- 左：有理数轴（穿孔） ----------
    \begin{scope}
      \draw[->,white!55] (-0.7,0)--(2.9,0) node[right,white!70] {$x$};
      \foreach \x in {0.1,0.35,0.62,0.9,1.15,1.62,1.85,2.1,2.4,2.65}
        \draw[second,fill=white] (\x,0) circle (0.042);
      \draw[second,very thick,fill=white] (1.414,0) circle (0.075);
      \node[second] at (1.414,-0.42) {\small $\sqrt{2}$（洞）};
      \draw[main,thick,domain=0.55:2.0,samples=70,smooth] plot(\x,{\x*\x-2});
      \node[main] at (1.45,2.30) {$y=x^2-2$};
      \node[white!70] at (1.05,-2.02) {\small $\mathbb{Q}$：看似连续，实则满布针眼};
      \node[second] at (1.05,-2.56) {\small 穿过空洞：无交点};
    \end{scope}
    % ---------- 右：实数轴（实心） ----------
    \begin{scope}[xshift=7.6cm]
      \draw[->,third,very thick] (-0.7,0)--(2.9,0) node[right,third] {$x$};
      \foreach \x in {0.1,0.35,0.62,0.9,1.15,1.62,1.85,2.1,2.4,2.65}
        \fill[third!70] (\x,0) circle (0.032);
      \draw[main,thick,domain=0.55:2.0,samples=70,smooth] plot(\x,{\x*\x-2});
      \fill[second] (1.414,0) circle (0.075);
      \node[second] at (1.414,-0.42) {$(\sqrt{2},0)$};
      \node[third] at (1.05,-2.02) {\small $\mathbb{R}$：无缝连续};
      \node[second] at (1.05,-2.56) {\small 洞被填满：零点被迫现身};
    \end{scope}
  \end{tikzpicture}
  \caption{同一函数 $f(x)=x^2-2$ 在两种「舞台」上的命运：左，定义域取有理数集 $\mathbb{Q}$，曲线从 $\sqrt{2}$ 的空心孔洞中漏过，变号却无零点；右，定义域取实数集 $\mathbb{R}$，孔洞被无理数填满，零点 $(\sqrt{2},0)$ 必然出现}
  \label{fig:qline}
\end{figure}

\subsubsection*{开场三问}
\begin{enumerate}
  \item 复数域 $\mathbb{C}$ 上能否定义「单调递增」？为什么？（提示：$\mathbb{C}$ 上有没有「小于」关系）\blank
  \item 在有理数集 $\mathbb{Q}$ 上考察 $f(x)=x^2-3$：它满足 $f(1)<0<f(2)$，却没有零点。请用「孔洞」的语言解释。\blank
  \item 狄利克雷函数 $D(x)=\begin{cases}1,&x\in\mathbb{Q}\\0,&x\in\mathbb{R}\setminus\mathbb{Q}\end{cases}$ 以任一非零有理数为周期，却没有最小正周期。它击穿了「周期函数一定有最小正周期」这句直觉——你能再举一句被它击穿的说法吗？\blank
\end{enumerate}

\begin{method}
「看集合选武器」的反射：拿到一道涉及函数的题，先扫一眼定义域是哪类集合，再决定调用哪套工具——
$\mathbb{N}^{*}$（离散有序）$\Rightarrow$ 差分、数学归纳法、不动点；
$[a,b]\subseteq\mathbb{R}$（连续完备）$\Rightarrow$ 单调性、零点定理、导数局部线性化；
$\mathbb{Q}$ 与 $\mathbb{R}\setminus\mathbb{Q}$ 混用 $\Rightarrow$ 立刻警觉稠密性与孔洞性，准备构造反例。
\end{method}

\begin{commonerror}{把「有界」当成「能取到最值」}
$f(x)>a$ 在 $(0,1)$ 上恒成立，即使 $f$ 在该区间单调递增且 $\lim\limits_{x\to0^{+}}f(x)=2$，也只能推出 $a\le 2$ 而不是 $a<2$：端点取不到时，边界值扮演的是\key{下确界}而非最小值。这类「端点开合」丢分，考的正是实数的完备性。
\end{commonerror}

四级进化完成后，$\mathbb{R}$ 与自身作笛卡尔积，舞台张成连续平面（图 \ref{fig:r2plane}）——后续的二维区域与函数图象都落在这块舞台上。

\begin{figure}[htbp]
  \centering
  \begin{tikzpicture}[>=stealth,font=\small,scale=1.02]
    \draw[third!35,dashed,thin] (-2.5,-2) grid (2.5,2);
    \draw[->,gray,very thick] (-2.75,0)--(2.75,0) node[right,gray] {$x$};
    \draw[->,gray,very thick] (0,-2.25)--(0,2.25) node[above,gray] {$y$};
    \foreach \t in {-2,-1,1,2}{
      \node[gray,font=\scriptsize] at (\t,-0.18) {$\t$};
      \node[gray,font=\scriptsize] at (-0.2,\t) {$\t$};
    }
    \fill[main!12] plot[smooth cycle] coordinates {(0.1,-1.1) (1.5,-0.7) (2.05,0.35) (1.25,1.35) (-0.15,1.55) (-1.35,0.65) (-1.55,-0.6)};
    \draw[main!70!black,very thick] plot[smooth cycle] coordinates {(0.1,-1.1) (1.5,-0.7) (2.05,0.35) (1.25,1.35) (-0.15,1.55) (-1.35,0.65) (-1.55,-0.6)};
    \node[main!80!black] at (0.45,0.35) {二维区域（「面」）};
    \fill[second] (1.414,0) circle (0.06);
    \node[second,below right,font=\footnotesize] at (1.414,0) {$(\sqrt2,0)$};
    \node[third,font=\large] at (2.05,1.95) {$\mathbb{R}^2$};
  \end{tikzpicture}
  \caption{舞台就绪：完备的 $\mathbb{R}$ 与自身作笛卡尔积张成连续平面 $\mathbb{R}^2$（坐标轴、网格与一块二维区域的定场）。零维的点已被连成无缝的线，线再铺成二维的面——不等式解集、函数图象等后续对象都以这块平面为舞台}
  \label{fig:r2plane}
\end{figure}

\subsection{实数的三种等价面孔：分割、基本列与确界}

前面用「针眼筛子」与「实心金条」画出 $\mathbb{Q}$ 与 $\mathbb{R}$ 的差别。本节把它变成可以动手验证的数学：孔洞究竟在哪里？怎么补？补完的实数，为什么恰好就是数轴上的点？
（数系为何还要从 $\mathbb{R}$ 继续走向 $\mathbb{C}$，见第 2 章「数系扩张与虚数单位 $i$ 的几何诞生」。）

\key{探究一：稠密也会漏——有理数的上确界会失踪}
设
\[
  S=\{x\in\mathbb{Q}\mid x^2<2\}.
\]
$S\neq\varnothing$（$1\in S$），也有上界（$2$ 就是一个上界）。若在有理数范围内「有上界必有最小上界」，$S$ 应当存在最小上界。事实是：不存在。

\begin{derivation}
任取正有理数 $r$，构造它的「平分修正」
\[
  t=\frac{r+\dfrac{2}{r}}{2}\in\mathbb{Q},\qquad
  t^2-2=\frac{(r^2-2)^2}{4r^2}\geq0,\qquad
  t-r=-\frac{r^2-2}{2r}.
\]
\begin{itemize}
  \item 若 $r^2>2$（此时 $r$ 是 $S$ 的上界）：仍有 $t^2>2$，即 $t$ 也是上界；但 $t<r$。于是任何有理上界都能被更小的有理上界取代，「最小上界」无从谈起。
  \item 若 $r^2<2$（此时 $r\in S$）：$t>r$ 且 $t^2>2$，即 $t$ 是一个比 $r$ 更大的上界。
\end{itemize}
两种情形合起来：$\mathbb{Q}$ 中既没有 $S$ 的最小上界，也没有恰好落在 $\sqrt2$ 处的点。孔洞不是文学比喻，而是确界原理在 $\mathbb{Q}$ 中失效。
\end{derivation}

\key{探究二：补洞的三条等价路线}
\begin{itemize}
  \item \key{戴德金分割}：把 $\mathbb{Q}$ 切成上下两堆 $(A\mid B)$，要求 $A$ 无最大元；每个分割定义一个实数。上面由 $S$ 诱导的分割 $\bigl(\{x\in\mathbb{Q}:x^2<2\}\mid\mathbb{Q}\setminus\{x\in\mathbb{Q}:x^2<2\}\bigr)$ 恰好补上缺失的 $\sqrt2$。
  \item \key{柯西基本列}：把「彼此无限靠近」的有理数列（如 $3,\,3.1,\,3.14,\,\dots$）视为一个实数；两列若差趋于 $0$ 就视为同一个实数。$\pi$ 是这样被「逼」出来的。
  \item \key{公理刻画}（只描述、不构造）：确界原理、单调有界定理、闭区间套定理、柯西收敛准则四者等价，说的是同一句话——$\mathbb{R}$ 没有孔洞。高中阶段的极限、零点存在定理、介值性全部建立在这条公理上。
\end{itemize}

\begin{center}
\footnotesize
\setlength{\tabcolsep}{3pt}
\renewcommand{\arraystretch}{1.4}
\begin{tabular}{p{2.4cm}p{3.6cm}p{3.0cm}p{3.8cm}}
\toprule
视角 & 怎么补洞 & 典型产物 & 高中落点 \\
\midrule
戴德金分割 & 用「上下两堆」的卡口位置定义新数 & $\sqrt2$ 这类不可公度长度 & 数轴上的点与实数一一对应 \\
柯西基本列 & 用「越来越近的有理列」定义新数 & $\pi,\ e$ 这类极限 & 数列极限、级数求和的合法性 \\
公理刻画 & 直接声明「有上界就有上确界」 & 完备性本身 & 零点定理、最值定理的理论根基 \\
\bottomrule
\end{tabular}
\end{center}

\key{探究三：数轴上的点恰好是实数}
稠密 $\neq$ 连续：$\mathbb{Q}$ 已经稠密（任意两点之间还有有理数），却仍漏点；只有补上完备性，才得到「每一条几何长度都对应唯一实数」的连续数轴。
十进制表示还带来两个容易含糊的点：
\begin{itemize}
  \item $0.999\dots=1$——无限小数本身要靠完备性才能定义清楚；
  \item 「无理数 $=$ 无限不循环小数」只是无理数的表征，不能当作定义：要说明无限小数的全体恰好铺满完备数轴，绕不开分割或基本列。
\end{itemize}
这两条留作课后思考（见课后作业拓展探究）。

\begin{conccard}[一级]{实数的三种等价面孔：分割、基本列与确界}{判断一个集合能否承载函数的单调性、零点定理与极限}{把「实数 $=$ 无限小数」当定义；把「稠密」误当「连续」}
实数的三种定义（戴德金分割、柯西基本列、确界公理）彼此等价，共同特征是「没有孔洞」。
$\mathbb{Q}$ 的致命伤不在稀疏而在漏点：$S=\{x\in\mathbb{Q}:x^2<2\}$ 有上界却无最小上界，任何有理数都填不进 $\sqrt2$ 的卡口。
高中阶段不必背构造细节，但必须知道：极限、零点定理、最值定理的合法性来自完备性，而不只是来自「能比大小」。
\end{conccard}

\begin{commonerror}{把「稠密」当成「连续」}
常见错误：认为「任意两个有理数之间还有有理数，所以有理数轴已经连续」。稠密只保证「找不到相邻两点」，不保证「不漏点」：$\sqrt2,\pi,e$ 都不在 $\mathbb{Q}$ 里，曲线 $y=x^2-2$ 照样从孔洞中穿过。
判断标准：能否说「有上界的非空集合必有最小上界」。$\mathbb{Q}$ 不能，$\mathbb{R}$ 能。
\end{commonerror}

\begin{method}
「看集合选武器」再补一条反射：遇到「变号却无零点／有界却无最值／递增却撞不到天花板」，先怀疑定义域是 $\mathbb{Q}$ 而非 $\mathbb{R}$；遇到「长度、面积、开方、$\pi$」，先确认舞台是 $\mathbb{R}$。
\end{method}

% ---------- 以上为背景铺垫，以下进入主线 ----------

初中我们说"$y$ 随 $x$ 的变化而变化"，高中要说得更精确：函数是一种特殊的对应，它的图象是平面上的点集。本导学案用一条主线重新组织"函数"概念——
\key{关系}（$R\subseteq A\times B$ 的三档舞台：离散$\times$离散的孤立点、离散$\times$连续的竖直线段族、连续$\times$连续的二维区域）$\xrightarrow{\text{方程／不等式＝条件的表达法}}$ \key{两条限制轴}（轴一·单值性 $\to$ 映射与函数；轴二·连续性与结构保持 $\to$ 基本初等函数），并在每一步用"\key{代数有序数对集合 $\leftrightarrow$ 几何坐标系点集 $\leftrightarrow$ 映射箭头图}"三种语言互相翻译。

\begin{learnobjective}
\begin{itemize}
  \item 从\key{点集视角}理解二元关系与函数：把对应看成 $\mathbb{R}^2$ 平面上的点集，说清"哪些点被选中"；
  \item 理解\key{单值性}（垂线检验）如何把任意点集过滤成函数图象：每条竖直线 $x=a$ 与点集至多一个交点；
  \item 掌握\key{三视角互译}：代数有序数对集合 $\leftrightarrow$ 几何坐标系点集 $\leftrightarrow$ 映射箭头图；
  \item 理解\key{两条限制轴}：轴一「单值性」回答它是不是函数，轴二「连续性／结构保持」回答它好不好算，二者互不替代；
  \item 说明\key{方程与不等式是关系的条件表达法}，并在大一统框架下说出二元方程、函数、一元方程（不等式）的解分别是几维对象。
\end{itemize}
\end{learnobjective}

\begin{preview}
\subsection*{知识回顾}
\begin{itemize}
  \item 笛卡尔积：$A\times B=\{(a,b)\mid a\in A,\ b\in B\}$。若 $A=\{1,2\}$，$B=\{3\}$，则 $A\times B=$ \blank；
  \item 二元关系：$A\times B$ 的\key{任意子集} $R$ 都叫 $A$ 到 $B$ 的二元关系。若 $R=\{(1,3),(2,3)\}$，则 $R$ 中元素个数是 \blank；
  \item 不等式解集：$x^2-1\leq0$ 的解集是 \blank，它表示数轴上的 \blank；
  \item 平面直角坐标系中，点 $(x,y)$ 与有序数对 $(x,y)$ 是 \blank（一一对应／没有关系）。
\end{itemize}

\subsection*{课前速读}
阅读「函数的概念」一节，完成下列问题：
\begin{enumerate}
  \item 函数定义中的两个关键词是"\key{任意}"和"\key{唯一}"；
  \item 函数的三要素是 \blank、\blank、\blank；
  \item 下列对象中哪些是函数：(1) 圆的面积 $S$ 与半径 $r$；(2) 一个人的年龄与身高；(3) 平方对应 $x\mapsto x^2$。（答：\blank）
\end{enumerate}
\end{preview}

\begin{explore}

\subsection{有限离散关系：从 $A\times B$ 出发}

\key{探究一：把所有可能的配对都列出来}
设 $A=\{1,2,3\}$，$B=\{2,4,6\}$。
\begin{itemize}
  \item $A\times B$ 共有 $3\times3=9$ 个有序数对：
        $(1,2),(1,4),(1,6),(2,2),(2,4),(2,6),(3,2),(3,4),(3,6)$。
  \item 把它看成\key{网格点集}：在平面直角坐标系中描出这 $9$ 个点（图 \ref{fig:grid9}）。
  \item $A\times B$ 的任意子集 $R$ 都叫 $A$ 到 $B$ 的一个\key{二元关系}。例如 $R=\{(1,2),(2,4),(3,6)\}$ 只是"挑出其中几个配对"。
\end{itemize}

\begin{figure}[htbp]
  \centering
  \begin{tikzpicture}
    \begin{axis}[
      width=0.55\textwidth, height=6cm,
      xmin=0.5, xmax=3.5, ymin=1, ymax=7,
      xtick={1,2,3}, ytick={2,4,6},
      xlabel=$x$, ylabel=$y$,
      grid=major, grid style={dashed,gray!35},
      tick label style={font=\footnotesize},
      title={\small $A\times B$ 的 9 个有序数对（网格点）},
      title style={font=\small},
    ]
      \addplot[only marks, mark=*, mark size=2pt, main!70!black]
        coordinates {(1,2)(1,4)(1,6)(2,2)(2,4)(2,6)(3,2)(3,4)(3,6)};
    \end{axis}
  \end{tikzpicture}
  \caption{$A\times B$ 的点集视角：平面直角坐标系中的 $3\times3$ 网格点集。竖直方向看，每一列对应同一个 $x\in A$，各含 $3$ 个点（一对多，尚未通过垂线检验）；代数清单见正文，映射箭头图视角见图 \ref{fig:arrow}}
  \label{fig:grid9}
\end{figure}

\key{同一母体，三档舞台：}$R\subseteq A\times B$ 长什么样，首先取决于 $A,B$ 是哪种集合（即前文「舞台四级进化」装了多少装备）。三档并列如下：
\begin{itemize}
  \item \key{离散$\times$离散}（$A=\{1,2,3\},\ B=\{2,4,6\}$）：$R$ 只能是有限个孤立点——图 \ref{fig:grid9} 的 $9$ 点网格；
  \item \key{离散$\times$连续}（$A=\{1,2,3\},\ B=[0,4]$）：每个 $x$ 上方可以出现一整段竖直线段（点集是「三根竖条」，图 \ref{fig:three-tiers} 中），也可以只取有限个点——具体形态仍由所加的条件决定；
  \item \key{连续$\times$连续}（$A=B=[0,3]$）：$R$ 可以是一整块二维区域——图 \ref{fig:three-tiers} 右的矩形。
\end{itemize}
三档共用同一套三语言：代数写 $R=\{(x,y)\mid\text{条件}\}$，几何描点，箭头图看方向。\textbf{舞台决定点能密到什么程度，条件决定点被约束成什么形状}；二者缺一不可。特别地，$A\times B$（$9$ 点网格）本身也是一个二元关系——「全取」的关系；它代表\textbf{两个变元完全自由}的起点，而不是终点，后面的全部内容都在给这份自由逐级上锁。

\begin{figure}[htbp]
  \centering
  \begin{minipage}{0.31\textwidth}
    \centering
    \begin{tikzpicture}[x=0.55cm,y=0.55cm,>=stealth,font=\scriptsize]
      \draw[->,gray] (0.5,0)--(3.6,0) node[right] {$x$};
      \draw[->,gray] (0,-0.2)--(0,4.6) node[above] {$y$};
      \foreach \x in {1,2,3}{ \foreach \y in {2,4,6}{ \fill[main!70!black] (\x,\y) circle (1.6pt); } }
      \node at (2,-0.8) {离散$\times$离散：$9$ 个孤立点};
    \end{tikzpicture}
  \end{minipage}\hfill
  \begin{minipage}{0.31\textwidth}
    \centering
    \begin{tikzpicture}[x=0.55cm,y=0.55cm,>=stealth,font=\scriptsize]
      \draw[->,gray] (0.5,0)--(3.6,0) node[right] {$x$};
      \draw[->,gray] (0,-0.2)--(0,4.6) node[above] {$y$};
      \foreach \x/\yl/\yu in {1/0.2/1.6, 2/0.6/3.4, 3/0/2.2}{
        \draw[second,very thick] (\x,\yl)--(\x,\yu);
        \fill[second] (\x,\yl) circle (1.6pt);
        \fill[second] (\x,\yu) circle (1.6pt);
      }
      \node at (2,-0.8) {离散$\times$连续：三根竖条};
    \end{tikzpicture}
  \end{minipage}\hfill
  \begin{minipage}{0.31\textwidth}
    \centering
    \begin{tikzpicture}[x=0.55cm,y=0.55cm,>=stealth,font=\scriptsize]
      \draw[->,gray] (0.5,0)--(3.6,0) node[right] {$x$};
      \draw[->,gray] (0,-0.2)--(0,4.6) node[above] {$y$};
      \fill[main!12] (0,0) rectangle (3,4);
      \draw[main!70!black,thick] (0,0) rectangle (3,4);
      \node at (2,-0.8) {连续$\times$连续：一整块区域};
    \end{tikzpicture}
  \end{minipage}
  \caption{关系 $R\subseteq A\times B$ 的三档形态：同一母体，因 $A,B$ 的装备不同，分别长成孤立点云、竖直线段族、二维区域。三档都可以同时用「有序数对集合／坐标点集／箭头图」三种语言描述}
  \label{fig:three-tiers}
\end{figure}

\key{追问：}把这三档点集都当作「从 $x$ 到 $y$ 的对应」，它们各有什么缺陷，迫使我们继续加限制条件？
（提示：$9$ 点网格中取 $x=1$，$y$ 可以是 $2,4,6$ 中任意一个——对应不唯一；竖条与矩形中一个 $x$ 拖着无穷多个 $y$；而且离散点集没有「连续变化」的资格，圆滑曲线、连续性都无从谈起。）

\begin{example}
写出 $A=\{1,2,3\}$、$B=\{2,4,6\}$ 时 $A\times B$ 的全部 $9$ 个有序数对并在网格中描点；再任取两个不同的二元关系 $R_1,R_2\subseteq A\times B$（各含 $2$ 个元素）。
\end{example}
\begin{solution}
$9$ 个有序数对为 $(1,2),(1,4),(1,6),(2,2),(2,4),(2,6),(3,2),(3,4),(3,6)$，描点见图 \ref{fig:grid9}。
任取 $R_1=\{(1,2),(2,6)\}$，$R_2=\{(1,4),(3,2)\}$ 即可（答案不唯一）。它们都是 $A\times B$ 的子集，所以都是二元关系。
\end{solution}

\key{探究二：单值性——每列至多一个点}
对比两个关系：
\[
  R_1=\{(1,2),(1,4),(2,6)\},\qquad R_2=\{(1,2),(2,4),(3,6)\}.
\]
用三视角分别表示：
\begin{itemize}
  \item \key{代数}（有序数对集合）：如上，$R_1$ 中 $x=1$ 出现了两次；
  \item \key{几何}（点集）：在网格中描点（图 \ref{fig:r1r2}）。$R_1$ 在竖直线 $x=1$ 上有两个点 $(1,2),(1,4)$；$R_2$ 在每条竖直线 $x=1,2,3$ 上各有一个点；
  \item \key{映射}（箭头图）：图 \ref{fig:arrow}。$R_1$ 中 $1\to2$ 与 $1\to4$ 同时出现（\key{一对多}）；$R_2$ 中 $1\to2$、$2\to4$、$3\to6$（一对一）。
\end{itemize}

\begin{figure}[htbp]
  \centering
  \begin{minipage}{0.48\textwidth}
    \centering
    \begin{tikzpicture}
      \begin{axis}[
        width=\textwidth, height=5cm,
        xmin=0.5, xmax=3.5, ymin=1, ymax=7,
        xtick={1,2,3}, ytick={2,4,6},
        grid=major, grid style={dashed,gray!35},
        tick label style={font=\footnotesize},
        title={\small $R_1$：$x=1$ 处有两个点（一对多）},
        title style={font=\small},
      ]
        \addplot[only marks, mark=*, mark size=2pt, second]
          coordinates {(1,2)(1,4)(2,6)};
      \end{axis}
    \end{tikzpicture}
  \end{minipage}\hfill
  \begin{minipage}{0.48\textwidth}
    \centering
    \begin{tikzpicture}
      \begin{axis}[
        width=\textwidth, height=5cm,
        xmin=0.5, xmax=3.5, ymin=1, ymax=7,
        xtick={1,2,3}, ytick={2,4,6},
        grid=major, grid style={dashed,gray!35},
        tick label style={font=\footnotesize},
        title={\small $R_2$：每列恰有一个点（一对一）},
        title style={font=\small},
      ]
        \addplot[only marks, mark=*, mark size=2pt, main]
          coordinates {(1,2)(2,4)(3,6)};
      \end{axis}
    \end{tikzpicture}
  \end{minipage}
  \caption{$R_1$ 与 $R_2$ 的点集对比：竖直线 $x=1$ 与 $R_1$ 有两个交点，与 $R_2$ 只有一个交点}
  \label{fig:r1r2}
\end{figure}

\begin{figure}[htbp]
  \centering
  \begin{minipage}{0.48\textwidth}
    \centering
    \begin{tikzpicture}[>=Stealth, semithick]
      \node (a1) at (0,2.2) {$1$}; \node (a2) at (0,1.1) {$2$}; \node (a3) at (0,0) {$3$};
      \node (b2) at (3,2.2) {$2$}; \node (b4) at (3,1.1) {$4$}; \node (b6) at (3,0) {$6$};
      \draw[->] (a1)--(b2); \draw[->] (a1)--(b4); \draw[->] (a2)--(b6);
      \node[font=\footnotesize] at (0,2.8) {$A$}; \node[font=\footnotesize] at (3,2.8) {$B$};
    \end{tikzpicture}
    \\[2pt]{\footnotesize $R_1$：$1$ 同时指向 $2$ 和 $4$}
  \end{minipage}\hfill
  \begin{minipage}{0.48\textwidth}
    \centering
    \begin{tikzpicture}[>=Stealth, semithick]
      \node (a1) at (0,2.2) {$1$}; \node (a2) at (0,1.1) {$2$}; \node (a3) at (0,0) {$3$};
      \node (b2) at (3,2.2) {$2$}; \node (b4) at (3,1.1) {$4$}; \node (b6) at (3,0) {$6$};
      \draw[->] (a1)--(b2); \draw[->] (a2)--(b4); \draw[->] (a3)--(b6);
      \node[font=\footnotesize] at (0,2.8) {$A$}; \node[font=\footnotesize] at (3,2.8) {$B$};
    \end{tikzpicture}
    \\[2pt]{\footnotesize $R_2$：每个 $x$ 只指向一个 $y$}
  \end{minipage}
  \caption{映射箭头图：$R_1$ 一对多（不是函数），$R_2$ 一对一（是函数）}
  \label{fig:arrow}
\end{figure}

\key{判断规则（垂线检验雏形）：}点集 $P\subseteq\mathbb{R}^2$ 能作为某个函数的图象，当且仅当\key{每条竖直线 $x=a$ 与 $P$ 至多有一个交点}（"每列至多 $1$ 个点"）。据此 $R_2$ 是函数，$R_1$ 不是函数。

为什么 $R_2$ 是函数：对 $A$ 中每个 $x$，都有唯一 $y$ 与之对应；而 $R_1$ 中 $x=1$ 同时对应 $2$ 和 $4$，违反唯一性。

\key{把对应「翻译」成有序数对（图象的第三种身份）：}若对应关系为 $f$，把每个 $x$ 与它的像 $y=f(x)$ 并排写成 $(x,y)=(x,f(x))$。于是「哪些数对被选中」就是集合
\[
  \Gamma_f=\{(x,f(x))\mid x\in A\}.
\]
它读作映射，就是箭头图；读作有序数对集合，就是 $\Gamma_f$；坐标化到平面直角坐标系，就是函数的\key{图象}。三种说法在此合流——后文说「图象」，指的就是这个数对集合，不是额外多出来的对象。

\key{追问：}如果允许"一对多"，还能叫函数吗？"一对多"的对象在数学里是什么？
（提示：它仍是\key{关系}。例如二元方程 $x^2+y^2=4$ 的解集就是一对多的点集——这为探究三埋下伏笔。）

\begin{example}
用三视角（代数、几何、箭头图）表示关系 $R=\{(1,2),(2,4),(3,6)\}$，并判断它是否为函数。
\end{example}
\begin{solution}
代数：$R=\{(1,2),(2,4),(3,6)\}$。
几何：描点 $(1,2),(2,4),(3,6)$，三点在同一直线 $y=2x$ 上，且每条竖直线 $x=1,2,3$ 上各有一个点。
箭头图：$1\to2$，$2\to4$，$3\to6$。
是函数：每个 $x\in A$ 都有唯一 $y$ 与之对应（事实上 $y=2x$）。
\end{solution}

\begin{problem}
在 $A=\{1,2,3\}$、$B=\{2,4,6\}$ 的网格中，任取一个二元关系 $R$，使它不是函数，并说明理由；再任取一个函数关系，写出它的对应法则。
\end{problem}
\begin{solution}
例：$R=\{(1,2),(1,6),(2,4)\}$。$x=1$ 同时对应 $2$ 和 $6$，违反单值性，不是函数（竖直线 $x=1$ 与点集有两个交点）。
例：$R=\{(1,4),(2,2),(3,6)\}$ 是函数（每列至多一个点），对应法则可写为 $f(1)=4$，$f(2)=2$，$f(3)=6$。
\end{solution}

\key{从「完全自由」到「规则登场」：一条收紧链}

回到 $A=\{1,2,3\}$、$B=\{2,4,6\}$：$A\times B$ 的 $9$ 个配对是\emph{两个变元完全自由}时的全部候选。逐步加限制，自由度这样消失：
\begin{itemize}
  \item \key{二元关系}：对 $9$ 个配对逐个决定「取／不取」，共 $2^9=512$ 种——仍几乎无规律；
  \item \key{单值性}：每列至多取一个。每列有「不取，或取 $2,4,6$ 之一」共 $4$ 种选择，合计 $4^3=64$ 种；若要求 $A$ 中每个 $x$ 都有对应，则只剩 $3^3=27$ 种；
  \item \key{统一法则}：要求用同一条公式（结构保持等式）表达所有对应——$27$ 种里只剩极少数，例如 $(1,2),(2,4),(3,6)$ 可被 $y=2x$ 一网打尽。
\end{itemize}
函数不是外加的定义，而是这条收紧链上的一环：\textbf{自由消失多少，规则就登场多少}；基本初等函数位于链条顶端——规则收得最紧、结构保存得最好。

\subsection{连续二维区域：不等式解集的几何本质}

\key{先约定：方程与不等式是「从 $A\times B$ 里挑子集」的条件句}

含变元的等式或不等式，读法统一为「满足条件的有序数对全体」：
\begin{itemize}
  \item 一元方程 $x^2=4$：只有一个变元，解集是数轴上的点 $\{-2,2\}$；
  \item 二元方程 $x^2+y^2=4$：两个变元，解集是平面曲线（圆周）；
  \item 二元不等式 $x^2+y^2\leq4$：同一条件放松为「$\leq$」，解集从曲线「鼓成」闭圆盘。
\end{itemize}
所以方程、不等式不是与函数、映射并列的另一条线，而是\key{关系的具体表达法}：变元个数决定几何维度，等号／不等号决定边界还是区域。下面用这把尺子重新看连续点集。

\key{探究三：把 $9$ 个点换成无穷多个点}
把取值范围从离散集合推广到连续区间，考察两个平面点集：
\[
  R_3=\{(x,y)\mid x^2+y^2\leq4\},\qquad
  R_4=\{(x,y)\mid x+y\leq3,\ x\geq0,\ y\geq0\}.
\]
\begin{itemize}
  \item $R_3$ 是圆心在原点、半径为 $2$ 的\key{闭圆盘}（一块二维区域），其边界恰是二元方程 $x^2+y^2=4$，内部由不等式 $x^2+y^2\leq4$ 给出——「方程给边界、不等式给区域」的第一个实例；
  \item $R_4$ 是由 $x$ 轴、$y$ 轴与直线 $x+y=3$ 围成的\key{三角形区域}（图 \ref{fig:region}）。
  \item 它们都是 $\mathbb{R}^2$ 中"一块黑乎乎的区域"，而不是孤立的点。
\end{itemize}

关键观察：取 $x=1$，在 $R_3$ 中 $y\in[-\sqrt3,\sqrt3]$，有无穷多个取值；在 $R_4$ 中 $y\in[0,2]$，也有无穷多个取值。一个 $x$ 对应无穷多个 $y$——\key{失控}，无法由 $x$ 唯一预测 $y$。

\key{追问：}二元不等式（组）的解集在几何上是什么？（答：$\mathbb{R}^2$ 中的二维区域。）这正是未来\key{线性规划可行域}的伏笔：满足一组约束的 $(x,y)$ 全体就是平面区域，而最优解往往出现在区域的顶点或边界上。

\begin{figure}[htbp]
  \centering
  \begin{minipage}{0.48\textwidth}
    \centering
    \begin{tikzpicture}[scale=1.1]
      \draw[->] (-2.4,0)--(2.4,0) node[right]{$x$};
      \draw[->] (0,-2.4)--(0,2.4) node[above]{$y$};
      \fill[main!15] (0,0) circle (2cm);
      \draw[main!70!black, thick] (0,0) circle (2cm);
      \draw[dashed, second, thick] (1,-1.732)--(1,1.732)
        node[above right, font=\footnotesize]{$x=1:y\in[-\sqrt3,\sqrt3]$};
      \fill[second] (1,1.732) circle (1.5pt);
      \fill[second] (1,-1.732) circle (1.5pt);
    \end{tikzpicture}
    \\[2pt]{\footnotesize $R_3$：闭圆盘，竖直线 $x=1$ 与它交于无穷多点}
  \end{minipage}\hfill
  \begin{minipage}{0.48\textwidth}
    \centering
    \begin{tikzpicture}[scale=1.1]
      \draw[->] (-0.5,0)--(3.6,0) node[right]{$x$};
      \draw[->] (0,-0.5)--(0,3.6) node[above]{$y$};
      \fill[main!15] (0,0)--(3,0)--(0,3)--cycle;
      \draw[main!70!black, thick] (0,0)--(3,0)--(0,3)--cycle;
      \draw[dashed, second, thick] (1,0)--(1,2)
        node[right, font=\footnotesize]{$x=1:y\in[0,2]$};
    \end{tikzpicture}
    \\[2pt]{\footnotesize $R_4$：三角形区域，竖直线 $x=1$ 与它交于无穷多点}
  \end{minipage}
  \caption{连续二维区域：二元不等式（组）的解集是 $\mathbb{R}^2$ 中的"面"，不是曲线}
  \label{fig:region}
\end{figure}

\begin{example}
在坐标系中画出 $R_4=\{(x,y)\mid x+y\leq3,\ x\geq0,\ y\geq0\}$，并求当 $x=1$ 时 $y$ 的取值范围。
\end{example}
\begin{solution}
画出直线 $x+y=3$ 与两条坐标轴，取左下方的三角形内部（含边界），见图 \ref{fig:region} 右图。
当 $x=1$ 时，条件化为 $1+y\leq3$ 且 $y\geq0$，即 $0\leq y\leq2$。$y$ 有无穷多个取值，无法由 $x=1$ 唯一确定——所以 $R_4$ 不是函数。
\end{solution}

\begin{problem}
设 $R=\{(x,y)\mid x^2+y^2\leq4\}$。判断 $R$ 是否为函数，并用垂线检验说明理由。
\end{problem}
\begin{solution}
不是函数。取 $x=0$，则 $y\in[-2,2]$，一个 $x$ 对应无穷多个 $y$；几何上看，竖直线 $x=0$ 与圆盘相交于无穷多个点，违反"每列至多一个点"。
\end{solution}

\subsection{两条限制轴：从区域到基本初等函数}

\key{探究四：单值性与正则性，是两条互相独立的轴}

任意点集（撒点）常常同时「一对多」又「不连续」。要把它变成能研究的函数，需要在\textbf{两条不同的轴}上分别加限制，二者不可互相替代：

\begin{itemize}
  \item \key{轴一·代数／规则——单值性。}\key{回扣：}这正是前文那条收紧链——从 $A\times B$ 的完全自由（$2^9$ 种取法），到单值（$4^3$；若处处有定义则 $3^3=27$），再到统一法则。具体做法：规定每条竖直线 $x=a$ 上\textbf{至多}保留一个点，并且\textbf{明确指定}保留哪一个（不能随手挑）；未指定的 $x$ 允许没有对应点——这些被指定的 $x$ 全体就是定义域 $A$。满足单值性的关系叫\key{映射} $f:A\to B$；当 $A,B$ 取为数集时，就是\key{函数} $y=f(x)$。判别标准只有一条：垂线检验（图 \ref{fig:orderline} 第 (3) 幅）。再要求每个对应都能用同一条解析式表达，才得到「写得出来的函数」。
  \item \key{轴二·几何／分析——正则性与结构保持。}单值性\emph{不附带}任何光滑承诺：轴一走到尽头，图象照样可以处处间断（开场三问的 Dirichlet 函数），也可以连续却处处不可微。要让它「越来越好算」，必须继续加条件：连续性 $\to$ 可微性 $\to$ 某类「运算封闭／结构保持」的等式，最终落在\key{基本初等函数}上：
  \begin{itemize}
    \item 指数函数 $y=a^x$：把定义域的加法变成值域的乘法，$f(x_1+x_2)=f(x_1)f(x_2)$；
    \item 对数函数 $y=\log_a x$：把指数函数的数对 $(x,y)$ 调换为 $(y,x)$ 的逆关系，图象关于直线 $y=x$ 对称；
    \item 幂函数 $y=x^{\alpha}$：刻画幂次结构（面积、体积、增长阶）；
    \item 三角函数：刻画周期结构（旋转、振动、波动）。
  \end{itemize}
\end{itemize}

\key{两条轴的分工：}轴一回答「它是不是函数」，轴二回答「它好不好算」。而方程与不等式（见本节开头的约定）是描述点集条件的语言，在任何一档舞台、任何一条轴上都可用：用二元方程 $F(x,y)=0$ 描述轴二下的平面曲线（可以一对多，如圆 $x^2+y^2=4$）；用一元方程 $f(x)=g(x)$ 求两条函数图象交点的横坐标；用一元不等式 $f(x)>g(x)$ 表示「$f$ 的点集位于 $g$ 的点集上方」的 $x$ 的范围。

\begin{figure}[htbp]
  \centering
  \begin{minipage}{0.235\textwidth}
    \centering
    \begin{tikzpicture}[x=0.6cm,y=0.6cm]
      \draw[main!60!black,thick,fill=main!4] (-2.2,-2) rectangle (2.2,2);
      \pgfmathsetseed{11}
      \foreach \i in {1,...,140}{
        \pgfmathsetmacro{\px}{-2.05+4.1*rnd}
        \pgfmathsetmacro{\py}{-1.85+3.7*rnd}
        \fill[main!45] (\px,\py) circle (0.8pt);
      }
    \end{tikzpicture}
    \\[2pt]{\footnotesize (1) 无序：随手撒点——一个 $x$ 拖着一大段 $y$}
  \end{minipage}\hfill
  \begin{minipage}{0.235\textwidth}
    \centering
    \begin{tikzpicture}[x=0.6cm,y=0.6cm]
      \draw[main!60!black,thick,fill=main!4] (-2.2,-2) rectangle (2.2,2);
      \pgfmathsetseed{11}
      \foreach \i in {1,...,140}{
        \pgfmathsetmacro{\px}{-2.05+4.1*rnd}
        \pgfmathsetmacro{\py}{-1.85+3.7*rnd}
        \fill[main!22] (\px,\py) circle (0.8pt);
      }
      \draw[second,very thick] (-0.6,-2)--(-0.6,2);
      \foreach \yy in {-1.75,-1.3,-0.85,-0.4,0.05,0.5,0.95,1.4,1.8}
        \fill[second] (-0.6,\yy) circle (1.5pt);
      \node[second,font=\scriptsize,right] at (-0.52,1.82) {$x=a$};
    \end{tikzpicture}
    \\[2pt]{\footnotesize (2) 竖截：竖线 $x=a$ 只截一段，仍是一整列点}
  \end{minipage}\hfill
  \begin{minipage}{0.235\textwidth}
    \centering
    \begin{tikzpicture}[x=0.6cm,y=0.6cm]
      \draw[main!60!black,thick,fill=main!4] (-2.2,-2) rectangle (2.2,2);
      \draw[third!55,thin] plot coordinates {(-2.05,1.82)(-1.89,1.54)(-1.72,1.15)(-1.56,-0.49)(-1.39,0.13)(-1.23,-0.37)(-1.07,-0.38)(-0.9,-1.22)(-0.74,-1.25)(-0.57,-0.93)(-0.41,-1.39)(-0.25,-1.05)(-0.08,-1.8)(0.08,-1.7)(0.25,-0.61)(0.41,-1.05)(0.57,-0.81)(0.74,-0.33)(0.9,-0.83)(1.07,-1.18)(1.23,0.24)(1.39,0.12)(1.56,0.48)(1.72,0.69)(1.89,1.58)(2.05,0.87)};
      \foreach \p in {(-2.05,1.82),(-1.89,1.54),(-1.72,1.15),(-1.56,-0.49),(-1.39,0.13),(-1.23,-0.37),(-1.07,-0.38),(-0.9,-1.22),(-0.74,-1.25),(-0.57,-0.93),(-0.41,-1.39),(-0.25,-1.05),(-0.08,-1.8),(0.08,-1.7),(0.25,-0.61),(0.41,-1.05),(0.57,-0.81),(0.74,-0.33),(0.9,-0.83),(1.07,-1.18),(1.23,0.24),(1.39,0.12),(1.56,0.48),(1.72,0.69),(1.89,1.58),(2.05,0.87)}
        \fill[third] \p circle (1.5pt);
    \end{tikzpicture}
    \\[2pt]{\footnotesize (3) 每列一点：被钉住了，但点列仍在乱跳}
  \end{minipage}\hfill
  \begin{minipage}{0.235\textwidth}
    \centering
    \begin{tikzpicture}[x=0.6cm,y=0.6cm]
      \draw[main!60!black,thick,fill=main!4] (-2.2,-2) rectangle (2.2,2);
      \pgfmathsetseed{37}
      \foreach \i in {1,...,90}{
        \pgfmathsetmacro{\bx}{-2.05+4.1*(\i-1)/89}
        \pgfmathsetmacro{\by}{0.5*\bx*\bx-1+0.55*(2*rnd-1)}
        \fill[main!40] (\bx,\by) circle (0.7pt);
      }
      \draw[second,very thick,domain=-2.05:2.05,samples=60,smooth] plot(\x,{0.5*\x*\x-1});
      \node[second,font=\scriptsize] at (0,-1.78) {$y=\tfrac12x^2-1$};
    \end{tikzpicture}
    \\[2pt]{\footnotesize (4) 收束：粗带压扁成唯一的曲线}
  \end{minipage}
  \caption{图像线是怎样从点集里长出来的：同一块平面区域上的四步图像实验。(1) 无序——随手撒点，一个 $x$ 的正上方压着一大段 $y$，一对多；(2) 竖截——用竖直线 $x=a$ 去截点云，截出的仍是一整列（一段）点，一个 $x$ 对应多个 $y$；(3) 立规矩——每条竖线只留一个点，得到一般函数的图象，但点列仍无规律地上下跳动；(4) 收束——点越密，跳动的点列糊成一条有厚度的粗带，厚度趋于 $0$ 时只剩唯一的曲线 $y=\tfrac12x^2-1$。图象就是「按 $x$ 排好队、且每个 $x$ 只剩唯一 $y$」的点列极限}
  \label{fig:orderline}
\end{figure}

\begin{conccard}[一级]{单值性与垂线检验}{任意平面点集（或图象）}{把"每列至多一个点"误写成"每行至多一个点"；忽略定义域}
点集 $P\subseteq\mathbb{R}^2$ 是某函数 $y=f(x)$ 的图象，当且仅当每条竖直线 $x=a$ 与 $P$ 至多有一个交点。代数表述：对每个 $x$，满足 $(x,y)\in P$ 的 $y$ 至多一个。这条判据把"任意点集"过滤成"函数图象"，是探究二、探究三共用的核心工具。
\end{conccard}

\begin{conccard}[二级]{三视角互译}{二元关系、函数、方程的解集}{只在一个视角里思考；把"数对集合"与"点集"割裂}
同一对象有三种等价表示：代数（有序数对集合）$\leftrightarrow$ 几何（坐标系点集）$\leftrightarrow$ 映射（箭头图）。代数精确、几何直观、箭头图体现对应方向。解题时按需要在三种语言间切换：例如解 $f(x)>g(x)$，就是看 $f$ 的点集在 $g$ 的点集上方时 $x$ 的范围。
\end{conccard}

\begin{method}
三视角互译法：遇到函数、方程、不等式问题时，先写出集合（代数），再画图（几何），必要时画箭头图（映射），最后回到代数表达。例如判断 $f(x)>g(x)$ 的解：画两条曲线，看 $f$ 在 $g$ 上方的 $x$ 区间，再写出集合 $\{x\mid f(x)>g(x)\}$。
\end{method}

\begin{commonerror}{把二维区域误当函数}
误以为"式子中有 $x,y$ 就是函数"或"给出 $x$ 就能算出 $y$"。反例：$x^2+y^2=4$ 中 $x=0$ 对应 $y=\pm2$，是一对多，不是函数；$\{(x,y)\mid x+y\leq3\}$ 中 $x=1$ 对应无穷多个 $y$。判别时必须回到单值性：每个 $x$ 是否唯一确定 $y$。
\end{commonerror}

\begin{example}
说明方程 $x^2=4$ 的解可以看成两个函数图象交点的横坐标，并求出解。
\end{example}
\begin{solution}
$x^2=4$ 即 $x^2-4=0$。令 $f(x)=x^2-4$，$g(x)=0$，交点横坐标即 $f(x)=g(x)$ 的解：$x=\pm2$。
也可看作曲线 $y=x^2$ 与直线 $y=4$ 交点的横坐标，结果同样是 $x=\pm2$。
\end{solution}

\begin{problem}
设 $f(x)=x^2$，$g(x)=2x+3$。求不等式 $f(x)>g(x)$ 的解集，并指出它对应"哪段 $x$ 上 $f$ 的图象在 $g$ 上方"。
\end{problem}
\begin{solution}
$x^2>2x+3\iff x^2-2x-3>0\iff(x-3)(x+1)>0$，解集为 $(-\infty,-1)\cup(3,+\infty)$。
几何意义：在 $(-\infty,-1)$ 与 $(3,+\infty)$ 上，抛物线 $y=x^2$ 位于直线 $y=2x+3$ 的上方；在 $[-1,3]$ 上则位于下方（或相交）。
\end{solution}

\subsection{大一统结构对照}

把本节对象放进同一张表：它们都是「集合」，并且位于同一条包含链上——关系 $\supseteq$ 映射 $\supseteq$ 函数 $\supseteq$ 基本初等函数；方程与不等式则是按变元个数给关系分类的条件表达法。

\begin{center}
\begin{tabular}{|c|p{3.0cm}|p{5.2cm}|p{4.2cm}|}
\hline
\textbf{数学对象} & \textbf{集合与点集本质} & \textbf{几何维度直观理解} \\
\hline
二元关系 $R\subseteq A\times B$ & $A\times B$ 的任意子集（条件句挑出） & 点云：孤立点／竖直线段族／区域／曲线皆可 \\
\hline
映射 $f:A\to B$ & 每列至多一个点的单值关系 & 箭头图；数集情形为一维曲线 \\
\hline
二元不等式（组） & $\{(x,y)\mid \text{约束条件}\}$ & $\mathbb{R}^2$ 中的二维区域（"面"） \\
\hline
二元方程 $F(x,y)=0$ & $\{(x,y)\mid F(x,y)=0\}$ & 平面上的一维曲线（可一对多） \\
\hline
函数 $y=f(x)$ & $\{(x,f(x))\mid x\in A\}$ & 一维曲线，且每列至多一个点 \\
\hline
一元方程的解 & $\{x\mid f(x)=g(x)\}$ & 数轴上的零维点（交点横坐标） \\
\hline
一元不等式的解 & $\{x\mid f(x)>g(x)\}$ & 数轴上的一维区间（上下位置） \\
\hline
\end{tabular}
\end{center}

\begin{summary}
本节主线：关系 $R\subseteq A\times B$ 有三档舞台——离散$\times$离散的孤立点、离散$\times$连续的竖直线段族、连续$\times$连续的二维区域；方程与不等式是把关系写成条件句的表达法（变元个数定维度，等号／不等号定边界与区域）。在关系之上加限制要分两条独立的轴：轴一加单值性，得到映射与函数（判别标准只有垂线检验）；轴二加连续性、可微性与结构保持等式，得到基本初等函数（指数、对数、幂、三角）。三种语言（代数、几何、箭头图）全程互译。
\end{summary}

\subsection{函数与图形变换完全图谱}

\begin{conccard}[一级]{第一性原理：一切图像变换都是有序数对 $(x,y)$ 的映射}{抽象函数图像辨析、变换类压轴题}{只背“左加右减”口诀，不追踪每个点的去向}
函数图像是 $\mathbb{R}^2$ 中的点集 $\{(x,y)\mid y=f(x)\}$。所谓“图像变换”，本质是对集合中每个点的坐标施加映射。
掌握下面三级体系与对称周期生成律，就能把“看图识函数”变成可推导的代数操作。
\end{conccard}

\begin{center}
\renewcommand{\arraystretch}{1.4}
\begin{tabular}{|p{3.2cm}|p{4.0cm}|p{4.8cm}|}
\hline
\textbf{变换} & \textbf{点集规则} & \textbf{表达式与几何效应} \\
\hline
\multicolumn{3}{|l|}{\textbf{一级：刚体等距（形状大小不变）}} \\
\hline
平移 & $(x,y)\mapsto(x+h,\,y+k)$ & $y-k=f(x-h)$（逆向代换） \\
\cline{2-3}
轴对称（$x=a$） & $(x,y)\mapsto(2a-x,\,y)$ & $y=f(2a-x)$ \\
\cline{2-3}
中心对称（$(a,b)$） & $(x,y)\mapsto(2a-x,\,2b-y)$ & $y=2b-f(2a-x)$（$a=b=0$ 即奇函数） \\
\hline
\multicolumn{3}{|l|}{\textbf{二级：伸缩与翻折}} \\
\hline
轴向伸缩 & $(x,y)\mapsto(\frac1\omega x,\,Ay)$ & $y=Af(\omega x)$ \\
\cline{2-3}
绝对值翻折 & $(x,y)\mapsto(x,\,|y|)$；$x\mapsto|x|$ & $y=|f(x)|$（下方翻上）；$y=f(|x|)$（右侧复制到左侧，强制偶化） \\
\hline
\multicolumn{3}{|l|}{\textbf{三级：非线性算子（图像拓扑突变）}} \\
\hline
自变量平方／倒数 & $u=g(x)$ 作为内层 & $y=f(x^2)$ 必为偶函数且在 $x=0$ 处切线水平 \\
\cline{2-3}
函数值平方／倒数 & $(x,y)\mapsto(x,\,y^2)$；$(x,y)\mapsto(x,\,\frac1y)$ & 零点变极小值点／垂直渐近线（详见下文铁律） \\
\hline
\end{tabular}
\end{center}

\begin{method}
\textbf{双对称推周期：连续两次平行镜像反射 = 平移}。看到“两个对称条件”或“半周期反对称”，直接锁定最小周期，不必画图：
\begin{enumerate}
\item \textbf{双轴对称}：$f(a+x)=f(a-x)$ 且 $f(b+x)=f(b-x)$ $\Rightarrow$ $T=2|b-a|$；
\item \textbf{双中心对称}：$f(a+x)+f(a-x)=2c$ 且 $f(b+x)+f(b-x)=2c$ $\Rightarrow$ $T=2|b-a|$；
\item \textbf{一轴一中心}（最高频）：$f(a+x)=f(a-x)$ 且 $f(b+x)+f(b-x)=2c$ $\Rightarrow$ $T=4|b-a|$；
\item \textbf{滑移反对称}：$f(x+a)=-f(x)$、$f(x+a)=\pm\frac1{f(x)}$ 或 $f(x+a)=\frac{1-f(x)}{1+f(x)}$，迭代两次即复原 $\Rightarrow$ $T=2a$。
\end{enumerate}
几何直观：点 $P(x)$ 先关于 $x=a$ 反射到 $P_1(2a-x)$，再关于 $x=b$ 反射到 $P_2\bigl(2b-(2a-x)\bigr)=P(x+2(b-a))$——两次翻转方向复原，合成一次纯平移。
\end{method}

\begin{center}
\begin{tikzpicture}[scale=1.0]
\draw[->, gray, thick] (-1.2,0) -- (7.2,0) node[right] {$x$};
\draw[main, dashed, very thick] (1.5,-0.6) -- (1.5,2.4) node[above] {$x=a$};
\draw[second, dashed, very thick] (4.0,-0.6) -- (4.0,2.4) node[above] {$x=b$};
\fill[black] (0.5,1.5) circle (2pt) node[above] {$P(x)$};
\fill[main] (2.5,1.5) circle (2pt) node[below] {$P_1(2a-x)$};
\fill[second] (5.5,1.5) circle (2pt) node[below] {$P_2\bigl(x+2(b-a)\bigr)$};
\draw[<->, gray, thick] (0.5,-0.35) -- (5.5,-0.35) node[midway, below, font=\small\bfseries] {合成平移 $T=2(b-a)$};
\end{tikzpicture}
\end{center}

\begin{conccard}[二级]{导数对称遗传律：周期全遗传，对称奇偶互换}{含参导数压轴、切线分析}{把 $f$ 与 $f'$ 的对称性混为一谈}
对 $f$ 的恒等式两边求导（链式法则）即得：
\begin{itemize}
\item $f$ 周期为 $T$ $\iff$ $f'$ 周期仍为 $T$；
\item $f$ 关于直线 $x=a$ 轴对称 $\Rightarrow$ $f'$ 关于点 $(a,0)$ 中心对称，且 $f'(a)=0$；
\item $f$ 关于点 $(a,b)$ 中心对称 $\Rightarrow$ $f'$ 关于直线 $x=a$ 轴对称。
\end{itemize}
\end{conccard}

\begin{method}
\textbf{非线性算子的五条铁律}（以 $g(x)=\dfrac1{f(x)}$ 为例，$y=[f(x)]^2$ 见表中）：
\begin{enumerate}
\item \key{符号不变}：$f>0$ 处 $g>0$，$f<0$ 处 $g<0$，上下位置不动；
\item \key{零点撕裂}：$f(x_0)=0$ 变成垂直渐近线 $x=x_0$；
\item \key{无穷变零}：$f\to\infty$ 处 $g\to0$（趋近 $x$ 轴）；
\item \key{单调反转}：同一同号区间内 $g$ 与 $f$ 单调性相反（不能跨渐近线谈单调）；
\item \key{极值互换}：正的极大值 $M$ 变成极小值 $\frac1M$；而 $y=\pm1$ 的点是\key{不动点}，位置纹丝不动。
\end{enumerate}
$y=f(x^2)$ 与 $y=f\!\left(\frac1x\right)$ 则作用在自变量上：$f(x^2)$ 是偶函数，
左半支与右半支互为镜像（同一 $|x|$ 处函数值相同）；在 $x\ge0$ 上它等于把
$y=f(u)\,(u\ge0)$ 的横坐标按 $u=x^2$ 重标得到的曲线——靠近原点被压缩、
远离原点被拉伸；$f\!\left(\frac1x\right)$ 则把 $x\to0$ 的行为与
$x\to\pm\infty$ 的行为\key{对调}。

\textbf{平方算子 $y=[f(x)]^2$ 的五条铁律}（与倒数算子对照记忆）：
\begin{itemize}
  \item \key{非负化}：整个图像被压进 $y\ge0$，负部沿 $x$ 轴\key{翻折}上来；
  \item \key{零点不变}：$f(x_0)=0\iff f^2(x_0)=0$，只是零点的「重数」加倍；
  \item \key{无穷仍无穷}：$f\to\infty$ 时 $f^2\to+\infty$ 增长更快；
  \item \key{单调分区}：由 $(f^2)'=2ff'$，$f>0$ 处 $f^2$ 与 $f'$ 同号、
        $f<0$ 处反号，故在 $f$ 的正区间保持单调方向、负区间反转；
  \item \key{极值缩放}：$|f|$ 的极大值 $M$ 变成 $f^2$ 的极大值 $M^2$
        （$M\ge1$ 时变大、$0<M<1$ 时变小），$|f|$ 的小极小值则被压得更小；
        $y=1$ 上的点是\key{不动点}。
\end{itemize}
\end{method}

\begin{center}
\begin{tikzpicture}[scale=1.15]
\draw[->, gray, thick] (-1.8,0) -- (3.8,0) node[right] {$x$};
\draw[->, gray, thick] (0,-2.2) -- (0,3.0) node[above] {$y$};
\draw[gray, dashed, thick] (2,-2.2) -- (2,3.0) node[above] {$x=2$};
\draw[main, thick, domain=-1.4:3.4, samples=60] plot(\x, {(\x)*(\x-2)});
\node[main] at (3.1,1.6) {$f(x)=x^2-2x$};
\draw[second, thick, domain=-1.5:-0.12, samples=60] plot(\x, {1/((\x)*(\x-2))});
\draw[second, thick, domain=0.12:1.9, samples=80] plot(\x, {1/((\x)*(\x-2))});
\draw[second, thick, domain=2.12:3.5, samples=60] plot(\x, {1/((\x)*(\x-2))});
\fill[second] (1,-1) circle (2pt) node[below left] {$(1,-1)$};
\end{tikzpicture}
{\small\itshape 抛物线与其倒数：零点 $0,2$ 变为竖直渐近线，极小值 $-1$ 变为极大值 $-1$。}
\end{center}

\begin{example}[双对称推周期与导数遗传]
已知 $f(x)$ 在 $\mathbb{R}$ 上可导，$f(x+1)$ 为偶函数，$f(x+2)$ 为奇函数，且 $f(1)=2$，求 $f(2027)+f'(2025)$。
\begin{solution}
$f(x+1)$ 偶 $\Rightarrow f(1-x)=f(1+x)$，故 $f$ 关于 $x=1$ 轴对称；
$f(x+2)$ 奇 $\Rightarrow f(2-x)=-f(2+x)$，故 $f$ 关于 $(2,0)$ 中心对称。
一轴一中心 $\Rightarrow T=4|2-1|=4$。
由中心对称取 $t=1$：$f(3)=-f(1)=-2$，于是 $f(2027)=f(4\times506+3)=f(3)=-2$。
$f$ 关于 $x=1$ 轴对称 $\Rightarrow f'$ 关于 $(1,0)$ 中心对称，故 $f'(1)=0$；又 $f'$ 周期仍为 $4$，
$f'(2025)=f'(4\times506+1)=f'(1)=0$。合计 $\key{-2}$。
\end{solution}
\end{example}

\begin{example}[用奇偶性速判函数图像]
函数 $g(x)=\dfrac{x}{e^x-1}$ 的大致图像如何？
\begin{solution}
先看定义域 $x\neq0$。由切线不等式 $e^x>x+1\ (x\neq0)$ 知 $e^x-1$ 与 $x$ 同号，故 $g(x)>0$ 恒成立（全在 $x$ 轴上方）；且 $\lim\limits_{x\to0}g(x)=1$。
当 $x\to+\infty$ 时指数增长远快于一次项，$g(x)\to0^+$；当 $x\to-\infty$ 时 $e^x\to0$，$g(x)\approx\dfrac{x}{-1}=-x\to+\infty$，即左端贴着斜率 $-1$ 的直线向上。无需二阶求导即可锁定图像。
\end{solution}
\end{example}

\begin{commonerror}{把“对称”当“周期”直接用}
$f(a+x)=f(a-x)$ 只说明关于 $x=a$ 对称，一个对称条件\key{产生不了周期}；必须有两个对称条件，
或用滑移反对称结构。此外 $y=f(x^2)$ 是偶函数，而 $y=[f(x)]^2$ 只保证非负——两者都“偶化”方向不同，不可互推。
\end{commonerror}

图形变换还有一种对偶记法，能把「左加右减」的口诀风险降到最低：动曲线与动坐标系是同一事实的两种记账方式（图 \ref{fig:activepassive}）。

\begin{methodbox}{主被动对偶律：曲线平移，等价于坐标系反向平移}
设曲线 $C:y=f(x)$ 主动向右平移 $h$ 个单位，得到新曲线 $y=f(x-h)$；若改为让坐标轴向左平移 $h$ 个单位、曲线在空间中保持不动，则新读数满足 $x'=x+h$、$y'=y$，同一物理状态在新系下的方程仍写作 $y'=f(x'-h)$。
两式描述的是同一幅终态图：顶点的读数相同，只是「谁在动」的归属不同。若条件按曲线平移给出（如「图象右移 $2$」），用主动记法；若条件按新原点、新读数给出，用被动记法。两者方向相反是必然，切勿把坐标轴的移动方向也套用「左加右减」。
\end{methodbox}

\begin{figure}[htbp]
  \centering
  \begin{tikzpicture}[>=stealth,font=\small]
    % ---------- 左：主动变换（坐标系不动，曲线右移） ----------
    \begin{scope}
      \draw[->,gray,thick] (-1.4,0)--(3.6,0) node[right,gray,font=\footnotesize] {$x$};
      \draw[->,gray,thick] (0,-0.7)--(0,3.0) node[above,gray,font=\footnotesize] {$y$};
      \draw[gray!70,dashed,thick,domain=-1.15:1.15,samples=50,smooth] plot(\x,{\x*\x});
      \draw[main,very thick,domain=0.85:3.15,samples=60,smooth] plot(\x,{(\x-2)*(\x-2)});
      \fill[gray!70] (0,0) circle (0.05);
      \node[gray!70,below left,font=\footnotesize] at (0,0) {$(0,0)$};
      \fill[main] (2,0) circle (0.05);
      \node[main,below right,font=\footnotesize] at (2,0) {$(2,0)$};
      % 平移箭头画在 x 轴下方空白处，避免压曲线
      \draw[<->,second,thick] (0.15,-0.35)--(1.85,-0.35)
        node[midway,above,font=\footnotesize] {右移 $2$};
      \node[main,font=\footnotesize] at (2.55,2.5) {$y=(x-2)^2$};
      \node[font=\footnotesize] at (1.1,-1.45) {主动：坐标系不动，曲线右移 $2$};
    \end{scope}
    % ---------- 右：被动变换（曲线不动，坐标轴左移） ----------
    \begin{scope}[xshift=7.7cm]
      \draw[gray!55,dashed,thick] (-1.4,0)--(2.2,0);
      \draw[gray!55,dashed,thick] (0,-0.7)--(0,3.0);
      \node[gray!55,font=\scriptsize,right] at (-1.35,2.55) {旧系（虚线）};
      \draw[->,second,thick] (-3.5,0)--(1.7,0) node[right,second,font=\footnotesize] {$x'$};
      \draw[->,second,thick] (-2,-0.7)--(-2,3.0) node[above,second,font=\footnotesize] {$y'$};
      \draw[main,very thick,domain=-1.15:1.15,samples=50,smooth] plot(\x,{\x*\x});
      \fill[main] (0,0) circle (0.05);
      \node[gray!60,font=\scriptsize,left] at (-0.06,0.28) {旧 $O$};
      \node[second,below left,font=\footnotesize] at (-2,0) {$O'$};
      \node[main,font=\footnotesize,right] at (0.12,-0.32) {顶点 $(2,0)$};
      \draw[<->,third,thick] (-2,-0.85)--(0,-0.85)
        node[midway,below,font=\footnotesize] {轴左移 $2$};
      \draw[third!70!black,->,shorten <=1pt,shorten >=1pt] (0.2,-0.68)--(0.06,-0.12);
      \node[third!70!black,font=\scriptsize,right] at (0.25,-0.72) {读数 $x'=2$};
      \node[main,font=\footnotesize] at (0.95,2.25) {$y'=(x'-2)^2$};
      \node[font=\footnotesize] at (-0.85,-1.45) {被动：曲线不动，坐标轴左移 $2$};
    \end{scope}
  \end{tikzpicture}
  \caption{主被动对偶：左图坐标系不动，抛物线 $y=x^2$ 主动右移 $2$ 得 $y=(x-2)^2$；右图抛物线原地不动，坐标轴被动左移 $2$（虚线为旧系），新读数 $x'=x+2$，顶点在新系下仍读作 $(2,0)$。两幅终态图完全相同——「动曲线」与「动坐标系」是同一平移的两种归属，不要混用方向}
  \label{fig:activepassive}
\end{figure}

% ---------- 收官微专题：$45^\circ$ 镜面反射定律 ----------

\begin{conccard}[二级]{唯有 $k=\pm1$ 的镜面，能保住「函数」的资格}{反函数对称、指对同构压轴选填、解析几何对称速算}{把斜线对称当成左右平移；漏掉 $b$ 的补偿项}
把对称轴从竖直线 $x=a$ 推广到一般直线 $y=kx+b$ 时，反射会让横、纵坐标\textbf{互相耦合}，反射后的曲线往往不再通过垂线检验，\key{跌出函数集合}，退回一般二元方程 $F(x,y)=0$。只有 $k=\pm1$（倾角 $45^\circ$ 或 $135^\circ$）时，耦合系数全部约尽，退化为干净的坐标互换——这正是反函数与指对同构压轴题的共同底座。
\end{conccard}

先把最典型的镜面 $l:y=x$ 看清楚（图 \ref{fig:expinvlog}）：反射把指数函数的点集 $\{(x,y)\}$ 整体翻成 $\{(y,x)\}$，对数曲线由此生成——这就是「反函数＝逆关系点集」的几何来源。

\begin{figure}[htbp]
  \centering
  \begin{tikzpicture}[x=0.9cm,y=0.9cm,>=stealth,font=\small]
    \draw[->,gray,thick] (-1.7,-1.7)--(3.6,-1.7) node[right] {$x$};
    \draw[->,gray,thick] (-1.7,-1.7)--(-1.7,3.8) node[above] {$y$};
    \draw[gray,dashed,thick,domain=-1.4:3.4,samples=40] plot(\x,\x) node[above right] {$y=x$};
    \draw[main,very thick,domain=-1.3:1.35,samples=80,smooth] plot(\x,{exp(\x)}) node[right] {$y=e^x$};
    \draw[second,very thick,domain=0.28:3.4,samples=90,smooth] plot(\x,{ln(\x)}) node[right] {$y=\ln x$};
    \fill[main] (0,1) circle (0.07);       \node[main,left] at (0,1) {$(0,1)$};
    \fill[main] (1,2.718) circle (0.07);   \node[main,above right] at (1,2.718) {$(1,e)$};
    \fill[main] (-1,0.368) circle (0.07);  \node[main,below left] at (-1,0.368) {$(-1,e^{-1})$};
    \fill[second] (1,0) circle (0.07);      \node[second,below] at (1,0) {$(1,0)$};
    \fill[second] (2.718,1) circle (0.07);  \node[second,right] at (2.718,1) {$(e,1)$};
    \fill[second] (0.368,-1) circle (0.07); \node[second,below left] at (0.368,-1) {$(e^{-1},-1)$};
    \draw[->,third,thick] (0,1) to[bend left=12] (1,0);
    \draw[->,third,thick] (1,2.718) to[bend left=18] (2.718,1);
    \draw[->,third,thick] (-1,0.368) to[bend right=12] (0.368,-1);
  \end{tikzpicture}
  \caption{同一坐标系中的互逆反射：$y=e^x$ 与 $y=\ln x$ 关于 $y=x$ 对称，每个点 $(a,b)$ 与它的像 $(b,a)$ 用跨越镜面的箭头一一相连，例如 $(0,1)\leftrightarrow(1,0)$、$(1,e)\leftrightarrow(e,1)$、$(-1,e^{-1})\leftrightarrow(e^{-1},-1)$。对数曲线就是指数曲线点集 $\{(x,y)\}$ 经 $y=x$ 翻转后的 $\{(y,x)\}$}
  \label{fig:expinvlog}
\end{figure}

\begin{method}
\textbf{$45^\circ$ 镜面反射替换速查表}（镜面记为 $l$，点 $P(x,y)$ 的像为 $P'(x',y')$）：
\begin{center}
\renewcommand{\arraystretch}{1.4}
\small
\begin{tabular}{|p{2.0cm}|p{3.6cm}|p{3.9cm}|p{4.2cm}|}
\hline
\textbf{镜面} & \textbf{点集映射} & \textbf{曲线替换口径} & \textbf{典型用途} \\
\hline
$l:y=x$ & $(x,y)\mapsto(y,\,x)$ & $F(y,\,x)=0$ & 反函数：$y=a^x$ 与 $y=\log_a x$ \\
\hline
$l:y=-x$ & $(x,y)\mapsto(-y,\,-x)$ & $F(-y,\,-x)=0$ & 双轴对称联动：$y=\dfrac1x$ \\
\hline
$l:y=x+b$ & $(x,y)\mapsto(y-b,\,x+b)$ & $F(y-b,\,x+b)=0$ & 平移型反函数、指对同构压轴 \\
\hline
$l:y=-x+b$ & $(x,y)\mapsto(b-y,\,b-x)$ & $F(b-y,\,b-x)=0$ & 斜率 $-1$ 直线族（$y=x$ 的垂线） \\
\hline
\end{tabular}
\end{center}
\vspace{0.3em}
\key{秒杀推论}：若两条曲线关于 $y=x$ 对称（如互为反函数），且被斜率为 $-1$ 的直线 $x+y=m$ 所截，则两交点 $A,B$ 必关于 $y=x$ 对称，中点 $\bigl(\tfrac m2,\tfrac m2\bigr)$ 落在 $y=x$ 上，于是
$$x_A+x_B=m,\qquad y_A+y_B=m .$$
\end{method}

\begin{figure}[htbp]
  \centering
  \begin{tikzpicture}[x=0.95cm,y=0.95cm,>=stealth,font=\small]
    \draw[->,gray,thick] (-1.2,-1.2)--(5.2,-1.2) node[right] {$x$};
    \draw[->,gray,thick] (-1.2,-1.2)--(-1.2,5.2) node[above] {$y$};
    \draw[main,dashed,very thick,domain=-1.1:3.9,samples=40] plot(\x,\x) node[above right] {$y=x$};
    \draw[second,very thick,domain=0.09:4.8,samples=90] plot(\x,{ln(\x)}) node[right] {$y=\ln x$};
    \draw[third,very thick,domain=-1:1.2,samples=70] plot(\x,{exp(\x)}) node[above] {$y=e^x$};
    \draw[red,very thick] (-1,5)--(5,-1) node[below right] {$x+y=4$};
    \fill[main] (2,2) circle (2.4pt) node[above left] {$M(2,2)$};
    \fill[red] (1.07,2.93) circle (2.4pt) node[above left] {$A(x_1,y_1)$};
    \fill[red] (2.93,1.07) circle (2.4pt) node[below right] {$B(y_1,x_1)$};
  \end{tikzpicture}
  \caption{反函数 $y=e^x$ 与 $y=\ln x$ 关于 $y=x$ 对称；斜率为 $-1$ 的直线 $x+y=4$ 与二者各交于 $A,B$，两点关于 $y=x$ 对称，中点 $M(2,2)$ 落在 $y=x$ 上，故 $x_1+x_2=4$}
  \label{fig:mirror45}
\end{figure}

\begin{example}[指对同构：不解方程，三步出答案]
设 $x_1$ 满足 $e^{x_1}+x_1=4$，$x_2$ 满足 $\ln x_2+x_2=4$，求 $x_1+x_2$。
\begin{solution}
两式分别改写为 $e^{x_1}=4-x_1$、$\ln x_2=4-x_2$。
于是 $A(x_1,e^{x_1})$ 是曲线 $y=e^x$ 与直线 $y=4-x$ 的交点，$B(x_2,\ln x_2)$ 是曲线 $y=\ln x$ 与同一直线的交点（图 \ref{fig:mirror45}）。
$y=e^x$ 与 $y=\ln x$ 互为反函数，关于 $y=x$ 对称；直线 $y=4-x$ 的斜率为 $-1$，与 $y=x$ 垂直，故 $A,B$ 关于 $y=x$ 对称，即 $x_2=e^{x_1}$。
所以 $x_1+x_2=x_1+e^{x_1}=\key{4}$。全程无需解超越方程。
\end{solution}
\end{example}

\begin{example}[平移型镜面：对称轴不是 $y=x$ 时]
求曲线 $C_1:y=e^x$ 关于直线 $l:y=x+1$ 对称的曲线 $C_2$ 的方程，并验证点 $P(1,e)$ 与其对称点 $P'$ 分别落在 $C_1,C_2$ 上。
\begin{solution}
用 $l:y=x+b$（$b=1$）的替换口径 $F(y-b,\,x+b)=0$，把 $C_1$ 的方程 $y=e^x$ 中的 $x$ 换成 $y-1$、$y$ 换成 $x+1$：
$$x+1=e^{\,y-1}\implies y=\ln(x+1)+1,$$
即 $C_2:y=\ln(x+1)+1$。
点 $P(1,e)$ 的对称点为 $P'=(y-b,\,x+b)=(e-1,\,2)$；代入 $C_2$：$\ln\bigl((e-1)+1\bigr)+1=\ln e+1=2$，恰好落在 $C_2$ 上。
（顺带一提：$x=0$ 时 $e^0=1=0+1$，两条曲线在镜面上的点 $(0,1)$ 处相遇。）
\end{solution}
\end{example}

\begin{commonerror}{把「$y=x+b$ 对称」误当成「左移 $b$」}
关于 $y=x+b$ 的对称\emph{不是}平移：它同时交换了横纵坐标的角色。若把 $C_1:y=e^x$ 关于 $y=x+1$ 的对称曲线误写成 $y=e^{x-1}$（左移思路），代入点 $P(1,e)$ 的对称点 $(e-1,2)$ 立即矛盾——正确的替换永远是 $x\leftrightarrow y-b$ 的\emph{互换}，而不是同向加减。
\end{commonerror}

\end{explore}

\begin{inclasstest}
\begin{enumerate}
  \item 用三视角（代数、几何、箭头图）表示关系 $R=\{(1,1),(2,4),(3,9)\}$，判断它是否为函数，并写出对应法则；\level{易}
  \item 判断点集 $P=\{(x,y)\mid x^2+y^2=1\}$ 是否为函数，用垂线检验说明理由；\level{易}
  \item 点集 $\{(x,y)\mid x^2+y^2\leq4\}$ 与 $\{(x,y)\mid x+y\leq3\}$ 的交集表示什么几何图形？它是否可能为某个函数的图象？说明理由；\level{中}
  \item 方程 $x^3=x$ 的解可看成哪两个函数图象交点的横坐标？求出所有解。\level{中}
\end{enumerate}
\end{inclasstest}

\begin{afterclass}
\subsection*{基础巩固}
\begin{enumerate}
  \item 写出 $A=\{1,2\}$、$B=\{3,4\}$ 的笛卡尔积 $A\times B$，并任取一个二元关系；
  \item 判断 $R=\{(1,2),(2,2),(3,2)\}$ 是否为函数，并说明理由。
\end{enumerate}

\subsection*{进阶提升}
\begin{enumerate}
  \item 画出 $R=\{(x,y)\mid y\geq x^2,\ x^2+y^2\leq4\}$ 表示的平面区域，并判断它是否为函数；
  \item 求不等式 $x^2<2x+3$ 的解集，并说明它对应两曲线怎样的位置关系。
\end{enumerate}

\subsection*{拓展探究}
\begin{enumerate}
  \item 查资料说明：为什么"对数能把乘除变成加减"？用 $y=a^x$ 与 $y=\log_a x$ 的互逆关系（$(x,y)\leftrightarrow(y,x)$）加以解释。
  \item （对应「实数的三种等价面孔」）设 $S=\{x\in\mathbb{Q}:x^2<2\}$。证明：$S$ 在 $\mathbb{Q}$ 中没有最小上界——要求对任意有理上界 $r$，构造出更小的有理上界；并由此说明曲线 $y=x^2-2$ 为什么能穿过有理数轴的孔洞。
  \item （表征与定义之别）某同学说：「$0.999\dots$ 和 $1$ 之间总差一点点，所以 $0.999\dots\neq1$。」请指出错在哪里，并用等比数列求和或方程法给出严格说明。
\end{enumerate}

\subsection*{教师版答案}
\begin{solution}
\textbf{基础 1：}$A\times B=\{(1,3),(1,4),(2,3),(2,4)\}$；任取一个子集即可，如 $R=\{(1,4),(2,3)\}$。
\end{solution}
\begin{solution}
\textbf{基础 2：}是函数。虽然三个 $x$ 都对应同一个 $y=2$，但每个 $x$ 对应的 $y$ 唯一（多对一是允许的），竖直线 $x=1,2,3$ 与点集各只有一个交点。
\end{solution}
\begin{solution}
\textbf{进阶 1：}区域是抛物线 $y=x^2$ 上方、圆盘 $x^2+y^2\leq4$ 内部（含边界）的公共部分，即圆盘被抛物线截出的"弓形"部分。它不是函数：取 $x=0$，$y\in[0,2]$，一个 $x$ 对应无穷多个 $y$（竖直线 $x=0$ 与区域有无穷多交点）。
\end{solution}
\begin{solution}
\textbf{进阶 2：}$x^2<2x+3\iff(x-3)(x+1)<0$，解集 $(-1,3)$。几何上，在 $(-1,3)$ 上抛物线 $y=x^2$ 位于直线 $y=2x+3$ 的下方。
\end{solution}
\begin{solution}
\textbf{拓展：}设 $x=a^{m}$，$y=a^{n}$，则 $m=\log_a x$，$n=\log_a y$。乘法 $xy=a^{m+n}$ 取对数得 $\log_a(xy)=m+n=\log_a x+\log_a y$，即乘法化为加法；除法 $x/y=a^{m-n}$ 同理化为减法。反函数关系体现在：指数函数把指数 $m$ 映为幂 $a^m$，对数函数把幂还原为指数，数对 $(x,y)$ 与 $(y,x)$ 关于 $y=x$ 对称。
\end{solution}
\begin{solution}
\textbf{拓展 2：}设 $r\in\mathbb{Q}$、$r>0$ 且 $r^2>2$（此时 $r$ 是上界），令 $t=\dfrac{r+\dfrac{2}{r}}{2}\in\mathbb{Q}$。由 $t-r=-\dfrac{r^2-2}{2r}<0$ 知 $t<r$，由 $t^2-2=\dfrac{(r^2-2)^2}{4r^2}>0$ 知 $t$ 仍是上界。故任何有理上界都不是最小，$S$ 在 $\mathbb{Q}$ 中无最小上界。几何上，$y=x^2-2$ 的零点恰是 $\sqrt2$；$\sqrt2\notin\mathbb{Q}$ 意味着这条连续曲线穿过数轴时，落点位置在有理数里没有任何对应——孔洞不仅存在，还被曲线「踩」了出来。
\end{solution}
\begin{solution}
\textbf{拓展 3：}设 $x=0.999\dots$，则 $10x=9.999\dots$，两式相减得 $9x=9$，即 $x=1$（方程法）；也可写成等比级数 $0.999\dots=0.9\bigl(1+0.1+0.01+\dots\bigr)=\dfrac{0.9}{1-0.1}=1$。「总差一点点」的错误在于：若 $1-x=\varepsilon>0$，则 $x\leq1-\varepsilon<1$，与 $x$ 的十进制表示各位都逼近 $1$ 矛盾——所谓「一点点」是一个比任何正实数都小的量，不存在。这也说明无限小数的含义要靠完备性才能定义清楚。
\end{solution}
\end{afterclass}
\reading{拓展阅读 07}{函数的结构特化}

% 高观点导学案：三角函数、解三角形与复平面旋转（第 2 章·高观点统一版）
% 主线：以「乘以 e^{iβ} = 逆时针旋转 β」为唯一母公理，让整套三角恒等变换从指数运算律自然生长
% 四模块：数系扩张与 i 的几何诞生 → 结构同构（旋转必为指数函数）→ 切线速度锁定底数 → 母公式引擎
% 另配一节：解三角形与边角互化。全程用「复平面旋转 ↔ 实虚部投影」互译，原生 TikZ 图可视化。

\section{三角函数、解三角形与复平面旋转（高观点统一版）}

\begin{methodbox}{多框架对读}
\small
\begin{center}
\begin{tabular}{c|p{2.9cm}|p{5.0cm}|p{3.1cm}}
\toprule
视角 & 本章看点 & 别的视角怎么说 & 哪侧最省力 \\
\midrule
代数 & 和差倍半公式、恒等变形 & 符号系统的机械推导 & 硬算与验算 \\
几何 & 单位圆上的旋转与投影 & 角度相加＝复数相乘 & 猜公式、看符号 \\
集合 & $\theta\mapsto(\cos\theta,\sin\theta)$ & 参数方程即一个对应 & 界定义域与值域 \\
映射 & 乘 $e^{i\beta}$＝逆时针转 $\beta$ & 旋转群与指数函数同构 & 统一推导所有公式 \\
\bottomrule
\end{tabular}
\end{center}
卡的使用：公式卡壳时回到单位圆「转一遍」；证明题优先用旋转映射（总图见第 0 章图~\ref{fig:four-view-map}）。
\end{methodbox}

本章不新增公式，而是给三角函数一章装上一个\key{母引擎}：
在复平面上，\key{乘以 $e^{i\beta}$ 就等于把向量逆时针旋转 $\beta$ 角}（模长不变）。
由这一条母公理出发，两角和差、二倍角三倍角、诱导公式、和差化积与积化和差，
全部只是指数运算律 $e^{i(\alpha+\beta)}=e^{i\alpha}e^{i\beta}$ 的代数推论。
四个模块环环相扣：
\key{数系扩张与 $i$ 的几何诞生}（为什么需要虚数）$\to$
\key{结构同构}（单位圆旋转必为指数函数）$\to$
\key{切线速度锁定底数}（欧拉公式 $e^{i\theta}=\cos\theta+i\sin\theta$）$\to$
\key{母公式引擎}（一条指数律长出全套公式），
最后用 \key{解三角形与边角互化} 把复数直觉落回高考主战场。
核心思想：\key{三角恒等变换不是几十个孤立公式，而是同一个旋转群的三种语言}（指数、坐标、三角形）。

\begin{learnobjective}
\begin{itemize}
  \item 理解虚数单位 $i$ 的几何意义（逆时针旋转 $90^\circ$），知道 $i^2=-1$ 是「转两次直角 $=$ 掉头」的代数翻译；
  \item 理解数系扩张 $\mathbb{N}\to\mathbb{Z}\to\mathbb{Q}\to\mathbb{R}\to\mathbb{C}$ 每一步都在为运算封闭与几何变换解锁新维度；
  \item 会用单位复数 $z(\theta)=\cos\theta+i\sin\theta$ 验证 $z(\alpha)z(\beta)=z(\alpha+\beta)$，理解「加法变乘法」是指数函数的基因；
  \item 会用切线速度 $(e^{k\theta})'\big|_{\theta=0}=k$ 锁定底数，写出欧拉公式 $e^{i\theta}=\cos\theta+i\sin\theta$，并得出 $e^{i\pi}+1=0$；
  \item 会用 $e^{i(\alpha+\beta)}$ 展开、棣莫弗定理与指数分解，现场推导两角和差、二倍角三倍角与和差化积／积化和差；
  \item 掌握正弦定理（外接圆半径 $R$）、余弦定理与面积公式，会按「一次齐次用正弦、二次齐次用余弦」做边角互化，并识别钝角三角形中的隐藏定义域。
\end{itemize}
\end{learnobjective}

\begin{preview}
\subsection*{知识回顾}
\begin{itemize}
  \item 实数轴上，乘以正数 $2$ 的几何意义是 \blank（拉伸／旋转）；乘以 $-1$ 的几何意义是 \blank；
  \item $(-1)\times(-1)=$ \blank；
  \item 指数函数 $a^x$ 满足的运算律是 $a^{\alpha+\beta}=$ \blank；
  \item 单位圆上辐角 $\theta$ 对应的点是 $(\cos\theta,\sin\theta)$，它到原点的距离是 \blank。
\end{itemize}

\subsection*{课前速读}
阅读三角函数一章及本章导学案，完成下列问题：
\begin{enumerate}
  \item $i^2=$ \blank；
  \item 欧拉公式 $e^{i\theta}=$ \blank；
  \item 正弦定理 $\dfrac{a}{\sin A}=$ \blank（用外接圆半径 $R$ 表示）；
  \item 余弦定理 $a^2=$ \blank。
\end{enumerate}
\end{preview}

\begin{explore}

\subsection{数系扩张与虚数单位 $i$ 的几何诞生}

\key{探究一：乘法即变换——从一维缩放、掉头到二维旋转}
在实数轴上盯住数字 $1$，看对它做乘法时的几何运动：
\begin{itemize}
  \item 乘以正数（如 $\times2$）：线段被\key{拉伸}为原来的 $2$ 倍，方向不变——纯\key{缩放}；
  \item 乘以 $-1$：把 $1$ 变成 $-1$，几何意义是绕原点逆时针旋转 \key{$180^\circ$}（半圈）；
  \item 连续乘两次 $-1$：$(-1)\times(-1)=1$，即转 $180^\circ+180^\circ=360^\circ$，恰好回到原位。
\end{itemize}
于是提出关键问题：\key{有没有一种操作，连续做两次恰好等于乘 $-1$（即一次只逆时针转 $90^\circ$）？}
设这个操作是 $i$，则
\[
  i\times i=-1\iff i^2=-1.
\]
一维数轴上根本容不下「转 $90^\circ$」后的落点，数轴被迫向垂直方向张开，出现\key{虚轴}，
一维直线跃升为二维\key{复平面}。

\begin{figure}[htbp]
  \centering
  \begin{tikzpicture}[>=Stealth, semithick, font=\footnotesize]
    % 左图：乘 -1 转 180°
    \draw[->, gray] (-2.5,0)--(2.5,0) node[below right] {实轴};
    \fill (0,0) circle (1pt);
    \draw[->, main!80!black, thick] (0,0)--(1,0) node[midway,above] {$1$};
    \draw[second, thick] (1,0) arc (0:180:1);
    \draw[->, second, thick] (0,0)--(-1,0) node[midway,above] {$-1$};
    \node[second] at (0,1.35) {$180^\circ$};
    \node[main!70!black] at (0,-2.05) {乘 $-1$};
    % 右图：乘 i 转 90°，张开虚轴
    \begin{scope}[xshift=6.2cm]
      \draw[->, gray] (-1.4,0)--(2.5,0) node[below right] {实轴};
      \draw[->, dashed, gray!70!black] (0,-1.7)--(0,2.3) node[above] {虚轴};
      \fill (0,0) circle (1pt);
      \draw[->, main!80!black, thick] (0,0)--(1,0) node[midway,below] {$1$};
      \draw[third, thick] (1,0) arc (0:90:1);
      \draw[->, third, thick] (0,0)--(0,1) node[midway,right] {$i$};
      \draw[dashed, gray] (0,1) arc (90:180:1);
      \draw[->, dashed, gray] (0,0)--(-1,0) node[midway,below] {$i^2=-1$};
      \node[third] at (0.62,0.30) {$90^\circ$};
      % 说明文字放在虚轴下端之下（虚轴只画到 $-1.7$），避免被竖直虚线穿过
      \node[main!70!black] at (0,-2.05) {乘 $i$（连乘两次得 $-1$）};
    \end{scope}
  \end{tikzpicture}
  \caption{乘法的几何：乘 $-1$ 逆时针转 $180^\circ$；乘 $i$ 只转 $90^\circ$，一维数轴被迫张开虚轴，成为复平面}
  \label{fig:rotate-i}
\end{figure}

点睛：实数 $1$ 负责向前走，$-1$ 负责向后转，$i$ 负责向左转；
连喊两次「向左转」正好就是「向后转」（$i^2=-1$）。

乘 $i$ 的作用不止于「转两次得 $-1$」：对任意辐角 $\theta$ 的单位向量，乘 $i$ 都把它的两个分量整体倒置（图~\ref{fig:i-component-flip}）。
\begin{figure}[htbp]
  \centering
  \begin{tikzpicture}[>=Stealth, semithick, font=\footnotesize]
    % ===== 左面板：转前 z=cosθ+isinθ =====
    \draw[->, gray] (-1.9,0)--(1.9,0) node[below right] {实轴};
    \draw[->, gray] (0,-0.9)--(0,1.95) node[above] {虚轴};
    \draw[gray!55, dashed] (0,0) circle (1.6);
    \draw[black!70, thick] (0,0)--(40:1.6) node[above right, font=\scriptsize] {$z$};
    \draw[blue!70!black, very thick] (0,0)--({1.6*cos(40)},0);
    \node[blue!70!black, below] at ({0.8*cos(40)},0) {$\cos\theta$};
    \draw[green!55!black, very thick] ({1.6*cos(40)},0)--(40:1.6);
    \node[green!55!black, left] at ({1.6*cos(40)+0.02},{0.8*sin(40)}) {$\sin\theta$};
    % 直角符号：两个投影分量互相垂直
    \draw[black!60, thin] ({1.6*cos(40)-0.22},0)--({1.6*cos(40)-0.22},0.22)--({1.6*cos(40)},0.22);
    \draw[black!70, thick] (0.42,0) arc (0:40:0.42) node[midway, above, font=\scriptsize] {$\theta$};
    \node[main!70!black] at (0,-0.7) {$z=\cos\theta+i\sin\theta$};
    % ===== 右面板：乘 i 后 iz=-sinθ+icosθ =====
    \begin{scope}[xshift=4.4cm]
      \draw[->, gray] (-1.9,0)--(1.9,0) node[below right] {实轴};
      \draw[->, gray] (0,-0.9)--(0,1.95) node[above] {虚轴};
      \draw[gray!55, dashed] (0,0) circle (1.6);
      \draw[dashed, gray!70!black, thick] (0,0)--(40:1.6);
      % 直角符号：z 与 iz 垂直（随 θ 取向）
      \draw[third, thin] ({0.24*cos(40)},{0.24*sin(40)})--({0.24*cos(40)-0.24*sin(40)},{0.24*sin(40)+0.24*cos(40)})--({-0.24*sin(40)},{0.24*cos(40)});
      % 90° 弧：从 θ 扫到 θ+90°
      \draw[->, third, thick] (40:0.5) arc (40:130:0.5);
      \node[third, left] at ({0.72*cos(95)},{0.72*sin(95)}) {$90^\circ$};
      \draw[->, second, very thick] (0,0)--({1.6*cos(130)},{1.6*sin(130)}) node[above left, font=\scriptsize] {$iz$};
      % 倒置后的两个分量：蓝竖（cosθ）、绿横（-sinθ）
      \draw[blue!70!black, very thick] (0,0)--(0,{1.6*cos(40)});
      \node[blue!70!black, right] at (0,{0.8*cos(40)}) {$\cos\theta$};
      \draw[green!55!black, very thick] (0,{1.6*cos(40)})--({-1.6*sin(40)},{1.6*cos(40)});
      \node[green!55!black, below] at ({-0.8*sin(40)},{1.6*cos(40)}) {$-\sin\theta$};
      \draw[dashed, gray] ({1.6*cos(130)},{1.6*sin(130)})--({1.6*cos(130)},0);
      \node[main!70!black] at (0,-0.7) {$iz=-\sin\theta+i\cos\theta$};
    \end{scope}
  \end{tikzpicture}
  \caption{乘 $i$ 的分量倒置：$z=\cos\theta+i\sin\theta$（横蓝 $\cos\theta$、纵绿 $\sin\theta$）乘 $i=e^{i\pi/2}$ 逆时针转 $90^\circ$ 后，蓝分量竖到虚轴成 $\cos\theta$、绿分量倒向负实轴成 $-\sin\theta$，即 $iz=-\sin\theta+i\cos\theta$}
  \label{fig:i-component-flip}
\end{figure}

\begin{example}
三次方程 $x^3=15x+4$ 显然有实根 $x=4$。用卡丹公式却会得到
\[
  x=\sqrt[3]{2+\sqrt{-121}}+\sqrt[3]{2-\sqrt{-121}}
  =\sqrt[3]{2+11i}+\sqrt[3]{2-11i}.
\]
请验证 $(2+i)^3=2+11i$，并说明虚部如何抵消得到 $x=4$。
\end{example}
\begin{solution}
\[
  (2+i)^3=8+12i+6i^2+i^3=8+12i-6-i=2+11i,
\]
同理 $(2-i)^3=2-11i$。于是
\[
  x=(2+i)+(2-i)=4,
\]
两个虚部一正一负抵消，还原出真实的实数解。
可见虚数不是凭空捏造的幻想，而是连接实数问题与实数解的「地下通道」：
引入 $i$ 后，$n$ 次代数方程必有 $n$ 个根（代数基本定理），数系运算彻底封闭。
\end{solution}

\begin{center}
\footnotesize
\setlength{\tabcolsep}{3pt}
\renewcommand{\arraystretch}{1.4}
\begin{tabular}{p{2.8cm}p{3.4cm}p{5.4cm}}
\toprule
扩张阶段 & 解决的方程障碍 & 解锁的几何空间与变换 \\
\midrule
$\mathbb{N}\to\mathbb{Z}$（引入负数） & $x+3=1$ 无解 & 数轴向负半轴延伸，解锁反向平移与 $180^\circ$ 掉头 \\
$\mathbb{Z}\to\mathbb{Q}$（引入分数） & $2x=1$ 无解 & 填满离散整数点之间的空隙，解锁任意比例缩放 \\
$\mathbb{Q}\to\mathbb{R}$（引入无理数） & $x^2=2$ 无解 & 消除数轴上的「针眼漏洞」，得到连续完备实数线 \\
$\mathbb{R}\to\mathbb{C}$（引入虚数 $i$） & $x^2+1=0$ 无解 & 一维直线跃升为二维平面，解锁任意角度旋转与代数封闭 \\
\bottomrule
\end{tabular}
\end{center}

\begin{conccard}[一级]{虚数单位 $i$ 的几何意义：$i^2=-1$}{理解复数的几何来源、判断复数乘法的旋转角与伸缩倍数}{把 $i$ 只当成「$\sqrt{-1}$」的代数符号，丢掉它的旋转几何；把 $i^2=-1$ 与「转 $180^\circ$」脱钩}
$1\times i$ 表示逆时针转 $90^\circ$，$1\times(-1)$ 表示转 $180^\circ$，而 $i\times i=-1$ 正是「转两次直角等于掉头」。
任何复数 $z$ 乘以 $i$，都是把 $z$ 的辐角加 $90^\circ$、模长不变。
\end{conccard}

\key{再换一个角度：每次数系扩充，保留了什么、失去了什么}
数系扩充不是「只增不减」地加数：每一步都在尽量保留旧运算律的前提下换来解题能力，同时不可避免地牺牲掉某些结构性性质。下表是四次扩充的准确清单（表中「失去」指该性质在新数系中不再整体成立）：

\begin{center}
\footnotesize
\setlength{\tabcolsep}{4pt}
\renewcommand{\arraystretch}{1.5}
\begin{tabular}{p{2.6cm}p{4.6cm}p{4.4cm}}
\toprule
扩充过程 & 保留的性质 & 不能继续保留的性质 \\
\midrule
$\mathbb{N}\to\mathbb{Z}$ & 加法、乘法及其运算律；相反数的出现使减法封闭 & 良序性（$\mathbb{N}$ 的任意非空子集有最小元；$\mathbb{Z}$ 中负整数全体就没有最小元） \\
$\mathbb{Z}\to\mathbb{Q}$ & 四则运算律与大小顺序；除法封闭、非零元都有逆元 & 离散性（不存在相邻的两个有理数，整数之间不再有空隙） \\
$\mathbb{Q}\to\mathbb{R}$ & 域结构与全序；开方、长度、极限有了合法舞台 & 可数性（$\mathbb{Q}$ 可逐个列出，$\mathbb{R}$ 不能——对角线法） \\
$\mathbb{R}\to\mathbb{C}$ & 加、减、乘、除四则运算律；代数封闭 & 与四则运算相容的大小顺序（不能像实数那样比大小） \\
\bottomrule
\end{tabular}
\end{center}

值得强调的是：前三次扩充都保留甚至加强了「能比大小」的能力，只有最后一步必须放弃它。这不是缺陷，而是换取代数封闭性的代价。

\key{为什么复数不能比较大小}
假设在 $\mathbb{C}$ 上存在一种全序 $<$，满足两条与运算相容的常规要求：($i$) $a<b\Rightarrow a+c<b+c$；($ii$) $a>0,\ b>0\Rightarrow ab>0$（于是非零数的平方必为正）。
\begin{derivation}
$i\neq0$，所以要么 $i>0$，要么 $i<0$。
\begin{itemize}
  \item 若 $i>0$：由 ($ii$) 得 $i^2>0$，但 $i^2=-1<0$，矛盾。
  \item 若 $i<0$：由 ($i$) 两边加 $-i$ 得 $0<-i$，即 $-i>0$；再由 ($ii$) 得 $(-i)^2>0$，而 $(-i)^2=i^2=-1<0$，同样矛盾。
\end{itemize}
两种情形都矛盾，故这样的全序不存在。
\end{derivation}
因此只能说「$z_1$ 与 $z_2$ 不相等」，不能说「$z_1$ 比 $z_2$ 大」；比较两个复数的唯一合法方式是比\key{模}（$|z|$ 是实数，可以比大小）或比\key{辐角}（方向）。这与实数数轴上的「谁左谁右」是完全不同的两回事。

\begin{commonerror}{把历史的顺序当成逻辑的唯一顺序}
常见讲法：「人类为了描述旋转 $90^\circ$ 才发明了 $i$。」这在数学史上不准确：虚数最早来自三次方程求根时出现的负数开平方（卡尔达诺《大术》1545 年已有记载，邦贝利随后系统研究），复平面的几何解释（把复数看成平面上的点、把乘法看成旋转伸缩）是更晚才发展出来的认识。
更严谨的教学表述是：人类先为突破实数范围内的代数求解限制而引入虚数，后来才发现它还赋予复数乘法漂亮的几何意义。
本节用「乘法即变换」作为引入路径，是\key{逻辑在先、直观优先}的教学设计，不是历史事实的顺序；讲课时应明确区分，避免学生形成「复数是为了画图方便而被发明」的错误印象。
\end{commonerror}

\begin{method}
判断一个数系够不够用的三个反射：
\begin{itemize}
  \item 方程层面：目标方程在当前数系内无解吗？（$x+3=1,\ 2x=1,\ x^2=2,\ x^2+1=0$ 各逼出一次扩充。）
  \item 运算层面：当前数系对四则运算封闭吗？对极限运算封闭吗？
  \item 结构层面：为了求解新问题，愿意放弃哪条旧性质？（序、离散性、可数性……）
\end{itemize}
\end{method}

\key{课堂讨论：数系扩充有没有终点}
\begin{itemize}
  \item \key{问题 1}：实数已经能表示数轴上的每一个点，为什么还不够用？
        （引导：完备性解决的是「连续」；方程 $x^2+1=0$ 需要的是「代数封闭」，两者是不同的要求。）
  \item \key{问题 2}：为什么复数不能比较大小？请用平方运算自己推出矛盾。
  \item \key{问题 3}：不能比大小，是不是说明复数比实数「更差」？
        （引导：失去序，换来代数封闭与二维几何；交流电路、振动、信号处理都靠复数把运算变简洁。）
\end{itemize}
给后续章节留一个伏笔：从 $\mathbb{R}$ 到 $\mathbb{C}$ 也改变了对函数的理解——实函数依赖大小顺序与数轴结构，复函数则可以借助复平面与几何变换建立新的理论。

\subsection{结构同构：单位圆旋转必为指数函数}

\key{探究二：角度相加，复数相乘}
把单位圆上辐角为 $\theta$ 的点写成复数
\[
  z(\theta)=\cos\theta+i\sin\theta.
\]
计算两个单位复数的乘积：
\begin{derivation}
\begin{align*}
  z(\alpha)z(\beta)&=(\cos\alpha+i\sin\alpha)(\cos\beta+i\sin\beta)\\
  &=(\cos\alpha\cos\beta-\sin\alpha\sin\beta)+i(\sin\alpha\cos\beta+\cos\alpha\sin\beta)\\
  &=\cos(\alpha+\beta)+i\sin(\alpha+\beta)=z(\alpha+\beta).
\end{align*}
\end{derivation}
（最后一步用了两角和公式——这正是我们稍后要用旋转重新解释的对象。）
惊人之处：\key{输入的角度相加，输出的复数相乘}。

唤醒旧知：哪类函数的核心基因是 $f(\alpha+\beta)=f(\alpha)f(\beta)$？
学生应脱口而出\key{指数函数} $f(x)=b^x$。
而 $z(0)=\cos0+i\sin0=1$，所以 $z(\theta)$ 在代数结构上就是一个底数为某复常数 $b$ 的指数函数 $z(\theta)=b^\theta$——
\key{单位圆旋转群与指数函数同构}。

\begin{conccard}[二级]{结构同构：$z(\alpha)z(\beta)=z(\alpha+\beta)$}{判断两个复数相乘的几何效果、理解三角公式的共同来源}{只记「模相乘、辐角相加」的结论，不从结构公理出发；把 $z(\theta)$ 与指数函数割裂}
单位复数 $z(\theta)=\cos\theta+i\sin\theta$ 满足「角度相加 $\to$ 复数相乘」，且 $z(0)=1$。
凡是满足 $f(\alpha+\beta)=f(\alpha)f(\beta)$ 的连续函数必为指数函数，故 $z(\theta)$ 一定可以写成 $b^\theta$ 的形式。
\end{conccard}

几何上看，$z(\alpha)\cdot e^{i\beta}$ 就是把 $z_1=z(\alpha)$ 绕原点再逆时针转 $\beta$，终点辐角为 $\alpha+\beta$（图~\ref{fig:rotate-sum}）。
\begin{figure}[htbp]
  \centering
  \begin{tikzpicture}[>=Stealth, semithick, font=\footnotesize]
    \draw[->, gray] (-1.8,0)--(3.5,0) node[below right] {$\mathrm{Re}$};
    \draw[->, gray] (0,-1.7)--(0,2.9) node[above] {$\mathrm{Im}$};
    \draw[gray!55, dashed] (0,0) circle (2.2);
    % 三段同心弧：α、β、α+β（半径与标注角错层，避免相叠）
    \draw[main, thick] (0.5,0) arc (0:20:0.5);
    \node[main] at (10:0.68) {$\alpha$};
    \draw[->, third, thick] (20:0.85) arc (20:75:0.85);
    \node[third] at (47:1.06) {$\beta$};
    \draw[->, gray!60!black, thick] (1.2,0) arc (0:75:1.2);
    \node[gray!60!black] at (12:1.45) {$\alpha+\beta$};
    % 旋转轨迹：z1 沿单位圆转到 z2
    \draw[->, third, thin] (20:2.2) arc (20:75:2.2);
    % z1=e^{iα} 与 z2=e^{i(α+β)}
    \draw[->, main!80!black, thick] (0,0)--(20:2.2) node[above right] {$z_1=e^{i\alpha}$};
    \draw[->, second, very thick] (0,0)--(75:2.2) node[above left] {$z_2=e^{i(\alpha+\beta)}$};
    % z2 的两轴投影：实部红、虚部青
    \draw[red!80!black, very thick] (0,0)--({2.2*cos(75)},0);
    \node[red!80!black, below] at ({1.1*cos(75)},0) {$\cos(\alpha+\beta)$};
    \draw[dashed, gray] (75:2.2)--(0,{2.2*sin(75)});
    \draw[cyan!80!black, very thick] ({2.2*cos(75)},0)--(75:2.2);
    \node[cyan!80!black, right] at ({2.2*cos(75)},1.2) {$\sin(\alpha+\beta)$};
    \node[main!70!black] at (0.9,-1.25) {$z_1\cdot e^{i\beta}=e^{i(\alpha+\beta)}=z_2$};
  \end{tikzpicture}
  \caption{旋转叠加：$z_1=e^{i\alpha}$ 再乘 $e^{i\beta}$ 即绕原点逆时针转 $\beta$，终点 $z_2=e^{i(\alpha+\beta)}$；$z_2$ 在实轴、虚轴上的投影 $\cos(\alpha+\beta)$、$\sin(\alpha+\beta)$ 与 $(\cos\alpha+i\sin\alpha)(\cos\beta+i\sin\beta)$ 的实部、虚部一一对应——两角和公式就是这么「转」出来的}
  \label{fig:rotate-sum}
\end{figure}

\subsection{切线速度锁定底数}

\key{探究三：切线速度——底数只能是 $e$}
对比两个「指数运动」的动力学：
\begin{itemize}
  \item 实指数 $e^{x}$：在 $x=0$ 处位置为 $1$，速度方向与位置\key{同向}（沿实轴向右），越跑越快，轨迹是直线爆炸；
  \item 复指数 $e^{i\theta}$：在 $\theta=0$ 处位置为 $1$，速度方向必须与位置\key{垂直}（沿虚轴向上，即 $+i$），只转弯不伸长，轨迹是单位圆。
\end{itemize}
把 $z(\theta)$ 看作质点位置，求导得切线速度：
\[
  z'(\theta)=-\sin\theta+i\cos\theta\implies z'(0)=i.
\]
而实指数 $f(\theta)=e^{k\theta}$ 满足
\[
  (e^{k\theta})'\Big|_{\theta=0}=k\cdot e^{0}=k.
\]
要让圆周运动在出发点的瞬时速度恰好是 $+i$，唯一可能就是 \key{$k=i$}。于是
\[
  e^{i\theta}=\cos\theta+i\sin\theta.
\]
代入 $\theta=\pi$：
\[
  e^{i\pi}=\cos\pi+i\sin\pi=-1\implies e^{i\pi}+1=0,
\]
即被誉为「上帝公式」的欧拉恒等式，把 $0,1,e,i,\pi$ 五个常数用最纯粹的圆周旋转焊在一起。

\begin{center}
\footnotesize
\setlength{\tabcolsep}{3pt}
\renewcommand{\arraystretch}{1.4}
\begin{tabular}{p{2.2cm}p{2.0cm}p{2.6cm}p{3.2cm}p{3.0cm}}
\toprule
函数 & 起点位置 & 初始速度方向 & 速度与位置关系 & 轨迹 \\
\midrule
$e^{x}$ & $1$ & $+1$（实轴向右） & 同向 & 沿实轴指数爆炸 \\
$e^{i\theta}$ & $1$ & $+i$（虚轴向上） & 垂直 & 绕单位圆匀速旋转 \\
\bottomrule
\end{tabular}
\end{center}

\begin{figure}[htbp]
  \centering
  \begin{tikzpicture}[>=Stealth, semithick, font=\footnotesize]
    \draw[->, gray] (-2.6,0)--(2.9,0) node[below right] {$\mathrm{Re}$};
    \draw[->, gray] (0,-2.3)--(0,2.5) node[above] {$\mathrm{Im}$};
    \draw[gray!55, dashed] (0,0) circle (2);
    % 一般位置 z(θ)：θ≈50°
    \draw[->, dashed, gray!80!black] (0,0)--(1.286,1.532) node[above] {$z(\theta)$};
    \draw[->, dashed, gray!80!black] (1.286,1.532)--(0.826,1.918) node[above left, font=\scriptsize] {$i\,z(\theta)$};
    % 起点 (1,0) 与切线速度 +i
    \fill (2,0) circle (1.2pt);
    \draw[->, main!80!black, thick] (0,0)--(2,0) node[midway,below] {$z(0)=1$};
    \draw[->, second, thick] (2,0)--(2,0.95) node[right] {$z'(0)=i$};
    \draw[gray, thin] (1.52,0)--(1.52,0.48)--(2,0.48);
    \draw[gray] (0.72,0) arc (0:50:0.72);
    \node[gray!80!black] at (0.98,0.22) {$\theta$};
    \node[main!70!black] at (0,-2.8) {位置向量与切线速度垂直：只转弯，不伸长};
  \end{tikzpicture}
  \caption{复平面单位圆：起点 $(1,0)$ 处位置向量为 $1$，切线速度 $+i$ 与其垂直；任意位置的速度 $i\,z(\theta)$ 始终垂直于位置}
  \label{fig:euler-tangent}
\end{figure}

\begin{conccard}[一级]{欧拉公式：$e^{i\theta}=\cos\theta+i\sin\theta$}{推导三角公式、判断复数旋转角、快速验算}{只把它当结论背诵；忘记「乘以 $e^{i\beta}$ 就是逆时针转 $\beta$」这一几何含义}
由 $z(\alpha)z(\beta)=z(\alpha+\beta)$ 与 $z(0)=1$ 知 $z(\theta)$ 必为指数函数 $b^\theta$；
再由切线速度 $(e^{k\theta})'\big|_{\theta=0}=k$ 与 $z'(0)=i$ 锁定 $k=i$。
于是乘以 $e^{i\beta}$ 等价于逆时针旋转 $\beta$，这就是本章全部三角公式的母公理。
\end{conccard}

\subsection{母公式引擎：一条指数律长出全套公式}

\key{探究四：$e^{i(\alpha+\beta)}=e^{i\alpha}e^{i\beta}$ 的展开与投影}

\key{（1）诱导公式——乘 $i=e^{i\pi/2}$ 即转 $90^\circ$。}
\[
  i(\cos\theta+i\sin\theta)=-\sin\theta+i\cos\theta,
\]
对比实部虚部，一眼读出
\[
  \cos\left(\theta+\frac{\pi}{2}\right)=-\sin\theta,\qquad
  \sin\left(\theta+\frac{\pi}{2}\right)=\cos\theta.
\]
同理乘 $-1=e^{i\pi}$ 得 $\cos(\theta+\pi)=-\cos\theta$、$\sin(\theta+\pi)=-\sin\theta$；
取共轭 $\overline{e^{i\theta}}=e^{-i\theta}$ 即关于实轴对称，这正是奇偶性
$\cos(-\theta)=\cos\theta$、$\sin(-\theta)=-\sin\theta$ 的几何本质。
「奇变偶不变，符号看象限」不再是口诀，而是「转了几个直角」的图形事实。

\key{（2）两角和与差——多项式展开一步到位。}
\begin{derivation}
\begin{align*}
  e^{i(\alpha+\beta)}&=e^{i\alpha}e^{i\beta}
  =(\cos\alpha+i\sin\alpha)(\cos\beta+i\sin\beta)\\
  &=(\cos\alpha\cos\beta-\sin\alpha\sin\beta)
  +i(\sin\alpha\cos\beta+\cos\alpha\sin\beta).
\end{align*}
\end{derivation}
左边又是 $\cos(\alpha+\beta)+i\sin(\alpha+\beta)$，对比实部与虚部：
\[
  \cos(\alpha+\beta)=\cos\alpha\cos\beta-\sin\alpha\sin\beta,\qquad
  \sin(\alpha+\beta)=\sin\alpha\cos\beta+\cos\alpha\sin\beta.
\]
把 $\beta$ 换成 $-\beta$（用 $\cos(-\beta)=\cos\beta$、$\sin(-\beta)=-\sin\beta$），立刻得到差角公式：
\[
  \cos(\alpha-\beta)=\cos\alpha\cos\beta+\sin\alpha\sin\beta,\qquad
  \sin(\alpha-\beta)=\sin\alpha\cos\beta-\cos\alpha\sin\beta.
\]

\key{（3）二倍角与三倍角——棣莫弗定理。}
由 $e^{in\theta}=(e^{i\theta})^n$ 得
\[
  (\cos\theta+i\sin\theta)^n=\cos n\theta+i\sin n\theta.
\]
$n=2$ 展开即得二倍角；$n=3$ 用二项式定理：
\[
  (\cos\theta+i\sin\theta)^3
  =(\cos^3\theta-3\cos\theta\sin^2\theta
  +i(3\cos^2\theta\sin\theta-\sin^3\theta),
\]
于是
\[
  \cos3\theta=4\cos^3\theta-3\cos\theta,\qquad
  \sin3\theta=3\sin\theta-4\sin^3\theta.
\]

\key{（4）和差化积与积化和差——同一指数分解。}
由欧拉公式解出
\[
  \cos\theta=\frac{e^{i\theta}+e^{-i\theta}}{2},\qquad
  \sin\theta=\frac{e^{i\theta}-e^{-i\theta}}{2i}.
\]
例如
\begin{align*}
  \cos\alpha\cos\beta
  &=\frac{(e^{i\alpha}+e^{-i\alpha})(e^{i\beta}+e^{-i\beta})}{4}\\
  &=\frac{e^{i(\alpha+\beta)}+e^{-i(\alpha+\beta)}+e^{i(\alpha-\beta)}+e^{-i(\alpha-\beta)}}{4}\\
  &=\frac{\cos(\alpha+\beta)+\cos(\alpha-\beta)}{2},
\end{align*}
这就是积化和差；把 $\alpha+\beta=A$、$\alpha-\beta=B$ 反解即得和差化积。
结论：和差、乘积、倍角只是同一组指数幂的不同代数外形。

\begin{example}
用棣莫弗定理把 $\cos3\theta$ 用 $\cos\theta$ 表示。
\end{example}
\begin{solution}
\[
  \cos3\theta=\operatorname{Re}(\cos\theta+i\sin\theta)^3
  =\cos^3\theta-3\cos\theta\sin^2\theta
  =\cos^3\theta-3\cos\theta(1-\cos^2\theta)
  =4\cos^3\theta-3\cos\theta.
\]
\end{solution}

\begin{example}
求 $\cos\theta+\cos2\theta+\cdots+\cos n\theta$（$\theta\neq2k\pi$）。
\end{example}
\begin{solution}
令 $q=e^{i\theta}$，则 $\cos k\theta=\operatorname{Re}q^k$，于是
\[
  S=\sum_{k=1}^{n}q^k=\frac{q(1-q^n)}{1-q}
  =\frac{e^{i\theta}(1-e^{in\theta})}{1-e^{i\theta}}
  =e^{i\frac{(n+1)\theta}{2}}\cdot\frac{\sin\frac{n\theta}{2}}{\sin\frac{\theta}{2}}.
\]
取实部即得
\[
  \cos\theta+\cos2\theta+\cdots+\cos n\theta
  =\frac{\sin\frac{n\theta}{2}}{\sin\frac{\theta}{2}}\cos\frac{(n+1)\theta}{2}.
\]
实数域里要凑 $\sin\frac{\theta}{2}$ 裂项，技巧性强；复数域里只是公比为 $q=e^{i\theta}$ 的等比数列求和。
\end{solution}

\begin{conccard}[一级]{考纲安全边界：复数法用于选填秒杀，解答题须还原为实数格式}{选择填空快速求值、验算；忘记公式时现场推导；快速判断符号}{在解答题中直接写 $e^{i\theta}$、$i$ 的推导过程而跳过实数三角变换步骤，造成阅卷扣分}
高考解答题默认在实数域 $\mathbb{R}$ 内评分。复数法与欧拉公式的正确用法是：
选填小题秒杀、快速验算、忘记公式时现场推导。
解答题书写必须还原为实数三角变换格式（和差倍半公式、正弦定理、余弦定理的标准步骤），
不能把 $e^{i\theta}$ 与 $i$ 直接写进解答题过程。
\end{conccard}

\begin{commonerror}{指数展开时丢符号：忘了 $i^2=-1$}
用 $e^{i(\alpha+\beta)}$ 展开时，交叉项来自 $i\sin\alpha\cdot i\sin\beta=i^2\sin\alpha\sin\beta=-\sin\alpha\sin\beta$。
若忘记 $i^2=-1$，就会把实部写成 $\cos\alpha\cos\beta+\sin\alpha\sin\beta$，得到错误的「和角公式」。
记住：\key{每出现一个 $i^2$ 就要换成 $-1$}，符号自然正确。
\end{commonerror}

\begin{method}
母公式引擎使用流程（复数法三步）：
\begin{enumerate}
  \item \key{翻译}：把三角式写成 $e^{i\theta}$ 的指数形式（$\cos\theta=\frac{e^{i\theta}+e^{-i\theta}}2$，$\sin\theta=\frac{e^{i\theta}-e^{-i\theta}}{2i}$）；
  \item \key{指数运算}：用 $e^{i(\alpha+\beta)}=e^{i\alpha}e^{i\beta}$、棣莫弗定理、等比数列求和等纯代数工具计算；
  \item \key{取实虚部还原}：需要余弦取实部，需要正弦取虚部，最后写回实数三角形式。
\end{enumerate}
\end{method}

\subsection{解三角形与边角互化}

\key{探究五：正余弦定理——把三角形放进外接圆}
正弦定理：
\[
  \frac{a}{\sin A}=\frac{b}{\sin B}=\frac{c}{\sin C}=2R,
\]
其中 $R$ 是 $\triangle ABC$ 外接圆半径。几何意义：作外接圆，由圆周角定理得 $a=2R\sin A$，其余同理。
余弦定理：
\[
  a^2=b^2+c^2-2bc\cos A
  \iff \cos A=\frac{b^2+c^2-a^2}{2bc}.
\]
面积公式：
\[
  S=\frac12bc\sin A=\frac12ca\sin B=\frac12ab\sin C=\frac{abc}{4R}.
\]

\key{边角互化的选择标准。}
\begin{itemize}
  \item 关于 $a,b,c$（或 $\sin A,\sin B,\sin C$）的\key{一次齐次式} $\to$ 用正弦定理，把边化角或把角化边；
  \item 关于 $a^2,b^2,c^2$（或 $\sin^2$）的\key{二次齐次式} $\to$ 用余弦定理。
\end{itemize}

\begin{example}
在 $\triangle ABC$ 中，已知 $a=2b\cos C$，判断三角形的形状。
\end{example}
\begin{solution}
由正弦定理把边化角：$\sin A=2\sin B\cos C$。
又 $A=\pi-(B+C)$，故 $\sin A=\sin(B+C)=\sin B\cos C+\cos B\sin C$，于是
\[
  \sin B\cos C+\cos B\sin C=2\sin B\cos C
  \implies \cos B\sin C-\sin B\cos C=0
  \implies \sin(C-B)=0.
\]
因为 $B,C\in(0,\pi)$，所以 $C-B=0$，即 $B=C$。故三角形为等腰三角形。
（一次齐次式，选正弦定理。）
\end{solution}

\begin{example}
在 $\triangle ABC$ 中，已知 $a^2=b^2+c^2-bc$，求角 $A$。
\end{example}
\begin{solution}
与余弦定理 $a^2=b^2+c^2-2bc\cos A$ 对比，得 $2\cos A=1$，即 $\cos A=\frac12$。
又 $A\in(0,\pi)$，故 $A=\dfrac{\pi}{3}$（$60^\circ$）。
（二次齐次式，选余弦定理。）
\end{solution}

\key{钝角三角形的隐藏定义域。}
当题目给出「钝角三角形」或隐含某角为钝角时，角的范围被额外压缩：
钝角 $\gamma\in\left(\dfrac{\pi}{2},\pi\right)$，两个锐角满足 $\alpha+\beta<\dfrac{\pi}{2}$。
用余弦定理处理时，$\cos\gamma<0$ 给出边的平方约束，而三角形存在条件给出另一侧约束，两者共同构成隐藏定义域。

\begin{example}
在 $\triangle ABC$ 中，已知 $a=2$，$b=3$，且 $C$ 为钝角，求边 $c$ 的取值范围。
\end{example}
\begin{solution}
由余弦定理
\[
  \cos C=\frac{a^2+b^2-c^2}{2ab}=\frac{13-c^2}{12}.
\]
$C$ 为钝角 $\iff \cos C<0\iff c^2>13\iff c>\sqrt{13}$。
同时三角形存在要求 $c<a+b=5$（且 $c>|b-a|=1$，已被 $c>\sqrt{13}$ 包含）。
故 $c\in(\sqrt{13},5)$。
若只写 $c^2>13$，就漏掉了三角形存在给出的上界 $c<5$——这正是「钝角三角形」隐藏的定义域。
\end{solution}

\begin{conccard}[二级]{边角互化：一次齐次用正弦，二次齐次用余弦}{解三角形中的求值、求角、判断形状、求范围}{齐次次数判断错误；用正弦定理求角时忽略两解；忽略钝角三角形对边角范围的额外约束}
正弦定理 $\frac{a}{\sin A}=2R$ 把边与角的正弦互相转化，适合一次齐次式；
余弦定理 $a^2=b^2+c^2-2bc\cos A$ 把三边与一角联系起来，适合二次齐次式。
涉及「钝角三角形」时，必须同时写出 $\cos\gamma<0$ 与三角形存在条件，得到完整的定义域。
\end{conccard}

\begin{commonerror}{解三角形忽略隐藏定义域}
两类常见失误：
\begin{itemize}
  \item 用正弦定理由 $\sin B=\frac{b\sin A}{a}$ 求角时，只取锐角解，漏掉钝角解；
  \item 题目说「钝角三角形」却只列 $\cos\gamma<0$，漏掉三角形存在条件（两边之和大于第三边）。
\end{itemize}
例如由 $a=2,b=3$ 且 $C$ 为钝角求 $c$ 时，只写 $c>\sqrt{13}$ 会漏掉 $c<5$。
正确做法是把「角的范围」与「三角形存在条件」同时列出，求交集。
\end{commonerror}

\begin{summary}
本章以「乘以 $e^{i\beta}=$ 逆时针旋转 $\beta$」为唯一母公理，把三角函数从零散公式还原为同一个旋转群的三种语言：
数系扩张 $\mathbb{N}\to\mathbb{Z}\to\mathbb{Q}\to\mathbb{R}\to\mathbb{C}$ 让 $i^2=-1$（转 $90^\circ$ 两次等于掉头）成为二维复平面的入场券；
单位复数 $z(\theta)=\cos\theta+i\sin\theta$ 满足 $z(\alpha)z(\beta)=z(\alpha+\beta)$ 且 $z(0)=1$，故必为指数函数；
由切线速度 $(e^{k\theta})'\big|_{\theta=0}=k$ 与 $z'(0)=i$ 锁定 $k=i$，得到欧拉公式 $e^{i\theta}=\cos\theta+i\sin\theta$ 与 $e^{i\pi}+1=0$；
再用 $e^{i(\alpha+\beta)}$ 展开、棣莫弗定理与 $\cos\theta=\frac{e^{i\theta}+e^{-i\theta}}2$、$\sin\theta=\frac{e^{i\theta}-e^{-i\theta}}{2i}$，
一次性长出两角和差、二倍角三倍角、诱导公式与和差化积／积化和差。
解三角形部分用正弦定理（外接圆半径 $R$）、余弦定理与面积公式，按「一次齐次用正弦、二次齐次用余弦」做边角互化，并警惕钝角三角形的隐藏定义域。
考纲边界：复数法用于选填秒杀与验算，解答题必须还原为实数三角变换格式。
\end{summary}

\end{explore}

\begin{inclasstest}
\begin{enumerate}
  \item 计算 $i\cdot(1+i)$，并说明它在复平面上表示怎样的旋转与伸缩；\level{易}
  \item 用 $e^{i(\alpha+\beta)}=e^{i\alpha}e^{i\beta}$ 展开并对比实部与虚部，写出 $\cos(\alpha-\beta)$ 与 $\sin(\alpha-\beta)$；\level{中}
  \item 用棣莫弗定理把 $\cos3\theta$ 用 $\cos\theta$ 表示；\level{中}
  \item 在 $\triangle ABC$ 中，$a=2$，$b=3$，$C=120^\circ$，求边 $c$ 与面积 $S$。\level{中}
\end{enumerate}
\end{inclasstest}

\begin{afterclass}
\subsection*{基础巩固}
\begin{enumerate}
  \item 计算 $i^3$、$i^4$ 与 $(1+i)^2$，并说出各自的几何意义；
  \item 由 $e^{i\theta}=\cos\theta+i\sin\theta$ 解出 $\cos\theta,\sin\theta$ 的指数表达式，并验证 $\cos^2\theta+\sin^2\theta=1$。
\end{enumerate}

\subsection*{进阶提升}
\begin{enumerate}
  \item 用 $i=e^{i\pi/2}$ 推导 $\cos\left(\theta-\dfrac{\pi}{2}\right)$ 与 $\sin\left(\theta-\dfrac{\pi}{2}\right)$；
  \item 在 $\triangle ABC$ 中，$a=2$，$b=3$，$c=4$，判断三角形的形状并求面积。
\end{enumerate}

\subsection*{拓展探究}
\begin{enumerate}
  \item 用复数等比数列求 $\cos\theta+\cos2\theta+\cdots+\cos n\theta$（$\theta\neq2k\pi$），写出结果并说明它为什么只是「取实部」。
\end{enumerate}

\subsection*{教师版答案}
\begin{solution}
\textbf{基础 1：}$i^3=i^2\cdot i=-i$（逆时针转 $270^\circ$）；$i^4=(i^2)^2=1$（转 $360^\circ$ 回到原位）；
$(1+i)^2=1+2i+i^2=2i$（模由 $\sqrt2$ 变为 $2$，辐角由 $45^\circ$ 变为 $90^\circ$）。
\end{solution}
\begin{solution}
\textbf{基础 2：}由 $e^{i\theta}=\cos\theta+i\sin\theta$、$e^{-i\theta}=\cos\theta-i\sin\theta$ 解得
\[
  \cos\theta=\frac{e^{i\theta}+e^{-i\theta}}{2},\qquad
  \sin\theta=\frac{e^{i\theta}-e^{-i\theta}}{2i}.
\]
于是
\[
  \cos^2\theta+\sin^2\theta
  =\frac{(e^{i\theta}+e^{-i\theta})^2}{4}+\frac{(e^{i\theta}-e^{-i\theta})^2}{(2i)^2}
  =\frac{e^{2i\theta}+2+e^{-2i\theta}-e^{2i\theta}+2-e^{-2i\theta}}{4}=1.
\]
\end{solution}
\begin{solution}
\textbf{进阶 1：}乘以 $-i=e^{-i\pi/2}$（顺时针转 $90^\circ$）：
\[
  (\cos\theta+i\sin\theta)(-i)=\sin\theta-i\cos\theta,
\]
对比实虚部得
\[
  \cos\left(\theta-\frac{\pi}{2}\right)=\sin\theta,\qquad
  \sin\left(\theta-\frac{\pi}{2}\right)=-\cos\theta.
\]
\end{solution}
\begin{solution}
\textbf{进阶 2：}最大边为 $c=4$，
\[
  \cos C=\frac{a^2+b^2-c^2}{2ab}=\frac{4+9-16}{12}=-\frac14<0,
\]
故 $C$ 为钝角，三角形是钝角三角形。
\[
  \sin C=\sqrt{1-\frac1{16}}=\frac{\sqrt{15}}4,\qquad
  S=\frac12ab\sin C=\frac12\cdot2\cdot3\cdot\frac{\sqrt{15}}4=\frac{3\sqrt{15}}4.
\]
\end{solution}
\begin{solution}
\textbf{拓展：}令 $q=e^{i\theta}$，则
\[
  \sum_{k=1}^{n}q^k=\frac{q(1-q^n)}{1-q}
  =\frac{e^{i\theta}(1-e^{in\theta})}{1-e^{i\theta}}
  =e^{i\frac{(n+1)\theta}{2}}\cdot\frac{\sin\frac{n\theta}{2}}{\sin\frac{\theta}{2}}.
\]
取实部得
\[
  \cos\theta+\cos2\theta+\cdots+\cos n\theta
  =\frac{\sin\frac{n\theta}{2}}{\sin\frac{\theta}{2}}\cos\frac{(n+1)\theta}{2}.
\]
因为 $\cos k\theta=\operatorname{Re}(e^{ik\theta})$，所以它只是等比数列求和后「取实部」。
\end{solution}
\reading{拓展阅读 08}{复数与三角的大一统}
\end{afterclass}

% 高观点导学案：平面向量与空间向量（第 3 章·高观点统一版）
% 主线：既有大小又有方向的量 → 点集的位移 → 基底下的坐标 → 内积投影算子 → 空间中的方向与法向
% 四模块：基底分解与坐标表示 → 数量积与投影算子 → 极化恒等式与面积 → 空间向量（方向与法向）
% 全程突出向量体系的独立地位：向量不是三角／几何的附庸，而是「几何空间代数化」的通用线性语言。

\section{平面向量与空间向量（高观点统一版）}

\begin{methodbox}{多框架对读}
\small
\begin{center}
\begin{tabular}{c|p{2.9cm}|p{5.0cm}|p{3.1cm}}
\toprule
视角 & 本章看点 & 别的视角怎么说 & 哪侧最省力 \\
\midrule
代数 & 坐标运算、数量积公式 & 有序数对的算术 & 硬算定坐标 \\
几何 & 分解、投影、法向量 & 图形中看方向与垂直 & 找关系、省计算 \\
集合 & 向量＝有向线段构成的集合 & 基底＝给集合装坐标系 & 界线性相关无关 \\
映射 & 点 $\to$ 坐标是同构；投影是算子 & 换基底＝换表示 & 整体搬迁结论 \\
\bottomrule
\end{tabular}
\end{center}
卡的使用：几何题先问「有没有现成的算子（投影／法向量）」，再决定建系还是综合（总图见第 0 章图~\ref{fig:four-view-map}）。
\end{methodbox}

第 1 章研究「对应」（函数），第 2 章研究「旋转」（三角函数）；本章登场的向量是第三条独立主线：
它研究\key{位移的合成与分解}。向量体系相对三角与函数具有显性独立地位——
三角函数是旋转群的语言，函数是对应的语言，而向量是\key{几何空间代数化}的通用线性语言：
不依赖角度公式、不依赖函数图象，只用「大小$+$方向」「基底分解」「数量积」三件工具，
就能把平面与空间中的距离、夹角、平行、垂直、面积全部翻译成代数。
四个模块层层展开：
\key{基底分解与坐标表示}（任一向量 $=$ 基底线性组合，唯一）$\to$
\key{数量积与投影算子}（数量积 $=$ 余弦定理的代数化）$\to$
\key{极化恒等式与面积}（共起点三向量乘积的秒杀通道）$\to$
\key{空间向量：方向与法向}（用法向量统一处理空间位置关系）。
核心思想：\key{选定一组基底，几何就变成了线性代数}。

\begin{learnobjective}
\begin{itemize}
  \item 理解向量「既有大小又有方向」的双重属性，区分向量与数量，理解自由向量（可平移）与零向量；
  \item 理解平面向量基本定理：不共线向量张成整个平面，任一向量有唯一分解；坐标即有序实数对，与第 1 章「一一对应」呼应；
  \item 掌握数量积的定义、坐标运算与投影公式，能把余弦定理与数量积互译（数量积就是余弦定理的代数化）；
  \item 掌握极化恒等式 $\overrightarrow{AB}\cdot\overrightarrow{AC}=|\overrightarrow{AM}|^2-|\overrightarrow{MB}|^2$（$M$ 为 $BC$ 中点），会用它秒杀共起点三向量乘积的最值问题；知道面积形式 $S=\frac12|\vec a\times\vec b|$；
  \item 会把空间向量坐标化，理解平面法向量、直线方向向量，用法向量统一处理线面／面面的平行、垂直与夹角（线面角、二面角）。
\end{itemize}
\end{learnobjective}

\begin{preview}
\subsection*{知识回顾}
\begin{itemize}
  \item 实数加法满足交换律：$a+b=$ \blank；类比思考：向量加法满足交换律吗？\blank；
  \item 平面直角坐标系中，点 $P$ 的坐标是 \blank（有序实数对／无序数组），它体现了点与 \blank 的一一对应；
  \item 两直线（斜率存在时）垂直的判定：斜率之积等于 \blank；
  \item 三角形面积公式 $S=\dfrac12\times$底$\times$ \blank；平行四边形面积 $=$ 底 $\times$ \blank。
\end{itemize}

\subsection*{课前速读}
阅读「平面向量」与「空间向量」两章内容，完成下列问题：
\begin{enumerate}
  \item 平面向量基本定理要求两个基底向量 \blank（不共线／相等）；
  \item 设 $\vec a=(x_1,y_1)$，$\vec b=(x_2,y_2)$，则 $\vec a\cdot\vec b=$ \blank；
  \item 与平面 \blank 的向量叫做该平面的法向量；
  \item 设直线方向向量为 $\vec s$、平面法向量为 $\vec n$，线面角 $\varphi$ 满足 $\sin\varphi=$ \blank。
\end{enumerate}
\end{preview}

\begin{explore}

\subsection{基底分解：平面向量基本定理与坐标表示}

\key{探究一：任一向量都能被基底「唯一拼装」出来}

\key{（1）从「既有大小又有方向的量」到「点集的位移」。}
向量是既有大小又有方向的量，用有向线段表示；
两个向量只要大小相等、方向相同就叫\key{相等向量}，与起点位置无关（\key{自由向量}，可平移）。
位移视角：$\overrightarrow{AB}$ 表示点 $A$ 到点 $B$ 的位移。把所有起点平移到原点 $O$，
向量 $\overrightarrow{OP}$ 与平面点 $P$ 一一对应——这就是「向量 $=$ 点集的位移」的准确含义。
规定：零向量 $\vec 0$（长度为 $0$，方向任意）；$\vec a$ 的相反向量记 $-\vec a$。

\key{（2）平面向量基本定理：不共线的两个向量张成整个平面。}
\begin{derivation}
定理内容：如果 $\vec e_1,\vec e_2$ 是同一平面内的两个\key{不共线}向量，
那么对这一平面内的任一向量 $\vec a$，存在\key{唯一}一对实数 $\lambda_1,\lambda_2$，使得
\[
  \vec a=\lambda_1\vec e_1+\lambda_2\vec e_2.
\]
存在性：过 $\vec a$ 的终点分别作 $\vec e_2,\vec e_1$ 方向的平行线，围成平行四边形，
两边即 $\lambda_1\vec e_1$ 与 $\lambda_2\vec e_2$（平行四边形法则）。
唯一性：设 $\vec a=\lambda_1\vec e_1+\lambda_2\vec e_2=\mu_1\vec e_1+\mu_2\vec e_2$，则
\[
  (\lambda_1-\mu_1)\vec e_1+(\lambda_2-\mu_2)\vec e_2=\vec 0.
\]
若 $\lambda_1\neq\mu_1$，则 $\vec e_1=\dfrac{\mu_2-\lambda_2}{\lambda_1-\mu_1}\vec e_2$，
与 $\vec e_1,\vec e_2$ 不共线矛盾；故 $\lambda_1=\mu_1$，同理 $\lambda_2=\mu_2$。唯一性得证。
\end{derivation}
$\{\vec e_1,\vec e_2\}$ 称为一组\key{基底}。注意基底不唯一、也不要求垂直——
这是与「坐标轴」的本质区别，也是「选得巧，计算就少」的前提。

\key{（3）坐标：基底下的有序实数对。}
取\key{正交单位基底} $\vec i,\vec j$（互相垂直、长度都为 $1$），则
\[
  \vec a=x\vec i+y\vec j \iff \vec a=(x,y),
\]
$(x,y)$ 就是 $\vec a$ 的\key{坐标}。当起点在原点时，向量 $(x,y)$ 恰好指向点 $(x,y)$——
几何对象（向量／点）、代数对象（有序实数对）、对应关系三者同构，
这正是第 1 章「笛卡尔积坐标化」思想的延续：\key{坐标 $=$ 向量在基底下的读数}。
坐标运算规则：
\[
  \vec a\pm\vec b=(x_1\pm x_2,\,y_1\pm y_2),\qquad
  \lambda\vec a=(\lambda x_1,\,\lambda y_1),\qquad
  |\vec a|=\sqrt{x_1^2+y_1^2}.
\]

\begin{figure}[htbp]
  \centering
  \begin{tikzpicture}[>=Stealth, semithick, font=\footnotesize]
    % ===== 左：斜基底的一般分解 =====
    \coordinate (O) at (0,0);
    \coordinate (E1) at (2,0.5);
    \coordinate (E2) at (0.6,1.6);
    \coordinate (P1) at (4,1);
    \coordinate (P2) at (4.6,2.6);
    \draw[->, gray] (-0.5,0)--(5,0);
    \draw[->, gray] (0,-0.5)--(0,3);
    \fill (O) circle (1pt);
    \draw[->, third, thick] (O)--(E1) node[right] {$\vec{e}_1$};
    \draw[->, third, thick] (O)--(E2) node[above left] {$\vec{e}_2$};
    \draw[->, main!80!black, very thick] (O)--(P2) node[above right] {$\vec{a}=x\vec{e}_1+y\vec{e}_2$};
    \draw[->, second, thick] (O)--(P1) node[midway,below] {$x\vec{e}_1$};
    \draw[->, second, thick] (P1)--(P2) node[midway,right] {$y\vec{e}_2$};
    \draw[dashed, second!70!black] (E2)--(P2);
    \node[main!70!black] at (0,-0.85) {一般基底：平行四边形法则唯一分解};
    % ===== 右：正交单位基底与坐标 =====
    \begin{scope}[xshift=6.6cm]
      \draw[->, gray] (-0.5,0)--(4.2,0) node[below right] {$x$};
      \draw[->, gray] (0,-0.5)--(0,3.2) node[above left] {$y$};
      \foreach \x in {1,2,3} \draw[gray!25, very thin] (\x,-0.4)--(\x,3);
      \foreach \y in {1,2} \draw[gray!25, very thin] (-0.4,\y)--(4,\y);
      \fill (0,0) circle (1pt);
      \draw[->, third, thick] (0,0)--(1,0) node[below] {$\vec{i}=\vec{e}_1$};
      \draw[->, third, thick] (0,0)--(0,1) node[left] {$\vec{j}=\vec{e}_2$};
      \draw[->, main!80!black, very thick] (0,0)--(3,2) node[above right] {$\vec{a}=(3,2)$};
      \draw[->, second, thick] (0,0)--(3,0) node[midway,below] {$3\vec{i}$};
      \draw[->, second, thick] (3,0)--(3,2) node[midway,right] {$2\vec{j}$};
      \draw[dashed, second!70!black] (0,2)--(3,2);
      \node[main!70!black] at (2,-0.85) {正交单位基底：读数即坐标};
    \end{scope}
  \end{tikzpicture}
  \caption{平面向量的基底分解与坐标：左图为任意不共线基底 $\{\vec e_1,\vec e_2\}$ 下 $\vec a=x\vec e_1+y\vec e_2$ 的唯一分解；右图取正交单位基底 $\{\vec i,\vec j\}$ 时，读数 $(x,y)$ 就是平面直角坐标}
  \label{fig:basis-coord}
\end{figure}

\key{（4）「基底选得巧，计算就少」。}
分解的目的不是形式，而是让未知关系落到已知方向上。看一个反复出现的模型：

\begin{example}
在 $\triangle ABC$ 中，$D$ 在 $BC$ 上且 $BD:DC=1:2$，用基底 $\{\overrightarrow{AB},\overrightarrow{AC}\}$ 表示 $\overrightarrow{AD}$。
\end{example}
\begin{solution}
取 $\overrightarrow{AB},\overrightarrow{AC}$ 为基底。由 $BD:DC=1:2$ 知 $\overrightarrow{BD}=\dfrac13\overrightarrow{BC}$，于是
\[
  \overrightarrow{AD}=\overrightarrow{AB}+\overrightarrow{BD}
  =\overrightarrow{AB}+\frac13\overrightarrow{BC}
  =\overrightarrow{AB}+\frac13(\overrightarrow{AC}-\overrightarrow{AB})
  =\frac23\overrightarrow{AB}+\frac13\overrightarrow{AC}.
\]
取 $D$ 为中点时即得常用结论 $\overrightarrow{AD}=\dfrac12(\overrightarrow{AB}+\overrightarrow{AC})$。
\key{一旦基底选定，所有向量都变成有序数对，几何题就变成解方程组。}
\end{solution}

\begin{conccard}[一级]{平面向量基本定理与坐标表示}{把任意向量用两个不共线向量线性表示；判断三点共线；坐标运算}{认为「基底必须互相垂直」；分解不唯一时怀疑定理；写出坐标却忘记它依赖所选基底}
$\{\vec e_1,\vec e_2\}$ 不共线 $\implies$ 平面内任一向量 $\vec a$ 都可唯一写成 $\lambda_1\vec e_1+\lambda_2\vec e_2$。
特别地，$\overrightarrow{OP}=x\overrightarrow{OA}+y\overrightarrow{OB}$ 且 $x+y=1\iff P,A,B$ 三点共线（等和线）。
\end{conccard}

\subsection{数量积与投影算子}

\key{探究二：数量积是「投影乘模长」，更是余弦定理的代数化}

\key{（1）定义与几何意义。}
\[
  \vec a\cdot\vec b=|\vec a|\,|\vec b|\cos\theta,\qquad \theta=\langle\vec a,\vec b\rangle\in[0,\pi].
\]
几何读法：先取 $\vec a$ 在 $\vec b$ 方向上的投影 $|\vec a|\cos\theta$（可正可负可为零），再乘以 $|\vec b|$。
特例：同向 $\vec a\cdot\vec b=|\vec a||\vec b|$；反向 $=-|\vec a||\vec b|$；垂直 $=0$。
显然 $\vec a\cdot\vec a=|\vec a|^2$——\key{数量积自带求模公式}。

\key{（2）坐标运算：把定义交给分配律。}
\begin{derivation}
\begin{align*}
  \vec a\cdot\vec b
  &=(x_1\vec i+y_1\vec j)\cdot(x_2\vec i+y_2\vec j)\\
  &=x_1x_2\,\underbrace{\vec i\cdot\vec i}_{=1}
    +(x_1y_2+x_2y_1)\,\underbrace{\vec i\cdot\vec j}_{=0}
    +y_1y_2\,\underbrace{\vec j\cdot\vec j}_{=1}
  =x_1x_2+y_1y_2.
\end{align*}
\end{derivation}
即「对应坐标相乘再相加」。由此还可直接验证交换律 $\vec a\cdot\vec b=\vec b\cdot\vec a$ 与分配律。
夹角公式：$\cos\theta=\dfrac{\vec a\cdot\vec b}{|\vec a||\vec b|}=\dfrac{x_1x_2+y_1y_2}{\sqrt{x_1^2+y_1^2}\sqrt{x_2^2+y_2^2}}$。

\key{（3）投影：数量投影与投影向量。}
\[
  \text{数量投影}=\frac{\vec a\cdot\vec b}{|\vec b|}=|\vec a|\cos\theta,\qquad
  \mathrm{proj}_{\vec b}\vec a=\frac{\vec a\cdot\vec b}{|\vec b|^2}\,\vec b.
\]
记忆钩子：\key{投影向量 $=$ 数量投影 $\times$ 单位向量}，即 $\mathrm{proj}_{\vec b}\vec a=\left(\vec a\cdot\dfrac{\vec b}{|\vec b|}\right)\dfrac{\vec b}{|\vec b|}$。
注意分母是 $|\vec b|^2$（保证量纲是长度），且投影方向可以逆着 $\vec b$（钝角时为负）。

\begin{figure}[htbp]
  \centering
  \begin{tikzpicture}[>=Stealth, semithick, font=\footnotesize]
    \coordinate (O) at (0,0);
    \coordinate (B) at (4,0);
    \coordinate (A) at (2.6,2.2);
    \coordinate (H) at (2.6,0);
    \draw[->, gray] (-0.4,0)--(4.6,0);
    \draw[->, gray] (0,-0.4)--(0,2.6);
    \fill (O) circle (1pt) node[below left] {$O$};
    \draw[->, main!80!black, very thick] (O)--(B) node[below right] {$\vec{b}$};
    \draw[->, main!80!black, very thick] (O)--(A) node[above] {$\vec{a}$};
    \draw[dashed, second, thick] (A)--(H);
    \draw[->, second, thick] (O)--(H);
    \draw[gray, thin] (2.35,0)--(2.35,0.25)--(2.6,0.25);
    \draw[gray] (0.9,0) arc (0:40:0.9);
    \node[gray!80!black] at (1.2,0.26) {$\theta$};
    % 两条公式统一放到底部居中说明，避免压到 x 轴与图形外撑宽图幅
    \node[second, align=center] at (2.1,-0.95)
      {向量投影 $\mathrm{proj}_{\vec b}\vec a=\dfrac{\vec a\cdot\vec b}{|\vec b|^2}\vec b$\\[3pt]
       数量投影 $=\dfrac{\vec a\cdot\vec b}{|\vec b|}=|\vec a|\cos\theta$};
  \end{tikzpicture}
  \caption{数量积的投影几何：$\vec a$ 在 $\vec b$ 方向上的垂足为 $H$，$OH$ 即投影向量；数量积 $=|\vec b|\times|\vec a|\cos\theta$，即「投影乘模长」}
  \label{fig:dot-proj}
\end{figure}

\key{（4）与余弦定理互译：数量积就是余弦定理的代数化。}
\begin{derivation}
在 $\triangle ABC$ 中记 $\overrightarrow{CA}=\vec u$，$\overrightarrow{CB}=\vec v$，则 $|\vec u|=b$，$|\vec v|=a$，$|\vec u-\vec v|=|\overrightarrow{BA}|=c$。于是
\[
  c^2=|\vec u-\vec v|^2=|\vec u|^2+|\vec v|^2-2\,\vec u\cdot\vec v
  =a^2+b^2-2ab\cos C,
\]
反解即 $\cos C=\dfrac{a^2+b^2-c^2}{2ab}$。
\end{derivation}
同一个事实的两种语言：\key{余弦定理是数量积的模长读法，数量积是余弦定理的代数化}。
第 2 章用「复平面投影」理解 $\cos$，本章用「向量投影」理解 $\cos$——两种语言指向同一个对象。

\begin{example}
已知 $|\vec a|=4$，$|\vec b|=5$，夹角 $\theta=120^\circ$，求 $\vec a\cdot\vec b$、数量投影与投影向量。
\end{example}
\begin{solution}
\[
  \vec a\cdot\vec b=4\times5\times\cos120^\circ=20\times\left(-\frac12\right)=-10.
\]
数量投影 $=\dfrac{\vec a\cdot\vec b}{|\vec b|}=\dfrac{-10}{5}=-2$（负号说明投影逆着 $\vec b$）；
投影向量 $=\dfrac{\vec a\cdot\vec b}{|\vec b|^2}\vec b=\dfrac{-10}{25}\vec b=-\dfrac25\vec b$。
\end{solution}

\begin{conccard}[一级]{数量积与投影}{求模长、判垂直、求夹角、做投影；坐标运算；与余弦定理互译}{把投影向量分母写成 $|\vec b|$；认为 $\vec a\cdot\vec b$ 是向量；默认投影非负}
$\vec a\cdot\vec b=|\vec a||\vec b|\cos\theta=x_1x_2+y_1y_2$；
$\mathrm{proj}_{\vec b}\vec a=\dfrac{\vec a\cdot\vec b}{|\vec b|^2}\vec b$，数量投影 $=\dfrac{\vec a\cdot\vec b}{|\vec b|}$。
$\vec a\perp\vec b\iff\vec a\cdot\vec b=0$（非零向量）。
\end{conccard}

\begin{commonerror}{投影运算的三个高频失误}
\paragraph{失误一：分母多除一次模长。}
投影向量是 $\dfrac{\vec a\cdot\vec b}{|\vec b|^2}\vec b$。若误写 $\dfrac{\vec a\cdot\vec b}{|\vec b|}\vec b$，
左边是「数量投影（长度）再乘向量」，量纲变成长度的平方，必然出错。
正确链条：数量投影 $\dfrac{\vec a\cdot\vec b}{|\vec b|}$（一个数）$\to$ 乘单位向量 $\dfrac{\vec b}{|\vec b|}$ $\to$ 投影向量。
\paragraph{失误二：把数量积当成向量。}
$\vec a\cdot\vec b$ 是一个\key{实数}（可正可负可为零），不是向量；
更不能写成 $\vec a\cdot\vec b\,\vec b$ 这种不伦不类的式子。
\paragraph{失误三：默认投影非负、夹角非钝。}
$\theta\in[0,\pi]$，钝角时 $\cos\theta<0$，数量投影与数量积都是负数；
例如上例 $\theta=120^\circ$ 时数量投影为 $-2$。符号携带方向信息，不能丢。
\end{commonerror}

\subsection{极化恒等式与面积}

\key{探究三：极化恒等式——共起点三向量乘积的秒杀通道}

（恒等式的核心是两向量之积 $\overrightarrow{AB}\cdot\overrightarrow{AC}$；
下面的中点形式由它配方推出。）

\key{（1）一般的极化恒等式。}
对任意向量 $\vec u,\vec v$：
\[
  \vec u\cdot\vec v=\frac14\left(|\vec u+\vec v|^2-|\vec u-\vec v|^2\right).
\]
推导只需展开右边：$|\vec u+\vec v|^2-|\vec u-\vec v|^2=4\,\vec u\cdot\vec v$。

\key{（2）三角形中的中点版本。}
设 $M$ 是 $BC$ 的中点，则 $\overrightarrow{AB}+\overrightarrow{AC}=2\overrightarrow{AM}$，$\overrightarrow{AB}-\overrightarrow{AC}=\overrightarrow{CB}$，代入得
\begin{derivation}
\[
  \overrightarrow{AB}\cdot\overrightarrow{AC}
  =\frac14\left(|2\overrightarrow{AM}|^2-|\overrightarrow{CB}|^2\right)
  =|\overrightarrow{AM}|^2-\frac14|\overrightarrow{CB}|^2
  =|\overrightarrow{AM}|^2-|\overrightarrow{MB}|^2.
\]
\end{derivation}
这个恒等式把「\key{共起点的两个向量之积}」改写成「\key{动点到中点的距离平方} $-$ 定值」，
是处理动态最值问题的标准武器。

\begin{example}
已知圆 $O$ 的半径为 $2$，弦 $BC=2$，点 $A$ 在圆上运动，求 $\overrightarrow{AB}\cdot\overrightarrow{AC}$ 的最大值。
\end{example}
\begin{solution}
取 $M$ 为 $BC$ 中点，则 $|\overrightarrow{MB}|=1$ 为定值，由极化恒等式
\[
  \overrightarrow{AB}\cdot\overrightarrow{AC}=|\overrightarrow{AM}|^2-1.
\]
问题转化为求 $|\overrightarrow{AM}|$ 的最大值。由垂径定理 $|OM|=\sqrt{2^2-1^2}=\sqrt3$，
$A$ 在圆上运动时 $|AM|$ 的最大值为 $R+|OM|=2+\sqrt3$（$A$ 位于射线 $MO$ 的延长线与圆的交点）。
因此
\[
  \left(\overrightarrow{AB}\cdot\overrightarrow{AC}\right)_{\max}=(2+\sqrt3)^2-1=6+4\sqrt3.
\]
若直接设坐标硬算，将陷入三角换元；\key{极化恒等式一步就把「双动向量乘积」化为「单动点距离」。}
\end{solution}

\key{（3）面积形式：叉积的模。}
与数量积配对的是\key{叉积（外积）}：$\vec a\times\vec b$ 是垂直于 $\vec a,\vec b$ 所在平面的向量，
其模
\[
  |\vec a\times\vec b|=|\vec a|\,|\vec b|\,|\sin\theta|
\]
恰等于以 $\vec a,\vec b$ 为邻边的平行四边形面积；坐标计算公式为
\[
  |\vec a\times\vec b|=|x_1y_2-x_2y_1|,\qquad
  S_{\triangle ABC}=\frac12|\overrightarrow{AB}\times\overrightarrow{AC}|
  =\frac12\,|x_1y_2-x_2y_1|.
\]
对照记忆：\key{点积管「角与投影」（余弦），叉积管「面积与正弦」}，二者互补。

\begin{example}
已知 $\vec a=(3,1)$，$\vec b=(2,4)$，求以 $\vec a,\vec b$ 为邻边的三角形面积。
\end{example}
\begin{solution}
\[
  S=\frac12|x_1y_2-x_2y_1|=\frac12|3\times4-1\times2|=\frac12\times10=5.
\]
\end{solution}

注记（高观点）：把三个极化恒等式按三角形面积加权求和，就得到著名的奔驰定理
$S_{\triangle BOC}\cdot\overrightarrow{OA}+S_{\triangle COA}\cdot\overrightarrow{OB}+S_{\triangle AOB}\cdot\overrightarrow{OC}=\vec 0$
（$O$ 为 $\triangle ABC$ 内一点）——它是「重心坐标」的面积版；第 7 章统计回归的样本中心点与之同源。

\begin{conccard}[二级]{极化恒等式：共起点三向量乘积}{共起点两向量数量积的最值／范围问题，尤其是「一定弦、一动顶点」模型}{认不出中点结构，仍去建系硬算；把 $|\overrightarrow{MB}|$ 误当 $|BC|$；忽略动点到中点距离的取值范围需按圆／线段讨论}
$\overrightarrow{AB}\cdot\overrightarrow{AC}=|\overrightarrow{AM}|^2-|\overrightarrow{MB}|^2$（$M$ 为 $BC$ 中点）。
更一般地 $\vec u\cdot\vec v=\frac14(|\vec u+\vec v|^2-|\vec u-\vec v|^2)$。
配套面积公式：$S_{\triangle ABC}=\frac12|\overrightarrow{AB}\times\overrightarrow{AC}|=\frac12|x_1y_2-x_2y_1|$。
\end{conccard}

\subsection{空间向量：方向向量与法向量}

\key{探究四：用法向量统一处理空间中的平行、垂直与夹角}

\key{（1）空间向量坐标化。}
以 $O$ 为原点、$\vec i,\vec j,\vec k$ 为三条两两垂直的单位向量建立空间直角坐标系，则
\[
  \vec a=x\vec i+y\vec j+z\vec k\iff \vec a=(x,y,z),\qquad
  |\vec a|=\sqrt{x^2+y^2+z^2},
\]
\[
  \vec a\cdot\vec b=x_1x_2+y_1y_2+z_1z_2,\qquad
  \cos\langle\vec a,\vec b\rangle=\frac{\vec a\cdot\vec b}{|\vec a||\vec b|}.
\]
几何空间代数化的关键一步：\key{建系 $=$ 给每个点分配坐标，之后所有几何量都是坐标运算}。

\key{（2）平面的法向量。}
与平面 $\alpha$ 垂直的向量 $\vec n$ 叫做 $\alpha$ 的\key{法向量}。
一个平面有无穷多个法向量，但它们彼此平行（共线），所以\key{法向量的方向就是平面的「朝向」}。
若 $\vec n=(A,B,C)$，$P_0(x_0,y_0,z_0)\in\alpha$，则平面方程为
\[
  A(x-x_0)+B(y-y_0)+C(z-z_0)=0\quad\text{即}\quad Ax+By+Cz+D=0.
\]
求法向量的通用方法：找平面内两个不共线向量 $\vec a,\vec b$，令 $\vec n\cdot\vec a=0$、$\vec n\cdot\vec b=0$ 解方程组。

\key{（3）线面角与二面角。}
设直线 $l$ 的方向向量为 $\vec s$，平面 $\alpha$ 的法向量为 $\vec n$：
\[
  \textbf{线面角 }\varphi:\quad \sin\varphi=\frac{|\vec s\cdot\vec n|}{|\vec s|\,|\vec n|}
  \quad(\varphi\in[0,\tfrac{\pi}{2}]).
\]
注意是\key{正弦}：直线与法向量的夹角与线面角互余。
设两平面法向量为 $\vec n_1,\vec n_2$：
\[
  \textbf{二面角 }\theta:\quad |\cos\theta|=\frac{|\vec n_1\cdot\vec n_2|}{|\vec n_1|\,|\vec n_2|},
\]
锐钝要结合图形判断（或用法向量方向调整符号）。

\begin{figure}[htbp]
  \centering
  \begin{tikzpicture}[>=Stealth, semithick, font=\small]
    % 两个半平面（平行四边形示意），公共棱为中间那条水平线
    \fill[main!8] (-2.6,0)--(2.6,0)--(1.9,2.4)--(-3.3,2.4)--cycle;
    \fill[second!10] (-2.6,0)--(2.6,0)--(1.9,-2.0)--(-3.3,-2.0)--cycle;
    \draw[thick, main!60!black] (-2.6,0)--(2.6,0)--(1.9,2.4)--(-3.3,2.4)--cycle;
    \draw[thick, second!50!black] (-2.6,0)--(2.6,0)--(1.9,-2.0)--(-3.3,-2.0)--cycle;
    % 棱与棱上一点
    \draw[very thick, gray] (-2.6,0)--(2.6,0);
    \fill (0,0) circle (1.5pt) node[below left] {$P$};
    % 二面角的平面角：在两个半平面内分别作垂直于棱的射线
    \draw[->, third, very thick] (0,0)--(-1.1,1.9) node[above left] {$\vec u\subset\alpha$};
    \draw[->, third, very thick] (0,0)--(-1.1,-1.6) node[below left] {$\vec v\subset\beta$};
    % 直角符号：u\perp l，v\perp l（放大并贴近射线）
    \draw[third!70!black, thin] (0.4,0)--(0.2,0.3464)--(-0.2,0.3464);
    \draw[third!70!black, thin] (0.4,0)--(0.1736,-0.3296)--(-0.2264,-0.3296);
    % 平面角 \theta：弧线从 \vec u 跨过 \vec v
    \draw[main, thick] (0,0) ++(120:0.95) arc (120:235:0.95);
    \node[main] at (-1.15,0) {$\theta$};
    % 法向量
    \draw[->, second, thick] (1.2,1.2)--(2.2,2.6) node[above right] {$\vec n_1$};
    \draw[->, main, thick] (1.2,-1.0)--(2.2,-2.4) node[below right] {$\vec n_2$};
    \node[main!70!black] at (-2.6,1.7) {$\alpha$};
    \node[second!60!black] at (-2.6,-1.4) {$\beta$};
    \node[gray!70!black] at (2.5,0.28) {$l$};
    \node[gray!70!black, align=center] at (0,-3.05) {沿棱 $l$ 取点 $P$，在 $\alpha,\beta$ 内分别作垂直于 $l$ 的射线，\\两条射线的夹角即二面角的平面角};
  \end{tikzpicture}
  \caption{二面角的几何图：两平面 $\alpha,\beta$ 交于棱 $l$；$\vec u,\vec v$ 分别躺在两个半平面内且都垂直于 $l$（直角符号），夹角就是二面角的平面角 $\theta$。法向量 $\vec n_1,\vec n_2$ 分别垂直于两平面，其夹角与 $\theta$ 相等或互补，故 $|\cos\theta|=\dfrac{|\vec n_1\cdot\vec n_2|}{|\vec n_1||\vec n_2|}$，锐钝需结合图形判断}
  \label{fig:dihedral-angle}
\end{figure}

\begin{center}
\footnotesize
\setlength{\tabcolsep}{3pt}
\renewcommand{\arraystretch}{1.4}
\begin{tabular}{p{2.4cm}p{6.8cm}p{3.6cm}}
\toprule
位置关系 & 向量判据 & 说明 \\
\midrule
线线平行 & $\vec s_1\parallel\vec s_2$（即 $\vec s_1\times\vec s_2=\vec 0$） & 方向向量共线 \\
线线垂直 & $\vec s_1\cdot\vec s_2=0$ & 方向向量垂直 \\
线面平行 & $\vec s\cdot\vec n=0$ 且 $l\not\subset\alpha$ & 方向向量与法向量垂直 \\
线面垂直 & $\vec s\parallel\vec n$ & 方向向量与法向量共线 \\
面面平行 & $\vec n_1\parallel\vec n_2$ & 法向量共线 \\
面面垂直 & $\vec n_1\cdot\vec n_2=0$ & 法向量垂直 \\
\bottomrule
\end{tabular}
\end{center}

\begin{figure}[htbp]
  \centering
  \begin{tikzpicture}[>=Stealth, semithick, font=\footnotesize]
    \coordinate (O) at (0,0);
    \coordinate (X) at (-2.3,-1.05);
    \coordinate (Y) at (2.9,-0.75);
    \coordinate (XY) at (0.6,-1.8);
    % 平面 α（xOy 平面，z=0）
    \fill[third!8] (O)--(X)--(XY)--(Y)--cycle;
    \draw[third!50!black, thin] (O)--(X)--(XY)--(Y)--cycle;
    \node[third!70!black] at (2.9,-1.05) {$\alpha$};
    % 坐标轴
    \draw[->, gray] (O)--(X) node[below left] {$x$};
    \draw[->, gray] (O)--(Y) node[below right] {$y$};
    \draw[->, gray] (O)--(0,3.1) node[above] {$z$};
    % 法向量 n（竖直向上，垂直于平面）
    \draw[->, third, very thick] (-0.35,-0.5)--(-0.35,1.3) node[pos=0.85,right] {$\vec{n}$};
    % 直线方向向量 s 穿过平面
    \coordinate (R) at (1.15,-0.85);
    \coordinate (S) at (1.7,1.5);
    \draw[->, main!80!black, very thick] (R)--(S) node[midway,right] {$\vec{s}$};
    \draw[dashed, gray] (S)--(1.7,0);
    % 线面角 φ
    \coordinate (I) at (1.349,0);
    \draw (1.799,0) arc (0:77:0.45);
    \node at (1.98,0.4) {$\varphi$};
    \node[main!70!black] at (0.6,-2.35) {方向向量 $\vec s$、法向量 $\vec n$ 与线面角 $\varphi$};
  \end{tikzpicture}
  \caption{空间向量与法向量：平面 $\alpha$ 的法向量 $\vec n$ 指示平面朝向；直线方向向量 $\vec s$ 穿过平面，$\varphi$ 为 $\vec s$ 与其在平面内投影的夹角（线面角）}
  \label{fig:space-normal}
\end{figure}

\begin{example}
正方体 $ABCD-A_1B_1C_1D_1$ 的棱长为 $1$，求二面角 $A_1-BD-C$ 的余弦值。
\end{example}
\begin{solution}
以 $D$ 为原点建系：$D(0,0,0)$，$A(1,0,0)$，$C(0,1,0)$，$B(1,1,0)$，$A_1(1,0,1)$。
平面 $BDC$ 即底面 $z=0$，取其法向量 $\vec n_1=(0,0,1)$。
平面 $A_1BD$ 内有 $\overrightarrow{DB}=(1,1,0)$，$\overrightarrow{DA_1}=(1,0,1)$，
\[
  \vec n_2=\overrightarrow{DB}\times\overrightarrow{DA_1}
  =(1\times1-0\times0,\;0\times1-1\times1,\;1\times0-1\times1)=(1,-1,-1).
\]
于是
\[
  |\cos\theta|=\frac{|\vec n_1\cdot\vec n_2|}{|\vec n_1|\,|\vec n_2|}
  =\frac{|-1|}{1\times\sqrt3}=\frac{\sqrt3}{3}.
\]
取 $BD$ 中点 $M$，$\overrightarrow{MA_1}=(-0.5,-0.5,1)$、$\overrightarrow{MC}=(-0.5,-0.5,0)$ 均与 $BD$ 垂直，
两者点积为正，故二面角 $A_1-BD-C$ 为锐角，其余弦值为 $\dfrac{\sqrt3}{3}$。
传统几何法需要作平面角、证垂直、再解三角形；向量法只需建系、算两个法向量、代公式。
\end{solution}

\begin{example}
平面 $\alpha,\beta$ 的法向量分别为 $\vec n_1=(2,-1,1)$，$\vec n_2=(1,1,-1)$，直线 $l$ 的方向向量为 $\vec s=(1,2,2)$。
（1）判断 $l$ 与 $\alpha$ 的位置关系并求线面角；（2）判断 $\alpha$ 与 $\beta$ 的位置关系。
\end{example}
\begin{solution}
（1）$\vec s\cdot\vec n_1=1\times2+2\times(-1)+2\times1=2\neq0$，
所以 $l$ 与 $\alpha$ \key{相交}（不平行也不在面内）。线面角 $\varphi$ 满足
\[
  \sin\varphi=\frac{|\vec s\cdot\vec n_1|}{|\vec s|\,|\vec n_1|}
  =\frac{2}{\sqrt{1+4+4}\cdot\sqrt{4+1+1}}=\frac{2}{3\sqrt6}=\frac{\sqrt6}{9}.
\]
（2）$\vec n_1\cdot\vec n_2=2-1-1=0$，所以两平面的法向量垂直，即 \key{$\alpha\perp\beta$}。
\end{solution}

\begin{conccard}[二级]{法向量统一处理空间位置关系}{线面／面面的平行、垂直、夹角问题；点到平面距离；存在性问题}{把线面角公式错写成余弦；二面角不判断锐钝；建系时没说明三条轴两两垂直或坐标写错}
直线方向向量 $\vec s$、法向量 $\vec n$：
线面平行 $\iff\vec s\cdot\vec n=0$；线面垂直 $\iff\vec s\parallel\vec n$；
线面角 $\sin\varphi=\dfrac{|\vec s\cdot\vec n|}{|\vec s||\vec n|}$；
二面角 $|\cos\theta|=\dfrac{|\vec n_1\cdot\vec n_2|}{|\vec n_1||\vec n_2|}$。
（第 4 章立体几何主干直接调用这套判据与下面的易错清单，不再重复。）
\end{conccard}

\begin{commonerror}{空间向量建系的三个高频失误}
\paragraph{失误一：建系不交代垂直性。}
使用空间直角坐标系前，必须说明所选三条轴（或从同一点出发的三条棱）两两垂直；
若题目只给出斜角或没有垂直条件，应先证明垂直关系再建系。
\paragraph{失误二：线面角套错公式。}
线面角 $\varphi$ 与「方向向量和法向量的夹角」互余，所以
$\sin\varphi=\dfrac{|\vec s\cdot\vec n|}{|\vec s||\vec n|}$，而不是余弦。
把 $\vec s$ 换成它在面内的投影，才是余弦关系，切勿混用。
\paragraph{失误三：二面角不判断锐钝。}
$|\cos\theta|$ 只给大小；必须结合图形（或取棱上一点，用两个半平面内的向量验号）确定二面角是锐角还是钝角，否则一半情形会丢。
\end{commonerror}

\begin{advancedonly}
\key{超纲延伸一：换基底与相似律——同一线性变换，在不同基下「读数不同，变换不动」。}

设线性变换 $T$ 在标准基 $\{\vec e_1,\vec e_2\}$ 下的矩阵为
$A=\begin{pmatrix}2&1\\0&1\end{pmatrix}$（即 $T(\vec e_1)=(2,0)$，$T(\vec e_2)=(1,1)$）。
取 $\vec v=(2,1)$，则
\[
  T(\vec v)=A\vec v=\begin{pmatrix}2&1\\0&1\end{pmatrix}\begin{pmatrix}2\\1\end{pmatrix}
  =\tk{\begin{pmatrix}5\\1\end{pmatrix}}.
\]
换一组新基 $\mathcal B'=\{\vec u_1,\vec u_2\}$：$\vec u_1=(1,1)$，$\vec u_2=(-1,1)$，
过渡矩阵 $P=(\vec u_1\ \vec u_2)=\begin{pmatrix}1&-1\\1&1\end{pmatrix}$。
同一个向量换基后读数要左乘 $P^{-1}$；由 $\det P=2$ 得
$P^{-1}=\dfrac12\begin{pmatrix}1&1\\-1&1\end{pmatrix}$，于是
\[
  [\vec v]_{\mathcal B'}=P^{-1}\begin{pmatrix}2\\1\end{pmatrix}
  =\tk{\begin{pmatrix}\tfrac32\\[2pt]-\tfrac12\end{pmatrix}},\qquad
  [T(\vec v)]_{\mathcal B'}=P^{-1}\begin{pmatrix}5\\1\end{pmatrix}
  =\tk{\begin{pmatrix}3\\-2\end{pmatrix}}.
\]
而 $T$ 在新基下的矩阵不再是 $A$，它由相似律给出：
\[
  A'=P^{-1}AP
  =\frac12\begin{pmatrix}1&1\\-1&1\end{pmatrix}
   \begin{pmatrix}2&1\\0&1\end{pmatrix}
   \begin{pmatrix}1&-1\\1&1\end{pmatrix}
  =\tk{\begin{pmatrix}2&0\\-1&1\end{pmatrix}}.
\]
验证（一步）：用新矩阵乘新坐标，必须还原像向量的读数，
\[
  A'\,[\vec v]_{\mathcal B'}
  =\begin{pmatrix}2&0\\-1&1\end{pmatrix}
   \begin{pmatrix}\tfrac32\\[2pt]-\tfrac12\end{pmatrix}
  =\tk{\begin{pmatrix}3\\-2\end{pmatrix}}=[T(\vec v)]_{\mathcal B'}.
\]
\key{$A'=P^{-1}AP$：换基底只是换一把尺子，坐标读数变了，变换本身纹丝不动。}
顺带一提：$\det A=\det A'=2$——相似矩阵的不变量，正是「换个角度，本质仍在」。
\end{advancedonly}

\begin{advancedonly}
\key{超纲延伸二：镜面反射公式——投影走多远，就回头走多远。}

以 $\vec b$ 所在的整条直线为镜面（不只是向量 $\vec b$ 那一截），
$\vec a$ 的反射像由投影直接写出：
\[
  \vec a'=2\,\mathrm{proj}_{\vec b}\vec a-\vec a.
\]
取 $\vec a=(2.5,1.5)$，$\vec b=(3,1)$。先投影：
\[
  \vec a\cdot\vec b=2.5\times3+1.5\times1=\tk{9},\qquad
  |\vec b|^2=3^2+1^2=\tk{10},
\]
\[
  \mathrm{proj}_{\vec b}\vec a=\frac{\vec a\cdot\vec b}{|\vec b|^2}\,\vec b
  =\frac{9}{10}(3,1)=\tk{(2.7,0.9)}.
\]
再越过镜面走同样远：
\[
  \vec a'=2(2.7,0.9)-(2.5,1.5)=\tk{(2.9,0.3)}.
\]
几何解释：垂足 $\mathrm{proj}_{\vec b}\vec a$ 恰好是 $\vec a$ 与 $\vec a'$ 连线的中点，
也就是说镜面是这条连线的垂直平分线——从 $\vec a$ 到镜面有多远，像就退到另一侧多远。

\begin{figure}[htbp]
  \centering
  \begin{tikzpicture}[>=Stealth, semithick, font=\footnotesize]
    \coordinate (O) at (0,0);
    \coordinate (A) at (2.5,1.5);
    \coordinate (B) at (3,1);
    \coordinate (F) at (2.7,0.9);
    \coordinate (Ap) at (2.9,0.3);
    % 镜面：\vec b 所在的整条直线（双向延长）
    \draw[dashed, second, thick] (-1.5,-0.5)--(4.2,1.4);
    \node[second, below left] at (-1.5,-0.5) {镜面};
    % 坐标轴
    \draw[->, gray] (-1.0,0)--(4.6,0);
    \draw[->, gray] (0,-0.9)--(0,2.6);
    \fill (O) circle (1pt);
    % 投影向量（恰好躺在镜面上）与垂足
    \draw[->, third, very thick] (O)--(F);
    \fill[third] (F) circle (1.5pt);
    \node[third] at (1.55,0.30) {$\mathrm{proj}_{\vec b}\vec a$};
    % 原向量、像向量与镜面方向向量
    \draw[->, main!80!black, very thick] (O)--(A) node[above left] {$\vec a=(2.5,1.5)$};
    \draw[->, second, very thick] (O)--(B) node[above right] {$\vec b=(3,1)$};
    \draw[->, main, very thick] (O)--(Ap) node[right] {$\vec a'=(2.9,0.3)$};
    % 垂线段与垂足处的直角符号
    \draw[dashed, gray] (A)--(F)--(Ap);
    \draw[third!70!black, thin] (2.85,0.95)--(2.80,1.10)--(2.65,1.05);
  \end{tikzpicture}
  \caption{镜面反射：以 $\vec b$ 所在直线（虚线）为镜面，$\mathrm{proj}_{\vec b}\vec a$ 是 $\vec a$ 的垂足，$\vec a'=2\,\mathrm{proj}_{\vec b}\vec a-\vec a$；垂足同时也是 $\vec a$ 与 $\vec a'$ 连线的中点}
  \label{fig:mirror-reflection}
\end{figure}
\end{advancedonly}

\begin{method}
向量解题三步法：
\begin{enumerate}
  \item \key{翻译}：把几何条件（中点、平行、垂直、夹角、面积）译成向量语言；
        不确定从哪下手时，先选一组基底或建立坐标系。有垂直关系优先建系。
  \item \key{运算}：坐标化或基底化之后，用数量积处理长度与夹角，
        用「中点公式 $+$ 极化恒等式」处理共起点向量乘积的最值，
        用方向向量与法向量处理空间位置关系；追求「选得巧，计算就少」。
  \item \key{回译}：把代数结果翻回几何结论，并检查退化情形
        （零向量、两向量共线、点在平面内、分母为零）。
\end{enumerate}
\end{method}

\begin{summary}
本章把向量当作「既有大小又有方向的量」$\to$「点集的位移」$\to$「基底下的坐标」$\to$「内积投影算子」$\to$「空间中的方向与法向」层层展开，凸显向量体系相对三角与函数的独立地位：它是几何空间代数化的通用线性语言。
平面向量基本定理保证任一向量被不共线基底唯一分解，坐标就是有序实数对；数量积 $\vec a\cdot\vec b=|\vec a||\vec b|\cos\theta$ 的几何意义是「投影乘模长」，投影向量 $\mathrm{proj}_{\vec b}\vec a=\frac{\vec a\cdot\vec b}{|\vec b|^2}\vec b$，它与余弦定理互为代数化与几何读法；
极化恒等式 $\overrightarrow{AB}\cdot\overrightarrow{AC}=|\overrightarrow{AM}|^2-|\overrightarrow{MB}|^2$ 把共起点三向量乘积化为动点到中点的距离平方差，是动态最值的秒杀通道，面积则由叉积的模 $S=\frac12|\vec a\times\vec b|$ 给出；
空间向量坐标化之后，平面的法向量成为统一语言：线面平行／垂直、线面角、二面角全部化为方向向量与法向量的数量积运算。
\end{summary}

\end{explore}

\begin{inclasstest}
\begin{enumerate}
  \item 已知 $\vec a=(1,2)$，$\vec b=(3,-1)$，求 $\vec a\cdot\vec b$、$|\vec a|$ 以及 $\vec a$ 在 $\vec b$ 上的投影向量；\level{易}
  \item 已知 $\vec e_1,\vec e_2$ 不共线，$\vec m=\vec e_1+2\vec e_2$，$\vec n=3\vec e_1+k\vec e_2$，若 $\vec m\parallel\vec n$，求实数 $k$；\level{易}
  \item 半径为 $2$ 的圆中，弦 $BC=2$，点 $A$ 在圆上运动，求 $\overrightarrow{AB}\cdot\overrightarrow{AC}$ 的取值范围；\level{中}
  \item 平面 $\alpha$ 的法向量 $\vec n_1=(1,2,-1)$，平面 $\beta$ 的法向量 $\vec n_2=(2,1,1)$，判断 $\alpha$ 与 $\beta$ 的位置关系，并求所成锐二面角的余弦值。\level{中}
\end{enumerate}
\end{inclasstest}

\begin{afterclass}
\subsection*{基础巩固}
\begin{enumerate}
  \item 已知 $\vec a=(2,-1)$，$\vec b=(-3,4)$，求 $\vec a+\vec b$、$\vec a\cdot\vec b$ 与两向量夹角的余弦值；
  \item 设 $\vec e_1,\vec e_2$ 不共线，$\vec a=\vec e_1+2\vec e_2$。若换用新基底 $\vec e_1'=\vec e_1+\vec e_2$，$\vec e_2'=\vec e_1-\vec e_2$，求 $\vec a$ 在新基底下的坐标，并说明坐标为什么随基底改变而改变。
\end{enumerate}

\subsection*{进阶提升}
\begin{enumerate}
  \item 在 $\triangle ABC$ 中，$AB=3$，$AC=5$，$BC=4$，用极化恒等式求 $\overrightarrow{AB}\cdot\overrightarrow{AC}$；
  \item 三棱锥 $D-ABC$ 中，$A(1,0,0)$，$B(0,1,0)$，$C(0,0,1)$，$D(0,0,0)$，求二面角 $A-BC-D$ 的余弦值。
\end{enumerate}

\subsection*{拓展探究}
\begin{enumerate}
  \item 用向量方法证明：三角形三条中线交于一点，且该点分每条中线为 $2:1$（提示：先用平面向量基本定理比较系数求交点分比）。
\end{enumerate}

\subsection*{教师版答案}
\begin{solution}
\textbf{基础 1：}$\vec a+\vec b=(2-3,\,-1+4)=(-1,3)$；
$\vec a\cdot\vec b=2\times(-3)+(-1)\times4=-6-4=-10$；
\[
  |\vec a|=\sqrt{2^2+(-1)^2}=\sqrt5,\quad |\vec b|=\sqrt{(-3)^2+4^2}=5,
\]
\[
  \cos\theta=\frac{\vec a\cdot\vec b}{|\vec a||\vec b|}=\frac{-10}{5\sqrt5}=-\frac{2\sqrt5}{5}.
\]
\end{solution}
\begin{solution}
\textbf{基础 2：}设 $\vec a=\lambda\vec e_1'+\mu\vec e_2'$，即
\[
  \vec a=\lambda(\vec e_1+\vec e_2)+\mu(\vec e_1-\vec e_2)
  =(\lambda+\mu)\vec e_1+(\lambda-\mu)\vec e_2.
\]
与 $\vec a=\vec e_1+2\vec e_2$ 比较系数（基本定理保证唯一）：
\[
  \lambda+\mu=1,\quad \lambda-\mu=2\implies\lambda=\frac32,\ \mu=-\frac12.
\]
故新坐标为 $\left(\dfrac32,-\dfrac12\right)$。
坐标是「向量在基底下的读数」，基底变了读数自然改变；变的只是坐标，向量本身不变。
\end{solution}
\begin{solution}
\textbf{进阶 1：}取 $BC$ 中点 $M$，$|MB|=2$，中线长
\[
  |AM|^2=\frac{2AB^2+2AC^2-BC^2}{4}
  =\frac{2\times9+2\times25-16}{4}=\frac{52}{4}=13.
\]
由极化恒等式
\[
  \overrightarrow{AB}\cdot\overrightarrow{AC}=|\overrightarrow{AM}|^2-|\overrightarrow{MB}|^2=13-4=9.
\]
（检验：$\cos A=\frac{9+25-16}{2\times3\times5}=\frac{3}{5}$，故 $\overrightarrow{AB}\cdot\overrightarrow{AC}=3\times5\times\frac35=9$。）
\end{solution}
\begin{solution}
\textbf{进阶 2：}平面 $ABC$：由 $A,B,C$ 的坐标知方程为 $x+y+z=1$，法向量 $\vec n_1=(1,1,1)$；
平面 $DBC$ 过原点且含 $B,C$，方程为 $x=0$，法向量 $\vec n_2=(1,0,0)$。
\[
  |\cos\theta|=\frac{|\vec n_1\cdot\vec n_2|}{|\vec n_1||\vec n_2|}
  =\frac{1}{\sqrt3\times1}=\frac{\sqrt3}{3}.
\]
取 $BC$ 中点 $M(0,\tfrac12,\tfrac12)$，$\overrightarrow{MA}=(1,-\tfrac12,-\tfrac12)$、$\overrightarrow{MD}=(0,-\tfrac12,-\tfrac12)$ 均与 $\overrightarrow{BC}=(0,1,-1)$ 垂直，且点积为正，故二面角 $A-BC-D$ 为锐角，余弦值为 $\dfrac{\sqrt3}{3}$。
\end{solution}
\begin{solution}
\textbf{拓展：}取 $A$ 为原点，记 $\overrightarrow{AB}=\vec b$，$\overrightarrow{AC}=\vec c$。
设 $D,E$ 分别为 $BC,CA$ 的中点，则
\[
  \overrightarrow{AD}=\frac{\vec b+\vec c}{2},\qquad
  \overrightarrow{BE}=\overrightarrow{AE}-\overrightarrow{AB}=\frac{\vec c}{2}-\vec b.
\]
设中线 $AD$ 与 $BE$ 交于点 $P$，令 $\overrightarrow{AP}=t\,\overrightarrow{AD}=\dfrac{t}{2}\vec b+\dfrac{t}{2}\vec c$；
另一方面 $\overrightarrow{AP}=\overrightarrow{AB}+s\,\overrightarrow{BE}=(1-s)\vec b+\dfrac{s}{2}\vec c$。
由平面向量基本定理比较系数：
\[
  \frac t2=1-s,\qquad \frac t2=\frac s2
  \implies s=t,\ \frac t2=1-t\implies t=\frac23.
\]
故 $AP:PD=BP:PE=2:1$。由对称性，$P$ 也在中线 $CF$（$F$ 为 $AB$ 中点）上，即三条中线共点，且该点（重心）分每条中线为 $2:1$。
这正是解析几何与统计中「均值点」的向量面孔（第 7 章样本中心与之同源）。
\end{solution}
\reading{拓展阅读 09}{向量：几何空间的通用线性语言}
\end{afterclass}

% 高观点导学案：立体几何与圆锥曲线大一统（第 4 章·高观点统一版）
% 主线：仿射变换（圆压缩为椭圆）→ 三维圆锥截面（Dandelin）→ 考纲安全落地双层法 → 解析几何主干 → 立体几何主干（向量法统一）
% 核心思想：圆锥曲线不是三类孤立的二次曲线，而是同一个三维对象（圆锥面）与不同倾角平面的相遇方式；
%           高观点负责「预判答案」，考纲通法负责「写出满分」。

\section{立体几何与圆锥曲线大一统（高观点版）}

\begin{methodbox}{多框架对读}
\small
\begin{center}
\begin{tabular}{c|p{2.9cm}|p{5.0cm}|p{3.1cm}}
\toprule
视角 & 本章看点 & 别的视角怎么说 & 哪侧最省力 \\
\midrule
代数 & 标准方程、判别式、弦长 & 坐标满足的条件 & 硬算与判别 \\
几何 & 轨迹形状、仿射变换、二面角 & 图形的位置与度量 & 猜结论、估范围 \\
集合 & 曲线＝点集；平行垂直＝关系 & 条件即约束的集合 & 界存在与唯一 \\
映射 & 仿射把圆映成椭圆；投影降维 & 换表示＝换坐标系 & 整体保持性质 \\
\bottomrule
\end{tabular}
\end{center}
卡的使用：解析几何先问「这是哪个变换下的不变量」，再决定代数硬算还是几何巧解（总图见第 0 章图~\ref{fig:four-view-map}）。
\end{methodbox}

第 3 章把向量确立为「几何空间代数化」的通用语言；本章再向前一步，打通立体几何与解析几何的维度壁垒：
椭圆不过是圆在\key{仿射变换}下的像，圆锥曲线不过是圆锥面在不同倾角平面上的\key{截线}——
维度不是障碍，变换才是本质。五个模块层层展开：
\key{仿射变换：从圆到椭圆的降维打击}（圆沿一个方向均匀压缩即得椭圆，共线、平行、中点、相切、面积比全部保持）$\to$
\key{Dandelin 视角}（椭圆／双曲线／抛物线是同一个圆锥面被不同倾角平面所截的结果）$\to$
\key{考纲安全落地的双层法}（高观点 10 秒预设定点定值，考纲通法写满分步骤）$\to$
\key{解析几何主干}（定义、标准方程、直线与圆锥曲线的位置关系、定点定值通法）$\to$
\key{立体几何主干}（线面关系、空间角与距离，向量法全面落地）。
核心思想：\key{圆锥曲线不是三类孤立的二次曲线，而是同一个三维对象与不同平面的相遇方式}。

\begin{learnobjective}
\begin{itemize}
  \item 理解仿射变换 $T:(x,y)\mapsto\left(x,\frac{b}{a}y\right)$，会把椭圆 $\frac{x^2}{a^2}+\frac{y^2}{b^2}=1$ 看作圆 $x^2+y^2=a^2$ 的压缩像；
  \item 掌握仿射不变量清单（共线、平行、中点、相切、面积比 $S_{\text{椭圆}}=\frac{b}{a}S_{\text{圆}}$），会用它秒算椭圆内接三角形面积最值；
  \item 知道圆锥曲线的 Dandelin 球证法：焦点—准线定义的统一，$e=\frac{\cos\beta}{\cos\alpha}$；
  \item 掌握「高观点预判 $+$ 考纲通法」双层解题策略：极点极线预设定点定值，设线—联立—韦达—判别式写满分步骤；
  \item 掌握椭圆／双曲线／抛物线的定义与标准方程，直线与圆锥曲线位置关系的判别式、弦长与中点弦公式，以及定点定值问题的通法步骤；
  \item 掌握线面平行／垂直的判定与性质、空间角（线面角、二面角）与点到平面距离，会用向量法统一处理（与第 3 章衔接）。
\end{itemize}
\end{learnobjective}

\begin{preview}
\subsection*{知识回顾}
\begin{itemize}
  \item 圆的标准方程是 \blank；把圆沿 $y$ 轴按比例 $\frac{b}{a}$ 均匀压缩后得到的图形是 \blank；
  \item 椭圆上任意一点到两焦点的距离之和等于 \blank；双曲线则对应距离之差的绝对值等于 \blank；
  \item 直线与一元二次方程联立后，判别式 $\Delta>0$ 表示有 \blank 个不同实根（交点）；
  \item 直线与平面平行的判定定理：平面外一条直线与此平面内一条直线 \blank，则线面平行。
\end{itemize}

\subsection*{教材速读}
阅读「圆锥曲线」与「立体几何」两章教材，完成下列问题：
\begin{enumerate}
  \item 椭圆定义要求 $2a$ \blank $2c$（填 $>$ 或 $<$）；双曲线要求 $2a$ \blank $2c$；
  \item 抛物线 $y^2=2px\,(p>0)$ 的焦点坐标是 \blank，准线方程是 \blank；
  \item 弦长公式 $|AB|=\sqrt{1+k^2}\cdot$ \blank；
  \item 点到平面的距离可用等体积法或 \blank 法求解。
\end{enumerate}
\end{preview}

\begin{explore}

\subsection{仿射变换：从圆到椭圆的降维打击}

\key{探究一：把椭圆「压」回圆，把难题「降」回熟题}

\key{（1）一个变换看出椭圆的身世。}
设圆 $C:x^2+y^2=a^2$。作变换
\[
  T:\begin{pmatrix}x\\y\end{pmatrix}\longmapsto\begin{pmatrix}x\\ \frac{b}{a}y\end{pmatrix},\qquad 0<b\le a.
\]
它把每一点沿 $y$ 轴方向按比例 $\frac{b}{a}$ \key{均匀压缩}：水平方向不动，竖直方向整体变扁。
设 $(x,y)\in C$，记 $(X,Y)=T(x,y)$，则 $x=X,\ y=\frac{a}{b}Y$，代入圆方程得
\[
  X^2+\left(\frac{a}{b}Y\right)^2=a^2
  \iff \frac{X^2}{a^2}+\frac{Y^2}{b^2}=1.
\]
于是 \key{椭圆就是圆在仿射变换 $T$ 下的像}。$T$ 是线性变换（矩阵 $\mathrm{diag}\left(1,\frac{b}{a}\right)$），所以它仍然把直线变成直线——这是下面所有「保持」的根源。

\key{（2）仿射不变量清单。}
因为 $T$ 把直线变成直线、且保持线段上的分比，以下事实全部保持：
\begin{itemize}
  \item \key{共线性}：三点共线 $\xrightarrow{T}$ 三点共线；
  \item \key{平行性}：$m\parallel m'\iff T(m)\parallel T(m')$；
  \item \key{中点}：线段的中点变为像线段的中点（$T$ 保持定比分点）；
  \item \key{相切}：直线与圆锥曲线相切（联立后 $\Delta=0$）的关系保持，因为 $T$ 可逆且把直线变直线；
  \item \key{面积}：任意平面图形的面积都乘以同一常数 $|\det T|=\frac{b}{a}$，故
  \[
    S_{\text{椭圆}}=\frac{b}{a}\,S_{\text{圆}}.
  \]
\end{itemize}
\key{面积比与图形的具体形状无关}，这是仿射变换最锋利的性质。

\begin{figure}[htbp]
  \centering
  \begin{tikzpicture}[>=Stealth, semithick, font=\footnotesize]
    % ---------- 左：圆与内接正三角形 ----------
    \draw[->, gray] (-2.5,0)--(2.7,0) node[below right] {$x$};
    \draw[->, gray] (0,-2.3)--(0,2.6) node[above left] {$y$};
    \fill[second!10] (0,0) circle (2);
    \draw[main!70!black, very thick] (0,0) circle (2);
    \fill[third!18] (0,2)--(-1.732,-1)--(1.732,-1)--cycle;
    \draw[third!70!black, very thick] (0,2)--(-1.732,-1)--(1.732,-1)--cycle;
    \fill (0,0) circle (1pt) node[below left] {$O$};
    \node[gray!70!black] at (1.15,-0.32) {$a$};
    \node[second!70!black] at (0,-2.02) {圆 $x^2+y^2=a^2$：$S_{\max}=\dfrac{3\sqrt3}{4}a^2$};
    % ---------- 右：压缩后的椭圆 ----------
    \begin{scope}[xshift=6.9cm]
      \draw[->, gray] (-2.5,0)--(2.7,0) node[below right] {$x$};
      \draw[->, gray] (0,-2.3)--(0,2.6) node[above left] {$y$};
      \fill[second!10] (0,0) ellipse (2 and 1.2);
      \draw[main!70!black, very thick] (0,0) ellipse (2 and 1.2);
      \fill[third!18] (0,1.2)--(-1.732,-0.6)--(1.732,-0.6)--cycle;
      \draw[third!70!black, very thick] (0,1.2)--(-1.732,-0.6)--(1.732,-0.6)--cycle;
      \fill (0,0) circle (1pt) node[below left] {$O$};
      \node[gray!70!black] at (1.15,-0.32) {$a$};
      \node[gray!70!black] at (-0.45,1.2) {$b$};
      \node[second!70!black] at (0,-2.02) {椭圆：$S_{\max}=\dfrac{b}{a}S_{\text{圆}}=\dfrac{3\sqrt3}{4}ab$};
    \end{scope}
    % ---------- 中间：压缩映射 ----------
    \draw[->, main!80!black, very thick] (3.0,2.0)--(3.7,2.0)
      node[midway,above] {$\displaystyle\binom{x}{y}\mapsto\binom{x}{\frac{b}{a}y}$};
  \end{tikzpicture}
  \caption{仿射压缩：左图为圆 $x^2+y^2=a^2$（取 $a=2$）及其内接正三角形；右图为沿 $y$ 轴按比例 $\frac{b}{a}=0.6$ 压缩后的椭圆 $\frac{x^2}{a^2}+\frac{y^2}{b^2}=1$（$b=1.2$），内接正三角形同步变形为内接三角形，面积同乘 $\frac{b}{a}$}
  \label{fig:affine-compress}
\end{figure}

\key{（3）降维打击：面积最值一步出答案。}
因为 $T$ 把所有三角形的面积同乘一个正常数，\key{「面积最大」在变换下保持}：
椭圆 $\frac{x^2}{a^2}+\frac{y^2}{b^2}=1$ 的内接三角形中，面积最大者正是圆 $x^2+y^2=a^2$ 的内接正三角形在 $T$ 下的像。
圆的内接正三角形面积为 $\frac{3\sqrt3}{4}a^2$，于是
\[
  S_{\max}=\frac{b}{a}\cdot\frac{3\sqrt3}{4}a^2=\frac{3\sqrt3}{4}ab.
\]
一个只含 $a,b$ 的干净答案——选填题可直接秒答。

\begin{example}
求椭圆 $E:\dfrac{x^2}{4}+\dfrac{y^2}{3}=1$ 的内接三角形面积的最大值。
\end{example}
\begin{solution}
这里 $a=2,\ b=\sqrt3$，压缩比为 $\frac{b}{a}=\frac{\sqrt3}{2}$。对应的圆半径为 $2$，其内接正三角形面积为
\[
  S_{\text{圆}}=\frac{3\sqrt3}{4}\cdot 2^2=3\sqrt3.
\]
故
\[
  S_{\max}=\frac{b}{a}S_{\text{圆}}=\frac{\sqrt3}{2}\cdot 3\sqrt3=\frac92
  \quad\left(\text{或直接用 }\frac{3\sqrt3}{4}ab=\frac{3\sqrt3}{4}\cdot2\cdot\sqrt3=\frac92\right).
\]
若按常规方法设三个动点坐标硬算，将陷入多元三角换元；\key{仿射变换一步就把「陌生题」变回「课本题」}。
\end{solution}

\begin{conccard}[一级]{仿射变换：椭圆是圆的压缩像}{椭圆内接三角形面积最值、共线／平行／中点类的比例问题、面积比的快速换算}{认为仿射变换保持角度与垂直；把「面积比恒定」误用于周长比；压缩后忘记新图形是椭圆而非相似形}
$T:(x,y)\mapsto\left(x,\frac{b}{a}y\right)$ 把圆 $x^2+y^2=a^2$ 变为椭圆 $\frac{x^2}{a^2}+\frac{y^2}{b^2}=1$；
共线、平行、中点、相切保持，面积同乘 $|\det T|=\frac{b}{a}$，故 $S_{\text{椭圆}}=\frac{b}{a}S_{\text{圆}}$，$S_{\triangle\max}=\frac{3\sqrt3}{4}ab$。
\end{conccard}

\begin{commonerror}{仿射变换「能保持什么、不能保持什么」}
\paragraph{失误一：以为角度也保持。}
仿射变换 $T$ 只保证「直线 $\to$ 直线、分比不变」，\key{不保持角度、不保持垂直、不保持长度}。
例如圆内接正三角形的一个内角是 $60^\circ$，压缩成椭圆的内接三角形后三个角一般不再是 $60^\circ$。
所以「椭圆 $\xrightarrow{\text{仿射}}$ 圆」可以搬用共线、平行、中点、面积比，却\textbf{不能}搬用直角与角度关系。
\paragraph{失误二：把面积比套到周长比上。}
周长在 $y$ 方向被压缩后会改变（弧长的积分并不同乘一个常数），只有\key{面积}这类二维测度才同乘 $|\det T|$。
\paragraph{失误三：忽略 $b>a$ 的情形。}
若沿 $y$ 轴拉伸（$\frac{b}{a}>1$），同样的推导给出 $\frac{x^2}{a^2}+\frac{y^2}{b^2}=1$（$b>a$），焦点在 $y$ 轴上。
公式 $S_{\max}=\frac{3\sqrt3}{4}ab$ 对两种情形都成立，但焦点位置要看 $a,b$ 谁大。
\end{commonerror}

\subsection{Dandelin 视角：圆锥曲线是三维截线}

\key{探究二：把解析几何的二次曲线放回立体几何的三维空间}

\key{（1）圆锥面与截面倾角。}
以 $z$ 轴为中心轴、半顶角为 $\alpha$ 的圆锥面方程为
\[
  x^2+y^2=z^2\tan^2\alpha .
\]
用一个平面去截它，截线的类型只取决于平面与中心轴的夹角 $\beta$（取锐角）：
\begin{itemize}
  \item $\beta>\alpha$：平面较「平」，只切断一侧锥面，得\key{椭圆}（$\beta=90^\circ$ 时退化为圆）；
  \item $\beta=\alpha$：平面恰好与一条母线平行，得\key{抛物线}；
  \item $\beta<\alpha$：平面较「陡」，同时切断两侧锥面，得\key{双曲线}。
\end{itemize}
并且离心率由这两个角完全决定：
\[
  e=\frac{\cos\beta}{\cos\alpha}.
\]
检验：$\beta=90^\circ\Rightarrow e=0$（圆）；$\beta=\alpha\Rightarrow e=1$（抛物线）；$\beta<\alpha\Rightarrow\cos\beta>\cos\alpha\Rightarrow e>1$（双曲线）。
\key{三类圆锥曲线从此不是三个定义，而是同一个对象的三种截法。}

\key{（2）Dandelin 球：把「和／差为定值」做成肉眼可见。}
在截面与锥面围成的空间内放入与圆锥面相切、又与截面相切的球（Dandelin 球），它们与截面的切点就是焦点（椭圆、双曲线用两个球，抛物线用一个球）。
设 $P$ 为截线上任一点，过 $P$ 的母线先后穿过两个切圆所在的平面。由「从一点向球引两条切线，切线长相等」：
\[
  |PF_1|=P\text{ 到第一个切圆的母线段长},\qquad |PF_2|=P\text{ 到第二个切圆的母线段长}.
\]
两段之和（椭圆）或之差（双曲线）恰为两切圆之间的母线长——这是一个与 $P$ 无关的定值。于是
\[
  \text{椭圆：}\ |PF_1|+|PF_2|=\text{定值};\qquad
  \text{双曲线：}\ \bigl||PF_1|-|PF_2|\bigr|=\text{定值}.
\]
\key{（3）统一的焦点—准线关系。}
抛物线情形只有一个球，但它已经足以给出统一定义。设 $\sigma$ 为球与圆锥的切圆所在平面，$l=\pi\cap\sigma$ 即准线（$\pi$ 为截面），$\delta$ 为 $\pi$ 与 $\sigma$ 的夹角。
设 $P$ 为截线上任一点，$h$ 为 $P$ 到 $\sigma$ 的距离，$T$ 为过 $P$ 的母线与切圆的交点。由「从一点向球引两条切线，切线长相等」：
\[
  |PF|=|PT|=\frac{h}{\cos\alpha},
  \qquad
  d(P,l)=\frac{h}{\sin\delta}.
\]
两式相除得
\[
  \frac{|PF|}{d(P,l)}=\frac{\sin\delta}{\cos\alpha}.
\]
而 $\sigma\perp$ 轴线，故 $\sin\delta=\cos\beta$，代入即得
\[
  \frac{|PF|}{d(P,l)}=\frac{\cos\beta}{\cos\alpha}=e.
\]
对椭圆、双曲线，把同一计算分别用到两个 Dandelin 球上，每一对「焦点—准线」都满足同一比值 $e$。
一般地，三种曲线统一为 \key{焦点—准线定义}：到定点（焦点）与定直线（准线）的距离之比为常数 $e$ 的点的轨迹，$e<1$ 为椭圆，$e=1$ 为抛物线，$e>1$ 为双曲线。
这正是「立体几何中的球与平面」给「解析几何中的二次曲线」发的身份证。

\begin{figure}[htbp]
  \centering
  \begin{tikzpicture}[>=Stealth, semithick, font=\footnotesize]
    % ================= 面板一：椭圆 =================
    \begin{scope}[xshift=0cm]
      \draw[dashed, gray!60] (0,-2.55)--(0,2.55);
      \draw[gray!70, thin] (0,0)--(1.162,2.492) (0,0)--(-1.162,2.492)
                         (0,0)--(1.162,-2.492) (0,0)--(-1.162,-2.492);
      \draw[second, very thick] (-2.068,0.685)--(2.068,1.415);
      \node[second!70!black] at (1.65,0.95) {截面};
      \draw[main!45!black, fill=main!8] (0,0.7347) circle (0.3106);
      \fill (0,0) circle (1pt);
      \fill[main] (-0.054,1.041) circle (1.2pt);
      \fill[main] (0.135,1.074) circle (1.2pt);
      % 两个焦点过近，合并为一个标注，并用引导线指出，避免标签互相叠字或压轴
      \node[main] at (1.05,0.45) {$\scriptstyle F_1,F_2$};
      \draw[main!55, thin] (0.78,0.52)--(0.18,1.0);
      \draw[main!80!black, thick] (1.62,2.42) ellipse (0.42 and 0.20);
      \node[main!70!black] at (0,-2.95) {$\beta>\alpha\Rightarrow$ 椭圆$\ (e<1)$};
    \end{scope}
    % ================= 面板二：抛物线 =================
    \begin{scope}[xshift=4.9cm]
      \draw[dashed, gray!60] (0,-2.55)--(0,2.55);
      \draw[gray!70, thin] (0,0)--(1.162,2.492) (0,0)--(-1.162,2.492)
                         (0,0)--(1.162,-2.492) (0,0)--(-1.162,-2.492);
      \draw[second, very thick] (-0.465,0.503)--(0.740,3.086);
      \node[second!70!black] at (0.72,0.62) {截面};
      \draw[main!45!black, fill=main!8] (0,0.75) circle (0.317);
      \fill (0,0) circle (1pt);
      \fill[main] (-0.287,0.884) circle (1.2pt) node[above] {$\scriptstyle F$};
      \draw[main!80!black, thick] (1.28,2.62) parabola bend (1.68,2.34) (2.08,2.62);
      \node[main!70!black] at (0,-2.95) {$\beta=\alpha\Rightarrow$ 抛物线$\ (e=1)$};
    \end{scope}
    % ================= 面板三：双曲线 =================
    \begin{scope}[xshift=9.8cm]
      \draw[dashed, gray!60] (0,-2.55)--(0,2.55);
      \draw[gray!70, thin] (0,0)--(1.162,2.492) (0,0)--(-1.162,2.492)
                         (0,0)--(1.162,-2.492) (0,0)--(-1.162,-2.492);
      \draw[second, very thick] (-0.330,-0.671)--(0.330,2.871);
      \node[second!70!black] at (0.42,2.55) {截面};
      \draw[main!45!black, fill=main!8] (0,0.3495) circle (0.1477);
      \draw[main!45!black, fill=main!8] (0,-0.8370) circle (0.3537);
      \fill (0,0) circle (1pt);
      \fill[main] (-0.026,0.495) circle (1.2pt) node[left] {$\scriptstyle F_1$};
      \fill[main] (-0.061,-0.489) circle (1.2pt) node[left] {$\scriptstyle F_2$};
      \draw[main!80!black, thick] (1.30,2.10) .. controls (1.78,2.32) and (1.78,2.52) .. (1.30,2.74);
      \draw[main!80!black, thick] (2.02,2.10) .. controls (1.54,2.32) and (1.54,2.52) .. (2.02,2.74);
      \node[main!70!black] at (0,-2.95) {$\beta<\alpha\Rightarrow$ 双曲线$\ (e>1)$};
    \end{scope}
  \end{tikzpicture}
  \caption{Dandelin 示意（轴截面）：同一个圆锥面被不同倾角的平面所截。$\beta>\alpha$ 得椭圆（两个 Dandelin 球各切出一个焦点），$\beta=\alpha$ 得抛物线（一个球切出焦点），$\beta<\alpha$ 得双曲线（两个球分置两侧锥面）；右上方小图标为截线实形}
  \label{fig:dandelin}
\end{figure}

\begin{example}
圆锥的半顶角 $\alpha=30^\circ$，分别用与轴线夹角 $\beta=60^\circ,30^\circ,15^\circ$ 的平面去截，判断截线类型并求离心率。
\end{example}
\begin{solution}
由 $e=\dfrac{\cos\beta}{\cos\alpha}$，其中 $\cos30^\circ=\frac{\sqrt3}{2}$：
\[
  \beta=60^\circ:\ e=\frac{\cos60^\circ}{\cos30^\circ}=\frac{1/2}{\sqrt3/2}=\frac{\sqrt3}{3}<1\Rightarrow\text{椭圆};
\]
\[
  \beta=30^\circ:\ e=\frac{\cos30^\circ}{\cos30^\circ}=1\Rightarrow\text{抛物线};
\]
\[
  \beta=15^\circ:\ e=\frac{\cos15^\circ}{\cos30^\circ}\approx\frac{0.966}{0.866}\approx1.12>1\Rightarrow\text{双曲线}.
\]
只需要比较两个角的大小，就能立刻判定类型——这是「维度打通」带来的预判力。
\end{solution}

\begin{conccard}[二级]{圆锥曲线的统一：截面与焦点—准线}{判断截线类型、求离心率、理解三类曲线定义的统一性}{把 $\beta$ 与 $\alpha$ 的大小关系记反；忘记 $\beta=90^\circ$ 是圆；以为双曲线需要「两个分开的圆锥」}
圆锥面 $x^2+y^2=z^2\tan^2\alpha$ 被与轴线夹角为 $\beta$ 的平面所截：
$\beta>\alpha$ 椭圆（$\beta=90^\circ$ 为圆），$\beta=\alpha$ 抛物线，$\beta<\alpha$ 双曲线；
统一离心率 $e=\dfrac{\cos\beta}{\cos\alpha}$，统一定义为 $\dfrac{|PF|}{d(P,l)}=e$。
\end{conccard}

\subsection{考纲安全落地：高观点预判与通法满分的双层法}

\key{探究三：上层用高观点「看透」，下层用考纲通法「写全」}

高考解析几何解答题的特点是：答案往往可以被高观点瞬间看穿，但阅卷只看考纲允许的推理链。
因此最稳的策略是\key{双层法}：先用高观点（极点极线、仿射变换）确定目标与答案，再用考纲通法（设线、联立、韦达、判别式）把过程写得滴水不漏。

\key{（1）高观点的两个预判工具。}
\begin{itemize}
  \item \key{极点极线预设定点定值}：椭圆 $\frac{x^2}{a^2}+\frac{y^2}{b^2}=1$ 外一点 $P(x_0,y_0)$ 的两条切线的切点连线（切点弦）方程为
  \[
    \frac{x_0x}{a^2}+\frac{y_0y}{b^2}=1,
  \]
  它被称为 $P$ 的\key{极线}。若 $P$ 在定直线 $x=m$ 上运动，则极线 $\frac{mx}{a^2}+\frac{y_0y}{b^2}=1$ 恒过定点 $\left(\frac{a^2}{m},0\right)$——不必动笔算，定点已出。
  \item \key{仿射变换秒算面积最值}：涉及椭圆内接图形面积的最值，先仿射回圆，用圆的结果乘 $\frac{b}{a}$ 即可。
\end{itemize}

\key{（2）考纲通法的标准四步。}
设线 $\to$ 联立 $\to$ 韦达 $\to$ 回代（并用判别式检验）。每一步都要写在卷面上，因为分步给分。

\begin{example}[两层对照示范]
椭圆 $E:\dfrac{x^2}{4}+\dfrac{y^2}{2}=1$，动点 $P$ 在定直线 $x=2$ 上，过 $P$ 作 $E$ 的两条切线，切点分别为 $A,B$。求证：直线 $AB$ 恒过定点，并求出该定点。
\end{example}
\begin{solution}
\key{第一步（高观点预判）。}设 $P(2,t)$。由极线公式，$AB$ 即 $P$ 的极线：
\[
  \frac{2x}{4}+\frac{ty}{2}=1\iff x+ty=2,
\]
它显然恒过 $(2,0)$。十秒钟，答案已经到手。

\key{第二步（考纲通法落地）。}设 $A(x_1,y_1),B(x_2,y_2)$。
由隐函数求导（第 6 章）：椭圆两边对 $x$ 求导得 $\frac{2x}{4}+\frac{2yy'}{2}=0$，即 $y'=-\frac{b^2x}{a^2y}=-\frac{x}{2y}$。
所以椭圆在 $A$ 处的切线方程为
\[
  y-y_1=-\frac{x_1}{2y_1}(x-x_1).
\]
利用 $\frac{x_1^2}{4}+\frac{y_1^2}{2}=1$ 化简，得切线的简洁形式
\[
  \frac{x_1x}{4}+\frac{y_1y}{2}=1.
\]
（若尚未学求导，也可设切线斜率 $k$，联立椭圆后令 $\Delta=0$，推出同一结果。）
因为 $P(2,t)$ 同时在两条切线上，代入得
\[
  \frac{2x_1}{4}+\frac{ty_1}{2}=1,\qquad \frac{2x_2}{4}+\frac{ty_2}{2}=1,
  \iff x_1+ty_1=2,\quad x_2+ty_2=2.
\]
这说明 $A,B$ 都在直线 $x+ty=2$ 上；两点确定一条直线，故
\[
  AB:\ x+ty=2.
\]
令 $y=0$ 得 $x=2$，与 $t$ 无关，所以 $AB$ \key{恒过定点 $(2,0)$}。

\key{对照结论}：高观点负责「知道往哪走」（预判定点 $(2,0)$），通法负责「写得干净」（求导得切线、两点定直线）。二者分工明确，考场最稳。
\end{solution}

\begin{figure}[htbp]
  \centering
  \begin{tikzpicture}[>=Stealth, semithick, font=\footnotesize]
    \draw[->, gray] (-2.9,0)--(4.4,0) node[below right] {$x$};
    \draw[->, gray] (0,-2.2)--(0,2.9) node[above left] {$y$};
    % 椭圆
    \fill[second!8] (0,0) ellipse (2 and 1.4142);
    \draw[main!70!black, very thick] (0,0) ellipse (2 and 1.4142);
    % 极线（切点弦）
    \draw[third!70!black, thick, dashed] (-0.5,2.75)--(2.0,-1.0);
    % 两条切线
    \draw[main!80!black, thick] (3,1)--($(1.771,-0.657)!-0.4!(3,1)$);
    \draw[main!80!black, thick] (3,1)--($(0.411,1.384)!-0.4!(3,1)$);
    % 过 P 的动弦
    \draw[second, thick] (-1.809,-0.603)--(3.75,1.25);
    % 点与标注
    \fill (0,0) circle (1pt);
    \fill[main] (3,1) circle (1.6pt) node[right] {$P(x_0,y_0)$};
    \fill[third] (1.771,-0.657) circle (1.6pt) node[below right] {$A$};
    \fill[third] (0.411,1.384) circle (1.6pt) node[above left] {$B$};
    \fill[second] (1.809,0.603) circle (1.6pt) node[above right] {$M$};
    \fill[second] (-1.809,-0.603) circle (1.6pt) node[below left] {$N$};
    \node[third!70!black, anchor=west] at (0.15,2.05) {$AB$：$P$ 的极线（切点弦）};
    \node[second!70!black, anchor=west] at (0.15,-1.7) {过 $P$ 的动弦 $MN$};
  \end{tikzpicture}
  \caption{极点极线与切点弦：椭圆 $\frac{x^2}{4}+\frac{y^2}{2}=1$ 外一点 $P(3,1)$，$PA,PB$ 为切线、$A,B$ 为切点，则 $AB$ 即 $P$ 的极线；过 $P$ 的任意弦 $MN$ 与极线的关系给出定点定值问题的预判。当 $P$ 在定直线 $x=m$ 上运动时，极线 $\frac{mx}{a^2}+\frac{y_0y}{b^2}=1$ 恒过定点 $\left(\frac{a^2}{m},0\right)$}
  \label{fig:pole-polar}
\end{figure}

注记（高观点对偶）：把上面的推理倒过来即得逆命题——过定点 $(2,0)$ 的直线与椭圆交于 $A,B$，则 $A,B$ 处切线的交点必在直线 $x=2$ 上。
这在极点极线理论中称为 \key{La Hire 对偶}：点 $P$ 在点 $Q$ 的极线上 $\iff$ 点 $Q$ 在点 $P$ 的极线上。

\begin{method}[双层法：高观点预判 $+$ 通法满分]
\begin{enumerate}
  \item \key{预判}：先问「答案大概长什么样」。定点定值类问题用极线公式扫一眼；面积最值类问题先仿射回圆算一遍；
  \item \key{设线}：设动直线（含参数）或设动点坐标，注意斜率不存在的情形要单独讨论；
  \item \key{联立}：与圆锥曲线方程联立，消元得到关于 $x$（或 $y$）的一元二次方程；
  \item \key{韦达}：先写 $\Delta>0$ 保证两交点，再写 $x_1+x_2,\ x_1x_2$；
  \item \key{回代}：把目标量（数量积、斜率、面积、中点）用韦达结果表示，整理成关于参数的式子；
  \item \key{定值}：令参数系数为零（或「对任意参数成立」列方程组）求出定点／定值，最后回代检验。
\end{enumerate}
\end{method}

\subsection{解析几何主干：定义、方程与直线的位置关系}

\key{探究四：把「二次曲线」拆成三件通用工具}

\key{（1）三大定义与标准方程。}
\begin{center}
\footnotesize
\setlength{\tabcolsep}{4pt}
\renewcommand{\arraystretch}{1.35}
\begin{tabular}{p{1.5cm}p{4.6cm}p{3.5cm}p{3.0cm}p{1.9cm}}
\toprule
曲线 & 定义 & 标准方程 & 焦点 & 离心率 \\
\midrule
椭圆 & $|PF_1|+|PF_2|=2a\ (2a>2c)$ & $\dfrac{x^2}{a^2}+\dfrac{y^2}{b^2}=1\ (a>b>0)$ & $(\pm c,0),\ c^2=a^2-b^2$ & $e=\dfrac ca\in(0,1)$ \\
双曲线 & $\bigl||PF_1|-|PF_2|\bigr|=2a\ (2a<2c)$ & $\dfrac{x^2}{a^2}-\dfrac{y^2}{b^2}=1$ & $(\pm c,0),\ c^2=a^2+b^2$ & $e=\dfrac ca>1$ \\
抛物线 & $|PF|=d(P,l)$ & $y^2=2px\ (p>0)$ & $\left(\dfrac p2,0\right)$，准线 $x=-\dfrac p2$ & $e=1$ \\
\bottomrule
\end{tabular}
\end{center}
三条定义其实是同一件事（到焦点与到准线的距离比为 $e$）的特例；$a,b,c$ 的关系只需记「椭圆 $c^2=a^2-b^2$、双曲线 $c^2=a^2+b^2$、抛物线只有一个 $p$」。

\key{（2）直线与圆锥曲线的位置关系：联立 $+$ 判别式。}
设直线 $l:y=kx+m$，与圆锥曲线方程联立、消去 $y$，得到关于 $x$ 的一元二次方程 $Ax^2+Bx+C=0$（$A$ 可能为 $0$），则
\[
  \Delta>0\iff \text{两个不同公共点};\quad
  \Delta=0\iff \text{相切（一个公共点）};\quad
  \Delta<0\iff \text{无公共点}.
\]
对抛物线等特殊情形，还要检查二次项系数 $A=0$（直线平行于对称轴）时的退化情况——此时方程退化为一次方程，最多一个交点。

\key{（3）弦长公式。}
设两交点为 $A(x_1,y_1),B(x_2,y_2)$，则
\[
  |AB|=\sqrt{1+k^2}\,|x_1-x_2|
  =\sqrt{1+k^2}\cdot\frac{\sqrt{\Delta}}{|A|}.
\]
若直线写成 $x=my+n$，则改用 $|AB|=\sqrt{1+m^2}\,|y_1-y_2|$。记忆钩子：\key{弦长 $=$ 水平（竖直）投影 $\times$ 放大因子}。

\key{（4）中点弦：点差法。}
设 $A(x_1,y_1),B(x_2,y_2)$ 在椭圆 $\frac{x^2}{a^2}+\frac{y^2}{b^2}=1$ 上，$M(x_0,y_0)$ 为 $AB$ 的中点。两式相减：
\[
  \frac{(x_1+x_2)(x_1-x_2)}{a^2}+\frac{(y_1+y_2)(y_1-y_2)}{b^2}=0.
\]
代入 $x_1+x_2=2x_0,\ y_1+y_2=2y_0$，得
\[
  k_{AB}=\frac{y_1-y_2}{x_1-x_2}=-\frac{b^2x_0}{a^2y_0},
  \qquad\text{等价地}\quad k_{AB}\cdot k_{OM}=-\frac{b^2}{a^2}.
\]
\key{点差法的实质是「中点 $\Rightarrow$ 斜率」}，但必须回代检验直线与椭圆确有两交点。

\begin{example}
椭圆 $E:\dfrac{x^2}{4}+\dfrac{y^2}{2}=1$ 的一条弦 $AB$ 以 $M(1,1)$ 为中点，求直线 $AB$ 的方程与弦长。
\end{example}
\begin{solution}
由点差法（$a^2=4,b^2=2,x_0=y_0=1$）：
\[
  k_{AB}=-\frac{b^2x_0}{a^2y_0}=-\frac{2}{4}=-\frac12.
\]
故直线 $AB:y-1=-\frac12(x-1)$，即
\[
  x+2y-3=0.
\]
检验（韦达）：代入椭圆得
\[
  \frac{x^2}{4}+\frac{\left(-\frac x2+\frac32\right)^2}{2}=1
  \iff 3x^2-6x+1=0,
\]
$\Delta=36-12=24>0$，确有两个交点。$x_1+x_2=2$（中点横坐标为 $1$，一致），
\[
  |x_1-x_2|=\frac{\sqrt{24}}{3}=\frac{2\sqrt6}{3},
  \qquad
  |AB|=\sqrt{1+\left(-\frac12\right)^2}\cdot\frac{2\sqrt6}{3}
  =\frac{\sqrt5}{2}\cdot\frac{2\sqrt6}{3}=\frac{\sqrt{30}}{3}.
\]
\end{solution}

\begin{conccard}[二级]{直线与圆锥曲线：判别式、弦长、中点弦}{位置关系判定、弦长计算、中点弦方程、弦中点轨迹}{忽略二次项系数为 $0$ 的退化；点差法不检验 $\Delta$；弦长公式漏乘 $\sqrt{1+k^2}$ 或把分母写成 $\sqrt\Delta$}
联立得 $Ax^2+Bx+C=0$：$\Delta>0$ 相交、$\Delta=0$ 相切、$\Delta<0$ 相离；
$|AB|=\sqrt{1+k^2}\cdot\frac{\sqrt\Delta}{|A|}$；
中点弦斜率 $k=-\dfrac{b^2x_0}{a^2y_0}$，且 $k\cdot k_{OM}=-\dfrac{b^2}{a^2}$。
\end{conccard}

\begin{commonerror}{直线与圆锥曲线联立的三个高频失误}
\paragraph{失误一：忘记讨论二次项系数为 $0$。}
消元后若 $x^2$ 项系数含参数（如直线与抛物线联立），它可能为 $0$，此时方程退化为一元一次方程，只有一个交点。
必须先声明「当 $A=0$ 时……；当 $A\ne0$ 时……」。
\paragraph{失误二：点差法求出斜率就直接写答案。}
点差法只保证「若 $M$ 是中点，则斜率必为 $k$」，并不保证这样的弦真的存在。
必须回代验证 $\Delta>0$（上面例题中 $\Delta=24>0$ 才可写）。
\paragraph{失误三：弦长公式张冠李戴。}
$|x_1-x_2|=\frac{\sqrt\Delta}{|A|}$ 只是水平投影；弦长还要乘 $\sqrt{1+k^2}$。
若直线设为 $x=my+n$，则应乘 $\sqrt{1+m^2}$ 并对 $y$ 用韦达，切勿混用。
\end{commonerror}

\begin{conccard}[三级]{定点定值问题的通法骨架}{含参数的动直线与圆锥曲线相交，证明某点／某量为定值}{只算特殊位置就下结论；忽略判别式；参数没有全部消去就说「定值」}
设线（含参数）$\to$ 联立 $\to$ 韦达（先 $\Delta>0$）$\to$ 把目标量表示成参数的式子 $\to$ 令各参数系数为零（或对任意参数成立）$\to$ 解出定点／定值并回代检验。
高观点（极线、仿射）用于提前预判答案，通法用于写满步骤。
\end{conccard}

\subsection{立体几何主干：位置关系、空间角与向量法}

\key{探究五：判定靠逻辑，计算靠向量}

\key{（1）线面平行与垂直：判定和性质是一体两面。}
\begin{itemize}
  \item \key{线面平行}：平面外一条直线与此平面内一条直线平行 $\implies$ 线面平行（判定）；
        线面平行 $\implies$ 过该直线的任一平面与原平面的交线与该直线平行（性质）。
  \item \key{线面垂直}：一条直线与平面内\key{两条相交}直线都垂直 $\implies$ 线面垂直（判定）；
        线面垂直 $\implies$ 垂直于平面内任意直线（性质）；垂直于同一平面的两直线互相平行。
  \item \key{面面关系}：线面平行 $+$ 线线平行可推面面平行；线面垂直可推面面垂直。它们是上述两个判定的「升级版」。
\end{itemize}
解题时的固定动作：\key{要证平行，先找「面内一条线与它平行」；要证垂直，先找「面内两条相交线与它垂直」}。

\key{（2）空间角与距离的定义。}
\begin{itemize}
  \item \key{线面角}：斜线与它在平面内的射影所成的角，范围 $\left[0,\frac\pi2\right]$；线面垂直时为 $\frac\pi2$，线面平行或在面内时为 $0$；
  \item \key{二面角}：在棱上取一点，在两个半平面内分别作垂直于棱的射线，两射线的夹角即二面角的平面角；
  \item \key{点到平面的距离}：点到平面的垂线段长。几何法常用\key{等体积法}（换底面，体积不变）；向量法用公式 $d=\frac{|\overrightarrow{AP}\cdot\vec n|}{|\vec n|}$。
\end{itemize}

\key{（3）向量法：把第 3 章的工具全部调用过来。}
以方向向量 $\vec s$ 与法向量 $\vec n$ 为「两个抓手」，空间位置关系与度量可全部坐标化。
线面平行／垂直、线面角、二面角四条判据与三个高频失误已在第 3 章
完整给出（点面距公式也出自同一套向量法），此处不重复；
本章只补一条应用——点到平面的距离
\[
  d=\frac{|\overrightarrow{AP}\cdot\vec n|}{|\vec n|}\quad(A\in\alpha),
\]
即 $\overrightarrow{AP}$ 在法向量方向上投影的绝对值；几何法常用等体积法对照。

\begin{example}
四棱锥 $P-ABCD$ 中，底面 $ABCD$ 是边长为 $2$ 的正方形，$PA\perp$ 底面 $ABCD$，$PA=2$。
（1）求证：$BD\perp$ 平面 $PAC$；（2）求点 $A$ 到平面 $PCD$ 的距离；（3）求二面角 $A-PC-D$ 的余弦值。
\end{example}
\begin{solution}
以 $A$ 为原点，$\overrightarrow{AB},\overrightarrow{AD},\overrightarrow{AP}$ 分别为 $x,y,z$ 轴建立空间直角坐标系（三条轴两两垂直），则
\[
  A(0,0,0),\ B(2,0,0),\ C(2,2,0),\ D(0,2,0),\ P(0,0,2).
\]
（1）$\overrightarrow{BD}=(-2,2,0)$，$\overrightarrow{AC}=(2,2,0)$，$\overrightarrow{AP}=(0,0,2)$。
\[
  \overrightarrow{BD}\cdot\overrightarrow{AC}=-4+4=0,\qquad
  \overrightarrow{BD}\cdot\overrightarrow{AP}=0,
\]
故 $BD\perp AC$ 且 $BD\perp AP$。又 $AC\cap AP=A$，所以 $BD\perp$ 平面 $PAC$。

（2）$\overrightarrow{PC}=(2,2,-2)$，$\overrightarrow{PD}=(0,2,-2)$，平面 $PCD$ 的法向量
\[
  \vec n=\overrightarrow{PC}\times\overrightarrow{PD}=(0,4,4)\parallel(0,1,1).
\]
点 $A$ 到平面 $PCD$ 的距离
\[
  d=\frac{|\overrightarrow{AD}\cdot\vec n|}{|\vec n|}
  =\frac{|(0,2,0)\cdot(0,1,1)|}{\sqrt2}=\frac{2}{\sqrt2}=\sqrt2.
\]
几何法对照（等体积法）：
$V_{P-ACD}=\frac13 S_{\triangle ACD}\cdot PA=\frac13\cdot2\cdot2=\frac43$；
又 $PD^2+CD^2=8+4=12=PC^2$，故 $\triangle PCD$ 在 $D$ 处为直角，$S_{\triangle PCD}=\frac12\cdot2\sqrt2\cdot2=2\sqrt2$。
由 $\frac13 S_{\triangle PCD}\cdot h=\frac43$ 得 $h=\sqrt2$，与向量法一致。

（3）平面 $APC$ 的法向量 $\vec n_1=\overrightarrow{AP}\times\overrightarrow{AC}=(-4,4,0)\parallel(-1,1,0)$；平面 $PCD$ 的法向量 $\vec n_2=(0,1,1)$。
\[
  |\cos\theta|=\frac{|\vec n_1\cdot\vec n_2|}{|\vec n_1|\,|\vec n_2|}
  =\frac{|0+1+0|}{\sqrt2\cdot\sqrt2}=\frac12.
\]
取 $PC$ 的中点 $M(1,1,1)$，把 $\overrightarrow{MA}=(-1,-1,-1)$、$\overrightarrow{MD}=(-1,1,-1)$ 去掉沿 $\overrightarrow{PC}$ 方向的分量，得到二面角平面角的两条边，计算得其夹角余弦为 $\frac12>0$，故二面角 $A-PC-D$ 为锐角，其余弦值为
\[
  \cos\theta=\frac12.
\]
传统几何法需要作平面角、证垂直、再解三角形；向量法只需建系、算两个法向量、代公式。
\end{solution}

\begin{conccard}[一级]{立体几何的「判定—性质」闭环}{证平行／垂直时的辅助线设计}{只会背定理不会用；找错「面内两条相交直线」；用缺少相交条件的向量}
证线面平行：面外线 $\parallel$ 面内线；证线面垂直：直线 $\perp$ 面内两条相交直线；
线面垂直 $\Rightarrow$ 线线垂直；线面平行 $+$ 相交平面 $\Rightarrow$ 线线平行。
几何法与向量法各有胜场：垂直关系明显时优先几何法（等体积、三垂线），坐标好写时优先向量法。
\end{conccard}

\begin{summary}
本章用「仿射变换 $+$ 三维截面 $+$ 向量法」三件武器打通了立体几何与解析几何的维度壁垒。
椭圆是圆沿 $y$ 轴按比例 $\frac{b}{a}$ 的均匀压缩像，仿射不变量（共线、平行、中点、相切、面积比）全部保持，故椭圆内接三角形最大面积为 $\frac{3\sqrt3}{4}ab$；
Dandelin 球说明椭圆／双曲线／抛物线是同一个圆锥面被不同倾角平面所截的结果，统一为焦点—准线定义，离心率 $e=\frac{\cos\beta}{\cos\alpha}$；
考场上层的「高观点预判」（极线公式预设定点、仿射变换秒算面积）与下层的「考纲通法」（设线—联立—韦达—判别式）合成双层法，既看得透又写得全；
解析几何主干落实为定义与标准方程表、判别式判位置、弦长公式与中点弦点差法，以及定点定值问题的六步通法；
立体几何主干落实为线面平行／垂直的判定与性质、空间角与距离，并由第 3 章的方向向量与法向量统一处理。
\key{维度不是障碍，变换才是本质；答案是预判出来的，步骤是通法写出来的。}
\end{summary}

\end{explore}

\begin{inclasstest}
\begin{enumerate}
  \item 把圆 $x^2+y^2=9$ 沿 $y$ 轴按比例 $\frac23$ 均匀压缩，写出所得椭圆方程，并求其内接三角形面积的最大值；\level{易}
  \item 求抛物线 $y^2=8x$ 的焦点坐标与准线方程；\level{易}
  \item 若直线 $y=x+m$ 与椭圆 $\dfrac{x^2}{4}+\dfrac{y^2}{2}=1$ 有两个不同的交点，求实数 $m$ 的取值范围；当两者相切时 $m$ 取何值？\level{中}
  \item 三棱锥 $O-ABC$ 中，$OA,OB,OC$ 两两垂直，$OA=1,OB=2,OC=3$，求直线 $OC$ 与平面 $ABC$ 所成角的正弦值。\level{中}
\end{enumerate}
\end{inclasstest}

\begin{afterclass}
\subsection*{基础巩固}
\begin{enumerate}
  \item 求椭圆 $\dfrac{x^2}{16}+\dfrac{y^2}{12}=1$ 的半长轴、半短轴、半焦距、焦点坐标与离心率；椭圆上任意一点到两焦点距离之和是多少？
  \item 求双曲线 $\dfrac{x^2}{9}-\dfrac{y^2}{16}=1$ 的实轴长、虚轴长、离心率与渐近线方程。
\end{enumerate}

\subsection*{进阶提升}
\begin{enumerate}
  \item 椭圆 $E:\dfrac{x^2}{4}+y^2=1$，动点 $P$ 在定直线 $x=4$ 上，过 $P$ 作 $E$ 的两条切线，切点分别为 $A,B$，求证：直线 $AB$ 恒过定点，并求出该定点。
  \item 椭圆 $E:\dfrac{x^2}{9}+\dfrac{y^2}{4}=1$ 中，以 $M\left(\dfrac32,1\right)$ 为中点的弦所在直线方程为何？该弦的长度是多少？
\end{enumerate}

\subsection*{拓展探究}
\begin{enumerate}
  \item （极点极线的一般化）已知椭圆 $E:\dfrac{x^2}{a^2}+\dfrac{y^2}{b^2}=1$，$P(x_0,y_0)$ 为 $E$ 外一点，过 $P$ 作两条切线，切点分别为 $A,B$。
  \begin{enumerate}
    \item 求证：切点弦（极线）$AB$ 的方程为 $\dfrac{x_0x}{a^2}+\dfrac{y_0y}{b^2}=1$；
    \item 若 $P$ 在定直线 $x=m\ (m\neq0)$ 上运动，求 $AB$ 恒过的定点，并用第 6 章的隐函数求导检验 $a=2,b=\sqrt2,x_0=2$ 的特例。
  \end{enumerate}
  （提示：设切点 $(x_1,y_1)$，用隐函数求导求切线斜率，化简时利用 $\frac{x_1^2}{a^2}+\frac{y_1^2}{b^2}=1$。）
\end{enumerate}

\subsection*{教师版答案}
\begin{solution}
\textbf{基础 1：}$a^2=16,b^2=12\Rightarrow a=4,b=2\sqrt3$；$c^2=a^2-b^2=4\Rightarrow c=2$。
焦点 $(\pm2,0)$，离心率 $e=\dfrac ca=\dfrac12$；椭圆上任意一点到两焦点距离之和为 $2a=8$。
\end{solution}
\begin{solution}
\textbf{基础 2：}$a^2=9,b^2=16\Rightarrow a=3,b=4$；$c^2=a^2+b^2=25\Rightarrow c=5$。
实轴长 $2a=6$，虚轴长 $2b=8$，离心率 $e=\dfrac53$，渐近线 $y=\pm\dfrac{b}{a}x=\pm\dfrac43x$。
\end{solution}
\begin{solution}
\textbf{进阶 1：}设 $P(4,t)$。由极线公式（或设切点用求导推导），切点弦 $AB$ 的方程为
\[
  \frac{4x}{4}+\frac{ty}{1}=1\iff x+ty=1.
\]
令 $y=0$ 得 $x=1$，与 $t$ 无关，故 $AB$ 恒过定点 $(1,0)$。
\end{solution}
\begin{solution}
\textbf{进阶 2：}由点差法，$k_{AB}=-\dfrac{b^2x_0}{a^2y_0}=-\dfrac{4\cdot\frac32}{9\cdot1}=-\dfrac23$。
直线 $AB:y-1=-\dfrac23\left(x-\dfrac32\right)$，即
\[
  2x+3y-6=0.
\]
代入椭圆：$\dfrac{x^2}{9}+\dfrac{\left(2-\frac23x\right)^2}{4}=1
\iff \dfrac29x^2-\dfrac23x=0$，得 $x=0$ 或 $x=3$。
对应交点为 $(0,2)$ 与 $(3,0)$（恰为椭圆的两个顶点），中点确为 $\left(\dfrac32,1\right)$，弦长
\[
  |AB|=\sqrt{(3-0)^2+(0-2)^2}=\sqrt{13}.
\]
\end{solution}
\begin{solution}
\textbf{拓展：（1）}设 $A(x_1,y_1)$。椭圆两边对 $x$ 求导：$\dfrac{2x}{a^2}+\dfrac{2yy'}{b^2}=0\Rightarrow y'=-\dfrac{b^2x}{a^2y}$。
故 $A$ 处切线为 $y-y_1=-\dfrac{b^2x_1}{a^2y_1}(x-x_1)$，整理：
\[
  a^2y_1y-a^2y_1^2=-b^2x_1x+b^2x_1^2
  \iff b^2x_1x+a^2y_1y=b^2x_1^2+a^2y_1^2
  =a^2b^2\left(\frac{x_1^2}{a^2}+\frac{y_1^2}{b^2}\right)=a^2b^2,
\]
两边除以 $a^2b^2$ 得切线方程 $\dfrac{x_1x}{a^2}+\dfrac{y_1y}{b^2}=1$。
因为 $P(x_0,y_0)$ 在切线上，所以 $\dfrac{x_1x_0}{a^2}+\dfrac{y_1y_0}{b^2}=1$；同理 $\dfrac{x_2x_0}{a^2}+\dfrac{y_2y_0}{b^2}=1$。
这表明 $A,B$ 都在直线 $\dfrac{x_0x}{a^2}+\dfrac{y_0y}{b^2}=1$ 上，两点定一直线，即为 $AB$ 的方程。

\textbf{（2）}令 $x_0=m$，极线为 $\dfrac{mx}{a^2}+\dfrac{y_0y}{b^2}=1$。令 $y=0$ 得 $x=\dfrac{a^2}{m}$，与 $y_0$ 无关，故 $AB$ 恒过定点 $\left(\dfrac{a^2}{m},0\right)$。
特例检验：$a=2,b=\sqrt2,m=2$，定点为 $\left(\dfrac{4}{2},0\right)=(2,0)$；极线方程 $\dfrac{2x}{4}+\dfrac{y_0y}{2}=1$ 即 $x+y_0y=2$，当 $y_0=1$ 时 $x+y=2$ 过 $(2,0)$，与课堂例题一致。
\end{solution}
\reading{拓展阅读 10}{圆锥曲线：同一个三维对象的三种截面}
\end{afterclass}

% 高观点导学案：数列与离散动态系统（第 5 章·高观点版）
% 主线：差分与求和互逆（离散微积分）→ 递推即离散动力系统（不动点与蛛网图）→ 线性递推与特征值（矩阵同构）
%       → Abel 变换（错位相减的机械化）→ 共轭变换（不动点为何生效）→ 生成函数（强基拓展，超出考纲）
% 核心思想：数列不是公式的仓库，而是离散动力系统在时间轴上的轨迹；解题的本质是换坐标系，
%           把陌生递推化为等比／等差／代数方程。

\section{数列与离散动态系统（高观点版）}

\begin{methodbox}{多框架对读}
\small
\begin{center}
\begin{tabular}{c|p{2.9cm}|p{5.0cm}|p{3.1cm}}
\toprule
视角 & 本章看点 & 别的视角怎么说 & 哪侧最省力 \\
\midrule
代数 & 通项公式、求和与差分 & 符号对象的恒等变形 & 硬算、裂项、Abel \\
几何 & 蛛网图与特征螺线 & 轨迹在坐标平面上的形状 & 看收敛、猜极限 \\
集合 & $\mathbb{N}\to\mathbb{R}$ 的离散函数 & 轨道是值的集合 & 审定义域与边界 \\
映射 & $x_{n+1}=F(x_n)$、共轭变换 & 换坐标系＝换表示 & 非线性 $\to$ 等比／等差 \\
\bottomrule
\end{tabular}
\end{center}
卡的使用：递推题先问「这是哪个映射的自迭代」，再决定搬去几何还是代数（总图见第 0 章图~\ref{fig:four-view-map}）。
\end{methodbox}

第 3 章把向量确立为「几何空间代数化」的通用语言，第 4 章用仿射变换打通了圆锥曲线的维度壁垒；
本章回到最朴素的对象——数列，把它从「套公式」升级为「离散动力系统 $+$ 降维变换」：
数列是定义域为 $\mathbb{N}^*$ 的离散函数，递推是离散时间的状态演化，求和与差分是离散的微积分，
而不动点、特征值、共轭、Abel、生成函数则是把复杂递推「降维」成等比／等差／代数方程的变换。
六个模块层层展开：
\key{差分与求和互逆}（离散版微积分基本定理，由 $S_n$ 反求通项、差分降次）$\to$
\key{递推即离散动力系统}（不动点、蛛网图、$|f'(x^*)|$ 稳定性判据）$\to$
\key{线性递推与特征值}（转移矩阵 $\begin{pmatrix}0&1\\q&p\end{pmatrix}$、特征方向即不动点方向）$\to$
\key{Abel 变换}（错位相减的机械化，与分部积分互译）$\to$
\key{共轭变换}（不动点为何生效：换坐标系、莫比乌斯分式）$\to$
\key{生成函数}（强基拓展，\key{超出高考考纲}）。
核心思想：\key{数列不是公式的仓库，而是离散动力系统在时间轴上的轨迹；解题的本质是换坐标系，把陌生递推化为等比／等差。}

\begin{learnobjective}
\begin{itemize}
  \item 理解前向差分 $\Delta a_n=a_{n+1}-a_n$ 与求和 $\sum$ 互为逆运算，掌握离散版微积分基本定理 $\sum_{k=1}^n\Delta a_k=a_{n+1}-a_1$，会由和式 $S_n$ 反求通项；
  \item 理解「差分降次」与导数降次的类比，会看差分表判断通项的次数；
  \item 会把一阶递推 $a_{n+1}=f(a_n)$ 看作离散动力系统，会解不动点 $x^*=f(x^*)$、画蛛网图，并用 $|f'(x^*)|$ 判稳定性（$<1$ 吸引、$>1$ 排斥）；
  \item 掌握二阶线性递推的特征方程 $\lambda^2-p\lambda-q=0$，理解它与转移矩阵 $\begin{pmatrix}0&1\\q&p\end{pmatrix}$ 的特征值同构，特征方向即分式递推的不动点方向；
  \item 会用 Abel 变换 $\sum_{k=1}^na_kb_k=a_nB_n-\sum_{k=1}^{n-1}(a_{k+1}-a_k)B_k$ 机械化求「等差 $\times$ 等比」型求和；
  \item 理解共轭变换（换坐标系）是不动点法生效的底层原因，会作 $b_n=a_n-x^*$ 的等比化与 $\dfrac{a_n-p}{a_n-q}$ 的分式化；
  \item 了解生成函数 $G(x)=\sum a_nx^n$（\key{超出高考考纲，仅竞赛／强基}）。
\end{itemize}
\end{learnobjective}

\begin{preview}
\subsection*{知识回顾}
\begin{itemize}
  \item 等差数列的通项公式是 \blank，前 $n$ 项和公式是 \blank；
  \item 等比数列的通项公式是 \blank，前 $n$ 项和公式是 \blank；
  \item 设 $S_n$ 是数列 $\{a_n\}$ 的前 $n$ 项和，则 $n\ge2$ 时 $a_n=$ \blank，且必须单独验证 \blank；
  \item 函数 $y=f(x)$ 的图像与直线 $y=x$ 的交点横坐标满足方程 \blank。
\end{itemize}

\subsection*{教材速读}
阅读「数列」一章教材，完成下列问题：
\begin{enumerate}
  \item 由 $S_n$ 求 $a_n$ 的「分解式」$a_n=\begin{cases}S_1,&n=1\\ S_n-S_{n-1},&n\ge2\end{cases}$ 中，$n=1$ 时 \blank（填「需要」或「不需要」）单独检验；
  \item 一阶线性递推 $a_{n+1}=pa_n+q\ (p\ne1)$ 的不动点 $x^*$ 满足 $x^*=$ \blank；
  \item 斐波那契数列 $F_1=F_2=1,\ F_{n+2}=F_{n+1}+F_n$ 的特征方程是 \blank；
  \item 「错位相减法」主要适用于 \blank 数列与 \blank 数列相乘型求和。
\end{enumerate}
\end{preview}

\begin{explore}

\subsection{差分与求和互逆：离散版微积分基本定理}

\key{探究一：$\Delta$ 与 $\sum$ 是同一个东西的两面}

\key{（1）前向差分：数列的「离散导数」。}
定义
\[
  \Delta a_n=a_{n+1}-a_n .
\]
$\Delta$ 把数列 $\{a_n\}$ 映成新数列 $\{\Delta a_n\}$；二阶差分 $\Delta^2a_n=\Delta(\Delta a_n)=a_{n+2}-2a_{n+1}+a_n$，以此类推。
以 $a_n=n^2$ 为例，列表如下：
\[
\begin{array}{c|ccccc}
  n & 1 & 2 & 3 & 4 & 5 \\ \hline
  a_n=n^2 & 1 & 4 & 9 & 16 & 25 \\
  \Delta a_n & 3 & 5 & 7 & 9 & \\
  \Delta^2a_n & 2 & 2 & 2 & &
\end{array}
\]
一阶差分 $\Delta a_n=2n+1$ 是一次等差数列，二阶差分 $\Delta^2a_n=2$ 是常数——每差一次，次数降一次。

\key{（2）差分与求和互逆：离散版微积分基本定理。}
因为求和时中间项全部抵消（telescoping）：
\[
  \sum_{k=1}^{n}\Delta a_k=\sum_{k=1}^{n}(a_{k+1}-a_k)=a_{n+1}-a_1 .
\]
这就是离散版的「牛顿—莱布尼茨公式」：$\sum$ 与 $\Delta$ 互为逆运算。
反过来，若已知前 $n$ 项和 $S_n=\sum_{k=1}^na_k$，则
\[
  a_n=S_n-S_{n-1}\quad(n\ge2),\qquad a_1=S_1 .
\]
所以「由和式反求通项」不是技巧，而是差分的逆运算；注意 $n=1$ 要单独检验。

\key{（3）差分降次：类比导数。}
若 $a_n$ 是 $k$ 次多项式，则 $\Delta a_n$ 是 $k-1$ 次多项式（类比 $(x^k)'=kx^{k-1}$）。由此得到判据：
\key{若 $\Delta^2a_n$ 恒为常数，则 $a_n$ 是二次式；若 $\Delta a_n$ 是等差数列，则 $a_n$ 是二次数列}。
差分表因此可以「降次」，把高次结构一层层剥成常数。

\begin{example}
已知 $S_n=n^2+2n$，求 $\{a_n\}$ 的通项公式。
\end{example}
\begin{solution}
$a_1=S_1=1+2=3$；当 $n\ge2$ 时
\[
  a_n=S_n-S_{n-1}=(n^2+2n)-\bigl((n-1)^2+2(n-1)\bigr)=2n+1 .
\]
$n=1$ 时 $2n+1=3=S_1$，可并入，故
\[
  a_n=2n+1 .
\]
差分表视角：$\Delta a_n=2$ 为常数，故 $a_n$ 是一次等差数列，首项 $3$、公差 $2$。
\end{solution}

\begin{example}
用「差分求和」求 $\sum_{k=1}^nk^2$。
\end{example}
\begin{solution}
利用恒等式
\[
  (k+1)^3-k^3=3k^2+3k+1 ,
\]
左边正是 $k^3$ 的差分结构。对 $k=1,2,\dots,n$ 求和，左边 telescoping：
\[
  \sum_{k=1}^n\bigl((k+1)^3-k^3\bigr)=(n+1)^3-1 ,
\]
右边为
\[
  \sum_{k=1}^n(3k^2+3k+1)=3\sum_{k=1}^nk^2+\frac{3n(n+1)}2+n .
\]
联立解得
\[
  \sum_{k=1}^nk^2=\frac{n(n+1)(2n+1)}6 .
\]
关键：把 $k^2$ 写成「差分的展开」，求和就自动相消。
\end{solution}

\begin{conccard}[一级]{差分与求和互逆：$\Delta\leftrightarrow\sum$}{由 $S_n$ 反求通项、裂项相消、判断通项次数}{忘记 $n=1$ 单独检验；把 $\Delta a_n$ 的符号与 $a_n$ 的符号混同；误以为 $\Delta$ 保持多项式次数}
$\Delta a_n=a_{n+1}-a_n$；$\sum_{k=1}^n\Delta a_k=a_{n+1}-a_1$（离散牛顿—莱布尼茨）；
$a_n=S_n-S_{n-1}\ (n\ge2)$，$a_1=S_1$；
$k$ 次多项式差分一次降一次，$\Delta^2a_n$ 为常数 $\iff a_n$ 为二次式。
\end{conccard}

\subsection{递推即离散动力系统：不动点与蛛网图}

（不动点的统一理论——连续情形、几何变换与射影看法——见第 9 章；
本节只取数列递推这一路应用。）

\key{探究二：把 $a_{n+1}=f(a_n)$ 看成状态演化}

\key{（1）离散动力系统。}
数列 $\{a_n\}$ 是定义域为 $\mathbb{N}^*$ 的函数；一阶递推
\[
  a_{n+1}=f(a_n)
\]
表示：已知当前状态 $a_n$，下一步状态由同一个规则 $f$ 给出。重复作用 $f$ 就得到轨迹 $a_1,a_2,a_3,\dots$。
这正是「离散时间、离散状态」的动力系统。

\key{（2）不动点：系统的平衡态。}
若
\[
  x^*=f(x^*),
\]
则从 $x^*$ 出发永远停在 $x^*$：$f(x^*)=x^*,f(f(x^*))=x^*,\dots$。
几何上，$x^*$ 是曲线 $y=f(x)$ 与直线 $y=x$ 的交点横坐标。
对 $a_{n+1}=pa_n+q\ (p\ne1)$，不动点为 $x^*=\dfrac{q}{1-p}$。

\key{（3）蛛网图：把迭代画成阶梯。}
画法：从 $x$ 轴上的 $a_1$ 出发，竖直画到曲线 $y=f(x)$ 得 $(a_1,a_2)$；水平画到直线 $y=x$ 得 $(a_2,a_2)$；
再竖直画到曲线得 $(a_2,a_3)$；如此反复。折线像「蜘蛛网」一样阶梯式逼近或远离不动点。

\key{（4）稳定性判据：$|f'(x^*)|$ 决定吸引还是排斥。}
在不动点附近写 $a_n=x^*+\varepsilon_n$，则
\[
  \varepsilon_{n+1}=f(x^*+\varepsilon_n)-x^*\approx f'(x^*)\,\varepsilon_n .
\]
于是偏差按因子 $\lambda=f'(x^*)$ 缩放：
\begin{itemize}
  \item $|f'(x^*)|<1$：\key{吸引（稳定）}，偏差逐次缩小，轨迹收敛到 $x^*$；
  \item $|f'(x^*)|>1$：\key{排斥（不稳定）}，偏差逐次放大，轨迹逃离 $x^*$；
  \item $f'(x^*)<0$ 时偏差每步变号，轨迹\key{摆动}收敛；$0<f'(x^*)<1$ 时偏差不变号，轨迹\key{单调}收敛。
\end{itemize}
对线性 $f(x)=px+q$，近似就是精确：$b_{n+1}=pb_n$，即等比数列。

\begin{figure}[htbp]
  \centering
  \begin{tikzpicture}[>=Stealth, semithick, font=\footnotesize]
    % ---------- 左：吸引不动点 ----------
    \begin{scope}[xshift=0cm]
      \draw[->, gray] (-0.2,0)--(5.6,0) node[below right] {$x$};
      \draw[->, gray] (0,-0.2)--(0,5.2) node[above left] {$y$};
      \draw[gray!60, dashed] (-0.2,-0.2)--(5.2,5.2) node[right] {$y=x$};
      \draw[main!70!black, very thick] (-0.1,1.95)--(5.4,4.7) node[right] {$y=\frac12x+2$};
      \fill[second] (4,4) circle (1.8pt) node[above left] {$x^*=4$};
      \draw[third!70!black, thick] (0.5,0)--(0.5,2.25)--(2.25,2.25)--(2.25,3.125)
        --(3.125,3.125)--(3.125,3.5625)--(3.5625,3.5625);
      \fill[third] (0.5,0) circle (1.2pt) node[below] {$a_1$};
      \fill[third] (2.25,2.25) circle (1.2pt) node[below right] {$a_2$};
      \fill[third] (3.125,3.125) circle (1.2pt) node[below right] {$a_3$};
      \node[main!70!black] at (2.6,-1.0) {$|f'(4)|=\frac12<1$：阶梯收缩，收敛};
    \end{scope}
    % ---------- 右：排斥不动点 ----------
    \begin{scope}[xshift=7.4cm]
      \draw[->, gray] (-0.2,0)--(5.6,0) node[below right] {$x$};
      \draw[->, gray] (0,-0.2)--(0,5.2) node[above left] {$y$};
      \draw[gray!60, dashed] (-0.2,-0.2)--(5.2,5.2) node[right] {$y=x$};
      \draw[main!70!black, very thick] (0.75,0)--(3.6,4.56) node[right] {$y=1.6x-1.2$};
      \fill[second] (2,2) circle (1.8pt) node[above left] {$x^*=2$};
      \draw[third!70!black, thick] (2.25,0)--(2.25,2.4)--(2.4,2.4)--(2.4,2.64)
        --(2.64,2.64)--(2.64,3.024)--(3.024,3.024);
      \fill[third] (2.25,0) circle (1.2pt) node[below] {$a_1$};
      \fill[third] (2.64,2.64) circle (1.2pt) node[below right] {$a_3$};
      \fill[third] (3.024,3.024) circle (1.2pt) node[below right] {$a_4$};
      \node[main!70!black] at (2.6,-1.0) {$|g'(2)|=1.6>1$：阶梯放大，发散};
    \end{scope}
  \end{tikzpicture}
  \caption{蛛网图：从 $a_1$ 出发，竖直打到 $y=f(x)$、水平反射到 $y=x$，如此阶梯迭代。左图 $|f'(x^*)|<1$，阶梯收缩到不动点；右图 $|g'(x^*)|>1$，阶梯放大逃离不动点}
  \label{fig:cobweb}
\end{figure}

\begin{example}
设 $a_1=0,\ a_{n+1}=\dfrac12a_n+2$，求通项并判断数列的收敛性。
\end{example}
\begin{solution}
不动点满足 $x^*=\frac12x^*+2$，解得 $x^*=4$。两式相减：
\[
  a_{n+1}-4=\frac12(a_n-4) .
\]
故 $\{a_n-4\}$ 是首项 $a_1-4=-4$、公比 $\frac12$ 的等比数列：
\[
  a_n-4=-4\left(\frac12\right)^{n-1}=-2^{3-n},
  \qquad
  a_n=4-2^{3-n} .
\]
因 $|f'(4)|=\frac12<1$，不动点 $4$ 吸引，$a_n\to4$（单调递增）。
\end{solution}

\begin{conccard}[二级]{不动点与稳定性：$|f'(x^*)|$ 判据}{一阶递推 $a_{n+1}=f(a_n)$ 的收敛性预判、蛛网图定性、线性化降维}{把「令 $a_{n+1}=a_n=x^*$」当成假设相邻项相等；解出 $x^*$ 不检验；对 $f'(x^*)$ 不存在或 $|f'(x^*)|=1$ 的临界情形乱下结论}
不动点 $x^*=f(x^*)$ 是曲线 $y=f(x)$ 与 $y=x$ 的交点；
$b_n=a_n-x^*$ 满足 $b_{n+1}\approx f'(x^*)b_n$；
$|f'(x^*)|<1$ 吸引（稳定），$|f'(x^*)|>1$ 排斥（不稳定）；$f'(x^*)<0$ 摆动收敛。
\end{conccard}

\begin{commonerror}{不动点与稳定性判据的三个高频失误}
\paragraph{失误一：以为「令 $a_{n+1}=a_n=x^*$」是假设相邻两项相等。}
不动点 $x^*$ 是\key{解方程 $x^*=f(x^*)$ 得到的平衡点}，不是把相邻两项强行看成相等。
它的真实作用是「作差消去常数项」（见探究五）。
\paragraph{失误二：不验证 $x^*$ 是否真是不动点。}
解出 $x^*$ 后必须回代 $f(x^*)=x^*$；尤其是分式递推，使分母为零的「解」要舍去。
\paragraph{失误三：乱用 $|f'(x^*)|$ 判据。}
该判据要求 $f$ 在 $x^*$ 附近可导且 $x^*$ 是不动点；对 $f'(x^*)$ 不存在（如 $f(x)=\sqrt x$ 在 $0$ 处）
或 $|f'(x^*)|=1$ 的临界情形，结论要另行分析。
\end{commonerror}

\subsection{线性递推与特征值：状态向量与矩阵同构}

\key{探究三：把递推升维成矩阵乘法}

\key{（1）状态向量：二阶递推 $=$ 二维线性系统。}
二阶线性递推
\[
  a_{n+2}=pa_{n+1}+qa_n
\]
等价于状态向量 $\mathbf v_n=\dbinom{a_n}{a_{n+1}}$ 的演化：
\[
  \binom{a_{n+1}}{a_{n+2}}
  =\begin{pmatrix}0&1\\ q&p\end{pmatrix}\binom{a_n}{a_{n+1}}
  =A\,\mathbf v_n .
\]
于是 $\mathbf v_n=A^{n-1}\mathbf v_1$：\key{递推的每一小步，都是矩阵 $A$ 的一次乘法}。
矩阵 $A$ 的特征方程为
\[
  \det(\lambda I-A)=\begin{vmatrix}\lambda&-1\\-q&\lambda-p\end{vmatrix}
  =\lambda^2-p\lambda-q=0 ,
\]
与递推的特征方程完全一致——这正是「递推特征方程」的矩阵来源。

\key{（2）特征值决定通项结构。}
设特征根为 $\lambda_1,\lambda_2$。
\begin{itemize}
  \item $\lambda_1\ne\lambda_2$：$a_n=c_1\lambda_1^{\,n-1}+c_2\lambda_2^{\,n-1}$；
  \item $\lambda_1=\lambda_2=\lambda$（重根）：$a_n=(c_1+c_2n)\lambda^{\,n-1}$。
\end{itemize}
原理：特征向量方向在 $A$ 作用下只伸缩不变，故沿特征方向的「坐标」按 $\lambda^n$ 增长；
两个特征方向张成状态空间，通项就是它们的线性组合。重根时出现 $n\lambda^{n-1}$ 项（类比 Jordan 块）。

\key{（3）特征方向即不动点方向（射影视角）。}
对分式递推（莫比乌斯变换）
\[
  x_{n+1}=\frac{ax_n+b}{cx_n+d}\quad(ad-bc\ne0),
\]
引入齐次坐标 $x_n=\dfrac{u_n}{v_n}$，则
\[
  \binom{u_{n+1}}{v_{n+1}}=\begin{pmatrix}a&b\\c&d\end{pmatrix}\binom{u_n}{v_n} .
\]
$x_n$ 不变（$x_{n+1}=x_n=x^*$）当且仅当 $\dbinom{x^*}{1}$ 是矩阵的特征向量：
\[
  \begin{pmatrix}a&b\\c&d\end{pmatrix}\binom{x^*}{1}=\lambda\binom{x^*}{1}
  \iff ax^*+b=\lambda x^*,\ \ cx^*+d=\lambda .
\]
消去 $\lambda$ 即得分式不动点方程 $x^*=\dfrac{ax^*+b}{cx^*+d}$，且 $\lambda=cx^*+d$。
若有两个相异特征值 $\lambda_1=cx_1+d,\lambda_2=cx_2+d$，则双不动点构造的公比正是 $\dfrac{\lambda_2}{\lambda_1}$（见探究五）。
\key{结论：矩阵的特征方向（射影点）就是分式递推的不动点；特征值之比就是等比构造的公比。}

\begin{figure}[htbp]
  \centering
  \begin{tikzpicture}[>=Stealth, semithick, font=\footnotesize]
    \draw[->, gray] (-4.6,0)--(4.6,0) node[below right] {$u$};
    \draw[->, gray] (0,-1.5)--(0,3.3) node[above left] {$v$};
    \draw[third, very thick] (-4.4,1)--(4.4,1) node[right] {$v=1$（投影数轴）};
    % 两条特征方向直线都收在 $|y|\le 2.3$ 内，避免穿进底部说明
    \draw[main!70!black, thick, dashed] (-2.0,-2.0)--(3.6,3.6) node[above right] {$v=u\ (\lambda_1=5)$};
    \draw[second!70!black, thick, dashed] (3.9,-1.3)--(-3.9,1.3) node[above left] {$v=-\frac u3\ (\lambda_2=1)$};
    \fill[main] (1,1) circle (1.8pt) node[above right] {$x_1=1$};
    \fill[second] (-3,1) circle (1.8pt) node[below left] {$x_2=-3$};
    \fill (0,0) circle (1pt) node[below left] {$O$};
    \node[gray!70!black] at (0,-2.75) {$M=\begin{pmatrix}2&3\\1&4\end{pmatrix}$ 的两条特征方向直线与投影线 $v=1$ 的交点};
  \end{tikzpicture}
  \caption{射影同构：矩阵 $M=\begin{pmatrix}2&3\\1&4\end{pmatrix}$ 的特征方向 $v=u$、$v=-\frac u3$ 与投影线 $v=1$ 的交点，正是分式递推 $x_{n+1}=\frac{2x_n+3}{x_n+4}$ 的两个不动点 $x_1=1,\ x_2=-3$；公比 $=\lambda_2/\lambda_1=1/5$}
  \label{fig:eigen-fixed}
\end{figure}

\begin{example}
求递推 $a_{n+2}=3a_{n+1}-2a_n$ 在 $a_1=1,\ a_2=3$ 下的通项。
\end{example}
\begin{solution}
这里 $p=3,\ q=-2$，转移矩阵
\[
  A=\begin{pmatrix}0&1\\-2&3\end{pmatrix},\qquad
  \det(\lambda I-A)=\lambda^2-3\lambda+2=0
  \Rightarrow \lambda_1=1,\ \lambda_2=2 .
\]
故设 $a_n=c_1\cdot1^{n-1}+c_2\cdot2^{n-1}$。由初值：
\[
  c_1+c_2=1,\qquad c_1+2c_2=3
  \ \Rightarrow\ c_2=2,\ c_1=-1 .
\]
所以
\[
  a_n=2^n-1 .
\]
检验：$a_3=3a_2-2a_1=9-2=7=2^3-1$。$\surd$
\end{solution}

\begin{conccard}[二级]{特征值法：通项的终极统一模型}{二阶（及高阶）线性常系数递推、斐波那契类递推、重根情形}{特征方程写错符号（应为 $\lambda^2-p\lambda-q=0$）；重根时仍套用两个相异根的公式；初值代入时指数起点 $n-1$ 与 $n$ 混用}
状态向量 $\mathbf v_n=\binom{a_n}{a_{n+1}}$，$\mathbf v_{n+1}=A\mathbf v_n$，$A=\begin{pmatrix}0&1\\q&p\end{pmatrix}$；
特征方程 $\lambda^2-p\lambda-q=0$；
$\lambda_1\ne\lambda_2$ 时 $a_n=c_1\lambda_1^{n-1}+c_2\lambda_2^{n-1}$；
$\lambda_1=\lambda_2=\lambda$ 时 $a_n=(c_1+c_2n)\lambda^{n-1}$。
\end{conccard}

\subsection{Abel 变换：错位相减的机械化}

\key{探究四：离散版分部积分}

\key{（1）Abel 恒等式。}
设 $B_k=\sum_{i=1}^kb_i$，$B_0=0$。则
\[
  \sum_{k=1}^na_kb_k=a_nB_n-\sum_{k=1}^{n-1}(a_{k+1}-a_k)B_k .
\]
推导如下：
\[
\begin{aligned}
  \sum_{k=1}^na_kb_k
  &=\sum_{k=1}^na_k(B_k-B_{k-1})
  =\sum_{k=1}^na_kB_k-\sum_{k=1}^na_kB_{k-1}\\
  &=a_nB_n+\sum_{k=1}^{n-1}a_kB_k-\sum_{k=2}^{n}a_kB_{k-1}
  =a_nB_n-\sum_{k=1}^{n-1}(a_{k+1}-a_k)B_k .
\end{aligned}
\]

\key{（2）与分部积分互译。}
连续版 $\int u\,dv=uv-\int v\,du$；离散版把 $u\to a_k,\ dv\to b_k,\ v\to B_k,\ du\to \Delta a_k$：
\[
\begin{array}{c|c}
  \text{分部积分} & \text{Abel 变换} \\ \hline
  uv & a_nB_n \\
  \displaystyle\int v\,du & \displaystyle\sum_{k=1}^{n-1}(a_{k+1}-a_k)B_k
\end{array}
\]
记忆：\key{「乘积求和 $=$ 末端乘积 $-$ 差分乘部分和的求和」}。

\key{（3）一次秒求「等差 $\times$ 等比」。}
当 $a_k$ 等差时，$\Delta a_k=a_{k+1}-a_k=d$ 是常数；当 $b_k$ 等比时，$B_k$ 有现成公式。代入即得，无需「错位」。

\begin{example}
求 $S_n=\sum_{k=1}^n(2k-1)\cdot2^k$。
\end{example}
\begin{solution}
取 $a_k=2k-1$（等差，$\Delta a_k=2$），$b_k=2^k$（等比，$B_k=2^{k+1}-2$）。由 Abel：
\[
\begin{aligned}
  S_n&=a_nB_n-\sum_{k=1}^{n-1}2\cdot B_k\\
  &=(2n-1)(2^{n+1}-2)-2\sum_{k=1}^{n-1}(2^{k+1}-2)\\
  &=(2n-1)(2^{n+1}-2)-2\bigl(2^{n+1}-4-2(n-1)\bigr)\\
  &=(2n-3)2^{n+1}+6 .
\end{aligned}
\]
检验：$n=1$ 时 $S_1=1\cdot2=2$，公式 $(2-3)\cdot4+6=2$ $\surd$；
$n=2$ 时 $S_2=2+12=14$，公式 $1\cdot8+6=14$ $\surd$。
\end{solution}

\begin{figure}[htbp]
  \centering
  \begin{tikzpicture}[>=Stealth, semithick, font=\footnotesize, xscale=0.55, yscale=0.75]
    % 坐标轴
    \draw[->, gray] (-0.3,0)--(16.2,0) node[below right] {$B_k$};
    \draw[->, gray] (0,-0.3)--(0,5.8) node[above left] {$a_k$};
    % 阶梯柱：$\sum a_k b_k$
    \fill[main!10] (0,0) rectangle (1,5);
    \fill[main!10] (1,0) rectangle (3,4);
    \fill[main!10] (3,0) rectangle (6,3);
    \fill[main!10] (6,0) rectangle (10,2);
    \fill[main!10] (10,0) rectangle (15,1);
    \draw[main!60!black, thick] (0,0) rectangle (1,5);
    \draw[main!60!black, thick] (1,0) rectangle (3,4);
    \draw[main!60!black, thick] (3,0) rectangle (6,3);
    \draw[main!60!black, thick] (6,0) rectangle (10,2);
    \draw[main!60!black, thick] (10,0) rectangle (15,1);
    % 大矩形 $a_nB_n$
    \draw[second!70!black, thick, dashed] (0,0) rectangle (15,1);
    \node[second!70!black] at (7.5,0.5) {$a_nB_n$};
    % 橙色差值区域
    \fill[second!14] (0,1) rectangle (1,5);
    \fill[second!14] (1,1) rectangle (3,4);
    \fill[second!14] (3,1) rectangle (6,3);
    \fill[second!14] (6,1) rectangle (10,2);
    % 标注
    \foreach \x/\lab in {0/$B_0$,1/$B_1$,3/$B_2$,6/$B_3$,10/$B_4$,15/$B_5$}
      \draw (\x,0.08)--(\x,-0.08) node[below] {\lab};
    \node[main!70!black] at (0.5,2.5) {$a_1b_1$};
    \node[main!70!black] at (2,2.5) {$a_2b_2$};
    \node[main!70!black] at (4.5,2.5) {$a_3b_3$};
    \node[main!70!black] at (8,2.5) {$a_4b_4$};
    \node[main!70!black] at (12.5,2.5) {$a_5b_5$};
    % 差值区域标注放在全图上方的空白带（阶梯最高点为 $y=5$）
    \node[second!70!black] at (7.5,5.45) {差值区域：面积 $=\sum\limits_{k=1}^{4}(a_k-a_5)b_k=\sum\limits_{k=1}^{4}B_k=20$};
  \end{tikzpicture}
  \caption{Abel 变换面积示意（取 $a_k=6-k$，$b_k=k$，则 $B_k=\frac{k(k+1)}2$）：蓝色阶梯面积为 $\sum_{k=1}^5a_kb_k=35$；虚线大矩形面积为 $a_5B_5=15$；橙色差值区域面积为 $\sum_{k=1}^{4}B_k=1+3+6+10=20$。三者关系即 $\sum a_kb_k=a_nB_n+\sum_{k=1}^{n-1}(a_k-a_{k+1})B_k$}
  \label{fig:abel-area}
\end{figure}

\begin{conccard}[一级]{Abel 变换：离散分部求和}{等差 $\times$ 等比求和、乘积型数列求和、数列不等式放缩}{把 $a_k$ 与 $b_k$ 角色用反；忘记 $B_0=0$；$n=1$ 时不检验公式退化}
设 $B_k=\sum_{i=1}^kb_i$，则
$\sum_{k=1}^na_kb_k=a_nB_n-\sum_{k=1}^{n-1}(a_{k+1}-a_k)B_k$；
它是 $\int u\,dv=uv-\int v\,du$ 的离散版本；
取 $a_k$ 为等差、$b_k$ 为等比，可三行写出求和。
\end{conccard}

\begin{commonerror}{错位相减与 Abel 变换的高频失误}
\paragraph{失误一：错位相减时项数对不齐。}
「$S_n$ 与 $qS_n$」相减时，错位后首尾会各多出一项，必须明确写出剩余的首项与末项再合并。
\paragraph{失误二：忘记 $B_0=0$ 与 $n=1$ 的检验。}
Abel 公式中 $B_0=0$；当 $n=1$ 时 $\sum_{k=1}^{n-1}$ 为空，公式退化为 $a_1B_1$，要单独验证。
\paragraph{失误三：把 $a_k$ 与 $b_k$ 的角色用反。}
公式中 $a_k$ 是被差分的那个（应选等差／简单数列），$b_k$ 是被求和的那个（应选等比／可求和的数列）。
若反过来，$\Delta a_k$ 不再简单，公式失效。
\end{commonerror}

\subsection{共轭变换：不动点为何生效}

\key{探究五：换坐标系，把复杂运动还原为缩放}

\key{（1）不是「假设相邻两项相等」，而是「换原点」。}
「令 $a_{n+1}=a_n=x^*$」只是记忆口诀。真实动作是：先解出平衡点 $x^*$，再令
\[
  b_n=a_n-x^* .
\]
则
\[
  b_{n+1}=a_{n+1}-x^*=f(a_n)-f(x^*)\approx f'(x^*)(a_n-x^*)=f'(x^*)b_n .
\]
对线性 $f(x)=px+q$，这是精确的：$b_{n+1}=pb_n$。
几何上相当于把数轴原点从 $0$ 平移到 $x^*$：原来的「拉伸 $p$ 倍 $+$ 平移 $q$」复合运动，
在新坐标系下变成纯粹的「以 $x^*$ 为中心的缩放」。这就是\key{共轭变换} $g=h^{-1}\circ f\circ h$（其中 $h(x)=x-x^*$）。

\key{（2）一般情形：$\lambda=f'(x^*)$ 是线性化乘子。}
即使 $f$ 非线性，在不动点附近也有
\[
  b_{n+1}=f(x^*+b_n)-x^*=f'(x^*)b_n+O(b_n^2),
\]
所以 $\lambda=f'(x^*)$ 控制偏差的衰减／放大——这正是探究二稳定性判据的来源。
\key{不动点法生效的底层原因，就是「减去不动点」把非线性递推线性化为等比。}

\key{（3）莫比乌斯分式：配 $\dfrac{a_n-p}{a_n-q}$。}
对 $a_{n+1}=\dfrac{aa_n+b}{ca_n+d}$，设两个相异不动点 $x_1,x_2$。计算偏差：
\[
  a_{n+1}-x_1=\frac{(ad-bc)(a_n-x_1)}{(ca_n+d)(cx_1+d)},\qquad
  a_{n+1}-x_2=\frac{(ad-bc)(a_n-x_2)}{(ca_n+d)(cx_2+d)} .
\]
两式相除，随 $n$ 变化的分母 $(ca_n+d)$ 恰好约去：
\[
  \frac{a_{n+1}-x_1}{a_{n+1}-x_2}
  =\frac{cx_2+d}{cx_1+d}\cdot\frac{a_n-x_1}{a_n-x_2} .
\]
于是 $\left\{\dfrac{a_n-x_1}{a_n-x_2}\right\}$ 是等比数列，公比为 $\dfrac{cx_2+d}{cx_1+d}$。
若只有一个重不动点 $x_0$，则取倒数：
\[
  \frac1{a_{n+1}-x_0}-\frac1{a_n-x_0}=\text{常数},
\]
化为等差数列。

\begin{example}
设 $a_1=2,\ a_{n+1}=\dfrac{2a_n+3}{a_n+4}$，求通项公式。
\end{example}
\begin{solution}
不动点方程 $x=\dfrac{2x+3}{x+4}\Rightarrow x^2+2x-3=0\Rightarrow x_1=1,\ x_2=-3$。
\[
  \frac{a_{n+1}-1}{a_{n+1}+3}
  =\frac{\dfrac{2a_n+3}{a_n+4}-1}{\dfrac{2a_n+3}{a_n+4}+3}
  =\frac{a_n-1}{5a_n+15}
  =\frac15\cdot\frac{a_n-1}{a_n+3} .
\]
故 $\left\{\dfrac{a_n-1}{a_n+3}\right\}$ 是首项 $\dfrac{2-1}{2+3}=\dfrac15$、公比 $\dfrac15$ 的等比数列：
\[
  \frac{a_n-1}{a_n+3}=\left(\frac15\right)^n
  \ \Rightarrow\
  a_n=\frac{5^n+3}{5^n-1} .
\]
检验：$n=1$ 时 $\dfrac{5+3}{5-1}=2$ $\surd$。
\end{solution}

\begin{conccard}[二级]{共轭变换：不动点为何生效}{一阶线性递推、分式（莫比乌斯）递推、嵌套迭代 $f^{(n)}(x^*)$}{把口诀「令 $a_{n+1}=a_n=x^*$」当成真的假设相邻项相等；双不动点时忘记检验 $x_1\ne x_2$；重不动点时仍强行作比}
一阶线性：$b_n=a_n-x^*$，$b_{n+1}=f'(x^*)b_n$（精确）；
一般情形：$b_{n+1}=f'(x^*)b_n+O(b_n^2)$，$\lambda=f'(x^*)$ 是线性化乘子；
莫比乌斯双不动点：$\dfrac{a_n-x_1}{a_n-x_2}$ 等比，公比 $\dfrac{cx_2+d}{cx_1+d}$；
重不动点：$\dfrac1{a_n-x_0}$ 等差。
\end{conccard}

\subsection{强基拓展：生成函数（超出高考考纲）}

\key{探究六：把整个数列打包成一个函数}

\key{温馨提示：本模块涉及形式幂级数，超出高考考纲，仅面向竞赛／强基与学有余力者。}

\key{（1）生成函数：序列 $\to$ 幂级数。}
定义
\[
  G(x)=\sum_{n\ge0}a_nx^n .
\]
它把无穷数列 $\{a_n\}$ 打包成一个形式幂级数：
\begin{itemize}
  \item 下标平移 $a_{n+1}$ 对应「除以 $x$ 并减去首项」：$\sum_{n\ge0}a_{n+1}x^n=\dfrac{G(x)-a_0}{x}$；
  \item 递推关系对应 $G(x)$ 的代数方程；
  \item 解出 $G(x)$ 再展开，$x^n$ 的系数就是 $a_n$。
\end{itemize}

\key{（2）例：斐波那契数列。}
设 $F_0=0,\ F_1=1,\ F_{n+2}=F_{n+1}+F_n$，$G(x)=\sum_{n\ge0}F_nx^n$。则
\[
\begin{aligned}
  G(x)&=x+\sum_{n\ge2}F_nx^n
  =x+\sum_{n\ge2}(F_{n-1}+F_{n-2})x^n\\
  &=x+x\sum_{n\ge2}F_{n-1}x^{n-1}+x^2\sum_{n\ge2}F_{n-2}x^{n-2}
  =x+xG(x)+x^2G(x) .
\end{aligned}
\]
故
\[
  G(x)=\frac{x}{1-x-x^2}
  =\frac{x}{(1-\varphi x)(1-\psi x)},\qquad
  \varphi=\frac{1+\sqrt5}{2},\ \psi=\frac{1-\sqrt5}{2} .
\]
部分分式分解：
\[
  G(x)=\frac1{\sqrt5}\left(\frac1{1-\varphi x}-\frac1{1-\psi x}\right)
  =\frac1{\sqrt5}\sum_{n\ge0}(\varphi^n-\psi^n)x^n ,
\]
比较系数即得比内公式
\[
  F_n=\frac{\varphi^n-\psi^n}{\sqrt5} .
\]
同一个 $G(x)$ 还能给出组合意义（第 7 章）：$\dfrac1{1-x-x^2}$ 的系数是「用 $1\times1$ 与 $1\times2$ 骨牌铺满 $1\times n$ 棋盘」的方法数。

\key{（3）适用边界。}
生成函数擅长：线性递推（齐次／非齐次）、卷积型求和、组合计数（卡特兰数等）。
它把「解递推」变成「解代数方程 $+$ 展开」，是强基与竞赛的核武器；
但解答题若用此法，需补写合法的代数推导。

\begin{conccard}[三级]{生成函数：递推 $\to$ 代数方程}{斐波那契等线性递推、非齐次递推、卷积与组合计数}{超出高考考纲，解答题不能直接套用；忘记首项 $a_0$ 导致平移错误；展开时忽略收敛域（形式幂级数只需形式运算）}
$G(x)=\sum_{n\ge0}a_nx^n$；$\sum_{n\ge0}a_{n+1}x^n=\dfrac{G(x)-a_0}{x}$；
由递推导出 $G(x)$ 的封闭形式，再展开读系数；
斐波那契：$G(x)=\dfrac{x}{1-x-x^2}$，$F_n=\dfrac{\varphi^n-\psi^n}{\sqrt5}$。
\end{conccard}

\begin{method}[变换—求解—逆变换：本章五件降维武器]
\begin{enumerate}
  \item \key{换坐标系}：看到 $a_{n+1}=f(a_n)$，先解不动点 $x^*=f(x^*)$，作 $b_n=a_n-x^*$；
  \item \key{升维}：看到高阶线性递推，写成状态向量 $\mathbf v_{n+1}=A\mathbf v_n$，求特征值；
  \item \key{互逆}：看到求和，想差分 $\Delta$ 与裂项；看到乘积求和，想 Abel 变换；
  \item \key{打包}：递推复杂时，用生成函数 $G(x)$ 把它变成代数方程（强基）；
  \item \key{还原}：在新坐标系算完后，用逆变换 $h^{-1}$ 把答案翻回原数列，并检验首项与边界。
\end{enumerate}
\end{method}

\begin{summary}
本章把数列从「套公式」升级为「离散动力系统 $+$ 降维变换」：
差分 $\Delta$ 与求和 $\sum$ 互为逆运算（离散版微积分基本定理），由 $S_n$ 反求 $a_n$、差分降次判断次数；
一阶递推 $a_{n+1}=f(a_n)$ 是不动点迭代，$x^*=f(x^*)$ 是平衡态，蛛网图把迭代画成阶梯，$|f'(x^*)|<1$ 吸引、$>1$ 排斥；
二阶线性递推对应转移矩阵 $\begin{pmatrix}0&1\\q&p\end{pmatrix}$，特征方程 $\lambda^2-p\lambda-q=0$ 决定通项结构，特征方向即分式递推的不动点方向；
Abel 变换是离散分部积分，把「等差 $\times$ 等比」求和机械化；
共轭变换（换坐标系）解释了不动点法为何生效：$b_n=a_n-x^*$ 把非线性递推线性化为等比；
生成函数（强基拓展）把递推打包成代数方程。

\begin{center}
\footnotesize
\setlength{\tabcolsep}{4pt}
\renewcommand{\arraystretch}{1.35}
\begin{tabular}{p{2.1cm}p{4.9cm}p{4.2cm}p{1.9cm}}
\toprule
降维变换 & 核心操作 & 把什么变成什么 & 对应探究 \\
\midrule
差分 $\Delta$ & $\Delta a_n=a_{n+1}-a_n$ & 求和 $\to$ 相消；高次 $\to$ 低次 & 探究一 \\
共轭（平移／分式） & $b_n=a_n-x^*$；$\dfrac{a_n-x_1}{a_n-x_2}$ & 拉伸$+$平移／非线性分式 $\to$ 等比／等差 & 探究二、五 \\
特征值对角化 & $\mathbf v_{n+1}=A\mathbf v_n$ & 递推 $\to$ 特征值幂 & 探究三 \\
Abel 变换 & $\sum a_kb_k=a_nB_n-\sum\Delta a_kB_k$ & 乘积求和 $\to$ 差分 $\times$ 部分和 & 探究四 \\
生成函数 & $G(x)=\sum a_nx^n$ & 递推 $\to$ 代数方程 & 探究六 \\
\bottomrule
\end{tabular}
\end{center}

\key{数列不是公式的仓库，而是离散动力系统的轨迹；解题的本质是换坐标系，把陌生递推化为等比／等差／代数方程。}
\end{summary}

\end{explore}

% ---------- 跨章联动：降维双引擎（正文收尾；原在答案区之后，已前移） ----------
\subsection{跨章方法论大一统：平移消常数，旋转消交叉项}

前面第 4 章用\key{仿射}压缩椭圆回圆、第 8 章用\key{齐次化}处理条件最值，
本章用\key{不动点}处理递推。它们看似分属三个板块，底层却是同一个动作：
\key{当前的坐标原点没选好，方程里才拖着讨厌的常数尾巴}。把原点精准搬到
锚点上，非齐次项会瞬间蒸发。

\begin{conccard}[二级]{换坐标系的两步走}{适用于所有「带常数项」「带一次项」「带交叉项」的结构}{只背结论，不问原点该搬去哪}
第一步\key{平移}（$x'=x-x_0,\ y'=y-y_0$）：把原点搬到不动点、对称中心或定点 $P$，利用「锚点自身满足方程」这一性质抵消常数项与一次项（第 8 章「手术刀之二」）；第二步\key{旋转}（乘以 $e^{i\theta}$）：把坐标轴对齐到主轴／特征方向，消灭二次交叉项 $Bxy$，并把复特征根的螺旋迭代还原为实数周期震荡。
\end{conccard}

\begin{center}
\renewcommand{\arraystretch}{1.4}
\begin{tabular}{|p{2.8cm}|p{5.0cm}|p{5.4cm}|}
\hline
\textbf{对比维度} & \textbf{数列：不动点法} & \textbf{条件最值／圆锥曲线：齐次化（第 8 章）} \\
\hline
非齐次障碍 & 递推拖着常数尾巴：$a_{n+1}=pa_n+q$ & 平移到定点后仍剩一次项与常数项 \\
\hline
锚点满足的方程 & $x^*=px^*+q$ & $\dfrac{x_0^2}{a^2}+\dfrac{y_0^2}{b^2}=1$ \\
\hline
平移后发生了什么 & 令 $b_n=a_n-x^*$，得 $b_{n+1}=pb_n$，$q$ 归零 & 常数项化为 $\dfrac{x_0^2}{a^2}+\dfrac{y_0^2}{b^2}-1=0$，自动抵消 \\
\hline
剩下的纯粹结构 & 只含公比 $p$ 的等比缩放 & 只含斜率 $k=\dfrac{y'}{x'}$ 的一元二次方程 \\
\hline
\end{tabular}
\end{center}

\begin{problem}[旋转 $45^\circ$：把 $xy=1$ 还原成标准双曲线]
证明反比例函数 $xy=1$ 的图象是等轴双曲线，并写出其焦点位置。
\end{problem}
\begin{solution}
将坐标轴逆时针旋转 $45^\circ$，即
\[
  x=\frac{\sqrt2}{2}(x'-y'),\qquad y=\frac{\sqrt2}{2}(x'+y') .
\]
代入 $xy=1$：
\[
  \frac{\sqrt2}{2}(x'-y')\cdot\frac{\sqrt2}{2}(x'+y')=1
  \ \Longrightarrow\
  \frac{x'^2}{2}-\frac{y'^2}{2}=1 .
\]
交叉项 $xy$ 变成了平方差：这正是以 $x'$ 轴为实轴的等轴双曲线，半实轴 $a=b=\sqrt2$，故 $c=\sqrt{a^2+b^2}=2$，焦点在新坐标系的 $(\pm2,0)$，也就是原坐标系的 $(\sqrt2,\sqrt2)$ 与 $(-\sqrt2,-\sqrt2)$。
\end{solution}

\begin{center}
  \begin{tikzpicture}[x=1.05cm,y=1.05cm,>=stealth,font=\small]
    \draw[->,white!55] (-3,0)--(3,0) node[right,white!70] {$x$};
    \draw[->,white!55] (0,-3)--(0,3) node[above,white!70] {$y$};
    % 旋转后的主轴
    \draw[->,second,dashed,very thick] (-2.3,-2.3)--(2.5,2.5) node[above right,second] {\small $x'$（新实轴）};
    \draw[->,second,dashed,very thick] (2.3,-2.3)--(-2.5,2.5) node[above left,second] {\small $y'$（新虚轴）};
    \draw[->,second,thick] (1.1,0) arc (0:45:1.1);
    \node[second] at (1.15,0.38) {\small $45^\circ$};
    % 双曲线 xy=1 的两支
    \draw[main,very thick,domain=0.42:2.6,samples=80,smooth] plot(\x,{1/\x});
    \draw[main,very thick,domain=-2.6:-0.42,samples=80,smooth] plot(\x,{1/\x});
    \node[main] at (1.9,1.15) {\small $xy=1$};
    \node[white!70] at (0,-3.4) {\small 交叉项 $xy$ 与旋转 $45^\circ$ 后的新主轴 $x',y'$};
  \end{tikzpicture}\\[0.3em]
  {\small\itshape 同一曲线的两种读法：原坐标系下是交叉项方程 $xy=1$；把轴转 $45^\circ$ 后成为标准等轴双曲线 $\frac{x'^2}{2}-\frac{y'^2}{2}=1$}
\end{center}

\paragraph{第二步：旋转也管数列。} 当二阶递推 $a_{n+2}=pa_{n+1}+qa_n$ 的特征方程判别式 $\Delta<0$ 时，两根为共轭复数 $\lambda_{1,2}=r(\cos\theta\pm i\sin\theta)$。此时几何图像不再是沿直线趋近不动点，而是每步\key{缩放 $r$ 倍并旋转 $\theta$ 角}的对数螺线；通项写成
\[
  a_n=r^{\,n-1}\bigl(c_1\cos n\theta+c_2\sin n\theta\bigr) .
\]
$r=1$ 且 $\theta=\frac{2\pi}{T}$ 时数列退化为纯周期，例如 $a_{n+2}=a_{n+1}-a_n$ 的 $\lambda=e^{\pm i\frac{\pi}{3}}$，周期为 $6$。这与「把轴旋转到主轴方向」是同一件事：都是把耦合项转没。

\begin{figure}[htbp]
  \centering
  \begin{tikzpicture}[>=Stealth, semithick, font=\footnotesize, x=1.15cm, y=1.15cm]
    \draw[->, gray] (-1.9,0)--(1.95,0);
    \draw[->, gray] (0,-1.35)--(0,1.55);
    % 迭代点 v0..v5（M = r R(40°)，r=0.8），相邻两点长度比 0.8、辐角差 40°
    \coordinate (v0) at (1.5,0.2);
    \coordinate (v1) at (0.8163,0.8942);
    \coordinate (v2) at (0.0402,0.9679);
    \coordinate (v3) at (-0.4733,0.6138);
    \coordinate (v4) at (-0.6058,0.1327);
    \coordinate (v5) at (-0.4394,-0.2302);
    \fill (0,0) circle (1.2pt) node[below right] {$O$};
    \draw[main, very thick] (v0)--(v1)--(v2)--(v3)--(v4)--(v5);
    \foreach \i/\p/\pos in {0/v0/above right,1/v1/right,2/v2/above left,3/v3/left,4/v4/below left,5/v5/below}
      \fill[second] (\p) circle (1.5pt) node[\pos] {$\mathbf v_{\i}$};
    \node[gray!70!black, align=center] at (0,-1.15) {$M=rR(\theta)$，$r=0.8,\ \theta=40^\circ$：每步长度乘 $r$、方向加 $\theta$；\\$r<1$ 螺线收拢到唯一不动点 $O$，$r>1$ 远离，$r=1$ 且 $\theta=2\pi/T$ 时为周期 $T$};
  \end{tikzpicture}
  \caption{复特征值的迭代图像：$v_{n+1}=Mv_n$ 且 $M=rR(\theta)$ 时，每一步都把向量缩放 $r$ 倍并旋转 $\theta$ 角，点列沿螺线运动——这就是「实部控伸缩、虚部控旋转」的离散版本}
  \label{fig:spiral-complex}
\end{figure}

\begin{commonerror}{卷面上禁止出现「由仿射变换知」「由极点极线知」}
高观点只用于草稿纸预判答案与取等位置，答题卡上必须写考纲认可的初等推导（联立、韦达定理、作差构造等比）。详见附录~\ref{app:legal}。
\end{commonerror}

\begin{inclasstest}
\begin{enumerate}
  \item 设 $S_n=2n^2-n$，求数列 $\{a_n\}$ 的通项公式；\level{易}
  \item 设 $a_1=1,\ a_{n+1}=2a_n+3$，求通项并判断 $\{a_n\}$ 是否收敛；\level{中}
  \item 用 Abel 变换求 $\sum_{k=1}^nk\cdot2^k$；\level{中}
  \item 求 $a_{n+2}=4a_{n+1}-4a_n$ 在 $a_1=1,\ a_2=4$ 下的通项（重根情形）。\level{中}
\end{enumerate}
\end{inclasstest}

\begin{afterclass}
\subsection*{基础巩固}
\begin{enumerate}
  \item 已知 $\Delta a_n=2n+1$，且 $a_1=1$，求 $a_n$；
  \item 证明 $\sum_{k=1}^n\Delta a_k=a_{n+1}-a_1$，并用它求 $\sum_{k=1}^n\dfrac1{k(k+1)}$。
\end{enumerate}

\subsection*{进阶提升}
\begin{enumerate}
  \item 用特征值法求 $a_{n+2}=2a_{n+1}+3a_n$ 在 $a_1=1,\ a_2=5$ 下的通项；
  \item 设 $a_1=0,\ a_{n+1}=\dfrac{3a_n+1}{a_n+3}$，求通项公式。
\end{enumerate}

\subsection*{拓展探究}
\begin{enumerate}
  \item （强基／竞赛，超出高考考纲）设 $a_0=1,\ a_{n+1}=2a_n+1$，用生成函数求通项公式。
\end{enumerate}

\subsection*{教师版答案}
\begin{solution}
\textbf{基础 1：}由离散版微积分基本定理，
\[
  a_n=a_1+\sum_{k=1}^{n-1}\Delta a_k
  =1+\sum_{k=1}^{n-1}(2k+1)
  =1+2\cdot\frac{(n-1)n}2+(n-1)=n^2 .
\]
（检验：$n=1$ 时 $a_1=1$。）
\end{solution}
\begin{solution}
\textbf{基础 2：}求和时中间项全部抵消：
\[
  \sum_{k=1}^n\Delta a_k=\sum_{k=1}^n(a_{k+1}-a_k)=a_{n+1}-a_1 .
\]
又 $\dfrac1{k(k+1)}=\dfrac1k-\dfrac1{k+1}=\Delta\!\left(-\dfrac1k\right)$，故
\[
  \sum_{k=1}^n\frac1{k(k+1)}
  =\sum_{k=1}^n\Delta\!\left(-\frac1k\right)
  =\left(-\frac1{n+1}\right)-(-1)=\frac n{n+1} .
\]
\end{solution}
\begin{solution}
\textbf{进阶 1：}$A=\begin{pmatrix}0&1\\3&2\end{pmatrix}$，特征方程
\[
  \lambda^2-2\lambda-3=0\Rightarrow\lambda_1=3,\ \lambda_2=-1 .
\]
设 $a_n=c_1\cdot3^{n-1}+c_2\cdot(-1)^{n-1}$，由 $a_1=1,\ a_2=5$：
\[
  c_1+c_2=1,\quad 3c_1-c_2=5
  \ \Rightarrow\ c_1=\frac32,\ c_2=-\frac12 .
\]
故
\[
  a_n=\frac{3^n+(-1)^n}2 .
\]
（检验：$a_3=2\cdot5+3\cdot1=13=\frac{27-1}2$。）
\end{solution}
\begin{solution}
\textbf{进阶 2：}不动点方程 $x=\dfrac{3x+1}{x+3}\Rightarrow x^2=1\Rightarrow x_1=1,\ x_2=-1$。
\[
  \frac{a_{n+1}-1}{a_{n+1}+1}
  =\frac{(3a_n+1)-(a_n+3)}{(3a_n+1)+(a_n+3)}
  =\frac{2a_n-2}{4a_n+4}
  =\frac12\cdot\frac{a_n-1}{a_n+1} .
\]
首项 $\dfrac{0-1}{0+1}=-1$，故 $\dfrac{a_n-1}{a_n+1}=-\left(\dfrac12\right)^{n-1}$，解得
\[
  a_n=\frac{1-\left(\frac12\right)^{n-1}}{1+\left(\frac12\right)^{n-1}}
  =\frac{2^{n-1}-1}{2^{n-1}+1} .
\]
（检验：$n=1$ 时 $a_1=0$；$n=2$ 时 $a_2=\frac13=\frac{3\cdot0+1}{0+3}$。）
\end{solution}
\begin{solution}
\textbf{拓展：}设 $G(x)=\sum_{n\ge0}a_nx^n$。由 $a_{n+1}=2a_n+1$：
\[
  \frac{G(x)-a_0}{x}=2G(x)+\sum_{n\ge0}x^n=2G(x)+\frac1{1-x} .
\]
代入 $a_0=1$：
\[
  G(x)-1=2xG(x)+\frac{x}{1-x}
  \ \Rightarrow\
  G(x)(1-2x)=\frac1{1-x}
  \ \Rightarrow\
  G(x)=\frac1{(1-x)(1-2x)} .
\]
部分分式：
\[
  G(x)=\frac2{1-2x}-\frac1{1-x}
  =\sum_{n\ge0}(2\cdot2^n-1)x^n
  =\sum_{n\ge0}(2^{n+1}-1)x^n ,
\]
故 $a_n=2^{n+1}-1$。（检验：$a_0=1,\ a_1=3$。）
\end{solution}

\reading{拓展阅读 11}{数列：离散动力系统的轨迹}
\end{afterclass}

% 高观点导学案：导数与局部线性化引擎（第 6 章·高观点版）
% 主线：局部线性化（以直代曲）→ 导数即变化率（差分的连续版）→ 五大不等式的结构本质（支撑边界）
%       → 切线放缩与隐零点 → 条件最值四大标准动作 → 含参恒成立/能成立的集合边界
% 核心思想：把「导数」与「不等式」统一在「以直代曲」这一个引擎下——曲线在一点附近用切线代替，
%           不等式是为点集划定不可逾越的边界；解题四步：识别结构 → 局部线性化 → 求最值 → 验取等。

\section{导数与局部线性化引擎（高观点版）}

\begin{methodbox}{多框架对读}
\small
\begin{center}
\begin{tabular}{c|p{2.9cm}|p{5.0cm}|p{3.1cm}}
\toprule
视角 & 本章看点 & 别的视角怎么说 & 哪侧最省力 \\
\midrule
代数 & 求导法则、恒等变形 & 符号对象的机械运算 & 硬算单调区间 \\
几何 & 切线、增减、凹凸的形状 & 曲线在放大镜下变直线 & 看图猜极值 \\
集合 & 零点集、符号区间 & 定义域值域的边界 & 界参数范围 \\
映射 & 局部线性化：曲线 $\to$ 直线 & 换到切线坐标系再比大小 & 五不等式一源 \\
\bottomrule
\end{tabular}
\end{center}
卡的使用：证不等式先想「是不是某个函数的单调性」，再把切线当替代函数（总图见第 0 章图~\ref{fig:four-view-map}）。
\end{methodbox}

第 1 章确立了函数与单调性，第 5 章把差分 $\Delta$ 确立为离散意义下的变化率；
本章把「导数」与「不等式」统一在同一个引擎下——\key{局部线性化（以直代曲）}：
任何光滑曲线在一点的无穷小邻域内都可以用它的切线近似代替，而这条切线的斜率就是导数。
于是「求导数、判符号」成为判定单调性／极值／最值的统一操作；
五大核心不等式分别是射影定理、内积公理、齐次化（完整展开见第 8 章）、度量三角公理与凸凹函数切线支撑这五种结构下的「支撑边界」；
条件最值与含参成立性问题，则归结为「边界值能否取到」。
六个模块层层展开：
\key{局部线性化}（放大镜下曲成直）$\to$
\key{导数即变化率}（差分的连续版，统一判定单调性／极值／最值）$\to$
\key{五大不等式的结构本质}（一张对比表：射影、内积、齐次、度量、凸支撑）$\to$
\key{切线放缩与隐零点}（$e^x\ge x+1$、$\ln x\le x-1$ 的严格证明与多步取等检验）$\to$
\key{条件最值四大标准动作}（1 的代换、消元主元、拆项分离、切线放缩）$\to$
\key{含参恒成立／能成立}（点集位置、集合包含与量词的值域边界）。
核心思想：\key{导数与不等式是同一个引擎的两面——以直代曲给出局部逼近，支撑边界给出整体控制；解题的本质是「识别结构 $\to$ 局部线性化 $\to$ 求最值 $\to$ 验取等」。}

\begin{learnobjective}
\begin{itemize}
  \item 理解局部线性化公式 $f(x)\approx f(x_0)+f'(x_0)(x-x_0)$ 与「以直代曲」的几何含义，会用切线作近似计算；同一个一阶逼近在第 5 章的离散递推中就是线性化乘子 $f'(x^*)$；
  \item 掌握导数的极限定义与导函数，理解它是差分的连续版本（$\Delta a_n$ 是 $a_{n+1}-a_n$，$f'(x)$ 是 $\frac{\Delta y}{\Delta x}$ 的极限）；
  \item 会用导数符号统一判定单调性、极值与最值，理解驻点与极值点的区别；
  \item 从射影定理、内积公理、齐次化、度量公理、凸凹性五个结构本质理解五大不等式，并记住各自的取等条件（齐次化专论见第 8 章）；
  \item 会用作差求导严格证明 $e^x\ge x+1$、$\ln x\le x-1$，能在多步放缩中检验各步取等条件是否同时成立；
  \item 掌握隐零点「设而不求、整体代换」的基本操作；
  \item 掌握条件最值的四大标准动作（1 的代换、消元主元、拆项分离常数、切线放缩），会检验取等的可行性；
  \item 会用点集上下位置与集合包含关系统一解释恒成立／能成立问题，掌握量词到最值的四类转化与含参二次不等式的分类讨论树。
\end{itemize}
\end{learnobjective}

\begin{preview}
\subsection*{知识回顾}
\begin{itemize}
  \item 函数 $y=f(x)$ 在区间 $[x_0,x_0+\Delta x]$ 上的平均变化率是 \blank；
  \item 导数的定义：$f'(x_0)=$ \blank（写成 $\Delta x\to0$ 时的极限）；
  \item 若 $f'(x)>0$ 在区间 $I$ 上恒成立，则 $f$ 在 $I$ 上 \blank；
  \item 基本不等式：$a,b>0$ 时 \blank $\ge$ \blank，当且仅当 \blank 时取等；
  \item 一元二次方程 $ax^2+bx+c=0$ 的判别式 $\Delta=$ \blank，$\Delta>0$ 时有两个 \blank 实根。
\end{itemize}

\subsection*{教材速读}
阅读「导数」一章教材，完成下列问题：
\begin{enumerate}
  \item 可导与连续的关系是：可导 \blank 连续，连续 \blank 可导（填「一定」或「不一定」）；
  \item 写出三个基本求导公式：$(x^n)'=$ \blank，$(e^x)'=$ \blank，$(\ln x)'=$ \blank；
  \item 函数极值的第一充分条件：在 $x_0$ 两侧 $f'(x)$ 异号，则 $x_0$ 是 \blank；
  \item 闭区间上连续函数的最值只可能在 \blank 与 \blank 处取得。
\end{enumerate}
\end{preview}

\begin{explore}

\subsection{局部线性化：放大镜下曲成直}

\key{探究一：切线是曲线在一点处的「最佳直线近似」}

\key{（1）从平均变化率到瞬时变化率。}
函数在 $x_0$ 附近的平均变化率是 $\dfrac{\Delta y}{\Delta x}=\dfrac{f(x_0+\Delta x)-f(x_0)}{\Delta x}$。
让 $\Delta x\to0$，若极限存在，就得到导数
\[
  f'(x_0)=\lim_{\Delta x\to0}\frac{f(x_0+\Delta x)-f(x_0)}{\Delta x}.
\]
这正是第 5 章差分 $\Delta a_n=a_{n+1}-a_n$ 的连续版：把「步长为 1 的差分」换成「步长趋于 0 的差商极限」。

\key{（2）局部线性化公式：以直代曲。}
过点 $P(x_0,f(x_0))$ 的切线方程为
\[
  y=f(x_0)+f'(x_0)(x-x_0),
\]
而在 $x_0$ 附近，曲线与切线的误差是关于 $(x-x_0)$ 的高阶无穷小：
\[
  f(x)=f(x_0)+f'(x_0)(x-x_0)+o(x-x_0).
\]
也就是说，\key{在 $x_0$ 的无穷小邻域内，曲线可以用切线代替——这就是「以直代曲」}。
几何显微镜放得越大，曲线看起来越像一条直线（图 \ref{fig:local-linear}）。

\begin{figure}[htbp]
  \centering
  \begin{tikzpicture}[>=Stealth, semithick, font=\footnotesize]
    % ---------- 左：全局视角 ----------
    \begin{scope}
      \draw[->, gray] (-1.3,0)--(3.0,0) node[below right] {$x$};
      \draw[->, gray] (0,-2.3)--(0,4.4) node[above left] {$y$};
      \draw[main!70!black, very thick, domain=-1:2, smooth, samples=60] plot (\x, {\x*\x});
      \node[main!70!black] at (1.72,3.15) {$y=x^2$};
      \draw[second, thick, domain=-0.4:2.5] plot (\x, {2*\x-1});
      \node[second] at (2.64,3.9) {$y=2x-1$};
      \fill (1,1) circle (1.6pt) node[above right] {$P(1,1)$};
      \draw[dashed, gray] (0.85,0.7) rectangle (1.15,1.3);
      % 说明文字整体右移，避开左图的 $y$ 轴
      \node[gray!70!black, align=left] at (0.2,-2.15) {虚线框为右图的放大区域};
    \end{scope}
    % ---------- 右：放大视角 ----------
    \begin{scope}[xshift=5.8cm]
      \draw[->, gray] (0.78,0.55)--(1.26,0.55) node[above right] {$x$};
      \draw[->, gray] (0.75,0.58)--(0.75,1.46) node[above left] {$y$};
      \draw[second, thick, dashed, domain=0.84:1.16] plot (\x, {2*\x-1});
      \draw[main!70!black, very thick, domain=0.84:1.16, smooth, samples=60] plot (\x, {\x*\x});
      \fill (1,1) circle (1.6pt) node[above right] {$P(1,1)$};
      \node[gray!70!black, align=center] at (1.01,-0.1) {放大约 $25$ 倍后，曲线（实线）\\与切线（虚线）几乎重合};
    \end{scope}
  \end{tikzpicture}
  \caption{以直代曲：抛物线 $y=x^2$ 在 $P(1,1)$ 处的切线为 $y=2x-1$。左图全局视角下曲线明显弯曲；右图把视野放大到 $P$ 附近，曲线与切线在视觉上无法区分——$f(x)\approx f(x_0)+f'(x_0)(x-x_0)$}
  \label{fig:local-linear}
\end{figure}

\begin{example}
用切线近似求 $\sqrt{4.1}$ 与 $\sin31^\circ$（精确到 $0.001$）。
\end{example}
\begin{solution}
取 $f(x)=\sqrt x$，在 $x_0=4$ 处：$f(4)=2$，$f'(x)=\frac1{2\sqrt x}$，$f'(4)=\frac14$，故
\[
  \sqrt{4.1}\approx 2+\frac14\cdot0.1=2.025 .
\]
（真值 $\sqrt{4.1}\approx2.0248$，误差约 $1.5\times10^{-4}$。）

取 $g(x)=\sin x$，在 $x_0=\frac\pi6$ 处：$g'(x)=\cos x$，$g'\!\left(\frac\pi6\right)=\frac{\sqrt3}{2}$，自变量增量
$\Delta x=\frac\pi{180}$，故
\[
  \sin31^\circ\approx \frac12+\frac{\sqrt3}{2}\cdot\frac{\pi}{180}
  \approx0.5+0.01511=0.515 .
\]
（真值 $\sin31^\circ\approx0.51504$。）
\end{solution}

\begin{conccard}[一级]{局部线性化：以直代曲}{光滑曲线的近似计算、误差估计、理解导数的几何意义}{把「近似」当成「相等」；放大倍数不够时强行替代；忘记公式只在 $x_0$ 附近有效}
$f'(x_0)=\lim\limits_{\Delta x\to0}\dfrac{f(x_0+\Delta x)-f(x_0)}{\Delta x}$；
切线 $y=f(x_0)+f'(x_0)(x-x_0)$；
$f(x)=f(x_0)+f'(x_0)(x-x_0)+o(x-x_0)$，即以直代曲；
导数是差分 $\Delta a_n$ 的连续版本。
\end{conccard}

\subsection{导数即变化率：统一判定单调性、极值、最值}

\key{探究二：$f'(x)$ 的正负号就是曲线的「上下坡」}

\key{（1）导函数与符号判定。}
把每一点的导数收集起来得到导函数 $f'(x)$。在区间 $I$ 上：
\begin{itemize}
  \item $f'(x)>0$ 恒成立 $\iff$ $f$ 在 $I$ 上单调递增；
  \item $f'(x)<0$ 恒成立 $\iff$ $f$ 在 $I$ 上单调递减。
\end{itemize}
这与第 5 章用差分判数列单调性完全同构：$\Delta a_n>0\iff a_{n+1}>a_n$（递增数列）。

\key{（2）极值：$f'$ 变号的点。}
费马引理：可导函数在极值点处必有 $f'(x_0)=0$（驻点）。
第一充分条件：$f'$ 在 $x_0$ 两侧异号 $\iff$ $x_0$ 是极值点。汇总如下：

\begin{center}\footnotesize\renewcommand{\arraystretch}{1.3}
\begin{tabular}{c|c|c}
$f'(x)$ 在 $x_0$ 附近的符号 & $f(x)$ 的形态 & 极值 \\\hline
$+$，然后 $-$ & 先增后减 & 极大值 \\
$-$，然后 $+$ & 先减后增 & 极小值 \\
不变号（如 $y=x^3$ 在 $0$ 处） & 单调 & 无极值 \\
\end{tabular}
\end{center}

\key{（3）最值：端点 $+$ 极值点。}
闭区间 $[a,b]$ 上的连续函数，最值只可能在端点 $a,b$ 与极值点处取得，比较这些点的函数值即可。

\begin{example}
求 $f(x)=x^3-3x$ 的单调区间与极值；若在 $[-2,2]$ 上求最值。
\end{example}
\begin{solution}
$f'(x)=3x^2-3=3(x+1)(x-1)$。
\begin{itemize}
  \item $x<-1$ 或 $x>1$ 时 $f'(x)>0$；$-1<x<1$ 时 $f'(x)<0$。
  \item 故 $f$ 在 $(-\infty,-1)$、$(1,+\infty)$ 上递增，在 $(-1,1)$ 上递减。
  \item 极大值 $f(-1)=-1+3=2$；极小值 $f(1)=1-3=-2$。
\end{itemize}
在 $[-2,2]$ 上额外比较端点：$f(-2)=-8+6=-2$，$f(2)=8-6=2$。
故最小值为 $-2$（在 $x=-2$ 与 $x=1$ 处），最大值为 $2$（在 $x=-1$ 与 $x=2$ 处）。
$\surd$
\end{solution}

\begin{commonerror}{驻点、极值点与最值的三个高频失误}
\paragraph{失误一：以为 $f'(x_0)=0$ 就一定是极值点。}
$f'(x_0)=0$ 只是必要条件。反例：$f(x)=x^3$ 在 $x=0$ 处 $f'(0)=0$，但两侧 $f'>0$ 不变号，不是极值点。
必须检验两侧符号是否改变。
\paragraph{失误二：忽略定义域。}
例如 $f(x)=x\ln x$ 的定义域是 $(0,+\infty)$，讨论单调性时不能把 $x\le0$ 也算进去；
求闭区间最值时也不能漏掉端点。
\paragraph{失误三：把「单点处 $f'(x_0)>0$」当成「区间上递增」。}
导数符号必须在整个区间上保持，才能推出单调性；个别点的符号说明不了整体趋势。
\end{commonerror}

\begin{conccard}[一级]{导数符号法：单调性、极值、最值一条龙}{含参数的函数性质讨论、恒成立问题中求最值}{驻点当成极值点；漏定义域；漏端点；把 $f'(x)\ge0$ 中的孤立零点误判为单调性破坏}
$f'(x)>0\Rightarrow$ 递增，$f'(x)<0\Rightarrow$ 递减；
极值点处 $f'(x_0)=0$ 且两侧异号；
闭区间最值 $=$ 端点值与极值值中的最大／最小者。
\end{conccard}

\subsection{五大不等式的结构本质}

\key{探究三：五条不等式，五个「支撑边界」}

同一个「$\ge$」在不同结构下有不同来历。看懂它们的结构本质，取等条件就会自己跳出来。

\begin{center}\footnotesize\setlength{\tabcolsep}{4pt}\renewcommand{\arraystretch}{1.4}
\begin{tabular}{p{1.9cm}p{3.5cm}p{5.2cm}p{2.6cm}}
\toprule
不等式 & 核心形式 & 结构本质 & 取等条件 \\
\midrule
基本不等式 & $\dfrac{a+b}{2}\ge\sqrt{ab}\ (a,b>0)$ & 半圆射影定理：直径 $a+b$ 的半径 $\ge$ 斜边上的高 $\sqrt{ab}$ & $a=b$ \\[2pt]
柯西不等式 & $(a_1^2+a_2^2)(b_1^2+b_2^2)\ge(a_1b_1+a_2b_2)^2$ & 向量内积公理 $|\vec u|^2|\vec v|^2\ge(\vec u\cdot\vec v)^2$（即 $|\cos\theta|\le1$） & $\vec u\parallel\vec v$（共线） \\[2pt]
权方和不等式 & $\dfrac{a_1^2}{b_1}+\dfrac{a_2^2}{b_2}\ge\dfrac{(a_1+a_2)^2}{b_1+b_2}\ (b_i>0)$ & 柯西在 $\left(\dfrac{a_i}{\sqrt{b_i}}\right)$ 与 $(\sqrt{b_i})$ 下的齐次化特例 & $\dfrac{a_1}{b_1}=\dfrac{a_2}{b_2}$ \\[2pt]
绝对值三角不等式 & $\bigl||a|-|b|\bigr|\le|a\pm b|\le|a|+|b|$ & 度量空间的距离三角不等式（三角形两边之和 $\ge$ 第三边） & 共线同向（右端）／反向（左端） \\[2pt]
切线放缩 & $e^x\ge x+1$，$\ln x\le x-1$ & 凸／凹函数的切线支撑边界（Jensen 不等式的一阶情形） & $x=0$／$x=1$ \\
\bottomrule
\end{tabular}
\end{center}

注：齐次化作为第一性原理的完整方法链见第 8 章；本章只把它用作
条件最值的一个入口，权方和则是柯西（内积公理）的齐次化特例。

\key{（1）基本不等式 $=$ 半圆射影定理。}
作直径 $AB=a+b$ 的半圆，圆心 $O$，则半径 $R=\dfrac{a+b}{2}$。
设分界点 $C$ 满足 $AC=a$、$CB=b$，过 $C$ 作 $CP\perp AB$ 交半圆于 $P$。
因为直径所对圆周角为直角，$\triangle APB$ 是直角三角形，由射影定理
\[
  PC^2=AC\cdot CB=ab\ \Rightarrow\ PC=\sqrt{ab}.
\]
又在直角三角形 $OPC$ 中 $OP$ 是斜边，故 $PC\le OP=R$，即
\[
  \sqrt{ab}\le\frac{a+b}{2},
\]
当且仅当 $C$ 与 $O$ 重合（即 $a=b$）时取等（图 \ref{fig:amgm-semicircle}）。

\begin{figure}[htbp]
  \centering
  \begin{tikzpicture}[>=Stealth, semithick, font=\footnotesize, scale=1.05]
    % 直径 AB=a+b，取 a=2.4,b=1.2 => R=1.8，C 位于 x=0.6，PC=sqrt(2.88)≈1.697
    \draw[->, gray] (-2.35,0)--(2.35,0);
    \draw[thick] (-1.8,0)--(1.8,0);
    \draw[thick] (1.8,0) arc (0:180:1.8);
    \fill (-1.8,0) circle (1.5pt) node[below left] {$A$};
    \fill (1.8,0) circle (1.5pt) node[below right] {$B$};
    \fill (0,0) circle (1.5pt) node[below] {$O$};
    \fill (0.6,0) circle (1.5pt) node[below] {$C$};
    \fill (0.6,1.697) circle (1.5pt) node[above right] {$P$};
    \draw[main!70!black, thick] (0.6,0)--(0.6,1.697);
    % 标签放在 $PC$ 右侧、斜边 $PB$ 下方的空档处
    \node[main!70!black] at (0.95,0.55) {$\sqrt{ab}$};
    \draw[second, thick, dashed] (0,0)--(0,1.8);
    % 半径标签锚定在虚线右侧起排，避免压线
    \node[second, anchor=west] at (0.72,1.25) {$\dfrac{a+b}{2}$};
    \draw[third!70!black, thick] (-1.8,0)--(0.6,1.697)--(1.8,0);
    \node at (-0.6,-0.4) {$a$};
    \node at (1.2,-0.4) {$b$};
    \node[gray!70!black, align=center] at (0,-2.15) {$PC^2=AC\cdot CB=ab\Rightarrow PC=\sqrt{ab}\le OP=\dfrac{a+b}{2}$\\（当且仅当 $C$ 与 $O$ 重合即 $a=b$ 时取等）};
  \end{tikzpicture}
  \caption{半圆射影定理：直径 $AB=a+b$，半径 $R=\frac{a+b}{2}$ 永远是斜边上的高 $PC=\sqrt{ab}$ 的上界——这就是基本不等式的几何模型}
  \label{fig:amgm-semicircle}
\end{figure}

\key{（2）柯西 $=$ 向量内积公理；权方和 $=$ 柯西的齐次化。}
$|\vec u\cdot\vec v|=|\vec u|\,|\vec v|\,|\cos\theta|\le|\vec u|\,|\vec v|$，平方即得二维柯西，
取等当且仅当两向量共线。代数证法（判别式法）：
\begin{derivation}
令 $\varphi(x)=(a_1x+b_1)^2+(a_2x+b_2)^2\ge0$ 对一切实数 $x$ 成立，展开得
\[
  (a_1^2+a_2^2)x^2+2(a_1b_1+a_2b_2)x+(b_1^2+b_2^2)\ge0 .
\]
二次函数恒非负，故判别式 $\Delta\le0$：
\[
  4(a_1b_1+a_2b_2)^2-4(a_1^2+a_2^2)(b_1^2+b_2^2)\le0,
\]
即 $(a_1^2+a_2^2)(b_1^2+b_2^2)\ge(a_1b_1+a_2b_2)^2$。
取等当且仅当存在 $x$ 使 $a_1x+b_1=a_2x+b_2=0$，即 $(a_1,a_2)\parallel(b_1,b_2)$。
\end{derivation}
权方和只是把 $\vec u=\left(\dfrac{a_1}{\sqrt{b_1}},\dfrac{a_2}{\sqrt{b_2}}\right)$、
$\vec v=(\sqrt{b_1},\sqrt{b_2})$ 代入柯西：
\[
  \left(\frac{a_1^2}{b_1}+\frac{a_2^2}{b_2}\right)(b_1+b_2)
  \ge\left(\frac{a_1}{\sqrt{b_1}}\cdot\sqrt{b_1}+\frac{a_2}{\sqrt{b_2}}\cdot\sqrt{b_2}\right)^2
  =(a_1+a_2)^2 .
\]

\begin{conccard}[二级]{五大不等式的结构本质与取等条件}{证明不等式、选择放缩工具、检验取等}{只背公式不看结构；张冠李戴（该用柯西时硬凑基本不等式）；忽视取等条件导致「最值取不到」}
基本不等式＝半圆射影定理（$a=b$ 取等）；
柯西＝内积公理（共线取等）；
权方和＝柯西的齐次化（$\frac{a_1}{b_1}=\frac{a_2}{b_2}$ 取等）；
绝对值三角＝度量三角公理（共线取等）；
切线放缩＝凸凹函数的切线支撑（切点取等）。
\end{conccard}

\subsection{切线放缩与隐零点}

\key{探究四：切线是全局支撑边界；隐零点靠「设而不求」}

\key{（1）$e^x\ge x+1$ 与 $\ln x\le x-1$ 的严格证明。}
\begin{derivation}
设 $g(x)=e^x-x-1$，则 $g'(x)=e^x-1$。
$x<0$ 时 $g'<0$，$x>0$ 时 $g'>0$，故 $g$ 在 $x=0$ 处取最小值 $g(0)=0$。
所以
\[
  e^x\ge x+1\quad(\text{当且仅当 }x=0\text{ 取等}).
\]
设 $h(x)=x-1-\ln x\ (x>0)$，则 $h'(x)=1-\frac1x$。
$0<x<1$ 时 $h'<0$，$x>1$ 时 $h'>0$，故 $h$ 在 $x=1$ 处取最小值 $h(1)=0$。
所以
\[
  \ln x\le x-1\quad(\text{当且仅当 }x=1\text{ 取等}).
\]
几何含义：$y=e^x$ 是凸函数，恒在切线 $y=x+1$ 上方；$y=\ln x$ 是凹函数，恒在切线 $y=x-1$ 下方（图 \ref{fig:tangent-exp}）。
\end{derivation}

\begin{figure}[htbp]
  \centering
  \begin{tikzpicture}[>=Stealth, semithick, font=\footnotesize]
    \draw[->, gray] (-2.2,0)--(1.95,0) node[below right] {$x$};
    \draw[->, gray] (0,-1.2)--(0,5.1) node[above left] {$y$};
    \fill[main!10] plot[domain=-2:1.6, samples=60] (\x, {exp(\x)})
      -- plot[domain=1.6:-2, samples=60] (\x, {\x+1}) -- cycle;
    \draw[main!70!black, very thick, domain=-2:1.6, smooth, samples=80] plot (\x, {exp(\x)});
    \draw[second, thick, domain=-2:1.6] plot (\x, {\x+1});
    \fill (0,1) circle (1.6pt) node[above right] {$(0,1)$};
    \node[main!70!black] at (0.45,3.25) {$y=e^x$};
    \node[second] at (-1.85,-0.55) {$y=x+1$};
    % 结论文字移到左上空白区，避免被 $e^x$ 曲线或切线穿过
    \node[gray!70!black] at (-0.35,4.35) {$e^x\ge x+1$};
    \node[gray!70!black] at (-0.35,3.95) {仅在 $x=0$ 取等};
  \end{tikzpicture}
  \caption{$y=e^x$ 恒在切线 $y=x+1$ 上方（凸函数的切线支撑）：阴影为差值区域 $e^x-(x+1)\ge0$，两者仅在切点 $(0,1)$ 接触}
  \label{fig:tangent-exp}
\end{figure}

\key{（2）多步放缩：各步取等条件必须同时成立。}
\begin{example}
证明：对 $x>0$，$e^x-\ln x>2$。
\end{example}
\begin{solution}
由 $e^x\ge x+1$（取等 $x=0$）与 $-\ln x\ge1-x$（取等 $x=1$）相加得
\[
  e^x-\ln x\ge (x+1)+(1-x)=2 .
\]
但两式取等条件分别是 $x=0$ 与 $x=1$，不可能同时成立，故等号永远取不到：
\[
  e^x-\ln x>2 .
\]
\key{多步放缩时，每一步的取等条件必须相容，否则最后的「$\ge$」其实是严格的「$>$」。}
\end{solution}

\key{（3）隐零点：设而不求，整体代换。}
当 $f'(x)=0$ 的解 $x_0$ 无法用初等形式写出时，不必硬解：设 $x_0$ 为零点，
利用 $f'(x_0)=0$ 提供的等量关系做整体代换，把超越结构化为代数结构。

\begin{example}
证明：对一切 $x>-2$，$e^x>\ln(x+2)$。
\end{example}
\begin{solution}
设 $f(x)=e^x-\ln(x+2)$，$f'(x)=e^x-\dfrac1{x+2}$，$f''(x)=e^x+\dfrac1{(x+2)^2}>0$，
故 $f'$ 在 $(-2,+\infty)$ 上严格递增。
又
\[
  f'\!\left(-\tfrac12\right)=e^{-1/2}-\frac23<0,\qquad f'(0)=1-\frac12>0,
\]
所以存在唯一的 $x_0\in\left(-\dfrac12,0\right)$ 使 $f'(x_0)=0$，即
\[
  e^{x_0}=\frac1{x_0+2}\quad\Longrightarrow\quad \ln(x_0+2)=-x_0 .
\]
$f$ 在 $(-2,x_0)$ 上递减、在 $(x_0,+\infty)$ 上递增，最小值为
\[
  f(x_0)=e^{x_0}-\ln(x_0+2)=\frac1{x_0+2}+x_0 .
\]
由 $x_0\in\left(-\dfrac12,0\right)$ 得 $\dfrac1{x_0+2}>\dfrac12$ 且 $x_0>-\dfrac12$，故
\[
  f(x_0)>\frac12-\frac12=0 .
\]
因此 $f(x)\ge f(x_0)>0$，即 $e^x>\ln(x+2)$ 对一切 $x>-2$ 成立。$\surd$
\end{solution}

\begin{conccard}[一级]{切线放缩与隐零点}{证明含 $e^x$、$\ln x$ 的不等式；求不可解方程的极值}{忘记检验各步取等条件；不敢把「$\ge$」收紧为「$>$」；找到 $x_0$ 后不会用 $f'(x_0)=0$ 代换；忽略 $x_0$ 的取值范围}
$e^x\ge x+1$（取等 $x=0$），$\ln x\le x-1$（取等 $x=1$）；
多步放缩取等条件不相容 $\Rightarrow$ 严格不等；
隐零点：设 $f'(x_0)=0$，用该等式整体代换，并确定 $x_0$ 的范围。
\end{conccard}

\subsection{条件最值：四大标准动作}

\key{探究五：把多元、超越的结构，逐个「降维」成对勾函数}

\key{动作一：齐次化与「1」的代换（入口；完整方法链见第 8 章）。}
已知 $ma+nb=k$ 求 $\dfrac pa+\dfrac qb$ 的最值：乘以 $\dfrac{ma+nb}k=1$，展开后用基本不等式。

\begin{example}
$a,b>0$，$a+b=1$，求 $\dfrac1a+\dfrac4b$ 的最小值。
\end{example}
\begin{solution}
「1 的代换」：
\[
  \frac1a+\frac4b=\left(\frac1a+\frac4b\right)(a+b)
  =5+\frac ba+\frac{4a}b\ge5+2\sqrt4=9,
\]
当且仅当 $\dfrac ba=\dfrac{4a}b$ 即 $b=2a$ 时取等，结合 $a+b=1$ 得 $a=\dfrac13,b=\dfrac23$。
「权方和秒杀」：
\[
  \frac1a+\frac4b=\frac{1^2}{a}+\frac{2^2}{b}\ge\frac{(1+2)^2}{a+b}=9,
\]
取等条件同为 $a:b=1:2$。
\end{solution}

\key{动作二：消元与主元转化（双变量 $\to$ 对勾函数 $t+\dfrac kt$）。}

\begin{example}
$x,y>0$ 满足 $x+2y+xy=30$，求 $x+2y$ 的最小值。
\end{example}
\begin{solution}
由等式解出 $y=\dfrac{30-x}{x+2}$。令 $t=x+2$，则 $t>2$，且 $y>0\Rightarrow t<32$，故 $t\in(2,32)$。
\[
  x+2y=t-2+\frac{2(32-t)}t=t+\frac{64}t-4\ge2\sqrt{64}-4=12,
\]
当且仅当 $t=\dfrac{64}t$ 即 $t=8\in(2,32)$ 时取等，此时 $x=6,y=3$（检验：$6+6+18=30$）。
\end{solution}

\key{动作三：拆项分离常数。}

\begin{example}
求 $f(x)=\dfrac{x^2+3x+5}{x+1}$ 在 $x>-1$ 时的最小值。
\end{example}
\begin{solution}
令 $t=x+1>0$，则 $x=t-1$，分子
\[
  x^2+3x+5=(t-1)^2+3(t-1)+5=t^2+t+3,
\]
故
\[
  f(x)=t+1+\frac3t\ge1+2\sqrt3,
\]
当且仅当 $t=\dfrac3t$ 即 $t=\sqrt3$（$x=\sqrt3-1$）时取等。
\end{solution}

\key{动作四：切线放缩放去超越函数。}

\begin{example}
证明：$x>0$ 时 $\ln(x+1)+\dfrac1{x+1}>1$。
\end{example}
\begin{solution}
在 $\ln u\le u-1$ 中取 $u=\dfrac1t\ (t>0)$，得 $\ln t\ge1-\dfrac1t$（取等 $t=1$）。
令 $t=x+1>1$，则
\[
  \ln(x+1)\ge1-\frac1{x+1}\quad\Longrightarrow\quad
  \ln(x+1)+\frac1{x+1}\ge1 .
\]
取等需要 $t=1$ 即 $x=0$，但 $x>0$，故严格成立：
\[
  \ln(x+1)+\frac1{x+1}>1 .
\]
\key{用切线（一阶 Taylor）把 $\ln$、$e^x$ 替换为整式，是消去超越函数的标准动作。}
\end{solution}

\begin{commonerror}{易错诊断室：连续放缩的取等冲突与换元定义域}
\paragraph{反例一：连续两次用基本不等式，取等条件互相打架。}
已知 $x,y>0$，$x+2y=1$，某同学求 $\dfrac1{x^2}+\dfrac1{y^2}$ 的「最小值」：
\[
  \frac1{x^2}+\frac1{y^2}\ge\frac{2}{xy},\qquad
  xy=\frac{x\cdot2y}{2}\le\frac12\left(\frac{x+2y}{2}\right)^2=\frac18,
\]
于是 $\dfrac1{x^2}+\dfrac1{y^2}\ge\dfrac{2}{1/8}=16$。
问题：第一步取等要求 $x=y$，第二步取等要求 $x=2y$，二者不能同时成立，
所以 $16$ 根本取不到（真实最小值大于 $16$）。
\key{每写一次放缩，就把取等条件记下来；所有条件必须相容，最值才可能取到。}

\paragraph{反例二：换元后忽略新元的取值范围。}
求 $f(x)=\dfrac{x^2+5}{\sqrt{x^2+4}}$ 的最小值。令 $t=\sqrt{x^2+4}$，则 $t\ge2$，且
\[
  f(x)=t+\frac1t .
\]
若直接写 $t+\dfrac1t\ge2$，取等要求 $t=1$，但 $t\ge2$，取不到！
正确做法：$\varphi(t)=t+\dfrac1t$ 在 $[2,+\infty)$ 上递增，故最小值为
\[
  \varphi(2)=2+\frac12=\frac52 .
\]
\key{换元必先写出新元的定义域；取等点若不在定义域内，必须改用单调性求最值。}
\end{commonerror}

\begin{conccard}[一级]{条件最值四大标准动作}{二元条件最值、分式函数最值、含超越函数的不等式}{取等条件冲突未检验；换元后忘定义域；该用权方和时硬展开；切线放缩用于非切点附近却当作等式}
（1）1 的代换：乘 $\frac{ma+nb}k=1$ 展开（或直接用权方和）；齐次化的完整展开见第 8 章，本节只作条件最值的入口；
（2）消元主元：解出一元，换 $t$ 凑 $t+\frac kt$；
（3）拆项分离常数：$t=x+1$ 类换元降幂；
（4）切线放缩：用 $e^x\ge x+1$ 等消去超越项。
每步都要检验：取等条件相容吗？在定义域内吗？
\end{conccard}

\subsection{含参恒成立与能成立：点集位置与集合包含}

\key{探究六：把「成立性」翻译成「边界值的大小关系」}

\key{（1）点集上下位置 $\to$ 恒成立。}
$f(x)\ge a$ 在 $D$ 上恒成立，几何上是直线 $y=a$ 始终位于曲线 $y=f(x)$ 的下方（允许相切），
等价于
\[
  a\le f(x)_{\min}\quad(x\in D).
\]
类比地，$f(x)>a$ 恒成立需 $a\le f_{\min}$ 且验证最小值点处是否严格（最小值能取到时需 $a<f_{\min}$）。

\key{（2）一元二次含参不等式：分类讨论树。}
以 $ax^2+px+q>0$ 为例，固定顺序：先看二次项系数 $a$（零／正／负），
再看判别式 $\Delta$（正／零／负），最后比较两根大小（图 \ref{fig:quad-tree}）。

\begin{figure}[htbp]
  \centering
  \begin{tikzpicture}[>=Stealth, semithick, font=\footnotesize,
    box/.style={draw, rounded corners, align=center, inner sep=3pt, text width=2.2cm}]
    \node[box, text width=3.8cm] (root) at (0,0) {解关于 $x$ 的含参不等式\\$ax^{2}+px+q>0$\\($a,p,q$ 为参数)};
    \node[box] (l) at (-4.6,-1.6) {$a=0$\\退化为一次};
    \node[box] (m) at (0,-1.6) {$a>0$\\开口向上};
    \node[box] (r) at (4.6,-1.6) {$a<0$\\开口向下};
    \node[box] (ll) at (-5.2,-3.3) {$px+q>0$\\按 $p$ 符号定向};
    \node[box] (m1) at (-2.6,-3.3) {$\Delta>0$\\两根之外};
    \node[box] (m2) at (0,-3.3) {$\Delta=0$\\$x\ne -\frac{p}{2a}$};
    \node[box] (m3) at (2.6,-3.3) {$\Delta<0$\\解集为 $\mathbb R$};
    \node[box] (rr) at (5.2,-3.3) {乘 $-1$ 变号\\化为 $a>0$};
    \draw (root)--(l) (root)--(m) (root)--(r);
    \draw (l)--(ll) (m)--(m1) (m)--(m2) (m)--(m3) (r)--(rr);
    \node[gray!70!black, align=center, text width=13cm] at (0,-4.9)
      {注：$\Delta>0$ 时两根为 $x_{1,2}=\dfrac{-p\pm\sqrt{\Delta}}{2a}$，\\还必须比较两根大小，然后才能写下解集区间。};
  \end{tikzpicture}
  \caption{含参一元二次不等式的分类讨论树：二次项系数（零／正／负）$\to$ 判别式 $\Delta$（正／零／负）$\to$ 两根大小}
  \label{fig:quad-tree}
\end{figure}

\begin{example}
解关于 $x$ 的不等式 $ax^2-(a+1)x+1>0$。
\end{example}
\begin{solution}
因式分解：$ax^2-(a+1)x+1=(x-1)(ax-1)>0$。
\begin{itemize}
  \item $a=0$：化为 $-(x-1)>0$，解集 $(-\infty,1)$；
  \item $a>0$：两根为 $1$ 与 $\dfrac1a$。
    $a>1$ 时 $\dfrac1a<1$，解集 $\left(-\infty,\dfrac1a\right)\cup(1,+\infty)$；
    $a=1$ 时为 $(x-1)^2>0$，解集 $\{x\mid x\ne1\}$；
    $0<a<1$ 时 $\dfrac1a>1$，解集 $(-\infty,1)\cup\left(\dfrac1a,+\infty\right)$；
  \item $a<0$：开口向下，两根 $\dfrac1a<0<1$，解集 $\left(\dfrac1a,1\right)$。
\end{itemize}
\end{solution}

\key{（3）量词问题：值域边界的对偶。}
「任意」「存在」混合时，统一翻译为两个函数值域的边界关系（设最值都能取到）：

\begin{center}\footnotesize\renewcommand{\arraystretch}{1.3}
\begin{tabular}{c|c}
量词结构 & 等价条件 \\\hline
$\forall x_1\in D_1,\ \forall x_2\in D_2:\ f(x_1)\ge g(x_2)$ & $f_{\min}\ge g_{\max}$ \\[2pt]
$\forall x_1\in D_1,\ \exists x_2\in D_2:\ f(x_1)\ge g(x_2)$ & $f_{\min}\ge g_{\min}$ \\[2pt]
$\exists x_1\in D_1,\ \forall x_2\in D_2:\ f(x_1)\ge g(x_2)$ & $f_{\max}\ge g_{\max}$ \\[2pt]
$\exists x_1\in D_1,\ \exists x_2\in D_2:\ f(x_1)\ge g(x_2)$ & $f_{\max}\ge g_{\min}$ \\
\end{tabular}
\end{center}

\begin{example}
设 $f(x)=e^x-x$，$g(x)=x^2-2x+a$，$x\in[0,2]$。若对任意 $x_1\in[0,2]$，都存在 $x_2\in[0,2]$ 使 $f(x_1)\ge g(x_2)$，求 $a$ 的取值范围。
\end{example}
\begin{solution}
条件等价于 $f_{\min}\ge g_{\min}$。
$f'(x)=e^x-1\ge0$ 在 $[0,2]$ 上成立，$f$ 递增，$f_{\min}=f(0)=1$。
$g(x)=(x-1)^2+a-1$，$g_{\min}=g(1)=a-1$。
故
\[
  1\ge a-1\ \Rightarrow\ a\le2 .
\]
\end{solution}

\begin{conccard}[二级]{恒成立／能成立：点集位置与值域边界}{含参不等式恒成立、双变量量词问题、参数范围求解}{把 $\forall$ 与 $\exists$ 的最值方向搞反；忽略最值是否能取到；严格不等式端点处理错误；二次不等式分类讨论顺序混乱}
第 0 章的集合版（值域子集包含）在这里换成最值语言：$f(x)\ge a$ 恒成立 $\iff a\le f_{\min}$；
$\forall x_1,\exists x_2:\ f(x_1)\ge g(x_2)\iff f_{\min}\ge g_{\min}$（其余三种对偶见表）；
含参二次不等式按「$a$ 的符号 $\to\Delta\to$ 两根大小」三步分类。
\end{conccard}

\begin{method}[局部线性化引擎的使用流程]
\begin{enumerate}
  \item \key{识别结构}：见 $e^x$、$\ln x$ 想切线放缩；见双变量等式想 1 的代换／消元；见含参成立性想最值与值域边界；
  \item \key{局部线性化}：在关键点求导、写切线，用 $f(x)\approx f(x_0)+f'(x_0)(x-x_0)$ 把非线性问题线性化；
  \item \key{求最值}：导数符号法找单调区间与极值，再比较端点；
  \item \key{验取等}：记录每一步的取等条件，检查是否相容、是否在定义域内——不相容则「$\ge$」改「$>$」，不在定义域内则改用单调性；
  \item \key{回代检验}：用特殊值（如 $x=0,1$）或量纲、符号快速验证结论。
\end{enumerate}
\end{method}

\begin{summary}
本章把「导数」与「不等式」统一在\key{局部线性化（以直代曲）}这一个引擎之下：
曲线在一点附近由切线 $f(x_0)+f'(x_0)(x-x_0)$ 近似代替，导数是差分的连续版本；
由 $f'(x)$ 的符号统一判定单调性、极值与最值；
五大核心不等式分别起源于射影定理、内积公理、齐次化、度量三角公理与凸凹函数的切线支撑，取等条件各不相同；
$e^x\ge x+1$、$\ln x\le x-1$ 的严格证明与多步取等的相容性检验，构成切线放缩的完整操作规程；
隐零点靠「设而不求、整体代换」；条件最值有四大标准动作（1 的代换、消元主元、拆项分离、切线放缩）；
含参恒成立／能成立问题则归结为点集的上下位置与值域的边界大小。

\begin{center}
\footnotesize
\setlength{\tabcolsep}{4pt}
\renewcommand{\arraystretch}{1.35}
\begin{tabular}{p{2.3cm}p{4.6cm}p{4.6cm}p{1.9cm}}
\toprule
引擎部件 & 核心操作 & 把什么变成什么 & 对应探究 \\
\midrule
局部线性化 & $f(x)\approx f(x_0)+f'(x_0)(x-x_0)$ & 非线性 $\to$ 线性；变化率 $\to$ 切线斜率 & 探究一 \\
导数符号法 & $f'(x)>0$ 增、$f'(x)<0$ 减 & 单调性／极值／最值 $\to$ 符号判定 & 探究二 \\
结构本质表 & 射影／内积／齐次／度量／凸支撑 & 不等式 $\to$ 几何与代数边界 & 探究三 \\
切线放缩 & $e^x\ge x+1$，$\ln x\le x-1$ & 超越 $\to$ 整式（须验取等） & 探究四 \\
隐零点代换 & 设 $x_0$，用 $f'(x_0)=0$ 整体代换 & 不可解方程 $\to$ 代数关系 & 探究四 \\
条件最值四动作 & 1 的代换／消元主元／拆项／放缩 & 多变量、超越 $\to$ 一元低次 & 探究五 \\
成立性量词化 & $\forall\exists$ 四类转化 & 参数成立性 $\to$ 值域边界 & 探究六 \\
\bottomrule
\end{tabular}
\end{center}

\key{导数与不等式是同一个引擎的两面：以直代曲给出局部逼近，支撑边界给出整体控制；解题的本质是「识别结构 $\to$ 局部线性化 $\to$ 求最值 $\to$ 验取等」。}
\end{summary}

\end{explore}

\begin{inclasstest}
\begin{enumerate}
  \item 用切线近似求 $\sqrt{4.02}$ 的近似值（精确到 $0.001$）；\level{易}
  \item 求 $f(x)=x^3-3x^2-9x+5$ 的单调区间与极值；\level{中}
  \item 证明：对任意 $x>0$，$e^x-\ln x>2$；\level{中}
  \item 解关于 $x$ 的不等式 $ax^2-(a+1)x+1>0$。\level{难}
\end{enumerate}
\end{inclasstest}

\begin{afterclass}
\subsection*{基础巩固}
\begin{enumerate}
  \item 求函数 $f(x)=x\ln x\ (x>0)$ 的最小值；
  \item 证明：对任意实数 $x$，$e^x\ge ex$（提示：求 $f(x)=e^x-ex$ 的最小值）。
\end{enumerate}

\subsection*{进阶提升}
\begin{enumerate}
  \item 已知 $a,b>0$ 且 $a+2b=1$，求 $\dfrac1a+\dfrac1b$ 的最小值；
  \item 设 $f(x)=e^x-ax\ (a>0)$，讨论其单调性；若 $f(x)\ge0$ 对一切 $x\in\mathbb R$ 恒成立，求 $a$ 的取值范围。
\end{enumerate}

\subsection*{拓展探究}
\begin{enumerate}
  \item （强基／竞赛）把权方和不等式推广到三元：证明
  $\dfrac{a_1^2}{b_1}+\dfrac{a_2^2}{b_2}+\dfrac{a_3^2}{b_3}\ge\dfrac{(a_1+a_2+a_3)^2}{b_1+b_2+b_3}\ (b_i>0)$，
  并用它求 $a,b,c>0$、$a+b+c=3$ 时 $\dfrac1a+\dfrac4b+\dfrac9c$ 的最小值。
\end{enumerate}

\subsection*{教师版答案}
\begin{solution}
\textbf{基础 1：}$f'(x)=\ln x+1$。令 $f'(x)=0$ 得 $x=\dfrac1e$。
$0<x<\dfrac1e$ 时 $f'<0$，$x>\dfrac1e$ 时 $f'>0$，故最小值为
\[
  f\!\left(\frac1e\right)=\frac1e\ln\frac1e=-\frac1e .
\]
\end{solution}
\begin{solution}
\textbf{基础 2：}设 $f(x)=e^x-ex$，$f'(x)=e^x-e$。
$x<1$ 时 $f'<0$，$x>1$ 时 $f'>0$，故 $f$ 在 $x=1$ 处取最小值 $f(1)=0$，
即 $e^x\ge ex$，当且仅当 $x=1$ 取等（这正是 $y=e^x$ 在 $(1,e)$ 处的切线 $y=ex$）。
\end{solution}
\begin{solution}
\textbf{进阶 1：}「1 的代换」：
\[
  \frac1a+\frac1b=\left(\frac1a+\frac1b\right)(a+2b)
  =3+\frac{2b}{a}+\frac ab\ge3+2\sqrt2,
\]
当且仅当 $\dfrac{2b}{a}=\dfrac ab$ 即 $a=\sqrt2\,b$ 时取等，结合 $a+2b=1$ 得
$b=\dfrac1{2+\sqrt2}$，$a=\dfrac{\sqrt2}{2+\sqrt2}$。最小值 $3+2\sqrt2$。
\end{solution}
\begin{solution}
\textbf{进阶 2：}$f'(x)=e^x-a$。$x<\ln a$ 时 $f'<0$，$x>\ln a$ 时 $f'>0$，
故 $f$ 在 $(-\infty,\ln a)$ 递减、在 $(\ln a,+\infty)$ 递增，最小值为
\[
  f(\ln a)=a-a\ln a=a(1-\ln a).
\]
$f(x)\ge0$ 恒成立 $\iff a(1-\ln a)\ge0\iff\ln a\le1\iff a\le e$。
故 $a$ 的取值范围是 $(0,e]$。
\end{solution}
\begin{solution}
\textbf{拓展：}由柯西不等式，取
$\vec u=\left(\dfrac{a_1}{\sqrt{b_1}},\dfrac{a_2}{\sqrt{b_2}},\dfrac{a_3}{\sqrt{b_3}}\right)$，
$\vec v=(\sqrt{b_1},\sqrt{b_2},\sqrt{b_3})$，得
\[
  \left(\frac{a_1^2}{b_1}+\frac{a_2^2}{b_2}+\frac{a_3^2}{b_3}\right)(b_1+b_2+b_3)
  \ge(a_1+a_2+a_3)^2,
\]
即所要的不等式，取等当且仅当 $\dfrac{a_1}{b_1}=\dfrac{a_2}{b_2}=\dfrac{a_3}{b_3}$。
应用：
\[
  \frac1a+\frac4b+\frac9c=\frac{1^2}{a}+\frac{2^2}{b}+\frac{3^2}{c}
  \ge\frac{(1+2+3)^2}{a+b+c}=\frac{36}{3}=12,
\]
当且仅当 $a:b:c=1:2:3$ 且 $a+b+c=3$，即 $a=\dfrac12,b=1,c=\dfrac32$ 时取等。最小值为 $12$。
\end{solution}
\reading{拓展阅读 12}{导数：以直代曲的局部线性化}
\end{afterclass}
% 高观点导学案：计数原理、概率测度空间与统计线性逼近（第 7 章·高观点版）
% 主线：把「不确定性」纳入集合框架——样本空间是集合，概率是其上的归一化计数测度；
%       条件概率 = 样本全集缩小后的空间缩放；全概率/贝叶斯 = 完备互斥划分；
%       回归直线 = 散点云的竖直最小二乘投影；独立性检验 = 小概率事件原理的统计反证法；
%       压轴收官：新定义情境题的四步破题法（翻译 → 特例 → 归纳 → 严格证明）。

\section{计数原理、概率测度空间与统计线性逼近（高观点版）}

\begin{methodbox}{多框架对读}
\small
\begin{center}
\begin{tabular}{c|p{2.9cm}|p{5.0cm}|p{3.1cm}}
\toprule
视角 & 本章看点 & 别的视角怎么说 & 哪侧最省力 \\
\midrule
代数 & 计数公式、期望方差 & 符号对象的机械运算 & 硬算分布列 \\
几何 & 散点、回归直线、面积测度 & 图形的位置与拟合 & 看趋势、估相关 \\
集合 & 样本空间、事件、完备划分 & 事件＝子集；概率＝测度 & 界互斥与独立 \\
映射 & $X:\Omega\to\mathbb R$；最小二乘＝投影 & 换到残差正交的坐标系 & 一图读懂相关 \\
\bottomrule
\end{tabular}
\end{center}
卡的使用：概率题先画维恩图与树图（集合），统计题先想「投影到哪个子空间」（总图见第 0 章图~\ref{fig:four-view-map}）。
\end{methodbox}

第 0 章确立了集合与逻辑这一全书语言，第 4 章把向量确立为空间的通用线性语言，第 6 章把曲线局部化为直线；
本章把「不确定性」也纳入同一个框架——随机试验的样本空间 $\Omega$ 是一个集合，概率是定义在其上的\key{归一化计数测度}；
条件概率 $P(B\mid A)$ 的本质是把样本全集 $\Omega$ 缩小为子集 $A$ 之后的\key{空间缩放}；
线性回归方程是二维散点云在竖直方向上的\key{最小二乘投影直线}，独立性检验则是「小概率事件原理」支配下的\key{统计学反证法}。
四个模块层层展开：
\key{计数原理的集合本质}（并集基数与笛卡尔积基数，排列＝有序单射、组合＝无序子集族，捆绑／插空／隔板三技法）$\to$
\key{概率空间的归一化测度}（古典与几何概型的统一除法，条件概率＝空间缩放，全概率与贝叶斯＝完备互斥划分）$\to$
\key{统计线性逼近}（百分位数规则，回归直线必过样本中心，独立性检验的统计反证）$\to$
\key{新定义情境题破题模板}（翻译符号集合 $\to$ 构造小规模特例 $\to$ 归纳猜想通项 $\to$ 反证法／数学归纳法严格证明）。
核心思想：\key{随机性不是结构的对立面，而是有限集合上的测度；统计不是死记公式，而是「投影逼近 $+$ 小概率反证」；压轴新定义题的本质是「把陌生符号翻译回熟悉的集合与结构」。}

\begin{learnobjective}
\begin{itemize}
  \item 理解分类加法原理即互斥子集并集的基数公式 $|A\cup B|=|A|+|B|$，分步乘法原理即笛卡尔积的基数公式 $|A\times B|=|A|\cdot|B|$；
  \item 理解排列是有序单射序列、组合是无序 $k$ 元子集族，掌握 $A_n^k$ 与 $C_n^k$ 的关系 $A_n^k=C_n^k\cdot k!$；
  \item 掌握捆绑法、插空法、隔板法的适用结构，会把「相邻／不相邻／相同元素分堆」问题翻译为标准计数模型；
  \item 会把古典概型与几何概型统一为归一化测度 $P(A)=\mu(A)/\mu(\Omega)$，说出概率三条公理（非负性、规范性、可加性）；
  \item 理解条件概率 $P(B\mid A)=\dfrac{P(A\cap B)}{P(A)}$ 是样本空间缩小为 $A$ 后的重新归一化，会用全概率公式与贝叶斯公式；
  \item 掌握百分位数的计算规则，理解回归直线 $\hat y=\hat bx+\hat a$ 是竖直残差平方和最小的投影直线且必过 $(\bar x,\bar y)$；
  \item 理解独立性检验 $K^2$ 统计量背后的「小概率事件原理」，能正确陈述检验结论；
  \item 掌握新定义情境题的四步破题法（读懂符号集合 $\to$ 小规模特例 $\to$ 归纳猜想 $\to$ 反证法／数学归纳法严格证明）。
\end{itemize}
\end{learnobjective}

\begin{preview}
\subsection*{知识回顾}
\begin{itemize}
  \item 若 $A\cap B=\varnothing$，则 $|A\cup B|=$ \blank；
  \item 完成一件事有两类互斥的途径，第一类有 $n_1$ 种方法、第二类有 $n_2$ 种方法，则共有 \blank 种方法；
  \item 古典概型：$P(A)=$ \blank；几何概型：$P(A)=$ \blank；
  \item 样本相关系数 $r$ 满足 $|r|$ \blank $1$，$r>0$ 表示两个变量 \blank 相关；
  \item 反证法的逻辑结构：假设结论不成立，推出 \blank，从而肯定原结论。
\end{itemize}

\subsection*{教材速读}
阅读「计数原理」「概率」「成对数据的统计」相关教材，完成下列问题：
\begin{enumerate}
  \item 排列数公式 $A_n^m=$ \blank；组合数公式 $C_n^m=$ \blank；
  \item 条件概率公式 $P(B\mid A)=$ \blank（要求 \blank $\ne0$）；
  \item 全概率公式：若 $A_1,\dots,A_n$ 是 $\Omega$ 的完备互斥划分，则 $P(B)=$ \blank；
  \item 回归直线 $\hat y=\hat bx+\hat a$ 必过样本中心 \blank，且 $\hat a=$ \blank（用 $\bar x,\bar y,\hat b$ 表示）；
  \item 独立性检验统计量 $K^2=$ \blank（设 $2\times2$ 列联表中四个数据为 $a,b,c,d$）。
\end{enumerate}
\end{preview}

\begin{explore}

\subsection{计数原理：并集、笛卡尔积与三类经典技法}

\key{探究一：计数不是数数，而是集合的度量}

\key{（1）两个原理的集合本质。}
设完成一件事的所有可能结果构成集合 $\Omega$。若 $\Omega$ 被分成两个互斥的子集 $\Omega=A\sqcup B$（即 $A\cap B=\varnothing$），则
\[
  |\Omega|=|A\sqcup B|=|A|+|B|,
\]
这就是\key{分类加法原理}：各类之间「不重不漏」，求总数就是求互斥子集并集的基数。
若完成一件事要分两步：第一步的结果集为 $A$、第二步的结果集为 $B$，则所有可能结果组成笛卡尔积
\[
  \Omega=A\times B=\{(a,b)\mid a\in A,\ b\in B\},\qquad |\Omega|=|A\times B|=|A|\cdot|B|,
\]
这就是\key{分步乘法原理}：每一步的选择共同确定一个有序对，求总数就是求积集的基数。
一句话对照：加法管「或」（不交并），乘法管「且」（有序积）。

\begin{example}
从 5 名男同学和 4 名女同学中选出 3 人组成学习小组（不分职务），有多少种选法？若从中选出 3 人分别担任语文、数学、英语三科课代表，又有多少种选法？
\end{example}
\begin{solution}
第一问只关心「选出的集合」，是无序的 $3$ 元子集族，用组合：
\[
  C_9^3=\frac{9!}{3!\,6!}=84 .
\]
第二问关心「谁对应哪一科」，是有序安排：先选后排，由分步乘法原理
\[
  C_9^3\cdot A_3^3=84\times6=504\quad(\text{或直接 }A_9^3=9\times7\times8=504).
\]
两问相差 $3!=6$ 倍——正是「无序」与「有序」的差别。
\end{solution}

\key{（2）排列与组合：有序单射与无序子集。}
从 $n$ 个不同元素中取出 $k$ 个排成一列，相当于构造一个单射 $\{1,2,\dots,k\}\hookrightarrow\{1,2,\dots,n\}$，其个数为排列数
\[
  A_n^k=n(n-1)\cdots(n-k+1)=\frac{n!}{(n-k)!}.
\]
从 $n$ 个不同元素中取出 $k$ 个组成子集（不计顺序），其个数为组合数
\[
  C_n^k=\frac{A_n^k}{k!}=\frac{n!}{k!\,(n-k)!},
\]
因为「先选后排」＝「选子集」乘「子集全排列」：$A_n^k=C_n^k\cdot k!$——这正是分步乘法原理的一次应用。
特别地，$C_n^k=C_n^{n-k}$（选 $k$ 个留下等价于选 $n-k$ 个淘汰）。

\key{（3）三类经典技法的结构对号。}
\begin{itemize}
  \item \key{捆绑法（必相邻）}：把必须相邻的 $m$ 个元素捆成一个「大元素」，先与其他元素一起排，再解开内部排列：
    $n$ 个元素中指定 $m$ 个相邻的排法 $=(n-m+1)!\cdot A_m^m$。
  \item \key{插空法（不相邻）}：先排其余 $n$ 个元素（产生 $n+1$ 个空位，含两端），再把 $m$ 个元素插入 $m$ 个不同的空位：$A_n^n\cdot A_{n+1}^{\,m}$。（此处的 $n$ 指「其余元素个数」，与上一行捆绑法中 $n$ 为总元素数的口径不同。）
  \item \key{隔板法（相同元素分堆）}：不定方程 $x_1+x_2+\cdots+x_m=n\ (x_i\in\mathbb N^*)$ 的正整数解，对应把 $n$ 个相同小球排成一排后在 $n-1$ 个空隙中插入 $m-1$ 块板，解数为
    \[
      C_{n-1}^{\,m-1}.
    \]
    若允许 $x_i=0$，先令 $x_i'=x_i+1$ 化为正整数解，得 $C_{n+m-1}^{\,m-1}$。
\end{itemize}

\begin{example}
（1）6 人站成一排，甲、乙必须相邻，有多少种站法？
（2）6 人站成一排，甲、乙互不相邻，有多少种站法？
（3）求 $x+y+z=10$ 的正整数解的组数。
\end{example}
\begin{solution}
（1）捆绑：甲乙捆成一个大元素，与其余 4 人共 5 个元素全排列，甲乙内部再交换：
\[
  A_5^5\cdot A_2^2=120\times2=240 .
\]
（2）插空：先排其余 4 人，产生 5 个空位，甲乙选 2 个空位排列：
\[
  A_4^4\cdot A_5^2=24\times20=480 .
\]
（3）隔板：10 个相同小球排一排，9 个空隙插 2 块板分成 3 份（每份至少 1）：
\[
  C_9^2=36 .
\]
\end{solution}

\begin{figure}[htbp]
  \centering
  \begin{tikzpicture}[>=Stealth, semithick, font=\footnotesize]
    % ---------- 左：文氏图 ----------
    \draw[gray] (-1.7,-1.6) rectangle (3.9,1.6);
    \begin{scope}
      \clip (0,0) circle (1.15);
      \fill[second!20] (-1.8,-1.7) rectangle (4.0,1.7);
    \end{scope}
    \begin{scope}
      \clip (1.4,0) circle (1.15);
      \fill[third!20] (-1.8,-1.7) rectangle (4.0,1.7);
    \end{scope}
    \begin{scope}
      \clip (0,0) circle (1.15);
      \fill[main!32] (1.4,0) circle (1.15);
    \end{scope}
    \draw (0,0) circle (1.15);
    \draw (1.4,0) circle (1.15);
    \node at (-0.62,0) {$A$};
    \node at (2.02,0) {$B$};
    \node[main!70!black] at (0.7,-0.05) {\small \color{white}$A\cap B$};
    \node[gray!70!black] at (0.7,1.4) {$A\cup B$};
    \node[gray!70!black] at (-1.5,1.4) {$U$};
    % ---------- 右：卡诺图式网格 ----------
    \begin{scope}[xshift=5.4cm, yshift=0.5cm]
      \draw (0,0) rectangle (3.0,2.4);
      \draw (1.5,0)--(1.5,2.4);
      \draw (0,1.2)--(3.0,1.2);
      \node at (0.75,1.8) {$A\cap B$};
      \node at (2.25,1.8) {$A\cap\bar B$};
      \node at (0.75,0.6) {$\bar A\cap B$};
      \node at (2.25,0.6) {$\bar A\cap\bar B$};
      \node[above] at (0.75,2.4) {$B$};
      \node[above] at (2.25,2.4) {$\bar B$};
      \node[left] at (0,1.8) {$A$};
      \node[left] at (0,0.6) {$\bar A$};
    \end{scope}
  \end{tikzpicture}
  \caption{集合运算的卡诺图式表示：左图文氏图中 $A\cup B$ 是「或」、$A\cap B$ 是「且」；右图 $2\times2$ 网格穷举两个事件的四种组合，对应概率论中样本空间的完备互斥划分}
  \label{fig:ch07-venn}
\end{figure}

\begin{commonerror}{易错诊断室：可重复与不可重复、有序与无序}
\paragraph{失误一：有序、无序混淆。}
从 5 人中选 3 人「组成代表团」用 $C_5^3=10$；若「分别担任三个不同职务」就要用 $A_5^3=60$。
题干中出现「分配、担任、排成一列、编号」等词，结果自带顺序，必须按有序处理。
\paragraph{失误二：平均分组忘记除序。}
把 6 人平均分成 3 组（每组 2 人）：若直接写 $C_6^2C_4^2C_2^2=90$，等于把三个无区别的组又按顺序数了一遍，
正确的组数为
\[
  \frac{C_6^2C_4^2C_2^2}{3!}=\frac{90}{6}=15 .
\]
\key{分组若无编号（与「分到三个不同地方」区别开），组间要除以组数的阶乘。}
\paragraph{失误三：可重复与不可重复混淆。}
3 个人从 5 种口味冰淇淋中各选一种：口味可被多人重复选择时，由分步乘法原理有 $5^3=125$ 种；
若规定三人口味互不相同，则相当于从 5 种中取 3 种的排列 $A_5^3=60$。
一字之差，原理全变——务必先确认「选过的还能不能再选」。
\end{commonerror}

\begin{conccard}[一级]{计数原理的集合本质与三技法}{一切计数问题，尤其是含「相邻／不相邻／平均分组」的排列组合题}{有序无序混淆；平均分组忘除序；可重复当成不可重复；隔板法忽略「每份至少一个」}
分类加法＝互斥子集并集基数：$|A\sqcup B|=|A|+|B|$；
分步乘法＝笛卡尔积基数：$|A\times B|=|A|\cdot|B|$；
排列＝有序单射序列 $A_n^k=\frac{n!}{(n-k)!}$，组合＝无序子集族 $C_n^k=\frac{A_n^k}{k!}$；
捆绑：$(n-m+1)!\,A_m^m$；插空：$A_n^nA_{n+1}^{\,m}$；
隔板：$x_1+\cdots+x_m=n$ 的正整数解数 $=C_{n-1}^{\,m-1}$。
\end{conccard}

\subsection{概率测度：归一化、空间缩放与完备划分}

\key{探究二：概率是定义在样本空间上的「归一化集合测度」}

\key{（1）古典概型与几何概型的统一。}
任何随机试验的样本空间 $\Omega$ 都是一个集合。在其上定义一个「大小」函数 $\mu$（测度）：
古典概型中 $\mu$ 是计数测度（$\mu(A)=|A|$），几何概型中 $\mu$ 是长度／面积／体积。
只要每个基本结果等可能，就有统一的归一化公式
\[
  P(A)=\frac{\mu(A)}{\mu(\Omega)}.
\]
由此可直接验证概率的三条公理：非负性 $P(A)\ge0$、规范性 $P(\Omega)=1$、互斥可加性 $P(A\sqcup B)=P(A)+P(B)$——它们不过是测度的归一化性质。

\begin{example}
在区间 $[0,1]$ 上任取两点 $X,Y$，求 $|X-Y|\le\dfrac12$ 的概率。
\end{example}
\begin{solution}
样本空间是单位正方形 $\Omega=[0,1]^2$，$\mu(\Omega)=1$。
满足 $|X-Y|>\dfrac12$ 的区域是直线 $Y=X\pm\dfrac12$ 之外的两个三角形，每个面积为 $\dfrac12\left(\dfrac12\right)^2=\dfrac18$。
故
\[
  P(|X-Y|\le\tfrac12)=1-2\times\tfrac18=\tfrac34 .
\]
古典（数点）与几何（量面积）在此只是测度 $\mu$ 的不同，算法同为「除以全集」。
\end{solution}

\key{（2）条件概率 $P(B\mid A)$：样本空间的缩放。}
事件 $A$ 发生后，样本全集由 $\Omega$ 缩小为子集 $A$，每个结果仍在 $A$ 内等可能。
新样本空间上的归一化给出
\[
  P(B\mid A)=\frac{P(A\cap B)}{P(A)}\qquad(P(A)>0).
\]
几何读法：\key{把 $A$ 当作新的「全宇宙」，概率质量重新按 $P(A)$ 缩放}。它仍满足三条公理（规范性 $P(\Omega\mid A)=1$、可加性），因此是以 $A$ 为底的合法概率。

\begin{example}
10 张考签中有 1 张中奖，5 人依次不放回抽取。证明：无论先抽后抽，每人中奖概率都等于 $\dfrac1{10}$（抽签公平性）。
\end{example}
\begin{solution}
第 1 人显然为 $\frac1{10}$。
第 2 人中奖 $=$「第 1 人未中且第 2 人中」，由乘法公式与缩小的样本空间：
\[
  P=P(\text{第 1 人未中})\cdot P(\text{第 2 人中}\mid\text{第 1 人未中})
  =\frac{9}{10}\times\frac{1}{9}=\frac{1}{10}.
\]
第 3 人同理：$\frac{9}{10}\times\frac{8}{9}\times\frac{1}{8}=\frac1{10}$。
第 4、5 人同理（依次缩小样本空间后概率仍为 $\frac1{10}$）。
可见「先抽优势」是错觉——条件概率把空间缩小后自动抹平了顺序差异。
\end{solution}

\key{（3）全概率公式与贝叶斯公式：完备互斥划分。}
若 $\Omega=A_1\sqcup A_2\sqcup\cdots\sqcup A_n$（$A_i$ 两两不交、并为全集，称为\key{完备互斥划分}），则
\[
  P(B)=\sum_{i=1}^{n}P(A_i)\,P(B\mid A_i)\quad(\text{全概率公式}),
\]
它是面积可加性 $\mu(B)=\sum\mu(A_i\cap B)$ 的直接翻译：把 $B$ 沿划分切成互不重叠的碎片再求和。
反过来，已知 $B$ 发生了，反推它是哪个碎片贡献的：
\[
  P(A_j\mid B)=\frac{P(A_j)\,P(B\mid A_j)}{\sum\limits_{i=1}^{n}P(A_i)\,P(B\mid A_i)}\quad(\text{贝叶斯公式}).
\]
全概率是「由因及果」，贝叶斯是「由果溯因」：$P(A_j)$ 是先验概率，$P(A_j\mid B)$ 是纳入证据 $B$ 后的后验概率。

\begin{example}
某厂三条流水线产量占比为 $50\%,30\%,20\%$，次品率分别为 $2\%,3\%,5\%$。任取一件产品，（1）求是次品的概率；（2）若它是次品，求它来自第一条流水线的概率。
\end{example}
\begin{solution}
记 $A_i$ 为「产品来自第 $i$ 条流水线」，$B$ 为「取到次品」。
（1）由全概率公式
\[
  P(B)=0.5\times0.02+0.3\times0.03+0.2\times0.05=0.010+0.009+0.010=0.029 .
\]
（2）由贝叶斯公式
\[
  P(A_1\mid B)=\frac{0.5\times0.02}{0.029}=\frac{0.010}{0.029}=\frac{10}{29}\approx0.345 .
\]
\end{solution}

\begin{figure}[htbp]
  \centering
  \begin{tikzpicture}[>=Stealth, semithick, font=\footnotesize]
    % ---------- 左：完备划分的面积分解 ----------
    \draw[thick] (0,0) rectangle (8,3);
    \draw (4,0)--(4,3);
    \draw (6.4,0)--(6.4,3);
    \fill[main!30] (0,2.4) rectangle (4,3);
    \fill[main!30] (4,2.1) rectangle (6.4,3);
    \fill[main!30] (6.4,1.5) rectangle (8,3);
    \draw[thick] (0,0) rectangle (8,3);
    \draw (4,0)--(4,3);
    \draw (6.4,0)--(6.4,3);
    \node[main!80!black] at (2,2.7) {$0.010$};
    \node[main!80!black] at (5.2,2.55) {$0.009$};
    \node[main!80!black] at (7.2,2.25) {$0.010$};
    \node[align=center] at (2,-0.42) {$A_1$\\ $P(A_1)=0.5$};
    \node[align=center] at (5.2,-0.42) {$A_2$\\ $P(A_2)=0.3$};
    \node[align=center] at (7.2,-0.42) {$A_3$\\ $P(A_3)=0.2$};
    \node[gray!70!black] at (-0.35,3.3) {$\Omega$};
    \node[align=center] at (4,-1.15) {$P(B)=0.010+0.009+0.010=0.029$\\(列宽 $\propto P(A_i)$，块高 $\propto P(B\mid A_i)$ 统一放大 $30$ 倍)\\(暗色块面积 $\propto P(A_i\cap B)$，数字为 $P(A_i\cap B)$)};
    % ---------- 右：概率树 ----------
    \begin{scope}[xshift=9.2cm]
      \node[circle,fill=black,inner sep=1.6pt] (r) at (0,1.85) {};
      \node[above left] at (r) {$\Omega$};
      \node[circle,fill=black,inner sep=1.6pt] (a1) at (1.8,3.25) {};
      \node[circle,fill=black,inner sep=1.6pt] (a2) at (1.8,1.85) {};
      \node[circle,fill=black,inner sep=1.6pt] (a3) at (1.8,0.45) {};
      \node[right] at (a1) {$A_1$};
      \node[right] at (a2) {$A_2$};
      \node[right] at (a3) {$A_3$};
      \node[circle,fill=main,inner sep=1.6pt] (b1) at (4.2,3.75) {};
      \node[circle,fill=gray,inner sep=1.6pt] (nb1) at (4.2,2.75) {};
      \node[circle,fill=main,inner sep=1.6pt] (b2) at (4.2,2.35) {};
      \node[circle,fill=gray,inner sep=1.6pt] (nb2) at (4.2,1.35) {};
      \node[circle,fill=main,inner sep=1.6pt] (b3) at (4.2,0.95) {};
      \node[circle,fill=gray,inner sep=1.6pt] (nb3) at (4.2,-0.05) {};
      \draw (r)--(a1) (r)--(a2) (r)--(a3);
      \draw[main] (a1)--(b1); \draw[gray] (a1)--(nb1);
      \draw[main] (a2)--(b2); \draw[gray] (a2)--(nb2);
      \draw[main] (a3)--(b3); \draw[gray] (a3)--(nb3);
      \node at (0.55,2.75) {$0.5$};
      \node at (0.9,2.15) {$0.3$};
      \node at (0.55,1.0) {$0.2$};
      \node[main,left] at (2.95,3.65) {$0.02$}; \node[gray,left] at (2.95,2.9) {$0.98$};
      \node[main,left] at (2.95,2.25) {$0.03$}; \node[gray,left] at (2.95,1.5) {$0.97$};
      \node[main,left] at (2.95,0.85) {$0.05$}; \node[gray,left] at (2.95,0.1) {$0.95$};
      \node[main,right=2mm] at (b1) {$B$: $0.010$};
      \node[gray,right=2mm] at (nb1) {$B^c$: $0.490$};
      \node[main,right=2mm] at (b2) {$B$: $0.009$};
      \node[gray,right=2mm] at (nb2) {$B^c$: $0.291$};
      \node[main,right=2mm] at (b3) {$B$: $0.010$};
      \node[gray,right=2mm] at (nb3) {$B^c$: $0.190$};
      \node[align=center] at (2.1,-0.8) {$P(B)=0.010+0.009+0.010=0.029$\\[2pt]$P(A_1\mid B)=\frac{0.010}{0.029}=\frac{10}{29}\approx0.345$};
    \end{scope}
  \end{tikzpicture}
  \caption{样本空间的完备互斥划分（全概率·贝叶斯）：左图中 $\Omega=A_1\sqcup A_2\sqcup A_3$，列宽正比于 $P(A_i)$、每列内暗色块高度正比于 $P(B\mid A_i)$（统一放大 $30$ 倍的示意图），故暗色块面积正比于 $P(A_i\cap B)$；右图沿树枝把每个 $A_i$ 再按条件概率分叉，路径概率求和得 $P(B)$，贝叶斯公式则由结果 $B$ 沿树反向回溯到原因 $A_1$}
  \label{fig:ch07-partition}
\end{figure}

\begin{conccard}[一级]{归一化测度：条件概率、全概率与贝叶斯}{几何概型、条件概率、由果溯因类应用题、医学筛查类情境}{把 $P(B\mid A)$ 与 $P(A\mid B)$ 混为一谈；划分不满足「不重不漏」就套全概率；只记住公式不理解「空间缩放」的几何意义}
$P(A)=\mu(A)/\mu(\Omega)$（古典＝计数测度，几何＝长度／面积／体积）；
条件概率＝空间缩放：$P(B\mid A)=P(A\cap B)/P(A)$；
全概率：$P(B)=\sum_i P(A_i)P(B\mid A_i)$；
贝叶斯（由果溯因）：$P(A_j\mid B)=\dfrac{P(A_j)P(B\mid A_j)}{\sum_i P(A_i)P(B\mid A_i)}$。
\end{conccard}

\subsection{统计线性逼近：最小二乘投影与统计反证}

\key{探究三：回归直线是散点云的「最小二乘投影」；独立性检验是统计反证法}

\key{（1）百分位数：位置的度量。}
把 $n$ 个数据从小到大排列，计算 $i=n\times p\%$：
\begin{itemize}
  \item 若 $i$ 不是整数，第 $p$ 百分位数是第 $\lceil i\rceil$ 项（向上取整的下一项）；
  \item 若 $i$ 是整数，第 $p$ 百分位数是第 $i$ 项与第 $i+1$ 项的平均数。
\end{itemize}

\begin{example}
求数据 $12,7,9,15,10,8,11,13,6,14$ 的第 $25$ 百分位数与第 $50$ 百分位数。
\end{example}
\begin{solution}
排序：$6,7,8,9,10,11,12,13,14,15$，共 $n=10$ 个。
\[
  i_1=10\times25\%=2.5\ (\text{非整数})\Rightarrow \text{第 }25\text{ 百分位数是第 }3\text{ 项 }=8;
\]
\[
  i_2=10\times50\%=5\ (\text{整数})\Rightarrow \text{第 }50\text{ 百分位数}=\frac{\text{第 }5\text{ 项}+\text{第 }6\text{ 项}}{2}=\frac{10+11}{2}=10.5 .
\]
\end{solution}

\key{（2）回归直线：竖直残差平方和最小的投影线。}
给定散点 $(x_1,y_1),\dots,(x_n,y_n)$，在直线族 $\hat y=bx+a$ 中挑选一条，使各点到直线的\key{竖直距离}（残差 $e_i=y_i-\hat y_i$）的平方和最小：
\[
  Q(a,b)=\sum_{i=1}^{n}(y_i-bx_i-a)^2\ \longrightarrow\ \min .
\]

\begin{derivation}
对 $a,b$ 分别求偏导并令其为零：
\[
  \frac{\partial Q}{\partial a}=-2\sum_{i=1}^n(y_i-bx_i-a)=0
  \ \Longrightarrow\ na+b\sum x_i=\sum y_i\ \Longrightarrow\ a=\bar y-b\bar x .
\]
\[
  \frac{\partial Q}{\partial b}=-2\sum_{i=1}^nx_i(y_i-bx_i-a)=0 .
\]
把 $a=\bar y-b\bar x$ 代入，利用 $\sum(x_i-\bar x)=0$，得
\[
  \sum_{i=1}^n(x_i-\bar x)(y_i-\bar y)-b\sum_{i=1}^n(x_i-\bar x)^2=0,
\]
故
\[
  \hat b=\frac{\sum\limits_{i=1}^n(x_i-\bar x)(y_i-\bar y)}{\sum\limits_{i=1}^n(x_i-\bar x)^2},
  \qquad \hat a=\bar y-\hat b\,\bar x .
\]
几何性质：$(\hat a,\hat b)$ 满足 $\hat a=\bar y-\hat b\,\bar x$，即点 $(\bar x,\bar y)$ 恒在直线上——\key{回归直线必过样本中心 $(\bar x,\bar y)$}。
这与第 6 章「以直代曲」、第 4 章「正交投影」同源：最小二乘就是把观测向量 $\vec y$ 向由 $\vec 1$ 与 $\vec x$ 张成的平面作投影，残差与平面正交。
\end{derivation}

\begin{example}
已知 5 组成对数据 $(1,2),(2,3),(3,5),(4,6),(5,8)$，求回归方程并计算样本相关系数 $r$。
\end{example}
\begin{solution}
\[
  \bar x=3,\quad \bar y=\frac{2+3+5+6+8}{5}=4.8 .
\]
\[
  \sum(x_i-\bar x)(y_i-\bar y)=(-2)(-2.8)+(-1)(-1.8)+0\times0.2+1\times1.2+2\times3.2=15,
\]
\[
  \sum(x_i-\bar x)^2=4+1+0+1+4=10,\quad
  \sum(y_i-\bar y)^2=7.84+3.24+0.04+1.44+10.24=22.8 .
\]
所以
\[
  \hat b=\frac{15}{10}=1.5,\qquad \hat a=4.8-1.5\times3=0.3,
  \qquad \hat y=1.5x+0.3 .
\]
\[
  r=\frac{15}{\sqrt{10\times22.8}}=\frac{15}{\sqrt{228}}\approx0.993 ,
\]
高度正相关（$|r|\le1$ 正是由第 6 章柯西不等式保证的）。
\end{solution}

\begin{figure}[htbp]
  \centering
  \begin{tikzpicture}[>=Stealth, semithick, font=\footnotesize]
    \draw[->,gray] (-0.4,0)--(5.9,0) node[below right] {$x$};
    \draw[->,gray] (0,-0.6)--(0,8.9) node[above left] {$y$};
    \foreach \x in {1,...,5} \draw (\x,0.06)--(\x,-0.06) node[below] {\x};
    \foreach \y in {2,4,6,8} \draw (0.06,\y)--(-0.06,\y) node[left] {\y};
    \draw[second,very thick,domain=-0.2:5.5,samples=2] plot (\x,{1.5*\x+0.3});
    \node[second] at (0.55,4.6) {$\hat y=1.5x+0.3$};
    \foreach \x/\y in {1/2,2/3,3/5,4/6,5/8}{
      \draw[second!80,dashed,very thick] (\x,\y)--(\x,{1.5*\x+0.3});
      \fill[main] (\x,\y) circle (1.7pt);
    }
    \node[second,right] at (1.04,1.9) {$e_1$};
    \node[second,right] at (2.04,3.15) {$e_2$};
    \node[second,left] at (2.96,4.98) {$e_3$};
    \node[second,right] at (4.04,6.15) {$e_4$};
    \node[second,right] at (5.04,7.9) {$e_5$};
    \draw[third,dashed,thick] (2.68,4.8)--(3.32,4.8);
    \draw[third,dashed,thick] (3,4.48)--(3,5.12);
    \fill[third] (3,4.8) circle (1.7pt);
    \node[third,right] at (3.35,4.86) {$(\bar x,\bar y)=(3,4.8)$};
    \fill[main] (2.9,8.55) circle (1.7pt); \node[right] at (2.98,8.55) {观测点};
    \draw[second,very thick] (4.1,8.55)--(4.5,8.55); \node[right] at (4.58,8.55) {回归直线};
  \end{tikzpicture}
  \caption{最小二乘回归：回归直线 $\hat y=1.5x+0.3$ 使各观测点到直线的竖直残差 $e_1,\dots,e_5$（橙色虚线）平方和最小；它必然穿过样本中心 $(\bar x,\bar y)=(3,4.8)$}
  \label{fig:ch07-leastsq}
\end{figure}

\key{（3）独立性检验：小概率事件原理的统计反证法。}
$2\times2$ 列联表
\[
  \begin{array}{c|cc|c}
    & y_1 & y_2 & \text{合计}\\\hline
    x_1 & a & b & a+b\\
    x_2 & c & d & c+d\\\hline
    \text{合计} & a+c & b+d & n
  \end{array}
\]
若两变量无关（零假设 $H_0$），应有 $ad\approx bc$。统计量
\[
  K^2=\frac{n(ad-bc)^2}{(a+b)(c+d)(a+c)(b+d)}
\]
度量实际数据偏离「无关」的程度。逻辑结构与第 0 章的反证法完全同构：先承认 $H_0$ 成立，则 $K^2$ 很大是\key{小概率事件}；它竟然发生了，就有理由拒绝 $H_0$。常用临界值：$K^2\ge3.841$ 对应犯错概率不超过 $5\%$，$K^2\ge6.635$ 对应不超过 $1\%$。

\begin{example}
调查 $200$ 人（男、女各 $100$ 人）对某项政策的支持情况：男性中 $70$ 人支持、$30$ 人反对；女性中 $50$ 人支持、$50$ 人反对。是否有 $99\%$ 以上的把握认为性别与支持态度有关？
\end{example}
\begin{solution}
\[
  K^2=\frac{200\times(70\times50-30\times50)^2}{100\times100\times120\times80}
  =\frac{200\times2000^2}{9.6\times10^7}
  =\frac{8\times10^8}{9.6\times10^7}\approx8.33 .
\]
因为 $8.33>6.635$，所以在犯错误的概率不超过 $0.01$ 的前提下，认为性别与支持态度有关——即有 $99\%$ 以上的把握。
\end{solution}

\begin{commonerror}{独立性检验的三个高频失误}
\paragraph{失误一：把结论说过头。}
正确表述是「在犯错误的概率不超过 $0.01$ 的前提下认为有关」或「有 $99\%$ 的把握认为有关」，
而不是「两变量有 $99\%$ 的关系」——$K^2$ 只衡量证据强度，不量化关联程度。
\paragraph{失误二：忽略样本量 $n$。}
公式分子含 $n$，同样比例的列联表，样本量扩大 $10$ 倍，$K^2$ 也扩大 $10$ 倍，结论可能完全不同；小样本下比例再悬殊也不宜下强结论。
\paragraph{失误三：把检验结论当成因果关系。}
「有关」不等于「导致」：独立性检验只能发现关联，因果方向须由实际背景另行论证。
\end{commonerror}

\begin{conccard}[一级]{最小二乘回归与独立性检验}{成对数据的回归计算、相关性判断、列联表分析}{把回归直线当成必过原点的比例函数；以为 $K^2$ 大就证明因果关系；计算 $\hat b$ 时忘除以 $\sum(x_i-\bar x)^2$；百分位数忽略 $i$ 是否为整数}
$\hat b=\frac{\sum(x_i-\bar x)(y_i-\bar y)}{\sum(x_i-\bar x)^2}$，$\hat a=\bar y-\hat b\bar x$，必过 $(\bar x,\bar y)$；
百分位数看 $i=n\cdot p\%$ 是否整数；
$K^2=\frac{n(ad-bc)^2}{(a+b)(c+d)(a+c)(b+d)}$，对照 $3.841$（$0.05$）、$6.635$（$0.01$）下结论。
\end{conccard}

\subsection{压轴收官：新定义情境题的四步破题法}

\key{探究四：新定义题不是考知识，是考「翻译」}

新高考压轴常以全新符号、新运算、新集合开场，陌生感是纸老虎。通用四步法：
\key{第一步，读懂符号集合}——把每个新名词逐字翻译为已有的集合语言（元素是什么？满足什么条件？运算封闭吗？）；
\key{第二步，构造小规模特例试探}——取 $n=1,2,3$ 等最小情形，逐一列举，列出数表；
\key{第三步，归纳猜想通项}——从特例中找规律、猜公式（归纳只指方向，不能当证明）；
\key{第四步，严格证明}——用双射（一一对应）、数学归纳法或反证法把猜想钉死。

\begin{method}[新定义情境题的四步破题流程]
\begin{enumerate}
  \item \key{翻译}：逐句把新定义的符号「转译」成集合、映射、计数或概率语言，画出韦恩图或数表；
  \item \key{特例}：取最小的参数（$n=1,2,3$ 或最小非平凡情形）把结果全部列出，观察结构；
  \item \key{归纳}：由特例数据猜想通项／公式／计数法则，并明确写出猜想；
  \item \key{证明}：优先寻找双射（把新对象一一对应到已知模型），否则用数学归纳法（与自然数相关）或反证法（唯一性、不存在性）。
\end{enumerate}
\end{method}

\begin{problem}
\key{压轴题一（新定义 $+$ 计数）}设 $n\ge k\ge1$。称 $\{1,2,\dots,n\}$ 的 $k$ 元子集 $A=\{a_1,a_2,\dots,a_k\}$（$a_1<a_2<\dots<a_k$）为「稀疏集」，若其中任意两个不同元素之差的绝对值不小于 $2$（即不含相邻整数）。记稀疏集的个数为 $f(n,k)$。
（1）列举 $n=5,\ k=2$ 的全部稀疏集，并求 $f(5,2)$；再求 $f(5,3)$；
（2）依（1）猜想 $f(n,k)$ 的表达式，并给出严格证明。
\end{problem}
\begin{solution}
\key{【第一步：读懂符号集合】}「稀疏集」＝$\{1,\dots,n\}$ 中不含相邻数的 $k$ 元子集，$f(n,k)$ 是这类子集的个数。
\key{【第二步：小规模特例】}（1）$n=5,k=2$：枚举得
\[
  \{1,3\},\{1,4\},\{1,5\},\{2,4\},\{2,5\},\{3,5\},
\]
共 $6$ 个，即 $f(5,2)=6=C_4^2$。
$n=5,k=3$：只能取 $\{1,3,5\}$，故 $f(5,3)=1=C_3^3$。
再试 $f(4,2)=3=C_3^2$、$f(6,2)=10=C_5^2$。
\key{【第三步：归纳猜想】}各值都等于 $C_{\,n-k+1}^{\,k}$，故猜想
\[
  f(n,k)=C_{\,n-k+1}^{\,k}.
\]
\key{【第四步：严格证明（双射法）】}设稀疏集 $A=\{a_1<\dots<a_k\}$，令
\[
  b_i=a_i-(i-1)\quad(i=1,\dots,k).
\]
由 $a_{i+1}-a_i\ge2$ 得 $b_{i+1}-b_i=(a_{i+1}-a_i)-1\ge1$，即 $b_1<b_2<\dots<b_k$；
又 $b_1=a_1\ge1$，$b_k=a_k-(k-1)\le n-k+1$，故 $\{b_1,\dots,b_k\}$ 是 $\{1,\dots,n-k+1\}$ 的任意 $k$ 元子集。
反之，任取 $1\le b_1<\dots<b_k\le n-k+1$，令 $a_i=b_i+(i-1)$，则 $a_{i+1}-a_i=b_{i+1}-b_i+1\ge2$，回到一个稀疏集。
两个方向互为逆映射，构成\key{一一对应}，故
\[
  f(n,k)=C_{\,n-k+1}^{\,k}.\ \surd
\]
\key{这一步同时解释了「为什么恰好是隔板模型」：减去的 $(i-1)$ 相当于预先在每个选中元素后面放好一个「隔离空位」，隔板数从 $n-1$ 缩为 $n-k+1$。}
\end{solution}

\begin{problem}
\key{压轴题二（新定义 $+$ 概率统计）}一批共 $6$ 件产品中有 $2$ 件次品，现逐个检查（不放回）。定义「转正等待次数」$X$ 为：首次抽到正品时已经检查过的产品件数。
（1）求 $X$ 的分布列与数学期望 $E(X)$；
（2）若改为有放回抽取，$X$ 服从什么分布？比较两种情形的期望并解释差异原因。
\end{problem}
\begin{solution}
\key{【第一步：翻译】}$X$ 是取值于 $\{1,2,3\}$ 的离散随机变量（最多查 3 件必出正品，因为只有 2 件次品）。
\key{【第二步：小特例手算】}
\[
  P(X=1)=\frac{4}{6}=\frac{2}{3};
\]
\[
  P(X=2)=P(\text{第 1 件次品})\cdot P(\text{第 2 件正品}\mid\text{第 1 件次品})
  =\frac{2}{6}\times\frac{4}{5}=\frac{4}{15};
\]
\[
  P(X=3)=P(\text{前两件均为次品})=\frac{2}{6}\times\frac{1}{5}=\frac{1}{15}.
\]
\key{【第三步：归纳校验】}$\frac23+\frac{4}{15}+\frac{1}{15}=1$，分布列为
\[
  \begin{array}{c|ccc}
    X & 1 & 2 & 3\\\hline
    P & \dfrac{2}{3} & \dfrac{4}{15} & \dfrac{1}{15}
  \end{array}
\]
\key{【第四步：期望与对比】}
\[
  E(X)=1\times\frac{2}{3}+2\times\frac{4}{15}+3\times\frac{1}{15}
  =\frac{10}{15}+\frac{8}{15}+\frac{3}{15}=\frac{21}{15}=\frac{7}{5}.
\]
（2）有放回时每次命中正品概率恒为 $\frac{4}{6}=\frac23$，$X$ 服从参数 $p=\frac23$ 的几何分布：
\[
  P(X=k)=\left(\frac{1}{3}\right)^{k-1}\cdot\frac{2}{3}\ (k\ge1),\qquad E(X)=\frac{1}{p}=\frac{3}{2}.
\]
$\frac{3}{2}>\frac{7}{5}$：不放回时次品一旦被抽出，剩下产品中正品浓度上升，等待时间自然缩短——这正是条件概率「空间缩放」在期望上的体现。
\end{solution}

\begin{conccard}[二级]{新定义情境题四步破题法}{强基、新高考压轴中「新运算、新集合、新统计量」类题目}{把归纳猜想的结论直接当证明使用；跳过特例直接硬算；翻译时漏掉隐含约束（如 $n\ge k$）}
第一步 翻译：逐句转译为集合／映射语言，画图表；
第二步 特例：$n=1,2,3$ 全部列举；
第三步 归纳：写出猜想公式；
第四步 证明：双射（计数）、数学归纳法（自然数命题）、反证法（唯一性与不存在）。
\end{conccard}

\begin{summary}
本章把「不确定性」彻底纳入第 0 章的集合框架：样本空间 $\Omega$ 是集合，概率是归一化计数测度 $P(A)=\mu(A)/\mu(\Omega)$；
计数问题的本质是并集基数（分类加法）与笛卡尔积基数（分步乘法），排列与组合分别是「有序单射」与「无序子集族」；
条件概率 $P(B\mid A)$ 是把样本全集缩小为 $A$ 后的空间缩放，全概率与贝叶斯公式源于样本空间的完备互斥划分（由因及果与由果溯因）；
线性回归直线是散点云的竖直最小二乘投影，必过样本中心 $(\bar x,\bar y)$；独立性检验则是「小概率事件原理」下的统计反证法；
压轴新定义题按「翻译符号 $\to$ 特例试探 $\to$ 归纳猜想 $\to$ 严格证明」四步破题。

\begin{center}
\footnotesize
\setlength{\tabcolsep}{4pt}
\renewcommand{\arraystretch}{1.35}
\begin{tabular}{p{2.5cm}p{4.7cm}p{4.5cm}p{1.4cm}}
\toprule
引擎部件 & 核心操作 & 把什么变成什么 & 对应探究 \\
\midrule
计数两原理 & $|A\sqcup B|=|A|+|B|$；$|A\times B|=|A|\,|B|$ & 分类（或）$\to$ 并；分步（且）$\to$ 积 & 探究一 \\
排列与组合 & $A_n^k=\frac{n!}{(n-k)!}$；$C_n^k=\frac{A_n^k}{k!}$ & 有序单射序列 $\to$ 无序 $k$ 元子集族 & 探究一 \\
三技法 & 捆绑／插空／隔板 $C_{n-1}^{m-1}$ & 相邻、不相邻、相同分堆 $\to$ 标准模型 & 探究一 \\
归一化测度 & $P(A)=\mu(A)/\mu(\Omega)$ & 古典与几何 $\to$ 同一除法 & 探究二 \\
条件概率 & $P(B\mid A)=\frac{P(A\cap B)}{P(A)}$ & 样本全集 $\Omega\to$ 子集 $A$：缩放归一 & 探究二 \\
全概率·贝叶斯 & $P(B)=\sum_iP(A_i)P(B\mid A_i)$ & 由因及果 $\leftrightarrow$ 由果溯因 & 探究二 \\
最小二乘回归 & $\min\sum(y_i-\hat bx_i-\hat a)^2$，过 $(\bar x,\bar y)$ & 散点云 $\to$ 竖直投影直线 & 探究三 \\
独立性检验 & $K^2$ 对临界值 & 关联判断 $\to$ 小概率反证 & 探究三 \\
四步破题 & 翻译 $\to$ 特例 $\to$ 归纳 $\to$ 严格证 & 陌生符号 $\to$ 熟悉结构 & 探究四 \\
\bottomrule
\end{tabular}
\end{center}

\key{随机性不是结构的对立面，而是有限集合上的测度；统计不是死记公式，而是「最小二乘投影 $+$ 小概率反证」的合体；新定义压轴的本质是把陌生符号翻译回熟悉的集合与结构。}
\end{summary}

\end{explore}

\begin{inclasstest}
\begin{enumerate}
  \item 从 6 名同学中选出 3 人，分别担任语文、数学、英语三科课代表，共有多少种不同的选法？\level{易}
  \item 把 10 个相同的小球放入 3 个不同的盒子，每个盒子至少放 1 个，共有多少种放法？\level{易}
  \item 口袋中装有 5 个红球、3 个白球，不放回地取 3 个球。已知取出的第一个球是红球，求取出的第三个球也是红球的概率。\level{中}
  \item 已知某组成对数据的样本中心为 $(\bar x,\bar y)=(3,5)$，且 $\sum(x_i-\bar x)(y_i-\bar y)=12$，$\sum(x_i-\bar x)^2=4$。求回归直线方程，并指出它必过的定点。\level{难}
\end{enumerate}
\end{inclasstest}

\begin{afterclass}
\subsection*{基础巩固}
\begin{enumerate}
  \item 6 人站成一排，甲、乙、丙三人必须站在一起，有多少种站法？
  \item 袋中有 3 个红球、2 个白球，不放回地依次取出 2 个球，求第二次取到红球的概率。
\end{enumerate}

\subsection*{进阶提升}
\begin{enumerate}
  \item 甲、乙两人射击的命中率分别为 $0.8$ 与 $0.7$，各射击一次。已知恰有一人命中，求命中者是甲的概率。
  \item 将 5 个不同的小球放入 3 个不同的盒子，每个盒子至少放 1 个，有多少种放法？
\end{enumerate}

\subsection*{拓展探究}
\begin{enumerate}
  \item （强基／竞赛）设 $n\ge k\ge1$。称 $\{1,2,\dots,n\}$ 的 $k$ 元子集 $A$ 为「强稀疏集」，若其中任意两个不同元素之差至少为 $3$。记其个数为 $g(n,k)$。
  （1）列出 $n=6,k=2$ 的全部强稀疏集，求 $g(6,2)$；
  （2）猜想 $g(n,k)$ 的表达式并用双射法证明；
  （3）比较 $g(n,k)$ 与本章 $f(n,k)$ 的大小，说明理由。
\end{enumerate}

\subsection*{教师版答案}
\begin{solution}
\textbf{基础 1：}把甲、乙、丙三人捆成一个大元素，与其余 3 人共 4 个元素全排列，三人内部再全排列：
\[
  A_4^4\cdot A_3^3=24\times6=144 .
\]
\end{solution}
\begin{solution}
\textbf{基础 2：}方法一（全概率）：
\[
  P=\frac35\times\frac24+\frac25\times\frac34=\frac{6}{20}+\frac{6}{20}=\frac{3}{5}.
\]
方法二（抽签公平性）：先后手对称，第二次取到红球的概率等于第一次取到红球的概率 $\frac35$。
\end{solution}
\begin{solution}
\textbf{进阶 1：}记 $A$ 为「甲命中」，$B$ 为「乙命中」，$A,B$ 独立。
\[
  P(\text{恰有一人命中})=P(A\bar B)+P(\bar A B)=0.8\times0.3+0.2\times0.7=0.38 .
\]
由条件概率（把样本空间缩小为「恰有一人命中」）：
\[
  P(A\mid \text{恰有一人命中})=\frac{0.24}{0.38}=\frac{12}{19}\approx0.632 .
\]
\end{solution}
\begin{solution}
\textbf{进阶 2：}由容斥原理（全集减去有空盒的放法）：
\[
  3^5-C_3^1\cdot2^5+C_3^2\cdot1^5=243-96+3=150 .
\]
\end{solution}
\begin{solution}
\textbf{拓展：}（1）$n=6,k=2$：枚举得
\[
  \{1,4\},\{1,5\},\{1,6\},\{2,5\},\{2,6\},\{3,6\},
\]
共 $6$ 个，即 $g(6,2)=6$。
（2）猜想：$g(n,k)=C_{\,n-2k+2}^{\,k}$（约定 $m<k$ 时 $C_m^k=0$，对应不可能情形）。
双射证明：设强稀疏集 $A=\{a_1<\dots<a_k\}$，$a_{i+1}-a_i\ge3$。令
\[
  b_i=a_i-2(i-1)\quad(i=1,\dots,k),
\]
则 $b_{i+1}-b_i=(a_{i+1}-a_i)-2\ge1$，即 $b_1<\dots<b_k$；且 $b_1\ge1$，$b_k=a_k-2(k-1)\le n-2k+2$，
故 $\{b_1,\dots,b_k\}$ 是 $\{1,\dots,n-2k+2\}$ 的任意 $k$ 元子集；反之 $a_i=b_i+2(i-1)$ 给出逆映射。一一对应，故 $g(n,k)=C_{\,n-2k+2}^{\,k}$。
（3）「强稀疏集」一定是「稀疏集」，反之不然（如 $\{1,3\}$ 稀疏但不强稀疏），前者集合族是后者的真子族，故除 $k=1$ 等平凡情形外恒有
\[
  g(n,k)<f(n,k).
\]
\end{solution}

\reading{拓展阅读 13}{概率：从计数测度到样本空间}
\end{afterclass}

% ==========================================================
% 第 8 章 齐次化：把二元问题压成一元比例
% 说明：本文件为纯正文片段，由 main_body.tex 挂载；
%       使用 versions.tex 的 \tk / \vchoice / advancedonly /
%       basiconly / solution / 四个盒子。
% ==========================================================

\section{齐次化：把二元问题压成一元比例（高观点结构版）}

\begin{methodbox}{多框架对读}
\small
\begin{center}
\begin{tabular}{c|p{2.9cm}|p{5.0cm}|p{3.1cm}}
\toprule
视角 & 本章看点 & 别的视角怎么说 & 哪侧最省力 \\
\midrule
代数 & 同除配凑、$1$ 的代换、换元 & 符号的恒等变形 & 硬算最值 \\
几何 & 比值 $t=\frac yx$＝射线斜率 & 点 $(x,y)$ 落在一条过原点射线上 & 看比例、判范围 \\
集合 & 齐次式按次数分层 & 同次＝缩放不变的性质 & 界次数与符号 \\
映射 & 齐次坐标 $\to$ 射影直线 & $(u,v)\sim\lambda(u,v)$，$x=u/v$ & 升到第 9 章接口 \\
\bottomrule
\end{tabular}
\end{center}
卡的使用：二元分式最值先数次数、再定「$1$ 的代换还是 $t=\frac yx$」，最后回代数收尾（总图见第 0 章图~\ref{fig:four-view-map}）。
\end{methodbox}

\begin{warmupbox}{读法（先看这里）}
齐次化不是「大招」，而是一句结构主义的第一性原理：\textbf{把
$n$ 个自由变量的绝对尺度约掉，只留下 $n-1$ 个比例。}
遇到二元（$x,y$）的对称式，先问一句：它们能不能合并成一个
分数 $t=\dfrac{y}{x}$？能，就降维；不能，就先配平次数。

（本章与第 5、6 章互为表里：第 6 章「条件最值四大标准动作」给出
「1 的代换」这一入口，第 5 章用矩阵特征方向处理分式递推；本章把
两者背后的原理——齐次化——上升为第一性原理并展开成完整方法链。）
\end{warmupbox}

\begin{corebox}{第一性原理}
设 $P(x,y)=0$ 是一条曲线方程，$Q(x,y)$ 是待求的目标式。若存在
正数 $\lambda$ 使得 $Q(\lambda x,\lambda y)=\lambda^{d}\,Q(x,y)$，
则 $Q$ 是 $d$ 次齐次的，它的取值只依赖于比例 $t=y/x$，与点的
绝对位置无关。非齐次的问题之所以难，是因为混进了不同次数的
项——\textbf{齐次化就是把这些项强行配到同一个次数上}。
\end{corebox}

\begin{methodbox}{次数扫描雷达（动手前先扫一遍）}
给每一项标一个「次数」，像雷达一样先扫描整个式子：
\begin{center}
\begin{tikzpicture}[scale=0.9, font=\small]
\draw[->] (-0.4,0)--(9.6,0) node[right] {$x$ 轴：项};
\draw[->] (0,-0.5)--(0,3.4);
\foreach \x/\h/\lab/\deg [count=\i] in {
  1.4/0.5/常数项/0,
  3.2/1.0/一次项/1,
  5.0/1.5/二次项/2,
  6.8/2.0/分式项/-1,
  8.6/2.5/根式项/{1/2}} {
  \draw[main, thick] (\x,0)--(\x,\h);
  \node[above] at (\x,\h) {\lab};
  \node[below] at (\x,-0.18) {$\deg$ 次};
}
\node[below] at (5.0,-0.66) {（柱高仅为扫描顺序，次数见下方标签）};
\end{tikzpicture}
\end{center}
（常数 $0$ 次、$x$ 与 $y$ 各 $1$ 次、$x^2/xy/y^2$ 为 $2$ 次、
$\dfrac1x$ 为 $-1$ 次、$\sqrt{x}$ 为 $\tfrac12$ 次。）
\textbf{判据}：一个式子是齐次的，当且仅当所有项次数相同。
\end{methodbox}

\begin{methodbox}{方法归纳：同除主元最高次幂（降维的钥匙）}\label{met:homo-divide}
若已经确认各项同次，就同除「主元」的最高次幂，把 $n$ 个变量
压成比例。例如二元齐次式同除 $x^2$（同除 $y^2$ 亦可，此时
比例换成 $x/y=1/t$），只剩
$t=\dfrac{y}{x}$：
\[
\frac{x^2+2xy+y^2}{x^2}=1+2t+t^2 .
\]
一元化之后，剩下的就是单变量函数的最值问题——\textbf{这就是
齐次化的全部收益：把二元压成一元}。
\end{methodbox}

\subsection{第 1 节 常规通法与齐次化的第一次对照}

\begin{example}[已知 $x>0,\ y>0,\ 2x+y=1$，求 $xy$ 的最大值]
\vspace{-1mm}
\begin{solution}
\textbf{考纲通法}：$y=1-2x$，故 $xy=x(1-2x)=-2x^2+x$，
开口向下，顶点在 $x=\dfrac14$，最大值为 $\tk{\dfrac18}$。

\textbf{齐次化视角}：目标与约束同时齐次化，把分母补成
\[
\frac{xy}{(2x+y)^2}=\frac{xy}{1}.
\]
令 $t=\dfrac{y}{x}>0$，则上式右端化为
\[
\frac{t}{(2+t)^2}.
\]
由 AM–GM（$(t+2)^2\ge 8t$），有 $\dfrac{t}{(2+t)^2}\le\dfrac18$，
等号在 $t=2$ 即 $y=2x$、$x=\dfrac14,\ y=\dfrac12$ 时成立。
两种算法给出同一个答案 $\tk{\dfrac18}$。
\end{solution}
\end{example}

\begin{conccard}{齐次化入口：把约束等式当 1 乘进目标式}{二元对称式的最值、范围题}{只会代入消元，不会「补次数」}
约束条件「$2x+y=1$」提供了 \textbf{一个} 齐次化入口：任何
目标式都可以乘上 $(2x+y)^k$ 而不改变取值，从而把混次的式子
全部抬到同一次数。\vchoice{乘方指数 $k$ 的选取原则是「让所有项次数相等」。}{看到等式条件，就想想能不能拿它当 1 来乘。}
\end{conccard}

\subsection{第 2 节 五大配平手术刀}

\begin{methodbox}{手术刀之一「1 的代换」升次}
条件等式本身等于 $1$（或可化成 $1$），把它乘进目标式，
低次项立刻被抬成齐次式。
\[
\left(\frac1x+\frac1y\right)\cdot 1
=\left(\frac1x+\frac1y\right)(x+y)
=2+\frac{y}{x}+\frac{x}{y}\ge 4 .
\]
\end{methodbox}

\begin{methodbox}{手术刀之二 平移原点，把方程归一}
把原点移到图形的中心（或对称中心），一次项/常数项会被
消掉，方程立刻变成标准型。
\[
x^2+y^2-2x=0\ \xrightarrow{\ X=x-1,\ Y=y\ }\ X^2+Y^2=1,
\]
于是 $x+y=X+Y+1\in[\tk{1-\sqrt2},\ \tk{1+\sqrt2}]$。
\end{methodbox}

\begin{methodbox}{手术刀之三 平方升次，剥掉根号}
$\sqrt{x}$ 的次数是 $\tfrac12$，平方一次就升到整次。
已知 $\sqrt{x}+\sqrt{y}=1$，则
\[
x+y=(\sqrt{x}+\sqrt{y})^2-2\sqrt{xy}=1-2\sqrt{xy},
\]
再用 $\sqrt{x}\cdot\sqrt{y}\le\left(\dfrac{\sqrt{x}+\sqrt{y}}{2}\right)^2=\dfrac14$
得 $x+y\ge\tk{\dfrac12}$。
\end{methodbox}

\begin{methodbox}{手术刀之四 分解与组合降次}
高次式先分解，往往能拆出可以约分的因式。
\[
\frac{1}{x(x+1)}=\frac1x-\frac1{x+1},\qquad
x^3+y^3=(x+y)(x^2-xy+y^2).
\]
\end{methodbox}

\begin{methodbox}{手术刀之五 旋转消交叉项}
二次曲线
\[
Ax^2+Bxy+Cy^2+\cdots=0
\]
的交叉项 $Bxy$ 用平移消不掉，换一招：把坐标轴旋转角度 $\theta$，
\[
x=x'\cos\theta-y'\sin\theta,\qquad y=x'\sin\theta+y'\cos\theta .
\]
代回后令 $x'y'$ 的系数为 $0$，即得消交叉项的条件
\[
\cot 2\theta=\frac{A-C}{B}
\qquad\text{（等价地 }\tan 2\theta=\frac{B}{A-C}\text{；当 }A=C\text{ 时用前者）}.
\]
特例 $xy=1$：这里 $A=C=0,\ B=1$，故 $\cot 2\theta=0$，取 $\theta=45^\circ$，
\[
xy=1\ \xrightarrow{\ \theta=45^\circ\ }\ \frac{x'^2}{2}-\frac{y'^2}{2}=1,
\]
等轴双曲线还原为标准型（$a=b=\sqrt2$）。
\textbf{完整推导、图与数列中的复特征值应用见第 5 章。}
\end{methodbox}

\begin{center}
\begin{tikzpicture}[scale=0.85, font=\small]
% 椭圆平移齐次化示意：从原点出发的斜率射线
\draw[gray!60, thick] (0,0) ellipse (3.2 and 2.0);
\draw[->, main, very thick] (0,0)--(3.6,2.4);
\draw[dashed, third] (0,0)--(3.6,-1.2);
\node[right] at (3.65,2.4) {$k=\dfrac{y}{x}>0$};
\node[right] at (3.65,-1.2) {$k<0$};
\fill[second] (0,0) circle (2pt) node[below left] {新原点};
\node at (0,-2.5) {$k=\dfrac{y}{x}$ 的符号决定直线落在哪个象限};
\end{tikzpicture}
\end{center}

\begin{conccard}{五把配平手术刀：升次、平移、平方、分解、旋转}{混次式、根式、分式迭代、含交叉项的二次曲线}{只背大招名称，不动手配平次数}
五把手术刀的共同目标是把方程「摆正」：前四把在\textbf{次数}上
配平，第五把在\textbf{坐标方向}上配平，消掉交叉项 $Bxy$。
判断该用哪一把，只需回看次数扫描的结果——缺几次就补几次，
多几次就降几次；遇到交叉项 $Bxy$ 就转坐标轴（第 5 章有完整推导），
直到方程露出标准型。
\vchoice{严格地说，配平后问题就化为关于 $t=y/x$ 的一元函数，导数或基本不等式都可直接上手。}{口诀：次数不齐先配平，同除主元变一元；方向不正先旋转。}
\end{conccard}

\begin{advancedonly}
\subsection{第 3 节 射影桥梁：为什么齐次化天然有效}

第 5 章探究三已经证明：分式迭代（莫比乌斯变换）
\[
x_{n+1}=\frac{a\,x_n+b}{c\,x_n+d}\qquad(ad-bc\neq0)
\]
的不动点恰是矩阵 $\begin{pmatrix}a&b\\c&d\end{pmatrix}$ 的特征方向，
特征值之比就是双不动点构造的公比（推导与图见第 5 章
图~\ref{fig:eigen-fixed}；不动点的一般理论与双不动点的几何图见第 9 章
图~\ref{fig:mobius-two-fixed}）。这里只补上齐次化的动机：

\textbf{齐次式天生属于射影。} 把 $(u,v)$ 与 $\lambda(u,v)$ 视为同一个
射影点，比值 $x=u/v$ 就是它的坐标；分子、分母同次的式子在这条射线
上取确定值。于是「找不动点」在射影语言里就是「求特征方向」——
齐次化把非线性的分式迭代重新线性化了。

\begin{example}[设 $x_{n+1}=\dfrac{2x_n+1}{x_n+2}$，$x_0=1$，求 $\displaystyle\lim_{n\to\infty}x_n$]
\vspace{-1mm}
\begin{solution}
矩阵 $M=\begin{pmatrix}2&1\\1&2\end{pmatrix}$ 的特征值为
$\lambda_1=3,\lambda_2=1$，特征向量分别沿
$\binom{1}{1}$ 与 $\binom{1}{-1}$，对应不动点 $x=1$ 与 $x=-1$。
$x_0=1$ 恰在 $x=1$ 处，故极限为 $\tk{1}$。
\end{solution}
\end{example}
\end{advancedonly}

\begin{basiconly}
\subsection{第 3 节 大白话版：比例不变，大小随意}

把 $(x,y)$ 想成一条过原点的射线上的点。把 $x,y$ 同时放大
$k$ 倍，点在射线上滑动，但「谁是谁的几倍」没变。齐次式
只认这条射线，不认点在射线上的位置——所以它可以只用一个
字母 $t$ 说话。

\textbf{照样子填空}：$\dfrac{xy}{(2x+y)^2}$ 的分子分母同除
$x^2$ 后得到 $\dfrac{\underline{\quad t\quad}}{(2+\underline{\quad t\quad})^2}$。
\end{basiconly}

\subsection{第 4 节 跨章实战：五题对照（含三道双解）}

\begin{problem}[闪电战 1] 已知 $x>0,\ y>0,\ x+2y=1$，求 $x^2+y^2$ 的最小值。\end{problem}
\begin{solution}
齐次化：$(x^2+y^2)(x+2y)^2=x^2+y^2$。令 $t=\dfrac{y}{x}$，
左边化为 $\dfrac{1+t^2}{(1+2t)^2}$。由柯西不等式
$(1+2t)^2\le(1^2+2^2)(1+t^2)$，得
$\dfrac{1+t^2}{(1+2t)^2}\ge\dfrac15$，等号在 $t=2$
即 $y=2x$、$x=\dfrac15,\ y=\dfrac25$ 处取得。最小值为 $\tk{\dfrac15}$。
\end{solution}

\begin{problem}[闪电战 2] 已知 $x^2+y^2-2x-4y=0$，求 $x+y$ 的取值范围。\end{problem}
\begin{solution}
配方：$(x-1)^2+(y-2)^2=5$，令
$X=x-1,\ Y=y-2$，则 $X^2+Y^2=5$，
$x+y=X+Y+3\in[3-\sqrt{10},\ 3+\sqrt{10}]$。
\end{solution}

\begin{problem}[双解对照 1] 已知 $x>0,\ y>0,\ \dfrac1x+\dfrac2y=1$，求 $x+y$ 的最小值。\end{problem}
\begin{solution}
\textbf{通法}：$x+y=(x+y)\left(\dfrac1x+\dfrac2y\right)
=3+\dfrac{y}{x}+\dfrac{2x}{y}\ge3+2\sqrt2$。

\textbf{齐次化}：令 $t=\dfrac{y}{x}$，约束化为
$\dfrac1x+\dfrac{2}{tx}=1\Rightarrow x=1+\dfrac2t$，
于是 $x+y=x(1+t)=\left(1+\dfrac2t\right)(1+t)
=1+t+\dfrac2t+2\ge3+2\sqrt2$。两种写法同解。
\end{solution}

\begin{problem}[双解对照 2] 曲线 $xy=1$ 与直线 $y=kx+2$ 有交点，求 $k$ 的范围。\end{problem}
\begin{solution}
代入得 $kx^2+2x-1=0$。有实根要求判别式
$\Delta=4+4k\ge0\Rightarrow k\ge-1$；$k=0$ 时方程为
$2x-1=0$，对应交点 $\left(\dfrac12,2\right)$，仍在曲线上。
故 $k\in[-1,+\infty)$，最小值在 $k=-1$ 处取到（相切）。
\end{solution}

\begin{problem}[双解对照 3] 已知 $\sqrt{x}+\sqrt{y}=2$，求 $x+y$ 的最小值与最大值。\end{problem}
\begin{solution}
平方：$x+y=4-2\sqrt{xy}$。由
$\sqrt{xy}\le\dfrac{\sqrt{x}+\sqrt{y}}{2}=1$ 得 $x+y\ge2$；
又 $x+y=(\sqrt{x}+\sqrt{y})^2-2\sqrt{xy}\le4$（当一项为 $0$ 时取等）。
故最小值 $\tk{2}$，最大值 $\tk{4}$。
\end{solution}

\begin{inclasstest}{当堂检测（4 题）}
\begin{enumerate}[label=\arabic*., itemsep=2pt]
\item 次数扫描：$\dfrac{x^2}{y}$ 是\underline{\hspace{2em}}次齐次式。
\item 已知 $x>0,\ y>0,\ x+y=1$，则 $xy$ 的最大值是\underline{\hspace{2em}}。
\item 曲线 $x^2+y^2-6x=0$ 平移后的标准方程是\underline{\hspace{2em}}。
\item 已知 $x>0,\ y>0$ 且 $\dfrac1x+\dfrac1y=1$，则 $x+y$ 的最小值是\underline{\hspace{2em}}。
\end{enumerate}
\begin{solution}
\tk{\text{一}}；\tk{\dfrac14}；\tk{(x-3)^2+y^2=9}；\tk{4}。
\end{solution}
\end{inclasstest}

\begin{homework}{分层作业}
\textbf{基础层（步步填空）}：已知 $x>0,\ y>0,\ x+3y=1$，则
$\dfrac1x+\dfrac3y=1+\underline{\quad}+\dfrac{3y}{x}+\underline{\quad}
\ge\underline{\quad}$。

\textbf{提高层（双解对照）}：已知 $x>0,\ y>0,\ 2x+3y=1$，
求 $\dfrac{1}{x}+\dfrac{1}{y}$ 的最小值（写出通法与齐次化两种）。

\begin{solution}
基础层：$\tk{\dfrac{3x}{y}}$，$\tk{9}$，$\tk{16}$。\\
提高层：$\left(\dfrac1x+\dfrac1y\right)(2x+3y)
=5+\dfrac{3y}{x}+\dfrac{2x}{y}\ge5+2\sqrt6$，最小值
$\tk{5+2\sqrt6}$。
\end{solution}
\end{homework}

\begin{commonerror}{易错点：把「同除」用反了方向}
常见的错误是把 $\dfrac{xy}{(2x+y)^2}$ 齐次化时只除一次，
错写成 $\dfrac{y/x}{2+y/x}$——\textbf{分母少除了一个 $x$}。
正确写法是分子分母\textbf{同除同一个量} $x^2$：
\[
  \dfrac{xy}{(2x+y)^2}=\dfrac{y/x}{(2+y/x)^2}.
\]
次数差才保持不变；除完应当立刻检查「分子还剩几次、分母还剩几次」。
\end{commonerror}

\begin{conccard}[二级]{什么时候必须齐次化，什么时候不必}{条件最值、曲线与直线的位置关系、根式最值}{不管题目长什么样都硬凑 $t=y/x$；或只会代入消元，从不配次数}
先看次数结构：目标式与约束能配成同次（或经平方、平移后可配平）时，
齐次化一步降为一元 $t$；消元后是一元二次、只问有无实根时，判别式
更直接；含 $\sqrt{\phantom{x}}$ 的半次项先平方升次，再决定走哪条路。
\end{conccard}

% ==========================================================
%  第 5 节 七问连环母题（解答即笔记 · Scaffolded Master Problem）
% ==========================================================

\subsection{第 5 节 七问连环母题：解答即笔记}

\begin{warmupbox}{母题怎么用（先看这里）}
本节两道大题，每道挂载七个递进小问。\vchoice{它们不是练习题，
而是「会说话的笔记」：第 (1)–(3) 问替代概念与规范书写笔记，
第 (4)–(6) 问替代大招、公式族与充要边界笔记，第 (7) 问替代
错题本。做完一道，这一章的框架就长在你的解题纸上。}{每道题顺着
七步走即可：先算具体数值，再翻译结构，然后规范推导、降维秒杀、
推广公式、逆向追问，最后排雷。}
\end{warmupbox}

\refstepcounter{hvmasterq}
\begin{masterq}{线性约束 $x+y=1$：从 $\dfrac1x+\dfrac1y$ 长出一族不等式}
设 $x>0,\ y>0,\ x+y=1$。
\begin{enumerate}[label=(\arabic*), itemsep=3pt]
\item \textbf{画图与数值感知}：计算 $(x,y)=\left(\dfrac14,\dfrac34\right)$、
$\left(\dfrac12,\dfrac12\right)$、$\left(\dfrac34,\dfrac14\right)$ 处
$\dfrac1x+\dfrac1y$ 的值，描点并补出草图。\logotag{核心公理}
\item \textbf{多视角结构翻译}：用集合语言写出点集；指出它在坐标系中的
几何形状；再用恒等式把 $\dfrac1x+\dfrac1y$ 改写成只含 $xy$ 的式子。\logotag{结构翻译}
\item \textbf{考纲规范通法}：用基本不等式求 $\dfrac1x+\dfrac1y$ 的最小值，
写出取等条件与关键步骤（用 \tk{} 填空）。\logotag{考纲通法}
\item \textbf{高观点降维秒杀}：令 $t=\dfrac{y}{x}$，把 $\dfrac1x+\dfrac1y$
写成 $t$ 的函数并用一行不等式收束，与第 (3) 问核对。\logotag{大招公式}
\item \textbf{公式变形家族生成}：推广到 $x+y=s$、$ax+by=1$，并推导
$\dfrac1{x^2}+\dfrac1{y^2}$ 的下界。\logotag{公式族}
\item \textbf{参数变式与逆向追问}：已知 $\dfrac1x+\dfrac1y=5$ 且 $x+y=1$，
反求 $x,y$；再回答：使 $\dfrac1x+\dfrac1y\ge m$ 恒成立的 $m$ 最大是多少？\logotag{充要边界}
\item \textbf{边界陷阱与错题排雷}：某同学删去「$x,y>0$」后仍写
$\dfrac1x+\dfrac1y\ge4$。请举反例推翻它，并说明原题有没有最大值。\logotag{避坑铁律}
\end{enumerate}
\end{masterq}

\begin{solution}
(1) $\dfrac{16}{3}$，$\tk{4}$，$\dfrac{16}{3}$；图象关于 $x=\dfrac12$ 对称，
两端趋于 $+\infty$。

(2) $\{(x,y)\mid x>0,\ y>0,\ x+y=1\}$；第一象限内连接 $(1,0)$ 与
$(0,1)$ 的线段；$\dfrac1x+\dfrac1y=\dfrac{x+y}{xy}=\tk{\dfrac{1}{xy}}$。

(3) $xy\le\left(\dfrac{x+y}{2}\right)^2=\dfrac14$，故
$\dfrac1x+\dfrac1y=\dfrac1{xy}\ge\tk{4}$，当且仅当 $x=y=\dfrac12$。

(4) $x=\dfrac{1}{1+t},\ y=\dfrac{t}{1+t}$，于是
$\dfrac1x+\dfrac1y=(1+t)+\dfrac{1+t}{t}=2+t+\dfrac1t\ge\tk{4}$
（$t=1$ 取等），与 (3) 一致。

(5) (a) $x+y=s$：$\dfrac1x+\dfrac1y=\dfrac{s}{xy}\ge\dfrac{s}{s^2/4}
=\tk{\dfrac4s}$；
(b) $ax+by=1$：$\dfrac1x+\dfrac1y=\left(\dfrac1x+\dfrac1y\right)(ax+by)
=a+b+\dfrac{ay}{x}+\dfrac{bx}{y}\ge a+b+2\sqrt{ab}=\tk{(\sqrt a+\sqrt b)^2}$；
(c) $\dfrac1{x^2}+\dfrac1{y^2}=\left(\dfrac1x+\dfrac1y\right)^2-\dfrac{2}{xy}
\ge16-8=\tk{8}$。

(6) $xy=\dfrac15$，故 $x,y$ 是 $z^2-z+\dfrac15=0$ 的两根，即
$\tk{\dfrac{5\pm\sqrt5}{10}}$；$m_{\max}=\tk{4}$（恰为下确界）。

(7) 反例：$x=2,\ y=-1$ 满足 $x+y=1$，但 $\dfrac1x+\dfrac1y=-\dfrac12<4$。
原题中 $x\to0^+$ 时 $\dfrac1x+\dfrac1y\to+\infty$，故只有最小值、无最大值。
\end{solution}

\refstepcounter{hvmasterq}
\begin{masterq}{二次约束 $x^2+y^2=1$：齐次化与「1 的代换」双剑合璧}
设 $x^2+y^2=1\ (x>0,\ y>0)$。
\begin{enumerate}[label=(\arabic*), itemsep=3pt]
\item \textbf{画图与数值感知}：计算 $\left(\dfrac{\sqrt2}{2},\dfrac{\sqrt2}{2}\right)$、
$\left(\dfrac35,\dfrac45\right)$、$\left(\dfrac12,\dfrac{\sqrt3}{2}\right)$ 处
$\dfrac1{x^2}+\dfrac1{y^2}$ 与 $\dfrac{x}{y}+\dfrac{y}{x}$ 的值，并在单位圆
第一象限弧上描点。\logotag{核心公理}
\item \textbf{多视角结构翻译}：用集合 $\{(\cos\theta,\sin\theta)\mid
0<\theta<\tfrac\pi2\}$ 重写点集；指出几何形状；说明 $x^2+y^2=1$ 为什么
可以当作「1」乘进任何式子。\logotag{结构翻译}
\item \textbf{考纲规范通法}：求 $\dfrac1{x^2}+\dfrac1{y^2}$ 的最小值，
写出取等条件（用 \tk{} 填空）。\logotag{考纲通法}
\item \textbf{高观点降维秒杀}：把 $\dfrac1{x^2}+\dfrac1{y^2}$ 配成分子分母
同次（乘「1」再除「1」），令 $t=\dfrac{y}{x}$ 一行出结果，并顺手得到
$\dfrac{x}{y}+\dfrac{y}{x}$ 的下界。\logotag{大招公式}
\item \textbf{公式变形家族生成}：推导 $x^2+y^2=r^2$、$x^2+2y^2=1$ 情形下
$\dfrac1{x^2}+\dfrac1{y^2}$、$\dfrac1{x^2}+\dfrac2{y^2}$ 与
$\dfrac1{x^4}+\dfrac1{y^4}$ 的下界。\logotag{公式族}
\item \textbf{参数变式与逆向追问}：已知 $\dfrac1{x^2}+\dfrac1{y^2}=5$ 且
$x^2+y^2=1$，求 $x^2,y^2$；若去掉「$x,y>0$」只保留 $x^2+y^2=1$，
最小值是否改变？\logotag{充要边界}
\item \textbf{边界陷阱与错题排雷}：错解「$xy\le\dfrac12\Rightarrow
\dfrac1{xy}\ge2\Rightarrow\dfrac1{x^2y^2}\ge4$」的哪一步依赖了符号条件？
写出防错铁律。\logotag{避坑铁律}
\end{enumerate}
\end{masterq}

\begin{solution}
(1) $\left(\dfrac{\sqrt2}{2},\dfrac{\sqrt2}{2}\right)$：$\dfrac1{x^2}+\dfrac1{y^2}=4$，
$\dfrac{x}{y}+\dfrac{y}{x}=2$；
$\left(\dfrac35,\dfrac45\right)$：$\dfrac{25}{9}+\dfrac{25}{16}
=\tk{\dfrac{625}{144}}$，$\dfrac{x}{y}+\dfrac{y}{x}=\dfrac34+\dfrac43
=\tk{\dfrac{25}{12}}$；
$\left(\dfrac12,\dfrac{\sqrt3}{2}\right)$：$4+\dfrac43=\tk{\dfrac{16}{3}}$。

(2) 几何上是单位圆在第一象限的四分之一圆弧；代数上 $x^2+y^2=1$ 是
「1 的代换」的合法来源（乘以它不改变任何式子的取值）。

(3) $\dfrac1{x^2}+\dfrac1{y^2}=\dfrac{x^2+y^2}{x^2y^2}=\dfrac1{x^2y^2}$，
而 $xy\le\dfrac{x^2+y^2}{2}=\dfrac12\Rightarrow x^2y^2\le\dfrac14$，
故原式 $\ge\tk{4}$，等号当且仅当 $x=y=\dfrac{\sqrt2}{2}$。

(4) 乘「1」：$\dfrac1{x^2}+\dfrac1{y^2}
=\dfrac{(x^2+y^2)^2}{x^2y^2}$；令 $t=\dfrac{y}{x}$ 得
$\dfrac{(1+t^2)^2}{t^2}=\left(t+\dfrac1t\right)^2\ge4$；同一行顺手得到
$\dfrac{x}{y}+\dfrac{y}{x}=t+\dfrac1t\ge\tk{2}$。

(5) (a) $x^2+y^2=r^2\Rightarrow\dfrac1{x^2}+\dfrac1{y^2}
=\dfrac{r^4}{x^2y^2}\ge\tk{\dfrac{4}{r^2}}$；
(b) $\left(\dfrac1{x^2}+\dfrac2{y^2}\right)(x^2+2y^2)
=5+2\left(\dfrac{y^2}{x^2}+\dfrac{x^2}{y^2}\right)\ge5+4=9$，
故 $\dfrac1{x^2}+\dfrac2{y^2}\ge\tk{9}$（等号 $x^2=y^2=\dfrac13$）；
(c) $\dfrac1{x^4}+\dfrac1{y^4}\ge\dfrac{2}{x^2y^2}\ge\tk{8}$。

(6) 令 $u=x^2,\ v=y^2$，则 $u+v=1,\ \dfrac{u+v}{uv}=5\Rightarrow uv=\dfrac15$，
故 $u,v=\tk{\dfrac{5\pm\sqrt5}{10}}$；去掉正数限制后结论不变（只含平方项），
但仍需 $x,y\neq0$。

(7) 「$\dfrac1{xy}\ge2$」这一步要求 $xy>0$；若 $xy<0$，不等号方向相反。
本题最终正确只是因为又平方了一次。防错铁律：取倒数或用基本不等式前，
先确认各项同号。
\end{solution}

\begin{basiconly}
\subsection{第 6 节 大白话版：两道母题的半成品脚手架}
母题一：两人分一个蛋糕，$x,y$ 是各拿多少。怎么分，倒数的和最小？
先试四等分点：$x=\dfrac14$ 时 $\dfrac1x+\dfrac1y=\dfrac{16}{3}$；
正中间 $x=\dfrac12$ 时等于 \underline{\hspace{2em}}，看起来最小。
于是猜想：\underline{\hspace{2em}} 时最小，最小值 \underline{\hspace{2em}}。

母题二：单位圆上找一点。试 $x=y=\dfrac{\sqrt2}{2}$，得
$\dfrac1{x^2}+\dfrac1{y^2}=4$；试 $x=\dfrac12,\ y=\dfrac{\sqrt3}{2}$，
得 $\dfrac{16}{3}$。故最小值是 \underline{\hspace{2em}}，
取到在 \underline{\hspace{2em}}。
\end{basiconly}

\begin{advancedonly}
\textbf{射影眼光看这两道母题}：第 (4) 问的比值 $t=y/x$ 是射影直线上的
坐标，齐次式只认射影点、不认代表元（同第 3 节的射影视角）。把母题一
的约束换成 $ax+by=1$ 后，(5)(b) 的取等条件 $\dfrac{ay}{x}=\dfrac{bx}{y}$
给出最优比例 $t^*=\sqrt{b/a}$，即最优射线的方向比为 $\sqrt b:\sqrt a$——
配平问题变成求一条射线与一条直线的交点。
\end{advancedonly}

\begin{conccard}{连环母题：做一题 = 学一节}{概念引入、性质翻译、规范书写、大招结论、公式族、充要边界、错题本七类笔记一次生成}{把母题当普通大题，只做最后一问}
\vchoice{做完这两道题，留白处就是一份《齐次化》随身笔记：第 (1)(2) 问
是概念与性质，第 (3) 问是考场书写，第 (4)(5) 问是大招与公式表，
第 (6) 问是逆向思维，第 (7) 问是错题本。}{基础较弱的同学先完成
第 (1)–(3) 问即可达标；第 (4)–(7) 问留给学有余力者。}
\end{conccard}

\begin{notezone}{章末归档：齐次化二级结论 + 七问笔记功能}
\textbf{二级结论}（方法链压缩版）：（1）齐次式的值只依赖比例 $t=y/x$；
（2）约束等式可以当 $1$ 乘进目标式完成升次（第 6 章入口、本章展开）；
（3）半次项平方可升为整次；（4）平移可以消去一次项，一般理论见
第 5 章；（5）分式迭代的不动点对应矩阵特征方向（第 5 章）。

\textbf{七问与笔记功能}：(1) 画图与数值感知 $\to$ 概念引入笔记；
(2) 三视角结构翻译 $\to$ 性质定义笔记；(3) 考纲规范通法 $\to$ 规范
书写笔记；(4) 高观点降维秒杀 $\to$ 二级结论与大招笔记；(5) 公式变形
家族 $\to$ 题型拓展笔记；(6) 参数变式与逆向追问 $\to$ 充要条件与逻辑
笔记；(7) 边界陷阱排雷 $\to$ 错题本易错点笔记。
\end{notezone}

\reading{拓展阅读 14}{齐次化：比例结构与射影桥梁}

% ==========================================================
% 第 9 章 不动点：迭代、方程与几何变换的交汇
% 说明：本文件为纯正文片段，由 main_body.tex 挂载；
%       使用 versions.tex 的 \tk / \vchoice / advancedonly /
%       basiconly / solution / 盒子与七问母题环境。
% ==========================================================

\section{不动点：迭代、方程与几何变换的交汇（高观点版）}

\begin{methodbox}{多框架对读}
\small
\begin{center}
\begin{tabular}{c|p{2.9cm}|p{5.0cm}|p{3.1cm}}
\toprule
视角 & 本章看点 & 别的视角怎么说 & 哪侧最省力 \\
\midrule
代数 & 解 $x=F(x)$、构造 $b_n=a_n-x^*$ & 把不动点当方程的根 & 求解析式 \\
几何 & 曲线与对角线的交点、变换的不动心 & 位置关系与轨迹形状 & 看收敛、判稳定 \\
集合 & 轨道 $\{a_n\}$ 与其聚点 & 值集、吸引域／排斥域 & 讨论存在与唯一 \\
映射 & 停留点 $F(p)=p$；Möbius 射影 & 换坐标系＝换表示 & 非线性 $\to$ 等比 \\
\bottomrule
\end{tabular}
\end{center}
卡的使用：四类不动点（数列／方程／变换／分式）共用同一句话——$F(p)=p$；算不下去就换表示（总图见第 0 章图~\ref{fig:four-view-map}）。
\end{methodbox}

\begin{warmupbox}{读法（先看这里）}
不动点只有一个定义：映射 $F$ 中满足 $F(p)=p$ 的点 $p$。
它有三副面孔：数列递推 $a_{n+1}=F(a_n)$（离散自迭代）、
方程 $f(x)=x$（曲线与对角线的交点）、几何变换（旋转中心、
位似中心、分式变换）。三者的共同动作是\key{换坐标系，把锚点
搬到原点}——这正是第 8 章齐次化、第 5 章平移／旋转的内核。
本章把这条线索单独展开，第 5 章与第 8 章的应用改为引用。
\end{warmupbox}

\begin{learnobjective}
\begin{itemize}
  \item 把数列看成 $\mathbb{N}\to\mathbb{R}$ 的（离散）函数，把递推看成映射的自迭代；
  \item 会用不动点方程 $x=F(x)$ 与稳定性判据 $|F'(x^*)|$ 处理一阶递推；
  \item 理解连续情形 $f(x)=x$ 的图解法与第 6 章隐零点、切线放缩的联系；
  \item 知道旋转中心、位似中心、仿射定点都是不动点的特例；
  \item 理解齐次坐标 $(u,v)$ 为什么能把分式迭代线性化（详见第 5 章）。
\end{itemize}
\end{learnobjective}

\subsection{第 1 节 数列是自然数集上的函数：递推即迭代}

数列 $a_n$ 本质是函数 $a:\mathbb{N}\to\mathbb{R}$，$n\mapsto a_n$。
一阶递推 $a_{n+1}=F(a_n)$ 就是映射 $F$ 的\key{自迭代}：
$a_1=F(a_0),\ a_2=F(F(a_0)),\dots$，数列是 $F$ 的一条轨道。
于是
\[
  \text{不动点}\ x^*=F(x^*)
  \iff \text{点 }(x^*,x^*)\text{ 同时在 }F\text{ 的图象与对角线 }y=x\text{ 上}.
\]
\textbf{稳定性判据}：在 $x^*$ 附近线性化 $F(x)\approx x^*+F'(x^*)(x-x^*)$，
故
\[
  |F'(x^*)|<1\Rightarrow\text{收缩（稳定）},\quad
  |F'(x^*)|>1\Rightarrow\text{扩张（不稳定）},\quad
  |F'(x^*)|=1\Rightarrow\text{看高阶项（临界）}.
\]

\begin{figure}[htbp]
  \centering
  \begin{minipage}{0.48\textwidth}
    \centering
    \begin{tikzpicture}[x=0.9cm,y=0.9cm,>=Stealth,semithick,font=\footnotesize]
      \draw[->,gray] (-1,0)--(3,0);
      \draw[->,gray] (0,-1)--(0,3);
      \draw[gray!60,dashed,domain=-1:3] plot(\x,\x);
      \draw[main,very thick,domain=-1:3] plot(\x,{0.5*\x+1});
      \fill[second] (2,2) circle (1.6pt) node[above right] {$x^*=2$};
      \draw[third,thick] (0,0)--(0,1)--(1,1)--(1,1.5)--(1.5,1.5)--(1.5,1.75)--(1.75,1.75);
      \node[gray!70!black] at (1.5,-0.65) {$F(x)=\tfrac12x+1$：收缩};
    \end{tikzpicture}
  \end{minipage}\hfill
  \begin{minipage}{0.48\textwidth}
    \centering
    \begin{tikzpicture}[x=0.9cm,y=0.9cm,>=Stealth,semithick,font=\footnotesize]
      \draw[->,gray] (-1,0)--(3,0);
      \draw[->,gray] (0,-1)--(0,3);
      \draw[gray!60,dashed,domain=-1:3] plot(\x,\x);
      \draw[main,very thick,domain=-0.5:2] plot(\x,{2*\x-1});
      \fill[second] (1,1) circle (1.6pt) node[below left] {$x^*=1$};
      \draw[third,thick] (1.2,1.2)--(1.2,1.4)--(1.4,1.4)--(1.4,1.8)--(1.8,1.8)--(1.8,2.6);
      \node[gray!70!black] at (1.5,-0.65) {$G(x)=2x-1$：扩张};
    \end{tikzpicture}
  \end{minipage}
  \caption{不动点的两种命运：同一张「$F$ 的图象＋对角线」图上，左图 $|F'|<1$，蛛网阶梯螺旋收拢到 $x^*$；右图 $|G'|>1$，阶梯被弹开。斜率与对角线 $y=x$ 的比较就是稳定性判据}
  \label{fig:cobweb-pair}
\end{figure}

\begin{methodbox}{不动点法三步}
\begin{enumerate}[label=(\arabic*), itemsep=2pt]
  \item \key{找}：解 $x^*=F(x^*)$（可能不止一个）；
  \item \key{判}：算 $|F'(x^*)|$ 定稳定／不稳定；
  \item \key{化}：令 $b_n=a_n-x^*$，则 $b_{n+1}\approx F'(x^*)\,b_n$；
        稳定时 $b_n\to0$，并可用等比模型写出通项。
\end{enumerate}
常见变体：分式递推用 $\dfrac{a_n-x_1}{a_n-x_2}$ 作商（第 5 章探究五），
双曲型递推用三角／指数代换。
\end{methodbox}

\subsection{第 2 节 连续情形：方程 $f(x)=x$ 与图解法}

把 $f(x)=x$ 改写成 $\varphi(x)=f(x)-x=0$：不动点就是 $\varphi$ 的零点。
几何上，它是曲线 $y=f(x)$ 与直线 $y=x$ 的交点。
\begin{example}[方程 $x=\cos x$ 在 $(0,\tfrac\pi2)$ 内有唯一解]
\vspace{-1mm}
\begin{solution}
$\varphi(x)=\cos x-x$，$\varphi(0)=1>0$，$\varphi(\tfrac\pi2)=-\tfrac\pi2<0$，
由零点定理至少有一解；又 $\varphi'(x)=-\sin x-1<0$ 严格递减，故唯一。
进一步 $|\varphi'|$ 与压缩性：$f(x)=\cos x$ 的导数 $|f'|=\sin x<1$ 在
$(0,\tfrac\pi2)$ 成立，故 $f$ 是压缩映射，迭代 $x_{n+1}=\cos x_n$
对任意初值都收敛到该唯一不动点——这就是第 6 章「隐零点＋单调性」
的几何版本。
\end{solution}
\end{example}
\textbf{与考纲的接口}：求超越方程的近似解（二分法）本质是不断把区间
压缩到不动点；证明 $f(x)\ge x$ 型不等式（如 $e^x\ge x+1$）则是在
比较曲线与对角线的位置（见第 6 章切线放缩）。

\subsection{第 3 节 几何变换的不动点：旋转中心、位似中心与仿射定点}

\begin{itemize}
  \item \key{旋转}：转角 $\theta\not\equiv0\pmod{2\pi}$ 的平面旋转只有
        \textbf{中心}一个不动点；
  \item \key{位似}：位似中心不动，其余点沿射线伸缩；
  \item \key{平移}：平面上没有不动点（它的「不动点」在无穷远处）；
  \item \key{仿射变换} $(x,y)\mapsto A\binom xy+\vec b$：不动点满足
        $(A-I)\binom xy=-\vec b$，当 $\det(A-I)\neq0$ 时唯一。
        椭圆的中心就是两方向伸缩变换的公共不动点（见第 4 章）。
\end{itemize}
直觉：任何「保结构」的变换都可以先找到它的不动点，再把坐标原点
搬过去——剩下的表达式往往立刻简化（第 8 章的平移手术刀、第 5 章
的旋转消交叉项都是这一招）。

\subsection{第 4 节 分式迭代与齐次坐标：射影视角}

一次分式变换 $x\mapsto\dfrac{ax+b}{cx+d}$ 本身不是线性的，但把 $x$
写成两个量的比 $x=u/v$ 之后，分子分母各自线性组合：
\[
  x_{n+1}=\frac{a\,x_n+b}{c\,x_n+d}
  \iff
  \binom{u_{n+1}}{v_{n+1}}
  =\begin{pmatrix}a&b\\ c&d\end{pmatrix}\binom{u_n}{v_n},
  \qquad x_n=\frac{u_n}{v_n}.
\]
关键在第二步：点 $(u,v)$ 与 $\lambda(u,v)$ 表示同一个 $x$（射影点），
所以「二维向量的线性作用」可以「投影」为一维分式作用。第 5 章
探究三已由此证明：\key{分式迭代的不动点 $=$ 矩阵的特征方向}
（推导与数值例子见第 5 章，此处不重复）。这里只看一件事：为什么
$(u,v)$ 与 $\lambda(u,v)$ 是同一点。

\begin{figure}[htbp]
  \centering
  \begin{tikzpicture}[>=Stealth,semithick,font=\footnotesize,x=1.1cm,y=1.1cm]
    \draw[->,gray] (-0.5,0)--(4.6,0) node[right] {$u$};
    \draw[->,gray] (0,-0.5)--(0,3.1) node[above] {$v$};
    \draw[main!60!black,dashed,very thick] (0,1)--(4.5,1) node[right] {$v=1$（投影线）};
    % 两条射线：同一条射线上取两个代表元
    \draw[second,very thick,->] (0,0)--(2.1,1.05);
    \draw[third,very thick,->] (0,0)--(1.05,2.1);
    \foreach \p/\lab/\pos in {(2,1)/{$x=2$}/above right,(1,2)/{$x=\tfrac12$}/above left}
      \fill[main] \p circle (1.6pt) node[\pos] {\lab};
    \fill (0,0) circle (1.2pt) node[below left] {$O$};
    \node[gray!70!black,align=center] at (2.2,-0.7) {同一条过原点的射线上的 $(u,v)$、$(2u,2v)$\\投影到 $v=1$ 上得到同一个 $x=u/v$};
  \end{tikzpicture}
  \caption{齐次坐标：用有序数对 $(u,v)$ 刻画比值 $x=u/v$。射线上的所有代表元落在投影线 $v=1$ 的同一点上，于是「一维分式迭代」升级为「二维线性迭代」}
  \label{fig:ray-projective}
\end{figure}

\textbf{双不动点的几何图像}：设 $M$ 的两个不动点为 $p,q$（射影直线上）。
取新坐标
\[
  X=\frac{x-p}{\,x-q\,},
\]
它把 $p$ 送到 $0$、$q$ 送到 $\infty$；在此坐标下 $M$ 退化为\emph{纯伸缩}
$X\mapsto\lambda X$，其中 $\lambda$ 正是转移矩阵两特征值之比（第 5 章探究五）。
于是
\[
  |\lambda|<1\Rightarrow X\to0\ \ (x\to p),\qquad
  |\lambda|>1\Rightarrow X\to\infty\ \ (x\to q).
\]
换言之：\key{两条特征方向在射影直线上就是两个不动点}；以 $p,q$ 为
$0,\infty$ 建立坐标后，分式迭代变成等比数列——这就是共轭变换的几何全貌。

\begin{figure}[htbp]
  \centering
  \begin{minipage}{0.48\textwidth}
    \centering
    \begin{tikzpicture}[>=Stealth,semithick,font=\footnotesize,x=1.2cm,y=1.2cm]
      \draw[->,gray] (-3.2,0)--(3.6,0);
      \node[gray!70!black,right] at (3.6,0) {$x$（射影直线）};
      \fill[main] (-2,0) circle (1.5pt) node[below] {$p$};
      \fill[second] (1.4,0) circle (1.5pt) node[below] {$q$};
      \draw[third,thick,->] (-1.2,0.28)--(-0.45,0.28);
      \draw[third,thick,->] (0.1,0.28)--(0.9,0.28);
      \node[gray!70!black,align=center] at (0.2,-1.0)
        {区间 $(p,q)$ 内 $X<0$，$|X|$ 越滚越大，\\$x$ 被拉向 $q$（$|\lambda|>1$）};
    \end{tikzpicture}
  \end{minipage}\hfill
  \begin{minipage}{0.48\textwidth}
    \centering
    \begin{tikzpicture}[>=Stealth,semithick,font=\footnotesize,x=1.0cm,y=1.0cm]
      \draw[->,gray] (-2.4,0)--(2.6,0) node[right] {$u$};
      \draw[->,gray] (0,-0.8)--(0,2.6) node[above] {$v$};
      \draw[main!60!black,dashed,thick] (-2.2,1)--(2.4,1) node[right] {$v=1$};
      \node[gray!70!black,left] at (-0.02,1) {$1$};
      \draw[main,very thick,->] (0,0)--(-1.05,0.55);
      \draw[main,very thick,dashed] (0,0)--(-2.1,1.05);
      \draw[second,very thick,->] (0,0)--(2.4,1.2);
      \fill[main] (-0.6,0.3) circle (1.4pt) node[below] {$(u,v)$};
      \node[main!70!black,above] at (-1.5,0.82) {$\times\lambda$};
      \fill[third] (-2,1) circle (1.4pt) node[below] {$p$};
      \fill[third] (2,1) circle (1.4pt) node[below] {$q$};
      \node[gray!70!black,align=center] at (0,-0.55)
        {两条特征直线 $\leftrightarrow$ 两个不动点};
    \end{tikzpicture}
  \end{minipage}
  \caption{双不动点的几何：左图在射影直线上，$X=(x-p)/(x-q)$ 使 $p\leftrightarrow0$、
  $q\leftrightarrow\infty$，迭代退化为 $X\mapsto\lambda X$（图示 $|\lambda|>1$，轨道被拉向 $q$）；
  右图回到 $(u,v)$ 平面，两条过原点的特征直线分别对应 $p,q$，作用即沿射线作伸缩}
  \label{fig:mobius-two-fixed}
\end{figure}

\subsection{第 5 节 七问连环母题：解答即笔记}

\refstepcounter{hvmasterq}
\begin{masterq}{不动点与收敛：$a_{n+1}=\dfrac12a_n+1$，$a_0=0$}
\begin{enumerate}[label=(\arabic*), itemsep=3pt]
\item \textbf{画图与数值感知}：算出 $a_1,a_2,a_3,a_4$，描点并猜极限。\logotag{核心公理}
\item \textbf{多视角结构翻译}：把该递推分别看成（i）$\mathbb{N}\to\mathbb{R}$
的函数迭代、（ii）平面直角坐标系中曲线 $y=\tfrac12x+1$ 与直线 $y=x$
的交点问题。\logotag{结构翻译}
\item \textbf{考纲规范通法}：用不动点法求通项，写出关键步骤
（令 $b_n=a_n-x^*$，\tk{} 填空）。\logotag{考纲通法}
\item \textbf{高观点降维秒杀}：用 $|F'(x^*)|$ 直接写出
$|a_n-x^*|=2\cdot2^{-n}$，并与第 6 章局部线性化对读。\logotag{大招公式}
\item \textbf{公式变形家族生成}：推导 $a_{n+1}=pa_n+q$ 的完整公式族
（$p\neq1$、$p=1$、$|p|>1$ 三种情形）。\logotag{公式族}
\item \textbf{参数变式与逆向追问}：由 $a_2=\tfrac32$、$a_3=\tfrac74$
反求 $p,q$；再回答：若 $\lim a_n=5$ 且 $p=\tfrac12$，求 $q$。\logotag{充要边界}
\item \textbf{边界陷阱与错题排雷}：某同学习惯「设 $\lim a_n=L$，代入得
$L=\tfrac12L+1$」，请指出缺少哪个前提，并举反例说明。\logotag{避坑铁律}
\end{enumerate}
\end{masterq}

\begin{solution}
(1) $a_1=1,\ a_2=\dfrac32,\ a_3=\dfrac74,\ a_4=\dfrac{15}8$，
趋向 $2$。

(2) （i）$F(x)=\tfrac12x+1$ 的自迭代；（ii）交点为
联立 $y=\tfrac12x+1$ 与 $y=x$ 的解 $(2,2)$。

(3) $x^*=\tfrac12x^*+1\Rightarrow x^*=2$。令 $b_n=a_n-2$，则
$b_{n+1}=a_{n+1}-2=\tfrac12(a_n-2)=\tfrac12 b_n$，$b_0=-2$，故
$b_n=-2\left(\tfrac12\right)^n$，
\[
  a_n=2-2^{\,1-n}.
\]
检验：$n=0$ 时 $a_0=0$，$n=1$ 时 $a_1=1$。

(4) $|F'(2)|=\tfrac12<1$，故 $|a_n-2|=|b_n|=2\cdot2^{-n}$ 指数衰减；
第 6 章的局部线性化 $F(x)\approx 2+\tfrac12(x-2)$ 正是这一步的连续版。

(5) $a_{n+1}=pa_n+q$：$x^*=px^*+q\Rightarrow x^*=\dfrac{q}{1-p}\ (p\neq1)$，
$a_n=x^*+p^n(a_0-x^*)$；$p=1$ 时退化为等差 $a_n=a_0+nq$；
$|p|>1$ 且 $a_0\neq x^*$ 时 $|a_n|\to\infty$。

(6) 由 $a_2=pa_1+q$、$a_3=pa_2+q$ 相减得
$p=\dfrac{a_3-a_2}{a_2-a_1}=\dfrac{1/4}{1/2}=\dfrac12$，
$q=a_2-pa_1=\dfrac32-\dfrac12=1$。
若 $\lim a_n=5$ 且 $p=\tfrac12$，则 $5=\dfrac q{1-1/2}=2q$，故 $q=\dfrac52$。

(7) 「取极限」需要先证明极限存在（单调有界或压缩映射）；反例：
$a_{n+1}=-a_n$、$a_0=1$ 振荡无极限，而方程 $x=-x$ 有唯一解 $0$——
\textbf{不动点存在不等于轨道收敛}。
\end{solution}

\refstepcounter{hvmasterq}
\begin{masterq}{分式迭代：$x_{n+1}=1+\dfrac1{x_n}$，$x_0=1$}
\begin{enumerate}[label=(\arabic*), itemsep=3pt]
\item \textbf{画图与数值感知}：算出 $x_1,x_2,x_3,x_4$，观察奇数项与
偶数项各自的走向。\logotag{核心公理}
\item \textbf{多视角结构翻译}：把它写成连分数 $x=1+\dfrac1{1+\dfrac1{1+\cdots}}$
与二次方程 $x^2-x-1=0$。\logotag{结构翻译}
\item \textbf{考纲规范通法}：证明奇数子列递增、偶数子列递减且有共同
极限 $\varphi=\dfrac{1+\sqrt5}{2}$（用 $\tk{}$ 搭骨架）。\logotag{考纲通法}
\item \textbf{高观点降维秒杀}：指出这是矩阵 $\begin{pmatrix}1&1\\1&0\end{pmatrix}$
的迭代，特征值之比给出收敛速率 $\dfrac{1}{\varphi^2}$（见第 5 章）。\logotag{大招公式}
\item \textbf{公式变形家族生成}：证明 $x_n=\dfrac{F_{n+2}}{F_{n+1}}$（$F_n$ 为
斐波那契数），并由此得到 $\varphi$ 的连分数与渐进分数。\logotag{公式族}
\item \textbf{参数变式与逆向追问}：若 $x_0=\varphi$ 会怎样？求使 $x_2=\varphi$
的 $x_0$；再把递推改成 $x_{n+1}=k+\dfrac1{x_n}$，求不动点与收敛范围。\logotag{充要边界}
\item \textbf{边界陷阱与错题排雷}：某同学由 $\lim x_n=L$ 直接写
$L=1+\dfrac1L$。请指出两处逻辑漏洞（存在性、正负号），并举出
$x_0<0$ 时的反例。\logotag{避坑铁律}
\end{enumerate}
\end{masterq}

\begin{solution}
(1) $x_1=2,\ x_2=\dfrac32,\ x_3=\dfrac53,\ x_4=\dfrac85$；奇数项递减、
偶数项递增，夹在 $1.6$ 附近。

(2) 迭代不变量：若 $x_n>0$，则 $x_{n+1}>1$；不动点满足
$x=1+\dfrac1x\Rightarrow x^2-x-1=0\Rightarrow \varphi=\dfrac{1+\sqrt5}{2}$
（负根 $-\dfrac1\varphi<0$ 舍去，因 $x_n>0$）。

(3) 设 $f(x)=1+\dfrac1x$，则 $x_{n+2}=f(f(x_n))=\dfrac{2x_n+1}{x_n+1}$。
直接证明：$x_{2k}$ 递增、$x_{2k+1}$ 递减，且
$|x_{n+1}-\varphi|=\dfrac{|x_n-\varphi|}{x_n\varphi}\le\dfrac1\varphi|x_n-\varphi|$
（因为 $x_n\ge1$），故两者夹逼到 $\varphi$。

(4) $\binom{u_{n+1}}{v_{n+1}}=\begin{pmatrix}1&1\\1&0\end{pmatrix}\binom{u_n}{v_n}$，
$x_n=u_n/v_n$；矩阵特征值为 $\varphi$ 与 $-1/\varphi$，故
$|x_n-\varphi|\sim C\varphi^{-2n}$（公比 $\left|\lambda_2/\lambda_1\right|
=\dfrac1{\varphi^2}$，见第 5 章探究五）。

(5) 归纳（取 $F_0=0,F_1=1$）：$x_0=\dfrac{F_2}{F_1}=\dfrac11=1$，且
$1+\dfrac1{x_n}=1+\dfrac{F_{n+1}}{F_{n+2}}
=\dfrac{F_{n+2}+F_{n+1}}{F_{n+2}}=\dfrac{F_{n+3}}{F_{n+2}}$，
故 $x_n=\dfrac{F_{n+2}}{F_{n+1}}$。

(6) $x_0=\varphi\Rightarrow$ 常数列 $\varphi$；反向递推
$x_n=\dfrac1{x_{n+1}-1}$，由 $x_2=\varphi$ 得 $x_1=\varphi$、$x_0=\varphi$。
一般地 $x_{n+1}=k+\dfrac1{x_n}$ 的不动点为
$\dfrac{k+\sqrt{k^2+4}}2$，$k>0$ 时正根稳定（$|f'|=1/x^{*2}<1$）。

(7) 漏洞一：未先证明极限存在；漏洞二：取极限时未说明保号。
方程 $L=1+\dfrac1L$ 的负根 $-\dfrac{1-\sqrt5}{2}$ 本身合法，
但它与 $x_n>0$ 的保号性矛盾，不能作为数列的极限。反例：
$x_0=-1$ 时 $x_1=0$，$x_2$ 无定义——定义域与符号必须先交代。
\end{solution}

\begin{conccard}{不动点：一把万能锚}{递推求通项与收敛性、超越方程近似解、几何变换的简化、分式迭代}{把「不动点存在」当成「轨道收敛」；未验证极限存在就取极限；忽略定义域与符号}
\vchoice{遇到「迭代／递推／方程 $f(x)=x$／变换简化」，先找不动点，
再判 $|F'|$，最后决定坐标系搬到哪里。}{记住三步：解 $x=F(x)$，
算 $|F'|$，搬原点。第 5、6、8 章都靠这一招联动。}
\end{conccard}

\begin{notezone}{章末归档：不动点二级结论 + 七问笔记功能}
\textbf{二级结论}：（1）一阶递推 $a_{n+1}=pa_n+q$ 的不动点
$x^*=\dfrac q{1-p}$，$a_n=x^*+p^n(a_0-x^*)$；
（2）稳定性 $|F'(x^*)|<1\Leftrightarrow$ 局部收缩；
（3）分式迭代的不动点＝矩阵特征方向（第 5 章）；
（4）齐次坐标 $(u,v)$ 把一维分式作用提升为二维线性作用；
（5）换坐标系（平移／旋转／齐次化）都是把不动点搬到原点。

\textbf{七问与笔记功能}：(1) 数值感知 $\to$ 概念引入；(2) 三视角
翻译 $\to$ 性质定义；(3) 考纲通法 $\to$ 规范书写；(4) 高观点秒杀
$\to$ 大招笔记；(5) 公式族 $\to$ 题型拓展；(6) 逆向追问 $\to$ 充要
边界；(7) 陷阱排雷 $\to$ 错题本。
\end{notezone}


% 示范章节（仅当 main.tex 中 \demomotrue 时挂载）
\ifdemomode
  % ==========================================================
% 示范章节：同一份源码，四种成品（约 40 行，只演示「怎么写」）
% 挂载：main.tex 中把 \demomodefalse 改为 \demomotrue，
%       或运行 python3 build_all.py --demo
% 覆盖写法：通用内容 + advancedonly + basiconly + \vchoice{}{}
%          + \tk{} + solution + warmupbox/trapbox/example
% ==========================================================
\section{示范：四合一多版本写法}

\begin{warmupbox}{读法（先看这里）}
本节不教新知识，只演示\textbf{版本分流写法}：修改 \verb|main.tex| 顶部的两行开关
（\verb|\teachertrue/\teacherfalse| 与 \verb|\basictrue/\basicfalse|），
即可从\emph{同一份源码}导出四种成品 PDF，不必复制、粘贴或维护任何副本。
\end{warmupbox}

设 $f(x)=x^{2}-2x$，考察它在区间 $[0,2]$ 上的结构。

\begin{learnobjective}
\begin{itemize}
  \item 会算 $f(0),f(1),f(2)$ 三个值并描点画草图；
  \item 会用集合数对、几何位置、代数方程三种语言描述同一条抛物线；
  \item 初步看懂「平方算子」$f\mapsto f^{2}$ 的零点行为。
\end{itemize}
\end{learnobjective}

\subsection*{第 (1) 问：具象感知与画图}
计算：$f(0)=\tk{0}$，$f(1)=\tk{-1}$，$f(2)=\tk{0}$；
描出 $(0,0),(1,-1),(2,0)$ 三点，并画出草图。

\vchoice{%
  三个函数值对应三个点；用集合数对写即
  $\{(x,y)\mid y=x^{2}-2x,\ x\in[0,2]\}$，几何上是开口向上、对称轴为
  $x=1$ 的一段抛物线。}{%
  三个点：$(0,0)$、$(1,-1)$、$(2,0)$；连起来是开口向上的小弧线，
  最低点在 $x=1$ 处。}

\subsection*{第 (2) 问：常规通法演算}
配方：$f(x)=x^{2}-2x=\tk{(x-1)^{2}-1}$，故顶点为 $\tk{(1,-1)}$，
对称轴为直线 $\tk{x=1}$。

\subsection*{第 (3) 问：三视角结构翻译}
同一条曲线可以写成三种语言：集合语言
$\{(x,f(x))\mid x\in\mathbb{R}\}$；几何语言「以 $x=1$ 为对称轴的抛物线」；
代数方程 $y=(x-1)^{2}-1$。三者说的是\emph{同一结构}。

\begin{advancedonly}
\subsection*{第 (4) 问：高观点降维秒杀（仅完整版可见）}
作平移 $u=x-1$，则 $y=u^{2}-1$：绝对尺度被消去，只剩一个比例变量 $u$，
这就是齐次化与平移归一的思想雏形。进一步看平方算子：
$f^{2}(x)=(x(x-2))^{2}$ 的零点是\textbf{水平相切的偶重零点}
（局部 $\approx 4(x-1)^{2}$ 的二重根平移），不是尖点——它解释了
「平方后符号不变、零点处压平」的现象。
\end{advancedonly}

\begin{basiconly}
\subsection*{第 (4)' 问：大白话口诀与半成品填空（仅薄弱生版可见）}
口诀：「平方不改号，零点压平地；倒数撕裂缝，无穷变成零。」
照样子填空：$f(1)=1-2=\tk{-1}$；顶点横坐标 $x=\tk{1}$。
\end{basiconly}

\begin{example}[用两种语言看同一结构]
\vchoice{双射同构映射}{一一对应的加工机器}的说法在本讲统一替换为
\vchoice{二次型结构}{抛物线的对称结构}。请用「顶点式」写出
$f(x)=x^{2}-4x+3$ 的对称轴：$\tk{x=2}$。
\end{example}
\begin{solution}
配方得 $f(x)=(x-2)^{2}-1$，故对称轴为 $x=2$，顶点 $(2,-1)$。
\end{solution}

\begin{trapbox}{顶点横坐标的符号}
配方时括号内是 $x-1$ 而不是 $x+1$：$y=(x-1)^{2}-1$ 的对称轴是 $x=1$。
若误写为 $x=-1$，是把一次项系数的符号弄反了。
\end{trapbox}

\fi

\appendix
% 附录 A：高观点结论的考纲合法化转译手册
\section{附录 A：高观点结论的考纲合法化转译手册}\label{app:legal}

\begin{commonerror}{铁律：草稿纸用高观点，答题卡写初等通法}
解答题中\textbf{严禁}出现「由仿射变换知」「由极点极线知」「由欧拉公式知」「由 Abel 变换知」「由洛必达法则知」等超纲字眼。这类词一旦出现，轻则扣步骤分，重则该问判零。高观点的正确用法是：在草稿纸上用它锁定答案、取等位置与参数边界，然后在卷面上用考纲认可的初等推导把结论\textbf{重新证一遍}——代数验算本身就是满分证明。
\end{commonerror}

\begin{conccard}[三级]{三类角色各就各位}{阅卷老师只看卷面}{把预判当成证明}
\textbf{草稿纸}：随便用，越快越准越好；\textbf{卷面}：只允许考纲工具（联立、韦达定理、作差构造、错位相减、分类讨论、切线不等式）；\textbf{采分点}：看关键等式有没有被显式写出。三者不冲突，缺的只是翻译。
\end{conccard}

下面六个模板覆盖高中压轴的高频场景，每个都按「预判 $\to$ 卷面 $\to$ 采分点」三段给出。

\subsection*{模板一：圆锥曲线中的仿射压缩}

\textbf{草稿纸预判}：椭圆沿短轴方向压缩 $b/a$ 倍即得圆，面积同比例缩小，故内接三角形最大面积为 $\dfrac{3\sqrt3}{4}ab$；斜率乘积定值 $-\dfrac{b^2}{a^2}$。

\textbf{卷面写法}：设过定点的直线为 $y=kx+m$，代入 $\dfrac{x^2}{a^2}+\dfrac{y^2}{b^2}=1$，写出韦达定理与判别式条件；由弦长公式 $|AB|=\sqrt{1+k^2}\,|x_1-x_2|$ 与点到直线距离写出面积
\[
  S=\frac12|AB|\cdot d ,
\]
整理成关于 $k,m$ 的函数，用基本不等式或求导在预判的取等点处取得最值，并验证此时满足判别式大于零。

\textbf{采分点}：韦达定理两式、弦长表达式、距离表达式、「当且仅当……时取等」以及取等条件的检验。仿射结论只用来决定往哪条直线上凑。

\subsection*{模板二：极点与极线的定点定值}

\textbf{草稿纸预判}：由极点 $P(x_0,y_0)$ 直接写出极线方程 $\dfrac{x_0x}{a^2}+\dfrac{y_0y}{b^2}=1$，目标直线即刻锁定，不必再猜。

\textbf{卷面写法}：设两交点 $A(x_1,y_1),B(x_2,y_2)$，分别代入曲线方程并作差：
\[
  \frac{x_1^2-x_2^2}{a^2}+\frac{y_1^2-y_2^2}{b^2}=0
  \ \Longrightarrow\
  \frac{(x_1+x_2)(x_1-x_2)}{a^2}+\frac{(y_1+y_2)(y_1-y_2)}{b^2}=0 .
\]
再结合 $A,B$ 都在已知直线上，整理出 $\dfrac{y_1-y_2}{x_1-x_2}$ 与 $\dfrac{y_1+y_2}{x_1+x_2}$ 的关系，最后代入要证的对称式化简即可。

\textbf{采分点}：两点式作差这一步必须显式写出；之后全是恒等变形。

\subsection*{模板三：不动点与特征根求通项}

\textbf{草稿纸预判}：解 $x^*=px^*+q$ 或 $\lambda^2-p\lambda-q=0$，读出公比与两根。

\textbf{卷面写法}：由原递推式恒等变形直接写出
\[
  a_{n+1}-x^*=p(a_n-x^*) ,
\]
或两根情形下
\[
  \frac{a_{n+1}-\lambda_1}{\,a_{n+1}-\lambda_2\,}=k\,\frac{a_n-\lambda_1}{\,a_n-\lambda_2\,} .
\]
\textbf{完全不需要解释} $x^*$、$\lambda_i$ 是怎么来的——作为「已知常数」使用，再用首项确定待定系数，最后验证 $n=1$ 成立。

\textbf{采分点}：构造式本身、首项代入、$n=1$ 的检验。

\subsection*{模板四：Abel 变换与错位相减}

\textbf{草稿纸预判}：用分部求和快速算出 $S_n$ 的闭式。

\textbf{卷面写法}：规范列出
\[
  S_n=a_1+a_2+\cdots+a_n,\qquad
  qS_n=\phantom{a_1+{}}a_1q+\cdots+a_nq ,
\]
错开一位对齐后写「两式相减得」，直接默写化简结果，并把 $n=1$ 与 $n\to\infty$（若涉及极限）单独验证。

\textbf{采分点}：两行原始展开、相减动作、最终合并式、首项检验。

\subsection*{模板五：三角恒等变换与欧拉公式}

\textbf{草稿纸预判}：忘了三倍角公式时，展开 $(\cos\theta+i\sin\theta)^3$ 取实部，秒得 $\cos3\theta=4\cos^3\theta-3\cos\theta$。

\textbf{卷面写法}：从两角和公式出发现场推导：
\[
  \cos3\theta=\cos(2\theta+\theta)=\cos2\theta\cos\theta-\sin2\theta\sin\theta=\cdots,
\]
再降幂为单角函数，最后写出所需形式。

\textbf{采分点}：起始公式明确、每一步降幂正确、最终形式与题目要求一致（平方关系、次数）。

\subsection*{模板六：导数压轴的隐零点与切线放缩}

\textbf{草稿纸预判}：用切线 $e^x\ge x+1$、$\ln x\le x-1$ 预判参数临界值 $a_0$；必要时用端点效应估出 $a_0$。

\textbf{卷面写法}：按 $a\le a_0$ 与 $a>a_0$ 分类讨论。对每一类，先证切线不等式（以 $e^x-x-1$ 求最小值的方式），再代入原函数完成放缩链；隐零点处写成「设 $x_0$ 满足 $f'(x_0)=0$，则……」，用整体代换避免解出具体数值。

\textbf{采分点}：分类标准清晰、切线不等式的完整证明、隐零点的设而不求与代换、端点值的检验。

\subsection*{自查清单}

\begin{itemize}
  \item 全篇是否出现超纲名词？若有，一律替换为等价代数操作。
  \item 关键等式是否都被显式写出（而不是「显然」「易得」）？
  \item 取等条件、判别式正负、$n=1$ 首项是否逐一检验？
  \item 参数分类是否有重叠或遗漏？
\end{itemize}

\begin{conccard}[一级]{一句话记住}{预判给方向，通法给分数}{两者都要}
\key{草稿纸负责快，卷面负责稳}。高观点决定你往哪条路上走，考纲通法决定你能拿多少分——只有前者，会丢步骤分；只有后者，会算不完。
\end{conccard}

% ==========================================================
% 附录 B：大学理工科荣誉课程（Honors Bridge）先修接口对照手册
% 说明：本文件为纯正文模块化片段（无 \documentclass），
%       由 main.tex 在 \appendix 之后 \input 挂载。
%       corebox / methodbox 现由 versions.tex 全局提供（四合一引擎），
%       配色与现有 main/third 主题一致，不再局部定义。
% ==========================================================

\section{附录 B：大学理工科荣誉课程先修接口对照手册}\label{app:honors}

\begin{corebox}{致准大学荣誉学院（Honors Program）学子}
普通高中生升入数学、物理、计算机、人工智能与电子信息等理工科专业时，普遍会经历一场剧烈的\textbf{“大一数学休克”}——例如看不懂《数学分析》严密的 $\varepsilon$-$\delta$ 极限语言，无法理解《线性代数》中抽象的基底变换与特征空间。

然而，如果你已经扎实地走完了本套讲义的第 0–9 章，\textbf{你实际上已经提前掌握了大学前两年硬核数学课的几何直觉与结构骨架}。本附录旨在为你完成最后一步\textbf{“学术暗号对接”}，将你熟悉的中学高观点，精准映射到大学专业课的正式定义与定理名称上。
\end{corebox}

\subsection*{一、 五大大学核心基础课“概念暗号”对接总表}

\begin{fastcard}{《数学分析》（Real Analysis / Honors Calculus）}
\textbf{你在本套讲义中学过的直觉与工具}：第 1 章有理数 $\mathbb{Q}$ 漏掉 $\sqrt{2}$ 孔洞，无理数填满实心数轴 $\mathbb{R}$；第 6 章局部切线近似（以直代曲）与切线支撑放缩。
\par\textbf{大学专业课正式学名与定理}：\textbf{实数完备性公理}（戴德金分割 / 确界原理）；\textbf{介值定理}与\textbf{微分中值定理}；\textbf{凸函数与 Jensen 不等式}。
\par\textbf{消除的“大一休克”痛点}：瞬间理解极限存在性证明，不再死背 $\varepsilon$-$\delta$ 符号，一眼看穿微分局部线性化本质。
\end{fastcard}

\begin{fastcard}{《线性代数》（Linear Algebra / 高等代数）}
\textbf{你在本套讲义中学过的直觉与工具}：第 5 章换基底对齐特征方向、旋转 $45^\circ$ 消除交叉项 $Bxy$；第 1 章函数变换中的主动（动图形）与被动（动坐标系）对偶；第 5 章转移矩阵特征向量锁定分式不动点。
\par\textbf{大学专业课正式学名与定理}：\textbf{线性空间与基底变换}（过渡矩阵 $P^{-1}AP$）；\textbf{实对称矩阵正交对角化}（谱定理）；\textbf{射影空间（Projective Space）}与齐次坐标。
\par\textbf{消除的“大一休克”痛点}：摆脱“矩阵只是数字表格”的机械计算误区，直接建立高维几何变换与特征空间直觉。
\end{fastcard}

\begin{fastcard}{《复变函数与信号系统》（Complex Analysis \& Signals）}
\textbf{你在本套讲义中学过的直觉与工具}：第 2 章虚数 $i$ 作为逆时针 $90^\circ$ 旋转算子；切线垂直速度演化出欧拉公式 $e^{i\theta}=\cos\theta+i\sin\theta$。
\par\textbf{大学专业课正式学名与定理}：\textbf{复数域的几何表示与保角映射}；\textbf{傅里叶级数与傅里叶变换}（正交旋转基展开）；工科电路\textbf{相量法（Phasor）}。
\par\textbf{消除的“大一休克”痛点}：在大学物理简谐振动、交流电路与信号处理中，直接把正弦量当成复平面旋转向量处理。
\end{fastcard}

\begin{fastcard}{《常微分方程与数值分析》（ODEs \& Numerical Methods）}
\textbf{你在本套讲义中学过的直觉与工具}：第 5 章差分 $\Delta a_n$ 与求和互逆、Abel 分部求和；蛛网图 $|f'(x^*)|<1$ 螺旋收敛于吸引不动点；生成函数 $G(x)=\sum a_n x^n$。
\par\textbf{大学专业课正式学名与定理}：\textbf{差分方程与常微分方程特征根法}；\textbf{巴拿赫压缩映射原理}（牛顿迭代法收敛性）；\textbf{Z 变换与拉普拉斯变换}。
\par\textbf{消除的“大一休克”痛点}：打通离散递归与连续微分方程，理解计算机迭代算法为何收敛以及混沌分岔的临界条件。
\end{fastcard}

\begin{fastcard}{《公理化概率论与机器学习》（Probability \& Machine Learning）}
\textbf{你在本套讲义中学过的直觉与工具}：第 7 章样本空间 $\Omega$ 归一化测度与条件子集缩放；二维散点云的竖直欧氏距离平方和最小投影线。
\par\textbf{大学专业课正式学名与定理}：\textbf{Kolmogorov 概率测度空间} $(\Omega, \mathcal{F}, P)$；\textbf{线性回归的最小二乘正交投影}（正规方程 $A^TA\hat{x}=A^Tb$）。
\par\textbf{消除的“大一休克”痛点}：跳过繁琐的古典计数，直接站在高维空间几何投影的高度理解统计学习与数据拟合。
\end{fastcard}

\subsection*{二、 三大核心概念的“中学直觉 $\to$ 大学严格化”跃迁示范}

为了让你在翻开大学教材的第一天就能写出严密的现代数学证明，以下展示三个最核心概念如何从本讲义的“几何直觉”平滑过渡到大学“公理化语言”：

\begin{methodbox}{跃迁 1：从“实心金条无漏洞”到《数学分析》的“确界原理”}
\begin{itemize}
    \item \textbf{中学高观点直觉}（第 1 章）：有理数轴 $\mathbb{Q}$ 像布满针眼的筛子，曲线 $y=x^2-2$ 会从 $\sqrt{2}$ 的孔洞漏过去；实数轴 $\mathbb{R}$ 填满了所有孔洞，因此有上界的递增序列必然会撞上一个确定的天花板。
    \item \textbf{大学《数学分析》严格表述（确界公理）}：设 $S \subset \mathbb{R}$ 为非空有上界实数集，则 $S$ 在 $\mathbb{R}$ 中必有\textbf{上确界（最小上界 $\sup S$）}。由此可直接推出“单调有界数列必有极限”与“闭区间套定理”。
\end{itemize}
\end{methodbox}

\begin{methodbox}{跃迁 2：从“旋转 $45^\circ$ 消交叉项 $xy$”到《线性代数》的“正交对角化”}
\begin{itemize}
    \item \textbf{中学高观点直觉}（第 5 章）：双曲线 $xy = 1$ 的主轴倾斜了 $45^\circ$，只要把坐标轴旋转 $45^\circ$ 对齐它的对称轴，交叉项就会消失，化为 $\frac{x'^2}{2} - \frac{y'^2}{2} = 1$。
    \item \textbf{大学《线性代数》严格表述（二次型主轴定理）}：将 $xy = 1$ 写成对称矩阵二次型：
    $$\begin{pmatrix} x & y \end{pmatrix} \begin{pmatrix} 0 & \frac{1}{2} \\ \frac{1}{2} & 0 \end{pmatrix} \begin{pmatrix} x \\ y \end{pmatrix} = 1$$
    求实对称矩阵 $A = \begin{pmatrix} 0 & 1/2 \\ 1/2 & 0 \end{pmatrix}$ 的特征值 $\lambda_1 = \frac{1}{2}, \lambda_2 = -\frac{1}{2}$ 及其单位正交特征向量（构成旋转 $45^\circ$ 的正交矩阵 $Q$），令 $\mathbf{x} = Q\mathbf{x}'$，立刻得到标准型 $\frac{1}{2}x'^2 - \frac{1}{2}y'^2 = 1$！
\end{itemize}
\end{methodbox}

\begin{methodbox}{跃迁 3：从“蛛网图台阶变小”到《泛函/数值分析》的“压缩映射原理”}
\begin{itemize}
    \item \textbf{中学高观点直觉}（第 5 章）：在蛛网图中，只要不动点处曲线斜率满足 $|f'(x^*)| < 1$，每次投影都会把误差缩短，序列最终螺旋收敛于 $x^*$。
    \item \textbf{大学严格表述（Banach Fixed-Point Theorem）}：设 $(X, d)$ 为完备度量空间，映射 $T: X \to X$ 满足存在常数 $0 \le L < 1$，使得对任意 $x, y \in X$ 均有 $d(T(x), T(y)) \le L \cdot d(x, y)$，则 $T$ 在 $X$ 中存在\textbf{唯一不动点 $x^*$}，且对任意初值 $x_0 \in X$，迭代序列 $x_{n+1} = T(x_n)$ 均以几何级数速度收敛于 $x^*$。
\end{itemize}
\end{methodbox}

\subsection*{三、 给荣誉学院新生的阅读进阶路线图}

当你顺利进入大学理工科荣誉班级后，推荐按以下书单无缝衔接本讲义的知识结构：
\begin{enumerate}
    \item \textbf{数学分析方向}：陶哲轩《陶哲轩实分析》（从集合与自然数公理重新严格建构整个实数与微积分体系，与本讲义第 0、1 章风格一脉相承）；
    \item \textbf{线性代数方向}：Sheldon Axler《Linear Algebra Done Right》（摒弃繁琐行列式，直接从向量空间与线性算子的不变子空间/特征向量切入，与本讲义换基底思想高度契合）；
    \item \textbf{复分析与几何直觉方向}：Tristan Needham《复分析：可视化方法》（用纯几何旋转与向量场讲透复变函数，完美延续本讲义第 2 章与 Manim 动画风格）。
\end{enumerate}


\end{document}
。
%
%  四种组合（改两行即可，无需复制任何文件）：
%    \teachertrue  \basicfalse   完整教师版（默认母版）
%    \teacherfalse \basicfalse   完整学生版（隐藏解析与填空答案）
%    \teachertrue  \basictrue    薄弱生教师版（显化脚手架 + 完整答案）
%    \teacherfalse \basictrue    薄弱生学生版（保留 \tk 留白横线）
%
%  本文件只依赖已加载的宏包（xcolor / tcolorbox / environ），
%  不修改 common/，故对 ../hv-bank 等其他分册零影响。
% ============================================================

% ---------- 布尔开关 ----------
\newif\ifteacher          % 教师解析版 / 学生留白版
\newif\ifbasic            % 基础通关版 / 完整高观点版
\newif\ifdemomode         % 是否挂载示范章节 chapters/demo_multiversion

% 安全默认（main.tex 的两行开关覆盖此处）
\teachertrue
\basicfalse
\demomodefalse

% ---------- 零残留隐藏所需 ----------
\usepackage{environ}      % \NewEnviron（隐藏时不留空白）

% ---------- 版本标签（封面 / 页眉自动显示） ----------
\newcommand{\versionlabel}{%
  \ifteacher 教师解析版\else 学生留白版\fi
  \,\textperiodcentered\,
  \ifbasic 基础通关版\else 完整高观点版\fi}

% ============================================================
%  一、版本内容分流环境
% ============================================================

% 仅完整版渲染；薄弱生版彻底忽略（不占位）
% 用途：莫比乌斯矩阵、射影同构、压轴第 3–4 阶梯难题等
\NewEnviron{advancedonly}{%
  \ifbasic\else
    \par\medskip
    \BODY
    \par\medskip
  \fi
}

% 仅薄弱生版渲染；完整版彻底忽略
% 用途：课前 15 分钟计算过关、大白话口诀、半成品步步填空
\NewEnviron{basiconly}{%
  \ifbasic
    \par\medskip
    \BODY
    \par\medskip
  \fi
}

% ============================================================
%  二、双模填空宏 \tk{...} 与内联版本选择 \vchoice{}{}
% ============================================================

% 教师版：红色下划线答案；学生版：等宽下划线留白（版面稳定）
% 内容一律写成裸数学式（不带 $）。\underline 继承当前模式，
% 而 \sbox 的盒子内容始终是文本模式、必须自带数学定界符。
\newcommand{\tk}[1]{%
  \ifteacher
    \ifmmode
      {\color{second}\underline{#1}}%
    \else
      {\color{second}\underline{$\displaystyle#1$}}%
    \fi
  \else
    \sbox0{$\displaystyle#1$}%
    \underline{\hspace{\wd0}}%
  \fi
}

% 同一句子里的措辞切换：完整版文本 / 薄弱生大白话文本
\newcommand{\vchoice}[2]{%
  \ifbasic #2\else #1\fi}

% ============================================================
%  三、解答环境重定义（教师版带边框【剖析与标准解答】）
% ============================================================
\RenewDocumentEnvironment{solution}{+b}{%
  \ifteacher
    \par\medskip
    \begin{tcolorbox}[
      breakable, enhanced, sharp corners,
      colback=main!3, colframe=main!60!black, boxrule=0.5pt,
      title={\bfseries 剖析与标准解答}, coltitle=black, fonttitle=\small,
      top=5pt,
    ]
    #1
    \end{tcolorbox}%
  \else
    % 学生版：不留内容，只保留少量书写留白
    \par\noteblank[6ex]%
  \fi
}{}

% ============================================================
%  四、语义化盒子（corebox / methodbox / trapbox / warmupbox）
% ============================================================

% 核心概念盒（绿）
\newtcolorbox{corebox}[1]{
  breakable, enhanced, sharp corners,
  colback=main!4, colframe=main!60!black, boxrule=0.6pt,
  title={\bfseries #1}, coltitle=black, fonttitle=\small,
  attach boxed title to top left={xshift=2mm, yshift=-2mm},
  boxed title style={sharp corners, colback=white,
    colframe=main!60!black, boxrule=0.4pt},
  top=7pt,
}

% 方法归纳盒（青）
\newtcolorbox{methodbox}[1]{
  breakable, enhanced, sharp corners,
  colback=third!5, colframe=third!60!black, boxrule=0.5pt,
  title={\bfseries #1}, coltitle=black, fonttitle=\small,
  attach boxed title to top left={xshift=2mm, yshift=-2mm},
  boxed title style={sharp corners, colback=white,
    colframe=third!60!black, boxrule=0.4pt},
  top=7pt,
}

% 易错陷阱盒（红）
\newtcolorbox{trapbox}[1]{
  breakable, enhanced, sharp corners,
  colback=second!4, colframe=second!65!black, boxrule=0.5pt,
  title={\bfseries 陷阱：#1}, coltitle=black, fonttitle=\small,
  attach boxed title to top left={xshift=2mm, yshift=-2mm},
  boxed title style={sharp corners, colback=white,
    colframe=second!65!black, boxrule=0.4pt},
  top=7pt,
}

% 课前热身盒（蓝）
\newtcolorbox{warmupbox}[1]{
  breakable, enhanced, sharp corners,
  colback=structurecolor!4, colframe=structurecolor!70!black, boxrule=0.5pt,
  title={\bfseries #1}, coltitle=black, fonttitle=\small,
  top=6pt,
}

% ============================================================
%  五、按章（节）自动编号的 example / exercise 环境
%  注：elegantnote 自带 example 计数器（\c@example），故另立
%      hvexample / hvexercise 计数，避免与模板冲突；编号仍按节累计。
% ============================================================
\newcounter{hvexample}[section]
\renewcommand{\thehvexample}{\thesection.\arabic{hvexample}}

\RenewDocumentEnvironment{example}{O{}}{%
  \refstepcounter{hvexample}%
  \par\medskip\noindent
  {\bfseries\color{main}例题\thehvexample\IfValueT{#1}{：#1}\quad}%
}{\par\medskip}

\newcounter{hvexercise}[section]
\renewcommand{\thehvexercise}{\thesection.\arabic{hvexercise}}

\NewDocumentEnvironment{exercise}{O{}}{%
  \refstepcounter{hvexercise}%
  \par\medskip\noindent
  {\bfseries\color{third}练习\thehvexercise\IfValueT{#1}{：#1}\quad}%
}{\par\medskip}

% ============================================================
%  六、七问连环母题（Scaffolded Master Problem）载体环境
% ------------------------------------------------------------
%  定位：一道母题挂载 (1)–(7) 个递进小问，学生做题留白处写下的
%        推理本身就是这一章的学习笔记。写作规范见 ARCHITECTURE.md。
% ============================================================

\newcounter{hvmasterq}[section]
\renewcommand{\thehvmasterq}{\thesection.\arabic{hvmasterq}}

% 母题盒：标题自动编号；盒内用 enumerate[label=(\arabic*)] 写七问
% 注意：计数器由调用方在 \begin{masterq} 前手动 \refstepcounter{hvmasterq}，
%       不要把 \refstepcounter 写进本定义的选项里（tcolorbox 会反复展开选项，
%       导致编译卡在分页计算）。
\newtcolorbox{masterq}[1]{
  breakable, enhanced, sharp corners,
  colback=structurecolor!3, colframe=structurecolor!70!black, boxrule=0.6pt,
  title={\bfseries 七问连环母题\thehvmasterq：#1},
  coltitle=black, fonttitle=\small,
  attach boxed title to top left={xshift=2mm, yshift=-2mm},
  boxed title style={sharp corners, colback=white,
    colframe=structurecolor!70!black, boxrule=0.4pt},
  top=7pt,
}

% 每问末尾的笔记归档标签：\logotag{存档类别}（如「核心公理」「大招公式」）
\newcommand{\logotag}[1]{%
  \hfill\textcolor{third}{\small\ttfamily【笔记归档·#1】}}
，
%        再设置两行开关，然后 % ============================================================
%  main_body.tex —— 正文主体（由 main.tex 在设置好版本开关后 \input）
%  拆分原因：四合一多版本编译时，包装文件只需复写「两行开关」，
%  再 % ============================================================
%  main_body.tex —— 正文主体（由 main.tex 在设置好版本开关后 \input）
%  拆分原因：四合一多版本编译时，包装文件只需复写「两行开关」，
%  再 % ============================================================
%  main_body.tex —— 正文主体（由 main.tex 在设置好版本开关后 \input）
%  拆分原因：四合一多版本编译时，包装文件只需复写「两行开关」，
%  再 \input{main_body}，避免导言区内容四处重复、彼此漂移。
% ============================================================

\title{数学导学案·高观点系列}
\author{Shawn Shang}
\date{2026 年 11 月}

\graphicspath{{../figures/}{../map/figures/}}

\begin{document}

\begin{titlepage}
  \thispagestyle{empty}
  \centering
  \makebox[\textwidth]{%
    \includegraphics[width=\paperwidth]{cover_cyan.png}}\par
  \vfill
  {\Huge\bfseries\color{ecolor!70!black} 数学导学案·高观点系列\par}
  \vskip 1.2em
  {\large 结构主义视角下的高中数学\par}
  \vskip 0.8em
  {\large 第 0 章 — 第 9 章\par}
  \vskip 0.8em
  {\large\versionlabel\par}
  \vskip 2.5em
  {\large Shawn Shang\par}
  \vfill
  {\normalsize 2026 年 11 月\par}
\end{titlepage}
\tableofcontents

% 让第一个 \section 编号为 0（第 0 章是总纲）
\setcounter{section}{-1}

\input{chapters/00}
\input{chapters/01}
\input{chapters/02}
\input{chapters/03}
\input{chapters/04}
\input{chapters/05}
\input{chapters/06}
\input{chapters/07}
\input{chapters/08_homogenization}
\input{chapters/09_fixed_point}

% 示范章节（仅当 main.tex 中 \demomotrue 时挂载）
\ifdemomode
  \input{chapters/demo_multiversion}
\fi

\appendix
\input{chapters/appA_legalization}
\input{chapters/appB_honors_bridge}

\end{document}
，避免导言区内容四处重复、彼此漂移。
% ============================================================

\title{数学导学案·高观点系列}
\author{Shawn Shang}
\date{2026 年 11 月}

\graphicspath{{../figures/}{../map/figures/}}

\begin{document}

\begin{titlepage}
  \thispagestyle{empty}
  \centering
  \makebox[\textwidth]{%
    \includegraphics[width=\paperwidth]{cover_cyan.png}}\par
  \vfill
  {\Huge\bfseries\color{ecolor!70!black} 数学导学案·高观点系列\par}
  \vskip 1.2em
  {\large 结构主义视角下的高中数学\par}
  \vskip 0.8em
  {\large 第 0 章 — 第 9 章\par}
  \vskip 0.8em
  {\large\versionlabel\par}
  \vskip 2.5em
  {\large Shawn Shang\par}
  \vfill
  {\normalsize 2026 年 11 月\par}
\end{titlepage}
\tableofcontents

% 让第一个 \section 编号为 0（第 0 章是总纲）
\setcounter{section}{-1}

% 高观点导学案：集合、逻辑演算与证明武器库（第 0 章·总纲）
% 主线：把“集合与逻辑”从概念背诵升级为贯穿高中解题的“题干同解变形与严格证明通用武器库”
% 四模块：集合语言与同解变形公理 → 命题演算三大变形引擎 → 两大间接证明支柱 → 武器库索引表与跨模块实战
% 全程用“集合包含 ↔ 逻辑等价 ↔ 量词否命题”三种语言互译，并用韦恩图与多米诺图把抽象逻辑可视化。

\section{集合、逻辑演算与证明武器库（高观点总纲）}

本章不引入新公式，而是把教材第 1 章的集合、常用逻辑用语重新锻造为一套“解题武器库”：
面对题干，先问它是哪一种\key{对象集合}（解集、值域、参数范围），再问它需要哪一种\key{变形引擎}
（同解变形、德摩根律、逆否等价、补集转化、反证法、数学归纳法）。四个模块环环相扣：
\key{集合语言与同解变形公理}（$A=B\iff A\subseteq B\land B\subseteq A$）$\to$
\key{命题演算三大变形引擎}（德摩根律、逆否等价、量词否定）$\to$
\key{两大间接证明支柱}（反证法、数学归纳法）$\to$
\key{武器库索引表与跨模块实战}。
核心思想：\key{题干中的每一个关键词，都对应一次集合运算或一次逻辑等价变形}；
解题失误往往不是计算错误，而是\key{没审清变形方向}（增根或漏根）或\key{没找对证明支柱}。

\begin{learnobjective}
\begin{itemize}
  \item 掌握集合相等判定公理 $A=B\iff(A\subseteq B)\land(B\subseteq A)$，能用它审查方程、不等式变形的同解性；
  \item 熟记同解变形防增漏根清单：平方（范围扩大易增根）、同除含参因式（范围缩小易漏根）、对数真数（定义域子集）；
  \item 掌握德摩根律 $\neg(p\land q)\iff\neg p\lor\neg q$ 与逆否等价 $(p\Rightarrow q)\iff(\neg q\Rightarrow\neg p)$，能用补集转化处理“至多、至少、不都、无解”；
  \item 掌握量词否定 $\neg\forall=\exists\neg$、$\neg\exists=\forall\neg$，会把“恒成立／能成立”翻译为值域的子集包含／交集非空；
  \item 理解反证法即证 $A\cap\complement_U B=\varnothing$，识别“唯一性、不存在、无理数、无限性”四类触发信号词；
  \item 理解数学归纳法的皮亚诺公理与良序原理本质，会用最小反例原理分析“归纳基础＋归纳递推”两步的缺一不可。
\end{itemize}
\end{learnobjective}

\begin{preview}
\subsection*{知识回顾}
\begin{itemize}
  \item $A\subseteq B$ 的定义是：对任意 $a\in A$，都有 \blank；
  \item 若 $A=\{x\mid x^2=1\}$，$B=\{-1,1\}$，则 $A$ 与 $B$ 的关系是 \blank（相等／互不包含）；
  \item 命题“若 $p$ 则 $q$”的逆否命题是 \blank；
  \item 德摩根律：$\neg(p\land q)\iff$ \blank。
\end{itemize}

\subsection*{教材速读}
阅读教材第 1 章“集合与常用逻辑用语”及本章导学案，完成下列问题：
\begin{enumerate}
  \item 补集 $\complement_U A$ 的定义是 $\{x\in U\mid \blank\}$；
  \item “$p$ 是 $q$ 的充分条件”用集合语言说是 \blank；
  \item “方程 $f(x)=0$ 无解”用集合语言说是 \blank；
  \item 数学归纳法证明的两步是 \blank 与 \blank。
\end{enumerate}
\end{preview}

\begin{explore}

\subsection{集合语言与同解变形公理}

\key{探究一：集合相等——互相包含才是相等}
集合相等不是“看起来一样”，而是\key{互相包含}：
\[
  A=B\iff(A\subseteq B)\land(B\subseteq A).
\]
证明 $A=B$ 必须证两个方向；只证 $A\subseteq B$ 只能得到包含，不能得到相等。
这条公理还是审查方程（不等式）变形的“合法性标准”：
把原方程的解集记为 $A$，变形后方程的解集记为 $B$，则
\begin{itemize}
  \item $A=B$：变形\key{同解}，解集不变，安全；
  \item $A\subsetneqq B$：变形后解集\key{变大}，多出的解是\key{增根}；
  \item $B\subsetneqq A$：变形后解集\key{变小}，丢掉的解是\key{漏根}。
\end{itemize}
因此解方程、解不等式的每一步都要问一句：\key{这一步把解集往哪个方向变了？}

\begin{center}
\footnotesize
\setlength{\tabcolsep}{3pt}
\renewcommand{\arraystretch}{1.4}
\begin{tabular}{p{2.6cm}p{2.8cm}p{3.2cm}p{3.2cm}}
\toprule
变形操作 & 解集关系 & 方向与风险 & 应对 \\
\midrule
两边加/减同一整式 & $A=B$ & 同解，安全 & 直接使用 \\
两边乘/除以非零常数 & $A=B$ & 同解，安全 & 先确认非零 \\
两边平方（或去绝对值） & $A\subsetneqq B$ & 范围扩大，易增根 & 必须验根 \\
两边同除含参/未知因式 & $B\subsetneqq A$ & 范围缩小，易漏根 & 移项分解，讨论零 \\
取对数/解对数方程 & $A\subsetneqq B$ & 真数 $>0$ 限制，易增根 & 先限定定义域 \\
\bottomrule
\end{tabular}
\end{center}

\key{风险一：平方变形范围扩大，易增根。}
由 $f(x)=g(x)$ 平方得 $f^2(x)=g^2(x)$，而
\[
  f^2=g^2\iff(f-g)(f+g)=0\iff f=g\ \text{或}\ f=-g,
\]
凭空多出 $f=-g$ 的解。必须验根。

\begin{example}
解方程 $\sqrt{x+1}=x-1$，并指出直接平方会多出哪个根。
\end{example}
\begin{solution}
两边平方（须先注意右边非负：$x-1\geq0$，即 $x\geq1$）：
\[
  x+1=(x-1)^2=x^2-2x+1\implies x^2-3x=0\implies x=0\ \text{或}\ x=3.
\]
检验：$x=0$ 时左边 $=1$，右边 $=-1$，不相等，是平方后产生的\key{增根}，舍去；
$x=3$ 时左边 $=2$，右边 $=2$，成立。故解集为 $\{3\}$。
\end{solution}

\key{风险二：同除含参因式范围缩小，易漏根。}
方程两边同除以含未知数（或参数）的因式，若该因式可能为 $0$，就会丢掉使因式为 $0$ 的解。
正确做法是移项因式分解，并单独讨论因式为 $0$ 的情形。

\begin{example}
解关于 $x$ 的方程 $x^2=ax$（$a$ 为参数），并说明“两边同除以 $x$”会漏掉什么。
\end{example}
\begin{solution}
移项分解：$x(x-a)=0$，故 $x=0$ 或 $x=a$。
当 $a\neq0$ 时解集为 $\{0,a\}$；当 $a=0$ 时方程为 $x^2=0$，解集为 $\{0\}$（两个相等实根）。
若直接两边同除以 $x$，得到 $x=a$，会漏掉 $x=0$（$a\neq0$ 时），
因为 $x$ 可能为 $0$，同除使解集由 $\{0,a\}$ 缩小为 $\{a\}$，属于“范围缩小易漏根”。
\end{solution}

\key{风险三：对数真数限制定义域子集，易增根。}
\[
  \log_a f(x)=\log_a g(x)\iff f(x)=g(x)>0,
\]
其中“$>0$”来自真数的定义域限制。定义域
$D=\{x\mid f(x)>0,\ g(x)>0\}$ 是一个\key{子集约束}，解必须落在 $D$ 内；忽略它就增根。
（同理，$a^{\log_a x}=x$ 要求 $x>0$，而 $\log_a a^x=x$ 对一切实数 $x$ 成立。）

\begin{example}
解方程 $\log_2(x^2-4)=\log_2(3x)$。
\end{example}
\begin{solution}
由对数定义，真数必须为正：$x^2-4>0$ 且 $3x>0$，即 $x>2$。
在此定义域内 $\log_2$ 是单射，故 $x^2-4=3x$，即 $x^2-3x-4=0$，解得 $x=4$ 或 $x=-1$。
$x=-1$ 不在定义域 $x>2$ 内，是增根，舍去。故解集为 $\{4\}$。
\end{solution}

\begin{conccard}[一级]{集合相等判定公理：$A=B\iff(A\subseteq B)\land(B\subseteq A)$}{证明集合相等、审查方程（不等式）同解变形}{只证 $A\subseteq B$ 一个方向就宣布相等；只算“怎么变形”，不算解集往哪变}
集合相等的本质是“互相包含”。证 $A=B$ 必须证两个方向：$A\subseteq B$ 且 $B\subseteq A$；
只证一个方向只能得到包含。把方程（不等式）的解集记为 $A$、变形后解集记为 $B$，
只有 $A=B$ 时变形才同解；$A\subsetneqq B$ 则增根，$B\subsetneqq A$ 则漏根。
这条公理是全部同解变形的“合法性审查标准”。
\end{conccard}

\begin{commonerror}{同解变形不审方向}
常见错误：平方后不验根、同除含参因式不讨论零、解对数方程不写定义域。
它们都源于同一个疏漏——只算“怎么变形”，不算“解集往哪个方向变”。
例如由 $\sqrt{x+1}=x-1$ 平方得 $x^2-3x=0$，解集由 $\{3\}$ 扩大为 $\{0,3\}$，多出的 $0$ 就是增根；
由 $x^2=ax$ 两边除以 $x$ 得 $x=a$，解集由 $\{0,a\}$ 缩小为 $\{a\}$，漏掉了 $0$。
\end{commonerror}

\subsection{命题演算三大变形引擎}

\key{探究二：德摩根律——“且”的否定是“或”}
\[
  \neg(p\land q)\iff\neg p\lor\neg q,\qquad
  \neg(p\lor q)\iff\neg p\land\neg q.
\]
对应的集合版本（补集的分配律）：
\[
  \complement_U(A\cap B)=\complement_U A\cup\complement_U B,\qquad
  \complement_U(A\cup B)=\complement_U A\cap\complement_U B.
\]
记忆：否定“且”要变“或”，否定“或”要变“且”，并且每个命题都要否定。

\key{引擎二：逆否等价——正难则反}
\[
  (p\Rightarrow q)\iff(\neg q\Rightarrow\neg p).
\]
原命题与逆否命题同真同假，这是“正难则反”的合法性来源。它还有集合版本：

\begin{derivation}
要证 $A\subseteq B\iff\complement_U B\subseteq\complement_U A$。
\begin{align*}
  A\subseteq B &\iff \forall x\,(x\in A\Rightarrow x\in B)\\
  &\iff \forall x\,(x\notin B\Rightarrow x\notin A) &&(\text{逆否等价})\\
  &\iff \forall x\,(x\in\complement_U B\Rightarrow x\in\complement_U A)\\
  &\iff \complement_U B\subseteq\complement_U A.
\end{align*}
\end{derivation}

可见：要证“小集合含于大集合”（$A\subseteq B$），可改证“大集合的补集含于小集合的补集”
（$\complement_U B\subseteq\complement_U A$），往往更容易。

\key{引擎三：量词否定与补集转化流程}
\[
  \neg(\forall x,P(x))\iff\exists x,\neg P(x),\qquad
  \neg(\exists x,P(x))\iff\forall x,\neg P(x).
\]
口诀：\key{量词互换，结论否定}。与德摩根律联用：
\[
  \neg\bigl(\forall x,\,p(x)\land q(x)\bigr)\iff\exists x,\,\neg p(x)\lor\neg q(x).
\]
题干关键词的补集转化流程：
\begin{itemize}
  \item “至少有一个”（$\exists$）的否定是“一个都没有”（$\forall,\neg$）；
  \item “至多有一个”= 否定“至少有两个”；
  \item “不都”= $\neg\forall$ = “至少有一个不”；
  \item “无解”= 解集 $\varnothing$ = 补集为全集 $U$。
\end{itemize}
流程：把“至多／至少／不都／无解”改写成“$\exists/\forall$ + 否定” $\to$ 用德摩根与量词否定化简
$\to$ 求补集（\key{补集法}）$\to$ 取补并检查端点。

\key{恒成立／能成立：值域的子集包含与交集非空。}
设 $f:D\to\mathbb{R}$ 的值域为 $R_f=\{f(x)\mid x\in D\}$，目标集合为 $M$，则
\[
  \forall x\in D,\ f(x)\in M \iff R_f\subseteq M,\qquad
  \exists x\in D,\ f(x)\in M \iff R_f\cap M\neq\varnothing.
\]
特例 $M=(a,+\infty)$：恒成立 $\iff\min R_f>a$（看最小值）；能成立 $\iff\max R_f>a$（看最大值）。
例如 $f(x)=x^2-2x$ 在 $[0,3]$ 上的值域为 $[-1,3]$，则 $f(x)>a$ 恒成立 $\iff a<-1$，
能成立 $\iff a<3$。

\begin{example}
写出命题“$\exists x\in\mathbb{R},\ x^2-2x+a\leq0$”的否定；
并说明“方程 $x^2-2x+a=0$ 无实根”用集合语言如何表述。
\end{example}
\begin{solution}
否定为“$\forall x\in\mathbb{R},\ x^2-2x+a>0$”（量词互换，结论否定）。
“无实根”即解集为空集：$\{x\mid x^2-2x+a=0\}=\varnothing$，
等价于“对任意 $x\in\mathbb{R}$，$x^2-2x+a\neq0$”，也等价于 $\Delta=4-4a<0$，即 $a>1$。
\end{solution}

\begin{example}
已知“存在 $x\in[1,2]$，使 $x^2-ax+1<0$”为真，求 $a$ 的取值范围（用补集法）。
\end{example}
\begin{solution}
因为 $x\in[1,2]$ 为正，原命题等价于
\[
  \exists x\in[1,2],\ a>x+\frac1x
  \iff a>\min_{x\in[1,2]}\Bigl(x+\frac1x\Bigr).
\]
由均值不等式 $x+\frac1x\geq2$，当 $x=1$ 时取等号，故最小值为 $2$，于是 $a>2$。

补集法验证：其否定为“$\forall x\in[1,2],\ x^2-ax+1\geq0$”，即 $a\leq x+\frac1x$ 恒成立，
得 $a\leq2$；取补集得 $a>2$，与上面一致。
\end{solution}

\begin{problem}
设 $f(x)=x^2-2x$ 在 $[0,3]$ 上的值域为 $R_f$。分别求“$f(x)>a$ 恒成立”与“$f(x)>a$ 能成立”时 $a$ 的范围。
\end{problem}
\begin{solution}
$R_f=[-1,3]$。恒成立 $\iff R_f\subseteq(a,+\infty)\iff a<-1$；
能成立 $\iff R_f\cap(a,+\infty)\neq\varnothing\iff a<3$。
\end{solution}

\begin{conccard}[二级]{逆否等价与补集反向包含：$p\Rightarrow q\iff\neg q\Rightarrow\neg p$；$A\subseteq B\iff\complement_U B\subseteq\complement_U A$}{证明充分条件、正难则反的题目}{把“逆命题”与“逆否命题”混同；把“否命题”当“命题的否定”}
原命题与逆否命题同真同假，这是“正难则反”的合法性来源。集合语言下，
要证 $A\subseteq B$ 可改证 $\complement_U B\subseteq\complement_U A$——
把“小推大”换成“大补推小补”，往往更容易。
\end{conccard}

\begin{conccard}[二级]{恒成立／能成立：值域的子集包含与交集非空}{含参恒成立、能成立求参数范围}{把“恒成立”与“能成立”方向用反；忽略值域端点是否取到}
设 $f$ 在 $D$ 上的值域为 $R_f$。$\forall x\in D,\ f(x)\in M\iff R_f\subseteq M$（看最小值）；
$\exists x\in D,\ f(x)\in M\iff R_f\cap M\neq\varnothing$（看最大值）。
特例 $M=(a,+\infty)$：恒成立 $\iff\min R_f>a$，能成立 $\iff\max R_f>a$。
\end{conccard}

\begin{commonerror}{把“否命题”当“命题的否定”}
命题“若 $p$ 则 $q$”的\key{否定}是“$p$ 且非 $q$”（$p\land\neg q$），它的\key{否命题}是“若非 $p$ 则非 $q$”（$\neg p\Rightarrow\neg q$）。
两者不同：否定只否定结论，否命题条件与结论都否定。
例如“若 $x>2$ 则 $x^2>4$”的否定是“存在 $x>2$ 使 $x^2\leq4$”；
其否命题“若 $x\leq2$ 则 $x^2\leq4$”是假命题（取 $x=-3$）。
\end{commonerror}

\begin{method}
补集转化法（正难则反）：遇到“至多、至少、不都、无解、并非”等关键词，先写出它的否定
（量词互换、联结词用德摩根），把“存在”转“任意”或把“任意”转“存在”，
求出否定命题成立的参数范围，再取补集。步骤：
\begin{enumerate}
  \item 定位量词（$\forall/\exists$）与联结词（$\land/\lor$）；
  \item 否定量词、否定结论，用德摩根律化简；
  \item 求否定命题对应的集合（补集法求参数）；
  \item 取补集，并检查端点。
\end{enumerate}
\end{method}

\subsection{两大间接证明支柱}

\key{探究三：反证法——证明反例集为空}
要证命题 $p\Rightarrow q$。设 $U$ 为论域，$A=\{x\mid p(x)\}$（条件成立），
$B=\{x\mid q(x)\}$（结论成立），则要证 $A\subseteq B$。
反证法不是直接证 $A\subseteq B$，而是证明\key{反例集为空}：
\[
  A\cap\complement_U B=\varnothing.
\]
做法：假设存在反例 $x_0\in A\cap\complement_U B$（满足条件却不满足结论），推出矛盾，
从而交集为空，即 $A\subseteq B$。

\begin{figure}[htbp]
  \centering
  \begin{tikzpicture}[>=Stealth, semithick, font=\footnotesize]
    \fill[gray!12] (0,0) rectangle (8,5);
    \draw[thick, main!70!black] (0,0) rectangle (8,5);
    \node[main!70!black] at (7.5,4.5) {$U$};
    \fill[main!10, main!70!black] (4.2,2.5) circle (1.8);
    \fill[second!20, second] (3.7,2.5) circle (0.9);
    \node[main!70!black] at (4.9,3.7) {$B$};
    \node[second] at (3.7,2.5) {$A$};
    \node[gray!60!black] at (1.1,2.6) {$\complement_U B$};
    \node[red!70!black] at (1.1,1.6) {$A\cap\complement_U B=\varnothing$};
  \end{tikzpicture}
  \caption{反证法的集合图景：条件集合 $A$ 完全落在结论集合 $B$ 内，反例集 $A\cap\complement_U B$ 为空}
  \label{fig:contradiction}
\end{figure}

反证法的第一步是明确写出“反设”：$A\not\subseteq B$，即存在 $x\in A$ 且 $x\notin B$。
图 \ref{fig:contradiction-assumption} 画出这一反设态；随后由条件与 $B$ 的定义推出矛盾，红色月牙收缩为空，回到图 \ref{fig:contradiction} 的结论态。

\begin{figure}[htbp]
  \centering
  \begin{tikzpicture}[>=Stealth, semithick, font=\footnotesize]
    \fill[gray!12] (0,0) rectangle (8,5);
    \draw[thick, main!70!black] (0,0) rectangle (8,5);
    \node[main!70!black] at (7.5,4.5) {$U$};

    \coordinate (B) at (2.7,2.6);
    \coordinate (A) at (4.9,2.6);

    % B（灰底）与 A；再把 A\B 涂成红色月牙
    \fill[main!14] (B) circle (1.6);
    \begin{scope}
      \clip (A) circle (1.4);
      \fill[red!50] (A) circle (1.4);
      \fill[main!14] (B) circle (1.6);
    \end{scope}

    \draw[main!60!black, thick] (B) circle (1.6);
    \draw[second!70!black, thick] (A) circle (1.4);

    \node[main!70!black] at (1.85,3.9) {$B$};
    \node[second!80!black] at (5.5,3.55) {$A$};
    \node[gray!60!black] at (0.75,4.4) {$\complement_U B$};

    \fill[red!80!black] (5.8,2.6) circle (2.2pt);
    \node[red!80!black, right] at (5.87,2.6) {$x$};
    \node[red!70!black, align=center] at (6.05,1.0) {$A\cap\complement_U B$\\ $\neq\varnothing$};
    \draw[red!60!black, ->, shorten >=1pt] (6.25,1.3) -- (5.9,2.1);

    \node[red!70!black] at (4.0,0.55) {反设：$\exists x\in A$ 且 $x\notin B$};
  \end{tikzpicture}
  \caption{反证法的反设态（正难则反）：假设 $A\not\subseteq B$，集合 $A$ 有一部分越出 $B$，红色月牙区域 $A\cap\complement_U B$ 非空，反例点 $x$ 落在其中}
  \label{fig:contradiction-assumption}
\end{figure}

\key{触发信号词：}
\begin{itemize}
  \item \key{唯一性}（“只有一个”“恰有一个”“唯一”）——假设存在两个不同的对象，推出矛盾（或证明任意两个相等）；
  \item \key{不存在}（“不存在”“没有”“无”）——假设存在，推出矛盾；
  \item \key{无理数}（“$\sqrt2$ 是无理数”）——假设可写成既约分数，推出分子分母同为偶数，与既约矛盾；
  \item \key{无限性}（“有无穷多个”）——假设只有有限个，推出还能构造出下一个，矛盾。
\end{itemize}

\begin{example}
用反证法证明：$\sqrt2$ 不是有理数。
\end{example}
\begin{solution}
假设 $\sqrt2$ 是有理数，则可写成既约分数 $\sqrt2=\dfrac pq$
（$p,q\in\mathbb{Z}^*$，$q>0$，$\gcd(p,q)=1$）。
两边平方得 $p^2=2q^2$，故 $p^2$ 为偶数，从而 $p$ 为偶数（奇数的平方仍为奇数）。
设 $p=2k$，代入得 $4k^2=2q^2$，即 $q^2=2k^2$，故 $q$ 也为偶数。
于是 $p,q$ 均为偶数，与 $\gcd(p,q)=1$ 矛盾。故假设不成立，$\sqrt2$ 是无理数。
\end{solution}

\key{探究四：数学归纳法——皮亚诺公理与最小反例原理}
数学归纳法不是“递推直觉”，而是自然数结构的公理推论。皮亚诺公理（$\mathbb{N}^*$）要点：
\begin{enumerate}
  \item $1\in\mathbb{N}^*$；
  \item 每个 $n\in\mathbb{N}^*$ 有唯一后继 $n+1$；
  \item $1$ 不是任何数的后继；
  \item 后继函数是单射：$m+1=n+1\Rightarrow m=n$；
  \item \key{归纳公理}：若 $S\subseteq\mathbb{N}^*$ 满足 $1\in S$ 且 $n\in S\Rightarrow n+1\in S$，则 $S=\mathbb{N}^*$。
\end{enumerate}
与归纳公理等价的是\key{良序原理}：$\mathbb{N}^*$ 的任一非空子集必有最小元。
由此得\key{最小反例原理}：若命题 $P(n)$ 不恒成立，则存在最小反例 $m$。
由 $m$ 最小知 $P(m-1)$ 成立，再用归纳递推得 $P(m)$ 成立，矛盾；故 $P(n)$ 恒成立。

因此数学归纳法两步缺一不可：
\begin{itemize}
  \item \key{归纳基础}：验证 $P(n_0)$ 成立（第一块骨牌倒下）；
  \item \key{归纳递推}：假设 $P(k)$ 成立，证 $P(k+1)$ 成立（前一块倒下必带倒后一块）。
\end{itemize}
只有基础是“只倒第一块”；只有递推是“没有启动”。归纳起点 $n_0$ 可以平移（不一定从 $1$ 开始）。

\begin{figure}[htbp]
  \centering
  \begin{tikzpicture}[>=Stealth, semithick, font=\footnotesize]
    \draw[->, gray] (0,0)--(10.8,0) node[right]{$n\in\mathbb{N}^*$};
    \foreach \i/\lab in {0/1,1/2,2/3,3/4,4/5,5/6} {
      \pgfmathsetmacro{\xx}{0.8+1.5*\i}
      \draw[main!70!black, fill=main!12, rounded corners=0.8pt] (\xx,0) rectangle ++(0.5,2.0);
      \node at (\xx+0.25,1.0) {$\lab$};
    }
    \foreach \i in {0,1,2,3,4} {
      \pgfmathsetmacro{\xx}{1.3+1.5*\i}
      \draw[->, second, thick] (\xx,1.0)--++(0.55,0);
    }
    \node[below, main] at (0.8,-0.15) {推倒第 $1$ 块};
    \node[above, second] at (5.0,2.3) {若第 $k$ 块倒下，则第 $k+1$ 块必倒下};
    \node[below, gray!70!black] at (8.2,-0.15) {$\cdots$};
    \node[below, gray!70!black] at (9.8,-0.15) {全部倒下};
  \end{tikzpicture}
  \caption{数学归纳法的多米诺模型：归纳基础推倒第 1 块，归纳递推保证“倒一块必带倒下一块”，于是全部倒下}
  \label{fig:induction}
\end{figure}

\begin{example}
用数学归纳法证明：$7^n-1$ 能被 $6$ 整除（$n\in\mathbb{N}^*$）。
\end{example}
\begin{solution}
(1) 基础：$n=1$ 时 $7^1-1=6$，能被 $6$ 整除。
(2) 递推：假设 $n=k$ 时 $7^k-1$ 能被 $6$ 整除，即 $7^k-1=6m$。则 $n=k+1$ 时
\[
  7^{k+1}-1=7\cdot7^k-1=7(7^k-1)+6=7\cdot6m+6=6(7m+1),
\]
能被 $6$ 整除。由(1)(2)，命题对一切 $n\in\mathbb{N}^*$ 成立。
\end{solution}

\subsection{武器库索引表与跨模块实战}

\key{《高中数学逻辑变形与证明触发信号对照表》}
把题干信号、集合／逻辑翻译、变形引擎与操作风险放在同一张表，解题时“对号入座”。

\begin{center}
\footnotesize
\setlength{\tabcolsep}{3pt}
\renewcommand{\arraystretch}{1.4}
\begin{tabular}{p{2.0cm}p{3.2cm}p{2.6cm}p{3.6cm}}
\toprule
题干信号 & 逻辑／集合翻译 & 变形引擎 & 操作与风险 \\
\midrule
集合相等 $A=B$ & 双向包含 & 集合相等公理 & 分证 $A\subseteq B$ 与 $B\subseteq A$ \\
\hline
平方、去绝对值 & $f^2=g^2\iff f=\pm g$ & 同解变形 & 解集扩大，必须验根 \\
\hline
同除含参因式 & 可能丢根 & 移项分解＋讨论 & 解集缩小，防漏根 \\
\hline
对数方程 & 真数 $>0$ & 同解变形＋定义域 & 忽略定义域则增根 \\
\hline
至少有一个 & $\exists x$ & 量词否定 & 否定为 $\forall x,\ \neg$ \\
\hline
至多有一个 & 否定“至少两个” & 量词否定＋德摩根 & 命题的否定改写 \\
\hline
不都 & $\neg\forall$ & 量词否定 & 否定为“至少有一个不” \\
\hline
无解、为空集 & $S=\varnothing$ & 补集 & 补集为全集 $U$ \\
\hline
恒成立 & $\forall x,\ f(x)\in M$ & 值域子集 & $R_f\subseteq M$，看最小值 \\
\hline
能成立 & $\exists x,\ f(x)\in M$ & 值域交集 & $R_f\cap M\neq\varnothing$，看最大值 \\
\hline
唯一性 & 存在且唯一 & 反证法 & 假设第二个，推出矛盾 \\
\hline
不存在、没有 & $\neg\exists$ & 反证法 & 假设存在，推出矛盾 \\
\hline
无理数 & 非有理数 & 反证法 & 设 $\sqrt2=\frac pq$ 既约 \\
\hline
无限性 & 非有限 & 反证法 & 假设有限，构造下一个 \\
\hline
对所有 $n\in\mathbb{N}^*$ & 自然数命题 & 数学归纳法 & 基础＋递推，缺一不可 \\
\hline
整除、离散不等式 & 自然数命题 & 数学归纳法／最小反例 & 归纳起点可平移至 $n_0$ \\
\bottomrule
\end{tabular}
\end{center}

\begin{example}
解不等式 $\dfrac{x^2-4}{x-2}\geq0$，并说明直接约去 $(x-2)$ 会丢掉什么。
\end{example}
\begin{solution}
定义域要求 $x-2\neq0$。在定义域内
\[
  \frac{x^2-4}{x-2}=\frac{(x-2)(x+2)}{x-2}=x+2,
\]
故原不等式等价于 $x+2\geq0$ 且 $x\neq2$，即 $x\geq-2$ 且 $x\neq2$，
解集为 $[-2,2)\cup(2,+\infty)$。
若直接约去 $(x-2)$ 得 $x+2\geq0$，解集为 $[-2,+\infty)$，把分母不能为 $0$ 的限制丢了（漏掉限制，不是漏根，而是扩大了解集），
所以必须保留 $x\neq2$。
\end{solution}

\begin{example}
若“存在 $x\in\mathbb{R}$，使 $x^2+2ax+4<0$”是假命题，求 $a$ 的取值范围。
\end{example}
\begin{solution}
该命题为假 $\iff$ 其否定“$\forall x\in\mathbb{R},\ x^2+2ax+4\geq0$”为真。
这是恒成立问题：$\Delta=(2a)^2-16\leq0\iff4a^2-16\leq0\iff-2\leq a\leq2$。
故 $a\in[-2,2]$。
\end{solution}

\begin{example}
证明：函数 $f(x)=x^3+x+1$ 有且只有一个零点。
\end{example}
\begin{solution}
\key{存在性}：$f(-1)=-1-1+1=-1<0$，$f(0)=1>0$，且 $f$ 是多项式（连续），
由零点存在定理，$(-1,0)$ 内至少有一个零点。
\key{唯一性（反证）}：假设有两个不相等的零点 $x_1,x_2$（设 $x_1<x_2$），则
\[
  x_1^3+x_1+1=0,\qquad x_2^3+x_2+1=0.
\]
两式相减得 $(x_1^3-x_2^3)+(x_1-x_2)=0$，即
\[
  (x_1-x_2)(x_1^2+x_1x_2+x_2^2+1)=0.
\]
因 $x_1\neq x_2$，需 $x_1^2+x_1x_2+x_2^2+1=0$。
但 $x_1^2+x_1x_2+x_2^2=\left(x_1+\frac{x_2}{2}\right)^2+\frac34x_2^2\geq0$，
故 $x_1^2+x_1x_2+x_2^2+1\geq1>0$，矛盾。故零点唯一。
\end{solution}

\begin{example}
用数学归纳法证明：$n^3+5n$ 能被 $6$ 整除（$n\in\mathbb{N}^*$）。
\end{example}
\begin{solution}
(1) 基础：$n=1$ 时 $1+5=6$，能被 $6$ 整除。
(2) 递推：假设 $n=k$ 时 $k^3+5k=6m$。则 $n=k+1$ 时
\begin{align*}
  (k+1)^3+5(k+1)&=k^3+3k^2+3k+1+5k+5\\
  &=(k^3+5k)+3k^2+3k+6\\
  &=6m+3k(k+1)+6.
\end{align*}
因为 $k(k+1)$ 是两个连续整数之积，必为偶数，故 $3k(k+1)$ 能被 $6$ 整除；又 $6$ 能被 $6$ 整除。
所以 $(k+1)^3+5(k+1)$ 能被 $6$ 整除。由(1)(2)，结论对一切 $n\in\mathbb{N}^*$ 成立。
\end{solution}

\begin{example}
用数学归纳法证明：对一切 $n\geq5$，$2^n>n^2$。
\end{example}
\begin{solution}
(1) 基础：$n=5$ 时 $2^5=32>25=5^2$，成立。
(2) 递推：假设 $n=k$（$k\geq5$）时 $2^k>k^2$。则 $n=k+1$ 时
\[
  2^{k+1}=2\cdot2^k>2k^2.
\]
要证 $2k^2>(k+1)^2=k^2+2k+1$，即 $k^2-2k-1>0$，即 $(k-1)^2>2$。
因 $k\geq5$，$(k-1)^2\geq16>2$，成立。故 $2^{k+1}>(k+1)^2$。
由数学归纳法，对一切 $n\geq5$，$2^n>n^2$（注意归纳起点平移到了 $5$）。
\end{solution}

\begin{example}
设命题 $p$：$\forall x\in\mathbb{R},\ x^2+ax+1>0$；命题 $q$：$\exists x\in\mathbb{R},\ x^2-2x+a=0$。
若 $p\lor q$ 为假，求 $a$ 的取值范围。
\end{example}
\begin{solution}
$p\lor q$ 为假 $\iff$ $p$ 假且 $q$ 假（德摩根律）。
$p$ 真 $\iff\Delta_1=a^2-4<0\iff-2<a<2$，故 $p$ 假 $\iff a\leq-2$ 或 $a\geq2$。
$q$ 真 $\iff\Delta_2=4-4a\geq0\iff a\leq1$，故 $q$ 假 $\iff a>1$。
取交：$(a\leq-2\ \text{或}\ a\geq2)\cap(a>1)$，得 $a\geq2$。
故 $a\in[2,+\infty)$。
\end{solution}

\begin{conccard}[三级]{反证法触发信号：唯一性、不存在、无理数、无限性}{结论含“只有一个／不存在／无理数／无穷多”的证明题}{把反证法当成“随便假设”；假设后没有推出矛盾就下结论}
反证法即证反例集 $A\cap\complement_U B=\varnothing$：假设反例存在，推出与公理、定理或临时假设矛盾的结果。
唯一性假设“至少两个”；不存在假设“存在”；无理数设既约分数 $\frac pq$；无限性设“只有有限个”。
矛盾必须明确写出，否则证明不完整。
\end{conccard}

\begin{conccard}[三级]{数学归纳法：归纳基础＋归纳递推}{证明对全体自然数（或 $n\geq n_0$）成立的命题}{只做递推不做基础；归纳假设没有真正用上；归纳起点误写为 $1$}
归纳公理：若 $1\in S$ 且 $n\in S\Rightarrow n+1\in S$，则 $S=\mathbb{N}^*$。
其根基是良序原理与最小反例原理：若存在反例，则存在最小反例，与递推矛盾。
使用时要检查两点：归纳基础是否验证；递推中是否用到了归纳假设。起点可按命题需要平移到 $n_0$。
\end{conccard}

\begin{summary}
本章把集合与逻辑锻造成一套“武器库”：集合相等公理 $A=B\iff(A\subseteq B)\land(B\subseteq A)$ 是审查一切同解变形的标准（$A\subsetneqq B$ 增根、$B\subsetneqq A$ 漏根）；德摩根律、逆否等价、量词否定与补集转化处理“至多、至少、不都、无解、恒成立、能成立”；反证法即证反例集 $A\cap\complement_U B=\varnothing$，由“唯一性、不存在、无理数、无限性”四类信号词触发；数学归纳法以皮亚诺公理与良序原理为根基，用最小反例原理保证“归纳基础＋归纳递推”两步缺一不可。解题时先翻译题干（集合／逻辑语言），再选引擎，最后审方向（防增漏根）。
\end{summary}

\end{explore}

\subsection{四视角等价地图：同一对象，多副面孔}

前四个模块其实都在做同一件事：给同一个数学对象换一副面孔。
代数、几何、集合、映射（对象本身）不是四门课，而是同一结构的四种\key{表示}；
只要两种表示之间存在\emph{保持结构的对应}（同构），命题就能互相翻译。
但「等价」不等于「一样难」：\textbf{同一问题在不同框架下的可解释性与计算量差别极大}，
所以解题的关键动作是\key{找一个同构映射，把问题搬到更顺的那一侧}。

\begin{center}
\footnotesize
\begin{tabular}{c|p{2.6cm}|p{3.2cm}|p{3.0cm}|p{3.4cm}}
\toprule
视角 & 看什么 & 擅长 & 短板 & 典型「搬运」 \\
\midrule
代数 & 符号、方程、恒等变形 & 机械化、可执行、可检验 & 缺直观，易迷失方向 & 几何 $\to$ 方程（解析几何）\\
几何 & 图形、变换、位置关系 & 直观、好找思路、便于猜答案 & 精度与严格性弱 & 方程 $\to$ 图形（数形结合）\\
集合 & 元素、包含、交并补、量词 & 严谨分类、盯住边界与反例 & 不给算法 & 命题 $\to$ 集合关系（充分必要）\\
映射／对象 & 对应关系、迭代、结构本身 & 统一语言、跨领域搬运结论 & 太抽象，需落地实例 & 矩阵 $=$ 线性映射；坐标 $=$ 同构 \\
\bottomrule
\end{tabular}
\end{center}

\begin{figure}[htbp]
\centering
\begin{tikzpicture}[>=Stealth,semithick,font=\footnotesize,
  box/.style={draw,rounded corners,inner sep=4pt,align=center,minimum width=2.1cm}]
  \node[box,draw=main] (A) at (-2.2,1.9) {代数\\方程};
  \node[box,draw=second] (G) at (2.2,1.9) {几何\\图形};
  \node[box,draw=third] (S) at (-2.2,-1.9) {集合\\逻辑};
  \node[box,draw=gray!60!black] (M) at (2.2,-1.9) {映射\\结构};
  \node[box,draw=main!70!black,minimum width=2.4cm] (C) at (0,0) {\textbf{同一个}\\数学对象};
  \draw[<->,main] (C)--(A) node[midway,above left,font=\scriptsize] {列方程};
  \draw[<->,second] (C)--(G) node[midway,above right,font=\scriptsize] {画出来};
  \draw[<->,third] (C)--(S) node[midway,below left,font=\scriptsize] {下定义};
  \draw[<->,gray!60!black] (C)--(M) node[midway,below right,font=\scriptsize] {建模型};
  \draw[->,gray!55,dashed] (A)--(G) node[midway,above,font=\scriptsize] {数形结合};
  \draw[->,gray!55,dashed] (S)--(M) node[midway,below,font=\scriptsize] {关系即映射};
\end{tikzpicture}
\caption{四视角等价地图：中心是对象本身，四个顶点是它的四种表示；
箭头都是「保持结构的映射」。解题时沿箭头换到顺手的一侧，算完再换回来}
\label{fig:four-view-map}
\end{figure}

\begin{figure}[htbp]
\centering
\begin{tikzpicture}[>=Stealth,semithick,font=\footnotesize,
  ch/.style={draw,rounded corners,inner sep=5pt,align=center,minimum width=2.0cm},
  lanename/.style={rotate=90,anchor=south,font=\scriptsize\bfseries,text=main!70!black}]
  % 三条主线 + 一条独立线
  \node[lanename] at (-1.1,6) {几何·结构线};
  \node[lanename] at (-1.1,3) {离散动力线};
  \node[lanename] at (-1.1,0.6) {连续分析线};
  % 第 0 层：语言底座
  \node[ch,draw=main,fill=main!8] (c0) at (1.6,7.6) {第 0 章\\集合·逻辑·证明};
  \node[ch,draw=gray!60!black,fill=gray!8] (c7) at (11.4,7.6) {第 7 章\\概率·统计};
  \draw[->,main] (c0)--(c7) node[midway,above,font=\scriptsize] {测度＝归一化计数};
  % 几何·结构线
  \node[ch,draw=second] (c1) at (1.6,6) {第 1 章\\函数};
  \node[ch,draw=second] (c2) at (5.0,6) {第 2 章\\三角·复数};
  \node[ch,draw=second] (c3) at (8.2,6) {第 3 章\\向量};
  \node[ch,draw=second] (c4) at (11.4,6) {第 4 章\\立几·圆锥曲线};
  \draw[->,second] (c1)--(c2);
  \draw[->,second] (c2)--(c3);
  \draw[->,second] (c3)--(c4);
  % 离散动力线
  \node[ch,draw=third] (c5) at (1.6,3) {第 5 章\\数列·递推};
  \node[ch,draw=third] (c9) at (11.4,3) {第 9 章\\不动点};
  \draw[->,third] (c5)--(c9) node[midway,above,font=\scriptsize] {递推求通项·收敛};
  % 连续分析线
  \node[ch,draw=main!60!black] (c6) at (5.0,0.6) {第 6 章\\导数·局部线性化};
  \node[ch,draw=main!60!black] (c8) at (8.2,0.6) {第 8 章\\齐次化};
  \draw[->,main!60!black] (c6)--(c9) node[midway,right,font=\scriptsize] {稳定性};
  \draw[->,main!60!black] (c8)--(c9) node[midway,right,font=\scriptsize] {射影};
  % 跨线接口
  \draw[->,gray!55,dashed] (c1)--(c5) node[midway,left,font=\scriptsize] {离散函数};
  \draw[->,gray!55,dashed] (c1)--(c6) node[midway,left,font=\scriptsize] {变化率};
  \draw[<->,gray!55,dashed] (c5)--(c6) node[midway,below,font=\scriptsize] {差分 $\leftrightarrow$ 导数};
  \draw[->,gray!55,dashed] (c4)--(c8) node[midway,right,font=\scriptsize] {条件最值};
  \draw[->,main!35] (c0)--(c1);
  \draw[->,gray!55,dashed] (c3)--(c9) node[midway,above right,font=\scriptsize] {特征方向};
\end{tikzpicture}
\caption{全书知识脉络图：实线是各自主线的推进，虚线是跨主线接口。
三条主线（几何·结构／离散动力／连续分析）最终在第 9 章不动点汇合；
第 8 章齐次化是射影观点的入口，第 0 章语言底座供全局调用}
\label{fig:book-map}
\end{figure}

\textbf{脉络怎么读}：几何·结构线从函数到圆锥曲线，解决「形的问题」；
离散动力线与连续分析线在 \key{差分 $\leftrightarrow$ 导数}处对接（同一个
「以直代曲／以比代差」思想），分别解决「列的问题」与「极限的问题」，
并共同汇入第 9 章的不动点理论；第 8 章齐次化把二元问题压到一元比例，
是通往射影／映射观点的接口。每一章章首的「多框架对读」卡，就是在这张
图上标明「你现在在哪条线、哪一侧」。

\textbf{怎么用它找思路}：拿到题先问「我现在站在哪个视角」——
卡在代数就画图（第 2、4 章的向量与圆锥曲线）；卡在几何就设坐标列方程
（第 4 章的仿射、第 8 章的齐次化）；被式子的次数与比值绕晕就升到射影／映射
观点（第 9 章的不动点与 Möbius）；讨论「有没有、有几个」就回到集合与量词
（本章的同解变形与反证法）。\textbf{同一结论往往能在四处各证一遍}，选最省力的那处书写。
此后各章章首统一附「多框架对读」卡，把这条暗线显性化。

\begin{inclasstest}
\begin{enumerate}
  \item 写出“$\exists x\in\mathbb{R},\ x^2-2x+1\leq0$”的否定，并判断其真假；\level{易}
  \item 解方程 $\sqrt{x+3}=x+1$，并说明验根的必要性；\level{中}
  \item 已知“存在 $x\in[1,3]$，使 $x^2-ax+2\leq0$”为真，求 $a$ 的取值范围；\level{中}
  \item 用反证法证明：若 $a,b\in\mathbb{Z}$ 且 $ab$ 能被 $3$ 整除，则 $a,b$ 中至少有一个能被 $3$ 整除。\level{中}
\end{enumerate}
\end{inclasstest}

\begin{afterclass}
\subsection*{基础巩固}

\begin{enumerate}
  \item 写出德摩根律与量词否定的四条公式：$\neg(p\land q)$、$\neg(p\lor q)$、$\neg(\forall x,P(x))$、$\neg(\exists x,P(x))$；
  \item 解方程 $x^2=2x$，并说明若两边同除以 $x$ 会漏掉哪个根。
\end{enumerate}

\subsection*{进阶提升}
\begin{enumerate}
  \item 已知“$\forall x\in\mathbb{R},\ x^2+2x+a>0$”为真，求 $a$；再写出它的否定命题为真时 $a$ 的范围；
  \item 用数学归纳法证明：$5^n-1$ 能被 $4$ 整除。
\end{enumerate}

\subsection*{拓展探究}
\begin{enumerate}
  \item 用反证法证明：$\sqrt3$ 是无理数。（提示：用到上面第 4 题的结论“若 $3\mid ab$，则 $3\mid a$ 或 $3\mid b$”。）
\end{enumerate}

\subsection*{教师版答案}
\begin{solution}
\textbf{基础 1：}$\neg(p\land q)\iff\neg p\lor\neg q$；$\neg(p\lor q)\iff\neg p\land\neg q$；
$\neg(\forall x,P(x))\iff\exists x,\neg P(x)$；$\neg(\exists x,P(x))\iff\forall x,\neg P(x)$。
\end{solution}
\begin{solution}
\textbf{基础 2：}移项分解 $x(x-2)=0$，得 $x=0$ 或 $x=2$，解集 $\{0,2\}$。
若两边同除以 $x$ 得 $x=2$，漏掉 $x=0$（同除使解集由 $\{0,2\}$ 缩小为 $\{2\}$）。
\end{solution}
\begin{solution}
\textbf{进阶 1：}$x^2+2x+a>0$ 恒成立 $\iff\Delta=4-4a<0\iff a>1$。
其否定“$\exists x\in\mathbb{R},\ x^2+2x+a\leq0$”为真 $\iff\Delta\geq0\iff a\leq1$。
\end{solution}
\begin{solution}
\textbf{进阶 2：}(1) $n=1$ 时 $5-1=4$，能被 $4$ 整除。
(2) 假设 $5^k-1=4m$，则 $5^{k+1}-1=5(5^k-1)+4=20m+4=4(5m+1)$，能被 $4$ 整除。
由数学归纳法，结论成立。
\end{solution}
\begin{solution}
\textbf{拓展：}假设 $\sqrt3$ 是有理数，写成既约分数 $\sqrt3=\frac pq$（$\gcd(p,q)=1$）。
平方得 $p^2=3q^2$，故 $3\mid p^2$；由提示（取 $a=b=p$）得 $3\mid p$，设 $p=3k$。
代入得 $9k^2=3q^2$，即 $q^2=3k^2$，故 $3\mid q$。于是 $p,q$ 均为 $3$ 的倍数，与既约矛盾。
故 $\sqrt3$ 是无理数。
\end{solution}
\reading{拓展阅读 00}{逻辑与证明的武器库}
\end{afterclass}

% 高观点导学案：函数的概念与基本初等函数
% 主线：关系 R⊆A×B（三档舞台）→ 方程/不等式＝关系的条件表达 → 两条限制轴（单值性→映射/函数；连续性与结构保持→基本初等函数）
% 全程三视角互译：代数有序数对集合 ↔ 几何坐标系点集 ↔ 映射箭头图

\section{函数的概念与基本初等函数（高观点结构版）}

\begin{methodbox}{多框架对读}
\small
\begin{center}
\begin{tabular}{c|p{2.9cm}|p{5.0cm}|p{3.1cm}}
\toprule
视角 & 本章看点 & 别的视角怎么说 & 哪侧最省力 \\
\midrule
代数 & 解析式、复合与解方程 & 有序数对集合的坐标运算 & 变形与恒等式 \\
几何 & 图象的平移／对称变换 & 坐标系中点集的移动 & 猜性质、定范围 \\
集合 & 定义域、值域、$f:A\to B$ & 取值范围即约束集合 & 审边界、找反例 \\
映射 & $y=f(x)$ 即对应法则 & 变换＝作用在图象上的映射 & 统一语言、整体搬迁 \\
\bottomrule
\end{tabular}
\end{center}
卡的使用：卡住时先判断自己站在哪一栏，再沿箭头换栏（总图见第 0 章图~\ref{fig:four-view-map}）。
\end{methodbox}

% ---------- 前置引入：舞台的诞生 ----------
\subsection{前置引入：舞台的诞生——集合自身的四级进化}

研究任何数学对象，现代数学靠两根支柱：\key{舞台}（集合自身的结构）与 \key{规则}（集合之间的映射）。本导学案主线一直在讲\key{规则}：$A\times B$ 中怎样挑点、怎样过滤出单值对应。但高中函数之所以拥有单调性、奇偶性、零点存在定理与导数，原因不在规则里，而在\key{舞台}上——定义域与值域选中的是哪种集合。

\begin{conccard}[一级]{集合装了多少「装备」，函数就有多少性质}{适用于所有「新定义」与概念辨析题}{只背性质条款，不问性质从何而来}
一般集合（例如 $\{\text{张三},\text{李四}\}$）只是一堆孤立元素：不能比大小，不能做加减，也没有「挨得很近」的概念。一旦把定义域与值域锁定在实数集 $\mathbb{R}$ 上，$\mathbb{R}$ 自带的三层结构便一次性注入函数。
\end{conccard}

\begin{itemize}
  \item \level{Lv.1} \key{原始集合}——一袋散装弹珠。只有「属于」与「元素个数」，无法比大小、不能运算：既谈不上单调性，也画不出坐标系中的图象。
  \item \level{Lv.2} \key{有序与代数集合}（$\mathbb{Z},\mathbb{Q}$）——排好队且能算账的刻度尺。\key{全序结构}（$x_1<x_2$ 可比大小）解锁\key{单调性与最值}；\key{代数运算与相反数}（$+,-,\times,\div$ 与 $-x$）解锁\key{奇偶性与周期性}。复数域 $\mathbb{C}$ 不能比大小，所以复变函数根本没有「单调递增」可言。
  \item \level{Lv.3} \key{稠密但有孔的集合}（$\mathbb{Q}$）——布满微观针眼的筛子。任意两点之间仍有无穷多个点，却漏掉了 $\sqrt{2},\pi,e$ 等无数无理数。致命后果见图~\ref{fig:qline}：$f(x)=x^2-2$ 从负值升到正值时，正好从 $\sqrt{2}$ 的孔洞中漏过去，\key{根本没有零点}。
  \item \level{Lv.4} \key{完备连续集合}（$\mathbb{R}$）——无缝焊接的实心金条。无理数填满全部孔洞，曲线无处可漏，\key{零点存在定理}（介值性）、极限以及「以直代曲」的导数才真正成为可能。
\end{itemize}

Lv.1 与 Lv.2 的差别不在元素本身，而在元素之间被装上了什么关系：前者只有「属于」与计数（图 \ref{fig:bag2axis} 左），后者多出了「谁大谁小」的全序与「怎么算账」的代数运算（图 \ref{fig:bag2axis} 右）。

\begin{figure}[htbp]
  \centering
  \begin{tikzpicture}[>=stealth,font=\small]
    % ---------- 左：Lv.1 一袋散装弹珠（无序点集） ----------
    \begin{scope}
      \draw[gray!60,very thick,fill=main!8] (0,0) ellipse (2.25cm and 1.6cm);
      \foreach \px/\py/\lab in {
        -1.45/0.62/-2, -0.95/-0.88/-1, -0.42/0.28/-\tfrac12,
        0.02/-0.55/0, 0.42/0.78/\tfrac12, 0.88/-0.18/1,
        1.35/0.5/\tfrac32, 1.12/-1.0/2, -1.78/-0.32/\tfrac52}
      {
        \fill[second] (\px,\py) circle (0.075);
        \node[gray!80,font=\scriptsize] at (\px+0.16,\py+0.16) {$\lab$};
      }
      \node at (0,-1.95) {\small 只有「属于」与元素个数};
      \node[second] at (0,-2.42) {\small 不能比大小，也不能加减};
    \end{scope}
    % ---------- 右：Lv.2 带刻度的数轴（有序集合 Q） ----------
    \begin{scope}[xshift=7.2cm]
      \draw[->,gray,very thick] (-2.7,0)--(3.1,0) node[right,gray] {$x$};
      \foreach \t in {-2,-1,0,1,2}
        \draw[gray,thick] (\t,0.09)--(\t,-0.09);
      \foreach \t in {-2,-1,0,1,2}
        \node[gray,font=\scriptsize] at (\t,-0.34) {$\t$};
      \foreach \t in {-2,-1,-0.5,0,0.5,1,1.5,2,2.5}
        \fill[main] (\t,0) circle (0.05);
      \node[third] at (0,1.15) {$(<)\quad +,\,-,\,\times,\,\div$};
      \node at (0,-1.35) {\small 排好队且能算账的刻度尺};
      \node[main] at (0,-1.82) {\small 全序解锁单调最值，运算解锁奇偶周期};
    \end{scope}
  \end{tikzpicture}
  \caption{从「属于」到「可比大小」：左为 Lv.1 原始集合——一袋散装弹珠，只有属于关系与元素个数，无法比大小、不能运算；右为 Lv.2 有序与代数集合——同一批元素排上带刻度的数轴 $\mathbb{Q}$，全序 $(<)$ 解锁单调性与最值，$+,-,\times,\div$ 与相反数解锁奇偶性与周期性}
  \label{fig:bag2axis}
\end{figure}

\begin{figure}[htbp]
  \centering
  \begin{tikzpicture}[x=0.92cm,y=0.92cm,>=stealth,font=\small]
    % ---------- 左：有理数轴（穿孔） ----------
    \begin{scope}
      \draw[->,white!55] (-0.7,0)--(2.9,0) node[right,white!70] {$x$};
      \foreach \x in {0.1,0.35,0.62,0.9,1.15,1.62,1.85,2.1,2.4,2.65}
        \draw[second,fill=white] (\x,0) circle (0.042);
      \draw[second,very thick,fill=white] (1.414,0) circle (0.075);
      \node[second] at (1.414,-0.42) {\small $\sqrt{2}$（洞）};
      \draw[main,thick,domain=0.55:2.0,samples=70,smooth] plot(\x,{\x*\x-2});
      \node[main] at (1.45,2.30) {$y=x^2-2$};
      \node[white!70] at (1.05,-2.02) {\small $\mathbb{Q}$：看似连续，实则满布针眼};
      \node[second] at (1.05,-2.56) {\small 穿过空洞：无交点};
    \end{scope}
    % ---------- 右：实数轴（实心） ----------
    \begin{scope}[xshift=7.6cm]
      \draw[->,third,very thick] (-0.7,0)--(2.9,0) node[right,third] {$x$};
      \foreach \x in {0.1,0.35,0.62,0.9,1.15,1.62,1.85,2.1,2.4,2.65}
        \fill[third!70] (\x,0) circle (0.032);
      \draw[main,thick,domain=0.55:2.0,samples=70,smooth] plot(\x,{\x*\x-2});
      \fill[second] (1.414,0) circle (0.075);
      \node[second] at (1.414,-0.42) {$(\sqrt{2},0)$};
      \node[third] at (1.05,-2.02) {\small $\mathbb{R}$：无缝连续};
      \node[second] at (1.05,-2.56) {\small 洞被填满：零点被迫现身};
    \end{scope}
  \end{tikzpicture}
  \caption{同一函数 $f(x)=x^2-2$ 在两种「舞台」上的命运：左，定义域取有理数集 $\mathbb{Q}$，曲线从 $\sqrt{2}$ 的空心孔洞中漏过，变号却无零点；右，定义域取实数集 $\mathbb{R}$，孔洞被无理数填满，零点 $(\sqrt{2},0)$ 必然出现}
  \label{fig:qline}
\end{figure}

\subsubsection*{开场三问}
\begin{enumerate}
  \item 复数域 $\mathbb{C}$ 上能否定义「单调递增」？为什么？（提示：$\mathbb{C}$ 上有没有「小于」关系）\blank
  \item 在有理数集 $\mathbb{Q}$ 上考察 $f(x)=x^2-3$：它满足 $f(1)<0<f(2)$，却没有零点。请用「孔洞」的语言解释。\blank
  \item 狄利克雷函数 $D(x)=\begin{cases}1,&x\in\mathbb{Q}\\0,&x\in\mathbb{R}\setminus\mathbb{Q}\end{cases}$ 以任一非零有理数为周期，却没有最小正周期。它击穿了「周期函数一定有最小正周期」这句直觉——你能再举一句被它击穿的说法吗？\blank
\end{enumerate}

\begin{method}
「看集合选武器」的反射：拿到一道涉及函数的题，先扫一眼定义域是哪类集合，再决定调用哪套工具——
$\mathbb{N}^{*}$（离散有序）$\Rightarrow$ 差分、数学归纳法、不动点；
$[a,b]\subseteq\mathbb{R}$（连续完备）$\Rightarrow$ 单调性、零点定理、导数局部线性化；
$\mathbb{Q}$ 与 $\mathbb{R}\setminus\mathbb{Q}$ 混用 $\Rightarrow$ 立刻警觉稠密性与孔洞性，准备构造反例。
\end{method}

\begin{commonerror}{把「有界」当成「能取到最值」}
$f(x)>a$ 在 $(0,1)$ 上恒成立，即使 $f$ 在该区间单调递增且 $\lim\limits_{x\to0^{+}}f(x)=2$，也只能推出 $a\le 2$ 而不是 $a<2$：端点取不到时，边界值扮演的是\key{下确界}而非最小值。这类「端点开合」丢分，考的正是实数的完备性。
\end{commonerror}

四级进化完成后，$\mathbb{R}$ 与自身作笛卡尔积，舞台张成连续平面（图 \ref{fig:r2plane}）——后续的二维区域与函数图象都落在这块舞台上。

\begin{figure}[htbp]
  \centering
  \begin{tikzpicture}[>=stealth,font=\small,scale=1.02]
    \draw[third!35,dashed,thin] (-2.5,-2) grid (2.5,2);
    \draw[->,gray,very thick] (-2.75,0)--(2.75,0) node[right,gray] {$x$};
    \draw[->,gray,very thick] (0,-2.25)--(0,2.25) node[above,gray] {$y$};
    \foreach \t in {-2,-1,1,2}{
      \node[gray,font=\scriptsize] at (\t,-0.18) {$\t$};
      \node[gray,font=\scriptsize] at (-0.2,\t) {$\t$};
    }
    \fill[main!12] plot[smooth cycle] coordinates {(0.1,-1.1) (1.5,-0.7) (2.05,0.35) (1.25,1.35) (-0.15,1.55) (-1.35,0.65) (-1.55,-0.6)};
    \draw[main!70!black,very thick] plot[smooth cycle] coordinates {(0.1,-1.1) (1.5,-0.7) (2.05,0.35) (1.25,1.35) (-0.15,1.55) (-1.35,0.65) (-1.55,-0.6)};
    \node[main!80!black] at (0.45,0.35) {二维区域（「面」）};
    \fill[second] (1.414,0) circle (0.06);
    \node[second,below right,font=\footnotesize] at (1.414,0) {$(\sqrt2,0)$};
    \node[third,font=\large] at (2.05,1.95) {$\mathbb{R}^2$};
  \end{tikzpicture}
  \caption{舞台就绪：完备的 $\mathbb{R}$ 与自身作笛卡尔积张成连续平面 $\mathbb{R}^2$（坐标轴、网格与一块二维区域的定场）。零维的点已被连成无缝的线，线再铺成二维的面——不等式解集、函数图象等后续对象都以这块平面为舞台}
  \label{fig:r2plane}
\end{figure}

\subsection{实数的三种等价面孔：分割、基本列与确界}

前面用「针眼筛子」与「实心金条」画出 $\mathbb{Q}$ 与 $\mathbb{R}$ 的差别。本节把它变成可以动手验证的数学：孔洞究竟在哪里？怎么补？补完的实数，为什么恰好就是数轴上的点？
（数系为何还要从 $\mathbb{R}$ 继续走向 $\mathbb{C}$，见第 2 章「数系扩张与虚数单位 $i$ 的几何诞生」。）

\key{探究一：稠密也会漏——有理数的上确界会失踪}
设
\[
  S=\{x\in\mathbb{Q}\mid x^2<2\}.
\]
$S\neq\varnothing$（$1\in S$），也有上界（$2$ 就是一个上界）。若在有理数范围内「有上界必有最小上界」，$S$ 应当存在最小上界。事实是：不存在。

\begin{derivation}
任取正有理数 $r$，构造它的「平分修正」
\[
  t=\frac{r+\dfrac{2}{r}}{2}\in\mathbb{Q},\qquad
  t^2-2=\frac{(r^2-2)^2}{4r^2}\geq0,\qquad
  t-r=-\frac{r^2-2}{2r}.
\]
\begin{itemize}
  \item 若 $r^2>2$（此时 $r$ 是 $S$ 的上界）：仍有 $t^2>2$，即 $t$ 也是上界；但 $t<r$。于是任何有理上界都能被更小的有理上界取代，「最小上界」无从谈起。
  \item 若 $r^2<2$（此时 $r\in S$）：$t>r$ 且 $t^2>2$，即 $t$ 是一个比 $r$ 更大的上界。
\end{itemize}
两种情形合起来：$\mathbb{Q}$ 中既没有 $S$ 的最小上界，也没有恰好落在 $\sqrt2$ 处的点。孔洞不是文学比喻，而是确界原理在 $\mathbb{Q}$ 中失效。
\end{derivation}

\key{探究二：补洞的三条等价路线}
\begin{itemize}
  \item \key{戴德金分割}：把 $\mathbb{Q}$ 切成上下两堆 $(A\mid B)$，要求 $A$ 无最大元；每个分割定义一个实数。上面由 $S$ 诱导的分割 $\bigl(\{x\in\mathbb{Q}:x^2<2\}\mid\mathbb{Q}\setminus\{x\in\mathbb{Q}:x^2<2\}\bigr)$ 恰好补上缺失的 $\sqrt2$。
  \item \key{柯西基本列}：把「彼此无限靠近」的有理数列（如 $3,\,3.1,\,3.14,\,\dots$）视为一个实数；两列若差趋于 $0$ 就视为同一个实数。$\pi$ 是这样被「逼」出来的。
  \item \key{公理刻画}（只描述、不构造）：确界原理、单调有界定理、闭区间套定理、柯西收敛准则四者等价，说的是同一句话——$\mathbb{R}$ 没有孔洞。高中阶段的极限、零点存在定理、介值性全部建立在这条公理上。
\end{itemize}

\begin{center}
\footnotesize
\setlength{\tabcolsep}{3pt}
\renewcommand{\arraystretch}{1.4}
\begin{tabular}{p{2.4cm}p{3.6cm}p{3.0cm}p{3.8cm}}
\toprule
视角 & 怎么补洞 & 典型产物 & 高中落点 \\
\midrule
戴德金分割 & 用「上下两堆」的卡口位置定义新数 & $\sqrt2$ 这类不可公度长度 & 数轴上的点与实数一一对应 \\
柯西基本列 & 用「越来越近的有理列」定义新数 & $\pi,\ e$ 这类极限 & 数列极限、级数求和的合法性 \\
公理刻画 & 直接声明「有上界就有上确界」 & 完备性本身 & 零点定理、最值定理的理论根基 \\
\bottomrule
\end{tabular}
\end{center}

\key{探究三：数轴上的点恰好是实数}
稠密 $\neq$ 连续：$\mathbb{Q}$ 已经稠密（任意两点之间还有有理数），却仍漏点；只有补上完备性，才得到「每一条几何长度都对应唯一实数」的连续数轴。
十进制表示还带来两个容易含糊的点：
\begin{itemize}
  \item $0.999\dots=1$——无限小数本身要靠完备性才能定义清楚；
  \item 「无理数 $=$ 无限不循环小数」只是无理数的表征，不能当作定义：要说明无限小数的全体恰好铺满完备数轴，绕不开分割或基本列。
\end{itemize}
这两条留作课后思考（见课后作业拓展探究）。

\begin{conccard}[一级]{实数的三种等价面孔：分割、基本列与确界}{判断一个集合能否承载函数的单调性、零点定理与极限}{把「实数 $=$ 无限小数」当定义；把「稠密」误当「连续」}
实数的三种定义（戴德金分割、柯西基本列、确界公理）彼此等价，共同特征是「没有孔洞」。
$\mathbb{Q}$ 的致命伤不在稀疏而在漏点：$S=\{x\in\mathbb{Q}:x^2<2\}$ 有上界却无最小上界，任何有理数都填不进 $\sqrt2$ 的卡口。
高中阶段不必背构造细节，但必须知道：极限、零点定理、最值定理的合法性来自完备性，而不只是来自「能比大小」。
\end{conccard}

\begin{commonerror}{把「稠密」当成「连续」}
常见错误：认为「任意两个有理数之间还有有理数，所以有理数轴已经连续」。稠密只保证「找不到相邻两点」，不保证「不漏点」：$\sqrt2,\pi,e$ 都不在 $\mathbb{Q}$ 里，曲线 $y=x^2-2$ 照样从孔洞中穿过。
判断标准：能否说「有上界的非空集合必有最小上界」。$\mathbb{Q}$ 不能，$\mathbb{R}$ 能。
\end{commonerror}

\begin{method}
「看集合选武器」再补一条反射：遇到「变号却无零点／有界却无最值／递增却撞不到天花板」，先怀疑定义域是 $\mathbb{Q}$ 而非 $\mathbb{R}$；遇到「长度、面积、开方、$\pi$」，先确认舞台是 $\mathbb{R}$。
\end{method}

% ---------- 以上为背景铺垫，以下进入主线 ----------

初中我们说"$y$ 随 $x$ 的变化而变化"，高中要说得更精确：函数是一种特殊的对应，它的图象是平面上的点集。本导学案用一条主线重新组织"函数"概念——
\key{关系}（$R\subseteq A\times B$ 的三档舞台：离散$\times$离散的孤立点、离散$\times$连续的竖直线段族、连续$\times$连续的二维区域）$\xrightarrow{\text{方程／不等式＝条件的表达法}}$ \key{两条限制轴}（轴一·单值性 $\to$ 映射与函数；轴二·连续性与结构保持 $\to$ 基本初等函数），并在每一步用"\key{代数有序数对集合 $\leftrightarrow$ 几何坐标系点集 $\leftrightarrow$ 映射箭头图}"三种语言互相翻译。

\begin{learnobjective}
\begin{itemize}
  \item 从\key{点集视角}理解二元关系与函数：把对应看成 $\mathbb{R}^2$ 平面上的点集，说清"哪些点被选中"；
  \item 理解\key{单值性}（垂线检验）如何把任意点集过滤成函数图象：每条竖直线 $x=a$ 与点集至多一个交点；
  \item 掌握\key{三视角互译}：代数有序数对集合 $\leftrightarrow$ 几何坐标系点集 $\leftrightarrow$ 映射箭头图；
  \item 理解\key{两条限制轴}：轴一「单值性」回答它是不是函数，轴二「连续性／结构保持」回答它好不好算，二者互不替代；
  \item 说明\key{方程与不等式是关系的条件表达法}，并在大一统框架下说出二元方程、函数、一元方程（不等式）的解分别是几维对象。
\end{itemize}
\end{learnobjective}

\begin{preview}
\subsection*{知识回顾}
\begin{itemize}
  \item 笛卡尔积：$A\times B=\{(a,b)\mid a\in A,\ b\in B\}$。若 $A=\{1,2\}$，$B=\{3\}$，则 $A\times B=$ \blank；
  \item 二元关系：$A\times B$ 的\key{任意子集} $R$ 都叫 $A$ 到 $B$ 的二元关系。若 $R=\{(1,3),(2,3)\}$，则 $R$ 中元素个数是 \blank；
  \item 不等式解集：$x^2-1\leq0$ 的解集是 \blank，它表示数轴上的 \blank；
  \item 平面直角坐标系中，点 $(x,y)$ 与有序数对 $(x,y)$ 是 \blank（一一对应／没有关系）。
\end{itemize}

\subsection*{课前速读}
阅读「函数的概念」一节，完成下列问题：
\begin{enumerate}
  \item 函数定义中的两个关键词是"\key{任意}"和"\key{唯一}"；
  \item 函数的三要素是 \blank、\blank、\blank；
  \item 下列对象中哪些是函数：(1) 圆的面积 $S$ 与半径 $r$；(2) 一个人的年龄与身高；(3) 平方对应 $x\mapsto x^2$。（答：\blank）
\end{enumerate}
\end{preview}

\begin{explore}

\subsection{有限离散关系：从 $A\times B$ 出发}

\key{探究一：把所有可能的配对都列出来}
设 $A=\{1,2,3\}$，$B=\{2,4,6\}$。
\begin{itemize}
  \item $A\times B$ 共有 $3\times3=9$ 个有序数对：
        $(1,2),(1,4),(1,6),(2,2),(2,4),(2,6),(3,2),(3,4),(3,6)$。
  \item 把它看成\key{网格点集}：在平面直角坐标系中描出这 $9$ 个点（图 \ref{fig:grid9}）。
  \item $A\times B$ 的任意子集 $R$ 都叫 $A$ 到 $B$ 的一个\key{二元关系}。例如 $R=\{(1,2),(2,4),(3,6)\}$ 只是"挑出其中几个配对"。
\end{itemize}

\begin{figure}[htbp]
  \centering
  \begin{tikzpicture}
    \begin{axis}[
      width=0.55\textwidth, height=6cm,
      xmin=0.5, xmax=3.5, ymin=1, ymax=7,
      xtick={1,2,3}, ytick={2,4,6},
      xlabel=$x$, ylabel=$y$,
      grid=major, grid style={dashed,gray!35},
      tick label style={font=\footnotesize},
      title={\small $A\times B$ 的 9 个有序数对（网格点）},
      title style={font=\small},
    ]
      \addplot[only marks, mark=*, mark size=2pt, main!70!black]
        coordinates {(1,2)(1,4)(1,6)(2,2)(2,4)(2,6)(3,2)(3,4)(3,6)};
    \end{axis}
  \end{tikzpicture}
  \caption{$A\times B$ 的点集视角：平面直角坐标系中的 $3\times3$ 网格点集。竖直方向看，每一列对应同一个 $x\in A$，各含 $3$ 个点（一对多，尚未通过垂线检验）；代数清单见正文，映射箭头图视角见图 \ref{fig:arrow}}
  \label{fig:grid9}
\end{figure}

\key{同一母体，三档舞台：}$R\subseteq A\times B$ 长什么样，首先取决于 $A,B$ 是哪种集合（即前文「舞台四级进化」装了多少装备）。三档并列如下：
\begin{itemize}
  \item \key{离散$\times$离散}（$A=\{1,2,3\},\ B=\{2,4,6\}$）：$R$ 只能是有限个孤立点——图 \ref{fig:grid9} 的 $9$ 点网格；
  \item \key{离散$\times$连续}（$A=\{1,2,3\},\ B=[0,4]$）：每个 $x$ 上方可以出现一整段竖直线段（点集是「三根竖条」，图 \ref{fig:three-tiers} 中），也可以只取有限个点——具体形态仍由所加的条件决定；
  \item \key{连续$\times$连续}（$A=B=[0,3]$）：$R$ 可以是一整块二维区域——图 \ref{fig:three-tiers} 右的矩形。
\end{itemize}
三档共用同一套三语言：代数写 $R=\{(x,y)\mid\text{条件}\}$，几何描点，箭头图看方向。\textbf{舞台决定点能密到什么程度，条件决定点被约束成什么形状}；二者缺一不可。特别地，$A\times B$（$9$ 点网格）本身也是一个二元关系——「全取」的关系；它代表\textbf{两个变元完全自由}的起点，而不是终点，后面的全部内容都在给这份自由逐级上锁。

\begin{figure}[htbp]
  \centering
  \begin{minipage}{0.31\textwidth}
    \centering
    \begin{tikzpicture}[x=0.55cm,y=0.55cm,>=stealth,font=\scriptsize]
      \draw[->,gray] (0.5,0)--(3.6,0) node[right] {$x$};
      \draw[->,gray] (0,-0.2)--(0,4.6) node[above] {$y$};
      \foreach \x in {1,2,3}{ \foreach \y in {2,4,6}{ \fill[main!70!black] (\x,\y) circle (1.6pt); } }
      \node at (2,-0.8) {离散$\times$离散：$9$ 个孤立点};
    \end{tikzpicture}
  \end{minipage}\hfill
  \begin{minipage}{0.31\textwidth}
    \centering
    \begin{tikzpicture}[x=0.55cm,y=0.55cm,>=stealth,font=\scriptsize]
      \draw[->,gray] (0.5,0)--(3.6,0) node[right] {$x$};
      \draw[->,gray] (0,-0.2)--(0,4.6) node[above] {$y$};
      \foreach \x/\yl/\yu in {1/0.2/1.6, 2/0.6/3.4, 3/0/2.2}{
        \draw[second,very thick] (\x,\yl)--(\x,\yu);
        \fill[second] (\x,\yl) circle (1.6pt);
        \fill[second] (\x,\yu) circle (1.6pt);
      }
      \node at (2,-0.8) {离散$\times$连续：三根竖条};
    \end{tikzpicture}
  \end{minipage}\hfill
  \begin{minipage}{0.31\textwidth}
    \centering
    \begin{tikzpicture}[x=0.55cm,y=0.55cm,>=stealth,font=\scriptsize]
      \draw[->,gray] (0.5,0)--(3.6,0) node[right] {$x$};
      \draw[->,gray] (0,-0.2)--(0,4.6) node[above] {$y$};
      \fill[main!12] (0,0) rectangle (3,4);
      \draw[main!70!black,thick] (0,0) rectangle (3,4);
      \node at (2,-0.8) {连续$\times$连续：一整块区域};
    \end{tikzpicture}
  \end{minipage}
  \caption{关系 $R\subseteq A\times B$ 的三档形态：同一母体，因 $A,B$ 的装备不同，分别长成孤立点云、竖直线段族、二维区域。三档都可以同时用「有序数对集合／坐标点集／箭头图」三种语言描述}
  \label{fig:three-tiers}
\end{figure}

\key{追问：}把这三档点集都当作「从 $x$ 到 $y$ 的对应」，它们各有什么缺陷，迫使我们继续加限制条件？
（提示：$9$ 点网格中取 $x=1$，$y$ 可以是 $2,4,6$ 中任意一个——对应不唯一；竖条与矩形中一个 $x$ 拖着无穷多个 $y$；而且离散点集没有「连续变化」的资格，圆滑曲线、连续性都无从谈起。）

\begin{example}
写出 $A=\{1,2,3\}$、$B=\{2,4,6\}$ 时 $A\times B$ 的全部 $9$ 个有序数对并在网格中描点；再任取两个不同的二元关系 $R_1,R_2\subseteq A\times B$（各含 $2$ 个元素）。
\end{example}
\begin{solution}
$9$ 个有序数对为 $(1,2),(1,4),(1,6),(2,2),(2,4),(2,6),(3,2),(3,4),(3,6)$，描点见图 \ref{fig:grid9}。
任取 $R_1=\{(1,2),(2,6)\}$，$R_2=\{(1,4),(3,2)\}$ 即可（答案不唯一）。它们都是 $A\times B$ 的子集，所以都是二元关系。
\end{solution}

\key{探究二：单值性——每列至多一个点}
对比两个关系：
\[
  R_1=\{(1,2),(1,4),(2,6)\},\qquad R_2=\{(1,2),(2,4),(3,6)\}.
\]
用三视角分别表示：
\begin{itemize}
  \item \key{代数}（有序数对集合）：如上，$R_1$ 中 $x=1$ 出现了两次；
  \item \key{几何}（点集）：在网格中描点（图 \ref{fig:r1r2}）。$R_1$ 在竖直线 $x=1$ 上有两个点 $(1,2),(1,4)$；$R_2$ 在每条竖直线 $x=1,2,3$ 上各有一个点；
  \item \key{映射}（箭头图）：图 \ref{fig:arrow}。$R_1$ 中 $1\to2$ 与 $1\to4$ 同时出现（\key{一对多}）；$R_2$ 中 $1\to2$、$2\to4$、$3\to6$（一对一）。
\end{itemize}

\begin{figure}[htbp]
  \centering
  \begin{minipage}{0.48\textwidth}
    \centering
    \begin{tikzpicture}
      \begin{axis}[
        width=\textwidth, height=5cm,
        xmin=0.5, xmax=3.5, ymin=1, ymax=7,
        xtick={1,2,3}, ytick={2,4,6},
        grid=major, grid style={dashed,gray!35},
        tick label style={font=\footnotesize},
        title={\small $R_1$：$x=1$ 处有两个点（一对多）},
        title style={font=\small},
      ]
        \addplot[only marks, mark=*, mark size=2pt, second]
          coordinates {(1,2)(1,4)(2,6)};
      \end{axis}
    \end{tikzpicture}
  \end{minipage}\hfill
  \begin{minipage}{0.48\textwidth}
    \centering
    \begin{tikzpicture}
      \begin{axis}[
        width=\textwidth, height=5cm,
        xmin=0.5, xmax=3.5, ymin=1, ymax=7,
        xtick={1,2,3}, ytick={2,4,6},
        grid=major, grid style={dashed,gray!35},
        tick label style={font=\footnotesize},
        title={\small $R_2$：每列恰有一个点（一对一）},
        title style={font=\small},
      ]
        \addplot[only marks, mark=*, mark size=2pt, main]
          coordinates {(1,2)(2,4)(3,6)};
      \end{axis}
    \end{tikzpicture}
  \end{minipage}
  \caption{$R_1$ 与 $R_2$ 的点集对比：竖直线 $x=1$ 与 $R_1$ 有两个交点，与 $R_2$ 只有一个交点}
  \label{fig:r1r2}
\end{figure}

\begin{figure}[htbp]
  \centering
  \begin{minipage}{0.48\textwidth}
    \centering
    \begin{tikzpicture}[>=Stealth, semithick]
      \node (a1) at (0,2.2) {$1$}; \node (a2) at (0,1.1) {$2$}; \node (a3) at (0,0) {$3$};
      \node (b2) at (3,2.2) {$2$}; \node (b4) at (3,1.1) {$4$}; \node (b6) at (3,0) {$6$};
      \draw[->] (a1)--(b2); \draw[->] (a1)--(b4); \draw[->] (a2)--(b6);
      \node[font=\footnotesize] at (0,2.8) {$A$}; \node[font=\footnotesize] at (3,2.8) {$B$};
    \end{tikzpicture}
    \\[2pt]{\footnotesize $R_1$：$1$ 同时指向 $2$ 和 $4$}
  \end{minipage}\hfill
  \begin{minipage}{0.48\textwidth}
    \centering
    \begin{tikzpicture}[>=Stealth, semithick]
      \node (a1) at (0,2.2) {$1$}; \node (a2) at (0,1.1) {$2$}; \node (a3) at (0,0) {$3$};
      \node (b2) at (3,2.2) {$2$}; \node (b4) at (3,1.1) {$4$}; \node (b6) at (3,0) {$6$};
      \draw[->] (a1)--(b2); \draw[->] (a2)--(b4); \draw[->] (a3)--(b6);
      \node[font=\footnotesize] at (0,2.8) {$A$}; \node[font=\footnotesize] at (3,2.8) {$B$};
    \end{tikzpicture}
    \\[2pt]{\footnotesize $R_2$：每个 $x$ 只指向一个 $y$}
  \end{minipage}
  \caption{映射箭头图：$R_1$ 一对多（不是函数），$R_2$ 一对一（是函数）}
  \label{fig:arrow}
\end{figure}

\key{判断规则（垂线检验雏形）：}点集 $P\subseteq\mathbb{R}^2$ 能作为某个函数的图象，当且仅当\key{每条竖直线 $x=a$ 与 $P$ 至多有一个交点}（"每列至多 $1$ 个点"）。据此 $R_2$ 是函数，$R_1$ 不是函数。

为什么 $R_2$ 是函数：对 $A$ 中每个 $x$，都有唯一 $y$ 与之对应；而 $R_1$ 中 $x=1$ 同时对应 $2$ 和 $4$，违反唯一性。

\key{把对应「翻译」成有序数对（图象的第三种身份）：}若对应关系为 $f$，把每个 $x$ 与它的像 $y=f(x)$ 并排写成 $(x,y)=(x,f(x))$。于是「哪些数对被选中」就是集合
\[
  \Gamma_f=\{(x,f(x))\mid x\in A\}.
\]
它读作映射，就是箭头图；读作有序数对集合，就是 $\Gamma_f$；坐标化到平面直角坐标系，就是函数的\key{图象}。三种说法在此合流——后文说「图象」，指的就是这个数对集合，不是额外多出来的对象。

\key{追问：}如果允许"一对多"，还能叫函数吗？"一对多"的对象在数学里是什么？
（提示：它仍是\key{关系}。例如二元方程 $x^2+y^2=4$ 的解集就是一对多的点集——这为探究三埋下伏笔。）

\begin{example}
用三视角（代数、几何、箭头图）表示关系 $R=\{(1,2),(2,4),(3,6)\}$，并判断它是否为函数。
\end{example}
\begin{solution}
代数：$R=\{(1,2),(2,4),(3,6)\}$。
几何：描点 $(1,2),(2,4),(3,6)$，三点在同一直线 $y=2x$ 上，且每条竖直线 $x=1,2,3$ 上各有一个点。
箭头图：$1\to2$，$2\to4$，$3\to6$。
是函数：每个 $x\in A$ 都有唯一 $y$ 与之对应（事实上 $y=2x$）。
\end{solution}

\begin{problem}
在 $A=\{1,2,3\}$、$B=\{2,4,6\}$ 的网格中，任取一个二元关系 $R$，使它不是函数，并说明理由；再任取一个函数关系，写出它的对应法则。
\end{problem}
\begin{solution}
例：$R=\{(1,2),(1,6),(2,4)\}$。$x=1$ 同时对应 $2$ 和 $6$，违反单值性，不是函数（竖直线 $x=1$ 与点集有两个交点）。
例：$R=\{(1,4),(2,2),(3,6)\}$ 是函数（每列至多一个点），对应法则可写为 $f(1)=4$，$f(2)=2$，$f(3)=6$。
\end{solution}

\key{从「完全自由」到「规则登场」：一条收紧链}

回到 $A=\{1,2,3\}$、$B=\{2,4,6\}$：$A\times B$ 的 $9$ 个配对是\emph{两个变元完全自由}时的全部候选。逐步加限制，自由度这样消失：
\begin{itemize}
  \item \key{二元关系}：对 $9$ 个配对逐个决定「取／不取」，共 $2^9=512$ 种——仍几乎无规律；
  \item \key{单值性}：每列至多取一个。每列有「不取，或取 $2,4,6$ 之一」共 $4$ 种选择，合计 $4^3=64$ 种；若要求 $A$ 中每个 $x$ 都有对应，则只剩 $3^3=27$ 种；
  \item \key{统一法则}：要求用同一条公式（结构保持等式）表达所有对应——$27$ 种里只剩极少数，例如 $(1,2),(2,4),(3,6)$ 可被 $y=2x$ 一网打尽。
\end{itemize}
函数不是外加的定义，而是这条收紧链上的一环：\textbf{自由消失多少，规则就登场多少}；基本初等函数位于链条顶端——规则收得最紧、结构保存得最好。

\subsection{连续二维区域：不等式解集的几何本质}

\key{先约定：方程与不等式是「从 $A\times B$ 里挑子集」的条件句}

含变元的等式或不等式，读法统一为「满足条件的有序数对全体」：
\begin{itemize}
  \item 一元方程 $x^2=4$：只有一个变元，解集是数轴上的点 $\{-2,2\}$；
  \item 二元方程 $x^2+y^2=4$：两个变元，解集是平面曲线（圆周）；
  \item 二元不等式 $x^2+y^2\leq4$：同一条件放松为「$\leq$」，解集从曲线「鼓成」闭圆盘。
\end{itemize}
所以方程、不等式不是与函数、映射并列的另一条线，而是\key{关系的具体表达法}：变元个数决定几何维度，等号／不等号决定边界还是区域。下面用这把尺子重新看连续点集。

\key{探究三：把 $9$ 个点换成无穷多个点}
把取值范围从离散集合推广到连续区间，考察两个平面点集：
\[
  R_3=\{(x,y)\mid x^2+y^2\leq4\},\qquad
  R_4=\{(x,y)\mid x+y\leq3,\ x\geq0,\ y\geq0\}.
\]
\begin{itemize}
  \item $R_3$ 是圆心在原点、半径为 $2$ 的\key{闭圆盘}（一块二维区域），其边界恰是二元方程 $x^2+y^2=4$，内部由不等式 $x^2+y^2\leq4$ 给出——「方程给边界、不等式给区域」的第一个实例；
  \item $R_4$ 是由 $x$ 轴、$y$ 轴与直线 $x+y=3$ 围成的\key{三角形区域}（图 \ref{fig:region}）。
  \item 它们都是 $\mathbb{R}^2$ 中"一块黑乎乎的区域"，而不是孤立的点。
\end{itemize}

关键观察：取 $x=1$，在 $R_3$ 中 $y\in[-\sqrt3,\sqrt3]$，有无穷多个取值；在 $R_4$ 中 $y\in[0,2]$，也有无穷多个取值。一个 $x$ 对应无穷多个 $y$——\key{失控}，无法由 $x$ 唯一预测 $y$。

\key{追问：}二元不等式（组）的解集在几何上是什么？（答：$\mathbb{R}^2$ 中的二维区域。）这正是未来\key{线性规划可行域}的伏笔：满足一组约束的 $(x,y)$ 全体就是平面区域，而最优解往往出现在区域的顶点或边界上。

\begin{figure}[htbp]
  \centering
  \begin{minipage}{0.48\textwidth}
    \centering
    \begin{tikzpicture}[scale=1.1]
      \draw[->] (-2.4,0)--(2.4,0) node[right]{$x$};
      \draw[->] (0,-2.4)--(0,2.4) node[above]{$y$};
      \fill[main!15] (0,0) circle (2cm);
      \draw[main!70!black, thick] (0,0) circle (2cm);
      \draw[dashed, second, thick] (1,-1.732)--(1,1.732)
        node[above right, font=\footnotesize]{$x=1:y\in[-\sqrt3,\sqrt3]$};
      \fill[second] (1,1.732) circle (1.5pt);
      \fill[second] (1,-1.732) circle (1.5pt);
    \end{tikzpicture}
    \\[2pt]{\footnotesize $R_3$：闭圆盘，竖直线 $x=1$ 与它交于无穷多点}
  \end{minipage}\hfill
  \begin{minipage}{0.48\textwidth}
    \centering
    \begin{tikzpicture}[scale=1.1]
      \draw[->] (-0.5,0)--(3.6,0) node[right]{$x$};
      \draw[->] (0,-0.5)--(0,3.6) node[above]{$y$};
      \fill[main!15] (0,0)--(3,0)--(0,3)--cycle;
      \draw[main!70!black, thick] (0,0)--(3,0)--(0,3)--cycle;
      \draw[dashed, second, thick] (1,0)--(1,2)
        node[right, font=\footnotesize]{$x=1:y\in[0,2]$};
    \end{tikzpicture}
    \\[2pt]{\footnotesize $R_4$：三角形区域，竖直线 $x=1$ 与它交于无穷多点}
  \end{minipage}
  \caption{连续二维区域：二元不等式（组）的解集是 $\mathbb{R}^2$ 中的"面"，不是曲线}
  \label{fig:region}
\end{figure}

\begin{example}
在坐标系中画出 $R_4=\{(x,y)\mid x+y\leq3,\ x\geq0,\ y\geq0\}$，并求当 $x=1$ 时 $y$ 的取值范围。
\end{example}
\begin{solution}
画出直线 $x+y=3$ 与两条坐标轴，取左下方的三角形内部（含边界），见图 \ref{fig:region} 右图。
当 $x=1$ 时，条件化为 $1+y\leq3$ 且 $y\geq0$，即 $0\leq y\leq2$。$y$ 有无穷多个取值，无法由 $x=1$ 唯一确定——所以 $R_4$ 不是函数。
\end{solution}

\begin{problem}
设 $R=\{(x,y)\mid x^2+y^2\leq4\}$。判断 $R$ 是否为函数，并用垂线检验说明理由。
\end{problem}
\begin{solution}
不是函数。取 $x=0$，则 $y\in[-2,2]$，一个 $x$ 对应无穷多个 $y$；几何上看，竖直线 $x=0$ 与圆盘相交于无穷多个点，违反"每列至多一个点"。
\end{solution}

\subsection{两条限制轴：从区域到基本初等函数}

\key{探究四：单值性与正则性，是两条互相独立的轴}

任意点集（撒点）常常同时「一对多」又「不连续」。要把它变成能研究的函数，需要在\textbf{两条不同的轴}上分别加限制，二者不可互相替代：

\begin{itemize}
  \item \key{轴一·代数／规则——单值性。}\key{回扣：}这正是前文那条收紧链——从 $A\times B$ 的完全自由（$2^9$ 种取法），到单值（$4^3$；若处处有定义则 $3^3=27$），再到统一法则。具体做法：规定每条竖直线 $x=a$ 上\textbf{至多}保留一个点，并且\textbf{明确指定}保留哪一个（不能随手挑）；未指定的 $x$ 允许没有对应点——这些被指定的 $x$ 全体就是定义域 $A$。满足单值性的关系叫\key{映射} $f:A\to B$；当 $A,B$ 取为数集时，就是\key{函数} $y=f(x)$。判别标准只有一条：垂线检验（图 \ref{fig:orderline} 第 (3) 幅）。再要求每个对应都能用同一条解析式表达，才得到「写得出来的函数」。
  \item \key{轴二·几何／分析——正则性与结构保持。}单值性\emph{不附带}任何光滑承诺：轴一走到尽头，图象照样可以处处间断（开场三问的 Dirichlet 函数），也可以连续却处处不可微。要让它「越来越好算」，必须继续加条件：连续性 $\to$ 可微性 $\to$ 某类「运算封闭／结构保持」的等式，最终落在\key{基本初等函数}上：
  \begin{itemize}
    \item 指数函数 $y=a^x$：把定义域的加法变成值域的乘法，$f(x_1+x_2)=f(x_1)f(x_2)$；
    \item 对数函数 $y=\log_a x$：把指数函数的数对 $(x,y)$ 调换为 $(y,x)$ 的逆关系，图象关于直线 $y=x$ 对称；
    \item 幂函数 $y=x^{\alpha}$：刻画幂次结构（面积、体积、增长阶）；
    \item 三角函数：刻画周期结构（旋转、振动、波动）。
  \end{itemize}
\end{itemize}

\key{两条轴的分工：}轴一回答「它是不是函数」，轴二回答「它好不好算」。而方程与不等式（见本节开头的约定）是描述点集条件的语言，在任何一档舞台、任何一条轴上都可用：用二元方程 $F(x,y)=0$ 描述轴二下的平面曲线（可以一对多，如圆 $x^2+y^2=4$）；用一元方程 $f(x)=g(x)$ 求两条函数图象交点的横坐标；用一元不等式 $f(x)>g(x)$ 表示「$f$ 的点集位于 $g$ 的点集上方」的 $x$ 的范围。

\begin{figure}[htbp]
  \centering
  \begin{minipage}{0.235\textwidth}
    \centering
    \begin{tikzpicture}[x=0.6cm,y=0.6cm]
      \draw[main!60!black,thick,fill=main!4] (-2.2,-2) rectangle (2.2,2);
      \pgfmathsetseed{11}
      \foreach \i in {1,...,140}{
        \pgfmathsetmacro{\px}{-2.05+4.1*rnd}
        \pgfmathsetmacro{\py}{-1.85+3.7*rnd}
        \fill[main!45] (\px,\py) circle (0.8pt);
      }
    \end{tikzpicture}
    \\[2pt]{\footnotesize (1) 无序：随手撒点——一个 $x$ 拖着一大段 $y$}
  \end{minipage}\hfill
  \begin{minipage}{0.235\textwidth}
    \centering
    \begin{tikzpicture}[x=0.6cm,y=0.6cm]
      \draw[main!60!black,thick,fill=main!4] (-2.2,-2) rectangle (2.2,2);
      \pgfmathsetseed{11}
      \foreach \i in {1,...,140}{
        \pgfmathsetmacro{\px}{-2.05+4.1*rnd}
        \pgfmathsetmacro{\py}{-1.85+3.7*rnd}
        \fill[main!22] (\px,\py) circle (0.8pt);
      }
      \draw[second,very thick] (-0.6,-2)--(-0.6,2);
      \foreach \yy in {-1.75,-1.3,-0.85,-0.4,0.05,0.5,0.95,1.4,1.8}
        \fill[second] (-0.6,\yy) circle (1.5pt);
      \node[second,font=\scriptsize,right] at (-0.52,1.82) {$x=a$};
    \end{tikzpicture}
    \\[2pt]{\footnotesize (2) 竖截：竖线 $x=a$ 只截一段，仍是一整列点}
  \end{minipage}\hfill
  \begin{minipage}{0.235\textwidth}
    \centering
    \begin{tikzpicture}[x=0.6cm,y=0.6cm]
      \draw[main!60!black,thick,fill=main!4] (-2.2,-2) rectangle (2.2,2);
      \draw[third!55,thin] plot coordinates {(-2.05,1.82)(-1.89,1.54)(-1.72,1.15)(-1.56,-0.49)(-1.39,0.13)(-1.23,-0.37)(-1.07,-0.38)(-0.9,-1.22)(-0.74,-1.25)(-0.57,-0.93)(-0.41,-1.39)(-0.25,-1.05)(-0.08,-1.8)(0.08,-1.7)(0.25,-0.61)(0.41,-1.05)(0.57,-0.81)(0.74,-0.33)(0.9,-0.83)(1.07,-1.18)(1.23,0.24)(1.39,0.12)(1.56,0.48)(1.72,0.69)(1.89,1.58)(2.05,0.87)};
      \foreach \p in {(-2.05,1.82),(-1.89,1.54),(-1.72,1.15),(-1.56,-0.49),(-1.39,0.13),(-1.23,-0.37),(-1.07,-0.38),(-0.9,-1.22),(-0.74,-1.25),(-0.57,-0.93),(-0.41,-1.39),(-0.25,-1.05),(-0.08,-1.8),(0.08,-1.7),(0.25,-0.61),(0.41,-1.05),(0.57,-0.81),(0.74,-0.33),(0.9,-0.83),(1.07,-1.18),(1.23,0.24),(1.39,0.12),(1.56,0.48),(1.72,0.69),(1.89,1.58),(2.05,0.87)}
        \fill[third] \p circle (1.5pt);
    \end{tikzpicture}
    \\[2pt]{\footnotesize (3) 每列一点：被钉住了，但点列仍在乱跳}
  \end{minipage}\hfill
  \begin{minipage}{0.235\textwidth}
    \centering
    \begin{tikzpicture}[x=0.6cm,y=0.6cm]
      \draw[main!60!black,thick,fill=main!4] (-2.2,-2) rectangle (2.2,2);
      \pgfmathsetseed{37}
      \foreach \i in {1,...,90}{
        \pgfmathsetmacro{\bx}{-2.05+4.1*(\i-1)/89}
        \pgfmathsetmacro{\by}{0.5*\bx*\bx-1+0.55*(2*rnd-1)}
        \fill[main!40] (\bx,\by) circle (0.7pt);
      }
      \draw[second,very thick,domain=-2.05:2.05,samples=60,smooth] plot(\x,{0.5*\x*\x-1});
      \node[second,font=\scriptsize] at (0,-1.78) {$y=\tfrac12x^2-1$};
    \end{tikzpicture}
    \\[2pt]{\footnotesize (4) 收束：粗带压扁成唯一的曲线}
  \end{minipage}
  \caption{图像线是怎样从点集里长出来的：同一块平面区域上的四步图像实验。(1) 无序——随手撒点，一个 $x$ 的正上方压着一大段 $y$，一对多；(2) 竖截——用竖直线 $x=a$ 去截点云，截出的仍是一整列（一段）点，一个 $x$ 对应多个 $y$；(3) 立规矩——每条竖线只留一个点，得到一般函数的图象，但点列仍无规律地上下跳动；(4) 收束——点越密，跳动的点列糊成一条有厚度的粗带，厚度趋于 $0$ 时只剩唯一的曲线 $y=\tfrac12x^2-1$。图象就是「按 $x$ 排好队、且每个 $x$ 只剩唯一 $y$」的点列极限}
  \label{fig:orderline}
\end{figure}

\begin{conccard}[一级]{单值性与垂线检验}{任意平面点集（或图象）}{把"每列至多一个点"误写成"每行至多一个点"；忽略定义域}
点集 $P\subseteq\mathbb{R}^2$ 是某函数 $y=f(x)$ 的图象，当且仅当每条竖直线 $x=a$ 与 $P$ 至多有一个交点。代数表述：对每个 $x$，满足 $(x,y)\in P$ 的 $y$ 至多一个。这条判据把"任意点集"过滤成"函数图象"，是探究二、探究三共用的核心工具。
\end{conccard}

\begin{conccard}[二级]{三视角互译}{二元关系、函数、方程的解集}{只在一个视角里思考；把"数对集合"与"点集"割裂}
同一对象有三种等价表示：代数（有序数对集合）$\leftrightarrow$ 几何（坐标系点集）$\leftrightarrow$ 映射（箭头图）。代数精确、几何直观、箭头图体现对应方向。解题时按需要在三种语言间切换：例如解 $f(x)>g(x)$，就是看 $f$ 的点集在 $g$ 的点集上方时 $x$ 的范围。
\end{conccard}

\begin{method}
三视角互译法：遇到函数、方程、不等式问题时，先写出集合（代数），再画图（几何），必要时画箭头图（映射），最后回到代数表达。例如判断 $f(x)>g(x)$ 的解：画两条曲线，看 $f$ 在 $g$ 上方的 $x$ 区间，再写出集合 $\{x\mid f(x)>g(x)\}$。
\end{method}

\begin{commonerror}{把二维区域误当函数}
误以为"式子中有 $x,y$ 就是函数"或"给出 $x$ 就能算出 $y$"。反例：$x^2+y^2=4$ 中 $x=0$ 对应 $y=\pm2$，是一对多，不是函数；$\{(x,y)\mid x+y\leq3\}$ 中 $x=1$ 对应无穷多个 $y$。判别时必须回到单值性：每个 $x$ 是否唯一确定 $y$。
\end{commonerror}

\begin{example}
说明方程 $x^2=4$ 的解可以看成两个函数图象交点的横坐标，并求出解。
\end{example}
\begin{solution}
$x^2=4$ 即 $x^2-4=0$。令 $f(x)=x^2-4$，$g(x)=0$，交点横坐标即 $f(x)=g(x)$ 的解：$x=\pm2$。
也可看作曲线 $y=x^2$ 与直线 $y=4$ 交点的横坐标，结果同样是 $x=\pm2$。
\end{solution}

\begin{problem}
设 $f(x)=x^2$，$g(x)=2x+3$。求不等式 $f(x)>g(x)$ 的解集，并指出它对应"哪段 $x$ 上 $f$ 的图象在 $g$ 上方"。
\end{problem}
\begin{solution}
$x^2>2x+3\iff x^2-2x-3>0\iff(x-3)(x+1)>0$，解集为 $(-\infty,-1)\cup(3,+\infty)$。
几何意义：在 $(-\infty,-1)$ 与 $(3,+\infty)$ 上，抛物线 $y=x^2$ 位于直线 $y=2x+3$ 的上方；在 $[-1,3]$ 上则位于下方（或相交）。
\end{solution}

\subsection{大一统结构对照}

把本节对象放进同一张表：它们都是「集合」，并且位于同一条包含链上——关系 $\supseteq$ 映射 $\supseteq$ 函数 $\supseteq$ 基本初等函数；方程与不等式则是按变元个数给关系分类的条件表达法。

\begin{center}
\begin{tabular}{|c|p{3.0cm}|p{5.2cm}|p{4.2cm}|}
\hline
\textbf{数学对象} & \textbf{集合与点集本质} & \textbf{几何维度直观理解} \\
\hline
二元关系 $R\subseteq A\times B$ & $A\times B$ 的任意子集（条件句挑出） & 点云：孤立点／竖直线段族／区域／曲线皆可 \\
\hline
映射 $f:A\to B$ & 每列至多一个点的单值关系 & 箭头图；数集情形为一维曲线 \\
\hline
二元不等式（组） & $\{(x,y)\mid \text{约束条件}\}$ & $\mathbb{R}^2$ 中的二维区域（"面"） \\
\hline
二元方程 $F(x,y)=0$ & $\{(x,y)\mid F(x,y)=0\}$ & 平面上的一维曲线（可一对多） \\
\hline
函数 $y=f(x)$ & $\{(x,f(x))\mid x\in A\}$ & 一维曲线，且每列至多一个点 \\
\hline
一元方程的解 & $\{x\mid f(x)=g(x)\}$ & 数轴上的零维点（交点横坐标） \\
\hline
一元不等式的解 & $\{x\mid f(x)>g(x)\}$ & 数轴上的一维区间（上下位置） \\
\hline
\end{tabular}
\end{center}

\begin{summary}
本节主线：关系 $R\subseteq A\times B$ 有三档舞台——离散$\times$离散的孤立点、离散$\times$连续的竖直线段族、连续$\times$连续的二维区域；方程与不等式是把关系写成条件句的表达法（变元个数定维度，等号／不等号定边界与区域）。在关系之上加限制要分两条独立的轴：轴一加单值性，得到映射与函数（判别标准只有垂线检验）；轴二加连续性、可微性与结构保持等式，得到基本初等函数（指数、对数、幂、三角）。三种语言（代数、几何、箭头图）全程互译。
\end{summary}

\subsection{函数与图形变换完全图谱}

\begin{conccard}[一级]{第一性原理：一切图像变换都是有序数对 $(x,y)$ 的映射}{抽象函数图像辨析、变换类压轴题}{只背“左加右减”口诀，不追踪每个点的去向}
函数图像是 $\mathbb{R}^2$ 中的点集 $\{(x,y)\mid y=f(x)\}$。所谓“图像变换”，本质是对集合中每个点的坐标施加映射。
掌握下面三级体系与对称周期生成律，就能把“看图识函数”变成可推导的代数操作。
\end{conccard}

\begin{center}
\renewcommand{\arraystretch}{1.4}
\begin{tabular}{|p{3.2cm}|p{4.0cm}|p{4.8cm}|}
\hline
\textbf{变换} & \textbf{点集规则} & \textbf{表达式与几何效应} \\
\hline
\multicolumn{3}{|l|}{\textbf{一级：刚体等距（形状大小不变）}} \\
\hline
平移 & $(x,y)\mapsto(x+h,\,y+k)$ & $y-k=f(x-h)$（逆向代换） \\
\cline{2-3}
轴对称（$x=a$） & $(x,y)\mapsto(2a-x,\,y)$ & $y=f(2a-x)$ \\
\cline{2-3}
中心对称（$(a,b)$） & $(x,y)\mapsto(2a-x,\,2b-y)$ & $y=2b-f(2a-x)$（$a=b=0$ 即奇函数） \\
\hline
\multicolumn{3}{|l|}{\textbf{二级：伸缩与翻折}} \\
\hline
轴向伸缩 & $(x,y)\mapsto(\frac1\omega x,\,Ay)$ & $y=Af(\omega x)$ \\
\cline{2-3}
绝对值翻折 & $(x,y)\mapsto(x,\,|y|)$；$x\mapsto|x|$ & $y=|f(x)|$（下方翻上）；$y=f(|x|)$（右侧复制到左侧，强制偶化） \\
\hline
\multicolumn{3}{|l|}{\textbf{三级：非线性算子（图像拓扑突变）}} \\
\hline
自变量平方／倒数 & $u=g(x)$ 作为内层 & $y=f(x^2)$ 必为偶函数且在 $x=0$ 处切线水平 \\
\cline{2-3}
函数值平方／倒数 & $(x,y)\mapsto(x,\,y^2)$；$(x,y)\mapsto(x,\,\frac1y)$ & 零点变极小值点／垂直渐近线（详见下文铁律） \\
\hline
\end{tabular}
\end{center}

\begin{method}
\textbf{双对称推周期：连续两次平行镜像反射 = 平移}。看到“两个对称条件”或“半周期反对称”，直接锁定最小周期，不必画图：
\begin{enumerate}
\item \textbf{双轴对称}：$f(a+x)=f(a-x)$ 且 $f(b+x)=f(b-x)$ $\Rightarrow$ $T=2|b-a|$；
\item \textbf{双中心对称}：$f(a+x)+f(a-x)=2c$ 且 $f(b+x)+f(b-x)=2c$ $\Rightarrow$ $T=2|b-a|$；
\item \textbf{一轴一中心}（最高频）：$f(a+x)=f(a-x)$ 且 $f(b+x)+f(b-x)=2c$ $\Rightarrow$ $T=4|b-a|$；
\item \textbf{滑移反对称}：$f(x+a)=-f(x)$、$f(x+a)=\pm\frac1{f(x)}$ 或 $f(x+a)=\frac{1-f(x)}{1+f(x)}$，迭代两次即复原 $\Rightarrow$ $T=2a$。
\end{enumerate}
几何直观：点 $P(x)$ 先关于 $x=a$ 反射到 $P_1(2a-x)$，再关于 $x=b$ 反射到 $P_2\bigl(2b-(2a-x)\bigr)=P(x+2(b-a))$——两次翻转方向复原，合成一次纯平移。
\end{method}

\begin{center}
\begin{tikzpicture}[scale=1.0]
\draw[->, gray, thick] (-1.2,0) -- (7.2,0) node[right] {$x$};
\draw[main, dashed, very thick] (1.5,-0.6) -- (1.5,2.4) node[above] {$x=a$};
\draw[second, dashed, very thick] (4.0,-0.6) -- (4.0,2.4) node[above] {$x=b$};
\fill[black] (0.5,1.5) circle (2pt) node[above] {$P(x)$};
\fill[main] (2.5,1.5) circle (2pt) node[below] {$P_1(2a-x)$};
\fill[second] (5.5,1.5) circle (2pt) node[below] {$P_2\bigl(x+2(b-a)\bigr)$};
\draw[<->, gray, thick] (0.5,-0.35) -- (5.5,-0.35) node[midway, below, font=\small\bfseries] {合成平移 $T=2(b-a)$};
\end{tikzpicture}
\end{center}

\begin{conccard}[二级]{导数对称遗传律：周期全遗传，对称奇偶互换}{含参导数压轴、切线分析}{把 $f$ 与 $f'$ 的对称性混为一谈}
对 $f$ 的恒等式两边求导（链式法则）即得：
\begin{itemize}
\item $f$ 周期为 $T$ $\iff$ $f'$ 周期仍为 $T$；
\item $f$ 关于直线 $x=a$ 轴对称 $\Rightarrow$ $f'$ 关于点 $(a,0)$ 中心对称，且 $f'(a)=0$；
\item $f$ 关于点 $(a,b)$ 中心对称 $\Rightarrow$ $f'$ 关于直线 $x=a$ 轴对称。
\end{itemize}
\end{conccard}

\begin{method}
\textbf{非线性算子的五条铁律}（以 $g(x)=\dfrac1{f(x)}$ 为例，$y=[f(x)]^2$ 见表中）：
\begin{enumerate}
\item \key{符号不变}：$f>0$ 处 $g>0$，$f<0$ 处 $g<0$，上下位置不动；
\item \key{零点撕裂}：$f(x_0)=0$ 变成垂直渐近线 $x=x_0$；
\item \key{无穷变零}：$f\to\infty$ 处 $g\to0$（趋近 $x$ 轴）；
\item \key{单调反转}：同一同号区间内 $g$ 与 $f$ 单调性相反（不能跨渐近线谈单调）；
\item \key{极值互换}：正的极大值 $M$ 变成极小值 $\frac1M$；而 $y=\pm1$ 的点是\key{不动点}，位置纹丝不动。
\end{enumerate}
$y=f(x^2)$ 与 $y=f\!\left(\frac1x\right)$ 则作用在自变量上：$f(x^2)$ 是偶函数，
左半支与右半支互为镜像（同一 $|x|$ 处函数值相同）；在 $x\ge0$ 上它等于把
$y=f(u)\,(u\ge0)$ 的横坐标按 $u=x^2$ 重标得到的曲线——靠近原点被压缩、
远离原点被拉伸；$f\!\left(\frac1x\right)$ 则把 $x\to0$ 的行为与
$x\to\pm\infty$ 的行为\key{对调}。

\textbf{平方算子 $y=[f(x)]^2$ 的五条铁律}（与倒数算子对照记忆）：
\begin{itemize}
  \item \key{非负化}：整个图像被压进 $y\ge0$，负部沿 $x$ 轴\key{翻折}上来；
  \item \key{零点不变}：$f(x_0)=0\iff f^2(x_0)=0$，只是零点的「重数」加倍；
  \item \key{无穷仍无穷}：$f\to\infty$ 时 $f^2\to+\infty$ 增长更快；
  \item \key{单调分区}：由 $(f^2)'=2ff'$，$f>0$ 处 $f^2$ 与 $f'$ 同号、
        $f<0$ 处反号，故在 $f$ 的正区间保持单调方向、负区间反转；
  \item \key{极值缩放}：$|f|$ 的极大值 $M$ 变成 $f^2$ 的极大值 $M^2$
        （$M\ge1$ 时变大、$0<M<1$ 时变小），$|f|$ 的小极小值则被压得更小；
        $y=1$ 上的点是\key{不动点}。
\end{itemize}
\end{method}

\begin{center}
\begin{tikzpicture}[scale=1.15]
\draw[->, gray, thick] (-1.8,0) -- (3.8,0) node[right] {$x$};
\draw[->, gray, thick] (0,-2.2) -- (0,3.0) node[above] {$y$};
\draw[gray, dashed, thick] (2,-2.2) -- (2,3.0) node[above] {$x=2$};
\draw[main, thick, domain=-1.4:3.4, samples=60] plot(\x, {(\x)*(\x-2)});
\node[main] at (3.1,1.6) {$f(x)=x^2-2x$};
\draw[second, thick, domain=-1.5:-0.12, samples=60] plot(\x, {1/((\x)*(\x-2))});
\draw[second, thick, domain=0.12:1.9, samples=80] plot(\x, {1/((\x)*(\x-2))});
\draw[second, thick, domain=2.12:3.5, samples=60] plot(\x, {1/((\x)*(\x-2))});
\fill[second] (1,-1) circle (2pt) node[below left] {$(1,-1)$};
\end{tikzpicture}
{\small\itshape 抛物线与其倒数：零点 $0,2$ 变为竖直渐近线，极小值 $-1$ 变为极大值 $-1$。}
\end{center}

\begin{example}[双对称推周期与导数遗传]
已知 $f(x)$ 在 $\mathbb{R}$ 上可导，$f(x+1)$ 为偶函数，$f(x+2)$ 为奇函数，且 $f(1)=2$，求 $f(2027)+f'(2025)$。
\begin{solution}
$f(x+1)$ 偶 $\Rightarrow f(1-x)=f(1+x)$，故 $f$ 关于 $x=1$ 轴对称；
$f(x+2)$ 奇 $\Rightarrow f(2-x)=-f(2+x)$，故 $f$ 关于 $(2,0)$ 中心对称。
一轴一中心 $\Rightarrow T=4|2-1|=4$。
由中心对称取 $t=1$：$f(3)=-f(1)=-2$，于是 $f(2027)=f(4\times506+3)=f(3)=-2$。
$f$ 关于 $x=1$ 轴对称 $\Rightarrow f'$ 关于 $(1,0)$ 中心对称，故 $f'(1)=0$；又 $f'$ 周期仍为 $4$，
$f'(2025)=f'(4\times506+1)=f'(1)=0$。合计 $\key{-2}$。
\end{solution}
\end{example}

\begin{example}[用奇偶性速判函数图像]
函数 $g(x)=\dfrac{x}{e^x-1}$ 的大致图像如何？
\begin{solution}
先看定义域 $x\neq0$。由切线不等式 $e^x>x+1\ (x\neq0)$ 知 $e^x-1$ 与 $x$ 同号，故 $g(x)>0$ 恒成立（全在 $x$ 轴上方）；且 $\lim\limits_{x\to0}g(x)=1$。
当 $x\to+\infty$ 时指数增长远快于一次项，$g(x)\to0^+$；当 $x\to-\infty$ 时 $e^x\to0$，$g(x)\approx\dfrac{x}{-1}=-x\to+\infty$，即左端贴着斜率 $-1$ 的直线向上。无需二阶求导即可锁定图像。
\end{solution}
\end{example}

\begin{commonerror}{把“对称”当“周期”直接用}
$f(a+x)=f(a-x)$ 只说明关于 $x=a$ 对称，一个对称条件\key{产生不了周期}；必须有两个对称条件，
或用滑移反对称结构。此外 $y=f(x^2)$ 是偶函数，而 $y=[f(x)]^2$ 只保证非负——两者都“偶化”方向不同，不可互推。
\end{commonerror}

图形变换还有一种对偶记法，能把「左加右减」的口诀风险降到最低：动曲线与动坐标系是同一事实的两种记账方式（图 \ref{fig:activepassive}）。

\begin{methodbox}{主被动对偶律：曲线平移，等价于坐标系反向平移}
设曲线 $C:y=f(x)$ 主动向右平移 $h$ 个单位，得到新曲线 $y=f(x-h)$；若改为让坐标轴向左平移 $h$ 个单位、曲线在空间中保持不动，则新读数满足 $x'=x+h$、$y'=y$，同一物理状态在新系下的方程仍写作 $y'=f(x'-h)$。
两式描述的是同一幅终态图：顶点的读数相同，只是「谁在动」的归属不同。若条件按曲线平移给出（如「图象右移 $2$」），用主动记法；若条件按新原点、新读数给出，用被动记法。两者方向相反是必然，切勿把坐标轴的移动方向也套用「左加右减」。
\end{methodbox}

\begin{figure}[htbp]
  \centering
  \begin{tikzpicture}[>=stealth,font=\small]
    % ---------- 左：主动变换（坐标系不动，曲线右移） ----------
    \begin{scope}
      \draw[->,gray,thick] (-1.4,0)--(3.6,0) node[right,gray,font=\footnotesize] {$x$};
      \draw[->,gray,thick] (0,-0.7)--(0,3.0) node[above,gray,font=\footnotesize] {$y$};
      \draw[gray!70,dashed,thick,domain=-1.15:1.15,samples=50,smooth] plot(\x,{\x*\x});
      \draw[main,very thick,domain=0.85:3.15,samples=60,smooth] plot(\x,{(\x-2)*(\x-2)});
      \fill[gray!70] (0,0) circle (0.05);
      \node[gray!70,below left,font=\footnotesize] at (0,0) {$(0,0)$};
      \fill[main] (2,0) circle (0.05);
      \node[main,below right,font=\footnotesize] at (2,0) {$(2,0)$};
      % 平移箭头画在 x 轴下方空白处，避免压曲线
      \draw[<->,second,thick] (0.15,-0.35)--(1.85,-0.35)
        node[midway,above,font=\footnotesize] {右移 $2$};
      \node[main,font=\footnotesize] at (2.55,2.5) {$y=(x-2)^2$};
      \node[font=\footnotesize] at (1.1,-1.45) {主动：坐标系不动，曲线右移 $2$};
    \end{scope}
    % ---------- 右：被动变换（曲线不动，坐标轴左移） ----------
    \begin{scope}[xshift=7.7cm]
      \draw[gray!55,dashed,thick] (-1.4,0)--(2.2,0);
      \draw[gray!55,dashed,thick] (0,-0.7)--(0,3.0);
      \node[gray!55,font=\scriptsize,right] at (-1.35,2.55) {旧系（虚线）};
      \draw[->,second,thick] (-3.5,0)--(1.7,0) node[right,second,font=\footnotesize] {$x'$};
      \draw[->,second,thick] (-2,-0.7)--(-2,3.0) node[above,second,font=\footnotesize] {$y'$};
      \draw[main,very thick,domain=-1.15:1.15,samples=50,smooth] plot(\x,{\x*\x});
      \fill[main] (0,0) circle (0.05);
      \node[gray!60,font=\scriptsize,left] at (-0.06,0.28) {旧 $O$};
      \node[second,below left,font=\footnotesize] at (-2,0) {$O'$};
      \node[main,font=\footnotesize,right] at (0.12,-0.32) {顶点 $(2,0)$};
      \draw[<->,third,thick] (-2,-0.85)--(0,-0.85)
        node[midway,below,font=\footnotesize] {轴左移 $2$};
      \draw[third!70!black,->,shorten <=1pt,shorten >=1pt] (0.2,-0.68)--(0.06,-0.12);
      \node[third!70!black,font=\scriptsize,right] at (0.25,-0.72) {读数 $x'=2$};
      \node[main,font=\footnotesize] at (0.95,2.25) {$y'=(x'-2)^2$};
      \node[font=\footnotesize] at (-0.85,-1.45) {被动：曲线不动，坐标轴左移 $2$};
    \end{scope}
  \end{tikzpicture}
  \caption{主被动对偶：左图坐标系不动，抛物线 $y=x^2$ 主动右移 $2$ 得 $y=(x-2)^2$；右图抛物线原地不动，坐标轴被动左移 $2$（虚线为旧系），新读数 $x'=x+2$，顶点在新系下仍读作 $(2,0)$。两幅终态图完全相同——「动曲线」与「动坐标系」是同一平移的两种归属，不要混用方向}
  \label{fig:activepassive}
\end{figure}

% ---------- 收官微专题：$45^\circ$ 镜面反射定律 ----------

\begin{conccard}[二级]{唯有 $k=\pm1$ 的镜面，能保住「函数」的资格}{反函数对称、指对同构压轴选填、解析几何对称速算}{把斜线对称当成左右平移；漏掉 $b$ 的补偿项}
把对称轴从竖直线 $x=a$ 推广到一般直线 $y=kx+b$ 时，反射会让横、纵坐标\textbf{互相耦合}，反射后的曲线往往不再通过垂线检验，\key{跌出函数集合}，退回一般二元方程 $F(x,y)=0$。只有 $k=\pm1$（倾角 $45^\circ$ 或 $135^\circ$）时，耦合系数全部约尽，退化为干净的坐标互换——这正是反函数与指对同构压轴题的共同底座。
\end{conccard}

先把最典型的镜面 $l:y=x$ 看清楚（图 \ref{fig:expinvlog}）：反射把指数函数的点集 $\{(x,y)\}$ 整体翻成 $\{(y,x)\}$，对数曲线由此生成——这就是「反函数＝逆关系点集」的几何来源。

\begin{figure}[htbp]
  \centering
  \begin{tikzpicture}[x=0.9cm,y=0.9cm,>=stealth,font=\small]
    \draw[->,gray,thick] (-1.7,-1.7)--(3.6,-1.7) node[right] {$x$};
    \draw[->,gray,thick] (-1.7,-1.7)--(-1.7,3.8) node[above] {$y$};
    \draw[gray,dashed,thick,domain=-1.4:3.4,samples=40] plot(\x,\x) node[above right] {$y=x$};
    \draw[main,very thick,domain=-1.3:1.35,samples=80,smooth] plot(\x,{exp(\x)}) node[right] {$y=e^x$};
    \draw[second,very thick,domain=0.28:3.4,samples=90,smooth] plot(\x,{ln(\x)}) node[right] {$y=\ln x$};
    \fill[main] (0,1) circle (0.07);       \node[main,left] at (0,1) {$(0,1)$};
    \fill[main] (1,2.718) circle (0.07);   \node[main,above right] at (1,2.718) {$(1,e)$};
    \fill[main] (-1,0.368) circle (0.07);  \node[main,below left] at (-1,0.368) {$(-1,e^{-1})$};
    \fill[second] (1,0) circle (0.07);      \node[second,below] at (1,0) {$(1,0)$};
    \fill[second] (2.718,1) circle (0.07);  \node[second,right] at (2.718,1) {$(e,1)$};
    \fill[second] (0.368,-1) circle (0.07); \node[second,below left] at (0.368,-1) {$(e^{-1},-1)$};
    \draw[->,third,thick] (0,1) to[bend left=12] (1,0);
    \draw[->,third,thick] (1,2.718) to[bend left=18] (2.718,1);
    \draw[->,third,thick] (-1,0.368) to[bend right=12] (0.368,-1);
  \end{tikzpicture}
  \caption{同一坐标系中的互逆反射：$y=e^x$ 与 $y=\ln x$ 关于 $y=x$ 对称，每个点 $(a,b)$ 与它的像 $(b,a)$ 用跨越镜面的箭头一一相连，例如 $(0,1)\leftrightarrow(1,0)$、$(1,e)\leftrightarrow(e,1)$、$(-1,e^{-1})\leftrightarrow(e^{-1},-1)$。对数曲线就是指数曲线点集 $\{(x,y)\}$ 经 $y=x$ 翻转后的 $\{(y,x)\}$}
  \label{fig:expinvlog}
\end{figure}

\begin{method}
\textbf{$45^\circ$ 镜面反射替换速查表}（镜面记为 $l$，点 $P(x,y)$ 的像为 $P'(x',y')$）：
\begin{center}
\renewcommand{\arraystretch}{1.4}
\small
\begin{tabular}{|p{2.0cm}|p{3.6cm}|p{3.9cm}|p{4.2cm}|}
\hline
\textbf{镜面} & \textbf{点集映射} & \textbf{曲线替换口径} & \textbf{典型用途} \\
\hline
$l:y=x$ & $(x,y)\mapsto(y,\,x)$ & $F(y,\,x)=0$ & 反函数：$y=a^x$ 与 $y=\log_a x$ \\
\hline
$l:y=-x$ & $(x,y)\mapsto(-y,\,-x)$ & $F(-y,\,-x)=0$ & 双轴对称联动：$y=\dfrac1x$ \\
\hline
$l:y=x+b$ & $(x,y)\mapsto(y-b,\,x+b)$ & $F(y-b,\,x+b)=0$ & 平移型反函数、指对同构压轴 \\
\hline
$l:y=-x+b$ & $(x,y)\mapsto(b-y,\,b-x)$ & $F(b-y,\,b-x)=0$ & 斜率 $-1$ 直线族（$y=x$ 的垂线） \\
\hline
\end{tabular}
\end{center}
\vspace{0.3em}
\key{秒杀推论}：若两条曲线关于 $y=x$ 对称（如互为反函数），且被斜率为 $-1$ 的直线 $x+y=m$ 所截，则两交点 $A,B$ 必关于 $y=x$ 对称，中点 $\bigl(\tfrac m2,\tfrac m2\bigr)$ 落在 $y=x$ 上，于是
$$x_A+x_B=m,\qquad y_A+y_B=m .$$
\end{method}

\begin{figure}[htbp]
  \centering
  \begin{tikzpicture}[x=0.95cm,y=0.95cm,>=stealth,font=\small]
    \draw[->,gray,thick] (-1.2,-1.2)--(5.2,-1.2) node[right] {$x$};
    \draw[->,gray,thick] (-1.2,-1.2)--(-1.2,5.2) node[above] {$y$};
    \draw[main,dashed,very thick,domain=-1.1:3.9,samples=40] plot(\x,\x) node[above right] {$y=x$};
    \draw[second,very thick,domain=0.09:4.8,samples=90] plot(\x,{ln(\x)}) node[right] {$y=\ln x$};
    \draw[third,very thick,domain=-1:1.2,samples=70] plot(\x,{exp(\x)}) node[above] {$y=e^x$};
    \draw[red,very thick] (-1,5)--(5,-1) node[below right] {$x+y=4$};
    \fill[main] (2,2) circle (2.4pt) node[above left] {$M(2,2)$};
    \fill[red] (1.07,2.93) circle (2.4pt) node[above left] {$A(x_1,y_1)$};
    \fill[red] (2.93,1.07) circle (2.4pt) node[below right] {$B(y_1,x_1)$};
  \end{tikzpicture}
  \caption{反函数 $y=e^x$ 与 $y=\ln x$ 关于 $y=x$ 对称；斜率为 $-1$ 的直线 $x+y=4$ 与二者各交于 $A,B$，两点关于 $y=x$ 对称，中点 $M(2,2)$ 落在 $y=x$ 上，故 $x_1+x_2=4$}
  \label{fig:mirror45}
\end{figure}

\begin{example}[指对同构：不解方程，三步出答案]
设 $x_1$ 满足 $e^{x_1}+x_1=4$，$x_2$ 满足 $\ln x_2+x_2=4$，求 $x_1+x_2$。
\begin{solution}
两式分别改写为 $e^{x_1}=4-x_1$、$\ln x_2=4-x_2$。
于是 $A(x_1,e^{x_1})$ 是曲线 $y=e^x$ 与直线 $y=4-x$ 的交点，$B(x_2,\ln x_2)$ 是曲线 $y=\ln x$ 与同一直线的交点（图 \ref{fig:mirror45}）。
$y=e^x$ 与 $y=\ln x$ 互为反函数，关于 $y=x$ 对称；直线 $y=4-x$ 的斜率为 $-1$，与 $y=x$ 垂直，故 $A,B$ 关于 $y=x$ 对称，即 $x_2=e^{x_1}$。
所以 $x_1+x_2=x_1+e^{x_1}=\key{4}$。全程无需解超越方程。
\end{solution}
\end{example}

\begin{example}[平移型镜面：对称轴不是 $y=x$ 时]
求曲线 $C_1:y=e^x$ 关于直线 $l:y=x+1$ 对称的曲线 $C_2$ 的方程，并验证点 $P(1,e)$ 与其对称点 $P'$ 分别落在 $C_1,C_2$ 上。
\begin{solution}
用 $l:y=x+b$（$b=1$）的替换口径 $F(y-b,\,x+b)=0$，把 $C_1$ 的方程 $y=e^x$ 中的 $x$ 换成 $y-1$、$y$ 换成 $x+1$：
$$x+1=e^{\,y-1}\implies y=\ln(x+1)+1,$$
即 $C_2:y=\ln(x+1)+1$。
点 $P(1,e)$ 的对称点为 $P'=(y-b,\,x+b)=(e-1,\,2)$；代入 $C_2$：$\ln\bigl((e-1)+1\bigr)+1=\ln e+1=2$，恰好落在 $C_2$ 上。
（顺带一提：$x=0$ 时 $e^0=1=0+1$，两条曲线在镜面上的点 $(0,1)$ 处相遇。）
\end{solution}
\end{example}

\begin{commonerror}{把「$y=x+b$ 对称」误当成「左移 $b$」}
关于 $y=x+b$ 的对称\emph{不是}平移：它同时交换了横纵坐标的角色。若把 $C_1:y=e^x$ 关于 $y=x+1$ 的对称曲线误写成 $y=e^{x-1}$（左移思路），代入点 $P(1,e)$ 的对称点 $(e-1,2)$ 立即矛盾——正确的替换永远是 $x\leftrightarrow y-b$ 的\emph{互换}，而不是同向加减。
\end{commonerror}

\end{explore}

\begin{inclasstest}
\begin{enumerate}
  \item 用三视角（代数、几何、箭头图）表示关系 $R=\{(1,1),(2,4),(3,9)\}$，判断它是否为函数，并写出对应法则；\level{易}
  \item 判断点集 $P=\{(x,y)\mid x^2+y^2=1\}$ 是否为函数，用垂线检验说明理由；\level{易}
  \item 点集 $\{(x,y)\mid x^2+y^2\leq4\}$ 与 $\{(x,y)\mid x+y\leq3\}$ 的交集表示什么几何图形？它是否可能为某个函数的图象？说明理由；\level{中}
  \item 方程 $x^3=x$ 的解可看成哪两个函数图象交点的横坐标？求出所有解。\level{中}
\end{enumerate}
\end{inclasstest}

\begin{afterclass}
\subsection*{基础巩固}
\begin{enumerate}
  \item 写出 $A=\{1,2\}$、$B=\{3,4\}$ 的笛卡尔积 $A\times B$，并任取一个二元关系；
  \item 判断 $R=\{(1,2),(2,2),(3,2)\}$ 是否为函数，并说明理由。
\end{enumerate}

\subsection*{进阶提升}
\begin{enumerate}
  \item 画出 $R=\{(x,y)\mid y\geq x^2,\ x^2+y^2\leq4\}$ 表示的平面区域，并判断它是否为函数；
  \item 求不等式 $x^2<2x+3$ 的解集，并说明它对应两曲线怎样的位置关系。
\end{enumerate}

\subsection*{拓展探究}
\begin{enumerate}
  \item 查资料说明：为什么"对数能把乘除变成加减"？用 $y=a^x$ 与 $y=\log_a x$ 的互逆关系（$(x,y)\leftrightarrow(y,x)$）加以解释。
  \item （对应「实数的三种等价面孔」）设 $S=\{x\in\mathbb{Q}:x^2<2\}$。证明：$S$ 在 $\mathbb{Q}$ 中没有最小上界——要求对任意有理上界 $r$，构造出更小的有理上界；并由此说明曲线 $y=x^2-2$ 为什么能穿过有理数轴的孔洞。
  \item （表征与定义之别）某同学说：「$0.999\dots$ 和 $1$ 之间总差一点点，所以 $0.999\dots\neq1$。」请指出错在哪里，并用等比数列求和或方程法给出严格说明。
\end{enumerate}

\subsection*{教师版答案}
\begin{solution}
\textbf{基础 1：}$A\times B=\{(1,3),(1,4),(2,3),(2,4)\}$；任取一个子集即可，如 $R=\{(1,4),(2,3)\}$。
\end{solution}
\begin{solution}
\textbf{基础 2：}是函数。虽然三个 $x$ 都对应同一个 $y=2$，但每个 $x$ 对应的 $y$ 唯一（多对一是允许的），竖直线 $x=1,2,3$ 与点集各只有一个交点。
\end{solution}
\begin{solution}
\textbf{进阶 1：}区域是抛物线 $y=x^2$ 上方、圆盘 $x^2+y^2\leq4$ 内部（含边界）的公共部分，即圆盘被抛物线截出的"弓形"部分。它不是函数：取 $x=0$，$y\in[0,2]$，一个 $x$ 对应无穷多个 $y$（竖直线 $x=0$ 与区域有无穷多交点）。
\end{solution}
\begin{solution}
\textbf{进阶 2：}$x^2<2x+3\iff(x-3)(x+1)<0$，解集 $(-1,3)$。几何上，在 $(-1,3)$ 上抛物线 $y=x^2$ 位于直线 $y=2x+3$ 的下方。
\end{solution}
\begin{solution}
\textbf{拓展：}设 $x=a^{m}$，$y=a^{n}$，则 $m=\log_a x$，$n=\log_a y$。乘法 $xy=a^{m+n}$ 取对数得 $\log_a(xy)=m+n=\log_a x+\log_a y$，即乘法化为加法；除法 $x/y=a^{m-n}$ 同理化为减法。反函数关系体现在：指数函数把指数 $m$ 映为幂 $a^m$，对数函数把幂还原为指数，数对 $(x,y)$ 与 $(y,x)$ 关于 $y=x$ 对称。
\end{solution}
\begin{solution}
\textbf{拓展 2：}设 $r\in\mathbb{Q}$、$r>0$ 且 $r^2>2$（此时 $r$ 是上界），令 $t=\dfrac{r+\dfrac{2}{r}}{2}\in\mathbb{Q}$。由 $t-r=-\dfrac{r^2-2}{2r}<0$ 知 $t<r$，由 $t^2-2=\dfrac{(r^2-2)^2}{4r^2}>0$ 知 $t$ 仍是上界。故任何有理上界都不是最小，$S$ 在 $\mathbb{Q}$ 中无最小上界。几何上，$y=x^2-2$ 的零点恰是 $\sqrt2$；$\sqrt2\notin\mathbb{Q}$ 意味着这条连续曲线穿过数轴时，落点位置在有理数里没有任何对应——孔洞不仅存在，还被曲线「踩」了出来。
\end{solution}
\begin{solution}
\textbf{拓展 3：}设 $x=0.999\dots$，则 $10x=9.999\dots$，两式相减得 $9x=9$，即 $x=1$（方程法）；也可写成等比级数 $0.999\dots=0.9\bigl(1+0.1+0.01+\dots\bigr)=\dfrac{0.9}{1-0.1}=1$。「总差一点点」的错误在于：若 $1-x=\varepsilon>0$，则 $x\leq1-\varepsilon<1$，与 $x$ 的十进制表示各位都逼近 $1$ 矛盾——所谓「一点点」是一个比任何正实数都小的量，不存在。这也说明无限小数的含义要靠完备性才能定义清楚。
\end{solution}
\end{afterclass}
\reading{拓展阅读 07}{函数的结构特化}

% 高观点导学案：三角函数、解三角形与复平面旋转（第 2 章·高观点统一版）
% 主线：以「乘以 e^{iβ} = 逆时针旋转 β」为唯一母公理，让整套三角恒等变换从指数运算律自然生长
% 四模块：数系扩张与 i 的几何诞生 → 结构同构（旋转必为指数函数）→ 切线速度锁定底数 → 母公式引擎
% 另配一节：解三角形与边角互化。全程用「复平面旋转 ↔ 实虚部投影」互译，原生 TikZ 图可视化。

\section{三角函数、解三角形与复平面旋转（高观点统一版）}

\begin{methodbox}{多框架对读}
\small
\begin{center}
\begin{tabular}{c|p{2.9cm}|p{5.0cm}|p{3.1cm}}
\toprule
视角 & 本章看点 & 别的视角怎么说 & 哪侧最省力 \\
\midrule
代数 & 和差倍半公式、恒等变形 & 符号系统的机械推导 & 硬算与验算 \\
几何 & 单位圆上的旋转与投影 & 角度相加＝复数相乘 & 猜公式、看符号 \\
集合 & $\theta\mapsto(\cos\theta,\sin\theta)$ & 参数方程即一个对应 & 界定义域与值域 \\
映射 & 乘 $e^{i\beta}$＝逆时针转 $\beta$ & 旋转群与指数函数同构 & 统一推导所有公式 \\
\bottomrule
\end{tabular}
\end{center}
卡的使用：公式卡壳时回到单位圆「转一遍」；证明题优先用旋转映射（总图见第 0 章图~\ref{fig:four-view-map}）。
\end{methodbox}

本章不新增公式，而是给三角函数一章装上一个\key{母引擎}：
在复平面上，\key{乘以 $e^{i\beta}$ 就等于把向量逆时针旋转 $\beta$ 角}（模长不变）。
由这一条母公理出发，两角和差、二倍角三倍角、诱导公式、和差化积与积化和差，
全部只是指数运算律 $e^{i(\alpha+\beta)}=e^{i\alpha}e^{i\beta}$ 的代数推论。
四个模块环环相扣：
\key{数系扩张与 $i$ 的几何诞生}（为什么需要虚数）$\to$
\key{结构同构}（单位圆旋转必为指数函数）$\to$
\key{切线速度锁定底数}（欧拉公式 $e^{i\theta}=\cos\theta+i\sin\theta$）$\to$
\key{母公式引擎}（一条指数律长出全套公式），
最后用 \key{解三角形与边角互化} 把复数直觉落回高考主战场。
核心思想：\key{三角恒等变换不是几十个孤立公式，而是同一个旋转群的三种语言}（指数、坐标、三角形）。

\begin{learnobjective}
\begin{itemize}
  \item 理解虚数单位 $i$ 的几何意义（逆时针旋转 $90^\circ$），知道 $i^2=-1$ 是「转两次直角 $=$ 掉头」的代数翻译；
  \item 理解数系扩张 $\mathbb{N}\to\mathbb{Z}\to\mathbb{Q}\to\mathbb{R}\to\mathbb{C}$ 每一步都在为运算封闭与几何变换解锁新维度；
  \item 会用单位复数 $z(\theta)=\cos\theta+i\sin\theta$ 验证 $z(\alpha)z(\beta)=z(\alpha+\beta)$，理解「加法变乘法」是指数函数的基因；
  \item 会用切线速度 $(e^{k\theta})'\big|_{\theta=0}=k$ 锁定底数，写出欧拉公式 $e^{i\theta}=\cos\theta+i\sin\theta$，并得出 $e^{i\pi}+1=0$；
  \item 会用 $e^{i(\alpha+\beta)}$ 展开、棣莫弗定理与指数分解，现场推导两角和差、二倍角三倍角与和差化积／积化和差；
  \item 掌握正弦定理（外接圆半径 $R$）、余弦定理与面积公式，会按「一次齐次用正弦、二次齐次用余弦」做边角互化，并识别钝角三角形中的隐藏定义域。
\end{itemize}
\end{learnobjective}

\begin{preview}
\subsection*{知识回顾}
\begin{itemize}
  \item 实数轴上，乘以正数 $2$ 的几何意义是 \blank（拉伸／旋转）；乘以 $-1$ 的几何意义是 \blank；
  \item $(-1)\times(-1)=$ \blank；
  \item 指数函数 $a^x$ 满足的运算律是 $a^{\alpha+\beta}=$ \blank；
  \item 单位圆上辐角 $\theta$ 对应的点是 $(\cos\theta,\sin\theta)$，它到原点的距离是 \blank。
\end{itemize}

\subsection*{课前速读}
阅读三角函数一章及本章导学案，完成下列问题：
\begin{enumerate}
  \item $i^2=$ \blank；
  \item 欧拉公式 $e^{i\theta}=$ \blank；
  \item 正弦定理 $\dfrac{a}{\sin A}=$ \blank（用外接圆半径 $R$ 表示）；
  \item 余弦定理 $a^2=$ \blank。
\end{enumerate}
\end{preview}

\begin{explore}

\subsection{数系扩张与虚数单位 $i$ 的几何诞生}

\key{探究一：乘法即变换——从一维缩放、掉头到二维旋转}
在实数轴上盯住数字 $1$，看对它做乘法时的几何运动：
\begin{itemize}
  \item 乘以正数（如 $\times2$）：线段被\key{拉伸}为原来的 $2$ 倍，方向不变——纯\key{缩放}；
  \item 乘以 $-1$：把 $1$ 变成 $-1$，几何意义是绕原点逆时针旋转 \key{$180^\circ$}（半圈）；
  \item 连续乘两次 $-1$：$(-1)\times(-1)=1$，即转 $180^\circ+180^\circ=360^\circ$，恰好回到原位。
\end{itemize}
于是提出关键问题：\key{有没有一种操作，连续做两次恰好等于乘 $-1$（即一次只逆时针转 $90^\circ$）？}
设这个操作是 $i$，则
\[
  i\times i=-1\iff i^2=-1.
\]
一维数轴上根本容不下「转 $90^\circ$」后的落点，数轴被迫向垂直方向张开，出现\key{虚轴}，
一维直线跃升为二维\key{复平面}。

\begin{figure}[htbp]
  \centering
  \begin{tikzpicture}[>=Stealth, semithick, font=\footnotesize]
    % 左图：乘 -1 转 180°
    \draw[->, gray] (-2.5,0)--(2.5,0) node[below right] {实轴};
    \fill (0,0) circle (1pt);
    \draw[->, main!80!black, thick] (0,0)--(1,0) node[midway,above] {$1$};
    \draw[second, thick] (1,0) arc (0:180:1);
    \draw[->, second, thick] (0,0)--(-1,0) node[midway,above] {$-1$};
    \node[second] at (0,1.35) {$180^\circ$};
    \node[main!70!black] at (0,-2.05) {乘 $-1$};
    % 右图：乘 i 转 90°，张开虚轴
    \begin{scope}[xshift=6.2cm]
      \draw[->, gray] (-1.4,0)--(2.5,0) node[below right] {实轴};
      \draw[->, dashed, gray!70!black] (0,-1.7)--(0,2.3) node[above] {虚轴};
      \fill (0,0) circle (1pt);
      \draw[->, main!80!black, thick] (0,0)--(1,0) node[midway,below] {$1$};
      \draw[third, thick] (1,0) arc (0:90:1);
      \draw[->, third, thick] (0,0)--(0,1) node[midway,right] {$i$};
      \draw[dashed, gray] (0,1) arc (90:180:1);
      \draw[->, dashed, gray] (0,0)--(-1,0) node[midway,below] {$i^2=-1$};
      \node[third] at (0.62,0.30) {$90^\circ$};
      % 说明文字放在虚轴下端之下（虚轴只画到 $-1.7$），避免被竖直虚线穿过
      \node[main!70!black] at (0,-2.05) {乘 $i$（连乘两次得 $-1$）};
    \end{scope}
  \end{tikzpicture}
  \caption{乘法的几何：乘 $-1$ 逆时针转 $180^\circ$；乘 $i$ 只转 $90^\circ$，一维数轴被迫张开虚轴，成为复平面}
  \label{fig:rotate-i}
\end{figure}

点睛：实数 $1$ 负责向前走，$-1$ 负责向后转，$i$ 负责向左转；
连喊两次「向左转」正好就是「向后转」（$i^2=-1$）。

乘 $i$ 的作用不止于「转两次得 $-1$」：对任意辐角 $\theta$ 的单位向量，乘 $i$ 都把它的两个分量整体倒置（图~\ref{fig:i-component-flip}）。
\begin{figure}[htbp]
  \centering
  \begin{tikzpicture}[>=Stealth, semithick, font=\footnotesize]
    % ===== 左面板：转前 z=cosθ+isinθ =====
    \draw[->, gray] (-1.9,0)--(1.9,0) node[below right] {实轴};
    \draw[->, gray] (0,-0.9)--(0,1.95) node[above] {虚轴};
    \draw[gray!55, dashed] (0,0) circle (1.6);
    \draw[black!70, thick] (0,0)--(40:1.6) node[above right, font=\scriptsize] {$z$};
    \draw[blue!70!black, very thick] (0,0)--({1.6*cos(40)},0);
    \node[blue!70!black, below] at ({0.8*cos(40)},0) {$\cos\theta$};
    \draw[green!55!black, very thick] ({1.6*cos(40)},0)--(40:1.6);
    \node[green!55!black, left] at ({1.6*cos(40)+0.02},{0.8*sin(40)}) {$\sin\theta$};
    % 直角符号：两个投影分量互相垂直
    \draw[black!60, thin] ({1.6*cos(40)-0.22},0)--({1.6*cos(40)-0.22},0.22)--({1.6*cos(40)},0.22);
    \draw[black!70, thick] (0.42,0) arc (0:40:0.42) node[midway, above, font=\scriptsize] {$\theta$};
    \node[main!70!black] at (0,-0.7) {$z=\cos\theta+i\sin\theta$};
    % ===== 右面板：乘 i 后 iz=-sinθ+icosθ =====
    \begin{scope}[xshift=4.4cm]
      \draw[->, gray] (-1.9,0)--(1.9,0) node[below right] {实轴};
      \draw[->, gray] (0,-0.9)--(0,1.95) node[above] {虚轴};
      \draw[gray!55, dashed] (0,0) circle (1.6);
      \draw[dashed, gray!70!black, thick] (0,0)--(40:1.6);
      % 直角符号：z 与 iz 垂直（随 θ 取向）
      \draw[third, thin] ({0.24*cos(40)},{0.24*sin(40)})--({0.24*cos(40)-0.24*sin(40)},{0.24*sin(40)+0.24*cos(40)})--({-0.24*sin(40)},{0.24*cos(40)});
      % 90° 弧：从 θ 扫到 θ+90°
      \draw[->, third, thick] (40:0.5) arc (40:130:0.5);
      \node[third, left] at ({0.72*cos(95)},{0.72*sin(95)}) {$90^\circ$};
      \draw[->, second, very thick] (0,0)--({1.6*cos(130)},{1.6*sin(130)}) node[above left, font=\scriptsize] {$iz$};
      % 倒置后的两个分量：蓝竖（cosθ）、绿横（-sinθ）
      \draw[blue!70!black, very thick] (0,0)--(0,{1.6*cos(40)});
      \node[blue!70!black, right] at (0,{0.8*cos(40)}) {$\cos\theta$};
      \draw[green!55!black, very thick] (0,{1.6*cos(40)})--({-1.6*sin(40)},{1.6*cos(40)});
      \node[green!55!black, below] at ({-0.8*sin(40)},{1.6*cos(40)}) {$-\sin\theta$};
      \draw[dashed, gray] ({1.6*cos(130)},{1.6*sin(130)})--({1.6*cos(130)},0);
      \node[main!70!black] at (0,-0.7) {$iz=-\sin\theta+i\cos\theta$};
    \end{scope}
  \end{tikzpicture}
  \caption{乘 $i$ 的分量倒置：$z=\cos\theta+i\sin\theta$（横蓝 $\cos\theta$、纵绿 $\sin\theta$）乘 $i=e^{i\pi/2}$ 逆时针转 $90^\circ$ 后，蓝分量竖到虚轴成 $\cos\theta$、绿分量倒向负实轴成 $-\sin\theta$，即 $iz=-\sin\theta+i\cos\theta$}
  \label{fig:i-component-flip}
\end{figure}

\begin{example}
三次方程 $x^3=15x+4$ 显然有实根 $x=4$。用卡丹公式却会得到
\[
  x=\sqrt[3]{2+\sqrt{-121}}+\sqrt[3]{2-\sqrt{-121}}
  =\sqrt[3]{2+11i}+\sqrt[3]{2-11i}.
\]
请验证 $(2+i)^3=2+11i$，并说明虚部如何抵消得到 $x=4$。
\end{example}
\begin{solution}
\[
  (2+i)^3=8+12i+6i^2+i^3=8+12i-6-i=2+11i,
\]
同理 $(2-i)^3=2-11i$。于是
\[
  x=(2+i)+(2-i)=4,
\]
两个虚部一正一负抵消，还原出真实的实数解。
可见虚数不是凭空捏造的幻想，而是连接实数问题与实数解的「地下通道」：
引入 $i$ 后，$n$ 次代数方程必有 $n$ 个根（代数基本定理），数系运算彻底封闭。
\end{solution}

\begin{center}
\footnotesize
\setlength{\tabcolsep}{3pt}
\renewcommand{\arraystretch}{1.4}
\begin{tabular}{p{2.8cm}p{3.4cm}p{5.4cm}}
\toprule
扩张阶段 & 解决的方程障碍 & 解锁的几何空间与变换 \\
\midrule
$\mathbb{N}\to\mathbb{Z}$（引入负数） & $x+3=1$ 无解 & 数轴向负半轴延伸，解锁反向平移与 $180^\circ$ 掉头 \\
$\mathbb{Z}\to\mathbb{Q}$（引入分数） & $2x=1$ 无解 & 填满离散整数点之间的空隙，解锁任意比例缩放 \\
$\mathbb{Q}\to\mathbb{R}$（引入无理数） & $x^2=2$ 无解 & 消除数轴上的「针眼漏洞」，得到连续完备实数线 \\
$\mathbb{R}\to\mathbb{C}$（引入虚数 $i$） & $x^2+1=0$ 无解 & 一维直线跃升为二维平面，解锁任意角度旋转与代数封闭 \\
\bottomrule
\end{tabular}
\end{center}

\begin{conccard}[一级]{虚数单位 $i$ 的几何意义：$i^2=-1$}{理解复数的几何来源、判断复数乘法的旋转角与伸缩倍数}{把 $i$ 只当成「$\sqrt{-1}$」的代数符号，丢掉它的旋转几何；把 $i^2=-1$ 与「转 $180^\circ$」脱钩}
$1\times i$ 表示逆时针转 $90^\circ$，$1\times(-1)$ 表示转 $180^\circ$，而 $i\times i=-1$ 正是「转两次直角等于掉头」。
任何复数 $z$ 乘以 $i$，都是把 $z$ 的辐角加 $90^\circ$、模长不变。
\end{conccard}

\key{再换一个角度：每次数系扩充，保留了什么、失去了什么}
数系扩充不是「只增不减」地加数：每一步都在尽量保留旧运算律的前提下换来解题能力，同时不可避免地牺牲掉某些结构性性质。下表是四次扩充的准确清单（表中「失去」指该性质在新数系中不再整体成立）：

\begin{center}
\footnotesize
\setlength{\tabcolsep}{4pt}
\renewcommand{\arraystretch}{1.5}
\begin{tabular}{p{2.6cm}p{4.6cm}p{4.4cm}}
\toprule
扩充过程 & 保留的性质 & 不能继续保留的性质 \\
\midrule
$\mathbb{N}\to\mathbb{Z}$ & 加法、乘法及其运算律；相反数的出现使减法封闭 & 良序性（$\mathbb{N}$ 的任意非空子集有最小元；$\mathbb{Z}$ 中负整数全体就没有最小元） \\
$\mathbb{Z}\to\mathbb{Q}$ & 四则运算律与大小顺序；除法封闭、非零元都有逆元 & 离散性（不存在相邻的两个有理数，整数之间不再有空隙） \\
$\mathbb{Q}\to\mathbb{R}$ & 域结构与全序；开方、长度、极限有了合法舞台 & 可数性（$\mathbb{Q}$ 可逐个列出，$\mathbb{R}$ 不能——对角线法） \\
$\mathbb{R}\to\mathbb{C}$ & 加、减、乘、除四则运算律；代数封闭 & 与四则运算相容的大小顺序（不能像实数那样比大小） \\
\bottomrule
\end{tabular}
\end{center}

值得强调的是：前三次扩充都保留甚至加强了「能比大小」的能力，只有最后一步必须放弃它。这不是缺陷，而是换取代数封闭性的代价。

\key{为什么复数不能比较大小}
假设在 $\mathbb{C}$ 上存在一种全序 $<$，满足两条与运算相容的常规要求：($i$) $a<b\Rightarrow a+c<b+c$；($ii$) $a>0,\ b>0\Rightarrow ab>0$（于是非零数的平方必为正）。
\begin{derivation}
$i\neq0$，所以要么 $i>0$，要么 $i<0$。
\begin{itemize}
  \item 若 $i>0$：由 ($ii$) 得 $i^2>0$，但 $i^2=-1<0$，矛盾。
  \item 若 $i<0$：由 ($i$) 两边加 $-i$ 得 $0<-i$，即 $-i>0$；再由 ($ii$) 得 $(-i)^2>0$，而 $(-i)^2=i^2=-1<0$，同样矛盾。
\end{itemize}
两种情形都矛盾，故这样的全序不存在。
\end{derivation}
因此只能说「$z_1$ 与 $z_2$ 不相等」，不能说「$z_1$ 比 $z_2$ 大」；比较两个复数的唯一合法方式是比\key{模}（$|z|$ 是实数，可以比大小）或比\key{辐角}（方向）。这与实数数轴上的「谁左谁右」是完全不同的两回事。

\begin{commonerror}{把历史的顺序当成逻辑的唯一顺序}
常见讲法：「人类为了描述旋转 $90^\circ$ 才发明了 $i$。」这在数学史上不准确：虚数最早来自三次方程求根时出现的负数开平方（卡尔达诺《大术》1545 年已有记载，邦贝利随后系统研究），复平面的几何解释（把复数看成平面上的点、把乘法看成旋转伸缩）是更晚才发展出来的认识。
更严谨的教学表述是：人类先为突破实数范围内的代数求解限制而引入虚数，后来才发现它还赋予复数乘法漂亮的几何意义。
本节用「乘法即变换」作为引入路径，是\key{逻辑在先、直观优先}的教学设计，不是历史事实的顺序；讲课时应明确区分，避免学生形成「复数是为了画图方便而被发明」的错误印象。
\end{commonerror}

\begin{method}
判断一个数系够不够用的三个反射：
\begin{itemize}
  \item 方程层面：目标方程在当前数系内无解吗？（$x+3=1,\ 2x=1,\ x^2=2,\ x^2+1=0$ 各逼出一次扩充。）
  \item 运算层面：当前数系对四则运算封闭吗？对极限运算封闭吗？
  \item 结构层面：为了求解新问题，愿意放弃哪条旧性质？（序、离散性、可数性……）
\end{itemize}
\end{method}

\key{课堂讨论：数系扩充有没有终点}
\begin{itemize}
  \item \key{问题 1}：实数已经能表示数轴上的每一个点，为什么还不够用？
        （引导：完备性解决的是「连续」；方程 $x^2+1=0$ 需要的是「代数封闭」，两者是不同的要求。）
  \item \key{问题 2}：为什么复数不能比较大小？请用平方运算自己推出矛盾。
  \item \key{问题 3}：不能比大小，是不是说明复数比实数「更差」？
        （引导：失去序，换来代数封闭与二维几何；交流电路、振动、信号处理都靠复数把运算变简洁。）
\end{itemize}
给后续章节留一个伏笔：从 $\mathbb{R}$ 到 $\mathbb{C}$ 也改变了对函数的理解——实函数依赖大小顺序与数轴结构，复函数则可以借助复平面与几何变换建立新的理论。

\subsection{结构同构：单位圆旋转必为指数函数}

\key{探究二：角度相加，复数相乘}
把单位圆上辐角为 $\theta$ 的点写成复数
\[
  z(\theta)=\cos\theta+i\sin\theta.
\]
计算两个单位复数的乘积：
\begin{derivation}
\begin{align*}
  z(\alpha)z(\beta)&=(\cos\alpha+i\sin\alpha)(\cos\beta+i\sin\beta)\\
  &=(\cos\alpha\cos\beta-\sin\alpha\sin\beta)+i(\sin\alpha\cos\beta+\cos\alpha\sin\beta)\\
  &=\cos(\alpha+\beta)+i\sin(\alpha+\beta)=z(\alpha+\beta).
\end{align*}
\end{derivation}
（最后一步用了两角和公式——这正是我们稍后要用旋转重新解释的对象。）
惊人之处：\key{输入的角度相加，输出的复数相乘}。

唤醒旧知：哪类函数的核心基因是 $f(\alpha+\beta)=f(\alpha)f(\beta)$？
学生应脱口而出\key{指数函数} $f(x)=b^x$。
而 $z(0)=\cos0+i\sin0=1$，所以 $z(\theta)$ 在代数结构上就是一个底数为某复常数 $b$ 的指数函数 $z(\theta)=b^\theta$——
\key{单位圆旋转群与指数函数同构}。

\begin{conccard}[二级]{结构同构：$z(\alpha)z(\beta)=z(\alpha+\beta)$}{判断两个复数相乘的几何效果、理解三角公式的共同来源}{只记「模相乘、辐角相加」的结论，不从结构公理出发；把 $z(\theta)$ 与指数函数割裂}
单位复数 $z(\theta)=\cos\theta+i\sin\theta$ 满足「角度相加 $\to$ 复数相乘」，且 $z(0)=1$。
凡是满足 $f(\alpha+\beta)=f(\alpha)f(\beta)$ 的连续函数必为指数函数，故 $z(\theta)$ 一定可以写成 $b^\theta$ 的形式。
\end{conccard}

几何上看，$z(\alpha)\cdot e^{i\beta}$ 就是把 $z_1=z(\alpha)$ 绕原点再逆时针转 $\beta$，终点辐角为 $\alpha+\beta$（图~\ref{fig:rotate-sum}）。
\begin{figure}[htbp]
  \centering
  \begin{tikzpicture}[>=Stealth, semithick, font=\footnotesize]
    \draw[->, gray] (-1.8,0)--(3.5,0) node[below right] {$\mathrm{Re}$};
    \draw[->, gray] (0,-1.7)--(0,2.9) node[above] {$\mathrm{Im}$};
    \draw[gray!55, dashed] (0,0) circle (2.2);
    % 三段同心弧：α、β、α+β（半径与标注角错层，避免相叠）
    \draw[main, thick] (0.5,0) arc (0:20:0.5);
    \node[main] at (10:0.68) {$\alpha$};
    \draw[->, third, thick] (20:0.85) arc (20:75:0.85);
    \node[third] at (47:1.06) {$\beta$};
    \draw[->, gray!60!black, thick] (1.2,0) arc (0:75:1.2);
    \node[gray!60!black] at (12:1.45) {$\alpha+\beta$};
    % 旋转轨迹：z1 沿单位圆转到 z2
    \draw[->, third, thin] (20:2.2) arc (20:75:2.2);
    % z1=e^{iα} 与 z2=e^{i(α+β)}
    \draw[->, main!80!black, thick] (0,0)--(20:2.2) node[above right] {$z_1=e^{i\alpha}$};
    \draw[->, second, very thick] (0,0)--(75:2.2) node[above left] {$z_2=e^{i(\alpha+\beta)}$};
    % z2 的两轴投影：实部红、虚部青
    \draw[red!80!black, very thick] (0,0)--({2.2*cos(75)},0);
    \node[red!80!black, below] at ({1.1*cos(75)},0) {$\cos(\alpha+\beta)$};
    \draw[dashed, gray] (75:2.2)--(0,{2.2*sin(75)});
    \draw[cyan!80!black, very thick] ({2.2*cos(75)},0)--(75:2.2);
    \node[cyan!80!black, right] at ({2.2*cos(75)},1.2) {$\sin(\alpha+\beta)$};
    \node[main!70!black] at (0.9,-1.25) {$z_1\cdot e^{i\beta}=e^{i(\alpha+\beta)}=z_2$};
  \end{tikzpicture}
  \caption{旋转叠加：$z_1=e^{i\alpha}$ 再乘 $e^{i\beta}$ 即绕原点逆时针转 $\beta$，终点 $z_2=e^{i(\alpha+\beta)}$；$z_2$ 在实轴、虚轴上的投影 $\cos(\alpha+\beta)$、$\sin(\alpha+\beta)$ 与 $(\cos\alpha+i\sin\alpha)(\cos\beta+i\sin\beta)$ 的实部、虚部一一对应——两角和公式就是这么「转」出来的}
  \label{fig:rotate-sum}
\end{figure}

\subsection{切线速度锁定底数}

\key{探究三：切线速度——底数只能是 $e$}
对比两个「指数运动」的动力学：
\begin{itemize}
  \item 实指数 $e^{x}$：在 $x=0$ 处位置为 $1$，速度方向与位置\key{同向}（沿实轴向右），越跑越快，轨迹是直线爆炸；
  \item 复指数 $e^{i\theta}$：在 $\theta=0$ 处位置为 $1$，速度方向必须与位置\key{垂直}（沿虚轴向上，即 $+i$），只转弯不伸长，轨迹是单位圆。
\end{itemize}
把 $z(\theta)$ 看作质点位置，求导得切线速度：
\[
  z'(\theta)=-\sin\theta+i\cos\theta\implies z'(0)=i.
\]
而实指数 $f(\theta)=e^{k\theta}$ 满足
\[
  (e^{k\theta})'\Big|_{\theta=0}=k\cdot e^{0}=k.
\]
要让圆周运动在出发点的瞬时速度恰好是 $+i$，唯一可能就是 \key{$k=i$}。于是
\[
  e^{i\theta}=\cos\theta+i\sin\theta.
\]
代入 $\theta=\pi$：
\[
  e^{i\pi}=\cos\pi+i\sin\pi=-1\implies e^{i\pi}+1=0,
\]
即被誉为「上帝公式」的欧拉恒等式，把 $0,1,e,i,\pi$ 五个常数用最纯粹的圆周旋转焊在一起。

\begin{center}
\footnotesize
\setlength{\tabcolsep}{3pt}
\renewcommand{\arraystretch}{1.4}
\begin{tabular}{p{2.2cm}p{2.0cm}p{2.6cm}p{3.2cm}p{3.0cm}}
\toprule
函数 & 起点位置 & 初始速度方向 & 速度与位置关系 & 轨迹 \\
\midrule
$e^{x}$ & $1$ & $+1$（实轴向右） & 同向 & 沿实轴指数爆炸 \\
$e^{i\theta}$ & $1$ & $+i$（虚轴向上） & 垂直 & 绕单位圆匀速旋转 \\
\bottomrule
\end{tabular}
\end{center}

\begin{figure}[htbp]
  \centering
  \begin{tikzpicture}[>=Stealth, semithick, font=\footnotesize]
    \draw[->, gray] (-2.6,0)--(2.9,0) node[below right] {$\mathrm{Re}$};
    \draw[->, gray] (0,-2.3)--(0,2.5) node[above] {$\mathrm{Im}$};
    \draw[gray!55, dashed] (0,0) circle (2);
    % 一般位置 z(θ)：θ≈50°
    \draw[->, dashed, gray!80!black] (0,0)--(1.286,1.532) node[above] {$z(\theta)$};
    \draw[->, dashed, gray!80!black] (1.286,1.532)--(0.826,1.918) node[above left, font=\scriptsize] {$i\,z(\theta)$};
    % 起点 (1,0) 与切线速度 +i
    \fill (2,0) circle (1.2pt);
    \draw[->, main!80!black, thick] (0,0)--(2,0) node[midway,below] {$z(0)=1$};
    \draw[->, second, thick] (2,0)--(2,0.95) node[right] {$z'(0)=i$};
    \draw[gray, thin] (1.52,0)--(1.52,0.48)--(2,0.48);
    \draw[gray] (0.72,0) arc (0:50:0.72);
    \node[gray!80!black] at (0.98,0.22) {$\theta$};
    \node[main!70!black] at (0,-2.8) {位置向量与切线速度垂直：只转弯，不伸长};
  \end{tikzpicture}
  \caption{复平面单位圆：起点 $(1,0)$ 处位置向量为 $1$，切线速度 $+i$ 与其垂直；任意位置的速度 $i\,z(\theta)$ 始终垂直于位置}
  \label{fig:euler-tangent}
\end{figure}

\begin{conccard}[一级]{欧拉公式：$e^{i\theta}=\cos\theta+i\sin\theta$}{推导三角公式、判断复数旋转角、快速验算}{只把它当结论背诵；忘记「乘以 $e^{i\beta}$ 就是逆时针转 $\beta$」这一几何含义}
由 $z(\alpha)z(\beta)=z(\alpha+\beta)$ 与 $z(0)=1$ 知 $z(\theta)$ 必为指数函数 $b^\theta$；
再由切线速度 $(e^{k\theta})'\big|_{\theta=0}=k$ 与 $z'(0)=i$ 锁定 $k=i$。
于是乘以 $e^{i\beta}$ 等价于逆时针旋转 $\beta$，这就是本章全部三角公式的母公理。
\end{conccard}

\subsection{母公式引擎：一条指数律长出全套公式}

\key{探究四：$e^{i(\alpha+\beta)}=e^{i\alpha}e^{i\beta}$ 的展开与投影}

\key{（1）诱导公式——乘 $i=e^{i\pi/2}$ 即转 $90^\circ$。}
\[
  i(\cos\theta+i\sin\theta)=-\sin\theta+i\cos\theta,
\]
对比实部虚部，一眼读出
\[
  \cos\left(\theta+\frac{\pi}{2}\right)=-\sin\theta,\qquad
  \sin\left(\theta+\frac{\pi}{2}\right)=\cos\theta.
\]
同理乘 $-1=e^{i\pi}$ 得 $\cos(\theta+\pi)=-\cos\theta$、$\sin(\theta+\pi)=-\sin\theta$；
取共轭 $\overline{e^{i\theta}}=e^{-i\theta}$ 即关于实轴对称，这正是奇偶性
$\cos(-\theta)=\cos\theta$、$\sin(-\theta)=-\sin\theta$ 的几何本质。
「奇变偶不变，符号看象限」不再是口诀，而是「转了几个直角」的图形事实。

\key{（2）两角和与差——多项式展开一步到位。}
\begin{derivation}
\begin{align*}
  e^{i(\alpha+\beta)}&=e^{i\alpha}e^{i\beta}
  =(\cos\alpha+i\sin\alpha)(\cos\beta+i\sin\beta)\\
  &=(\cos\alpha\cos\beta-\sin\alpha\sin\beta)
  +i(\sin\alpha\cos\beta+\cos\alpha\sin\beta).
\end{align*}
\end{derivation}
左边又是 $\cos(\alpha+\beta)+i\sin(\alpha+\beta)$，对比实部与虚部：
\[
  \cos(\alpha+\beta)=\cos\alpha\cos\beta-\sin\alpha\sin\beta,\qquad
  \sin(\alpha+\beta)=\sin\alpha\cos\beta+\cos\alpha\sin\beta.
\]
把 $\beta$ 换成 $-\beta$（用 $\cos(-\beta)=\cos\beta$、$\sin(-\beta)=-\sin\beta$），立刻得到差角公式：
\[
  \cos(\alpha-\beta)=\cos\alpha\cos\beta+\sin\alpha\sin\beta,\qquad
  \sin(\alpha-\beta)=\sin\alpha\cos\beta-\cos\alpha\sin\beta.
\]

\key{（3）二倍角与三倍角——棣莫弗定理。}
由 $e^{in\theta}=(e^{i\theta})^n$ 得
\[
  (\cos\theta+i\sin\theta)^n=\cos n\theta+i\sin n\theta.
\]
$n=2$ 展开即得二倍角；$n=3$ 用二项式定理：
\[
  (\cos\theta+i\sin\theta)^3
  =(\cos^3\theta-3\cos\theta\sin^2\theta
  +i(3\cos^2\theta\sin\theta-\sin^3\theta),
\]
于是
\[
  \cos3\theta=4\cos^3\theta-3\cos\theta,\qquad
  \sin3\theta=3\sin\theta-4\sin^3\theta.
\]

\key{（4）和差化积与积化和差——同一指数分解。}
由欧拉公式解出
\[
  \cos\theta=\frac{e^{i\theta}+e^{-i\theta}}{2},\qquad
  \sin\theta=\frac{e^{i\theta}-e^{-i\theta}}{2i}.
\]
例如
\begin{align*}
  \cos\alpha\cos\beta
  &=\frac{(e^{i\alpha}+e^{-i\alpha})(e^{i\beta}+e^{-i\beta})}{4}\\
  &=\frac{e^{i(\alpha+\beta)}+e^{-i(\alpha+\beta)}+e^{i(\alpha-\beta)}+e^{-i(\alpha-\beta)}}{4}\\
  &=\frac{\cos(\alpha+\beta)+\cos(\alpha-\beta)}{2},
\end{align*}
这就是积化和差；把 $\alpha+\beta=A$、$\alpha-\beta=B$ 反解即得和差化积。
结论：和差、乘积、倍角只是同一组指数幂的不同代数外形。

\begin{example}
用棣莫弗定理把 $\cos3\theta$ 用 $\cos\theta$ 表示。
\end{example}
\begin{solution}
\[
  \cos3\theta=\operatorname{Re}(\cos\theta+i\sin\theta)^3
  =\cos^3\theta-3\cos\theta\sin^2\theta
  =\cos^3\theta-3\cos\theta(1-\cos^2\theta)
  =4\cos^3\theta-3\cos\theta.
\]
\end{solution}

\begin{example}
求 $\cos\theta+\cos2\theta+\cdots+\cos n\theta$（$\theta\neq2k\pi$）。
\end{example}
\begin{solution}
令 $q=e^{i\theta}$，则 $\cos k\theta=\operatorname{Re}q^k$，于是
\[
  S=\sum_{k=1}^{n}q^k=\frac{q(1-q^n)}{1-q}
  =\frac{e^{i\theta}(1-e^{in\theta})}{1-e^{i\theta}}
  =e^{i\frac{(n+1)\theta}{2}}\cdot\frac{\sin\frac{n\theta}{2}}{\sin\frac{\theta}{2}}.
\]
取实部即得
\[
  \cos\theta+\cos2\theta+\cdots+\cos n\theta
  =\frac{\sin\frac{n\theta}{2}}{\sin\frac{\theta}{2}}\cos\frac{(n+1)\theta}{2}.
\]
实数域里要凑 $\sin\frac{\theta}{2}$ 裂项，技巧性强；复数域里只是公比为 $q=e^{i\theta}$ 的等比数列求和。
\end{solution}

\begin{conccard}[一级]{考纲安全边界：复数法用于选填秒杀，解答题须还原为实数格式}{选择填空快速求值、验算；忘记公式时现场推导；快速判断符号}{在解答题中直接写 $e^{i\theta}$、$i$ 的推导过程而跳过实数三角变换步骤，造成阅卷扣分}
高考解答题默认在实数域 $\mathbb{R}$ 内评分。复数法与欧拉公式的正确用法是：
选填小题秒杀、快速验算、忘记公式时现场推导。
解答题书写必须还原为实数三角变换格式（和差倍半公式、正弦定理、余弦定理的标准步骤），
不能把 $e^{i\theta}$ 与 $i$ 直接写进解答题过程。
\end{conccard}

\begin{commonerror}{指数展开时丢符号：忘了 $i^2=-1$}
用 $e^{i(\alpha+\beta)}$ 展开时，交叉项来自 $i\sin\alpha\cdot i\sin\beta=i^2\sin\alpha\sin\beta=-\sin\alpha\sin\beta$。
若忘记 $i^2=-1$，就会把实部写成 $\cos\alpha\cos\beta+\sin\alpha\sin\beta$，得到错误的「和角公式」。
记住：\key{每出现一个 $i^2$ 就要换成 $-1$}，符号自然正确。
\end{commonerror}

\begin{method}
母公式引擎使用流程（复数法三步）：
\begin{enumerate}
  \item \key{翻译}：把三角式写成 $e^{i\theta}$ 的指数形式（$\cos\theta=\frac{e^{i\theta}+e^{-i\theta}}2$，$\sin\theta=\frac{e^{i\theta}-e^{-i\theta}}{2i}$）；
  \item \key{指数运算}：用 $e^{i(\alpha+\beta)}=e^{i\alpha}e^{i\beta}$、棣莫弗定理、等比数列求和等纯代数工具计算；
  \item \key{取实虚部还原}：需要余弦取实部，需要正弦取虚部，最后写回实数三角形式。
\end{enumerate}
\end{method}

\subsection{解三角形与边角互化}

\key{探究五：正余弦定理——把三角形放进外接圆}
正弦定理：
\[
  \frac{a}{\sin A}=\frac{b}{\sin B}=\frac{c}{\sin C}=2R,
\]
其中 $R$ 是 $\triangle ABC$ 外接圆半径。几何意义：作外接圆，由圆周角定理得 $a=2R\sin A$，其余同理。
余弦定理：
\[
  a^2=b^2+c^2-2bc\cos A
  \iff \cos A=\frac{b^2+c^2-a^2}{2bc}.
\]
面积公式：
\[
  S=\frac12bc\sin A=\frac12ca\sin B=\frac12ab\sin C=\frac{abc}{4R}.
\]

\key{边角互化的选择标准。}
\begin{itemize}
  \item 关于 $a,b,c$（或 $\sin A,\sin B,\sin C$）的\key{一次齐次式} $\to$ 用正弦定理，把边化角或把角化边；
  \item 关于 $a^2,b^2,c^2$（或 $\sin^2$）的\key{二次齐次式} $\to$ 用余弦定理。
\end{itemize}

\begin{example}
在 $\triangle ABC$ 中，已知 $a=2b\cos C$，判断三角形的形状。
\end{example}
\begin{solution}
由正弦定理把边化角：$\sin A=2\sin B\cos C$。
又 $A=\pi-(B+C)$，故 $\sin A=\sin(B+C)=\sin B\cos C+\cos B\sin C$，于是
\[
  \sin B\cos C+\cos B\sin C=2\sin B\cos C
  \implies \cos B\sin C-\sin B\cos C=0
  \implies \sin(C-B)=0.
\]
因为 $B,C\in(0,\pi)$，所以 $C-B=0$，即 $B=C$。故三角形为等腰三角形。
（一次齐次式，选正弦定理。）
\end{solution}

\begin{example}
在 $\triangle ABC$ 中，已知 $a^2=b^2+c^2-bc$，求角 $A$。
\end{example}
\begin{solution}
与余弦定理 $a^2=b^2+c^2-2bc\cos A$ 对比，得 $2\cos A=1$，即 $\cos A=\frac12$。
又 $A\in(0,\pi)$，故 $A=\dfrac{\pi}{3}$（$60^\circ$）。
（二次齐次式，选余弦定理。）
\end{solution}

\key{钝角三角形的隐藏定义域。}
当题目给出「钝角三角形」或隐含某角为钝角时，角的范围被额外压缩：
钝角 $\gamma\in\left(\dfrac{\pi}{2},\pi\right)$，两个锐角满足 $\alpha+\beta<\dfrac{\pi}{2}$。
用余弦定理处理时，$\cos\gamma<0$ 给出边的平方约束，而三角形存在条件给出另一侧约束，两者共同构成隐藏定义域。

\begin{example}
在 $\triangle ABC$ 中，已知 $a=2$，$b=3$，且 $C$ 为钝角，求边 $c$ 的取值范围。
\end{example}
\begin{solution}
由余弦定理
\[
  \cos C=\frac{a^2+b^2-c^2}{2ab}=\frac{13-c^2}{12}.
\]
$C$ 为钝角 $\iff \cos C<0\iff c^2>13\iff c>\sqrt{13}$。
同时三角形存在要求 $c<a+b=5$（且 $c>|b-a|=1$，已被 $c>\sqrt{13}$ 包含）。
故 $c\in(\sqrt{13},5)$。
若只写 $c^2>13$，就漏掉了三角形存在给出的上界 $c<5$——这正是「钝角三角形」隐藏的定义域。
\end{solution}

\begin{conccard}[二级]{边角互化：一次齐次用正弦，二次齐次用余弦}{解三角形中的求值、求角、判断形状、求范围}{齐次次数判断错误；用正弦定理求角时忽略两解；忽略钝角三角形对边角范围的额外约束}
正弦定理 $\frac{a}{\sin A}=2R$ 把边与角的正弦互相转化，适合一次齐次式；
余弦定理 $a^2=b^2+c^2-2bc\cos A$ 把三边与一角联系起来，适合二次齐次式。
涉及「钝角三角形」时，必须同时写出 $\cos\gamma<0$ 与三角形存在条件，得到完整的定义域。
\end{conccard}

\begin{commonerror}{解三角形忽略隐藏定义域}
两类常见失误：
\begin{itemize}
  \item 用正弦定理由 $\sin B=\frac{b\sin A}{a}$ 求角时，只取锐角解，漏掉钝角解；
  \item 题目说「钝角三角形」却只列 $\cos\gamma<0$，漏掉三角形存在条件（两边之和大于第三边）。
\end{itemize}
例如由 $a=2,b=3$ 且 $C$ 为钝角求 $c$ 时，只写 $c>\sqrt{13}$ 会漏掉 $c<5$。
正确做法是把「角的范围」与「三角形存在条件」同时列出，求交集。
\end{commonerror}

\begin{summary}
本章以「乘以 $e^{i\beta}=$ 逆时针旋转 $\beta$」为唯一母公理，把三角函数从零散公式还原为同一个旋转群的三种语言：
数系扩张 $\mathbb{N}\to\mathbb{Z}\to\mathbb{Q}\to\mathbb{R}\to\mathbb{C}$ 让 $i^2=-1$（转 $90^\circ$ 两次等于掉头）成为二维复平面的入场券；
单位复数 $z(\theta)=\cos\theta+i\sin\theta$ 满足 $z(\alpha)z(\beta)=z(\alpha+\beta)$ 且 $z(0)=1$，故必为指数函数；
由切线速度 $(e^{k\theta})'\big|_{\theta=0}=k$ 与 $z'(0)=i$ 锁定 $k=i$，得到欧拉公式 $e^{i\theta}=\cos\theta+i\sin\theta$ 与 $e^{i\pi}+1=0$；
再用 $e^{i(\alpha+\beta)}$ 展开、棣莫弗定理与 $\cos\theta=\frac{e^{i\theta}+e^{-i\theta}}2$、$\sin\theta=\frac{e^{i\theta}-e^{-i\theta}}{2i}$，
一次性长出两角和差、二倍角三倍角、诱导公式与和差化积／积化和差。
解三角形部分用正弦定理（外接圆半径 $R$）、余弦定理与面积公式，按「一次齐次用正弦、二次齐次用余弦」做边角互化，并警惕钝角三角形的隐藏定义域。
考纲边界：复数法用于选填秒杀与验算，解答题必须还原为实数三角变换格式。
\end{summary}

\end{explore}

\begin{inclasstest}
\begin{enumerate}
  \item 计算 $i\cdot(1+i)$，并说明它在复平面上表示怎样的旋转与伸缩；\level{易}
  \item 用 $e^{i(\alpha+\beta)}=e^{i\alpha}e^{i\beta}$ 展开并对比实部与虚部，写出 $\cos(\alpha-\beta)$ 与 $\sin(\alpha-\beta)$；\level{中}
  \item 用棣莫弗定理把 $\cos3\theta$ 用 $\cos\theta$ 表示；\level{中}
  \item 在 $\triangle ABC$ 中，$a=2$，$b=3$，$C=120^\circ$，求边 $c$ 与面积 $S$。\level{中}
\end{enumerate}
\end{inclasstest}

\begin{afterclass}
\subsection*{基础巩固}
\begin{enumerate}
  \item 计算 $i^3$、$i^4$ 与 $(1+i)^2$，并说出各自的几何意义；
  \item 由 $e^{i\theta}=\cos\theta+i\sin\theta$ 解出 $\cos\theta,\sin\theta$ 的指数表达式，并验证 $\cos^2\theta+\sin^2\theta=1$。
\end{enumerate}

\subsection*{进阶提升}
\begin{enumerate}
  \item 用 $i=e^{i\pi/2}$ 推导 $\cos\left(\theta-\dfrac{\pi}{2}\right)$ 与 $\sin\left(\theta-\dfrac{\pi}{2}\right)$；
  \item 在 $\triangle ABC$ 中，$a=2$，$b=3$，$c=4$，判断三角形的形状并求面积。
\end{enumerate}

\subsection*{拓展探究}
\begin{enumerate}
  \item 用复数等比数列求 $\cos\theta+\cos2\theta+\cdots+\cos n\theta$（$\theta\neq2k\pi$），写出结果并说明它为什么只是「取实部」。
\end{enumerate}

\subsection*{教师版答案}
\begin{solution}
\textbf{基础 1：}$i^3=i^2\cdot i=-i$（逆时针转 $270^\circ$）；$i^4=(i^2)^2=1$（转 $360^\circ$ 回到原位）；
$(1+i)^2=1+2i+i^2=2i$（模由 $\sqrt2$ 变为 $2$，辐角由 $45^\circ$ 变为 $90^\circ$）。
\end{solution}
\begin{solution}
\textbf{基础 2：}由 $e^{i\theta}=\cos\theta+i\sin\theta$、$e^{-i\theta}=\cos\theta-i\sin\theta$ 解得
\[
  \cos\theta=\frac{e^{i\theta}+e^{-i\theta}}{2},\qquad
  \sin\theta=\frac{e^{i\theta}-e^{-i\theta}}{2i}.
\]
于是
\[
  \cos^2\theta+\sin^2\theta
  =\frac{(e^{i\theta}+e^{-i\theta})^2}{4}+\frac{(e^{i\theta}-e^{-i\theta})^2}{(2i)^2}
  =\frac{e^{2i\theta}+2+e^{-2i\theta}-e^{2i\theta}+2-e^{-2i\theta}}{4}=1.
\]
\end{solution}
\begin{solution}
\textbf{进阶 1：}乘以 $-i=e^{-i\pi/2}$（顺时针转 $90^\circ$）：
\[
  (\cos\theta+i\sin\theta)(-i)=\sin\theta-i\cos\theta,
\]
对比实虚部得
\[
  \cos\left(\theta-\frac{\pi}{2}\right)=\sin\theta,\qquad
  \sin\left(\theta-\frac{\pi}{2}\right)=-\cos\theta.
\]
\end{solution}
\begin{solution}
\textbf{进阶 2：}最大边为 $c=4$，
\[
  \cos C=\frac{a^2+b^2-c^2}{2ab}=\frac{4+9-16}{12}=-\frac14<0,
\]
故 $C$ 为钝角，三角形是钝角三角形。
\[
  \sin C=\sqrt{1-\frac1{16}}=\frac{\sqrt{15}}4,\qquad
  S=\frac12ab\sin C=\frac12\cdot2\cdot3\cdot\frac{\sqrt{15}}4=\frac{3\sqrt{15}}4.
\]
\end{solution}
\begin{solution}
\textbf{拓展：}令 $q=e^{i\theta}$，则
\[
  \sum_{k=1}^{n}q^k=\frac{q(1-q^n)}{1-q}
  =\frac{e^{i\theta}(1-e^{in\theta})}{1-e^{i\theta}}
  =e^{i\frac{(n+1)\theta}{2}}\cdot\frac{\sin\frac{n\theta}{2}}{\sin\frac{\theta}{2}}.
\]
取实部得
\[
  \cos\theta+\cos2\theta+\cdots+\cos n\theta
  =\frac{\sin\frac{n\theta}{2}}{\sin\frac{\theta}{2}}\cos\frac{(n+1)\theta}{2}.
\]
因为 $\cos k\theta=\operatorname{Re}(e^{ik\theta})$，所以它只是等比数列求和后「取实部」。
\end{solution}
\reading{拓展阅读 08}{复数与三角的大一统}
\end{afterclass}

% 高观点导学案：平面向量与空间向量（第 3 章·高观点统一版）
% 主线：既有大小又有方向的量 → 点集的位移 → 基底下的坐标 → 内积投影算子 → 空间中的方向与法向
% 四模块：基底分解与坐标表示 → 数量积与投影算子 → 极化恒等式与面积 → 空间向量（方向与法向）
% 全程突出向量体系的独立地位：向量不是三角／几何的附庸，而是「几何空间代数化」的通用线性语言。

\section{平面向量与空间向量（高观点统一版）}

\begin{methodbox}{多框架对读}
\small
\begin{center}
\begin{tabular}{c|p{2.9cm}|p{5.0cm}|p{3.1cm}}
\toprule
视角 & 本章看点 & 别的视角怎么说 & 哪侧最省力 \\
\midrule
代数 & 坐标运算、数量积公式 & 有序数对的算术 & 硬算定坐标 \\
几何 & 分解、投影、法向量 & 图形中看方向与垂直 & 找关系、省计算 \\
集合 & 向量＝有向线段构成的集合 & 基底＝给集合装坐标系 & 界线性相关无关 \\
映射 & 点 $\to$ 坐标是同构；投影是算子 & 换基底＝换表示 & 整体搬迁结论 \\
\bottomrule
\end{tabular}
\end{center}
卡的使用：几何题先问「有没有现成的算子（投影／法向量）」，再决定建系还是综合（总图见第 0 章图~\ref{fig:four-view-map}）。
\end{methodbox}

第 1 章研究「对应」（函数），第 2 章研究「旋转」（三角函数）；本章登场的向量是第三条独立主线：
它研究\key{位移的合成与分解}。向量体系相对三角与函数具有显性独立地位——
三角函数是旋转群的语言，函数是对应的语言，而向量是\key{几何空间代数化}的通用线性语言：
不依赖角度公式、不依赖函数图象，只用「大小$+$方向」「基底分解」「数量积」三件工具，
就能把平面与空间中的距离、夹角、平行、垂直、面积全部翻译成代数。
四个模块层层展开：
\key{基底分解与坐标表示}（任一向量 $=$ 基底线性组合，唯一）$\to$
\key{数量积与投影算子}（数量积 $=$ 余弦定理的代数化）$\to$
\key{极化恒等式与面积}（共起点三向量乘积的秒杀通道）$\to$
\key{空间向量：方向与法向}（用法向量统一处理空间位置关系）。
核心思想：\key{选定一组基底，几何就变成了线性代数}。

\begin{learnobjective}
\begin{itemize}
  \item 理解向量「既有大小又有方向」的双重属性，区分向量与数量，理解自由向量（可平移）与零向量；
  \item 理解平面向量基本定理：不共线向量张成整个平面，任一向量有唯一分解；坐标即有序实数对，与第 1 章「一一对应」呼应；
  \item 掌握数量积的定义、坐标运算与投影公式，能把余弦定理与数量积互译（数量积就是余弦定理的代数化）；
  \item 掌握极化恒等式 $\overrightarrow{AB}\cdot\overrightarrow{AC}=|\overrightarrow{AM}|^2-|\overrightarrow{MB}|^2$（$M$ 为 $BC$ 中点），会用它秒杀共起点三向量乘积的最值问题；知道面积形式 $S=\frac12|\vec a\times\vec b|$；
  \item 会把空间向量坐标化，理解平面法向量、直线方向向量，用法向量统一处理线面／面面的平行、垂直与夹角（线面角、二面角）。
\end{itemize}
\end{learnobjective}

\begin{preview}
\subsection*{知识回顾}
\begin{itemize}
  \item 实数加法满足交换律：$a+b=$ \blank；类比思考：向量加法满足交换律吗？\blank；
  \item 平面直角坐标系中，点 $P$ 的坐标是 \blank（有序实数对／无序数组），它体现了点与 \blank 的一一对应；
  \item 两直线（斜率存在时）垂直的判定：斜率之积等于 \blank；
  \item 三角形面积公式 $S=\dfrac12\times$底$\times$ \blank；平行四边形面积 $=$ 底 $\times$ \blank。
\end{itemize}

\subsection*{课前速读}
阅读「平面向量」与「空间向量」两章内容，完成下列问题：
\begin{enumerate}
  \item 平面向量基本定理要求两个基底向量 \blank（不共线／相等）；
  \item 设 $\vec a=(x_1,y_1)$，$\vec b=(x_2,y_2)$，则 $\vec a\cdot\vec b=$ \blank；
  \item 与平面 \blank 的向量叫做该平面的法向量；
  \item 设直线方向向量为 $\vec s$、平面法向量为 $\vec n$，线面角 $\varphi$ 满足 $\sin\varphi=$ \blank。
\end{enumerate}
\end{preview}

\begin{explore}

\subsection{基底分解：平面向量基本定理与坐标表示}

\key{探究一：任一向量都能被基底「唯一拼装」出来}

\key{（1）从「既有大小又有方向的量」到「点集的位移」。}
向量是既有大小又有方向的量，用有向线段表示；
两个向量只要大小相等、方向相同就叫\key{相等向量}，与起点位置无关（\key{自由向量}，可平移）。
位移视角：$\overrightarrow{AB}$ 表示点 $A$ 到点 $B$ 的位移。把所有起点平移到原点 $O$，
向量 $\overrightarrow{OP}$ 与平面点 $P$ 一一对应——这就是「向量 $=$ 点集的位移」的准确含义。
规定：零向量 $\vec 0$（长度为 $0$，方向任意）；$\vec a$ 的相反向量记 $-\vec a$。

\key{（2）平面向量基本定理：不共线的两个向量张成整个平面。}
\begin{derivation}
定理内容：如果 $\vec e_1,\vec e_2$ 是同一平面内的两个\key{不共线}向量，
那么对这一平面内的任一向量 $\vec a$，存在\key{唯一}一对实数 $\lambda_1,\lambda_2$，使得
\[
  \vec a=\lambda_1\vec e_1+\lambda_2\vec e_2.
\]
存在性：过 $\vec a$ 的终点分别作 $\vec e_2,\vec e_1$ 方向的平行线，围成平行四边形，
两边即 $\lambda_1\vec e_1$ 与 $\lambda_2\vec e_2$（平行四边形法则）。
唯一性：设 $\vec a=\lambda_1\vec e_1+\lambda_2\vec e_2=\mu_1\vec e_1+\mu_2\vec e_2$，则
\[
  (\lambda_1-\mu_1)\vec e_1+(\lambda_2-\mu_2)\vec e_2=\vec 0.
\]
若 $\lambda_1\neq\mu_1$，则 $\vec e_1=\dfrac{\mu_2-\lambda_2}{\lambda_1-\mu_1}\vec e_2$，
与 $\vec e_1,\vec e_2$ 不共线矛盾；故 $\lambda_1=\mu_1$，同理 $\lambda_2=\mu_2$。唯一性得证。
\end{derivation}
$\{\vec e_1,\vec e_2\}$ 称为一组\key{基底}。注意基底不唯一、也不要求垂直——
这是与「坐标轴」的本质区别，也是「选得巧，计算就少」的前提。

\key{（3）坐标：基底下的有序实数对。}
取\key{正交单位基底} $\vec i,\vec j$（互相垂直、长度都为 $1$），则
\[
  \vec a=x\vec i+y\vec j \iff \vec a=(x,y),
\]
$(x,y)$ 就是 $\vec a$ 的\key{坐标}。当起点在原点时，向量 $(x,y)$ 恰好指向点 $(x,y)$——
几何对象（向量／点）、代数对象（有序实数对）、对应关系三者同构，
这正是第 1 章「笛卡尔积坐标化」思想的延续：\key{坐标 $=$ 向量在基底下的读数}。
坐标运算规则：
\[
  \vec a\pm\vec b=(x_1\pm x_2,\,y_1\pm y_2),\qquad
  \lambda\vec a=(\lambda x_1,\,\lambda y_1),\qquad
  |\vec a|=\sqrt{x_1^2+y_1^2}.
\]

\begin{figure}[htbp]
  \centering
  \begin{tikzpicture}[>=Stealth, semithick, font=\footnotesize]
    % ===== 左：斜基底的一般分解 =====
    \coordinate (O) at (0,0);
    \coordinate (E1) at (2,0.5);
    \coordinate (E2) at (0.6,1.6);
    \coordinate (P1) at (4,1);
    \coordinate (P2) at (4.6,2.6);
    \draw[->, gray] (-0.5,0)--(5,0);
    \draw[->, gray] (0,-0.5)--(0,3);
    \fill (O) circle (1pt);
    \draw[->, third, thick] (O)--(E1) node[right] {$\vec{e}_1$};
    \draw[->, third, thick] (O)--(E2) node[above left] {$\vec{e}_2$};
    \draw[->, main!80!black, very thick] (O)--(P2) node[above right] {$\vec{a}=x\vec{e}_1+y\vec{e}_2$};
    \draw[->, second, thick] (O)--(P1) node[midway,below] {$x\vec{e}_1$};
    \draw[->, second, thick] (P1)--(P2) node[midway,right] {$y\vec{e}_2$};
    \draw[dashed, second!70!black] (E2)--(P2);
    \node[main!70!black] at (0,-0.85) {一般基底：平行四边形法则唯一分解};
    % ===== 右：正交单位基底与坐标 =====
    \begin{scope}[xshift=6.6cm]
      \draw[->, gray] (-0.5,0)--(4.2,0) node[below right] {$x$};
      \draw[->, gray] (0,-0.5)--(0,3.2) node[above left] {$y$};
      \foreach \x in {1,2,3} \draw[gray!25, very thin] (\x,-0.4)--(\x,3);
      \foreach \y in {1,2} \draw[gray!25, very thin] (-0.4,\y)--(4,\y);
      \fill (0,0) circle (1pt);
      \draw[->, third, thick] (0,0)--(1,0) node[below] {$\vec{i}=\vec{e}_1$};
      \draw[->, third, thick] (0,0)--(0,1) node[left] {$\vec{j}=\vec{e}_2$};
      \draw[->, main!80!black, very thick] (0,0)--(3,2) node[above right] {$\vec{a}=(3,2)$};
      \draw[->, second, thick] (0,0)--(3,0) node[midway,below] {$3\vec{i}$};
      \draw[->, second, thick] (3,0)--(3,2) node[midway,right] {$2\vec{j}$};
      \draw[dashed, second!70!black] (0,2)--(3,2);
      \node[main!70!black] at (2,-0.85) {正交单位基底：读数即坐标};
    \end{scope}
  \end{tikzpicture}
  \caption{平面向量的基底分解与坐标：左图为任意不共线基底 $\{\vec e_1,\vec e_2\}$ 下 $\vec a=x\vec e_1+y\vec e_2$ 的唯一分解；右图取正交单位基底 $\{\vec i,\vec j\}$ 时，读数 $(x,y)$ 就是平面直角坐标}
  \label{fig:basis-coord}
\end{figure}

\key{（4）「基底选得巧，计算就少」。}
分解的目的不是形式，而是让未知关系落到已知方向上。看一个反复出现的模型：

\begin{example}
在 $\triangle ABC$ 中，$D$ 在 $BC$ 上且 $BD:DC=1:2$，用基底 $\{\overrightarrow{AB},\overrightarrow{AC}\}$ 表示 $\overrightarrow{AD}$。
\end{example}
\begin{solution}
取 $\overrightarrow{AB},\overrightarrow{AC}$ 为基底。由 $BD:DC=1:2$ 知 $\overrightarrow{BD}=\dfrac13\overrightarrow{BC}$，于是
\[
  \overrightarrow{AD}=\overrightarrow{AB}+\overrightarrow{BD}
  =\overrightarrow{AB}+\frac13\overrightarrow{BC}
  =\overrightarrow{AB}+\frac13(\overrightarrow{AC}-\overrightarrow{AB})
  =\frac23\overrightarrow{AB}+\frac13\overrightarrow{AC}.
\]
取 $D$ 为中点时即得常用结论 $\overrightarrow{AD}=\dfrac12(\overrightarrow{AB}+\overrightarrow{AC})$。
\key{一旦基底选定，所有向量都变成有序数对，几何题就变成解方程组。}
\end{solution}

\begin{conccard}[一级]{平面向量基本定理与坐标表示}{把任意向量用两个不共线向量线性表示；判断三点共线；坐标运算}{认为「基底必须互相垂直」；分解不唯一时怀疑定理；写出坐标却忘记它依赖所选基底}
$\{\vec e_1,\vec e_2\}$ 不共线 $\implies$ 平面内任一向量 $\vec a$ 都可唯一写成 $\lambda_1\vec e_1+\lambda_2\vec e_2$。
特别地，$\overrightarrow{OP}=x\overrightarrow{OA}+y\overrightarrow{OB}$ 且 $x+y=1\iff P,A,B$ 三点共线（等和线）。
\end{conccard}

\subsection{数量积与投影算子}

\key{探究二：数量积是「投影乘模长」，更是余弦定理的代数化}

\key{（1）定义与几何意义。}
\[
  \vec a\cdot\vec b=|\vec a|\,|\vec b|\cos\theta,\qquad \theta=\langle\vec a,\vec b\rangle\in[0,\pi].
\]
几何读法：先取 $\vec a$ 在 $\vec b$ 方向上的投影 $|\vec a|\cos\theta$（可正可负可为零），再乘以 $|\vec b|$。
特例：同向 $\vec a\cdot\vec b=|\vec a||\vec b|$；反向 $=-|\vec a||\vec b|$；垂直 $=0$。
显然 $\vec a\cdot\vec a=|\vec a|^2$——\key{数量积自带求模公式}。

\key{（2）坐标运算：把定义交给分配律。}
\begin{derivation}
\begin{align*}
  \vec a\cdot\vec b
  &=(x_1\vec i+y_1\vec j)\cdot(x_2\vec i+y_2\vec j)\\
  &=x_1x_2\,\underbrace{\vec i\cdot\vec i}_{=1}
    +(x_1y_2+x_2y_1)\,\underbrace{\vec i\cdot\vec j}_{=0}
    +y_1y_2\,\underbrace{\vec j\cdot\vec j}_{=1}
  =x_1x_2+y_1y_2.
\end{align*}
\end{derivation}
即「对应坐标相乘再相加」。由此还可直接验证交换律 $\vec a\cdot\vec b=\vec b\cdot\vec a$ 与分配律。
夹角公式：$\cos\theta=\dfrac{\vec a\cdot\vec b}{|\vec a||\vec b|}=\dfrac{x_1x_2+y_1y_2}{\sqrt{x_1^2+y_1^2}\sqrt{x_2^2+y_2^2}}$。

\key{（3）投影：数量投影与投影向量。}
\[
  \text{数量投影}=\frac{\vec a\cdot\vec b}{|\vec b|}=|\vec a|\cos\theta,\qquad
  \mathrm{proj}_{\vec b}\vec a=\frac{\vec a\cdot\vec b}{|\vec b|^2}\,\vec b.
\]
记忆钩子：\key{投影向量 $=$ 数量投影 $\times$ 单位向量}，即 $\mathrm{proj}_{\vec b}\vec a=\left(\vec a\cdot\dfrac{\vec b}{|\vec b|}\right)\dfrac{\vec b}{|\vec b|}$。
注意分母是 $|\vec b|^2$（保证量纲是长度），且投影方向可以逆着 $\vec b$（钝角时为负）。

\begin{figure}[htbp]
  \centering
  \begin{tikzpicture}[>=Stealth, semithick, font=\footnotesize]
    \coordinate (O) at (0,0);
    \coordinate (B) at (4,0);
    \coordinate (A) at (2.6,2.2);
    \coordinate (H) at (2.6,0);
    \draw[->, gray] (-0.4,0)--(4.6,0);
    \draw[->, gray] (0,-0.4)--(0,2.6);
    \fill (O) circle (1pt) node[below left] {$O$};
    \draw[->, main!80!black, very thick] (O)--(B) node[below right] {$\vec{b}$};
    \draw[->, main!80!black, very thick] (O)--(A) node[above] {$\vec{a}$};
    \draw[dashed, second, thick] (A)--(H);
    \draw[->, second, thick] (O)--(H);
    \draw[gray, thin] (2.35,0)--(2.35,0.25)--(2.6,0.25);
    \draw[gray] (0.9,0) arc (0:40:0.9);
    \node[gray!80!black] at (1.2,0.26) {$\theta$};
    % 两条公式统一放到底部居中说明，避免压到 x 轴与图形外撑宽图幅
    \node[second, align=center] at (2.1,-0.95)
      {向量投影 $\mathrm{proj}_{\vec b}\vec a=\dfrac{\vec a\cdot\vec b}{|\vec b|^2}\vec b$\\[3pt]
       数量投影 $=\dfrac{\vec a\cdot\vec b}{|\vec b|}=|\vec a|\cos\theta$};
  \end{tikzpicture}
  \caption{数量积的投影几何：$\vec a$ 在 $\vec b$ 方向上的垂足为 $H$，$OH$ 即投影向量；数量积 $=|\vec b|\times|\vec a|\cos\theta$，即「投影乘模长」}
  \label{fig:dot-proj}
\end{figure}

\key{（4）与余弦定理互译：数量积就是余弦定理的代数化。}
\begin{derivation}
在 $\triangle ABC$ 中记 $\overrightarrow{CA}=\vec u$，$\overrightarrow{CB}=\vec v$，则 $|\vec u|=b$，$|\vec v|=a$，$|\vec u-\vec v|=|\overrightarrow{BA}|=c$。于是
\[
  c^2=|\vec u-\vec v|^2=|\vec u|^2+|\vec v|^2-2\,\vec u\cdot\vec v
  =a^2+b^2-2ab\cos C,
\]
反解即 $\cos C=\dfrac{a^2+b^2-c^2}{2ab}$。
\end{derivation}
同一个事实的两种语言：\key{余弦定理是数量积的模长读法，数量积是余弦定理的代数化}。
第 2 章用「复平面投影」理解 $\cos$，本章用「向量投影」理解 $\cos$——两种语言指向同一个对象。

\begin{example}
已知 $|\vec a|=4$，$|\vec b|=5$，夹角 $\theta=120^\circ$，求 $\vec a\cdot\vec b$、数量投影与投影向量。
\end{example}
\begin{solution}
\[
  \vec a\cdot\vec b=4\times5\times\cos120^\circ=20\times\left(-\frac12\right)=-10.
\]
数量投影 $=\dfrac{\vec a\cdot\vec b}{|\vec b|}=\dfrac{-10}{5}=-2$（负号说明投影逆着 $\vec b$）；
投影向量 $=\dfrac{\vec a\cdot\vec b}{|\vec b|^2}\vec b=\dfrac{-10}{25}\vec b=-\dfrac25\vec b$。
\end{solution}

\begin{conccard}[一级]{数量积与投影}{求模长、判垂直、求夹角、做投影；坐标运算；与余弦定理互译}{把投影向量分母写成 $|\vec b|$；认为 $\vec a\cdot\vec b$ 是向量；默认投影非负}
$\vec a\cdot\vec b=|\vec a||\vec b|\cos\theta=x_1x_2+y_1y_2$；
$\mathrm{proj}_{\vec b}\vec a=\dfrac{\vec a\cdot\vec b}{|\vec b|^2}\vec b$，数量投影 $=\dfrac{\vec a\cdot\vec b}{|\vec b|}$。
$\vec a\perp\vec b\iff\vec a\cdot\vec b=0$（非零向量）。
\end{conccard}

\begin{commonerror}{投影运算的三个高频失误}
\paragraph{失误一：分母多除一次模长。}
投影向量是 $\dfrac{\vec a\cdot\vec b}{|\vec b|^2}\vec b$。若误写 $\dfrac{\vec a\cdot\vec b}{|\vec b|}\vec b$，
左边是「数量投影（长度）再乘向量」，量纲变成长度的平方，必然出错。
正确链条：数量投影 $\dfrac{\vec a\cdot\vec b}{|\vec b|}$（一个数）$\to$ 乘单位向量 $\dfrac{\vec b}{|\vec b|}$ $\to$ 投影向量。
\paragraph{失误二：把数量积当成向量。}
$\vec a\cdot\vec b$ 是一个\key{实数}（可正可负可为零），不是向量；
更不能写成 $\vec a\cdot\vec b\,\vec b$ 这种不伦不类的式子。
\paragraph{失误三：默认投影非负、夹角非钝。}
$\theta\in[0,\pi]$，钝角时 $\cos\theta<0$，数量投影与数量积都是负数；
例如上例 $\theta=120^\circ$ 时数量投影为 $-2$。符号携带方向信息，不能丢。
\end{commonerror}

\subsection{极化恒等式与面积}

\key{探究三：极化恒等式——共起点三向量乘积的秒杀通道}

（恒等式的核心是两向量之积 $\overrightarrow{AB}\cdot\overrightarrow{AC}$；
下面的中点形式由它配方推出。）

\key{（1）一般的极化恒等式。}
对任意向量 $\vec u,\vec v$：
\[
  \vec u\cdot\vec v=\frac14\left(|\vec u+\vec v|^2-|\vec u-\vec v|^2\right).
\]
推导只需展开右边：$|\vec u+\vec v|^2-|\vec u-\vec v|^2=4\,\vec u\cdot\vec v$。

\key{（2）三角形中的中点版本。}
设 $M$ 是 $BC$ 的中点，则 $\overrightarrow{AB}+\overrightarrow{AC}=2\overrightarrow{AM}$，$\overrightarrow{AB}-\overrightarrow{AC}=\overrightarrow{CB}$，代入得
\begin{derivation}
\[
  \overrightarrow{AB}\cdot\overrightarrow{AC}
  =\frac14\left(|2\overrightarrow{AM}|^2-|\overrightarrow{CB}|^2\right)
  =|\overrightarrow{AM}|^2-\frac14|\overrightarrow{CB}|^2
  =|\overrightarrow{AM}|^2-|\overrightarrow{MB}|^2.
\]
\end{derivation}
这个恒等式把「\key{共起点的两个向量之积}」改写成「\key{动点到中点的距离平方} $-$ 定值」，
是处理动态最值问题的标准武器。

\begin{example}
已知圆 $O$ 的半径为 $2$，弦 $BC=2$，点 $A$ 在圆上运动，求 $\overrightarrow{AB}\cdot\overrightarrow{AC}$ 的最大值。
\end{example}
\begin{solution}
取 $M$ 为 $BC$ 中点，则 $|\overrightarrow{MB}|=1$ 为定值，由极化恒等式
\[
  \overrightarrow{AB}\cdot\overrightarrow{AC}=|\overrightarrow{AM}|^2-1.
\]
问题转化为求 $|\overrightarrow{AM}|$ 的最大值。由垂径定理 $|OM|=\sqrt{2^2-1^2}=\sqrt3$，
$A$ 在圆上运动时 $|AM|$ 的最大值为 $R+|OM|=2+\sqrt3$（$A$ 位于射线 $MO$ 的延长线与圆的交点）。
因此
\[
  \left(\overrightarrow{AB}\cdot\overrightarrow{AC}\right)_{\max}=(2+\sqrt3)^2-1=6+4\sqrt3.
\]
若直接设坐标硬算，将陷入三角换元；\key{极化恒等式一步就把「双动向量乘积」化为「单动点距离」。}
\end{solution}

\key{（3）面积形式：叉积的模。}
与数量积配对的是\key{叉积（外积）}：$\vec a\times\vec b$ 是垂直于 $\vec a,\vec b$ 所在平面的向量，
其模
\[
  |\vec a\times\vec b|=|\vec a|\,|\vec b|\,|\sin\theta|
\]
恰等于以 $\vec a,\vec b$ 为邻边的平行四边形面积；坐标计算公式为
\[
  |\vec a\times\vec b|=|x_1y_2-x_2y_1|,\qquad
  S_{\triangle ABC}=\frac12|\overrightarrow{AB}\times\overrightarrow{AC}|
  =\frac12\,|x_1y_2-x_2y_1|.
\]
对照记忆：\key{点积管「角与投影」（余弦），叉积管「面积与正弦」}，二者互补。

\begin{example}
已知 $\vec a=(3,1)$，$\vec b=(2,4)$，求以 $\vec a,\vec b$ 为邻边的三角形面积。
\end{example}
\begin{solution}
\[
  S=\frac12|x_1y_2-x_2y_1|=\frac12|3\times4-1\times2|=\frac12\times10=5.
\]
\end{solution}

注记（高观点）：把三个极化恒等式按三角形面积加权求和，就得到著名的奔驰定理
$S_{\triangle BOC}\cdot\overrightarrow{OA}+S_{\triangle COA}\cdot\overrightarrow{OB}+S_{\triangle AOB}\cdot\overrightarrow{OC}=\vec 0$
（$O$ 为 $\triangle ABC$ 内一点）——它是「重心坐标」的面积版；第 7 章统计回归的样本中心点与之同源。

\begin{conccard}[二级]{极化恒等式：共起点三向量乘积}{共起点两向量数量积的最值／范围问题，尤其是「一定弦、一动顶点」模型}{认不出中点结构，仍去建系硬算；把 $|\overrightarrow{MB}|$ 误当 $|BC|$；忽略动点到中点距离的取值范围需按圆／线段讨论}
$\overrightarrow{AB}\cdot\overrightarrow{AC}=|\overrightarrow{AM}|^2-|\overrightarrow{MB}|^2$（$M$ 为 $BC$ 中点）。
更一般地 $\vec u\cdot\vec v=\frac14(|\vec u+\vec v|^2-|\vec u-\vec v|^2)$。
配套面积公式：$S_{\triangle ABC}=\frac12|\overrightarrow{AB}\times\overrightarrow{AC}|=\frac12|x_1y_2-x_2y_1|$。
\end{conccard}

\subsection{空间向量：方向向量与法向量}

\key{探究四：用法向量统一处理空间中的平行、垂直与夹角}

\key{（1）空间向量坐标化。}
以 $O$ 为原点、$\vec i,\vec j,\vec k$ 为三条两两垂直的单位向量建立空间直角坐标系，则
\[
  \vec a=x\vec i+y\vec j+z\vec k\iff \vec a=(x,y,z),\qquad
  |\vec a|=\sqrt{x^2+y^2+z^2},
\]
\[
  \vec a\cdot\vec b=x_1x_2+y_1y_2+z_1z_2,\qquad
  \cos\langle\vec a,\vec b\rangle=\frac{\vec a\cdot\vec b}{|\vec a||\vec b|}.
\]
几何空间代数化的关键一步：\key{建系 $=$ 给每个点分配坐标，之后所有几何量都是坐标运算}。

\key{（2）平面的法向量。}
与平面 $\alpha$ 垂直的向量 $\vec n$ 叫做 $\alpha$ 的\key{法向量}。
一个平面有无穷多个法向量，但它们彼此平行（共线），所以\key{法向量的方向就是平面的「朝向」}。
若 $\vec n=(A,B,C)$，$P_0(x_0,y_0,z_0)\in\alpha$，则平面方程为
\[
  A(x-x_0)+B(y-y_0)+C(z-z_0)=0\quad\text{即}\quad Ax+By+Cz+D=0.
\]
求法向量的通用方法：找平面内两个不共线向量 $\vec a,\vec b$，令 $\vec n\cdot\vec a=0$、$\vec n\cdot\vec b=0$ 解方程组。

\key{（3）线面角与二面角。}
设直线 $l$ 的方向向量为 $\vec s$，平面 $\alpha$ 的法向量为 $\vec n$：
\[
  \textbf{线面角 }\varphi:\quad \sin\varphi=\frac{|\vec s\cdot\vec n|}{|\vec s|\,|\vec n|}
  \quad(\varphi\in[0,\tfrac{\pi}{2}]).
\]
注意是\key{正弦}：直线与法向量的夹角与线面角互余。
设两平面法向量为 $\vec n_1,\vec n_2$：
\[
  \textbf{二面角 }\theta:\quad |\cos\theta|=\frac{|\vec n_1\cdot\vec n_2|}{|\vec n_1|\,|\vec n_2|},
\]
锐钝要结合图形判断（或用法向量方向调整符号）。

\begin{figure}[htbp]
  \centering
  \begin{tikzpicture}[>=Stealth, semithick, font=\small]
    % 两个半平面（平行四边形示意），公共棱为中间那条水平线
    \fill[main!8] (-2.6,0)--(2.6,0)--(1.9,2.4)--(-3.3,2.4)--cycle;
    \fill[second!10] (-2.6,0)--(2.6,0)--(1.9,-2.0)--(-3.3,-2.0)--cycle;
    \draw[thick, main!60!black] (-2.6,0)--(2.6,0)--(1.9,2.4)--(-3.3,2.4)--cycle;
    \draw[thick, second!50!black] (-2.6,0)--(2.6,0)--(1.9,-2.0)--(-3.3,-2.0)--cycle;
    % 棱与棱上一点
    \draw[very thick, gray] (-2.6,0)--(2.6,0);
    \fill (0,0) circle (1.5pt) node[below left] {$P$};
    % 二面角的平面角：在两个半平面内分别作垂直于棱的射线
    \draw[->, third, very thick] (0,0)--(-1.1,1.9) node[above left] {$\vec u\subset\alpha$};
    \draw[->, third, very thick] (0,0)--(-1.1,-1.6) node[below left] {$\vec v\subset\beta$};
    % 直角符号：u\perp l，v\perp l（放大并贴近射线）
    \draw[third!70!black, thin] (0.4,0)--(0.2,0.3464)--(-0.2,0.3464);
    \draw[third!70!black, thin] (0.4,0)--(0.1736,-0.3296)--(-0.2264,-0.3296);
    % 平面角 \theta：弧线从 \vec u 跨过 \vec v
    \draw[main, thick] (0,0) ++(120:0.95) arc (120:235:0.95);
    \node[main] at (-1.15,0) {$\theta$};
    % 法向量
    \draw[->, second, thick] (1.2,1.2)--(2.2,2.6) node[above right] {$\vec n_1$};
    \draw[->, main, thick] (1.2,-1.0)--(2.2,-2.4) node[below right] {$\vec n_2$};
    \node[main!70!black] at (-2.6,1.7) {$\alpha$};
    \node[second!60!black] at (-2.6,-1.4) {$\beta$};
    \node[gray!70!black] at (2.5,0.28) {$l$};
    \node[gray!70!black, align=center] at (0,-3.05) {沿棱 $l$ 取点 $P$，在 $\alpha,\beta$ 内分别作垂直于 $l$ 的射线，\\两条射线的夹角即二面角的平面角};
  \end{tikzpicture}
  \caption{二面角的几何图：两平面 $\alpha,\beta$ 交于棱 $l$；$\vec u,\vec v$ 分别躺在两个半平面内且都垂直于 $l$（直角符号），夹角就是二面角的平面角 $\theta$。法向量 $\vec n_1,\vec n_2$ 分别垂直于两平面，其夹角与 $\theta$ 相等或互补，故 $|\cos\theta|=\dfrac{|\vec n_1\cdot\vec n_2|}{|\vec n_1||\vec n_2|}$，锐钝需结合图形判断}
  \label{fig:dihedral-angle}
\end{figure}

\begin{center}
\footnotesize
\setlength{\tabcolsep}{3pt}
\renewcommand{\arraystretch}{1.4}
\begin{tabular}{p{2.4cm}p{6.8cm}p{3.6cm}}
\toprule
位置关系 & 向量判据 & 说明 \\
\midrule
线线平行 & $\vec s_1\parallel\vec s_2$（即 $\vec s_1\times\vec s_2=\vec 0$） & 方向向量共线 \\
线线垂直 & $\vec s_1\cdot\vec s_2=0$ & 方向向量垂直 \\
线面平行 & $\vec s\cdot\vec n=0$ 且 $l\not\subset\alpha$ & 方向向量与法向量垂直 \\
线面垂直 & $\vec s\parallel\vec n$ & 方向向量与法向量共线 \\
面面平行 & $\vec n_1\parallel\vec n_2$ & 法向量共线 \\
面面垂直 & $\vec n_1\cdot\vec n_2=0$ & 法向量垂直 \\
\bottomrule
\end{tabular}
\end{center}

\begin{figure}[htbp]
  \centering
  \begin{tikzpicture}[>=Stealth, semithick, font=\footnotesize]
    \coordinate (O) at (0,0);
    \coordinate (X) at (-2.3,-1.05);
    \coordinate (Y) at (2.9,-0.75);
    \coordinate (XY) at (0.6,-1.8);
    % 平面 α（xOy 平面，z=0）
    \fill[third!8] (O)--(X)--(XY)--(Y)--cycle;
    \draw[third!50!black, thin] (O)--(X)--(XY)--(Y)--cycle;
    \node[third!70!black] at (2.9,-1.05) {$\alpha$};
    % 坐标轴
    \draw[->, gray] (O)--(X) node[below left] {$x$};
    \draw[->, gray] (O)--(Y) node[below right] {$y$};
    \draw[->, gray] (O)--(0,3.1) node[above] {$z$};
    % 法向量 n（竖直向上，垂直于平面）
    \draw[->, third, very thick] (-0.35,-0.5)--(-0.35,1.3) node[pos=0.85,right] {$\vec{n}$};
    % 直线方向向量 s 穿过平面
    \coordinate (R) at (1.15,-0.85);
    \coordinate (S) at (1.7,1.5);
    \draw[->, main!80!black, very thick] (R)--(S) node[midway,right] {$\vec{s}$};
    \draw[dashed, gray] (S)--(1.7,0);
    % 线面角 φ
    \coordinate (I) at (1.349,0);
    \draw (1.799,0) arc (0:77:0.45);
    \node at (1.98,0.4) {$\varphi$};
    \node[main!70!black] at (0.6,-2.35) {方向向量 $\vec s$、法向量 $\vec n$ 与线面角 $\varphi$};
  \end{tikzpicture}
  \caption{空间向量与法向量：平面 $\alpha$ 的法向量 $\vec n$ 指示平面朝向；直线方向向量 $\vec s$ 穿过平面，$\varphi$ 为 $\vec s$ 与其在平面内投影的夹角（线面角）}
  \label{fig:space-normal}
\end{figure}

\begin{example}
正方体 $ABCD-A_1B_1C_1D_1$ 的棱长为 $1$，求二面角 $A_1-BD-C$ 的余弦值。
\end{example}
\begin{solution}
以 $D$ 为原点建系：$D(0,0,0)$，$A(1,0,0)$，$C(0,1,0)$，$B(1,1,0)$，$A_1(1,0,1)$。
平面 $BDC$ 即底面 $z=0$，取其法向量 $\vec n_1=(0,0,1)$。
平面 $A_1BD$ 内有 $\overrightarrow{DB}=(1,1,0)$，$\overrightarrow{DA_1}=(1,0,1)$，
\[
  \vec n_2=\overrightarrow{DB}\times\overrightarrow{DA_1}
  =(1\times1-0\times0,\;0\times1-1\times1,\;1\times0-1\times1)=(1,-1,-1).
\]
于是
\[
  |\cos\theta|=\frac{|\vec n_1\cdot\vec n_2|}{|\vec n_1|\,|\vec n_2|}
  =\frac{|-1|}{1\times\sqrt3}=\frac{\sqrt3}{3}.
\]
取 $BD$ 中点 $M$，$\overrightarrow{MA_1}=(-0.5,-0.5,1)$、$\overrightarrow{MC}=(-0.5,-0.5,0)$ 均与 $BD$ 垂直，
两者点积为正，故二面角 $A_1-BD-C$ 为锐角，其余弦值为 $\dfrac{\sqrt3}{3}$。
传统几何法需要作平面角、证垂直、再解三角形；向量法只需建系、算两个法向量、代公式。
\end{solution}

\begin{example}
平面 $\alpha,\beta$ 的法向量分别为 $\vec n_1=(2,-1,1)$，$\vec n_2=(1,1,-1)$，直线 $l$ 的方向向量为 $\vec s=(1,2,2)$。
（1）判断 $l$ 与 $\alpha$ 的位置关系并求线面角；（2）判断 $\alpha$ 与 $\beta$ 的位置关系。
\end{example}
\begin{solution}
（1）$\vec s\cdot\vec n_1=1\times2+2\times(-1)+2\times1=2\neq0$，
所以 $l$ 与 $\alpha$ \key{相交}（不平行也不在面内）。线面角 $\varphi$ 满足
\[
  \sin\varphi=\frac{|\vec s\cdot\vec n_1|}{|\vec s|\,|\vec n_1|}
  =\frac{2}{\sqrt{1+4+4}\cdot\sqrt{4+1+1}}=\frac{2}{3\sqrt6}=\frac{\sqrt6}{9}.
\]
（2）$\vec n_1\cdot\vec n_2=2-1-1=0$，所以两平面的法向量垂直，即 \key{$\alpha\perp\beta$}。
\end{solution}

\begin{conccard}[二级]{法向量统一处理空间位置关系}{线面／面面的平行、垂直、夹角问题；点到平面距离；存在性问题}{把线面角公式错写成余弦；二面角不判断锐钝；建系时没说明三条轴两两垂直或坐标写错}
直线方向向量 $\vec s$、法向量 $\vec n$：
线面平行 $\iff\vec s\cdot\vec n=0$；线面垂直 $\iff\vec s\parallel\vec n$；
线面角 $\sin\varphi=\dfrac{|\vec s\cdot\vec n|}{|\vec s||\vec n|}$；
二面角 $|\cos\theta|=\dfrac{|\vec n_1\cdot\vec n_2|}{|\vec n_1||\vec n_2|}$。
（第 4 章立体几何主干直接调用这套判据与下面的易错清单，不再重复。）
\end{conccard}

\begin{commonerror}{空间向量建系的三个高频失误}
\paragraph{失误一：建系不交代垂直性。}
使用空间直角坐标系前，必须说明所选三条轴（或从同一点出发的三条棱）两两垂直；
若题目只给出斜角或没有垂直条件，应先证明垂直关系再建系。
\paragraph{失误二：线面角套错公式。}
线面角 $\varphi$ 与「方向向量和法向量的夹角」互余，所以
$\sin\varphi=\dfrac{|\vec s\cdot\vec n|}{|\vec s||\vec n|}$，而不是余弦。
把 $\vec s$ 换成它在面内的投影，才是余弦关系，切勿混用。
\paragraph{失误三：二面角不判断锐钝。}
$|\cos\theta|$ 只给大小；必须结合图形（或取棱上一点，用两个半平面内的向量验号）确定二面角是锐角还是钝角，否则一半情形会丢。
\end{commonerror}

\begin{advancedonly}
\key{超纲延伸一：换基底与相似律——同一线性变换，在不同基下「读数不同，变换不动」。}

设线性变换 $T$ 在标准基 $\{\vec e_1,\vec e_2\}$ 下的矩阵为
$A=\begin{pmatrix}2&1\\0&1\end{pmatrix}$（即 $T(\vec e_1)=(2,0)$，$T(\vec e_2)=(1,1)$）。
取 $\vec v=(2,1)$，则
\[
  T(\vec v)=A\vec v=\begin{pmatrix}2&1\\0&1\end{pmatrix}\begin{pmatrix}2\\1\end{pmatrix}
  =\tk{\begin{pmatrix}5\\1\end{pmatrix}}.
\]
换一组新基 $\mathcal B'=\{\vec u_1,\vec u_2\}$：$\vec u_1=(1,1)$，$\vec u_2=(-1,1)$，
过渡矩阵 $P=(\vec u_1\ \vec u_2)=\begin{pmatrix}1&-1\\1&1\end{pmatrix}$。
同一个向量换基后读数要左乘 $P^{-1}$；由 $\det P=2$ 得
$P^{-1}=\dfrac12\begin{pmatrix}1&1\\-1&1\end{pmatrix}$，于是
\[
  [\vec v]_{\mathcal B'}=P^{-1}\begin{pmatrix}2\\1\end{pmatrix}
  =\tk{\begin{pmatrix}\tfrac32\\[2pt]-\tfrac12\end{pmatrix}},\qquad
  [T(\vec v)]_{\mathcal B'}=P^{-1}\begin{pmatrix}5\\1\end{pmatrix}
  =\tk{\begin{pmatrix}3\\-2\end{pmatrix}}.
\]
而 $T$ 在新基下的矩阵不再是 $A$，它由相似律给出：
\[
  A'=P^{-1}AP
  =\frac12\begin{pmatrix}1&1\\-1&1\end{pmatrix}
   \begin{pmatrix}2&1\\0&1\end{pmatrix}
   \begin{pmatrix}1&-1\\1&1\end{pmatrix}
  =\tk{\begin{pmatrix}2&0\\-1&1\end{pmatrix}}.
\]
验证（一步）：用新矩阵乘新坐标，必须还原像向量的读数，
\[
  A'\,[\vec v]_{\mathcal B'}
  =\begin{pmatrix}2&0\\-1&1\end{pmatrix}
   \begin{pmatrix}\tfrac32\\[2pt]-\tfrac12\end{pmatrix}
  =\tk{\begin{pmatrix}3\\-2\end{pmatrix}}=[T(\vec v)]_{\mathcal B'}.
\]
\key{$A'=P^{-1}AP$：换基底只是换一把尺子，坐标读数变了，变换本身纹丝不动。}
顺带一提：$\det A=\det A'=2$——相似矩阵的不变量，正是「换个角度，本质仍在」。
\end{advancedonly}

\begin{advancedonly}
\key{超纲延伸二：镜面反射公式——投影走多远，就回头走多远。}

以 $\vec b$ 所在的整条直线为镜面（不只是向量 $\vec b$ 那一截），
$\vec a$ 的反射像由投影直接写出：
\[
  \vec a'=2\,\mathrm{proj}_{\vec b}\vec a-\vec a.
\]
取 $\vec a=(2.5,1.5)$，$\vec b=(3,1)$。先投影：
\[
  \vec a\cdot\vec b=2.5\times3+1.5\times1=\tk{9},\qquad
  |\vec b|^2=3^2+1^2=\tk{10},
\]
\[
  \mathrm{proj}_{\vec b}\vec a=\frac{\vec a\cdot\vec b}{|\vec b|^2}\,\vec b
  =\frac{9}{10}(3,1)=\tk{(2.7,0.9)}.
\]
再越过镜面走同样远：
\[
  \vec a'=2(2.7,0.9)-(2.5,1.5)=\tk{(2.9,0.3)}.
\]
几何解释：垂足 $\mathrm{proj}_{\vec b}\vec a$ 恰好是 $\vec a$ 与 $\vec a'$ 连线的中点，
也就是说镜面是这条连线的垂直平分线——从 $\vec a$ 到镜面有多远，像就退到另一侧多远。

\begin{figure}[htbp]
  \centering
  \begin{tikzpicture}[>=Stealth, semithick, font=\footnotesize]
    \coordinate (O) at (0,0);
    \coordinate (A) at (2.5,1.5);
    \coordinate (B) at (3,1);
    \coordinate (F) at (2.7,0.9);
    \coordinate (Ap) at (2.9,0.3);
    % 镜面：\vec b 所在的整条直线（双向延长）
    \draw[dashed, second, thick] (-1.5,-0.5)--(4.2,1.4);
    \node[second, below left] at (-1.5,-0.5) {镜面};
    % 坐标轴
    \draw[->, gray] (-1.0,0)--(4.6,0);
    \draw[->, gray] (0,-0.9)--(0,2.6);
    \fill (O) circle (1pt);
    % 投影向量（恰好躺在镜面上）与垂足
    \draw[->, third, very thick] (O)--(F);
    \fill[third] (F) circle (1.5pt);
    \node[third] at (1.55,0.30) {$\mathrm{proj}_{\vec b}\vec a$};
    % 原向量、像向量与镜面方向向量
    \draw[->, main!80!black, very thick] (O)--(A) node[above left] {$\vec a=(2.5,1.5)$};
    \draw[->, second, very thick] (O)--(B) node[above right] {$\vec b=(3,1)$};
    \draw[->, main, very thick] (O)--(Ap) node[right] {$\vec a'=(2.9,0.3)$};
    % 垂线段与垂足处的直角符号
    \draw[dashed, gray] (A)--(F)--(Ap);
    \draw[third!70!black, thin] (2.85,0.95)--(2.80,1.10)--(2.65,1.05);
  \end{tikzpicture}
  \caption{镜面反射：以 $\vec b$ 所在直线（虚线）为镜面，$\mathrm{proj}_{\vec b}\vec a$ 是 $\vec a$ 的垂足，$\vec a'=2\,\mathrm{proj}_{\vec b}\vec a-\vec a$；垂足同时也是 $\vec a$ 与 $\vec a'$ 连线的中点}
  \label{fig:mirror-reflection}
\end{figure}
\end{advancedonly}

\begin{method}
向量解题三步法：
\begin{enumerate}
  \item \key{翻译}：把几何条件（中点、平行、垂直、夹角、面积）译成向量语言；
        不确定从哪下手时，先选一组基底或建立坐标系。有垂直关系优先建系。
  \item \key{运算}：坐标化或基底化之后，用数量积处理长度与夹角，
        用「中点公式 $+$ 极化恒等式」处理共起点向量乘积的最值，
        用方向向量与法向量处理空间位置关系；追求「选得巧，计算就少」。
  \item \key{回译}：把代数结果翻回几何结论，并检查退化情形
        （零向量、两向量共线、点在平面内、分母为零）。
\end{enumerate}
\end{method}

\begin{summary}
本章把向量当作「既有大小又有方向的量」$\to$「点集的位移」$\to$「基底下的坐标」$\to$「内积投影算子」$\to$「空间中的方向与法向」层层展开，凸显向量体系相对三角与函数的独立地位：它是几何空间代数化的通用线性语言。
平面向量基本定理保证任一向量被不共线基底唯一分解，坐标就是有序实数对；数量积 $\vec a\cdot\vec b=|\vec a||\vec b|\cos\theta$ 的几何意义是「投影乘模长」，投影向量 $\mathrm{proj}_{\vec b}\vec a=\frac{\vec a\cdot\vec b}{|\vec b|^2}\vec b$，它与余弦定理互为代数化与几何读法；
极化恒等式 $\overrightarrow{AB}\cdot\overrightarrow{AC}=|\overrightarrow{AM}|^2-|\overrightarrow{MB}|^2$ 把共起点三向量乘积化为动点到中点的距离平方差，是动态最值的秒杀通道，面积则由叉积的模 $S=\frac12|\vec a\times\vec b|$ 给出；
空间向量坐标化之后，平面的法向量成为统一语言：线面平行／垂直、线面角、二面角全部化为方向向量与法向量的数量积运算。
\end{summary}

\end{explore}

\begin{inclasstest}
\begin{enumerate}
  \item 已知 $\vec a=(1,2)$，$\vec b=(3,-1)$，求 $\vec a\cdot\vec b$、$|\vec a|$ 以及 $\vec a$ 在 $\vec b$ 上的投影向量；\level{易}
  \item 已知 $\vec e_1,\vec e_2$ 不共线，$\vec m=\vec e_1+2\vec e_2$，$\vec n=3\vec e_1+k\vec e_2$，若 $\vec m\parallel\vec n$，求实数 $k$；\level{易}
  \item 半径为 $2$ 的圆中，弦 $BC=2$，点 $A$ 在圆上运动，求 $\overrightarrow{AB}\cdot\overrightarrow{AC}$ 的取值范围；\level{中}
  \item 平面 $\alpha$ 的法向量 $\vec n_1=(1,2,-1)$，平面 $\beta$ 的法向量 $\vec n_2=(2,1,1)$，判断 $\alpha$ 与 $\beta$ 的位置关系，并求所成锐二面角的余弦值。\level{中}
\end{enumerate}
\end{inclasstest}

\begin{afterclass}
\subsection*{基础巩固}
\begin{enumerate}
  \item 已知 $\vec a=(2,-1)$，$\vec b=(-3,4)$，求 $\vec a+\vec b$、$\vec a\cdot\vec b$ 与两向量夹角的余弦值；
  \item 设 $\vec e_1,\vec e_2$ 不共线，$\vec a=\vec e_1+2\vec e_2$。若换用新基底 $\vec e_1'=\vec e_1+\vec e_2$，$\vec e_2'=\vec e_1-\vec e_2$，求 $\vec a$ 在新基底下的坐标，并说明坐标为什么随基底改变而改变。
\end{enumerate}

\subsection*{进阶提升}
\begin{enumerate}
  \item 在 $\triangle ABC$ 中，$AB=3$，$AC=5$，$BC=4$，用极化恒等式求 $\overrightarrow{AB}\cdot\overrightarrow{AC}$；
  \item 三棱锥 $D-ABC$ 中，$A(1,0,0)$，$B(0,1,0)$，$C(0,0,1)$，$D(0,0,0)$，求二面角 $A-BC-D$ 的余弦值。
\end{enumerate}

\subsection*{拓展探究}
\begin{enumerate}
  \item 用向量方法证明：三角形三条中线交于一点，且该点分每条中线为 $2:1$（提示：先用平面向量基本定理比较系数求交点分比）。
\end{enumerate}

\subsection*{教师版答案}
\begin{solution}
\textbf{基础 1：}$\vec a+\vec b=(2-3,\,-1+4)=(-1,3)$；
$\vec a\cdot\vec b=2\times(-3)+(-1)\times4=-6-4=-10$；
\[
  |\vec a|=\sqrt{2^2+(-1)^2}=\sqrt5,\quad |\vec b|=\sqrt{(-3)^2+4^2}=5,
\]
\[
  \cos\theta=\frac{\vec a\cdot\vec b}{|\vec a||\vec b|}=\frac{-10}{5\sqrt5}=-\frac{2\sqrt5}{5}.
\]
\end{solution}
\begin{solution}
\textbf{基础 2：}设 $\vec a=\lambda\vec e_1'+\mu\vec e_2'$，即
\[
  \vec a=\lambda(\vec e_1+\vec e_2)+\mu(\vec e_1-\vec e_2)
  =(\lambda+\mu)\vec e_1+(\lambda-\mu)\vec e_2.
\]
与 $\vec a=\vec e_1+2\vec e_2$ 比较系数（基本定理保证唯一）：
\[
  \lambda+\mu=1,\quad \lambda-\mu=2\implies\lambda=\frac32,\ \mu=-\frac12.
\]
故新坐标为 $\left(\dfrac32,-\dfrac12\right)$。
坐标是「向量在基底下的读数」，基底变了读数自然改变；变的只是坐标，向量本身不变。
\end{solution}
\begin{solution}
\textbf{进阶 1：}取 $BC$ 中点 $M$，$|MB|=2$，中线长
\[
  |AM|^2=\frac{2AB^2+2AC^2-BC^2}{4}
  =\frac{2\times9+2\times25-16}{4}=\frac{52}{4}=13.
\]
由极化恒等式
\[
  \overrightarrow{AB}\cdot\overrightarrow{AC}=|\overrightarrow{AM}|^2-|\overrightarrow{MB}|^2=13-4=9.
\]
（检验：$\cos A=\frac{9+25-16}{2\times3\times5}=\frac{3}{5}$，故 $\overrightarrow{AB}\cdot\overrightarrow{AC}=3\times5\times\frac35=9$。）
\end{solution}
\begin{solution}
\textbf{进阶 2：}平面 $ABC$：由 $A,B,C$ 的坐标知方程为 $x+y+z=1$，法向量 $\vec n_1=(1,1,1)$；
平面 $DBC$ 过原点且含 $B,C$，方程为 $x=0$，法向量 $\vec n_2=(1,0,0)$。
\[
  |\cos\theta|=\frac{|\vec n_1\cdot\vec n_2|}{|\vec n_1||\vec n_2|}
  =\frac{1}{\sqrt3\times1}=\frac{\sqrt3}{3}.
\]
取 $BC$ 中点 $M(0,\tfrac12,\tfrac12)$，$\overrightarrow{MA}=(1,-\tfrac12,-\tfrac12)$、$\overrightarrow{MD}=(0,-\tfrac12,-\tfrac12)$ 均与 $\overrightarrow{BC}=(0,1,-1)$ 垂直，且点积为正，故二面角 $A-BC-D$ 为锐角，余弦值为 $\dfrac{\sqrt3}{3}$。
\end{solution}
\begin{solution}
\textbf{拓展：}取 $A$ 为原点，记 $\overrightarrow{AB}=\vec b$，$\overrightarrow{AC}=\vec c$。
设 $D,E$ 分别为 $BC,CA$ 的中点，则
\[
  \overrightarrow{AD}=\frac{\vec b+\vec c}{2},\qquad
  \overrightarrow{BE}=\overrightarrow{AE}-\overrightarrow{AB}=\frac{\vec c}{2}-\vec b.
\]
设中线 $AD$ 与 $BE$ 交于点 $P$，令 $\overrightarrow{AP}=t\,\overrightarrow{AD}=\dfrac{t}{2}\vec b+\dfrac{t}{2}\vec c$；
另一方面 $\overrightarrow{AP}=\overrightarrow{AB}+s\,\overrightarrow{BE}=(1-s)\vec b+\dfrac{s}{2}\vec c$。
由平面向量基本定理比较系数：
\[
  \frac t2=1-s,\qquad \frac t2=\frac s2
  \implies s=t,\ \frac t2=1-t\implies t=\frac23.
\]
故 $AP:PD=BP:PE=2:1$。由对称性，$P$ 也在中线 $CF$（$F$ 为 $AB$ 中点）上，即三条中线共点，且该点（重心）分每条中线为 $2:1$。
这正是解析几何与统计中「均值点」的向量面孔（第 7 章样本中心与之同源）。
\end{solution}
\reading{拓展阅读 09}{向量：几何空间的通用线性语言}
\end{afterclass}

% 高观点导学案：立体几何与圆锥曲线大一统（第 4 章·高观点统一版）
% 主线：仿射变换（圆压缩为椭圆）→ 三维圆锥截面（Dandelin）→ 考纲安全落地双层法 → 解析几何主干 → 立体几何主干（向量法统一）
% 核心思想：圆锥曲线不是三类孤立的二次曲线，而是同一个三维对象（圆锥面）与不同倾角平面的相遇方式；
%           高观点负责「预判答案」，考纲通法负责「写出满分」。

\section{立体几何与圆锥曲线大一统（高观点版）}

\begin{methodbox}{多框架对读}
\small
\begin{center}
\begin{tabular}{c|p{2.9cm}|p{5.0cm}|p{3.1cm}}
\toprule
视角 & 本章看点 & 别的视角怎么说 & 哪侧最省力 \\
\midrule
代数 & 标准方程、判别式、弦长 & 坐标满足的条件 & 硬算与判别 \\
几何 & 轨迹形状、仿射变换、二面角 & 图形的位置与度量 & 猜结论、估范围 \\
集合 & 曲线＝点集；平行垂直＝关系 & 条件即约束的集合 & 界存在与唯一 \\
映射 & 仿射把圆映成椭圆；投影降维 & 换表示＝换坐标系 & 整体保持性质 \\
\bottomrule
\end{tabular}
\end{center}
卡的使用：解析几何先问「这是哪个变换下的不变量」，再决定代数硬算还是几何巧解（总图见第 0 章图~\ref{fig:four-view-map}）。
\end{methodbox}

第 3 章把向量确立为「几何空间代数化」的通用语言；本章再向前一步，打通立体几何与解析几何的维度壁垒：
椭圆不过是圆在\key{仿射变换}下的像，圆锥曲线不过是圆锥面在不同倾角平面上的\key{截线}——
维度不是障碍，变换才是本质。五个模块层层展开：
\key{仿射变换：从圆到椭圆的降维打击}（圆沿一个方向均匀压缩即得椭圆，共线、平行、中点、相切、面积比全部保持）$\to$
\key{Dandelin 视角}（椭圆／双曲线／抛物线是同一个圆锥面被不同倾角平面所截的结果）$\to$
\key{考纲安全落地的双层法}（高观点 10 秒预设定点定值，考纲通法写满分步骤）$\to$
\key{解析几何主干}（定义、标准方程、直线与圆锥曲线的位置关系、定点定值通法）$\to$
\key{立体几何主干}（线面关系、空间角与距离，向量法全面落地）。
核心思想：\key{圆锥曲线不是三类孤立的二次曲线，而是同一个三维对象与不同平面的相遇方式}。

\begin{learnobjective}
\begin{itemize}
  \item 理解仿射变换 $T:(x,y)\mapsto\left(x,\frac{b}{a}y\right)$，会把椭圆 $\frac{x^2}{a^2}+\frac{y^2}{b^2}=1$ 看作圆 $x^2+y^2=a^2$ 的压缩像；
  \item 掌握仿射不变量清单（共线、平行、中点、相切、面积比 $S_{\text{椭圆}}=\frac{b}{a}S_{\text{圆}}$），会用它秒算椭圆内接三角形面积最值；
  \item 知道圆锥曲线的 Dandelin 球证法：焦点—准线定义的统一，$e=\frac{\cos\beta}{\cos\alpha}$；
  \item 掌握「高观点预判 $+$ 考纲通法」双层解题策略：极点极线预设定点定值，设线—联立—韦达—判别式写满分步骤；
  \item 掌握椭圆／双曲线／抛物线的定义与标准方程，直线与圆锥曲线位置关系的判别式、弦长与中点弦公式，以及定点定值问题的通法步骤；
  \item 掌握线面平行／垂直的判定与性质、空间角（线面角、二面角）与点到平面距离，会用向量法统一处理（与第 3 章衔接）。
\end{itemize}
\end{learnobjective}

\begin{preview}
\subsection*{知识回顾}
\begin{itemize}
  \item 圆的标准方程是 \blank；把圆沿 $y$ 轴按比例 $\frac{b}{a}$ 均匀压缩后得到的图形是 \blank；
  \item 椭圆上任意一点到两焦点的距离之和等于 \blank；双曲线则对应距离之差的绝对值等于 \blank；
  \item 直线与一元二次方程联立后，判别式 $\Delta>0$ 表示有 \blank 个不同实根（交点）；
  \item 直线与平面平行的判定定理：平面外一条直线与此平面内一条直线 \blank，则线面平行。
\end{itemize}

\subsection*{教材速读}
阅读「圆锥曲线」与「立体几何」两章教材，完成下列问题：
\begin{enumerate}
  \item 椭圆定义要求 $2a$ \blank $2c$（填 $>$ 或 $<$）；双曲线要求 $2a$ \blank $2c$；
  \item 抛物线 $y^2=2px\,(p>0)$ 的焦点坐标是 \blank，准线方程是 \blank；
  \item 弦长公式 $|AB|=\sqrt{1+k^2}\cdot$ \blank；
  \item 点到平面的距离可用等体积法或 \blank 法求解。
\end{enumerate}
\end{preview}

\begin{explore}

\subsection{仿射变换：从圆到椭圆的降维打击}

\key{探究一：把椭圆「压」回圆，把难题「降」回熟题}

\key{（1）一个变换看出椭圆的身世。}
设圆 $C:x^2+y^2=a^2$。作变换
\[
  T:\begin{pmatrix}x\\y\end{pmatrix}\longmapsto\begin{pmatrix}x\\ \frac{b}{a}y\end{pmatrix},\qquad 0<b\le a.
\]
它把每一点沿 $y$ 轴方向按比例 $\frac{b}{a}$ \key{均匀压缩}：水平方向不动，竖直方向整体变扁。
设 $(x,y)\in C$，记 $(X,Y)=T(x,y)$，则 $x=X,\ y=\frac{a}{b}Y$，代入圆方程得
\[
  X^2+\left(\frac{a}{b}Y\right)^2=a^2
  \iff \frac{X^2}{a^2}+\frac{Y^2}{b^2}=1.
\]
于是 \key{椭圆就是圆在仿射变换 $T$ 下的像}。$T$ 是线性变换（矩阵 $\mathrm{diag}\left(1,\frac{b}{a}\right)$），所以它仍然把直线变成直线——这是下面所有「保持」的根源。

\key{（2）仿射不变量清单。}
因为 $T$ 把直线变成直线、且保持线段上的分比，以下事实全部保持：
\begin{itemize}
  \item \key{共线性}：三点共线 $\xrightarrow{T}$ 三点共线；
  \item \key{平行性}：$m\parallel m'\iff T(m)\parallel T(m')$；
  \item \key{中点}：线段的中点变为像线段的中点（$T$ 保持定比分点）；
  \item \key{相切}：直线与圆锥曲线相切（联立后 $\Delta=0$）的关系保持，因为 $T$ 可逆且把直线变直线；
  \item \key{面积}：任意平面图形的面积都乘以同一常数 $|\det T|=\frac{b}{a}$，故
  \[
    S_{\text{椭圆}}=\frac{b}{a}\,S_{\text{圆}}.
  \]
\end{itemize}
\key{面积比与图形的具体形状无关}，这是仿射变换最锋利的性质。

\begin{figure}[htbp]
  \centering
  \begin{tikzpicture}[>=Stealth, semithick, font=\footnotesize]
    % ---------- 左：圆与内接正三角形 ----------
    \draw[->, gray] (-2.5,0)--(2.7,0) node[below right] {$x$};
    \draw[->, gray] (0,-2.3)--(0,2.6) node[above left] {$y$};
    \fill[second!10] (0,0) circle (2);
    \draw[main!70!black, very thick] (0,0) circle (2);
    \fill[third!18] (0,2)--(-1.732,-1)--(1.732,-1)--cycle;
    \draw[third!70!black, very thick] (0,2)--(-1.732,-1)--(1.732,-1)--cycle;
    \fill (0,0) circle (1pt) node[below left] {$O$};
    \node[gray!70!black] at (1.15,-0.32) {$a$};
    \node[second!70!black] at (0,-2.02) {圆 $x^2+y^2=a^2$：$S_{\max}=\dfrac{3\sqrt3}{4}a^2$};
    % ---------- 右：压缩后的椭圆 ----------
    \begin{scope}[xshift=6.9cm]
      \draw[->, gray] (-2.5,0)--(2.7,0) node[below right] {$x$};
      \draw[->, gray] (0,-2.3)--(0,2.6) node[above left] {$y$};
      \fill[second!10] (0,0) ellipse (2 and 1.2);
      \draw[main!70!black, very thick] (0,0) ellipse (2 and 1.2);
      \fill[third!18] (0,1.2)--(-1.732,-0.6)--(1.732,-0.6)--cycle;
      \draw[third!70!black, very thick] (0,1.2)--(-1.732,-0.6)--(1.732,-0.6)--cycle;
      \fill (0,0) circle (1pt) node[below left] {$O$};
      \node[gray!70!black] at (1.15,-0.32) {$a$};
      \node[gray!70!black] at (-0.45,1.2) {$b$};
      \node[second!70!black] at (0,-2.02) {椭圆：$S_{\max}=\dfrac{b}{a}S_{\text{圆}}=\dfrac{3\sqrt3}{4}ab$};
    \end{scope}
    % ---------- 中间：压缩映射 ----------
    \draw[->, main!80!black, very thick] (3.0,2.0)--(3.7,2.0)
      node[midway,above] {$\displaystyle\binom{x}{y}\mapsto\binom{x}{\frac{b}{a}y}$};
  \end{tikzpicture}
  \caption{仿射压缩：左图为圆 $x^2+y^2=a^2$（取 $a=2$）及其内接正三角形；右图为沿 $y$ 轴按比例 $\frac{b}{a}=0.6$ 压缩后的椭圆 $\frac{x^2}{a^2}+\frac{y^2}{b^2}=1$（$b=1.2$），内接正三角形同步变形为内接三角形，面积同乘 $\frac{b}{a}$}
  \label{fig:affine-compress}
\end{figure}

\key{（3）降维打击：面积最值一步出答案。}
因为 $T$ 把所有三角形的面积同乘一个正常数，\key{「面积最大」在变换下保持}：
椭圆 $\frac{x^2}{a^2}+\frac{y^2}{b^2}=1$ 的内接三角形中，面积最大者正是圆 $x^2+y^2=a^2$ 的内接正三角形在 $T$ 下的像。
圆的内接正三角形面积为 $\frac{3\sqrt3}{4}a^2$，于是
\[
  S_{\max}=\frac{b}{a}\cdot\frac{3\sqrt3}{4}a^2=\frac{3\sqrt3}{4}ab.
\]
一个只含 $a,b$ 的干净答案——选填题可直接秒答。

\begin{example}
求椭圆 $E:\dfrac{x^2}{4}+\dfrac{y^2}{3}=1$ 的内接三角形面积的最大值。
\end{example}
\begin{solution}
这里 $a=2,\ b=\sqrt3$，压缩比为 $\frac{b}{a}=\frac{\sqrt3}{2}$。对应的圆半径为 $2$，其内接正三角形面积为
\[
  S_{\text{圆}}=\frac{3\sqrt3}{4}\cdot 2^2=3\sqrt3.
\]
故
\[
  S_{\max}=\frac{b}{a}S_{\text{圆}}=\frac{\sqrt3}{2}\cdot 3\sqrt3=\frac92
  \quad\left(\text{或直接用 }\frac{3\sqrt3}{4}ab=\frac{3\sqrt3}{4}\cdot2\cdot\sqrt3=\frac92\right).
\]
若按常规方法设三个动点坐标硬算，将陷入多元三角换元；\key{仿射变换一步就把「陌生题」变回「课本题」}。
\end{solution}

\begin{conccard}[一级]{仿射变换：椭圆是圆的压缩像}{椭圆内接三角形面积最值、共线／平行／中点类的比例问题、面积比的快速换算}{认为仿射变换保持角度与垂直；把「面积比恒定」误用于周长比；压缩后忘记新图形是椭圆而非相似形}
$T:(x,y)\mapsto\left(x,\frac{b}{a}y\right)$ 把圆 $x^2+y^2=a^2$ 变为椭圆 $\frac{x^2}{a^2}+\frac{y^2}{b^2}=1$；
共线、平行、中点、相切保持，面积同乘 $|\det T|=\frac{b}{a}$，故 $S_{\text{椭圆}}=\frac{b}{a}S_{\text{圆}}$，$S_{\triangle\max}=\frac{3\sqrt3}{4}ab$。
\end{conccard}

\begin{commonerror}{仿射变换「能保持什么、不能保持什么」}
\paragraph{失误一：以为角度也保持。}
仿射变换 $T$ 只保证「直线 $\to$ 直线、分比不变」，\key{不保持角度、不保持垂直、不保持长度}。
例如圆内接正三角形的一个内角是 $60^\circ$，压缩成椭圆的内接三角形后三个角一般不再是 $60^\circ$。
所以「椭圆 $\xrightarrow{\text{仿射}}$ 圆」可以搬用共线、平行、中点、面积比，却\textbf{不能}搬用直角与角度关系。
\paragraph{失误二：把面积比套到周长比上。}
周长在 $y$ 方向被压缩后会改变（弧长的积分并不同乘一个常数），只有\key{面积}这类二维测度才同乘 $|\det T|$。
\paragraph{失误三：忽略 $b>a$ 的情形。}
若沿 $y$ 轴拉伸（$\frac{b}{a}>1$），同样的推导给出 $\frac{x^2}{a^2}+\frac{y^2}{b^2}=1$（$b>a$），焦点在 $y$ 轴上。
公式 $S_{\max}=\frac{3\sqrt3}{4}ab$ 对两种情形都成立，但焦点位置要看 $a,b$ 谁大。
\end{commonerror}

\subsection{Dandelin 视角：圆锥曲线是三维截线}

\key{探究二：把解析几何的二次曲线放回立体几何的三维空间}

\key{（1）圆锥面与截面倾角。}
以 $z$ 轴为中心轴、半顶角为 $\alpha$ 的圆锥面方程为
\[
  x^2+y^2=z^2\tan^2\alpha .
\]
用一个平面去截它，截线的类型只取决于平面与中心轴的夹角 $\beta$（取锐角）：
\begin{itemize}
  \item $\beta>\alpha$：平面较「平」，只切断一侧锥面，得\key{椭圆}（$\beta=90^\circ$ 时退化为圆）；
  \item $\beta=\alpha$：平面恰好与一条母线平行，得\key{抛物线}；
  \item $\beta<\alpha$：平面较「陡」，同时切断两侧锥面，得\key{双曲线}。
\end{itemize}
并且离心率由这两个角完全决定：
\[
  e=\frac{\cos\beta}{\cos\alpha}.
\]
检验：$\beta=90^\circ\Rightarrow e=0$（圆）；$\beta=\alpha\Rightarrow e=1$（抛物线）；$\beta<\alpha\Rightarrow\cos\beta>\cos\alpha\Rightarrow e>1$（双曲线）。
\key{三类圆锥曲线从此不是三个定义，而是同一个对象的三种截法。}

\key{（2）Dandelin 球：把「和／差为定值」做成肉眼可见。}
在截面与锥面围成的空间内放入与圆锥面相切、又与截面相切的球（Dandelin 球），它们与截面的切点就是焦点（椭圆、双曲线用两个球，抛物线用一个球）。
设 $P$ 为截线上任一点，过 $P$ 的母线先后穿过两个切圆所在的平面。由「从一点向球引两条切线，切线长相等」：
\[
  |PF_1|=P\text{ 到第一个切圆的母线段长},\qquad |PF_2|=P\text{ 到第二个切圆的母线段长}.
\]
两段之和（椭圆）或之差（双曲线）恰为两切圆之间的母线长——这是一个与 $P$ 无关的定值。于是
\[
  \text{椭圆：}\ |PF_1|+|PF_2|=\text{定值};\qquad
  \text{双曲线：}\ \bigl||PF_1|-|PF_2|\bigr|=\text{定值}.
\]
\key{（3）统一的焦点—准线关系。}
抛物线情形只有一个球，但它已经足以给出统一定义。设 $\sigma$ 为球与圆锥的切圆所在平面，$l=\pi\cap\sigma$ 即准线（$\pi$ 为截面），$\delta$ 为 $\pi$ 与 $\sigma$ 的夹角。
设 $P$ 为截线上任一点，$h$ 为 $P$ 到 $\sigma$ 的距离，$T$ 为过 $P$ 的母线与切圆的交点。由「从一点向球引两条切线，切线长相等」：
\[
  |PF|=|PT|=\frac{h}{\cos\alpha},
  \qquad
  d(P,l)=\frac{h}{\sin\delta}.
\]
两式相除得
\[
  \frac{|PF|}{d(P,l)}=\frac{\sin\delta}{\cos\alpha}.
\]
而 $\sigma\perp$ 轴线，故 $\sin\delta=\cos\beta$，代入即得
\[
  \frac{|PF|}{d(P,l)}=\frac{\cos\beta}{\cos\alpha}=e.
\]
对椭圆、双曲线，把同一计算分别用到两个 Dandelin 球上，每一对「焦点—准线」都满足同一比值 $e$。
一般地，三种曲线统一为 \key{焦点—准线定义}：到定点（焦点）与定直线（准线）的距离之比为常数 $e$ 的点的轨迹，$e<1$ 为椭圆，$e=1$ 为抛物线，$e>1$ 为双曲线。
这正是「立体几何中的球与平面」给「解析几何中的二次曲线」发的身份证。

\begin{figure}[htbp]
  \centering
  \begin{tikzpicture}[>=Stealth, semithick, font=\footnotesize]
    % ================= 面板一：椭圆 =================
    \begin{scope}[xshift=0cm]
      \draw[dashed, gray!60] (0,-2.55)--(0,2.55);
      \draw[gray!70, thin] (0,0)--(1.162,2.492) (0,0)--(-1.162,2.492)
                         (0,0)--(1.162,-2.492) (0,0)--(-1.162,-2.492);
      \draw[second, very thick] (-2.068,0.685)--(2.068,1.415);
      \node[second!70!black] at (1.65,0.95) {截面};
      \draw[main!45!black, fill=main!8] (0,0.7347) circle (0.3106);
      \fill (0,0) circle (1pt);
      \fill[main] (-0.054,1.041) circle (1.2pt);
      \fill[main] (0.135,1.074) circle (1.2pt);
      % 两个焦点过近，合并为一个标注，并用引导线指出，避免标签互相叠字或压轴
      \node[main] at (1.05,0.45) {$\scriptstyle F_1,F_2$};
      \draw[main!55, thin] (0.78,0.52)--(0.18,1.0);
      \draw[main!80!black, thick] (1.62,2.42) ellipse (0.42 and 0.20);
      \node[main!70!black] at (0,-2.95) {$\beta>\alpha\Rightarrow$ 椭圆$\ (e<1)$};
    \end{scope}
    % ================= 面板二：抛物线 =================
    \begin{scope}[xshift=4.9cm]
      \draw[dashed, gray!60] (0,-2.55)--(0,2.55);
      \draw[gray!70, thin] (0,0)--(1.162,2.492) (0,0)--(-1.162,2.492)
                         (0,0)--(1.162,-2.492) (0,0)--(-1.162,-2.492);
      \draw[second, very thick] (-0.465,0.503)--(0.740,3.086);
      \node[second!70!black] at (0.72,0.62) {截面};
      \draw[main!45!black, fill=main!8] (0,0.75) circle (0.317);
      \fill (0,0) circle (1pt);
      \fill[main] (-0.287,0.884) circle (1.2pt) node[above] {$\scriptstyle F$};
      \draw[main!80!black, thick] (1.28,2.62) parabola bend (1.68,2.34) (2.08,2.62);
      \node[main!70!black] at (0,-2.95) {$\beta=\alpha\Rightarrow$ 抛物线$\ (e=1)$};
    \end{scope}
    % ================= 面板三：双曲线 =================
    \begin{scope}[xshift=9.8cm]
      \draw[dashed, gray!60] (0,-2.55)--(0,2.55);
      \draw[gray!70, thin] (0,0)--(1.162,2.492) (0,0)--(-1.162,2.492)
                         (0,0)--(1.162,-2.492) (0,0)--(-1.162,-2.492);
      \draw[second, very thick] (-0.330,-0.671)--(0.330,2.871);
      \node[second!70!black] at (0.42,2.55) {截面};
      \draw[main!45!black, fill=main!8] (0,0.3495) circle (0.1477);
      \draw[main!45!black, fill=main!8] (0,-0.8370) circle (0.3537);
      \fill (0,0) circle (1pt);
      \fill[main] (-0.026,0.495) circle (1.2pt) node[left] {$\scriptstyle F_1$};
      \fill[main] (-0.061,-0.489) circle (1.2pt) node[left] {$\scriptstyle F_2$};
      \draw[main!80!black, thick] (1.30,2.10) .. controls (1.78,2.32) and (1.78,2.52) .. (1.30,2.74);
      \draw[main!80!black, thick] (2.02,2.10) .. controls (1.54,2.32) and (1.54,2.52) .. (2.02,2.74);
      \node[main!70!black] at (0,-2.95) {$\beta<\alpha\Rightarrow$ 双曲线$\ (e>1)$};
    \end{scope}
  \end{tikzpicture}
  \caption{Dandelin 示意（轴截面）：同一个圆锥面被不同倾角的平面所截。$\beta>\alpha$ 得椭圆（两个 Dandelin 球各切出一个焦点），$\beta=\alpha$ 得抛物线（一个球切出焦点），$\beta<\alpha$ 得双曲线（两个球分置两侧锥面）；右上方小图标为截线实形}
  \label{fig:dandelin}
\end{figure}

\begin{example}
圆锥的半顶角 $\alpha=30^\circ$，分别用与轴线夹角 $\beta=60^\circ,30^\circ,15^\circ$ 的平面去截，判断截线类型并求离心率。
\end{example}
\begin{solution}
由 $e=\dfrac{\cos\beta}{\cos\alpha}$，其中 $\cos30^\circ=\frac{\sqrt3}{2}$：
\[
  \beta=60^\circ:\ e=\frac{\cos60^\circ}{\cos30^\circ}=\frac{1/2}{\sqrt3/2}=\frac{\sqrt3}{3}<1\Rightarrow\text{椭圆};
\]
\[
  \beta=30^\circ:\ e=\frac{\cos30^\circ}{\cos30^\circ}=1\Rightarrow\text{抛物线};
\]
\[
  \beta=15^\circ:\ e=\frac{\cos15^\circ}{\cos30^\circ}\approx\frac{0.966}{0.866}\approx1.12>1\Rightarrow\text{双曲线}.
\]
只需要比较两个角的大小，就能立刻判定类型——这是「维度打通」带来的预判力。
\end{solution}

\begin{conccard}[二级]{圆锥曲线的统一：截面与焦点—准线}{判断截线类型、求离心率、理解三类曲线定义的统一性}{把 $\beta$ 与 $\alpha$ 的大小关系记反；忘记 $\beta=90^\circ$ 是圆；以为双曲线需要「两个分开的圆锥」}
圆锥面 $x^2+y^2=z^2\tan^2\alpha$ 被与轴线夹角为 $\beta$ 的平面所截：
$\beta>\alpha$ 椭圆（$\beta=90^\circ$ 为圆），$\beta=\alpha$ 抛物线，$\beta<\alpha$ 双曲线；
统一离心率 $e=\dfrac{\cos\beta}{\cos\alpha}$，统一定义为 $\dfrac{|PF|}{d(P,l)}=e$。
\end{conccard}

\subsection{考纲安全落地：高观点预判与通法满分的双层法}

\key{探究三：上层用高观点「看透」，下层用考纲通法「写全」}

高考解析几何解答题的特点是：答案往往可以被高观点瞬间看穿，但阅卷只看考纲允许的推理链。
因此最稳的策略是\key{双层法}：先用高观点（极点极线、仿射变换）确定目标与答案，再用考纲通法（设线、联立、韦达、判别式）把过程写得滴水不漏。

\key{（1）高观点的两个预判工具。}
\begin{itemize}
  \item \key{极点极线预设定点定值}：椭圆 $\frac{x^2}{a^2}+\frac{y^2}{b^2}=1$ 外一点 $P(x_0,y_0)$ 的两条切线的切点连线（切点弦）方程为
  \[
    \frac{x_0x}{a^2}+\frac{y_0y}{b^2}=1,
  \]
  它被称为 $P$ 的\key{极线}。若 $P$ 在定直线 $x=m$ 上运动，则极线 $\frac{mx}{a^2}+\frac{y_0y}{b^2}=1$ 恒过定点 $\left(\frac{a^2}{m},0\right)$——不必动笔算，定点已出。
  \item \key{仿射变换秒算面积最值}：涉及椭圆内接图形面积的最值，先仿射回圆，用圆的结果乘 $\frac{b}{a}$ 即可。
\end{itemize}

\key{（2）考纲通法的标准四步。}
设线 $\to$ 联立 $\to$ 韦达 $\to$ 回代（并用判别式检验）。每一步都要写在卷面上，因为分步给分。

\begin{example}[两层对照示范]
椭圆 $E:\dfrac{x^2}{4}+\dfrac{y^2}{2}=1$，动点 $P$ 在定直线 $x=2$ 上，过 $P$ 作 $E$ 的两条切线，切点分别为 $A,B$。求证：直线 $AB$ 恒过定点，并求出该定点。
\end{example}
\begin{solution}
\key{第一步（高观点预判）。}设 $P(2,t)$。由极线公式，$AB$ 即 $P$ 的极线：
\[
  \frac{2x}{4}+\frac{ty}{2}=1\iff x+ty=2,
\]
它显然恒过 $(2,0)$。十秒钟，答案已经到手。

\key{第二步（考纲通法落地）。}设 $A(x_1,y_1),B(x_2,y_2)$。
由隐函数求导（第 6 章）：椭圆两边对 $x$ 求导得 $\frac{2x}{4}+\frac{2yy'}{2}=0$，即 $y'=-\frac{b^2x}{a^2y}=-\frac{x}{2y}$。
所以椭圆在 $A$ 处的切线方程为
\[
  y-y_1=-\frac{x_1}{2y_1}(x-x_1).
\]
利用 $\frac{x_1^2}{4}+\frac{y_1^2}{2}=1$ 化简，得切线的简洁形式
\[
  \frac{x_1x}{4}+\frac{y_1y}{2}=1.
\]
（若尚未学求导，也可设切线斜率 $k$，联立椭圆后令 $\Delta=0$，推出同一结果。）
因为 $P(2,t)$ 同时在两条切线上，代入得
\[
  \frac{2x_1}{4}+\frac{ty_1}{2}=1,\qquad \frac{2x_2}{4}+\frac{ty_2}{2}=1,
  \iff x_1+ty_1=2,\quad x_2+ty_2=2.
\]
这说明 $A,B$ 都在直线 $x+ty=2$ 上；两点确定一条直线，故
\[
  AB:\ x+ty=2.
\]
令 $y=0$ 得 $x=2$，与 $t$ 无关，所以 $AB$ \key{恒过定点 $(2,0)$}。

\key{对照结论}：高观点负责「知道往哪走」（预判定点 $(2,0)$），通法负责「写得干净」（求导得切线、两点定直线）。二者分工明确，考场最稳。
\end{solution}

\begin{figure}[htbp]
  \centering
  \begin{tikzpicture}[>=Stealth, semithick, font=\footnotesize]
    \draw[->, gray] (-2.9,0)--(4.4,0) node[below right] {$x$};
    \draw[->, gray] (0,-2.2)--(0,2.9) node[above left] {$y$};
    % 椭圆
    \fill[second!8] (0,0) ellipse (2 and 1.4142);
    \draw[main!70!black, very thick] (0,0) ellipse (2 and 1.4142);
    % 极线（切点弦）
    \draw[third!70!black, thick, dashed] (-0.5,2.75)--(2.0,-1.0);
    % 两条切线
    \draw[main!80!black, thick] (3,1)--($(1.771,-0.657)!-0.4!(3,1)$);
    \draw[main!80!black, thick] (3,1)--($(0.411,1.384)!-0.4!(3,1)$);
    % 过 P 的动弦
    \draw[second, thick] (-1.809,-0.603)--(3.75,1.25);
    % 点与标注
    \fill (0,0) circle (1pt);
    \fill[main] (3,1) circle (1.6pt) node[right] {$P(x_0,y_0)$};
    \fill[third] (1.771,-0.657) circle (1.6pt) node[below right] {$A$};
    \fill[third] (0.411,1.384) circle (1.6pt) node[above left] {$B$};
    \fill[second] (1.809,0.603) circle (1.6pt) node[above right] {$M$};
    \fill[second] (-1.809,-0.603) circle (1.6pt) node[below left] {$N$};
    \node[third!70!black, anchor=west] at (0.15,2.05) {$AB$：$P$ 的极线（切点弦）};
    \node[second!70!black, anchor=west] at (0.15,-1.7) {过 $P$ 的动弦 $MN$};
  \end{tikzpicture}
  \caption{极点极线与切点弦：椭圆 $\frac{x^2}{4}+\frac{y^2}{2}=1$ 外一点 $P(3,1)$，$PA,PB$ 为切线、$A,B$ 为切点，则 $AB$ 即 $P$ 的极线；过 $P$ 的任意弦 $MN$ 与极线的关系给出定点定值问题的预判。当 $P$ 在定直线 $x=m$ 上运动时，极线 $\frac{mx}{a^2}+\frac{y_0y}{b^2}=1$ 恒过定点 $\left(\frac{a^2}{m},0\right)$}
  \label{fig:pole-polar}
\end{figure}

注记（高观点对偶）：把上面的推理倒过来即得逆命题——过定点 $(2,0)$ 的直线与椭圆交于 $A,B$，则 $A,B$ 处切线的交点必在直线 $x=2$ 上。
这在极点极线理论中称为 \key{La Hire 对偶}：点 $P$ 在点 $Q$ 的极线上 $\iff$ 点 $Q$ 在点 $P$ 的极线上。

\begin{method}[双层法：高观点预判 $+$ 通法满分]
\begin{enumerate}
  \item \key{预判}：先问「答案大概长什么样」。定点定值类问题用极线公式扫一眼；面积最值类问题先仿射回圆算一遍；
  \item \key{设线}：设动直线（含参数）或设动点坐标，注意斜率不存在的情形要单独讨论；
  \item \key{联立}：与圆锥曲线方程联立，消元得到关于 $x$（或 $y$）的一元二次方程；
  \item \key{韦达}：先写 $\Delta>0$ 保证两交点，再写 $x_1+x_2,\ x_1x_2$；
  \item \key{回代}：把目标量（数量积、斜率、面积、中点）用韦达结果表示，整理成关于参数的式子；
  \item \key{定值}：令参数系数为零（或「对任意参数成立」列方程组）求出定点／定值，最后回代检验。
\end{enumerate}
\end{method}

\subsection{解析几何主干：定义、方程与直线的位置关系}

\key{探究四：把「二次曲线」拆成三件通用工具}

\key{（1）三大定义与标准方程。}
\begin{center}
\footnotesize
\setlength{\tabcolsep}{4pt}
\renewcommand{\arraystretch}{1.35}
\begin{tabular}{p{1.5cm}p{4.6cm}p{3.5cm}p{3.0cm}p{1.9cm}}
\toprule
曲线 & 定义 & 标准方程 & 焦点 & 离心率 \\
\midrule
椭圆 & $|PF_1|+|PF_2|=2a\ (2a>2c)$ & $\dfrac{x^2}{a^2}+\dfrac{y^2}{b^2}=1\ (a>b>0)$ & $(\pm c,0),\ c^2=a^2-b^2$ & $e=\dfrac ca\in(0,1)$ \\
双曲线 & $\bigl||PF_1|-|PF_2|\bigr|=2a\ (2a<2c)$ & $\dfrac{x^2}{a^2}-\dfrac{y^2}{b^2}=1$ & $(\pm c,0),\ c^2=a^2+b^2$ & $e=\dfrac ca>1$ \\
抛物线 & $|PF|=d(P,l)$ & $y^2=2px\ (p>0)$ & $\left(\dfrac p2,0\right)$，准线 $x=-\dfrac p2$ & $e=1$ \\
\bottomrule
\end{tabular}
\end{center}
三条定义其实是同一件事（到焦点与到准线的距离比为 $e$）的特例；$a,b,c$ 的关系只需记「椭圆 $c^2=a^2-b^2$、双曲线 $c^2=a^2+b^2$、抛物线只有一个 $p$」。

\key{（2）直线与圆锥曲线的位置关系：联立 $+$ 判别式。}
设直线 $l:y=kx+m$，与圆锥曲线方程联立、消去 $y$，得到关于 $x$ 的一元二次方程 $Ax^2+Bx+C=0$（$A$ 可能为 $0$），则
\[
  \Delta>0\iff \text{两个不同公共点};\quad
  \Delta=0\iff \text{相切（一个公共点）};\quad
  \Delta<0\iff \text{无公共点}.
\]
对抛物线等特殊情形，还要检查二次项系数 $A=0$（直线平行于对称轴）时的退化情况——此时方程退化为一次方程，最多一个交点。

\key{（3）弦长公式。}
设两交点为 $A(x_1,y_1),B(x_2,y_2)$，则
\[
  |AB|=\sqrt{1+k^2}\,|x_1-x_2|
  =\sqrt{1+k^2}\cdot\frac{\sqrt{\Delta}}{|A|}.
\]
若直线写成 $x=my+n$，则改用 $|AB|=\sqrt{1+m^2}\,|y_1-y_2|$。记忆钩子：\key{弦长 $=$ 水平（竖直）投影 $\times$ 放大因子}。

\key{（4）中点弦：点差法。}
设 $A(x_1,y_1),B(x_2,y_2)$ 在椭圆 $\frac{x^2}{a^2}+\frac{y^2}{b^2}=1$ 上，$M(x_0,y_0)$ 为 $AB$ 的中点。两式相减：
\[
  \frac{(x_1+x_2)(x_1-x_2)}{a^2}+\frac{(y_1+y_2)(y_1-y_2)}{b^2}=0.
\]
代入 $x_1+x_2=2x_0,\ y_1+y_2=2y_0$，得
\[
  k_{AB}=\frac{y_1-y_2}{x_1-x_2}=-\frac{b^2x_0}{a^2y_0},
  \qquad\text{等价地}\quad k_{AB}\cdot k_{OM}=-\frac{b^2}{a^2}.
\]
\key{点差法的实质是「中点 $\Rightarrow$ 斜率」}，但必须回代检验直线与椭圆确有两交点。

\begin{example}
椭圆 $E:\dfrac{x^2}{4}+\dfrac{y^2}{2}=1$ 的一条弦 $AB$ 以 $M(1,1)$ 为中点，求直线 $AB$ 的方程与弦长。
\end{example}
\begin{solution}
由点差法（$a^2=4,b^2=2,x_0=y_0=1$）：
\[
  k_{AB}=-\frac{b^2x_0}{a^2y_0}=-\frac{2}{4}=-\frac12.
\]
故直线 $AB:y-1=-\frac12(x-1)$，即
\[
  x+2y-3=0.
\]
检验（韦达）：代入椭圆得
\[
  \frac{x^2}{4}+\frac{\left(-\frac x2+\frac32\right)^2}{2}=1
  \iff 3x^2-6x+1=0,
\]
$\Delta=36-12=24>0$，确有两个交点。$x_1+x_2=2$（中点横坐标为 $1$，一致），
\[
  |x_1-x_2|=\frac{\sqrt{24}}{3}=\frac{2\sqrt6}{3},
  \qquad
  |AB|=\sqrt{1+\left(-\frac12\right)^2}\cdot\frac{2\sqrt6}{3}
  =\frac{\sqrt5}{2}\cdot\frac{2\sqrt6}{3}=\frac{\sqrt{30}}{3}.
\]
\end{solution}

\begin{conccard}[二级]{直线与圆锥曲线：判别式、弦长、中点弦}{位置关系判定、弦长计算、中点弦方程、弦中点轨迹}{忽略二次项系数为 $0$ 的退化；点差法不检验 $\Delta$；弦长公式漏乘 $\sqrt{1+k^2}$ 或把分母写成 $\sqrt\Delta$}
联立得 $Ax^2+Bx+C=0$：$\Delta>0$ 相交、$\Delta=0$ 相切、$\Delta<0$ 相离；
$|AB|=\sqrt{1+k^2}\cdot\frac{\sqrt\Delta}{|A|}$；
中点弦斜率 $k=-\dfrac{b^2x_0}{a^2y_0}$，且 $k\cdot k_{OM}=-\dfrac{b^2}{a^2}$。
\end{conccard}

\begin{commonerror}{直线与圆锥曲线联立的三个高频失误}
\paragraph{失误一：忘记讨论二次项系数为 $0$。}
消元后若 $x^2$ 项系数含参数（如直线与抛物线联立），它可能为 $0$，此时方程退化为一元一次方程，只有一个交点。
必须先声明「当 $A=0$ 时……；当 $A\ne0$ 时……」。
\paragraph{失误二：点差法求出斜率就直接写答案。}
点差法只保证「若 $M$ 是中点，则斜率必为 $k$」，并不保证这样的弦真的存在。
必须回代验证 $\Delta>0$（上面例题中 $\Delta=24>0$ 才可写）。
\paragraph{失误三：弦长公式张冠李戴。}
$|x_1-x_2|=\frac{\sqrt\Delta}{|A|}$ 只是水平投影；弦长还要乘 $\sqrt{1+k^2}$。
若直线设为 $x=my+n$，则应乘 $\sqrt{1+m^2}$ 并对 $y$ 用韦达，切勿混用。
\end{commonerror}

\begin{conccard}[三级]{定点定值问题的通法骨架}{含参数的动直线与圆锥曲线相交，证明某点／某量为定值}{只算特殊位置就下结论；忽略判别式；参数没有全部消去就说「定值」}
设线（含参数）$\to$ 联立 $\to$ 韦达（先 $\Delta>0$）$\to$ 把目标量表示成参数的式子 $\to$ 令各参数系数为零（或对任意参数成立）$\to$ 解出定点／定值并回代检验。
高观点（极线、仿射）用于提前预判答案，通法用于写满步骤。
\end{conccard}

\subsection{立体几何主干：位置关系、空间角与向量法}

\key{探究五：判定靠逻辑，计算靠向量}

\key{（1）线面平行与垂直：判定和性质是一体两面。}
\begin{itemize}
  \item \key{线面平行}：平面外一条直线与此平面内一条直线平行 $\implies$ 线面平行（判定）；
        线面平行 $\implies$ 过该直线的任一平面与原平面的交线与该直线平行（性质）。
  \item \key{线面垂直}：一条直线与平面内\key{两条相交}直线都垂直 $\implies$ 线面垂直（判定）；
        线面垂直 $\implies$ 垂直于平面内任意直线（性质）；垂直于同一平面的两直线互相平行。
  \item \key{面面关系}：线面平行 $+$ 线线平行可推面面平行；线面垂直可推面面垂直。它们是上述两个判定的「升级版」。
\end{itemize}
解题时的固定动作：\key{要证平行，先找「面内一条线与它平行」；要证垂直，先找「面内两条相交线与它垂直」}。

\key{（2）空间角与距离的定义。}
\begin{itemize}
  \item \key{线面角}：斜线与它在平面内的射影所成的角，范围 $\left[0,\frac\pi2\right]$；线面垂直时为 $\frac\pi2$，线面平行或在面内时为 $0$；
  \item \key{二面角}：在棱上取一点，在两个半平面内分别作垂直于棱的射线，两射线的夹角即二面角的平面角；
  \item \key{点到平面的距离}：点到平面的垂线段长。几何法常用\key{等体积法}（换底面，体积不变）；向量法用公式 $d=\frac{|\overrightarrow{AP}\cdot\vec n|}{|\vec n|}$。
\end{itemize}

\key{（3）向量法：把第 3 章的工具全部调用过来。}
以方向向量 $\vec s$ 与法向量 $\vec n$ 为「两个抓手」，空间位置关系与度量可全部坐标化。
线面平行／垂直、线面角、二面角四条判据与三个高频失误已在第 3 章
完整给出（点面距公式也出自同一套向量法），此处不重复；
本章只补一条应用——点到平面的距离
\[
  d=\frac{|\overrightarrow{AP}\cdot\vec n|}{|\vec n|}\quad(A\in\alpha),
\]
即 $\overrightarrow{AP}$ 在法向量方向上投影的绝对值；几何法常用等体积法对照。

\begin{example}
四棱锥 $P-ABCD$ 中，底面 $ABCD$ 是边长为 $2$ 的正方形，$PA\perp$ 底面 $ABCD$，$PA=2$。
（1）求证：$BD\perp$ 平面 $PAC$；（2）求点 $A$ 到平面 $PCD$ 的距离；（3）求二面角 $A-PC-D$ 的余弦值。
\end{example}
\begin{solution}
以 $A$ 为原点，$\overrightarrow{AB},\overrightarrow{AD},\overrightarrow{AP}$ 分别为 $x,y,z$ 轴建立空间直角坐标系（三条轴两两垂直），则
\[
  A(0,0,0),\ B(2,0,0),\ C(2,2,0),\ D(0,2,0),\ P(0,0,2).
\]
（1）$\overrightarrow{BD}=(-2,2,0)$，$\overrightarrow{AC}=(2,2,0)$，$\overrightarrow{AP}=(0,0,2)$。
\[
  \overrightarrow{BD}\cdot\overrightarrow{AC}=-4+4=0,\qquad
  \overrightarrow{BD}\cdot\overrightarrow{AP}=0,
\]
故 $BD\perp AC$ 且 $BD\perp AP$。又 $AC\cap AP=A$，所以 $BD\perp$ 平面 $PAC$。

（2）$\overrightarrow{PC}=(2,2,-2)$，$\overrightarrow{PD}=(0,2,-2)$，平面 $PCD$ 的法向量
\[
  \vec n=\overrightarrow{PC}\times\overrightarrow{PD}=(0,4,4)\parallel(0,1,1).
\]
点 $A$ 到平面 $PCD$ 的距离
\[
  d=\frac{|\overrightarrow{AD}\cdot\vec n|}{|\vec n|}
  =\frac{|(0,2,0)\cdot(0,1,1)|}{\sqrt2}=\frac{2}{\sqrt2}=\sqrt2.
\]
几何法对照（等体积法）：
$V_{P-ACD}=\frac13 S_{\triangle ACD}\cdot PA=\frac13\cdot2\cdot2=\frac43$；
又 $PD^2+CD^2=8+4=12=PC^2$，故 $\triangle PCD$ 在 $D$ 处为直角，$S_{\triangle PCD}=\frac12\cdot2\sqrt2\cdot2=2\sqrt2$。
由 $\frac13 S_{\triangle PCD}\cdot h=\frac43$ 得 $h=\sqrt2$，与向量法一致。

（3）平面 $APC$ 的法向量 $\vec n_1=\overrightarrow{AP}\times\overrightarrow{AC}=(-4,4,0)\parallel(-1,1,0)$；平面 $PCD$ 的法向量 $\vec n_2=(0,1,1)$。
\[
  |\cos\theta|=\frac{|\vec n_1\cdot\vec n_2|}{|\vec n_1|\,|\vec n_2|}
  =\frac{|0+1+0|}{\sqrt2\cdot\sqrt2}=\frac12.
\]
取 $PC$ 的中点 $M(1,1,1)$，把 $\overrightarrow{MA}=(-1,-1,-1)$、$\overrightarrow{MD}=(-1,1,-1)$ 去掉沿 $\overrightarrow{PC}$ 方向的分量，得到二面角平面角的两条边，计算得其夹角余弦为 $\frac12>0$，故二面角 $A-PC-D$ 为锐角，其余弦值为
\[
  \cos\theta=\frac12.
\]
传统几何法需要作平面角、证垂直、再解三角形；向量法只需建系、算两个法向量、代公式。
\end{solution}

\begin{conccard}[一级]{立体几何的「判定—性质」闭环}{证平行／垂直时的辅助线设计}{只会背定理不会用；找错「面内两条相交直线」；用缺少相交条件的向量}
证线面平行：面外线 $\parallel$ 面内线；证线面垂直：直线 $\perp$ 面内两条相交直线；
线面垂直 $\Rightarrow$ 线线垂直；线面平行 $+$ 相交平面 $\Rightarrow$ 线线平行。
几何法与向量法各有胜场：垂直关系明显时优先几何法（等体积、三垂线），坐标好写时优先向量法。
\end{conccard}

\begin{summary}
本章用「仿射变换 $+$ 三维截面 $+$ 向量法」三件武器打通了立体几何与解析几何的维度壁垒。
椭圆是圆沿 $y$ 轴按比例 $\frac{b}{a}$ 的均匀压缩像，仿射不变量（共线、平行、中点、相切、面积比）全部保持，故椭圆内接三角形最大面积为 $\frac{3\sqrt3}{4}ab$；
Dandelin 球说明椭圆／双曲线／抛物线是同一个圆锥面被不同倾角平面所截的结果，统一为焦点—准线定义，离心率 $e=\frac{\cos\beta}{\cos\alpha}$；
考场上层的「高观点预判」（极线公式预设定点、仿射变换秒算面积）与下层的「考纲通法」（设线—联立—韦达—判别式）合成双层法，既看得透又写得全；
解析几何主干落实为定义与标准方程表、判别式判位置、弦长公式与中点弦点差法，以及定点定值问题的六步通法；
立体几何主干落实为线面平行／垂直的判定与性质、空间角与距离，并由第 3 章的方向向量与法向量统一处理。
\key{维度不是障碍，变换才是本质；答案是预判出来的，步骤是通法写出来的。}
\end{summary}

\end{explore}

\begin{inclasstest}
\begin{enumerate}
  \item 把圆 $x^2+y^2=9$ 沿 $y$ 轴按比例 $\frac23$ 均匀压缩，写出所得椭圆方程，并求其内接三角形面积的最大值；\level{易}
  \item 求抛物线 $y^2=8x$ 的焦点坐标与准线方程；\level{易}
  \item 若直线 $y=x+m$ 与椭圆 $\dfrac{x^2}{4}+\dfrac{y^2}{2}=1$ 有两个不同的交点，求实数 $m$ 的取值范围；当两者相切时 $m$ 取何值？\level{中}
  \item 三棱锥 $O-ABC$ 中，$OA,OB,OC$ 两两垂直，$OA=1,OB=2,OC=3$，求直线 $OC$ 与平面 $ABC$ 所成角的正弦值。\level{中}
\end{enumerate}
\end{inclasstest}

\begin{afterclass}
\subsection*{基础巩固}
\begin{enumerate}
  \item 求椭圆 $\dfrac{x^2}{16}+\dfrac{y^2}{12}=1$ 的半长轴、半短轴、半焦距、焦点坐标与离心率；椭圆上任意一点到两焦点距离之和是多少？
  \item 求双曲线 $\dfrac{x^2}{9}-\dfrac{y^2}{16}=1$ 的实轴长、虚轴长、离心率与渐近线方程。
\end{enumerate}

\subsection*{进阶提升}
\begin{enumerate}
  \item 椭圆 $E:\dfrac{x^2}{4}+y^2=1$，动点 $P$ 在定直线 $x=4$ 上，过 $P$ 作 $E$ 的两条切线，切点分别为 $A,B$，求证：直线 $AB$ 恒过定点，并求出该定点。
  \item 椭圆 $E:\dfrac{x^2}{9}+\dfrac{y^2}{4}=1$ 中，以 $M\left(\dfrac32,1\right)$ 为中点的弦所在直线方程为何？该弦的长度是多少？
\end{enumerate}

\subsection*{拓展探究}
\begin{enumerate}
  \item （极点极线的一般化）已知椭圆 $E:\dfrac{x^2}{a^2}+\dfrac{y^2}{b^2}=1$，$P(x_0,y_0)$ 为 $E$ 外一点，过 $P$ 作两条切线，切点分别为 $A,B$。
  \begin{enumerate}
    \item 求证：切点弦（极线）$AB$ 的方程为 $\dfrac{x_0x}{a^2}+\dfrac{y_0y}{b^2}=1$；
    \item 若 $P$ 在定直线 $x=m\ (m\neq0)$ 上运动，求 $AB$ 恒过的定点，并用第 6 章的隐函数求导检验 $a=2,b=\sqrt2,x_0=2$ 的特例。
  \end{enumerate}
  （提示：设切点 $(x_1,y_1)$，用隐函数求导求切线斜率，化简时利用 $\frac{x_1^2}{a^2}+\frac{y_1^2}{b^2}=1$。）
\end{enumerate}

\subsection*{教师版答案}
\begin{solution}
\textbf{基础 1：}$a^2=16,b^2=12\Rightarrow a=4,b=2\sqrt3$；$c^2=a^2-b^2=4\Rightarrow c=2$。
焦点 $(\pm2,0)$，离心率 $e=\dfrac ca=\dfrac12$；椭圆上任意一点到两焦点距离之和为 $2a=8$。
\end{solution}
\begin{solution}
\textbf{基础 2：}$a^2=9,b^2=16\Rightarrow a=3,b=4$；$c^2=a^2+b^2=25\Rightarrow c=5$。
实轴长 $2a=6$，虚轴长 $2b=8$，离心率 $e=\dfrac53$，渐近线 $y=\pm\dfrac{b}{a}x=\pm\dfrac43x$。
\end{solution}
\begin{solution}
\textbf{进阶 1：}设 $P(4,t)$。由极线公式（或设切点用求导推导），切点弦 $AB$ 的方程为
\[
  \frac{4x}{4}+\frac{ty}{1}=1\iff x+ty=1.
\]
令 $y=0$ 得 $x=1$，与 $t$ 无关，故 $AB$ 恒过定点 $(1,0)$。
\end{solution}
\begin{solution}
\textbf{进阶 2：}由点差法，$k_{AB}=-\dfrac{b^2x_0}{a^2y_0}=-\dfrac{4\cdot\frac32}{9\cdot1}=-\dfrac23$。
直线 $AB:y-1=-\dfrac23\left(x-\dfrac32\right)$，即
\[
  2x+3y-6=0.
\]
代入椭圆：$\dfrac{x^2}{9}+\dfrac{\left(2-\frac23x\right)^2}{4}=1
\iff \dfrac29x^2-\dfrac23x=0$，得 $x=0$ 或 $x=3$。
对应交点为 $(0,2)$ 与 $(3,0)$（恰为椭圆的两个顶点），中点确为 $\left(\dfrac32,1\right)$，弦长
\[
  |AB|=\sqrt{(3-0)^2+(0-2)^2}=\sqrt{13}.
\]
\end{solution}
\begin{solution}
\textbf{拓展：（1）}设 $A(x_1,y_1)$。椭圆两边对 $x$ 求导：$\dfrac{2x}{a^2}+\dfrac{2yy'}{b^2}=0\Rightarrow y'=-\dfrac{b^2x}{a^2y}$。
故 $A$ 处切线为 $y-y_1=-\dfrac{b^2x_1}{a^2y_1}(x-x_1)$，整理：
\[
  a^2y_1y-a^2y_1^2=-b^2x_1x+b^2x_1^2
  \iff b^2x_1x+a^2y_1y=b^2x_1^2+a^2y_1^2
  =a^2b^2\left(\frac{x_1^2}{a^2}+\frac{y_1^2}{b^2}\right)=a^2b^2,
\]
两边除以 $a^2b^2$ 得切线方程 $\dfrac{x_1x}{a^2}+\dfrac{y_1y}{b^2}=1$。
因为 $P(x_0,y_0)$ 在切线上，所以 $\dfrac{x_1x_0}{a^2}+\dfrac{y_1y_0}{b^2}=1$；同理 $\dfrac{x_2x_0}{a^2}+\dfrac{y_2y_0}{b^2}=1$。
这表明 $A,B$ 都在直线 $\dfrac{x_0x}{a^2}+\dfrac{y_0y}{b^2}=1$ 上，两点定一直线，即为 $AB$ 的方程。

\textbf{（2）}令 $x_0=m$，极线为 $\dfrac{mx}{a^2}+\dfrac{y_0y}{b^2}=1$。令 $y=0$ 得 $x=\dfrac{a^2}{m}$，与 $y_0$ 无关，故 $AB$ 恒过定点 $\left(\dfrac{a^2}{m},0\right)$。
特例检验：$a=2,b=\sqrt2,m=2$，定点为 $\left(\dfrac{4}{2},0\right)=(2,0)$；极线方程 $\dfrac{2x}{4}+\dfrac{y_0y}{2}=1$ 即 $x+y_0y=2$，当 $y_0=1$ 时 $x+y=2$ 过 $(2,0)$，与课堂例题一致。
\end{solution}
\reading{拓展阅读 10}{圆锥曲线：同一个三维对象的三种截面}
\end{afterclass}

% 高观点导学案：数列与离散动态系统（第 5 章·高观点版）
% 主线：差分与求和互逆（离散微积分）→ 递推即离散动力系统（不动点与蛛网图）→ 线性递推与特征值（矩阵同构）
%       → Abel 变换（错位相减的机械化）→ 共轭变换（不动点为何生效）→ 生成函数（强基拓展，超出考纲）
% 核心思想：数列不是公式的仓库，而是离散动力系统在时间轴上的轨迹；解题的本质是换坐标系，
%           把陌生递推化为等比／等差／代数方程。

\section{数列与离散动态系统（高观点版）}

\begin{methodbox}{多框架对读}
\small
\begin{center}
\begin{tabular}{c|p{2.9cm}|p{5.0cm}|p{3.1cm}}
\toprule
视角 & 本章看点 & 别的视角怎么说 & 哪侧最省力 \\
\midrule
代数 & 通项公式、求和与差分 & 符号对象的恒等变形 & 硬算、裂项、Abel \\
几何 & 蛛网图与特征螺线 & 轨迹在坐标平面上的形状 & 看收敛、猜极限 \\
集合 & $\mathbb{N}\to\mathbb{R}$ 的离散函数 & 轨道是值的集合 & 审定义域与边界 \\
映射 & $x_{n+1}=F(x_n)$、共轭变换 & 换坐标系＝换表示 & 非线性 $\to$ 等比／等差 \\
\bottomrule
\end{tabular}
\end{center}
卡的使用：递推题先问「这是哪个映射的自迭代」，再决定搬去几何还是代数（总图见第 0 章图~\ref{fig:four-view-map}）。
\end{methodbox}

第 3 章把向量确立为「几何空间代数化」的通用语言，第 4 章用仿射变换打通了圆锥曲线的维度壁垒；
本章回到最朴素的对象——数列，把它从「套公式」升级为「离散动力系统 $+$ 降维变换」：
数列是定义域为 $\mathbb{N}^*$ 的离散函数，递推是离散时间的状态演化，求和与差分是离散的微积分，
而不动点、特征值、共轭、Abel、生成函数则是把复杂递推「降维」成等比／等差／代数方程的变换。
六个模块层层展开：
\key{差分与求和互逆}（离散版微积分基本定理，由 $S_n$ 反求通项、差分降次）$\to$
\key{递推即离散动力系统}（不动点、蛛网图、$|f'(x^*)|$ 稳定性判据）$\to$
\key{线性递推与特征值}（转移矩阵 $\begin{pmatrix}0&1\\q&p\end{pmatrix}$、特征方向即不动点方向）$\to$
\key{Abel 变换}（错位相减的机械化，与分部积分互译）$\to$
\key{共轭变换}（不动点为何生效：换坐标系、莫比乌斯分式）$\to$
\key{生成函数}（强基拓展，\key{超出高考考纲}）。
核心思想：\key{数列不是公式的仓库，而是离散动力系统在时间轴上的轨迹；解题的本质是换坐标系，把陌生递推化为等比／等差。}

\begin{learnobjective}
\begin{itemize}
  \item 理解前向差分 $\Delta a_n=a_{n+1}-a_n$ 与求和 $\sum$ 互为逆运算，掌握离散版微积分基本定理 $\sum_{k=1}^n\Delta a_k=a_{n+1}-a_1$，会由和式 $S_n$ 反求通项；
  \item 理解「差分降次」与导数降次的类比，会看差分表判断通项的次数；
  \item 会把一阶递推 $a_{n+1}=f(a_n)$ 看作离散动力系统，会解不动点 $x^*=f(x^*)$、画蛛网图，并用 $|f'(x^*)|$ 判稳定性（$<1$ 吸引、$>1$ 排斥）；
  \item 掌握二阶线性递推的特征方程 $\lambda^2-p\lambda-q=0$，理解它与转移矩阵 $\begin{pmatrix}0&1\\q&p\end{pmatrix}$ 的特征值同构，特征方向即分式递推的不动点方向；
  \item 会用 Abel 变换 $\sum_{k=1}^na_kb_k=a_nB_n-\sum_{k=1}^{n-1}(a_{k+1}-a_k)B_k$ 机械化求「等差 $\times$ 等比」型求和；
  \item 理解共轭变换（换坐标系）是不动点法生效的底层原因，会作 $b_n=a_n-x^*$ 的等比化与 $\dfrac{a_n-p}{a_n-q}$ 的分式化；
  \item 了解生成函数 $G(x)=\sum a_nx^n$（\key{超出高考考纲，仅竞赛／强基}）。
\end{itemize}
\end{learnobjective}

\begin{preview}
\subsection*{知识回顾}
\begin{itemize}
  \item 等差数列的通项公式是 \blank，前 $n$ 项和公式是 \blank；
  \item 等比数列的通项公式是 \blank，前 $n$ 项和公式是 \blank；
  \item 设 $S_n$ 是数列 $\{a_n\}$ 的前 $n$ 项和，则 $n\ge2$ 时 $a_n=$ \blank，且必须单独验证 \blank；
  \item 函数 $y=f(x)$ 的图像与直线 $y=x$ 的交点横坐标满足方程 \blank。
\end{itemize}

\subsection*{教材速读}
阅读「数列」一章教材，完成下列问题：
\begin{enumerate}
  \item 由 $S_n$ 求 $a_n$ 的「分解式」$a_n=\begin{cases}S_1,&n=1\\ S_n-S_{n-1},&n\ge2\end{cases}$ 中，$n=1$ 时 \blank（填「需要」或「不需要」）单独检验；
  \item 一阶线性递推 $a_{n+1}=pa_n+q\ (p\ne1)$ 的不动点 $x^*$ 满足 $x^*=$ \blank；
  \item 斐波那契数列 $F_1=F_2=1,\ F_{n+2}=F_{n+1}+F_n$ 的特征方程是 \blank；
  \item 「错位相减法」主要适用于 \blank 数列与 \blank 数列相乘型求和。
\end{enumerate}
\end{preview}

\begin{explore}

\subsection{差分与求和互逆：离散版微积分基本定理}

\key{探究一：$\Delta$ 与 $\sum$ 是同一个东西的两面}

\key{（1）前向差分：数列的「离散导数」。}
定义
\[
  \Delta a_n=a_{n+1}-a_n .
\]
$\Delta$ 把数列 $\{a_n\}$ 映成新数列 $\{\Delta a_n\}$；二阶差分 $\Delta^2a_n=\Delta(\Delta a_n)=a_{n+2}-2a_{n+1}+a_n$，以此类推。
以 $a_n=n^2$ 为例，列表如下：
\[
\begin{array}{c|ccccc}
  n & 1 & 2 & 3 & 4 & 5 \\ \hline
  a_n=n^2 & 1 & 4 & 9 & 16 & 25 \\
  \Delta a_n & 3 & 5 & 7 & 9 & \\
  \Delta^2a_n & 2 & 2 & 2 & &
\end{array}
\]
一阶差分 $\Delta a_n=2n+1$ 是一次等差数列，二阶差分 $\Delta^2a_n=2$ 是常数——每差一次，次数降一次。

\key{（2）差分与求和互逆：离散版微积分基本定理。}
因为求和时中间项全部抵消（telescoping）：
\[
  \sum_{k=1}^{n}\Delta a_k=\sum_{k=1}^{n}(a_{k+1}-a_k)=a_{n+1}-a_1 .
\]
这就是离散版的「牛顿—莱布尼茨公式」：$\sum$ 与 $\Delta$ 互为逆运算。
反过来，若已知前 $n$ 项和 $S_n=\sum_{k=1}^na_k$，则
\[
  a_n=S_n-S_{n-1}\quad(n\ge2),\qquad a_1=S_1 .
\]
所以「由和式反求通项」不是技巧，而是差分的逆运算；注意 $n=1$ 要单独检验。

\key{（3）差分降次：类比导数。}
若 $a_n$ 是 $k$ 次多项式，则 $\Delta a_n$ 是 $k-1$ 次多项式（类比 $(x^k)'=kx^{k-1}$）。由此得到判据：
\key{若 $\Delta^2a_n$ 恒为常数，则 $a_n$ 是二次式；若 $\Delta a_n$ 是等差数列，则 $a_n$ 是二次数列}。
差分表因此可以「降次」，把高次结构一层层剥成常数。

\begin{example}
已知 $S_n=n^2+2n$，求 $\{a_n\}$ 的通项公式。
\end{example}
\begin{solution}
$a_1=S_1=1+2=3$；当 $n\ge2$ 时
\[
  a_n=S_n-S_{n-1}=(n^2+2n)-\bigl((n-1)^2+2(n-1)\bigr)=2n+1 .
\]
$n=1$ 时 $2n+1=3=S_1$，可并入，故
\[
  a_n=2n+1 .
\]
差分表视角：$\Delta a_n=2$ 为常数，故 $a_n$ 是一次等差数列，首项 $3$、公差 $2$。
\end{solution}

\begin{example}
用「差分求和」求 $\sum_{k=1}^nk^2$。
\end{example}
\begin{solution}
利用恒等式
\[
  (k+1)^3-k^3=3k^2+3k+1 ,
\]
左边正是 $k^3$ 的差分结构。对 $k=1,2,\dots,n$ 求和，左边 telescoping：
\[
  \sum_{k=1}^n\bigl((k+1)^3-k^3\bigr)=(n+1)^3-1 ,
\]
右边为
\[
  \sum_{k=1}^n(3k^2+3k+1)=3\sum_{k=1}^nk^2+\frac{3n(n+1)}2+n .
\]
联立解得
\[
  \sum_{k=1}^nk^2=\frac{n(n+1)(2n+1)}6 .
\]
关键：把 $k^2$ 写成「差分的展开」，求和就自动相消。
\end{solution}

\begin{conccard}[一级]{差分与求和互逆：$\Delta\leftrightarrow\sum$}{由 $S_n$ 反求通项、裂项相消、判断通项次数}{忘记 $n=1$ 单独检验；把 $\Delta a_n$ 的符号与 $a_n$ 的符号混同；误以为 $\Delta$ 保持多项式次数}
$\Delta a_n=a_{n+1}-a_n$；$\sum_{k=1}^n\Delta a_k=a_{n+1}-a_1$（离散牛顿—莱布尼茨）；
$a_n=S_n-S_{n-1}\ (n\ge2)$，$a_1=S_1$；
$k$ 次多项式差分一次降一次，$\Delta^2a_n$ 为常数 $\iff a_n$ 为二次式。
\end{conccard}

\subsection{递推即离散动力系统：不动点与蛛网图}

（不动点的统一理论——连续情形、几何变换与射影看法——见第 9 章；
本节只取数列递推这一路应用。）

\key{探究二：把 $a_{n+1}=f(a_n)$ 看成状态演化}

\key{（1）离散动力系统。}
数列 $\{a_n\}$ 是定义域为 $\mathbb{N}^*$ 的函数；一阶递推
\[
  a_{n+1}=f(a_n)
\]
表示：已知当前状态 $a_n$，下一步状态由同一个规则 $f$ 给出。重复作用 $f$ 就得到轨迹 $a_1,a_2,a_3,\dots$。
这正是「离散时间、离散状态」的动力系统。

\key{（2）不动点：系统的平衡态。}
若
\[
  x^*=f(x^*),
\]
则从 $x^*$ 出发永远停在 $x^*$：$f(x^*)=x^*,f(f(x^*))=x^*,\dots$。
几何上，$x^*$ 是曲线 $y=f(x)$ 与直线 $y=x$ 的交点横坐标。
对 $a_{n+1}=pa_n+q\ (p\ne1)$，不动点为 $x^*=\dfrac{q}{1-p}$。

\key{（3）蛛网图：把迭代画成阶梯。}
画法：从 $x$ 轴上的 $a_1$ 出发，竖直画到曲线 $y=f(x)$ 得 $(a_1,a_2)$；水平画到直线 $y=x$ 得 $(a_2,a_2)$；
再竖直画到曲线得 $(a_2,a_3)$；如此反复。折线像「蜘蛛网」一样阶梯式逼近或远离不动点。

\key{（4）稳定性判据：$|f'(x^*)|$ 决定吸引还是排斥。}
在不动点附近写 $a_n=x^*+\varepsilon_n$，则
\[
  \varepsilon_{n+1}=f(x^*+\varepsilon_n)-x^*\approx f'(x^*)\,\varepsilon_n .
\]
于是偏差按因子 $\lambda=f'(x^*)$ 缩放：
\begin{itemize}
  \item $|f'(x^*)|<1$：\key{吸引（稳定）}，偏差逐次缩小，轨迹收敛到 $x^*$；
  \item $|f'(x^*)|>1$：\key{排斥（不稳定）}，偏差逐次放大，轨迹逃离 $x^*$；
  \item $f'(x^*)<0$ 时偏差每步变号，轨迹\key{摆动}收敛；$0<f'(x^*)<1$ 时偏差不变号，轨迹\key{单调}收敛。
\end{itemize}
对线性 $f(x)=px+q$，近似就是精确：$b_{n+1}=pb_n$，即等比数列。

\begin{figure}[htbp]
  \centering
  \begin{tikzpicture}[>=Stealth, semithick, font=\footnotesize]
    % ---------- 左：吸引不动点 ----------
    \begin{scope}[xshift=0cm]
      \draw[->, gray] (-0.2,0)--(5.6,0) node[below right] {$x$};
      \draw[->, gray] (0,-0.2)--(0,5.2) node[above left] {$y$};
      \draw[gray!60, dashed] (-0.2,-0.2)--(5.2,5.2) node[right] {$y=x$};
      \draw[main!70!black, very thick] (-0.1,1.95)--(5.4,4.7) node[right] {$y=\frac12x+2$};
      \fill[second] (4,4) circle (1.8pt) node[above left] {$x^*=4$};
      \draw[third!70!black, thick] (0.5,0)--(0.5,2.25)--(2.25,2.25)--(2.25,3.125)
        --(3.125,3.125)--(3.125,3.5625)--(3.5625,3.5625);
      \fill[third] (0.5,0) circle (1.2pt) node[below] {$a_1$};
      \fill[third] (2.25,2.25) circle (1.2pt) node[below right] {$a_2$};
      \fill[third] (3.125,3.125) circle (1.2pt) node[below right] {$a_3$};
      \node[main!70!black] at (2.6,-1.0) {$|f'(4)|=\frac12<1$：阶梯收缩，收敛};
    \end{scope}
    % ---------- 右：排斥不动点 ----------
    \begin{scope}[xshift=7.4cm]
      \draw[->, gray] (-0.2,0)--(5.6,0) node[below right] {$x$};
      \draw[->, gray] (0,-0.2)--(0,5.2) node[above left] {$y$};
      \draw[gray!60, dashed] (-0.2,-0.2)--(5.2,5.2) node[right] {$y=x$};
      \draw[main!70!black, very thick] (0.75,0)--(3.6,4.56) node[right] {$y=1.6x-1.2$};
      \fill[second] (2,2) circle (1.8pt) node[above left] {$x^*=2$};
      \draw[third!70!black, thick] (2.25,0)--(2.25,2.4)--(2.4,2.4)--(2.4,2.64)
        --(2.64,2.64)--(2.64,3.024)--(3.024,3.024);
      \fill[third] (2.25,0) circle (1.2pt) node[below] {$a_1$};
      \fill[third] (2.64,2.64) circle (1.2pt) node[below right] {$a_3$};
      \fill[third] (3.024,3.024) circle (1.2pt) node[below right] {$a_4$};
      \node[main!70!black] at (2.6,-1.0) {$|g'(2)|=1.6>1$：阶梯放大，发散};
    \end{scope}
  \end{tikzpicture}
  \caption{蛛网图：从 $a_1$ 出发，竖直打到 $y=f(x)$、水平反射到 $y=x$，如此阶梯迭代。左图 $|f'(x^*)|<1$，阶梯收缩到不动点；右图 $|g'(x^*)|>1$，阶梯放大逃离不动点}
  \label{fig:cobweb}
\end{figure}

\begin{example}
设 $a_1=0,\ a_{n+1}=\dfrac12a_n+2$，求通项并判断数列的收敛性。
\end{example}
\begin{solution}
不动点满足 $x^*=\frac12x^*+2$，解得 $x^*=4$。两式相减：
\[
  a_{n+1}-4=\frac12(a_n-4) .
\]
故 $\{a_n-4\}$ 是首项 $a_1-4=-4$、公比 $\frac12$ 的等比数列：
\[
  a_n-4=-4\left(\frac12\right)^{n-1}=-2^{3-n},
  \qquad
  a_n=4-2^{3-n} .
\]
因 $|f'(4)|=\frac12<1$，不动点 $4$ 吸引，$a_n\to4$（单调递增）。
\end{solution}

\begin{conccard}[二级]{不动点与稳定性：$|f'(x^*)|$ 判据}{一阶递推 $a_{n+1}=f(a_n)$ 的收敛性预判、蛛网图定性、线性化降维}{把「令 $a_{n+1}=a_n=x^*$」当成假设相邻项相等；解出 $x^*$ 不检验；对 $f'(x^*)$ 不存在或 $|f'(x^*)|=1$ 的临界情形乱下结论}
不动点 $x^*=f(x^*)$ 是曲线 $y=f(x)$ 与 $y=x$ 的交点；
$b_n=a_n-x^*$ 满足 $b_{n+1}\approx f'(x^*)b_n$；
$|f'(x^*)|<1$ 吸引（稳定），$|f'(x^*)|>1$ 排斥（不稳定）；$f'(x^*)<0$ 摆动收敛。
\end{conccard}

\begin{commonerror}{不动点与稳定性判据的三个高频失误}
\paragraph{失误一：以为「令 $a_{n+1}=a_n=x^*$」是假设相邻两项相等。}
不动点 $x^*$ 是\key{解方程 $x^*=f(x^*)$ 得到的平衡点}，不是把相邻两项强行看成相等。
它的真实作用是「作差消去常数项」（见探究五）。
\paragraph{失误二：不验证 $x^*$ 是否真是不动点。}
解出 $x^*$ 后必须回代 $f(x^*)=x^*$；尤其是分式递推，使分母为零的「解」要舍去。
\paragraph{失误三：乱用 $|f'(x^*)|$ 判据。}
该判据要求 $f$ 在 $x^*$ 附近可导且 $x^*$ 是不动点；对 $f'(x^*)$ 不存在（如 $f(x)=\sqrt x$ 在 $0$ 处）
或 $|f'(x^*)|=1$ 的临界情形，结论要另行分析。
\end{commonerror}

\subsection{线性递推与特征值：状态向量与矩阵同构}

\key{探究三：把递推升维成矩阵乘法}

\key{（1）状态向量：二阶递推 $=$ 二维线性系统。}
二阶线性递推
\[
  a_{n+2}=pa_{n+1}+qa_n
\]
等价于状态向量 $\mathbf v_n=\dbinom{a_n}{a_{n+1}}$ 的演化：
\[
  \binom{a_{n+1}}{a_{n+2}}
  =\begin{pmatrix}0&1\\ q&p\end{pmatrix}\binom{a_n}{a_{n+1}}
  =A\,\mathbf v_n .
\]
于是 $\mathbf v_n=A^{n-1}\mathbf v_1$：\key{递推的每一小步，都是矩阵 $A$ 的一次乘法}。
矩阵 $A$ 的特征方程为
\[
  \det(\lambda I-A)=\begin{vmatrix}\lambda&-1\\-q&\lambda-p\end{vmatrix}
  =\lambda^2-p\lambda-q=0 ,
\]
与递推的特征方程完全一致——这正是「递推特征方程」的矩阵来源。

\key{（2）特征值决定通项结构。}
设特征根为 $\lambda_1,\lambda_2$。
\begin{itemize}
  \item $\lambda_1\ne\lambda_2$：$a_n=c_1\lambda_1^{\,n-1}+c_2\lambda_2^{\,n-1}$；
  \item $\lambda_1=\lambda_2=\lambda$（重根）：$a_n=(c_1+c_2n)\lambda^{\,n-1}$。
\end{itemize}
原理：特征向量方向在 $A$ 作用下只伸缩不变，故沿特征方向的「坐标」按 $\lambda^n$ 增长；
两个特征方向张成状态空间，通项就是它们的线性组合。重根时出现 $n\lambda^{n-1}$ 项（类比 Jordan 块）。

\key{（3）特征方向即不动点方向（射影视角）。}
对分式递推（莫比乌斯变换）
\[
  x_{n+1}=\frac{ax_n+b}{cx_n+d}\quad(ad-bc\ne0),
\]
引入齐次坐标 $x_n=\dfrac{u_n}{v_n}$，则
\[
  \binom{u_{n+1}}{v_{n+1}}=\begin{pmatrix}a&b\\c&d\end{pmatrix}\binom{u_n}{v_n} .
\]
$x_n$ 不变（$x_{n+1}=x_n=x^*$）当且仅当 $\dbinom{x^*}{1}$ 是矩阵的特征向量：
\[
  \begin{pmatrix}a&b\\c&d\end{pmatrix}\binom{x^*}{1}=\lambda\binom{x^*}{1}
  \iff ax^*+b=\lambda x^*,\ \ cx^*+d=\lambda .
\]
消去 $\lambda$ 即得分式不动点方程 $x^*=\dfrac{ax^*+b}{cx^*+d}$，且 $\lambda=cx^*+d$。
若有两个相异特征值 $\lambda_1=cx_1+d,\lambda_2=cx_2+d$，则双不动点构造的公比正是 $\dfrac{\lambda_2}{\lambda_1}$（见探究五）。
\key{结论：矩阵的特征方向（射影点）就是分式递推的不动点；特征值之比就是等比构造的公比。}

\begin{figure}[htbp]
  \centering
  \begin{tikzpicture}[>=Stealth, semithick, font=\footnotesize]
    \draw[->, gray] (-4.6,0)--(4.6,0) node[below right] {$u$};
    \draw[->, gray] (0,-1.5)--(0,3.3) node[above left] {$v$};
    \draw[third, very thick] (-4.4,1)--(4.4,1) node[right] {$v=1$（投影数轴）};
    % 两条特征方向直线都收在 $|y|\le 2.3$ 内，避免穿进底部说明
    \draw[main!70!black, thick, dashed] (-2.0,-2.0)--(3.6,3.6) node[above right] {$v=u\ (\lambda_1=5)$};
    \draw[second!70!black, thick, dashed] (3.9,-1.3)--(-3.9,1.3) node[above left] {$v=-\frac u3\ (\lambda_2=1)$};
    \fill[main] (1,1) circle (1.8pt) node[above right] {$x_1=1$};
    \fill[second] (-3,1) circle (1.8pt) node[below left] {$x_2=-3$};
    \fill (0,0) circle (1pt) node[below left] {$O$};
    \node[gray!70!black] at (0,-2.75) {$M=\begin{pmatrix}2&3\\1&4\end{pmatrix}$ 的两条特征方向直线与投影线 $v=1$ 的交点};
  \end{tikzpicture}
  \caption{射影同构：矩阵 $M=\begin{pmatrix}2&3\\1&4\end{pmatrix}$ 的特征方向 $v=u$、$v=-\frac u3$ 与投影线 $v=1$ 的交点，正是分式递推 $x_{n+1}=\frac{2x_n+3}{x_n+4}$ 的两个不动点 $x_1=1,\ x_2=-3$；公比 $=\lambda_2/\lambda_1=1/5$}
  \label{fig:eigen-fixed}
\end{figure}

\begin{example}
求递推 $a_{n+2}=3a_{n+1}-2a_n$ 在 $a_1=1,\ a_2=3$ 下的通项。
\end{example}
\begin{solution}
这里 $p=3,\ q=-2$，转移矩阵
\[
  A=\begin{pmatrix}0&1\\-2&3\end{pmatrix},\qquad
  \det(\lambda I-A)=\lambda^2-3\lambda+2=0
  \Rightarrow \lambda_1=1,\ \lambda_2=2 .
\]
故设 $a_n=c_1\cdot1^{n-1}+c_2\cdot2^{n-1}$。由初值：
\[
  c_1+c_2=1,\qquad c_1+2c_2=3
  \ \Rightarrow\ c_2=2,\ c_1=-1 .
\]
所以
\[
  a_n=2^n-1 .
\]
检验：$a_3=3a_2-2a_1=9-2=7=2^3-1$。$\surd$
\end{solution}

\begin{conccard}[二级]{特征值法：通项的终极统一模型}{二阶（及高阶）线性常系数递推、斐波那契类递推、重根情形}{特征方程写错符号（应为 $\lambda^2-p\lambda-q=0$）；重根时仍套用两个相异根的公式；初值代入时指数起点 $n-1$ 与 $n$ 混用}
状态向量 $\mathbf v_n=\binom{a_n}{a_{n+1}}$，$\mathbf v_{n+1}=A\mathbf v_n$，$A=\begin{pmatrix}0&1\\q&p\end{pmatrix}$；
特征方程 $\lambda^2-p\lambda-q=0$；
$\lambda_1\ne\lambda_2$ 时 $a_n=c_1\lambda_1^{n-1}+c_2\lambda_2^{n-1}$；
$\lambda_1=\lambda_2=\lambda$ 时 $a_n=(c_1+c_2n)\lambda^{n-1}$。
\end{conccard}

\subsection{Abel 变换：错位相减的机械化}

\key{探究四：离散版分部积分}

\key{（1）Abel 恒等式。}
设 $B_k=\sum_{i=1}^kb_i$，$B_0=0$。则
\[
  \sum_{k=1}^na_kb_k=a_nB_n-\sum_{k=1}^{n-1}(a_{k+1}-a_k)B_k .
\]
推导如下：
\[
\begin{aligned}
  \sum_{k=1}^na_kb_k
  &=\sum_{k=1}^na_k(B_k-B_{k-1})
  =\sum_{k=1}^na_kB_k-\sum_{k=1}^na_kB_{k-1}\\
  &=a_nB_n+\sum_{k=1}^{n-1}a_kB_k-\sum_{k=2}^{n}a_kB_{k-1}
  =a_nB_n-\sum_{k=1}^{n-1}(a_{k+1}-a_k)B_k .
\end{aligned}
\]

\key{（2）与分部积分互译。}
连续版 $\int u\,dv=uv-\int v\,du$；离散版把 $u\to a_k,\ dv\to b_k,\ v\to B_k,\ du\to \Delta a_k$：
\[
\begin{array}{c|c}
  \text{分部积分} & \text{Abel 变换} \\ \hline
  uv & a_nB_n \\
  \displaystyle\int v\,du & \displaystyle\sum_{k=1}^{n-1}(a_{k+1}-a_k)B_k
\end{array}
\]
记忆：\key{「乘积求和 $=$ 末端乘积 $-$ 差分乘部分和的求和」}。

\key{（3）一次秒求「等差 $\times$ 等比」。}
当 $a_k$ 等差时，$\Delta a_k=a_{k+1}-a_k=d$ 是常数；当 $b_k$ 等比时，$B_k$ 有现成公式。代入即得，无需「错位」。

\begin{example}
求 $S_n=\sum_{k=1}^n(2k-1)\cdot2^k$。
\end{example}
\begin{solution}
取 $a_k=2k-1$（等差，$\Delta a_k=2$），$b_k=2^k$（等比，$B_k=2^{k+1}-2$）。由 Abel：
\[
\begin{aligned}
  S_n&=a_nB_n-\sum_{k=1}^{n-1}2\cdot B_k\\
  &=(2n-1)(2^{n+1}-2)-2\sum_{k=1}^{n-1}(2^{k+1}-2)\\
  &=(2n-1)(2^{n+1}-2)-2\bigl(2^{n+1}-4-2(n-1)\bigr)\\
  &=(2n-3)2^{n+1}+6 .
\end{aligned}
\]
检验：$n=1$ 时 $S_1=1\cdot2=2$，公式 $(2-3)\cdot4+6=2$ $\surd$；
$n=2$ 时 $S_2=2+12=14$，公式 $1\cdot8+6=14$ $\surd$。
\end{solution}

\begin{figure}[htbp]
  \centering
  \begin{tikzpicture}[>=Stealth, semithick, font=\footnotesize, xscale=0.55, yscale=0.75]
    % 坐标轴
    \draw[->, gray] (-0.3,0)--(16.2,0) node[below right] {$B_k$};
    \draw[->, gray] (0,-0.3)--(0,5.8) node[above left] {$a_k$};
    % 阶梯柱：$\sum a_k b_k$
    \fill[main!10] (0,0) rectangle (1,5);
    \fill[main!10] (1,0) rectangle (3,4);
    \fill[main!10] (3,0) rectangle (6,3);
    \fill[main!10] (6,0) rectangle (10,2);
    \fill[main!10] (10,0) rectangle (15,1);
    \draw[main!60!black, thick] (0,0) rectangle (1,5);
    \draw[main!60!black, thick] (1,0) rectangle (3,4);
    \draw[main!60!black, thick] (3,0) rectangle (6,3);
    \draw[main!60!black, thick] (6,0) rectangle (10,2);
    \draw[main!60!black, thick] (10,0) rectangle (15,1);
    % 大矩形 $a_nB_n$
    \draw[second!70!black, thick, dashed] (0,0) rectangle (15,1);
    \node[second!70!black] at (7.5,0.5) {$a_nB_n$};
    % 橙色差值区域
    \fill[second!14] (0,1) rectangle (1,5);
    \fill[second!14] (1,1) rectangle (3,4);
    \fill[second!14] (3,1) rectangle (6,3);
    \fill[second!14] (6,1) rectangle (10,2);
    % 标注
    \foreach \x/\lab in {0/$B_0$,1/$B_1$,3/$B_2$,6/$B_3$,10/$B_4$,15/$B_5$}
      \draw (\x,0.08)--(\x,-0.08) node[below] {\lab};
    \node[main!70!black] at (0.5,2.5) {$a_1b_1$};
    \node[main!70!black] at (2,2.5) {$a_2b_2$};
    \node[main!70!black] at (4.5,2.5) {$a_3b_3$};
    \node[main!70!black] at (8,2.5) {$a_4b_4$};
    \node[main!70!black] at (12.5,2.5) {$a_5b_5$};
    % 差值区域标注放在全图上方的空白带（阶梯最高点为 $y=5$）
    \node[second!70!black] at (7.5,5.45) {差值区域：面积 $=\sum\limits_{k=1}^{4}(a_k-a_5)b_k=\sum\limits_{k=1}^{4}B_k=20$};
  \end{tikzpicture}
  \caption{Abel 变换面积示意（取 $a_k=6-k$，$b_k=k$，则 $B_k=\frac{k(k+1)}2$）：蓝色阶梯面积为 $\sum_{k=1}^5a_kb_k=35$；虚线大矩形面积为 $a_5B_5=15$；橙色差值区域面积为 $\sum_{k=1}^{4}B_k=1+3+6+10=20$。三者关系即 $\sum a_kb_k=a_nB_n+\sum_{k=1}^{n-1}(a_k-a_{k+1})B_k$}
  \label{fig:abel-area}
\end{figure}

\begin{conccard}[一级]{Abel 变换：离散分部求和}{等差 $\times$ 等比求和、乘积型数列求和、数列不等式放缩}{把 $a_k$ 与 $b_k$ 角色用反；忘记 $B_0=0$；$n=1$ 时不检验公式退化}
设 $B_k=\sum_{i=1}^kb_i$，则
$\sum_{k=1}^na_kb_k=a_nB_n-\sum_{k=1}^{n-1}(a_{k+1}-a_k)B_k$；
它是 $\int u\,dv=uv-\int v\,du$ 的离散版本；
取 $a_k$ 为等差、$b_k$ 为等比，可三行写出求和。
\end{conccard}

\begin{commonerror}{错位相减与 Abel 变换的高频失误}
\paragraph{失误一：错位相减时项数对不齐。}
「$S_n$ 与 $qS_n$」相减时，错位后首尾会各多出一项，必须明确写出剩余的首项与末项再合并。
\paragraph{失误二：忘记 $B_0=0$ 与 $n=1$ 的检验。}
Abel 公式中 $B_0=0$；当 $n=1$ 时 $\sum_{k=1}^{n-1}$ 为空，公式退化为 $a_1B_1$，要单独验证。
\paragraph{失误三：把 $a_k$ 与 $b_k$ 的角色用反。}
公式中 $a_k$ 是被差分的那个（应选等差／简单数列），$b_k$ 是被求和的那个（应选等比／可求和的数列）。
若反过来，$\Delta a_k$ 不再简单，公式失效。
\end{commonerror}

\subsection{共轭变换：不动点为何生效}

\key{探究五：换坐标系，把复杂运动还原为缩放}

\key{（1）不是「假设相邻两项相等」，而是「换原点」。}
「令 $a_{n+1}=a_n=x^*$」只是记忆口诀。真实动作是：先解出平衡点 $x^*$，再令
\[
  b_n=a_n-x^* .
\]
则
\[
  b_{n+1}=a_{n+1}-x^*=f(a_n)-f(x^*)\approx f'(x^*)(a_n-x^*)=f'(x^*)b_n .
\]
对线性 $f(x)=px+q$，这是精确的：$b_{n+1}=pb_n$。
几何上相当于把数轴原点从 $0$ 平移到 $x^*$：原来的「拉伸 $p$ 倍 $+$ 平移 $q$」复合运动，
在新坐标系下变成纯粹的「以 $x^*$ 为中心的缩放」。这就是\key{共轭变换} $g=h^{-1}\circ f\circ h$（其中 $h(x)=x-x^*$）。

\key{（2）一般情形：$\lambda=f'(x^*)$ 是线性化乘子。}
即使 $f$ 非线性，在不动点附近也有
\[
  b_{n+1}=f(x^*+b_n)-x^*=f'(x^*)b_n+O(b_n^2),
\]
所以 $\lambda=f'(x^*)$ 控制偏差的衰减／放大——这正是探究二稳定性判据的来源。
\key{不动点法生效的底层原因，就是「减去不动点」把非线性递推线性化为等比。}

\key{（3）莫比乌斯分式：配 $\dfrac{a_n-p}{a_n-q}$。}
对 $a_{n+1}=\dfrac{aa_n+b}{ca_n+d}$，设两个相异不动点 $x_1,x_2$。计算偏差：
\[
  a_{n+1}-x_1=\frac{(ad-bc)(a_n-x_1)}{(ca_n+d)(cx_1+d)},\qquad
  a_{n+1}-x_2=\frac{(ad-bc)(a_n-x_2)}{(ca_n+d)(cx_2+d)} .
\]
两式相除，随 $n$ 变化的分母 $(ca_n+d)$ 恰好约去：
\[
  \frac{a_{n+1}-x_1}{a_{n+1}-x_2}
  =\frac{cx_2+d}{cx_1+d}\cdot\frac{a_n-x_1}{a_n-x_2} .
\]
于是 $\left\{\dfrac{a_n-x_1}{a_n-x_2}\right\}$ 是等比数列，公比为 $\dfrac{cx_2+d}{cx_1+d}$。
若只有一个重不动点 $x_0$，则取倒数：
\[
  \frac1{a_{n+1}-x_0}-\frac1{a_n-x_0}=\text{常数},
\]
化为等差数列。

\begin{example}
设 $a_1=2,\ a_{n+1}=\dfrac{2a_n+3}{a_n+4}$，求通项公式。
\end{example}
\begin{solution}
不动点方程 $x=\dfrac{2x+3}{x+4}\Rightarrow x^2+2x-3=0\Rightarrow x_1=1,\ x_2=-3$。
\[
  \frac{a_{n+1}-1}{a_{n+1}+3}
  =\frac{\dfrac{2a_n+3}{a_n+4}-1}{\dfrac{2a_n+3}{a_n+4}+3}
  =\frac{a_n-1}{5a_n+15}
  =\frac15\cdot\frac{a_n-1}{a_n+3} .
\]
故 $\left\{\dfrac{a_n-1}{a_n+3}\right\}$ 是首项 $\dfrac{2-1}{2+3}=\dfrac15$、公比 $\dfrac15$ 的等比数列：
\[
  \frac{a_n-1}{a_n+3}=\left(\frac15\right)^n
  \ \Rightarrow\
  a_n=\frac{5^n+3}{5^n-1} .
\]
检验：$n=1$ 时 $\dfrac{5+3}{5-1}=2$ $\surd$。
\end{solution}

\begin{conccard}[二级]{共轭变换：不动点为何生效}{一阶线性递推、分式（莫比乌斯）递推、嵌套迭代 $f^{(n)}(x^*)$}{把口诀「令 $a_{n+1}=a_n=x^*$」当成真的假设相邻项相等；双不动点时忘记检验 $x_1\ne x_2$；重不动点时仍强行作比}
一阶线性：$b_n=a_n-x^*$，$b_{n+1}=f'(x^*)b_n$（精确）；
一般情形：$b_{n+1}=f'(x^*)b_n+O(b_n^2)$，$\lambda=f'(x^*)$ 是线性化乘子；
莫比乌斯双不动点：$\dfrac{a_n-x_1}{a_n-x_2}$ 等比，公比 $\dfrac{cx_2+d}{cx_1+d}$；
重不动点：$\dfrac1{a_n-x_0}$ 等差。
\end{conccard}

\subsection{强基拓展：生成函数（超出高考考纲）}

\key{探究六：把整个数列打包成一个函数}

\key{温馨提示：本模块涉及形式幂级数，超出高考考纲，仅面向竞赛／强基与学有余力者。}

\key{（1）生成函数：序列 $\to$ 幂级数。}
定义
\[
  G(x)=\sum_{n\ge0}a_nx^n .
\]
它把无穷数列 $\{a_n\}$ 打包成一个形式幂级数：
\begin{itemize}
  \item 下标平移 $a_{n+1}$ 对应「除以 $x$ 并减去首项」：$\sum_{n\ge0}a_{n+1}x^n=\dfrac{G(x)-a_0}{x}$；
  \item 递推关系对应 $G(x)$ 的代数方程；
  \item 解出 $G(x)$ 再展开，$x^n$ 的系数就是 $a_n$。
\end{itemize}

\key{（2）例：斐波那契数列。}
设 $F_0=0,\ F_1=1,\ F_{n+2}=F_{n+1}+F_n$，$G(x)=\sum_{n\ge0}F_nx^n$。则
\[
\begin{aligned}
  G(x)&=x+\sum_{n\ge2}F_nx^n
  =x+\sum_{n\ge2}(F_{n-1}+F_{n-2})x^n\\
  &=x+x\sum_{n\ge2}F_{n-1}x^{n-1}+x^2\sum_{n\ge2}F_{n-2}x^{n-2}
  =x+xG(x)+x^2G(x) .
\end{aligned}
\]
故
\[
  G(x)=\frac{x}{1-x-x^2}
  =\frac{x}{(1-\varphi x)(1-\psi x)},\qquad
  \varphi=\frac{1+\sqrt5}{2},\ \psi=\frac{1-\sqrt5}{2} .
\]
部分分式分解：
\[
  G(x)=\frac1{\sqrt5}\left(\frac1{1-\varphi x}-\frac1{1-\psi x}\right)
  =\frac1{\sqrt5}\sum_{n\ge0}(\varphi^n-\psi^n)x^n ,
\]
比较系数即得比内公式
\[
  F_n=\frac{\varphi^n-\psi^n}{\sqrt5} .
\]
同一个 $G(x)$ 还能给出组合意义（第 7 章）：$\dfrac1{1-x-x^2}$ 的系数是「用 $1\times1$ 与 $1\times2$ 骨牌铺满 $1\times n$ 棋盘」的方法数。

\key{（3）适用边界。}
生成函数擅长：线性递推（齐次／非齐次）、卷积型求和、组合计数（卡特兰数等）。
它把「解递推」变成「解代数方程 $+$ 展开」，是强基与竞赛的核武器；
但解答题若用此法，需补写合法的代数推导。

\begin{conccard}[三级]{生成函数：递推 $\to$ 代数方程}{斐波那契等线性递推、非齐次递推、卷积与组合计数}{超出高考考纲，解答题不能直接套用；忘记首项 $a_0$ 导致平移错误；展开时忽略收敛域（形式幂级数只需形式运算）}
$G(x)=\sum_{n\ge0}a_nx^n$；$\sum_{n\ge0}a_{n+1}x^n=\dfrac{G(x)-a_0}{x}$；
由递推导出 $G(x)$ 的封闭形式，再展开读系数；
斐波那契：$G(x)=\dfrac{x}{1-x-x^2}$，$F_n=\dfrac{\varphi^n-\psi^n}{\sqrt5}$。
\end{conccard}

\begin{method}[变换—求解—逆变换：本章五件降维武器]
\begin{enumerate}
  \item \key{换坐标系}：看到 $a_{n+1}=f(a_n)$，先解不动点 $x^*=f(x^*)$，作 $b_n=a_n-x^*$；
  \item \key{升维}：看到高阶线性递推，写成状态向量 $\mathbf v_{n+1}=A\mathbf v_n$，求特征值；
  \item \key{互逆}：看到求和，想差分 $\Delta$ 与裂项；看到乘积求和，想 Abel 变换；
  \item \key{打包}：递推复杂时，用生成函数 $G(x)$ 把它变成代数方程（强基）；
  \item \key{还原}：在新坐标系算完后，用逆变换 $h^{-1}$ 把答案翻回原数列，并检验首项与边界。
\end{enumerate}
\end{method}

\begin{summary}
本章把数列从「套公式」升级为「离散动力系统 $+$ 降维变换」：
差分 $\Delta$ 与求和 $\sum$ 互为逆运算（离散版微积分基本定理），由 $S_n$ 反求 $a_n$、差分降次判断次数；
一阶递推 $a_{n+1}=f(a_n)$ 是不动点迭代，$x^*=f(x^*)$ 是平衡态，蛛网图把迭代画成阶梯，$|f'(x^*)|<1$ 吸引、$>1$ 排斥；
二阶线性递推对应转移矩阵 $\begin{pmatrix}0&1\\q&p\end{pmatrix}$，特征方程 $\lambda^2-p\lambda-q=0$ 决定通项结构，特征方向即分式递推的不动点方向；
Abel 变换是离散分部积分，把「等差 $\times$ 等比」求和机械化；
共轭变换（换坐标系）解释了不动点法为何生效：$b_n=a_n-x^*$ 把非线性递推线性化为等比；
生成函数（强基拓展）把递推打包成代数方程。

\begin{center}
\footnotesize
\setlength{\tabcolsep}{4pt}
\renewcommand{\arraystretch}{1.35}
\begin{tabular}{p{2.1cm}p{4.9cm}p{4.2cm}p{1.9cm}}
\toprule
降维变换 & 核心操作 & 把什么变成什么 & 对应探究 \\
\midrule
差分 $\Delta$ & $\Delta a_n=a_{n+1}-a_n$ & 求和 $\to$ 相消；高次 $\to$ 低次 & 探究一 \\
共轭（平移／分式） & $b_n=a_n-x^*$；$\dfrac{a_n-x_1}{a_n-x_2}$ & 拉伸$+$平移／非线性分式 $\to$ 等比／等差 & 探究二、五 \\
特征值对角化 & $\mathbf v_{n+1}=A\mathbf v_n$ & 递推 $\to$ 特征值幂 & 探究三 \\
Abel 变换 & $\sum a_kb_k=a_nB_n-\sum\Delta a_kB_k$ & 乘积求和 $\to$ 差分 $\times$ 部分和 & 探究四 \\
生成函数 & $G(x)=\sum a_nx^n$ & 递推 $\to$ 代数方程 & 探究六 \\
\bottomrule
\end{tabular}
\end{center}

\key{数列不是公式的仓库，而是离散动力系统的轨迹；解题的本质是换坐标系，把陌生递推化为等比／等差／代数方程。}
\end{summary}

\end{explore}

% ---------- 跨章联动：降维双引擎（正文收尾；原在答案区之后，已前移） ----------
\subsection{跨章方法论大一统：平移消常数，旋转消交叉项}

前面第 4 章用\key{仿射}压缩椭圆回圆、第 8 章用\key{齐次化}处理条件最值，
本章用\key{不动点}处理递推。它们看似分属三个板块，底层却是同一个动作：
\key{当前的坐标原点没选好，方程里才拖着讨厌的常数尾巴}。把原点精准搬到
锚点上，非齐次项会瞬间蒸发。

\begin{conccard}[二级]{换坐标系的两步走}{适用于所有「带常数项」「带一次项」「带交叉项」的结构}{只背结论，不问原点该搬去哪}
第一步\key{平移}（$x'=x-x_0,\ y'=y-y_0$）：把原点搬到不动点、对称中心或定点 $P$，利用「锚点自身满足方程」这一性质抵消常数项与一次项（第 8 章「手术刀之二」）；第二步\key{旋转}（乘以 $e^{i\theta}$）：把坐标轴对齐到主轴／特征方向，消灭二次交叉项 $Bxy$，并把复特征根的螺旋迭代还原为实数周期震荡。
\end{conccard}

\begin{center}
\renewcommand{\arraystretch}{1.4}
\begin{tabular}{|p{2.8cm}|p{5.0cm}|p{5.4cm}|}
\hline
\textbf{对比维度} & \textbf{数列：不动点法} & \textbf{条件最值／圆锥曲线：齐次化（第 8 章）} \\
\hline
非齐次障碍 & 递推拖着常数尾巴：$a_{n+1}=pa_n+q$ & 平移到定点后仍剩一次项与常数项 \\
\hline
锚点满足的方程 & $x^*=px^*+q$ & $\dfrac{x_0^2}{a^2}+\dfrac{y_0^2}{b^2}=1$ \\
\hline
平移后发生了什么 & 令 $b_n=a_n-x^*$，得 $b_{n+1}=pb_n$，$q$ 归零 & 常数项化为 $\dfrac{x_0^2}{a^2}+\dfrac{y_0^2}{b^2}-1=0$，自动抵消 \\
\hline
剩下的纯粹结构 & 只含公比 $p$ 的等比缩放 & 只含斜率 $k=\dfrac{y'}{x'}$ 的一元二次方程 \\
\hline
\end{tabular}
\end{center}

\begin{problem}[旋转 $45^\circ$：把 $xy=1$ 还原成标准双曲线]
证明反比例函数 $xy=1$ 的图象是等轴双曲线，并写出其焦点位置。
\end{problem}
\begin{solution}
将坐标轴逆时针旋转 $45^\circ$，即
\[
  x=\frac{\sqrt2}{2}(x'-y'),\qquad y=\frac{\sqrt2}{2}(x'+y') .
\]
代入 $xy=1$：
\[
  \frac{\sqrt2}{2}(x'-y')\cdot\frac{\sqrt2}{2}(x'+y')=1
  \ \Longrightarrow\
  \frac{x'^2}{2}-\frac{y'^2}{2}=1 .
\]
交叉项 $xy$ 变成了平方差：这正是以 $x'$ 轴为实轴的等轴双曲线，半实轴 $a=b=\sqrt2$，故 $c=\sqrt{a^2+b^2}=2$，焦点在新坐标系的 $(\pm2,0)$，也就是原坐标系的 $(\sqrt2,\sqrt2)$ 与 $(-\sqrt2,-\sqrt2)$。
\end{solution}

\begin{center}
  \begin{tikzpicture}[x=1.05cm,y=1.05cm,>=stealth,font=\small]
    \draw[->,white!55] (-3,0)--(3,0) node[right,white!70] {$x$};
    \draw[->,white!55] (0,-3)--(0,3) node[above,white!70] {$y$};
    % 旋转后的主轴
    \draw[->,second,dashed,very thick] (-2.3,-2.3)--(2.5,2.5) node[above right,second] {\small $x'$（新实轴）};
    \draw[->,second,dashed,very thick] (2.3,-2.3)--(-2.5,2.5) node[above left,second] {\small $y'$（新虚轴）};
    \draw[->,second,thick] (1.1,0) arc (0:45:1.1);
    \node[second] at (1.15,0.38) {\small $45^\circ$};
    % 双曲线 xy=1 的两支
    \draw[main,very thick,domain=0.42:2.6,samples=80,smooth] plot(\x,{1/\x});
    \draw[main,very thick,domain=-2.6:-0.42,samples=80,smooth] plot(\x,{1/\x});
    \node[main] at (1.9,1.15) {\small $xy=1$};
    \node[white!70] at (0,-3.4) {\small 交叉项 $xy$ 与旋转 $45^\circ$ 后的新主轴 $x',y'$};
  \end{tikzpicture}\\[0.3em]
  {\small\itshape 同一曲线的两种读法：原坐标系下是交叉项方程 $xy=1$；把轴转 $45^\circ$ 后成为标准等轴双曲线 $\frac{x'^2}{2}-\frac{y'^2}{2}=1$}
\end{center}

\paragraph{第二步：旋转也管数列。} 当二阶递推 $a_{n+2}=pa_{n+1}+qa_n$ 的特征方程判别式 $\Delta<0$ 时，两根为共轭复数 $\lambda_{1,2}=r(\cos\theta\pm i\sin\theta)$。此时几何图像不再是沿直线趋近不动点，而是每步\key{缩放 $r$ 倍并旋转 $\theta$ 角}的对数螺线；通项写成
\[
  a_n=r^{\,n-1}\bigl(c_1\cos n\theta+c_2\sin n\theta\bigr) .
\]
$r=1$ 且 $\theta=\frac{2\pi}{T}$ 时数列退化为纯周期，例如 $a_{n+2}=a_{n+1}-a_n$ 的 $\lambda=e^{\pm i\frac{\pi}{3}}$，周期为 $6$。这与「把轴旋转到主轴方向」是同一件事：都是把耦合项转没。

\begin{figure}[htbp]
  \centering
  \begin{tikzpicture}[>=Stealth, semithick, font=\footnotesize, x=1.15cm, y=1.15cm]
    \draw[->, gray] (-1.9,0)--(1.95,0);
    \draw[->, gray] (0,-1.35)--(0,1.55);
    % 迭代点 v0..v5（M = r R(40°)，r=0.8），相邻两点长度比 0.8、辐角差 40°
    \coordinate (v0) at (1.5,0.2);
    \coordinate (v1) at (0.8163,0.8942);
    \coordinate (v2) at (0.0402,0.9679);
    \coordinate (v3) at (-0.4733,0.6138);
    \coordinate (v4) at (-0.6058,0.1327);
    \coordinate (v5) at (-0.4394,-0.2302);
    \fill (0,0) circle (1.2pt) node[below right] {$O$};
    \draw[main, very thick] (v0)--(v1)--(v2)--(v3)--(v4)--(v5);
    \foreach \i/\p/\pos in {0/v0/above right,1/v1/right,2/v2/above left,3/v3/left,4/v4/below left,5/v5/below}
      \fill[second] (\p) circle (1.5pt) node[\pos] {$\mathbf v_{\i}$};
    \node[gray!70!black, align=center] at (0,-1.15) {$M=rR(\theta)$，$r=0.8,\ \theta=40^\circ$：每步长度乘 $r$、方向加 $\theta$；\\$r<1$ 螺线收拢到唯一不动点 $O$，$r>1$ 远离，$r=1$ 且 $\theta=2\pi/T$ 时为周期 $T$};
  \end{tikzpicture}
  \caption{复特征值的迭代图像：$v_{n+1}=Mv_n$ 且 $M=rR(\theta)$ 时，每一步都把向量缩放 $r$ 倍并旋转 $\theta$ 角，点列沿螺线运动——这就是「实部控伸缩、虚部控旋转」的离散版本}
  \label{fig:spiral-complex}
\end{figure}

\begin{commonerror}{卷面上禁止出现「由仿射变换知」「由极点极线知」}
高观点只用于草稿纸预判答案与取等位置，答题卡上必须写考纲认可的初等推导（联立、韦达定理、作差构造等比）。详见附录~\ref{app:legal}。
\end{commonerror}

\begin{inclasstest}
\begin{enumerate}
  \item 设 $S_n=2n^2-n$，求数列 $\{a_n\}$ 的通项公式；\level{易}
  \item 设 $a_1=1,\ a_{n+1}=2a_n+3$，求通项并判断 $\{a_n\}$ 是否收敛；\level{中}
  \item 用 Abel 变换求 $\sum_{k=1}^nk\cdot2^k$；\level{中}
  \item 求 $a_{n+2}=4a_{n+1}-4a_n$ 在 $a_1=1,\ a_2=4$ 下的通项（重根情形）。\level{中}
\end{enumerate}
\end{inclasstest}

\begin{afterclass}
\subsection*{基础巩固}
\begin{enumerate}
  \item 已知 $\Delta a_n=2n+1$，且 $a_1=1$，求 $a_n$；
  \item 证明 $\sum_{k=1}^n\Delta a_k=a_{n+1}-a_1$，并用它求 $\sum_{k=1}^n\dfrac1{k(k+1)}$。
\end{enumerate}

\subsection*{进阶提升}
\begin{enumerate}
  \item 用特征值法求 $a_{n+2}=2a_{n+1}+3a_n$ 在 $a_1=1,\ a_2=5$ 下的通项；
  \item 设 $a_1=0,\ a_{n+1}=\dfrac{3a_n+1}{a_n+3}$，求通项公式。
\end{enumerate}

\subsection*{拓展探究}
\begin{enumerate}
  \item （强基／竞赛，超出高考考纲）设 $a_0=1,\ a_{n+1}=2a_n+1$，用生成函数求通项公式。
\end{enumerate}

\subsection*{教师版答案}
\begin{solution}
\textbf{基础 1：}由离散版微积分基本定理，
\[
  a_n=a_1+\sum_{k=1}^{n-1}\Delta a_k
  =1+\sum_{k=1}^{n-1}(2k+1)
  =1+2\cdot\frac{(n-1)n}2+(n-1)=n^2 .
\]
（检验：$n=1$ 时 $a_1=1$。）
\end{solution}
\begin{solution}
\textbf{基础 2：}求和时中间项全部抵消：
\[
  \sum_{k=1}^n\Delta a_k=\sum_{k=1}^n(a_{k+1}-a_k)=a_{n+1}-a_1 .
\]
又 $\dfrac1{k(k+1)}=\dfrac1k-\dfrac1{k+1}=\Delta\!\left(-\dfrac1k\right)$，故
\[
  \sum_{k=1}^n\frac1{k(k+1)}
  =\sum_{k=1}^n\Delta\!\left(-\frac1k\right)
  =\left(-\frac1{n+1}\right)-(-1)=\frac n{n+1} .
\]
\end{solution}
\begin{solution}
\textbf{进阶 1：}$A=\begin{pmatrix}0&1\\3&2\end{pmatrix}$，特征方程
\[
  \lambda^2-2\lambda-3=0\Rightarrow\lambda_1=3,\ \lambda_2=-1 .
\]
设 $a_n=c_1\cdot3^{n-1}+c_2\cdot(-1)^{n-1}$，由 $a_1=1,\ a_2=5$：
\[
  c_1+c_2=1,\quad 3c_1-c_2=5
  \ \Rightarrow\ c_1=\frac32,\ c_2=-\frac12 .
\]
故
\[
  a_n=\frac{3^n+(-1)^n}2 .
\]
（检验：$a_3=2\cdot5+3\cdot1=13=\frac{27-1}2$。）
\end{solution}
\begin{solution}
\textbf{进阶 2：}不动点方程 $x=\dfrac{3x+1}{x+3}\Rightarrow x^2=1\Rightarrow x_1=1,\ x_2=-1$。
\[
  \frac{a_{n+1}-1}{a_{n+1}+1}
  =\frac{(3a_n+1)-(a_n+3)}{(3a_n+1)+(a_n+3)}
  =\frac{2a_n-2}{4a_n+4}
  =\frac12\cdot\frac{a_n-1}{a_n+1} .
\]
首项 $\dfrac{0-1}{0+1}=-1$，故 $\dfrac{a_n-1}{a_n+1}=-\left(\dfrac12\right)^{n-1}$，解得
\[
  a_n=\frac{1-\left(\frac12\right)^{n-1}}{1+\left(\frac12\right)^{n-1}}
  =\frac{2^{n-1}-1}{2^{n-1}+1} .
\]
（检验：$n=1$ 时 $a_1=0$；$n=2$ 时 $a_2=\frac13=\frac{3\cdot0+1}{0+3}$。）
\end{solution}
\begin{solution}
\textbf{拓展：}设 $G(x)=\sum_{n\ge0}a_nx^n$。由 $a_{n+1}=2a_n+1$：
\[
  \frac{G(x)-a_0}{x}=2G(x)+\sum_{n\ge0}x^n=2G(x)+\frac1{1-x} .
\]
代入 $a_0=1$：
\[
  G(x)-1=2xG(x)+\frac{x}{1-x}
  \ \Rightarrow\
  G(x)(1-2x)=\frac1{1-x}
  \ \Rightarrow\
  G(x)=\frac1{(1-x)(1-2x)} .
\]
部分分式：
\[
  G(x)=\frac2{1-2x}-\frac1{1-x}
  =\sum_{n\ge0}(2\cdot2^n-1)x^n
  =\sum_{n\ge0}(2^{n+1}-1)x^n ,
\]
故 $a_n=2^{n+1}-1$。（检验：$a_0=1,\ a_1=3$。）
\end{solution}

\reading{拓展阅读 11}{数列：离散动力系统的轨迹}
\end{afterclass}

% 高观点导学案：导数与局部线性化引擎（第 6 章·高观点版）
% 主线：局部线性化（以直代曲）→ 导数即变化率（差分的连续版）→ 五大不等式的结构本质（支撑边界）
%       → 切线放缩与隐零点 → 条件最值四大标准动作 → 含参恒成立/能成立的集合边界
% 核心思想：把「导数」与「不等式」统一在「以直代曲」这一个引擎下——曲线在一点附近用切线代替，
%           不等式是为点集划定不可逾越的边界；解题四步：识别结构 → 局部线性化 → 求最值 → 验取等。

\section{导数与局部线性化引擎（高观点版）}

\begin{methodbox}{多框架对读}
\small
\begin{center}
\begin{tabular}{c|p{2.9cm}|p{5.0cm}|p{3.1cm}}
\toprule
视角 & 本章看点 & 别的视角怎么说 & 哪侧最省力 \\
\midrule
代数 & 求导法则、恒等变形 & 符号对象的机械运算 & 硬算单调区间 \\
几何 & 切线、增减、凹凸的形状 & 曲线在放大镜下变直线 & 看图猜极值 \\
集合 & 零点集、符号区间 & 定义域值域的边界 & 界参数范围 \\
映射 & 局部线性化：曲线 $\to$ 直线 & 换到切线坐标系再比大小 & 五不等式一源 \\
\bottomrule
\end{tabular}
\end{center}
卡的使用：证不等式先想「是不是某个函数的单调性」，再把切线当替代函数（总图见第 0 章图~\ref{fig:four-view-map}）。
\end{methodbox}

第 1 章确立了函数与单调性，第 5 章把差分 $\Delta$ 确立为离散意义下的变化率；
本章把「导数」与「不等式」统一在同一个引擎下——\key{局部线性化（以直代曲）}：
任何光滑曲线在一点的无穷小邻域内都可以用它的切线近似代替，而这条切线的斜率就是导数。
于是「求导数、判符号」成为判定单调性／极值／最值的统一操作；
五大核心不等式分别是射影定理、内积公理、齐次化（完整展开见第 8 章）、度量三角公理与凸凹函数切线支撑这五种结构下的「支撑边界」；
条件最值与含参成立性问题，则归结为「边界值能否取到」。
六个模块层层展开：
\key{局部线性化}（放大镜下曲成直）$\to$
\key{导数即变化率}（差分的连续版，统一判定单调性／极值／最值）$\to$
\key{五大不等式的结构本质}（一张对比表：射影、内积、齐次、度量、凸支撑）$\to$
\key{切线放缩与隐零点}（$e^x\ge x+1$、$\ln x\le x-1$ 的严格证明与多步取等检验）$\to$
\key{条件最值四大标准动作}（1 的代换、消元主元、拆项分离、切线放缩）$\to$
\key{含参恒成立／能成立}（点集位置、集合包含与量词的值域边界）。
核心思想：\key{导数与不等式是同一个引擎的两面——以直代曲给出局部逼近，支撑边界给出整体控制；解题的本质是「识别结构 $\to$ 局部线性化 $\to$ 求最值 $\to$ 验取等」。}

\begin{learnobjective}
\begin{itemize}
  \item 理解局部线性化公式 $f(x)\approx f(x_0)+f'(x_0)(x-x_0)$ 与「以直代曲」的几何含义，会用切线作近似计算；同一个一阶逼近在第 5 章的离散递推中就是线性化乘子 $f'(x^*)$；
  \item 掌握导数的极限定义与导函数，理解它是差分的连续版本（$\Delta a_n$ 是 $a_{n+1}-a_n$，$f'(x)$ 是 $\frac{\Delta y}{\Delta x}$ 的极限）；
  \item 会用导数符号统一判定单调性、极值与最值，理解驻点与极值点的区别；
  \item 从射影定理、内积公理、齐次化、度量公理、凸凹性五个结构本质理解五大不等式，并记住各自的取等条件（齐次化专论见第 8 章）；
  \item 会用作差求导严格证明 $e^x\ge x+1$、$\ln x\le x-1$，能在多步放缩中检验各步取等条件是否同时成立；
  \item 掌握隐零点「设而不求、整体代换」的基本操作；
  \item 掌握条件最值的四大标准动作（1 的代换、消元主元、拆项分离常数、切线放缩），会检验取等的可行性；
  \item 会用点集上下位置与集合包含关系统一解释恒成立／能成立问题，掌握量词到最值的四类转化与含参二次不等式的分类讨论树。
\end{itemize}
\end{learnobjective}

\begin{preview}
\subsection*{知识回顾}
\begin{itemize}
  \item 函数 $y=f(x)$ 在区间 $[x_0,x_0+\Delta x]$ 上的平均变化率是 \blank；
  \item 导数的定义：$f'(x_0)=$ \blank（写成 $\Delta x\to0$ 时的极限）；
  \item 若 $f'(x)>0$ 在区间 $I$ 上恒成立，则 $f$ 在 $I$ 上 \blank；
  \item 基本不等式：$a,b>0$ 时 \blank $\ge$ \blank，当且仅当 \blank 时取等；
  \item 一元二次方程 $ax^2+bx+c=0$ 的判别式 $\Delta=$ \blank，$\Delta>0$ 时有两个 \blank 实根。
\end{itemize}

\subsection*{教材速读}
阅读「导数」一章教材，完成下列问题：
\begin{enumerate}
  \item 可导与连续的关系是：可导 \blank 连续，连续 \blank 可导（填「一定」或「不一定」）；
  \item 写出三个基本求导公式：$(x^n)'=$ \blank，$(e^x)'=$ \blank，$(\ln x)'=$ \blank；
  \item 函数极值的第一充分条件：在 $x_0$ 两侧 $f'(x)$ 异号，则 $x_0$ 是 \blank；
  \item 闭区间上连续函数的最值只可能在 \blank 与 \blank 处取得。
\end{enumerate}
\end{preview}

\begin{explore}

\subsection{局部线性化：放大镜下曲成直}

\key{探究一：切线是曲线在一点处的「最佳直线近似」}

\key{（1）从平均变化率到瞬时变化率。}
函数在 $x_0$ 附近的平均变化率是 $\dfrac{\Delta y}{\Delta x}=\dfrac{f(x_0+\Delta x)-f(x_0)}{\Delta x}$。
让 $\Delta x\to0$，若极限存在，就得到导数
\[
  f'(x_0)=\lim_{\Delta x\to0}\frac{f(x_0+\Delta x)-f(x_0)}{\Delta x}.
\]
这正是第 5 章差分 $\Delta a_n=a_{n+1}-a_n$ 的连续版：把「步长为 1 的差分」换成「步长趋于 0 的差商极限」。

\key{（2）局部线性化公式：以直代曲。}
过点 $P(x_0,f(x_0))$ 的切线方程为
\[
  y=f(x_0)+f'(x_0)(x-x_0),
\]
而在 $x_0$ 附近，曲线与切线的误差是关于 $(x-x_0)$ 的高阶无穷小：
\[
  f(x)=f(x_0)+f'(x_0)(x-x_0)+o(x-x_0).
\]
也就是说，\key{在 $x_0$ 的无穷小邻域内，曲线可以用切线代替——这就是「以直代曲」}。
几何显微镜放得越大，曲线看起来越像一条直线（图 \ref{fig:local-linear}）。

\begin{figure}[htbp]
  \centering
  \begin{tikzpicture}[>=Stealth, semithick, font=\footnotesize]
    % ---------- 左：全局视角 ----------
    \begin{scope}
      \draw[->, gray] (-1.3,0)--(3.0,0) node[below right] {$x$};
      \draw[->, gray] (0,-2.3)--(0,4.4) node[above left] {$y$};
      \draw[main!70!black, very thick, domain=-1:2, smooth, samples=60] plot (\x, {\x*\x});
      \node[main!70!black] at (1.72,3.15) {$y=x^2$};
      \draw[second, thick, domain=-0.4:2.5] plot (\x, {2*\x-1});
      \node[second] at (2.64,3.9) {$y=2x-1$};
      \fill (1,1) circle (1.6pt) node[above right] {$P(1,1)$};
      \draw[dashed, gray] (0.85,0.7) rectangle (1.15,1.3);
      % 说明文字整体右移，避开左图的 $y$ 轴
      \node[gray!70!black, align=left] at (0.2,-2.15) {虚线框为右图的放大区域};
    \end{scope}
    % ---------- 右：放大视角 ----------
    \begin{scope}[xshift=5.8cm]
      \draw[->, gray] (0.78,0.55)--(1.26,0.55) node[above right] {$x$};
      \draw[->, gray] (0.75,0.58)--(0.75,1.46) node[above left] {$y$};
      \draw[second, thick, dashed, domain=0.84:1.16] plot (\x, {2*\x-1});
      \draw[main!70!black, very thick, domain=0.84:1.16, smooth, samples=60] plot (\x, {\x*\x});
      \fill (1,1) circle (1.6pt) node[above right] {$P(1,1)$};
      \node[gray!70!black, align=center] at (1.01,-0.1) {放大约 $25$ 倍后，曲线（实线）\\与切线（虚线）几乎重合};
    \end{scope}
  \end{tikzpicture}
  \caption{以直代曲：抛物线 $y=x^2$ 在 $P(1,1)$ 处的切线为 $y=2x-1$。左图全局视角下曲线明显弯曲；右图把视野放大到 $P$ 附近，曲线与切线在视觉上无法区分——$f(x)\approx f(x_0)+f'(x_0)(x-x_0)$}
  \label{fig:local-linear}
\end{figure}

\begin{example}
用切线近似求 $\sqrt{4.1}$ 与 $\sin31^\circ$（精确到 $0.001$）。
\end{example}
\begin{solution}
取 $f(x)=\sqrt x$，在 $x_0=4$ 处：$f(4)=2$，$f'(x)=\frac1{2\sqrt x}$，$f'(4)=\frac14$，故
\[
  \sqrt{4.1}\approx 2+\frac14\cdot0.1=2.025 .
\]
（真值 $\sqrt{4.1}\approx2.0248$，误差约 $1.5\times10^{-4}$。）

取 $g(x)=\sin x$，在 $x_0=\frac\pi6$ 处：$g'(x)=\cos x$，$g'\!\left(\frac\pi6\right)=\frac{\sqrt3}{2}$，自变量增量
$\Delta x=\frac\pi{180}$，故
\[
  \sin31^\circ\approx \frac12+\frac{\sqrt3}{2}\cdot\frac{\pi}{180}
  \approx0.5+0.01511=0.515 .
\]
（真值 $\sin31^\circ\approx0.51504$。）
\end{solution}

\begin{conccard}[一级]{局部线性化：以直代曲}{光滑曲线的近似计算、误差估计、理解导数的几何意义}{把「近似」当成「相等」；放大倍数不够时强行替代；忘记公式只在 $x_0$ 附近有效}
$f'(x_0)=\lim\limits_{\Delta x\to0}\dfrac{f(x_0+\Delta x)-f(x_0)}{\Delta x}$；
切线 $y=f(x_0)+f'(x_0)(x-x_0)$；
$f(x)=f(x_0)+f'(x_0)(x-x_0)+o(x-x_0)$，即以直代曲；
导数是差分 $\Delta a_n$ 的连续版本。
\end{conccard}

\subsection{导数即变化率：统一判定单调性、极值、最值}

\key{探究二：$f'(x)$ 的正负号就是曲线的「上下坡」}

\key{（1）导函数与符号判定。}
把每一点的导数收集起来得到导函数 $f'(x)$。在区间 $I$ 上：
\begin{itemize}
  \item $f'(x)>0$ 恒成立 $\iff$ $f$ 在 $I$ 上单调递增；
  \item $f'(x)<0$ 恒成立 $\iff$ $f$ 在 $I$ 上单调递减。
\end{itemize}
这与第 5 章用差分判数列单调性完全同构：$\Delta a_n>0\iff a_{n+1}>a_n$（递增数列）。

\key{（2）极值：$f'$ 变号的点。}
费马引理：可导函数在极值点处必有 $f'(x_0)=0$（驻点）。
第一充分条件：$f'$ 在 $x_0$ 两侧异号 $\iff$ $x_0$ 是极值点。汇总如下：

\begin{center}\footnotesize\renewcommand{\arraystretch}{1.3}
\begin{tabular}{c|c|c}
$f'(x)$ 在 $x_0$ 附近的符号 & $f(x)$ 的形态 & 极值 \\\hline
$+$，然后 $-$ & 先增后减 & 极大值 \\
$-$，然后 $+$ & 先减后增 & 极小值 \\
不变号（如 $y=x^3$ 在 $0$ 处） & 单调 & 无极值 \\
\end{tabular}
\end{center}

\key{（3）最值：端点 $+$ 极值点。}
闭区间 $[a,b]$ 上的连续函数，最值只可能在端点 $a,b$ 与极值点处取得，比较这些点的函数值即可。

\begin{example}
求 $f(x)=x^3-3x$ 的单调区间与极值；若在 $[-2,2]$ 上求最值。
\end{example}
\begin{solution}
$f'(x)=3x^2-3=3(x+1)(x-1)$。
\begin{itemize}
  \item $x<-1$ 或 $x>1$ 时 $f'(x)>0$；$-1<x<1$ 时 $f'(x)<0$。
  \item 故 $f$ 在 $(-\infty,-1)$、$(1,+\infty)$ 上递增，在 $(-1,1)$ 上递减。
  \item 极大值 $f(-1)=-1+3=2$；极小值 $f(1)=1-3=-2$。
\end{itemize}
在 $[-2,2]$ 上额外比较端点：$f(-2)=-8+6=-2$，$f(2)=8-6=2$。
故最小值为 $-2$（在 $x=-2$ 与 $x=1$ 处），最大值为 $2$（在 $x=-1$ 与 $x=2$ 处）。
$\surd$
\end{solution}

\begin{commonerror}{驻点、极值点与最值的三个高频失误}
\paragraph{失误一：以为 $f'(x_0)=0$ 就一定是极值点。}
$f'(x_0)=0$ 只是必要条件。反例：$f(x)=x^3$ 在 $x=0$ 处 $f'(0)=0$，但两侧 $f'>0$ 不变号，不是极值点。
必须检验两侧符号是否改变。
\paragraph{失误二：忽略定义域。}
例如 $f(x)=x\ln x$ 的定义域是 $(0,+\infty)$，讨论单调性时不能把 $x\le0$ 也算进去；
求闭区间最值时也不能漏掉端点。
\paragraph{失误三：把「单点处 $f'(x_0)>0$」当成「区间上递增」。}
导数符号必须在整个区间上保持，才能推出单调性；个别点的符号说明不了整体趋势。
\end{commonerror}

\begin{conccard}[一级]{导数符号法：单调性、极值、最值一条龙}{含参数的函数性质讨论、恒成立问题中求最值}{驻点当成极值点；漏定义域；漏端点；把 $f'(x)\ge0$ 中的孤立零点误判为单调性破坏}
$f'(x)>0\Rightarrow$ 递增，$f'(x)<0\Rightarrow$ 递减；
极值点处 $f'(x_0)=0$ 且两侧异号；
闭区间最值 $=$ 端点值与极值值中的最大／最小者。
\end{conccard}

\subsection{五大不等式的结构本质}

\key{探究三：五条不等式，五个「支撑边界」}

同一个「$\ge$」在不同结构下有不同来历。看懂它们的结构本质，取等条件就会自己跳出来。

\begin{center}\footnotesize\setlength{\tabcolsep}{4pt}\renewcommand{\arraystretch}{1.4}
\begin{tabular}{p{1.9cm}p{3.5cm}p{5.2cm}p{2.6cm}}
\toprule
不等式 & 核心形式 & 结构本质 & 取等条件 \\
\midrule
基本不等式 & $\dfrac{a+b}{2}\ge\sqrt{ab}\ (a,b>0)$ & 半圆射影定理：直径 $a+b$ 的半径 $\ge$ 斜边上的高 $\sqrt{ab}$ & $a=b$ \\[2pt]
柯西不等式 & $(a_1^2+a_2^2)(b_1^2+b_2^2)\ge(a_1b_1+a_2b_2)^2$ & 向量内积公理 $|\vec u|^2|\vec v|^2\ge(\vec u\cdot\vec v)^2$（即 $|\cos\theta|\le1$） & $\vec u\parallel\vec v$（共线） \\[2pt]
权方和不等式 & $\dfrac{a_1^2}{b_1}+\dfrac{a_2^2}{b_2}\ge\dfrac{(a_1+a_2)^2}{b_1+b_2}\ (b_i>0)$ & 柯西在 $\left(\dfrac{a_i}{\sqrt{b_i}}\right)$ 与 $(\sqrt{b_i})$ 下的齐次化特例 & $\dfrac{a_1}{b_1}=\dfrac{a_2}{b_2}$ \\[2pt]
绝对值三角不等式 & $\bigl||a|-|b|\bigr|\le|a\pm b|\le|a|+|b|$ & 度量空间的距离三角不等式（三角形两边之和 $\ge$ 第三边） & 共线同向（右端）／反向（左端） \\[2pt]
切线放缩 & $e^x\ge x+1$，$\ln x\le x-1$ & 凸／凹函数的切线支撑边界（Jensen 不等式的一阶情形） & $x=0$／$x=1$ \\
\bottomrule
\end{tabular}
\end{center}

注：齐次化作为第一性原理的完整方法链见第 8 章；本章只把它用作
条件最值的一个入口，权方和则是柯西（内积公理）的齐次化特例。

\key{（1）基本不等式 $=$ 半圆射影定理。}
作直径 $AB=a+b$ 的半圆，圆心 $O$，则半径 $R=\dfrac{a+b}{2}$。
设分界点 $C$ 满足 $AC=a$、$CB=b$，过 $C$ 作 $CP\perp AB$ 交半圆于 $P$。
因为直径所对圆周角为直角，$\triangle APB$ 是直角三角形，由射影定理
\[
  PC^2=AC\cdot CB=ab\ \Rightarrow\ PC=\sqrt{ab}.
\]
又在直角三角形 $OPC$ 中 $OP$ 是斜边，故 $PC\le OP=R$，即
\[
  \sqrt{ab}\le\frac{a+b}{2},
\]
当且仅当 $C$ 与 $O$ 重合（即 $a=b$）时取等（图 \ref{fig:amgm-semicircle}）。

\begin{figure}[htbp]
  \centering
  \begin{tikzpicture}[>=Stealth, semithick, font=\footnotesize, scale=1.05]
    % 直径 AB=a+b，取 a=2.4,b=1.2 => R=1.8，C 位于 x=0.6，PC=sqrt(2.88)≈1.697
    \draw[->, gray] (-2.35,0)--(2.35,0);
    \draw[thick] (-1.8,0)--(1.8,0);
    \draw[thick] (1.8,0) arc (0:180:1.8);
    \fill (-1.8,0) circle (1.5pt) node[below left] {$A$};
    \fill (1.8,0) circle (1.5pt) node[below right] {$B$};
    \fill (0,0) circle (1.5pt) node[below] {$O$};
    \fill (0.6,0) circle (1.5pt) node[below] {$C$};
    \fill (0.6,1.697) circle (1.5pt) node[above right] {$P$};
    \draw[main!70!black, thick] (0.6,0)--(0.6,1.697);
    % 标签放在 $PC$ 右侧、斜边 $PB$ 下方的空档处
    \node[main!70!black] at (0.95,0.55) {$\sqrt{ab}$};
    \draw[second, thick, dashed] (0,0)--(0,1.8);
    % 半径标签锚定在虚线右侧起排，避免压线
    \node[second, anchor=west] at (0.72,1.25) {$\dfrac{a+b}{2}$};
    \draw[third!70!black, thick] (-1.8,0)--(0.6,1.697)--(1.8,0);
    \node at (-0.6,-0.4) {$a$};
    \node at (1.2,-0.4) {$b$};
    \node[gray!70!black, align=center] at (0,-2.15) {$PC^2=AC\cdot CB=ab\Rightarrow PC=\sqrt{ab}\le OP=\dfrac{a+b}{2}$\\（当且仅当 $C$ 与 $O$ 重合即 $a=b$ 时取等）};
  \end{tikzpicture}
  \caption{半圆射影定理：直径 $AB=a+b$，半径 $R=\frac{a+b}{2}$ 永远是斜边上的高 $PC=\sqrt{ab}$ 的上界——这就是基本不等式的几何模型}
  \label{fig:amgm-semicircle}
\end{figure}

\key{（2）柯西 $=$ 向量内积公理；权方和 $=$ 柯西的齐次化。}
$|\vec u\cdot\vec v|=|\vec u|\,|\vec v|\,|\cos\theta|\le|\vec u|\,|\vec v|$，平方即得二维柯西，
取等当且仅当两向量共线。代数证法（判别式法）：
\begin{derivation}
令 $\varphi(x)=(a_1x+b_1)^2+(a_2x+b_2)^2\ge0$ 对一切实数 $x$ 成立，展开得
\[
  (a_1^2+a_2^2)x^2+2(a_1b_1+a_2b_2)x+(b_1^2+b_2^2)\ge0 .
\]
二次函数恒非负，故判别式 $\Delta\le0$：
\[
  4(a_1b_1+a_2b_2)^2-4(a_1^2+a_2^2)(b_1^2+b_2^2)\le0,
\]
即 $(a_1^2+a_2^2)(b_1^2+b_2^2)\ge(a_1b_1+a_2b_2)^2$。
取等当且仅当存在 $x$ 使 $a_1x+b_1=a_2x+b_2=0$，即 $(a_1,a_2)\parallel(b_1,b_2)$。
\end{derivation}
权方和只是把 $\vec u=\left(\dfrac{a_1}{\sqrt{b_1}},\dfrac{a_2}{\sqrt{b_2}}\right)$、
$\vec v=(\sqrt{b_1},\sqrt{b_2})$ 代入柯西：
\[
  \left(\frac{a_1^2}{b_1}+\frac{a_2^2}{b_2}\right)(b_1+b_2)
  \ge\left(\frac{a_1}{\sqrt{b_1}}\cdot\sqrt{b_1}+\frac{a_2}{\sqrt{b_2}}\cdot\sqrt{b_2}\right)^2
  =(a_1+a_2)^2 .
\]

\begin{conccard}[二级]{五大不等式的结构本质与取等条件}{证明不等式、选择放缩工具、检验取等}{只背公式不看结构；张冠李戴（该用柯西时硬凑基本不等式）；忽视取等条件导致「最值取不到」}
基本不等式＝半圆射影定理（$a=b$ 取等）；
柯西＝内积公理（共线取等）；
权方和＝柯西的齐次化（$\frac{a_1}{b_1}=\frac{a_2}{b_2}$ 取等）；
绝对值三角＝度量三角公理（共线取等）；
切线放缩＝凸凹函数的切线支撑（切点取等）。
\end{conccard}

\subsection{切线放缩与隐零点}

\key{探究四：切线是全局支撑边界；隐零点靠「设而不求」}

\key{（1）$e^x\ge x+1$ 与 $\ln x\le x-1$ 的严格证明。}
\begin{derivation}
设 $g(x)=e^x-x-1$，则 $g'(x)=e^x-1$。
$x<0$ 时 $g'<0$，$x>0$ 时 $g'>0$，故 $g$ 在 $x=0$ 处取最小值 $g(0)=0$。
所以
\[
  e^x\ge x+1\quad(\text{当且仅当 }x=0\text{ 取等}).
\]
设 $h(x)=x-1-\ln x\ (x>0)$，则 $h'(x)=1-\frac1x$。
$0<x<1$ 时 $h'<0$，$x>1$ 时 $h'>0$，故 $h$ 在 $x=1$ 处取最小值 $h(1)=0$。
所以
\[
  \ln x\le x-1\quad(\text{当且仅当 }x=1\text{ 取等}).
\]
几何含义：$y=e^x$ 是凸函数，恒在切线 $y=x+1$ 上方；$y=\ln x$ 是凹函数，恒在切线 $y=x-1$ 下方（图 \ref{fig:tangent-exp}）。
\end{derivation}

\begin{figure}[htbp]
  \centering
  \begin{tikzpicture}[>=Stealth, semithick, font=\footnotesize]
    \draw[->, gray] (-2.2,0)--(1.95,0) node[below right] {$x$};
    \draw[->, gray] (0,-1.2)--(0,5.1) node[above left] {$y$};
    \fill[main!10] plot[domain=-2:1.6, samples=60] (\x, {exp(\x)})
      -- plot[domain=1.6:-2, samples=60] (\x, {\x+1}) -- cycle;
    \draw[main!70!black, very thick, domain=-2:1.6, smooth, samples=80] plot (\x, {exp(\x)});
    \draw[second, thick, domain=-2:1.6] plot (\x, {\x+1});
    \fill (0,1) circle (1.6pt) node[above right] {$(0,1)$};
    \node[main!70!black] at (0.45,3.25) {$y=e^x$};
    \node[second] at (-1.85,-0.55) {$y=x+1$};
    % 结论文字移到左上空白区，避免被 $e^x$ 曲线或切线穿过
    \node[gray!70!black] at (-0.35,4.35) {$e^x\ge x+1$};
    \node[gray!70!black] at (-0.35,3.95) {仅在 $x=0$ 取等};
  \end{tikzpicture}
  \caption{$y=e^x$ 恒在切线 $y=x+1$ 上方（凸函数的切线支撑）：阴影为差值区域 $e^x-(x+1)\ge0$，两者仅在切点 $(0,1)$ 接触}
  \label{fig:tangent-exp}
\end{figure}

\key{（2）多步放缩：各步取等条件必须同时成立。}
\begin{example}
证明：对 $x>0$，$e^x-\ln x>2$。
\end{example}
\begin{solution}
由 $e^x\ge x+1$（取等 $x=0$）与 $-\ln x\ge1-x$（取等 $x=1$）相加得
\[
  e^x-\ln x\ge (x+1)+(1-x)=2 .
\]
但两式取等条件分别是 $x=0$ 与 $x=1$，不可能同时成立，故等号永远取不到：
\[
  e^x-\ln x>2 .
\]
\key{多步放缩时，每一步的取等条件必须相容，否则最后的「$\ge$」其实是严格的「$>$」。}
\end{solution}

\key{（3）隐零点：设而不求，整体代换。}
当 $f'(x)=0$ 的解 $x_0$ 无法用初等形式写出时，不必硬解：设 $x_0$ 为零点，
利用 $f'(x_0)=0$ 提供的等量关系做整体代换，把超越结构化为代数结构。

\begin{example}
证明：对一切 $x>-2$，$e^x>\ln(x+2)$。
\end{example}
\begin{solution}
设 $f(x)=e^x-\ln(x+2)$，$f'(x)=e^x-\dfrac1{x+2}$，$f''(x)=e^x+\dfrac1{(x+2)^2}>0$，
故 $f'$ 在 $(-2,+\infty)$ 上严格递增。
又
\[
  f'\!\left(-\tfrac12\right)=e^{-1/2}-\frac23<0,\qquad f'(0)=1-\frac12>0,
\]
所以存在唯一的 $x_0\in\left(-\dfrac12,0\right)$ 使 $f'(x_0)=0$，即
\[
  e^{x_0}=\frac1{x_0+2}\quad\Longrightarrow\quad \ln(x_0+2)=-x_0 .
\]
$f$ 在 $(-2,x_0)$ 上递减、在 $(x_0,+\infty)$ 上递增，最小值为
\[
  f(x_0)=e^{x_0}-\ln(x_0+2)=\frac1{x_0+2}+x_0 .
\]
由 $x_0\in\left(-\dfrac12,0\right)$ 得 $\dfrac1{x_0+2}>\dfrac12$ 且 $x_0>-\dfrac12$，故
\[
  f(x_0)>\frac12-\frac12=0 .
\]
因此 $f(x)\ge f(x_0)>0$，即 $e^x>\ln(x+2)$ 对一切 $x>-2$ 成立。$\surd$
\end{solution}

\begin{conccard}[一级]{切线放缩与隐零点}{证明含 $e^x$、$\ln x$ 的不等式；求不可解方程的极值}{忘记检验各步取等条件；不敢把「$\ge$」收紧为「$>$」；找到 $x_0$ 后不会用 $f'(x_0)=0$ 代换；忽略 $x_0$ 的取值范围}
$e^x\ge x+1$（取等 $x=0$），$\ln x\le x-1$（取等 $x=1$）；
多步放缩取等条件不相容 $\Rightarrow$ 严格不等；
隐零点：设 $f'(x_0)=0$，用该等式整体代换，并确定 $x_0$ 的范围。
\end{conccard}

\subsection{条件最值：四大标准动作}

\key{探究五：把多元、超越的结构，逐个「降维」成对勾函数}

\key{动作一：齐次化与「1」的代换（入口；完整方法链见第 8 章）。}
已知 $ma+nb=k$ 求 $\dfrac pa+\dfrac qb$ 的最值：乘以 $\dfrac{ma+nb}k=1$，展开后用基本不等式。

\begin{example}
$a,b>0$，$a+b=1$，求 $\dfrac1a+\dfrac4b$ 的最小值。
\end{example}
\begin{solution}
「1 的代换」：
\[
  \frac1a+\frac4b=\left(\frac1a+\frac4b\right)(a+b)
  =5+\frac ba+\frac{4a}b\ge5+2\sqrt4=9,
\]
当且仅当 $\dfrac ba=\dfrac{4a}b$ 即 $b=2a$ 时取等，结合 $a+b=1$ 得 $a=\dfrac13,b=\dfrac23$。
「权方和秒杀」：
\[
  \frac1a+\frac4b=\frac{1^2}{a}+\frac{2^2}{b}\ge\frac{(1+2)^2}{a+b}=9,
\]
取等条件同为 $a:b=1:2$。
\end{solution}

\key{动作二：消元与主元转化（双变量 $\to$ 对勾函数 $t+\dfrac kt$）。}

\begin{example}
$x,y>0$ 满足 $x+2y+xy=30$，求 $x+2y$ 的最小值。
\end{example}
\begin{solution}
由等式解出 $y=\dfrac{30-x}{x+2}$。令 $t=x+2$，则 $t>2$，且 $y>0\Rightarrow t<32$，故 $t\in(2,32)$。
\[
  x+2y=t-2+\frac{2(32-t)}t=t+\frac{64}t-4\ge2\sqrt{64}-4=12,
\]
当且仅当 $t=\dfrac{64}t$ 即 $t=8\in(2,32)$ 时取等，此时 $x=6,y=3$（检验：$6+6+18=30$）。
\end{solution}

\key{动作三：拆项分离常数。}

\begin{example}
求 $f(x)=\dfrac{x^2+3x+5}{x+1}$ 在 $x>-1$ 时的最小值。
\end{example}
\begin{solution}
令 $t=x+1>0$，则 $x=t-1$，分子
\[
  x^2+3x+5=(t-1)^2+3(t-1)+5=t^2+t+3,
\]
故
\[
  f(x)=t+1+\frac3t\ge1+2\sqrt3,
\]
当且仅当 $t=\dfrac3t$ 即 $t=\sqrt3$（$x=\sqrt3-1$）时取等。
\end{solution}

\key{动作四：切线放缩放去超越函数。}

\begin{example}
证明：$x>0$ 时 $\ln(x+1)+\dfrac1{x+1}>1$。
\end{example}
\begin{solution}
在 $\ln u\le u-1$ 中取 $u=\dfrac1t\ (t>0)$，得 $\ln t\ge1-\dfrac1t$（取等 $t=1$）。
令 $t=x+1>1$，则
\[
  \ln(x+1)\ge1-\frac1{x+1}\quad\Longrightarrow\quad
  \ln(x+1)+\frac1{x+1}\ge1 .
\]
取等需要 $t=1$ 即 $x=0$，但 $x>0$，故严格成立：
\[
  \ln(x+1)+\frac1{x+1}>1 .
\]
\key{用切线（一阶 Taylor）把 $\ln$、$e^x$ 替换为整式，是消去超越函数的标准动作。}
\end{solution}

\begin{commonerror}{易错诊断室：连续放缩的取等冲突与换元定义域}
\paragraph{反例一：连续两次用基本不等式，取等条件互相打架。}
已知 $x,y>0$，$x+2y=1$，某同学求 $\dfrac1{x^2}+\dfrac1{y^2}$ 的「最小值」：
\[
  \frac1{x^2}+\frac1{y^2}\ge\frac{2}{xy},\qquad
  xy=\frac{x\cdot2y}{2}\le\frac12\left(\frac{x+2y}{2}\right)^2=\frac18,
\]
于是 $\dfrac1{x^2}+\dfrac1{y^2}\ge\dfrac{2}{1/8}=16$。
问题：第一步取等要求 $x=y$，第二步取等要求 $x=2y$，二者不能同时成立，
所以 $16$ 根本取不到（真实最小值大于 $16$）。
\key{每写一次放缩，就把取等条件记下来；所有条件必须相容，最值才可能取到。}

\paragraph{反例二：换元后忽略新元的取值范围。}
求 $f(x)=\dfrac{x^2+5}{\sqrt{x^2+4}}$ 的最小值。令 $t=\sqrt{x^2+4}$，则 $t\ge2$，且
\[
  f(x)=t+\frac1t .
\]
若直接写 $t+\dfrac1t\ge2$，取等要求 $t=1$，但 $t\ge2$，取不到！
正确做法：$\varphi(t)=t+\dfrac1t$ 在 $[2,+\infty)$ 上递增，故最小值为
\[
  \varphi(2)=2+\frac12=\frac52 .
\]
\key{换元必先写出新元的定义域；取等点若不在定义域内，必须改用单调性求最值。}
\end{commonerror}

\begin{conccard}[一级]{条件最值四大标准动作}{二元条件最值、分式函数最值、含超越函数的不等式}{取等条件冲突未检验；换元后忘定义域；该用权方和时硬展开；切线放缩用于非切点附近却当作等式}
（1）1 的代换：乘 $\frac{ma+nb}k=1$ 展开（或直接用权方和）；齐次化的完整展开见第 8 章，本节只作条件最值的入口；
（2）消元主元：解出一元，换 $t$ 凑 $t+\frac kt$；
（3）拆项分离常数：$t=x+1$ 类换元降幂；
（4）切线放缩：用 $e^x\ge x+1$ 等消去超越项。
每步都要检验：取等条件相容吗？在定义域内吗？
\end{conccard}

\subsection{含参恒成立与能成立：点集位置与集合包含}

\key{探究六：把「成立性」翻译成「边界值的大小关系」}

\key{（1）点集上下位置 $\to$ 恒成立。}
$f(x)\ge a$ 在 $D$ 上恒成立，几何上是直线 $y=a$ 始终位于曲线 $y=f(x)$ 的下方（允许相切），
等价于
\[
  a\le f(x)_{\min}\quad(x\in D).
\]
类比地，$f(x)>a$ 恒成立需 $a\le f_{\min}$ 且验证最小值点处是否严格（最小值能取到时需 $a<f_{\min}$）。

\key{（2）一元二次含参不等式：分类讨论树。}
以 $ax^2+px+q>0$ 为例，固定顺序：先看二次项系数 $a$（零／正／负），
再看判别式 $\Delta$（正／零／负），最后比较两根大小（图 \ref{fig:quad-tree}）。

\begin{figure}[htbp]
  \centering
  \begin{tikzpicture}[>=Stealth, semithick, font=\footnotesize,
    box/.style={draw, rounded corners, align=center, inner sep=3pt, text width=2.2cm}]
    \node[box, text width=3.8cm] (root) at (0,0) {解关于 $x$ 的含参不等式\\$ax^{2}+px+q>0$\\($a,p,q$ 为参数)};
    \node[box] (l) at (-4.6,-1.6) {$a=0$\\退化为一次};
    \node[box] (m) at (0,-1.6) {$a>0$\\开口向上};
    \node[box] (r) at (4.6,-1.6) {$a<0$\\开口向下};
    \node[box] (ll) at (-5.2,-3.3) {$px+q>0$\\按 $p$ 符号定向};
    \node[box] (m1) at (-2.6,-3.3) {$\Delta>0$\\两根之外};
    \node[box] (m2) at (0,-3.3) {$\Delta=0$\\$x\ne -\frac{p}{2a}$};
    \node[box] (m3) at (2.6,-3.3) {$\Delta<0$\\解集为 $\mathbb R$};
    \node[box] (rr) at (5.2,-3.3) {乘 $-1$ 变号\\化为 $a>0$};
    \draw (root)--(l) (root)--(m) (root)--(r);
    \draw (l)--(ll) (m)--(m1) (m)--(m2) (m)--(m3) (r)--(rr);
    \node[gray!70!black, align=center, text width=13cm] at (0,-4.9)
      {注：$\Delta>0$ 时两根为 $x_{1,2}=\dfrac{-p\pm\sqrt{\Delta}}{2a}$，\\还必须比较两根大小，然后才能写下解集区间。};
  \end{tikzpicture}
  \caption{含参一元二次不等式的分类讨论树：二次项系数（零／正／负）$\to$ 判别式 $\Delta$（正／零／负）$\to$ 两根大小}
  \label{fig:quad-tree}
\end{figure}

\begin{example}
解关于 $x$ 的不等式 $ax^2-(a+1)x+1>0$。
\end{example}
\begin{solution}
因式分解：$ax^2-(a+1)x+1=(x-1)(ax-1)>0$。
\begin{itemize}
  \item $a=0$：化为 $-(x-1)>0$，解集 $(-\infty,1)$；
  \item $a>0$：两根为 $1$ 与 $\dfrac1a$。
    $a>1$ 时 $\dfrac1a<1$，解集 $\left(-\infty,\dfrac1a\right)\cup(1,+\infty)$；
    $a=1$ 时为 $(x-1)^2>0$，解集 $\{x\mid x\ne1\}$；
    $0<a<1$ 时 $\dfrac1a>1$，解集 $(-\infty,1)\cup\left(\dfrac1a,+\infty\right)$；
  \item $a<0$：开口向下，两根 $\dfrac1a<0<1$，解集 $\left(\dfrac1a,1\right)$。
\end{itemize}
\end{solution}

\key{（3）量词问题：值域边界的对偶。}
「任意」「存在」混合时，统一翻译为两个函数值域的边界关系（设最值都能取到）：

\begin{center}\footnotesize\renewcommand{\arraystretch}{1.3}
\begin{tabular}{c|c}
量词结构 & 等价条件 \\\hline
$\forall x_1\in D_1,\ \forall x_2\in D_2:\ f(x_1)\ge g(x_2)$ & $f_{\min}\ge g_{\max}$ \\[2pt]
$\forall x_1\in D_1,\ \exists x_2\in D_2:\ f(x_1)\ge g(x_2)$ & $f_{\min}\ge g_{\min}$ \\[2pt]
$\exists x_1\in D_1,\ \forall x_2\in D_2:\ f(x_1)\ge g(x_2)$ & $f_{\max}\ge g_{\max}$ \\[2pt]
$\exists x_1\in D_1,\ \exists x_2\in D_2:\ f(x_1)\ge g(x_2)$ & $f_{\max}\ge g_{\min}$ \\
\end{tabular}
\end{center}

\begin{example}
设 $f(x)=e^x-x$，$g(x)=x^2-2x+a$，$x\in[0,2]$。若对任意 $x_1\in[0,2]$，都存在 $x_2\in[0,2]$ 使 $f(x_1)\ge g(x_2)$，求 $a$ 的取值范围。
\end{example}
\begin{solution}
条件等价于 $f_{\min}\ge g_{\min}$。
$f'(x)=e^x-1\ge0$ 在 $[0,2]$ 上成立，$f$ 递增，$f_{\min}=f(0)=1$。
$g(x)=(x-1)^2+a-1$，$g_{\min}=g(1)=a-1$。
故
\[
  1\ge a-1\ \Rightarrow\ a\le2 .
\]
\end{solution}

\begin{conccard}[二级]{恒成立／能成立：点集位置与值域边界}{含参不等式恒成立、双变量量词问题、参数范围求解}{把 $\forall$ 与 $\exists$ 的最值方向搞反；忽略最值是否能取到；严格不等式端点处理错误；二次不等式分类讨论顺序混乱}
第 0 章的集合版（值域子集包含）在这里换成最值语言：$f(x)\ge a$ 恒成立 $\iff a\le f_{\min}$；
$\forall x_1,\exists x_2:\ f(x_1)\ge g(x_2)\iff f_{\min}\ge g_{\min}$（其余三种对偶见表）；
含参二次不等式按「$a$ 的符号 $\to\Delta\to$ 两根大小」三步分类。
\end{conccard}

\begin{method}[局部线性化引擎的使用流程]
\begin{enumerate}
  \item \key{识别结构}：见 $e^x$、$\ln x$ 想切线放缩；见双变量等式想 1 的代换／消元；见含参成立性想最值与值域边界；
  \item \key{局部线性化}：在关键点求导、写切线，用 $f(x)\approx f(x_0)+f'(x_0)(x-x_0)$ 把非线性问题线性化；
  \item \key{求最值}：导数符号法找单调区间与极值，再比较端点；
  \item \key{验取等}：记录每一步的取等条件，检查是否相容、是否在定义域内——不相容则「$\ge$」改「$>$」，不在定义域内则改用单调性；
  \item \key{回代检验}：用特殊值（如 $x=0,1$）或量纲、符号快速验证结论。
\end{enumerate}
\end{method}

\begin{summary}
本章把「导数」与「不等式」统一在\key{局部线性化（以直代曲）}这一个引擎之下：
曲线在一点附近由切线 $f(x_0)+f'(x_0)(x-x_0)$ 近似代替，导数是差分的连续版本；
由 $f'(x)$ 的符号统一判定单调性、极值与最值；
五大核心不等式分别起源于射影定理、内积公理、齐次化、度量三角公理与凸凹函数的切线支撑，取等条件各不相同；
$e^x\ge x+1$、$\ln x\le x-1$ 的严格证明与多步取等的相容性检验，构成切线放缩的完整操作规程；
隐零点靠「设而不求、整体代换」；条件最值有四大标准动作（1 的代换、消元主元、拆项分离、切线放缩）；
含参恒成立／能成立问题则归结为点集的上下位置与值域的边界大小。

\begin{center}
\footnotesize
\setlength{\tabcolsep}{4pt}
\renewcommand{\arraystretch}{1.35}
\begin{tabular}{p{2.3cm}p{4.6cm}p{4.6cm}p{1.9cm}}
\toprule
引擎部件 & 核心操作 & 把什么变成什么 & 对应探究 \\
\midrule
局部线性化 & $f(x)\approx f(x_0)+f'(x_0)(x-x_0)$ & 非线性 $\to$ 线性；变化率 $\to$ 切线斜率 & 探究一 \\
导数符号法 & $f'(x)>0$ 增、$f'(x)<0$ 减 & 单调性／极值／最值 $\to$ 符号判定 & 探究二 \\
结构本质表 & 射影／内积／齐次／度量／凸支撑 & 不等式 $\to$ 几何与代数边界 & 探究三 \\
切线放缩 & $e^x\ge x+1$，$\ln x\le x-1$ & 超越 $\to$ 整式（须验取等） & 探究四 \\
隐零点代换 & 设 $x_0$，用 $f'(x_0)=0$ 整体代换 & 不可解方程 $\to$ 代数关系 & 探究四 \\
条件最值四动作 & 1 的代换／消元主元／拆项／放缩 & 多变量、超越 $\to$ 一元低次 & 探究五 \\
成立性量词化 & $\forall\exists$ 四类转化 & 参数成立性 $\to$ 值域边界 & 探究六 \\
\bottomrule
\end{tabular}
\end{center}

\key{导数与不等式是同一个引擎的两面：以直代曲给出局部逼近，支撑边界给出整体控制；解题的本质是「识别结构 $\to$ 局部线性化 $\to$ 求最值 $\to$ 验取等」。}
\end{summary}

\end{explore}

\begin{inclasstest}
\begin{enumerate}
  \item 用切线近似求 $\sqrt{4.02}$ 的近似值（精确到 $0.001$）；\level{易}
  \item 求 $f(x)=x^3-3x^2-9x+5$ 的单调区间与极值；\level{中}
  \item 证明：对任意 $x>0$，$e^x-\ln x>2$；\level{中}
  \item 解关于 $x$ 的不等式 $ax^2-(a+1)x+1>0$。\level{难}
\end{enumerate}
\end{inclasstest}

\begin{afterclass}
\subsection*{基础巩固}
\begin{enumerate}
  \item 求函数 $f(x)=x\ln x\ (x>0)$ 的最小值；
  \item 证明：对任意实数 $x$，$e^x\ge ex$（提示：求 $f(x)=e^x-ex$ 的最小值）。
\end{enumerate}

\subsection*{进阶提升}
\begin{enumerate}
  \item 已知 $a,b>0$ 且 $a+2b=1$，求 $\dfrac1a+\dfrac1b$ 的最小值；
  \item 设 $f(x)=e^x-ax\ (a>0)$，讨论其单调性；若 $f(x)\ge0$ 对一切 $x\in\mathbb R$ 恒成立，求 $a$ 的取值范围。
\end{enumerate}

\subsection*{拓展探究}
\begin{enumerate}
  \item （强基／竞赛）把权方和不等式推广到三元：证明
  $\dfrac{a_1^2}{b_1}+\dfrac{a_2^2}{b_2}+\dfrac{a_3^2}{b_3}\ge\dfrac{(a_1+a_2+a_3)^2}{b_1+b_2+b_3}\ (b_i>0)$，
  并用它求 $a,b,c>0$、$a+b+c=3$ 时 $\dfrac1a+\dfrac4b+\dfrac9c$ 的最小值。
\end{enumerate}

\subsection*{教师版答案}
\begin{solution}
\textbf{基础 1：}$f'(x)=\ln x+1$。令 $f'(x)=0$ 得 $x=\dfrac1e$。
$0<x<\dfrac1e$ 时 $f'<0$，$x>\dfrac1e$ 时 $f'>0$，故最小值为
\[
  f\!\left(\frac1e\right)=\frac1e\ln\frac1e=-\frac1e .
\]
\end{solution}
\begin{solution}
\textbf{基础 2：}设 $f(x)=e^x-ex$，$f'(x)=e^x-e$。
$x<1$ 时 $f'<0$，$x>1$ 时 $f'>0$，故 $f$ 在 $x=1$ 处取最小值 $f(1)=0$，
即 $e^x\ge ex$，当且仅当 $x=1$ 取等（这正是 $y=e^x$ 在 $(1,e)$ 处的切线 $y=ex$）。
\end{solution}
\begin{solution}
\textbf{进阶 1：}「1 的代换」：
\[
  \frac1a+\frac1b=\left(\frac1a+\frac1b\right)(a+2b)
  =3+\frac{2b}{a}+\frac ab\ge3+2\sqrt2,
\]
当且仅当 $\dfrac{2b}{a}=\dfrac ab$ 即 $a=\sqrt2\,b$ 时取等，结合 $a+2b=1$ 得
$b=\dfrac1{2+\sqrt2}$，$a=\dfrac{\sqrt2}{2+\sqrt2}$。最小值 $3+2\sqrt2$。
\end{solution}
\begin{solution}
\textbf{进阶 2：}$f'(x)=e^x-a$。$x<\ln a$ 时 $f'<0$，$x>\ln a$ 时 $f'>0$，
故 $f$ 在 $(-\infty,\ln a)$ 递减、在 $(\ln a,+\infty)$ 递增，最小值为
\[
  f(\ln a)=a-a\ln a=a(1-\ln a).
\]
$f(x)\ge0$ 恒成立 $\iff a(1-\ln a)\ge0\iff\ln a\le1\iff a\le e$。
故 $a$ 的取值范围是 $(0,e]$。
\end{solution}
\begin{solution}
\textbf{拓展：}由柯西不等式，取
$\vec u=\left(\dfrac{a_1}{\sqrt{b_1}},\dfrac{a_2}{\sqrt{b_2}},\dfrac{a_3}{\sqrt{b_3}}\right)$，
$\vec v=(\sqrt{b_1},\sqrt{b_2},\sqrt{b_3})$，得
\[
  \left(\frac{a_1^2}{b_1}+\frac{a_2^2}{b_2}+\frac{a_3^2}{b_3}\right)(b_1+b_2+b_3)
  \ge(a_1+a_2+a_3)^2,
\]
即所要的不等式，取等当且仅当 $\dfrac{a_1}{b_1}=\dfrac{a_2}{b_2}=\dfrac{a_3}{b_3}$。
应用：
\[
  \frac1a+\frac4b+\frac9c=\frac{1^2}{a}+\frac{2^2}{b}+\frac{3^2}{c}
  \ge\frac{(1+2+3)^2}{a+b+c}=\frac{36}{3}=12,
\]
当且仅当 $a:b:c=1:2:3$ 且 $a+b+c=3$，即 $a=\dfrac12,b=1,c=\dfrac32$ 时取等。最小值为 $12$。
\end{solution}
\reading{拓展阅读 12}{导数：以直代曲的局部线性化}
\end{afterclass}
% 高观点导学案：计数原理、概率测度空间与统计线性逼近（第 7 章·高观点版）
% 主线：把「不确定性」纳入集合框架——样本空间是集合，概率是其上的归一化计数测度；
%       条件概率 = 样本全集缩小后的空间缩放；全概率/贝叶斯 = 完备互斥划分；
%       回归直线 = 散点云的竖直最小二乘投影；独立性检验 = 小概率事件原理的统计反证法；
%       压轴收官：新定义情境题的四步破题法（翻译 → 特例 → 归纳 → 严格证明）。

\section{计数原理、概率测度空间与统计线性逼近（高观点版）}

\begin{methodbox}{多框架对读}
\small
\begin{center}
\begin{tabular}{c|p{2.9cm}|p{5.0cm}|p{3.1cm}}
\toprule
视角 & 本章看点 & 别的视角怎么说 & 哪侧最省力 \\
\midrule
代数 & 计数公式、期望方差 & 符号对象的机械运算 & 硬算分布列 \\
几何 & 散点、回归直线、面积测度 & 图形的位置与拟合 & 看趋势、估相关 \\
集合 & 样本空间、事件、完备划分 & 事件＝子集；概率＝测度 & 界互斥与独立 \\
映射 & $X:\Omega\to\mathbb R$；最小二乘＝投影 & 换到残差正交的坐标系 & 一图读懂相关 \\
\bottomrule
\end{tabular}
\end{center}
卡的使用：概率题先画维恩图与树图（集合），统计题先想「投影到哪个子空间」（总图见第 0 章图~\ref{fig:four-view-map}）。
\end{methodbox}

第 0 章确立了集合与逻辑这一全书语言，第 4 章把向量确立为空间的通用线性语言，第 6 章把曲线局部化为直线；
本章把「不确定性」也纳入同一个框架——随机试验的样本空间 $\Omega$ 是一个集合，概率是定义在其上的\key{归一化计数测度}；
条件概率 $P(B\mid A)$ 的本质是把样本全集 $\Omega$ 缩小为子集 $A$ 之后的\key{空间缩放}；
线性回归方程是二维散点云在竖直方向上的\key{最小二乘投影直线}，独立性检验则是「小概率事件原理」支配下的\key{统计学反证法}。
四个模块层层展开：
\key{计数原理的集合本质}（并集基数与笛卡尔积基数，排列＝有序单射、组合＝无序子集族，捆绑／插空／隔板三技法）$\to$
\key{概率空间的归一化测度}（古典与几何概型的统一除法，条件概率＝空间缩放，全概率与贝叶斯＝完备互斥划分）$\to$
\key{统计线性逼近}（百分位数规则，回归直线必过样本中心，独立性检验的统计反证）$\to$
\key{新定义情境题破题模板}（翻译符号集合 $\to$ 构造小规模特例 $\to$ 归纳猜想通项 $\to$ 反证法／数学归纳法严格证明）。
核心思想：\key{随机性不是结构的对立面，而是有限集合上的测度；统计不是死记公式，而是「投影逼近 $+$ 小概率反证」；压轴新定义题的本质是「把陌生符号翻译回熟悉的集合与结构」。}

\begin{learnobjective}
\begin{itemize}
  \item 理解分类加法原理即互斥子集并集的基数公式 $|A\cup B|=|A|+|B|$，分步乘法原理即笛卡尔积的基数公式 $|A\times B|=|A|\cdot|B|$；
  \item 理解排列是有序单射序列、组合是无序 $k$ 元子集族，掌握 $A_n^k$ 与 $C_n^k$ 的关系 $A_n^k=C_n^k\cdot k!$；
  \item 掌握捆绑法、插空法、隔板法的适用结构，会把「相邻／不相邻／相同元素分堆」问题翻译为标准计数模型；
  \item 会把古典概型与几何概型统一为归一化测度 $P(A)=\mu(A)/\mu(\Omega)$，说出概率三条公理（非负性、规范性、可加性）；
  \item 理解条件概率 $P(B\mid A)=\dfrac{P(A\cap B)}{P(A)}$ 是样本空间缩小为 $A$ 后的重新归一化，会用全概率公式与贝叶斯公式；
  \item 掌握百分位数的计算规则，理解回归直线 $\hat y=\hat bx+\hat a$ 是竖直残差平方和最小的投影直线且必过 $(\bar x,\bar y)$；
  \item 理解独立性检验 $K^2$ 统计量背后的「小概率事件原理」，能正确陈述检验结论；
  \item 掌握新定义情境题的四步破题法（读懂符号集合 $\to$ 小规模特例 $\to$ 归纳猜想 $\to$ 反证法／数学归纳法严格证明）。
\end{itemize}
\end{learnobjective}

\begin{preview}
\subsection*{知识回顾}
\begin{itemize}
  \item 若 $A\cap B=\varnothing$，则 $|A\cup B|=$ \blank；
  \item 完成一件事有两类互斥的途径，第一类有 $n_1$ 种方法、第二类有 $n_2$ 种方法，则共有 \blank 种方法；
  \item 古典概型：$P(A)=$ \blank；几何概型：$P(A)=$ \blank；
  \item 样本相关系数 $r$ 满足 $|r|$ \blank $1$，$r>0$ 表示两个变量 \blank 相关；
  \item 反证法的逻辑结构：假设结论不成立，推出 \blank，从而肯定原结论。
\end{itemize}

\subsection*{教材速读}
阅读「计数原理」「概率」「成对数据的统计」相关教材，完成下列问题：
\begin{enumerate}
  \item 排列数公式 $A_n^m=$ \blank；组合数公式 $C_n^m=$ \blank；
  \item 条件概率公式 $P(B\mid A)=$ \blank（要求 \blank $\ne0$）；
  \item 全概率公式：若 $A_1,\dots,A_n$ 是 $\Omega$ 的完备互斥划分，则 $P(B)=$ \blank；
  \item 回归直线 $\hat y=\hat bx+\hat a$ 必过样本中心 \blank，且 $\hat a=$ \blank（用 $\bar x,\bar y,\hat b$ 表示）；
  \item 独立性检验统计量 $K^2=$ \blank（设 $2\times2$ 列联表中四个数据为 $a,b,c,d$）。
\end{enumerate}
\end{preview}

\begin{explore}

\subsection{计数原理：并集、笛卡尔积与三类经典技法}

\key{探究一：计数不是数数，而是集合的度量}

\key{（1）两个原理的集合本质。}
设完成一件事的所有可能结果构成集合 $\Omega$。若 $\Omega$ 被分成两个互斥的子集 $\Omega=A\sqcup B$（即 $A\cap B=\varnothing$），则
\[
  |\Omega|=|A\sqcup B|=|A|+|B|,
\]
这就是\key{分类加法原理}：各类之间「不重不漏」，求总数就是求互斥子集并集的基数。
若完成一件事要分两步：第一步的结果集为 $A$、第二步的结果集为 $B$，则所有可能结果组成笛卡尔积
\[
  \Omega=A\times B=\{(a,b)\mid a\in A,\ b\in B\},\qquad |\Omega|=|A\times B|=|A|\cdot|B|,
\]
这就是\key{分步乘法原理}：每一步的选择共同确定一个有序对，求总数就是求积集的基数。
一句话对照：加法管「或」（不交并），乘法管「且」（有序积）。

\begin{example}
从 5 名男同学和 4 名女同学中选出 3 人组成学习小组（不分职务），有多少种选法？若从中选出 3 人分别担任语文、数学、英语三科课代表，又有多少种选法？
\end{example}
\begin{solution}
第一问只关心「选出的集合」，是无序的 $3$ 元子集族，用组合：
\[
  C_9^3=\frac{9!}{3!\,6!}=84 .
\]
第二问关心「谁对应哪一科」，是有序安排：先选后排，由分步乘法原理
\[
  C_9^3\cdot A_3^3=84\times6=504\quad(\text{或直接 }A_9^3=9\times7\times8=504).
\]
两问相差 $3!=6$ 倍——正是「无序」与「有序」的差别。
\end{solution}

\key{（2）排列与组合：有序单射与无序子集。}
从 $n$ 个不同元素中取出 $k$ 个排成一列，相当于构造一个单射 $\{1,2,\dots,k\}\hookrightarrow\{1,2,\dots,n\}$，其个数为排列数
\[
  A_n^k=n(n-1)\cdots(n-k+1)=\frac{n!}{(n-k)!}.
\]
从 $n$ 个不同元素中取出 $k$ 个组成子集（不计顺序），其个数为组合数
\[
  C_n^k=\frac{A_n^k}{k!}=\frac{n!}{k!\,(n-k)!},
\]
因为「先选后排」＝「选子集」乘「子集全排列」：$A_n^k=C_n^k\cdot k!$——这正是分步乘法原理的一次应用。
特别地，$C_n^k=C_n^{n-k}$（选 $k$ 个留下等价于选 $n-k$ 个淘汰）。

\key{（3）三类经典技法的结构对号。}
\begin{itemize}
  \item \key{捆绑法（必相邻）}：把必须相邻的 $m$ 个元素捆成一个「大元素」，先与其他元素一起排，再解开内部排列：
    $n$ 个元素中指定 $m$ 个相邻的排法 $=(n-m+1)!\cdot A_m^m$。
  \item \key{插空法（不相邻）}：先排其余 $n$ 个元素（产生 $n+1$ 个空位，含两端），再把 $m$ 个元素插入 $m$ 个不同的空位：$A_n^n\cdot A_{n+1}^{\,m}$。（此处的 $n$ 指「其余元素个数」，与上一行捆绑法中 $n$ 为总元素数的口径不同。）
  \item \key{隔板法（相同元素分堆）}：不定方程 $x_1+x_2+\cdots+x_m=n\ (x_i\in\mathbb N^*)$ 的正整数解，对应把 $n$ 个相同小球排成一排后在 $n-1$ 个空隙中插入 $m-1$ 块板，解数为
    \[
      C_{n-1}^{\,m-1}.
    \]
    若允许 $x_i=0$，先令 $x_i'=x_i+1$ 化为正整数解，得 $C_{n+m-1}^{\,m-1}$。
\end{itemize}

\begin{example}
（1）6 人站成一排，甲、乙必须相邻，有多少种站法？
（2）6 人站成一排，甲、乙互不相邻，有多少种站法？
（3）求 $x+y+z=10$ 的正整数解的组数。
\end{example}
\begin{solution}
（1）捆绑：甲乙捆成一个大元素，与其余 4 人共 5 个元素全排列，甲乙内部再交换：
\[
  A_5^5\cdot A_2^2=120\times2=240 .
\]
（2）插空：先排其余 4 人，产生 5 个空位，甲乙选 2 个空位排列：
\[
  A_4^4\cdot A_5^2=24\times20=480 .
\]
（3）隔板：10 个相同小球排一排，9 个空隙插 2 块板分成 3 份（每份至少 1）：
\[
  C_9^2=36 .
\]
\end{solution}

\begin{figure}[htbp]
  \centering
  \begin{tikzpicture}[>=Stealth, semithick, font=\footnotesize]
    % ---------- 左：文氏图 ----------
    \draw[gray] (-1.7,-1.6) rectangle (3.9,1.6);
    \begin{scope}
      \clip (0,0) circle (1.15);
      \fill[second!20] (-1.8,-1.7) rectangle (4.0,1.7);
    \end{scope}
    \begin{scope}
      \clip (1.4,0) circle (1.15);
      \fill[third!20] (-1.8,-1.7) rectangle (4.0,1.7);
    \end{scope}
    \begin{scope}
      \clip (0,0) circle (1.15);
      \fill[main!32] (1.4,0) circle (1.15);
    \end{scope}
    \draw (0,0) circle (1.15);
    \draw (1.4,0) circle (1.15);
    \node at (-0.62,0) {$A$};
    \node at (2.02,0) {$B$};
    \node[main!70!black] at (0.7,-0.05) {\small \color{white}$A\cap B$};
    \node[gray!70!black] at (0.7,1.4) {$A\cup B$};
    \node[gray!70!black] at (-1.5,1.4) {$U$};
    % ---------- 右：卡诺图式网格 ----------
    \begin{scope}[xshift=5.4cm, yshift=0.5cm]
      \draw (0,0) rectangle (3.0,2.4);
      \draw (1.5,0)--(1.5,2.4);
      \draw (0,1.2)--(3.0,1.2);
      \node at (0.75,1.8) {$A\cap B$};
      \node at (2.25,1.8) {$A\cap\bar B$};
      \node at (0.75,0.6) {$\bar A\cap B$};
      \node at (2.25,0.6) {$\bar A\cap\bar B$};
      \node[above] at (0.75,2.4) {$B$};
      \node[above] at (2.25,2.4) {$\bar B$};
      \node[left] at (0,1.8) {$A$};
      \node[left] at (0,0.6) {$\bar A$};
    \end{scope}
  \end{tikzpicture}
  \caption{集合运算的卡诺图式表示：左图文氏图中 $A\cup B$ 是「或」、$A\cap B$ 是「且」；右图 $2\times2$ 网格穷举两个事件的四种组合，对应概率论中样本空间的完备互斥划分}
  \label{fig:ch07-venn}
\end{figure}

\begin{commonerror}{易错诊断室：可重复与不可重复、有序与无序}
\paragraph{失误一：有序、无序混淆。}
从 5 人中选 3 人「组成代表团」用 $C_5^3=10$；若「分别担任三个不同职务」就要用 $A_5^3=60$。
题干中出现「分配、担任、排成一列、编号」等词，结果自带顺序，必须按有序处理。
\paragraph{失误二：平均分组忘记除序。}
把 6 人平均分成 3 组（每组 2 人）：若直接写 $C_6^2C_4^2C_2^2=90$，等于把三个无区别的组又按顺序数了一遍，
正确的组数为
\[
  \frac{C_6^2C_4^2C_2^2}{3!}=\frac{90}{6}=15 .
\]
\key{分组若无编号（与「分到三个不同地方」区别开），组间要除以组数的阶乘。}
\paragraph{失误三：可重复与不可重复混淆。}
3 个人从 5 种口味冰淇淋中各选一种：口味可被多人重复选择时，由分步乘法原理有 $5^3=125$ 种；
若规定三人口味互不相同，则相当于从 5 种中取 3 种的排列 $A_5^3=60$。
一字之差，原理全变——务必先确认「选过的还能不能再选」。
\end{commonerror}

\begin{conccard}[一级]{计数原理的集合本质与三技法}{一切计数问题，尤其是含「相邻／不相邻／平均分组」的排列组合题}{有序无序混淆；平均分组忘除序；可重复当成不可重复；隔板法忽略「每份至少一个」}
分类加法＝互斥子集并集基数：$|A\sqcup B|=|A|+|B|$；
分步乘法＝笛卡尔积基数：$|A\times B|=|A|\cdot|B|$；
排列＝有序单射序列 $A_n^k=\frac{n!}{(n-k)!}$，组合＝无序子集族 $C_n^k=\frac{A_n^k}{k!}$；
捆绑：$(n-m+1)!\,A_m^m$；插空：$A_n^nA_{n+1}^{\,m}$；
隔板：$x_1+\cdots+x_m=n$ 的正整数解数 $=C_{n-1}^{\,m-1}$。
\end{conccard}

\subsection{概率测度：归一化、空间缩放与完备划分}

\key{探究二：概率是定义在样本空间上的「归一化集合测度」}

\key{（1）古典概型与几何概型的统一。}
任何随机试验的样本空间 $\Omega$ 都是一个集合。在其上定义一个「大小」函数 $\mu$（测度）：
古典概型中 $\mu$ 是计数测度（$\mu(A)=|A|$），几何概型中 $\mu$ 是长度／面积／体积。
只要每个基本结果等可能，就有统一的归一化公式
\[
  P(A)=\frac{\mu(A)}{\mu(\Omega)}.
\]
由此可直接验证概率的三条公理：非负性 $P(A)\ge0$、规范性 $P(\Omega)=1$、互斥可加性 $P(A\sqcup B)=P(A)+P(B)$——它们不过是测度的归一化性质。

\begin{example}
在区间 $[0,1]$ 上任取两点 $X,Y$，求 $|X-Y|\le\dfrac12$ 的概率。
\end{example}
\begin{solution}
样本空间是单位正方形 $\Omega=[0,1]^2$，$\mu(\Omega)=1$。
满足 $|X-Y|>\dfrac12$ 的区域是直线 $Y=X\pm\dfrac12$ 之外的两个三角形，每个面积为 $\dfrac12\left(\dfrac12\right)^2=\dfrac18$。
故
\[
  P(|X-Y|\le\tfrac12)=1-2\times\tfrac18=\tfrac34 .
\]
古典（数点）与几何（量面积）在此只是测度 $\mu$ 的不同，算法同为「除以全集」。
\end{solution}

\key{（2）条件概率 $P(B\mid A)$：样本空间的缩放。}
事件 $A$ 发生后，样本全集由 $\Omega$ 缩小为子集 $A$，每个结果仍在 $A$ 内等可能。
新样本空间上的归一化给出
\[
  P(B\mid A)=\frac{P(A\cap B)}{P(A)}\qquad(P(A)>0).
\]
几何读法：\key{把 $A$ 当作新的「全宇宙」，概率质量重新按 $P(A)$ 缩放}。它仍满足三条公理（规范性 $P(\Omega\mid A)=1$、可加性），因此是以 $A$ 为底的合法概率。

\begin{example}
10 张考签中有 1 张中奖，5 人依次不放回抽取。证明：无论先抽后抽，每人中奖概率都等于 $\dfrac1{10}$（抽签公平性）。
\end{example}
\begin{solution}
第 1 人显然为 $\frac1{10}$。
第 2 人中奖 $=$「第 1 人未中且第 2 人中」，由乘法公式与缩小的样本空间：
\[
  P=P(\text{第 1 人未中})\cdot P(\text{第 2 人中}\mid\text{第 1 人未中})
  =\frac{9}{10}\times\frac{1}{9}=\frac{1}{10}.
\]
第 3 人同理：$\frac{9}{10}\times\frac{8}{9}\times\frac{1}{8}=\frac1{10}$。
第 4、5 人同理（依次缩小样本空间后概率仍为 $\frac1{10}$）。
可见「先抽优势」是错觉——条件概率把空间缩小后自动抹平了顺序差异。
\end{solution}

\key{（3）全概率公式与贝叶斯公式：完备互斥划分。}
若 $\Omega=A_1\sqcup A_2\sqcup\cdots\sqcup A_n$（$A_i$ 两两不交、并为全集，称为\key{完备互斥划分}），则
\[
  P(B)=\sum_{i=1}^{n}P(A_i)\,P(B\mid A_i)\quad(\text{全概率公式}),
\]
它是面积可加性 $\mu(B)=\sum\mu(A_i\cap B)$ 的直接翻译：把 $B$ 沿划分切成互不重叠的碎片再求和。
反过来，已知 $B$ 发生了，反推它是哪个碎片贡献的：
\[
  P(A_j\mid B)=\frac{P(A_j)\,P(B\mid A_j)}{\sum\limits_{i=1}^{n}P(A_i)\,P(B\mid A_i)}\quad(\text{贝叶斯公式}).
\]
全概率是「由因及果」，贝叶斯是「由果溯因」：$P(A_j)$ 是先验概率，$P(A_j\mid B)$ 是纳入证据 $B$ 后的后验概率。

\begin{example}
某厂三条流水线产量占比为 $50\%,30\%,20\%$，次品率分别为 $2\%,3\%,5\%$。任取一件产品，（1）求是次品的概率；（2）若它是次品，求它来自第一条流水线的概率。
\end{example}
\begin{solution}
记 $A_i$ 为「产品来自第 $i$ 条流水线」，$B$ 为「取到次品」。
（1）由全概率公式
\[
  P(B)=0.5\times0.02+0.3\times0.03+0.2\times0.05=0.010+0.009+0.010=0.029 .
\]
（2）由贝叶斯公式
\[
  P(A_1\mid B)=\frac{0.5\times0.02}{0.029}=\frac{0.010}{0.029}=\frac{10}{29}\approx0.345 .
\]
\end{solution}

\begin{figure}[htbp]
  \centering
  \begin{tikzpicture}[>=Stealth, semithick, font=\footnotesize]
    % ---------- 左：完备划分的面积分解 ----------
    \draw[thick] (0,0) rectangle (8,3);
    \draw (4,0)--(4,3);
    \draw (6.4,0)--(6.4,3);
    \fill[main!30] (0,2.4) rectangle (4,3);
    \fill[main!30] (4,2.1) rectangle (6.4,3);
    \fill[main!30] (6.4,1.5) rectangle (8,3);
    \draw[thick] (0,0) rectangle (8,3);
    \draw (4,0)--(4,3);
    \draw (6.4,0)--(6.4,3);
    \node[main!80!black] at (2,2.7) {$0.010$};
    \node[main!80!black] at (5.2,2.55) {$0.009$};
    \node[main!80!black] at (7.2,2.25) {$0.010$};
    \node[align=center] at (2,-0.42) {$A_1$\\ $P(A_1)=0.5$};
    \node[align=center] at (5.2,-0.42) {$A_2$\\ $P(A_2)=0.3$};
    \node[align=center] at (7.2,-0.42) {$A_3$\\ $P(A_3)=0.2$};
    \node[gray!70!black] at (-0.35,3.3) {$\Omega$};
    \node[align=center] at (4,-1.15) {$P(B)=0.010+0.009+0.010=0.029$\\(列宽 $\propto P(A_i)$，块高 $\propto P(B\mid A_i)$ 统一放大 $30$ 倍)\\(暗色块面积 $\propto P(A_i\cap B)$，数字为 $P(A_i\cap B)$)};
    % ---------- 右：概率树 ----------
    \begin{scope}[xshift=9.2cm]
      \node[circle,fill=black,inner sep=1.6pt] (r) at (0,1.85) {};
      \node[above left] at (r) {$\Omega$};
      \node[circle,fill=black,inner sep=1.6pt] (a1) at (1.8,3.25) {};
      \node[circle,fill=black,inner sep=1.6pt] (a2) at (1.8,1.85) {};
      \node[circle,fill=black,inner sep=1.6pt] (a3) at (1.8,0.45) {};
      \node[right] at (a1) {$A_1$};
      \node[right] at (a2) {$A_2$};
      \node[right] at (a3) {$A_3$};
      \node[circle,fill=main,inner sep=1.6pt] (b1) at (4.2,3.75) {};
      \node[circle,fill=gray,inner sep=1.6pt] (nb1) at (4.2,2.75) {};
      \node[circle,fill=main,inner sep=1.6pt] (b2) at (4.2,2.35) {};
      \node[circle,fill=gray,inner sep=1.6pt] (nb2) at (4.2,1.35) {};
      \node[circle,fill=main,inner sep=1.6pt] (b3) at (4.2,0.95) {};
      \node[circle,fill=gray,inner sep=1.6pt] (nb3) at (4.2,-0.05) {};
      \draw (r)--(a1) (r)--(a2) (r)--(a3);
      \draw[main] (a1)--(b1); \draw[gray] (a1)--(nb1);
      \draw[main] (a2)--(b2); \draw[gray] (a2)--(nb2);
      \draw[main] (a3)--(b3); \draw[gray] (a3)--(nb3);
      \node at (0.55,2.75) {$0.5$};
      \node at (0.9,2.15) {$0.3$};
      \node at (0.55,1.0) {$0.2$};
      \node[main,left] at (2.95,3.65) {$0.02$}; \node[gray,left] at (2.95,2.9) {$0.98$};
      \node[main,left] at (2.95,2.25) {$0.03$}; \node[gray,left] at (2.95,1.5) {$0.97$};
      \node[main,left] at (2.95,0.85) {$0.05$}; \node[gray,left] at (2.95,0.1) {$0.95$};
      \node[main,right=2mm] at (b1) {$B$: $0.010$};
      \node[gray,right=2mm] at (nb1) {$B^c$: $0.490$};
      \node[main,right=2mm] at (b2) {$B$: $0.009$};
      \node[gray,right=2mm] at (nb2) {$B^c$: $0.291$};
      \node[main,right=2mm] at (b3) {$B$: $0.010$};
      \node[gray,right=2mm] at (nb3) {$B^c$: $0.190$};
      \node[align=center] at (2.1,-0.8) {$P(B)=0.010+0.009+0.010=0.029$\\[2pt]$P(A_1\mid B)=\frac{0.010}{0.029}=\frac{10}{29}\approx0.345$};
    \end{scope}
  \end{tikzpicture}
  \caption{样本空间的完备互斥划分（全概率·贝叶斯）：左图中 $\Omega=A_1\sqcup A_2\sqcup A_3$，列宽正比于 $P(A_i)$、每列内暗色块高度正比于 $P(B\mid A_i)$（统一放大 $30$ 倍的示意图），故暗色块面积正比于 $P(A_i\cap B)$；右图沿树枝把每个 $A_i$ 再按条件概率分叉，路径概率求和得 $P(B)$，贝叶斯公式则由结果 $B$ 沿树反向回溯到原因 $A_1$}
  \label{fig:ch07-partition}
\end{figure}

\begin{conccard}[一级]{归一化测度：条件概率、全概率与贝叶斯}{几何概型、条件概率、由果溯因类应用题、医学筛查类情境}{把 $P(B\mid A)$ 与 $P(A\mid B)$ 混为一谈；划分不满足「不重不漏」就套全概率；只记住公式不理解「空间缩放」的几何意义}
$P(A)=\mu(A)/\mu(\Omega)$（古典＝计数测度，几何＝长度／面积／体积）；
条件概率＝空间缩放：$P(B\mid A)=P(A\cap B)/P(A)$；
全概率：$P(B)=\sum_i P(A_i)P(B\mid A_i)$；
贝叶斯（由果溯因）：$P(A_j\mid B)=\dfrac{P(A_j)P(B\mid A_j)}{\sum_i P(A_i)P(B\mid A_i)}$。
\end{conccard}

\subsection{统计线性逼近：最小二乘投影与统计反证}

\key{探究三：回归直线是散点云的「最小二乘投影」；独立性检验是统计反证法}

\key{（1）百分位数：位置的度量。}
把 $n$ 个数据从小到大排列，计算 $i=n\times p\%$：
\begin{itemize}
  \item 若 $i$ 不是整数，第 $p$ 百分位数是第 $\lceil i\rceil$ 项（向上取整的下一项）；
  \item 若 $i$ 是整数，第 $p$ 百分位数是第 $i$ 项与第 $i+1$ 项的平均数。
\end{itemize}

\begin{example}
求数据 $12,7,9,15,10,8,11,13,6,14$ 的第 $25$ 百分位数与第 $50$ 百分位数。
\end{example}
\begin{solution}
排序：$6,7,8,9,10,11,12,13,14,15$，共 $n=10$ 个。
\[
  i_1=10\times25\%=2.5\ (\text{非整数})\Rightarrow \text{第 }25\text{ 百分位数是第 }3\text{ 项 }=8;
\]
\[
  i_2=10\times50\%=5\ (\text{整数})\Rightarrow \text{第 }50\text{ 百分位数}=\frac{\text{第 }5\text{ 项}+\text{第 }6\text{ 项}}{2}=\frac{10+11}{2}=10.5 .
\]
\end{solution}

\key{（2）回归直线：竖直残差平方和最小的投影线。}
给定散点 $(x_1,y_1),\dots,(x_n,y_n)$，在直线族 $\hat y=bx+a$ 中挑选一条，使各点到直线的\key{竖直距离}（残差 $e_i=y_i-\hat y_i$）的平方和最小：
\[
  Q(a,b)=\sum_{i=1}^{n}(y_i-bx_i-a)^2\ \longrightarrow\ \min .
\]

\begin{derivation}
对 $a,b$ 分别求偏导并令其为零：
\[
  \frac{\partial Q}{\partial a}=-2\sum_{i=1}^n(y_i-bx_i-a)=0
  \ \Longrightarrow\ na+b\sum x_i=\sum y_i\ \Longrightarrow\ a=\bar y-b\bar x .
\]
\[
  \frac{\partial Q}{\partial b}=-2\sum_{i=1}^nx_i(y_i-bx_i-a)=0 .
\]
把 $a=\bar y-b\bar x$ 代入，利用 $\sum(x_i-\bar x)=0$，得
\[
  \sum_{i=1}^n(x_i-\bar x)(y_i-\bar y)-b\sum_{i=1}^n(x_i-\bar x)^2=0,
\]
故
\[
  \hat b=\frac{\sum\limits_{i=1}^n(x_i-\bar x)(y_i-\bar y)}{\sum\limits_{i=1}^n(x_i-\bar x)^2},
  \qquad \hat a=\bar y-\hat b\,\bar x .
\]
几何性质：$(\hat a,\hat b)$ 满足 $\hat a=\bar y-\hat b\,\bar x$，即点 $(\bar x,\bar y)$ 恒在直线上——\key{回归直线必过样本中心 $(\bar x,\bar y)$}。
这与第 6 章「以直代曲」、第 4 章「正交投影」同源：最小二乘就是把观测向量 $\vec y$ 向由 $\vec 1$ 与 $\vec x$ 张成的平面作投影，残差与平面正交。
\end{derivation}

\begin{example}
已知 5 组成对数据 $(1,2),(2,3),(3,5),(4,6),(5,8)$，求回归方程并计算样本相关系数 $r$。
\end{example}
\begin{solution}
\[
  \bar x=3,\quad \bar y=\frac{2+3+5+6+8}{5}=4.8 .
\]
\[
  \sum(x_i-\bar x)(y_i-\bar y)=(-2)(-2.8)+(-1)(-1.8)+0\times0.2+1\times1.2+2\times3.2=15,
\]
\[
  \sum(x_i-\bar x)^2=4+1+0+1+4=10,\quad
  \sum(y_i-\bar y)^2=7.84+3.24+0.04+1.44+10.24=22.8 .
\]
所以
\[
  \hat b=\frac{15}{10}=1.5,\qquad \hat a=4.8-1.5\times3=0.3,
  \qquad \hat y=1.5x+0.3 .
\]
\[
  r=\frac{15}{\sqrt{10\times22.8}}=\frac{15}{\sqrt{228}}\approx0.993 ,
\]
高度正相关（$|r|\le1$ 正是由第 6 章柯西不等式保证的）。
\end{solution}

\begin{figure}[htbp]
  \centering
  \begin{tikzpicture}[>=Stealth, semithick, font=\footnotesize]
    \draw[->,gray] (-0.4,0)--(5.9,0) node[below right] {$x$};
    \draw[->,gray] (0,-0.6)--(0,8.9) node[above left] {$y$};
    \foreach \x in {1,...,5} \draw (\x,0.06)--(\x,-0.06) node[below] {\x};
    \foreach \y in {2,4,6,8} \draw (0.06,\y)--(-0.06,\y) node[left] {\y};
    \draw[second,very thick,domain=-0.2:5.5,samples=2] plot (\x,{1.5*\x+0.3});
    \node[second] at (0.55,4.6) {$\hat y=1.5x+0.3$};
    \foreach \x/\y in {1/2,2/3,3/5,4/6,5/8}{
      \draw[second!80,dashed,very thick] (\x,\y)--(\x,{1.5*\x+0.3});
      \fill[main] (\x,\y) circle (1.7pt);
    }
    \node[second,right] at (1.04,1.9) {$e_1$};
    \node[second,right] at (2.04,3.15) {$e_2$};
    \node[second,left] at (2.96,4.98) {$e_3$};
    \node[second,right] at (4.04,6.15) {$e_4$};
    \node[second,right] at (5.04,7.9) {$e_5$};
    \draw[third,dashed,thick] (2.68,4.8)--(3.32,4.8);
    \draw[third,dashed,thick] (3,4.48)--(3,5.12);
    \fill[third] (3,4.8) circle (1.7pt);
    \node[third,right] at (3.35,4.86) {$(\bar x,\bar y)=(3,4.8)$};
    \fill[main] (2.9,8.55) circle (1.7pt); \node[right] at (2.98,8.55) {观测点};
    \draw[second,very thick] (4.1,8.55)--(4.5,8.55); \node[right] at (4.58,8.55) {回归直线};
  \end{tikzpicture}
  \caption{最小二乘回归：回归直线 $\hat y=1.5x+0.3$ 使各观测点到直线的竖直残差 $e_1,\dots,e_5$（橙色虚线）平方和最小；它必然穿过样本中心 $(\bar x,\bar y)=(3,4.8)$}
  \label{fig:ch07-leastsq}
\end{figure}

\key{（3）独立性检验：小概率事件原理的统计反证法。}
$2\times2$ 列联表
\[
  \begin{array}{c|cc|c}
    & y_1 & y_2 & \text{合计}\\\hline
    x_1 & a & b & a+b\\
    x_2 & c & d & c+d\\\hline
    \text{合计} & a+c & b+d & n
  \end{array}
\]
若两变量无关（零假设 $H_0$），应有 $ad\approx bc$。统计量
\[
  K^2=\frac{n(ad-bc)^2}{(a+b)(c+d)(a+c)(b+d)}
\]
度量实际数据偏离「无关」的程度。逻辑结构与第 0 章的反证法完全同构：先承认 $H_0$ 成立，则 $K^2$ 很大是\key{小概率事件}；它竟然发生了，就有理由拒绝 $H_0$。常用临界值：$K^2\ge3.841$ 对应犯错概率不超过 $5\%$，$K^2\ge6.635$ 对应不超过 $1\%$。

\begin{example}
调查 $200$ 人（男、女各 $100$ 人）对某项政策的支持情况：男性中 $70$ 人支持、$30$ 人反对；女性中 $50$ 人支持、$50$ 人反对。是否有 $99\%$ 以上的把握认为性别与支持态度有关？
\end{example}
\begin{solution}
\[
  K^2=\frac{200\times(70\times50-30\times50)^2}{100\times100\times120\times80}
  =\frac{200\times2000^2}{9.6\times10^7}
  =\frac{8\times10^8}{9.6\times10^7}\approx8.33 .
\]
因为 $8.33>6.635$，所以在犯错误的概率不超过 $0.01$ 的前提下，认为性别与支持态度有关——即有 $99\%$ 以上的把握。
\end{solution}

\begin{commonerror}{独立性检验的三个高频失误}
\paragraph{失误一：把结论说过头。}
正确表述是「在犯错误的概率不超过 $0.01$ 的前提下认为有关」或「有 $99\%$ 的把握认为有关」，
而不是「两变量有 $99\%$ 的关系」——$K^2$ 只衡量证据强度，不量化关联程度。
\paragraph{失误二：忽略样本量 $n$。}
公式分子含 $n$，同样比例的列联表，样本量扩大 $10$ 倍，$K^2$ 也扩大 $10$ 倍，结论可能完全不同；小样本下比例再悬殊也不宜下强结论。
\paragraph{失误三：把检验结论当成因果关系。}
「有关」不等于「导致」：独立性检验只能发现关联，因果方向须由实际背景另行论证。
\end{commonerror}

\begin{conccard}[一级]{最小二乘回归与独立性检验}{成对数据的回归计算、相关性判断、列联表分析}{把回归直线当成必过原点的比例函数；以为 $K^2$ 大就证明因果关系；计算 $\hat b$ 时忘除以 $\sum(x_i-\bar x)^2$；百分位数忽略 $i$ 是否为整数}
$\hat b=\frac{\sum(x_i-\bar x)(y_i-\bar y)}{\sum(x_i-\bar x)^2}$，$\hat a=\bar y-\hat b\bar x$，必过 $(\bar x,\bar y)$；
百分位数看 $i=n\cdot p\%$ 是否整数；
$K^2=\frac{n(ad-bc)^2}{(a+b)(c+d)(a+c)(b+d)}$，对照 $3.841$（$0.05$）、$6.635$（$0.01$）下结论。
\end{conccard}

\subsection{压轴收官：新定义情境题的四步破题法}

\key{探究四：新定义题不是考知识，是考「翻译」}

新高考压轴常以全新符号、新运算、新集合开场，陌生感是纸老虎。通用四步法：
\key{第一步，读懂符号集合}——把每个新名词逐字翻译为已有的集合语言（元素是什么？满足什么条件？运算封闭吗？）；
\key{第二步，构造小规模特例试探}——取 $n=1,2,3$ 等最小情形，逐一列举，列出数表；
\key{第三步，归纳猜想通项}——从特例中找规律、猜公式（归纳只指方向，不能当证明）；
\key{第四步，严格证明}——用双射（一一对应）、数学归纳法或反证法把猜想钉死。

\begin{method}[新定义情境题的四步破题流程]
\begin{enumerate}
  \item \key{翻译}：逐句把新定义的符号「转译」成集合、映射、计数或概率语言，画出韦恩图或数表；
  \item \key{特例}：取最小的参数（$n=1,2,3$ 或最小非平凡情形）把结果全部列出，观察结构；
  \item \key{归纳}：由特例数据猜想通项／公式／计数法则，并明确写出猜想；
  \item \key{证明}：优先寻找双射（把新对象一一对应到已知模型），否则用数学归纳法（与自然数相关）或反证法（唯一性、不存在性）。
\end{enumerate}
\end{method}

\begin{problem}
\key{压轴题一（新定义 $+$ 计数）}设 $n\ge k\ge1$。称 $\{1,2,\dots,n\}$ 的 $k$ 元子集 $A=\{a_1,a_2,\dots,a_k\}$（$a_1<a_2<\dots<a_k$）为「稀疏集」，若其中任意两个不同元素之差的绝对值不小于 $2$（即不含相邻整数）。记稀疏集的个数为 $f(n,k)$。
（1）列举 $n=5,\ k=2$ 的全部稀疏集，并求 $f(5,2)$；再求 $f(5,3)$；
（2）依（1）猜想 $f(n,k)$ 的表达式，并给出严格证明。
\end{problem}
\begin{solution}
\key{【第一步：读懂符号集合】}「稀疏集」＝$\{1,\dots,n\}$ 中不含相邻数的 $k$ 元子集，$f(n,k)$ 是这类子集的个数。
\key{【第二步：小规模特例】}（1）$n=5,k=2$：枚举得
\[
  \{1,3\},\{1,4\},\{1,5\},\{2,4\},\{2,5\},\{3,5\},
\]
共 $6$ 个，即 $f(5,2)=6=C_4^2$。
$n=5,k=3$：只能取 $\{1,3,5\}$，故 $f(5,3)=1=C_3^3$。
再试 $f(4,2)=3=C_3^2$、$f(6,2)=10=C_5^2$。
\key{【第三步：归纳猜想】}各值都等于 $C_{\,n-k+1}^{\,k}$，故猜想
\[
  f(n,k)=C_{\,n-k+1}^{\,k}.
\]
\key{【第四步：严格证明（双射法）】}设稀疏集 $A=\{a_1<\dots<a_k\}$，令
\[
  b_i=a_i-(i-1)\quad(i=1,\dots,k).
\]
由 $a_{i+1}-a_i\ge2$ 得 $b_{i+1}-b_i=(a_{i+1}-a_i)-1\ge1$，即 $b_1<b_2<\dots<b_k$；
又 $b_1=a_1\ge1$，$b_k=a_k-(k-1)\le n-k+1$，故 $\{b_1,\dots,b_k\}$ 是 $\{1,\dots,n-k+1\}$ 的任意 $k$ 元子集。
反之，任取 $1\le b_1<\dots<b_k\le n-k+1$，令 $a_i=b_i+(i-1)$，则 $a_{i+1}-a_i=b_{i+1}-b_i+1\ge2$，回到一个稀疏集。
两个方向互为逆映射，构成\key{一一对应}，故
\[
  f(n,k)=C_{\,n-k+1}^{\,k}.\ \surd
\]
\key{这一步同时解释了「为什么恰好是隔板模型」：减去的 $(i-1)$ 相当于预先在每个选中元素后面放好一个「隔离空位」，隔板数从 $n-1$ 缩为 $n-k+1$。}
\end{solution}

\begin{problem}
\key{压轴题二（新定义 $+$ 概率统计）}一批共 $6$ 件产品中有 $2$ 件次品，现逐个检查（不放回）。定义「转正等待次数」$X$ 为：首次抽到正品时已经检查过的产品件数。
（1）求 $X$ 的分布列与数学期望 $E(X)$；
（2）若改为有放回抽取，$X$ 服从什么分布？比较两种情形的期望并解释差异原因。
\end{problem}
\begin{solution}
\key{【第一步：翻译】}$X$ 是取值于 $\{1,2,3\}$ 的离散随机变量（最多查 3 件必出正品，因为只有 2 件次品）。
\key{【第二步：小特例手算】}
\[
  P(X=1)=\frac{4}{6}=\frac{2}{3};
\]
\[
  P(X=2)=P(\text{第 1 件次品})\cdot P(\text{第 2 件正品}\mid\text{第 1 件次品})
  =\frac{2}{6}\times\frac{4}{5}=\frac{4}{15};
\]
\[
  P(X=3)=P(\text{前两件均为次品})=\frac{2}{6}\times\frac{1}{5}=\frac{1}{15}.
\]
\key{【第三步：归纳校验】}$\frac23+\frac{4}{15}+\frac{1}{15}=1$，分布列为
\[
  \begin{array}{c|ccc}
    X & 1 & 2 & 3\\\hline
    P & \dfrac{2}{3} & \dfrac{4}{15} & \dfrac{1}{15}
  \end{array}
\]
\key{【第四步：期望与对比】}
\[
  E(X)=1\times\frac{2}{3}+2\times\frac{4}{15}+3\times\frac{1}{15}
  =\frac{10}{15}+\frac{8}{15}+\frac{3}{15}=\frac{21}{15}=\frac{7}{5}.
\]
（2）有放回时每次命中正品概率恒为 $\frac{4}{6}=\frac23$，$X$ 服从参数 $p=\frac23$ 的几何分布：
\[
  P(X=k)=\left(\frac{1}{3}\right)^{k-1}\cdot\frac{2}{3}\ (k\ge1),\qquad E(X)=\frac{1}{p}=\frac{3}{2}.
\]
$\frac{3}{2}>\frac{7}{5}$：不放回时次品一旦被抽出，剩下产品中正品浓度上升，等待时间自然缩短——这正是条件概率「空间缩放」在期望上的体现。
\end{solution}

\begin{conccard}[二级]{新定义情境题四步破题法}{强基、新高考压轴中「新运算、新集合、新统计量」类题目}{把归纳猜想的结论直接当证明使用；跳过特例直接硬算；翻译时漏掉隐含约束（如 $n\ge k$）}
第一步 翻译：逐句转译为集合／映射语言，画图表；
第二步 特例：$n=1,2,3$ 全部列举；
第三步 归纳：写出猜想公式；
第四步 证明：双射（计数）、数学归纳法（自然数命题）、反证法（唯一性与不存在）。
\end{conccard}

\begin{summary}
本章把「不确定性」彻底纳入第 0 章的集合框架：样本空间 $\Omega$ 是集合，概率是归一化计数测度 $P(A)=\mu(A)/\mu(\Omega)$；
计数问题的本质是并集基数（分类加法）与笛卡尔积基数（分步乘法），排列与组合分别是「有序单射」与「无序子集族」；
条件概率 $P(B\mid A)$ 是把样本全集缩小为 $A$ 后的空间缩放，全概率与贝叶斯公式源于样本空间的完备互斥划分（由因及果与由果溯因）；
线性回归直线是散点云的竖直最小二乘投影，必过样本中心 $(\bar x,\bar y)$；独立性检验则是「小概率事件原理」下的统计反证法；
压轴新定义题按「翻译符号 $\to$ 特例试探 $\to$ 归纳猜想 $\to$ 严格证明」四步破题。

\begin{center}
\footnotesize
\setlength{\tabcolsep}{4pt}
\renewcommand{\arraystretch}{1.35}
\begin{tabular}{p{2.5cm}p{4.7cm}p{4.5cm}p{1.4cm}}
\toprule
引擎部件 & 核心操作 & 把什么变成什么 & 对应探究 \\
\midrule
计数两原理 & $|A\sqcup B|=|A|+|B|$；$|A\times B|=|A|\,|B|$ & 分类（或）$\to$ 并；分步（且）$\to$ 积 & 探究一 \\
排列与组合 & $A_n^k=\frac{n!}{(n-k)!}$；$C_n^k=\frac{A_n^k}{k!}$ & 有序单射序列 $\to$ 无序 $k$ 元子集族 & 探究一 \\
三技法 & 捆绑／插空／隔板 $C_{n-1}^{m-1}$ & 相邻、不相邻、相同分堆 $\to$ 标准模型 & 探究一 \\
归一化测度 & $P(A)=\mu(A)/\mu(\Omega)$ & 古典与几何 $\to$ 同一除法 & 探究二 \\
条件概率 & $P(B\mid A)=\frac{P(A\cap B)}{P(A)}$ & 样本全集 $\Omega\to$ 子集 $A$：缩放归一 & 探究二 \\
全概率·贝叶斯 & $P(B)=\sum_iP(A_i)P(B\mid A_i)$ & 由因及果 $\leftrightarrow$ 由果溯因 & 探究二 \\
最小二乘回归 & $\min\sum(y_i-\hat bx_i-\hat a)^2$，过 $(\bar x,\bar y)$ & 散点云 $\to$ 竖直投影直线 & 探究三 \\
独立性检验 & $K^2$ 对临界值 & 关联判断 $\to$ 小概率反证 & 探究三 \\
四步破题 & 翻译 $\to$ 特例 $\to$ 归纳 $\to$ 严格证 & 陌生符号 $\to$ 熟悉结构 & 探究四 \\
\bottomrule
\end{tabular}
\end{center}

\key{随机性不是结构的对立面，而是有限集合上的测度；统计不是死记公式，而是「最小二乘投影 $+$ 小概率反证」的合体；新定义压轴的本质是把陌生符号翻译回熟悉的集合与结构。}
\end{summary}

\end{explore}

\begin{inclasstest}
\begin{enumerate}
  \item 从 6 名同学中选出 3 人，分别担任语文、数学、英语三科课代表，共有多少种不同的选法？\level{易}
  \item 把 10 个相同的小球放入 3 个不同的盒子，每个盒子至少放 1 个，共有多少种放法？\level{易}
  \item 口袋中装有 5 个红球、3 个白球，不放回地取 3 个球。已知取出的第一个球是红球，求取出的第三个球也是红球的概率。\level{中}
  \item 已知某组成对数据的样本中心为 $(\bar x,\bar y)=(3,5)$，且 $\sum(x_i-\bar x)(y_i-\bar y)=12$，$\sum(x_i-\bar x)^2=4$。求回归直线方程，并指出它必过的定点。\level{难}
\end{enumerate}
\end{inclasstest}

\begin{afterclass}
\subsection*{基础巩固}
\begin{enumerate}
  \item 6 人站成一排，甲、乙、丙三人必须站在一起，有多少种站法？
  \item 袋中有 3 个红球、2 个白球，不放回地依次取出 2 个球，求第二次取到红球的概率。
\end{enumerate}

\subsection*{进阶提升}
\begin{enumerate}
  \item 甲、乙两人射击的命中率分别为 $0.8$ 与 $0.7$，各射击一次。已知恰有一人命中，求命中者是甲的概率。
  \item 将 5 个不同的小球放入 3 个不同的盒子，每个盒子至少放 1 个，有多少种放法？
\end{enumerate}

\subsection*{拓展探究}
\begin{enumerate}
  \item （强基／竞赛）设 $n\ge k\ge1$。称 $\{1,2,\dots,n\}$ 的 $k$ 元子集 $A$ 为「强稀疏集」，若其中任意两个不同元素之差至少为 $3$。记其个数为 $g(n,k)$。
  （1）列出 $n=6,k=2$ 的全部强稀疏集，求 $g(6,2)$；
  （2）猜想 $g(n,k)$ 的表达式并用双射法证明；
  （3）比较 $g(n,k)$ 与本章 $f(n,k)$ 的大小，说明理由。
\end{enumerate}

\subsection*{教师版答案}
\begin{solution}
\textbf{基础 1：}把甲、乙、丙三人捆成一个大元素，与其余 3 人共 4 个元素全排列，三人内部再全排列：
\[
  A_4^4\cdot A_3^3=24\times6=144 .
\]
\end{solution}
\begin{solution}
\textbf{基础 2：}方法一（全概率）：
\[
  P=\frac35\times\frac24+\frac25\times\frac34=\frac{6}{20}+\frac{6}{20}=\frac{3}{5}.
\]
方法二（抽签公平性）：先后手对称，第二次取到红球的概率等于第一次取到红球的概率 $\frac35$。
\end{solution}
\begin{solution}
\textbf{进阶 1：}记 $A$ 为「甲命中」，$B$ 为「乙命中」，$A,B$ 独立。
\[
  P(\text{恰有一人命中})=P(A\bar B)+P(\bar A B)=0.8\times0.3+0.2\times0.7=0.38 .
\]
由条件概率（把样本空间缩小为「恰有一人命中」）：
\[
  P(A\mid \text{恰有一人命中})=\frac{0.24}{0.38}=\frac{12}{19}\approx0.632 .
\]
\end{solution}
\begin{solution}
\textbf{进阶 2：}由容斥原理（全集减去有空盒的放法）：
\[
  3^5-C_3^1\cdot2^5+C_3^2\cdot1^5=243-96+3=150 .
\]
\end{solution}
\begin{solution}
\textbf{拓展：}（1）$n=6,k=2$：枚举得
\[
  \{1,4\},\{1,5\},\{1,6\},\{2,5\},\{2,6\},\{3,6\},
\]
共 $6$ 个，即 $g(6,2)=6$。
（2）猜想：$g(n,k)=C_{\,n-2k+2}^{\,k}$（约定 $m<k$ 时 $C_m^k=0$，对应不可能情形）。
双射证明：设强稀疏集 $A=\{a_1<\dots<a_k\}$，$a_{i+1}-a_i\ge3$。令
\[
  b_i=a_i-2(i-1)\quad(i=1,\dots,k),
\]
则 $b_{i+1}-b_i=(a_{i+1}-a_i)-2\ge1$，即 $b_1<\dots<b_k$；且 $b_1\ge1$，$b_k=a_k-2(k-1)\le n-2k+2$，
故 $\{b_1,\dots,b_k\}$ 是 $\{1,\dots,n-2k+2\}$ 的任意 $k$ 元子集；反之 $a_i=b_i+2(i-1)$ 给出逆映射。一一对应，故 $g(n,k)=C_{\,n-2k+2}^{\,k}$。
（3）「强稀疏集」一定是「稀疏集」，反之不然（如 $\{1,3\}$ 稀疏但不强稀疏），前者集合族是后者的真子族，故除 $k=1$ 等平凡情形外恒有
\[
  g(n,k)<f(n,k).
\]
\end{solution}

\reading{拓展阅读 13}{概率：从计数测度到样本空间}
\end{afterclass}

% ==========================================================
% 第 8 章 齐次化：把二元问题压成一元比例
% 说明：本文件为纯正文片段，由 main_body.tex 挂载；
%       使用 versions.tex 的 \tk / \vchoice / advancedonly /
%       basiconly / solution / 四个盒子。
% ==========================================================

\section{齐次化：把二元问题压成一元比例（高观点结构版）}

\begin{methodbox}{多框架对读}
\small
\begin{center}
\begin{tabular}{c|p{2.9cm}|p{5.0cm}|p{3.1cm}}
\toprule
视角 & 本章看点 & 别的视角怎么说 & 哪侧最省力 \\
\midrule
代数 & 同除配凑、$1$ 的代换、换元 & 符号的恒等变形 & 硬算最值 \\
几何 & 比值 $t=\frac yx$＝射线斜率 & 点 $(x,y)$ 落在一条过原点射线上 & 看比例、判范围 \\
集合 & 齐次式按次数分层 & 同次＝缩放不变的性质 & 界次数与符号 \\
映射 & 齐次坐标 $\to$ 射影直线 & $(u,v)\sim\lambda(u,v)$，$x=u/v$ & 升到第 9 章接口 \\
\bottomrule
\end{tabular}
\end{center}
卡的使用：二元分式最值先数次数、再定「$1$ 的代换还是 $t=\frac yx$」，最后回代数收尾（总图见第 0 章图~\ref{fig:four-view-map}）。
\end{methodbox}

\begin{warmupbox}{读法（先看这里）}
齐次化不是「大招」，而是一句结构主义的第一性原理：\textbf{把
$n$ 个自由变量的绝对尺度约掉，只留下 $n-1$ 个比例。}
遇到二元（$x,y$）的对称式，先问一句：它们能不能合并成一个
分数 $t=\dfrac{y}{x}$？能，就降维；不能，就先配平次数。

（本章与第 5、6 章互为表里：第 6 章「条件最值四大标准动作」给出
「1 的代换」这一入口，第 5 章用矩阵特征方向处理分式递推；本章把
两者背后的原理——齐次化——上升为第一性原理并展开成完整方法链。）
\end{warmupbox}

\begin{corebox}{第一性原理}
设 $P(x,y)=0$ 是一条曲线方程，$Q(x,y)$ 是待求的目标式。若存在
正数 $\lambda$ 使得 $Q(\lambda x,\lambda y)=\lambda^{d}\,Q(x,y)$，
则 $Q$ 是 $d$ 次齐次的，它的取值只依赖于比例 $t=y/x$，与点的
绝对位置无关。非齐次的问题之所以难，是因为混进了不同次数的
项——\textbf{齐次化就是把这些项强行配到同一个次数上}。
\end{corebox}

\begin{methodbox}{次数扫描雷达（动手前先扫一遍）}
给每一项标一个「次数」，像雷达一样先扫描整个式子：
\begin{center}
\begin{tikzpicture}[scale=0.9, font=\small]
\draw[->] (-0.4,0)--(9.6,0) node[right] {$x$ 轴：项};
\draw[->] (0,-0.5)--(0,3.4);
\foreach \x/\h/\lab/\deg [count=\i] in {
  1.4/0.5/常数项/0,
  3.2/1.0/一次项/1,
  5.0/1.5/二次项/2,
  6.8/2.0/分式项/-1,
  8.6/2.5/根式项/{1/2}} {
  \draw[main, thick] (\x,0)--(\x,\h);
  \node[above] at (\x,\h) {\lab};
  \node[below] at (\x,-0.18) {$\deg$ 次};
}
\node[below] at (5.0,-0.66) {（柱高仅为扫描顺序，次数见下方标签）};
\end{tikzpicture}
\end{center}
（常数 $0$ 次、$x$ 与 $y$ 各 $1$ 次、$x^2/xy/y^2$ 为 $2$ 次、
$\dfrac1x$ 为 $-1$ 次、$\sqrt{x}$ 为 $\tfrac12$ 次。）
\textbf{判据}：一个式子是齐次的，当且仅当所有项次数相同。
\end{methodbox}

\begin{methodbox}{方法归纳：同除主元最高次幂（降维的钥匙）}\label{met:homo-divide}
若已经确认各项同次，就同除「主元」的最高次幂，把 $n$ 个变量
压成比例。例如二元齐次式同除 $x^2$（同除 $y^2$ 亦可，此时
比例换成 $x/y=1/t$），只剩
$t=\dfrac{y}{x}$：
\[
\frac{x^2+2xy+y^2}{x^2}=1+2t+t^2 .
\]
一元化之后，剩下的就是单变量函数的最值问题——\textbf{这就是
齐次化的全部收益：把二元压成一元}。
\end{methodbox}

\subsection{第 1 节 常规通法与齐次化的第一次对照}

\begin{example}[已知 $x>0,\ y>0,\ 2x+y=1$，求 $xy$ 的最大值]
\vspace{-1mm}
\begin{solution}
\textbf{考纲通法}：$y=1-2x$，故 $xy=x(1-2x)=-2x^2+x$，
开口向下，顶点在 $x=\dfrac14$，最大值为 $\tk{\dfrac18}$。

\textbf{齐次化视角}：目标与约束同时齐次化，把分母补成
\[
\frac{xy}{(2x+y)^2}=\frac{xy}{1}.
\]
令 $t=\dfrac{y}{x}>0$，则上式右端化为
\[
\frac{t}{(2+t)^2}.
\]
由 AM–GM（$(t+2)^2\ge 8t$），有 $\dfrac{t}{(2+t)^2}\le\dfrac18$，
等号在 $t=2$ 即 $y=2x$、$x=\dfrac14,\ y=\dfrac12$ 时成立。
两种算法给出同一个答案 $\tk{\dfrac18}$。
\end{solution}
\end{example}

\begin{conccard}{齐次化入口：把约束等式当 1 乘进目标式}{二元对称式的最值、范围题}{只会代入消元，不会「补次数」}
约束条件「$2x+y=1$」提供了 \textbf{一个} 齐次化入口：任何
目标式都可以乘上 $(2x+y)^k$ 而不改变取值，从而把混次的式子
全部抬到同一次数。\vchoice{乘方指数 $k$ 的选取原则是「让所有项次数相等」。}{看到等式条件，就想想能不能拿它当 1 来乘。}
\end{conccard}

\subsection{第 2 节 五大配平手术刀}

\begin{methodbox}{手术刀之一「1 的代换」升次}
条件等式本身等于 $1$（或可化成 $1$），把它乘进目标式，
低次项立刻被抬成齐次式。
\[
\left(\frac1x+\frac1y\right)\cdot 1
=\left(\frac1x+\frac1y\right)(x+y)
=2+\frac{y}{x}+\frac{x}{y}\ge 4 .
\]
\end{methodbox}

\begin{methodbox}{手术刀之二 平移原点，把方程归一}
把原点移到图形的中心（或对称中心），一次项/常数项会被
消掉，方程立刻变成标准型。
\[
x^2+y^2-2x=0\ \xrightarrow{\ X=x-1,\ Y=y\ }\ X^2+Y^2=1,
\]
于是 $x+y=X+Y+1\in[\tk{1-\sqrt2},\ \tk{1+\sqrt2}]$。
\end{methodbox}

\begin{methodbox}{手术刀之三 平方升次，剥掉根号}
$\sqrt{x}$ 的次数是 $\tfrac12$，平方一次就升到整次。
已知 $\sqrt{x}+\sqrt{y}=1$，则
\[
x+y=(\sqrt{x}+\sqrt{y})^2-2\sqrt{xy}=1-2\sqrt{xy},
\]
再用 $\sqrt{x}\cdot\sqrt{y}\le\left(\dfrac{\sqrt{x}+\sqrt{y}}{2}\right)^2=\dfrac14$
得 $x+y\ge\tk{\dfrac12}$。
\end{methodbox}

\begin{methodbox}{手术刀之四 分解与组合降次}
高次式先分解，往往能拆出可以约分的因式。
\[
\frac{1}{x(x+1)}=\frac1x-\frac1{x+1},\qquad
x^3+y^3=(x+y)(x^2-xy+y^2).
\]
\end{methodbox}

\begin{methodbox}{手术刀之五 旋转消交叉项}
二次曲线
\[
Ax^2+Bxy+Cy^2+\cdots=0
\]
的交叉项 $Bxy$ 用平移消不掉，换一招：把坐标轴旋转角度 $\theta$，
\[
x=x'\cos\theta-y'\sin\theta,\qquad y=x'\sin\theta+y'\cos\theta .
\]
代回后令 $x'y'$ 的系数为 $0$，即得消交叉项的条件
\[
\cot 2\theta=\frac{A-C}{B}
\qquad\text{（等价地 }\tan 2\theta=\frac{B}{A-C}\text{；当 }A=C\text{ 时用前者）}.
\]
特例 $xy=1$：这里 $A=C=0,\ B=1$，故 $\cot 2\theta=0$，取 $\theta=45^\circ$，
\[
xy=1\ \xrightarrow{\ \theta=45^\circ\ }\ \frac{x'^2}{2}-\frac{y'^2}{2}=1,
\]
等轴双曲线还原为标准型（$a=b=\sqrt2$）。
\textbf{完整推导、图与数列中的复特征值应用见第 5 章。}
\end{methodbox}

\begin{center}
\begin{tikzpicture}[scale=0.85, font=\small]
% 椭圆平移齐次化示意：从原点出发的斜率射线
\draw[gray!60, thick] (0,0) ellipse (3.2 and 2.0);
\draw[->, main, very thick] (0,0)--(3.6,2.4);
\draw[dashed, third] (0,0)--(3.6,-1.2);
\node[right] at (3.65,2.4) {$k=\dfrac{y}{x}>0$};
\node[right] at (3.65,-1.2) {$k<0$};
\fill[second] (0,0) circle (2pt) node[below left] {新原点};
\node at (0,-2.5) {$k=\dfrac{y}{x}$ 的符号决定直线落在哪个象限};
\end{tikzpicture}
\end{center}

\begin{conccard}{五把配平手术刀：升次、平移、平方、分解、旋转}{混次式、根式、分式迭代、含交叉项的二次曲线}{只背大招名称，不动手配平次数}
五把手术刀的共同目标是把方程「摆正」：前四把在\textbf{次数}上
配平，第五把在\textbf{坐标方向}上配平，消掉交叉项 $Bxy$。
判断该用哪一把，只需回看次数扫描的结果——缺几次就补几次，
多几次就降几次；遇到交叉项 $Bxy$ 就转坐标轴（第 5 章有完整推导），
直到方程露出标准型。
\vchoice{严格地说，配平后问题就化为关于 $t=y/x$ 的一元函数，导数或基本不等式都可直接上手。}{口诀：次数不齐先配平，同除主元变一元；方向不正先旋转。}
\end{conccard}

\begin{advancedonly}
\subsection{第 3 节 射影桥梁：为什么齐次化天然有效}

第 5 章探究三已经证明：分式迭代（莫比乌斯变换）
\[
x_{n+1}=\frac{a\,x_n+b}{c\,x_n+d}\qquad(ad-bc\neq0)
\]
的不动点恰是矩阵 $\begin{pmatrix}a&b\\c&d\end{pmatrix}$ 的特征方向，
特征值之比就是双不动点构造的公比（推导与图见第 5 章
图~\ref{fig:eigen-fixed}；不动点的一般理论与双不动点的几何图见第 9 章
图~\ref{fig:mobius-two-fixed}）。这里只补上齐次化的动机：

\textbf{齐次式天生属于射影。} 把 $(u,v)$ 与 $\lambda(u,v)$ 视为同一个
射影点，比值 $x=u/v$ 就是它的坐标；分子、分母同次的式子在这条射线
上取确定值。于是「找不动点」在射影语言里就是「求特征方向」——
齐次化把非线性的分式迭代重新线性化了。

\begin{example}[设 $x_{n+1}=\dfrac{2x_n+1}{x_n+2}$，$x_0=1$，求 $\displaystyle\lim_{n\to\infty}x_n$]
\vspace{-1mm}
\begin{solution}
矩阵 $M=\begin{pmatrix}2&1\\1&2\end{pmatrix}$ 的特征值为
$\lambda_1=3,\lambda_2=1$，特征向量分别沿
$\binom{1}{1}$ 与 $\binom{1}{-1}$，对应不动点 $x=1$ 与 $x=-1$。
$x_0=1$ 恰在 $x=1$ 处，故极限为 $\tk{1}$。
\end{solution}
\end{example}
\end{advancedonly}

\begin{basiconly}
\subsection{第 3 节 大白话版：比例不变，大小随意}

把 $(x,y)$ 想成一条过原点的射线上的点。把 $x,y$ 同时放大
$k$ 倍，点在射线上滑动，但「谁是谁的几倍」没变。齐次式
只认这条射线，不认点在射线上的位置——所以它可以只用一个
字母 $t$ 说话。

\textbf{照样子填空}：$\dfrac{xy}{(2x+y)^2}$ 的分子分母同除
$x^2$ 后得到 $\dfrac{\underline{\quad t\quad}}{(2+\underline{\quad t\quad})^2}$。
\end{basiconly}

\subsection{第 4 节 跨章实战：五题对照（含三道双解）}

\begin{problem}[闪电战 1] 已知 $x>0,\ y>0,\ x+2y=1$，求 $x^2+y^2$ 的最小值。\end{problem}
\begin{solution}
齐次化：$(x^2+y^2)(x+2y)^2=x^2+y^2$。令 $t=\dfrac{y}{x}$，
左边化为 $\dfrac{1+t^2}{(1+2t)^2}$。由柯西不等式
$(1+2t)^2\le(1^2+2^2)(1+t^2)$，得
$\dfrac{1+t^2}{(1+2t)^2}\ge\dfrac15$，等号在 $t=2$
即 $y=2x$、$x=\dfrac15,\ y=\dfrac25$ 处取得。最小值为 $\tk{\dfrac15}$。
\end{solution}

\begin{problem}[闪电战 2] 已知 $x^2+y^2-2x-4y=0$，求 $x+y$ 的取值范围。\end{problem}
\begin{solution}
配方：$(x-1)^2+(y-2)^2=5$，令
$X=x-1,\ Y=y-2$，则 $X^2+Y^2=5$，
$x+y=X+Y+3\in[3-\sqrt{10},\ 3+\sqrt{10}]$。
\end{solution}

\begin{problem}[双解对照 1] 已知 $x>0,\ y>0,\ \dfrac1x+\dfrac2y=1$，求 $x+y$ 的最小值。\end{problem}
\begin{solution}
\textbf{通法}：$x+y=(x+y)\left(\dfrac1x+\dfrac2y\right)
=3+\dfrac{y}{x}+\dfrac{2x}{y}\ge3+2\sqrt2$。

\textbf{齐次化}：令 $t=\dfrac{y}{x}$，约束化为
$\dfrac1x+\dfrac{2}{tx}=1\Rightarrow x=1+\dfrac2t$，
于是 $x+y=x(1+t)=\left(1+\dfrac2t\right)(1+t)
=1+t+\dfrac2t+2\ge3+2\sqrt2$。两种写法同解。
\end{solution}

\begin{problem}[双解对照 2] 曲线 $xy=1$ 与直线 $y=kx+2$ 有交点，求 $k$ 的范围。\end{problem}
\begin{solution}
代入得 $kx^2+2x-1=0$。有实根要求判别式
$\Delta=4+4k\ge0\Rightarrow k\ge-1$；$k=0$ 时方程为
$2x-1=0$，对应交点 $\left(\dfrac12,2\right)$，仍在曲线上。
故 $k\in[-1,+\infty)$，最小值在 $k=-1$ 处取到（相切）。
\end{solution}

\begin{problem}[双解对照 3] 已知 $\sqrt{x}+\sqrt{y}=2$，求 $x+y$ 的最小值与最大值。\end{problem}
\begin{solution}
平方：$x+y=4-2\sqrt{xy}$。由
$\sqrt{xy}\le\dfrac{\sqrt{x}+\sqrt{y}}{2}=1$ 得 $x+y\ge2$；
又 $x+y=(\sqrt{x}+\sqrt{y})^2-2\sqrt{xy}\le4$（当一项为 $0$ 时取等）。
故最小值 $\tk{2}$，最大值 $\tk{4}$。
\end{solution}

\begin{inclasstest}{当堂检测（4 题）}
\begin{enumerate}[label=\arabic*., itemsep=2pt]
\item 次数扫描：$\dfrac{x^2}{y}$ 是\underline{\hspace{2em}}次齐次式。
\item 已知 $x>0,\ y>0,\ x+y=1$，则 $xy$ 的最大值是\underline{\hspace{2em}}。
\item 曲线 $x^2+y^2-6x=0$ 平移后的标准方程是\underline{\hspace{2em}}。
\item 已知 $x>0,\ y>0$ 且 $\dfrac1x+\dfrac1y=1$，则 $x+y$ 的最小值是\underline{\hspace{2em}}。
\end{enumerate}
\begin{solution}
\tk{\text{一}}；\tk{\dfrac14}；\tk{(x-3)^2+y^2=9}；\tk{4}。
\end{solution}
\end{inclasstest}

\begin{homework}{分层作业}
\textbf{基础层（步步填空）}：已知 $x>0,\ y>0,\ x+3y=1$，则
$\dfrac1x+\dfrac3y=1+\underline{\quad}+\dfrac{3y}{x}+\underline{\quad}
\ge\underline{\quad}$。

\textbf{提高层（双解对照）}：已知 $x>0,\ y>0,\ 2x+3y=1$，
求 $\dfrac{1}{x}+\dfrac{1}{y}$ 的最小值（写出通法与齐次化两种）。

\begin{solution}
基础层：$\tk{\dfrac{3x}{y}}$，$\tk{9}$，$\tk{16}$。\\
提高层：$\left(\dfrac1x+\dfrac1y\right)(2x+3y)
=5+\dfrac{3y}{x}+\dfrac{2x}{y}\ge5+2\sqrt6$，最小值
$\tk{5+2\sqrt6}$。
\end{solution}
\end{homework}

\begin{commonerror}{易错点：把「同除」用反了方向}
常见的错误是把 $\dfrac{xy}{(2x+y)^2}$ 齐次化时只除一次，
错写成 $\dfrac{y/x}{2+y/x}$——\textbf{分母少除了一个 $x$}。
正确写法是分子分母\textbf{同除同一个量} $x^2$：
\[
  \dfrac{xy}{(2x+y)^2}=\dfrac{y/x}{(2+y/x)^2}.
\]
次数差才保持不变；除完应当立刻检查「分子还剩几次、分母还剩几次」。
\end{commonerror}

\begin{conccard}[二级]{什么时候必须齐次化，什么时候不必}{条件最值、曲线与直线的位置关系、根式最值}{不管题目长什么样都硬凑 $t=y/x$；或只会代入消元，从不配次数}
先看次数结构：目标式与约束能配成同次（或经平方、平移后可配平）时，
齐次化一步降为一元 $t$；消元后是一元二次、只问有无实根时，判别式
更直接；含 $\sqrt{\phantom{x}}$ 的半次项先平方升次，再决定走哪条路。
\end{conccard}

% ==========================================================
%  第 5 节 七问连环母题（解答即笔记 · Scaffolded Master Problem）
% ==========================================================

\subsection{第 5 节 七问连环母题：解答即笔记}

\begin{warmupbox}{母题怎么用（先看这里）}
本节两道大题，每道挂载七个递进小问。\vchoice{它们不是练习题，
而是「会说话的笔记」：第 (1)–(3) 问替代概念与规范书写笔记，
第 (4)–(6) 问替代大招、公式族与充要边界笔记，第 (7) 问替代
错题本。做完一道，这一章的框架就长在你的解题纸上。}{每道题顺着
七步走即可：先算具体数值，再翻译结构，然后规范推导、降维秒杀、
推广公式、逆向追问，最后排雷。}
\end{warmupbox}

\refstepcounter{hvmasterq}
\begin{masterq}{线性约束 $x+y=1$：从 $\dfrac1x+\dfrac1y$ 长出一族不等式}
设 $x>0,\ y>0,\ x+y=1$。
\begin{enumerate}[label=(\arabic*), itemsep=3pt]
\item \textbf{画图与数值感知}：计算 $(x,y)=\left(\dfrac14,\dfrac34\right)$、
$\left(\dfrac12,\dfrac12\right)$、$\left(\dfrac34,\dfrac14\right)$ 处
$\dfrac1x+\dfrac1y$ 的值，描点并补出草图。\logotag{核心公理}
\item \textbf{多视角结构翻译}：用集合语言写出点集；指出它在坐标系中的
几何形状；再用恒等式把 $\dfrac1x+\dfrac1y$ 改写成只含 $xy$ 的式子。\logotag{结构翻译}
\item \textbf{考纲规范通法}：用基本不等式求 $\dfrac1x+\dfrac1y$ 的最小值，
写出取等条件与关键步骤（用 \tk{} 填空）。\logotag{考纲通法}
\item \textbf{高观点降维秒杀}：令 $t=\dfrac{y}{x}$，把 $\dfrac1x+\dfrac1y$
写成 $t$ 的函数并用一行不等式收束，与第 (3) 问核对。\logotag{大招公式}
\item \textbf{公式变形家族生成}：推广到 $x+y=s$、$ax+by=1$，并推导
$\dfrac1{x^2}+\dfrac1{y^2}$ 的下界。\logotag{公式族}
\item \textbf{参数变式与逆向追问}：已知 $\dfrac1x+\dfrac1y=5$ 且 $x+y=1$，
反求 $x,y$；再回答：使 $\dfrac1x+\dfrac1y\ge m$ 恒成立的 $m$ 最大是多少？\logotag{充要边界}
\item \textbf{边界陷阱与错题排雷}：某同学删去「$x,y>0$」后仍写
$\dfrac1x+\dfrac1y\ge4$。请举反例推翻它，并说明原题有没有最大值。\logotag{避坑铁律}
\end{enumerate}
\end{masterq}

\begin{solution}
(1) $\dfrac{16}{3}$，$\tk{4}$，$\dfrac{16}{3}$；图象关于 $x=\dfrac12$ 对称，
两端趋于 $+\infty$。

(2) $\{(x,y)\mid x>0,\ y>0,\ x+y=1\}$；第一象限内连接 $(1,0)$ 与
$(0,1)$ 的线段；$\dfrac1x+\dfrac1y=\dfrac{x+y}{xy}=\tk{\dfrac{1}{xy}}$。

(3) $xy\le\left(\dfrac{x+y}{2}\right)^2=\dfrac14$，故
$\dfrac1x+\dfrac1y=\dfrac1{xy}\ge\tk{4}$，当且仅当 $x=y=\dfrac12$。

(4) $x=\dfrac{1}{1+t},\ y=\dfrac{t}{1+t}$，于是
$\dfrac1x+\dfrac1y=(1+t)+\dfrac{1+t}{t}=2+t+\dfrac1t\ge\tk{4}$
（$t=1$ 取等），与 (3) 一致。

(5) (a) $x+y=s$：$\dfrac1x+\dfrac1y=\dfrac{s}{xy}\ge\dfrac{s}{s^2/4}
=\tk{\dfrac4s}$；
(b) $ax+by=1$：$\dfrac1x+\dfrac1y=\left(\dfrac1x+\dfrac1y\right)(ax+by)
=a+b+\dfrac{ay}{x}+\dfrac{bx}{y}\ge a+b+2\sqrt{ab}=\tk{(\sqrt a+\sqrt b)^2}$；
(c) $\dfrac1{x^2}+\dfrac1{y^2}=\left(\dfrac1x+\dfrac1y\right)^2-\dfrac{2}{xy}
\ge16-8=\tk{8}$。

(6) $xy=\dfrac15$，故 $x,y$ 是 $z^2-z+\dfrac15=0$ 的两根，即
$\tk{\dfrac{5\pm\sqrt5}{10}}$；$m_{\max}=\tk{4}$（恰为下确界）。

(7) 反例：$x=2,\ y=-1$ 满足 $x+y=1$，但 $\dfrac1x+\dfrac1y=-\dfrac12<4$。
原题中 $x\to0^+$ 时 $\dfrac1x+\dfrac1y\to+\infty$，故只有最小值、无最大值。
\end{solution}

\refstepcounter{hvmasterq}
\begin{masterq}{二次约束 $x^2+y^2=1$：齐次化与「1 的代换」双剑合璧}
设 $x^2+y^2=1\ (x>0,\ y>0)$。
\begin{enumerate}[label=(\arabic*), itemsep=3pt]
\item \textbf{画图与数值感知}：计算 $\left(\dfrac{\sqrt2}{2},\dfrac{\sqrt2}{2}\right)$、
$\left(\dfrac35,\dfrac45\right)$、$\left(\dfrac12,\dfrac{\sqrt3}{2}\right)$ 处
$\dfrac1{x^2}+\dfrac1{y^2}$ 与 $\dfrac{x}{y}+\dfrac{y}{x}$ 的值，并在单位圆
第一象限弧上描点。\logotag{核心公理}
\item \textbf{多视角结构翻译}：用集合 $\{(\cos\theta,\sin\theta)\mid
0<\theta<\tfrac\pi2\}$ 重写点集；指出几何形状；说明 $x^2+y^2=1$ 为什么
可以当作「1」乘进任何式子。\logotag{结构翻译}
\item \textbf{考纲规范通法}：求 $\dfrac1{x^2}+\dfrac1{y^2}$ 的最小值，
写出取等条件（用 \tk{} 填空）。\logotag{考纲通法}
\item \textbf{高观点降维秒杀}：把 $\dfrac1{x^2}+\dfrac1{y^2}$ 配成分子分母
同次（乘「1」再除「1」），令 $t=\dfrac{y}{x}$ 一行出结果，并顺手得到
$\dfrac{x}{y}+\dfrac{y}{x}$ 的下界。\logotag{大招公式}
\item \textbf{公式变形家族生成}：推导 $x^2+y^2=r^2$、$x^2+2y^2=1$ 情形下
$\dfrac1{x^2}+\dfrac1{y^2}$、$\dfrac1{x^2}+\dfrac2{y^2}$ 与
$\dfrac1{x^4}+\dfrac1{y^4}$ 的下界。\logotag{公式族}
\item \textbf{参数变式与逆向追问}：已知 $\dfrac1{x^2}+\dfrac1{y^2}=5$ 且
$x^2+y^2=1$，求 $x^2,y^2$；若去掉「$x,y>0$」只保留 $x^2+y^2=1$，
最小值是否改变？\logotag{充要边界}
\item \textbf{边界陷阱与错题排雷}：错解「$xy\le\dfrac12\Rightarrow
\dfrac1{xy}\ge2\Rightarrow\dfrac1{x^2y^2}\ge4$」的哪一步依赖了符号条件？
写出防错铁律。\logotag{避坑铁律}
\end{enumerate}
\end{masterq}

\begin{solution}
(1) $\left(\dfrac{\sqrt2}{2},\dfrac{\sqrt2}{2}\right)$：$\dfrac1{x^2}+\dfrac1{y^2}=4$，
$\dfrac{x}{y}+\dfrac{y}{x}=2$；
$\left(\dfrac35,\dfrac45\right)$：$\dfrac{25}{9}+\dfrac{25}{16}
=\tk{\dfrac{625}{144}}$，$\dfrac{x}{y}+\dfrac{y}{x}=\dfrac34+\dfrac43
=\tk{\dfrac{25}{12}}$；
$\left(\dfrac12,\dfrac{\sqrt3}{2}\right)$：$4+\dfrac43=\tk{\dfrac{16}{3}}$。

(2) 几何上是单位圆在第一象限的四分之一圆弧；代数上 $x^2+y^2=1$ 是
「1 的代换」的合法来源（乘以它不改变任何式子的取值）。

(3) $\dfrac1{x^2}+\dfrac1{y^2}=\dfrac{x^2+y^2}{x^2y^2}=\dfrac1{x^2y^2}$，
而 $xy\le\dfrac{x^2+y^2}{2}=\dfrac12\Rightarrow x^2y^2\le\dfrac14$，
故原式 $\ge\tk{4}$，等号当且仅当 $x=y=\dfrac{\sqrt2}{2}$。

(4) 乘「1」：$\dfrac1{x^2}+\dfrac1{y^2}
=\dfrac{(x^2+y^2)^2}{x^2y^2}$；令 $t=\dfrac{y}{x}$ 得
$\dfrac{(1+t^2)^2}{t^2}=\left(t+\dfrac1t\right)^2\ge4$；同一行顺手得到
$\dfrac{x}{y}+\dfrac{y}{x}=t+\dfrac1t\ge\tk{2}$。

(5) (a) $x^2+y^2=r^2\Rightarrow\dfrac1{x^2}+\dfrac1{y^2}
=\dfrac{r^4}{x^2y^2}\ge\tk{\dfrac{4}{r^2}}$；
(b) $\left(\dfrac1{x^2}+\dfrac2{y^2}\right)(x^2+2y^2)
=5+2\left(\dfrac{y^2}{x^2}+\dfrac{x^2}{y^2}\right)\ge5+4=9$，
故 $\dfrac1{x^2}+\dfrac2{y^2}\ge\tk{9}$（等号 $x^2=y^2=\dfrac13$）；
(c) $\dfrac1{x^4}+\dfrac1{y^4}\ge\dfrac{2}{x^2y^2}\ge\tk{8}$。

(6) 令 $u=x^2,\ v=y^2$，则 $u+v=1,\ \dfrac{u+v}{uv}=5\Rightarrow uv=\dfrac15$，
故 $u,v=\tk{\dfrac{5\pm\sqrt5}{10}}$；去掉正数限制后结论不变（只含平方项），
但仍需 $x,y\neq0$。

(7) 「$\dfrac1{xy}\ge2$」这一步要求 $xy>0$；若 $xy<0$，不等号方向相反。
本题最终正确只是因为又平方了一次。防错铁律：取倒数或用基本不等式前，
先确认各项同号。
\end{solution}

\begin{basiconly}
\subsection{第 6 节 大白话版：两道母题的半成品脚手架}
母题一：两人分一个蛋糕，$x,y$ 是各拿多少。怎么分，倒数的和最小？
先试四等分点：$x=\dfrac14$ 时 $\dfrac1x+\dfrac1y=\dfrac{16}{3}$；
正中间 $x=\dfrac12$ 时等于 \underline{\hspace{2em}}，看起来最小。
于是猜想：\underline{\hspace{2em}} 时最小，最小值 \underline{\hspace{2em}}。

母题二：单位圆上找一点。试 $x=y=\dfrac{\sqrt2}{2}$，得
$\dfrac1{x^2}+\dfrac1{y^2}=4$；试 $x=\dfrac12,\ y=\dfrac{\sqrt3}{2}$，
得 $\dfrac{16}{3}$。故最小值是 \underline{\hspace{2em}}，
取到在 \underline{\hspace{2em}}。
\end{basiconly}

\begin{advancedonly}
\textbf{射影眼光看这两道母题}：第 (4) 问的比值 $t=y/x$ 是射影直线上的
坐标，齐次式只认射影点、不认代表元（同第 3 节的射影视角）。把母题一
的约束换成 $ax+by=1$ 后，(5)(b) 的取等条件 $\dfrac{ay}{x}=\dfrac{bx}{y}$
给出最优比例 $t^*=\sqrt{b/a}$，即最优射线的方向比为 $\sqrt b:\sqrt a$——
配平问题变成求一条射线与一条直线的交点。
\end{advancedonly}

\begin{conccard}{连环母题：做一题 = 学一节}{概念引入、性质翻译、规范书写、大招结论、公式族、充要边界、错题本七类笔记一次生成}{把母题当普通大题，只做最后一问}
\vchoice{做完这两道题，留白处就是一份《齐次化》随身笔记：第 (1)(2) 问
是概念与性质，第 (3) 问是考场书写，第 (4)(5) 问是大招与公式表，
第 (6) 问是逆向思维，第 (7) 问是错题本。}{基础较弱的同学先完成
第 (1)–(3) 问即可达标；第 (4)–(7) 问留给学有余力者。}
\end{conccard}

\begin{notezone}{章末归档：齐次化二级结论 + 七问笔记功能}
\textbf{二级结论}（方法链压缩版）：（1）齐次式的值只依赖比例 $t=y/x$；
（2）约束等式可以当 $1$ 乘进目标式完成升次（第 6 章入口、本章展开）；
（3）半次项平方可升为整次；（4）平移可以消去一次项，一般理论见
第 5 章；（5）分式迭代的不动点对应矩阵特征方向（第 5 章）。

\textbf{七问与笔记功能}：(1) 画图与数值感知 $\to$ 概念引入笔记；
(2) 三视角结构翻译 $\to$ 性质定义笔记；(3) 考纲规范通法 $\to$ 规范
书写笔记；(4) 高观点降维秒杀 $\to$ 二级结论与大招笔记；(5) 公式变形
家族 $\to$ 题型拓展笔记；(6) 参数变式与逆向追问 $\to$ 充要条件与逻辑
笔记；(7) 边界陷阱排雷 $\to$ 错题本易错点笔记。
\end{notezone}

\reading{拓展阅读 14}{齐次化：比例结构与射影桥梁}

% ==========================================================
% 第 9 章 不动点：迭代、方程与几何变换的交汇
% 说明：本文件为纯正文片段，由 main_body.tex 挂载；
%       使用 versions.tex 的 \tk / \vchoice / advancedonly /
%       basiconly / solution / 盒子与七问母题环境。
% ==========================================================

\section{不动点：迭代、方程与几何变换的交汇（高观点版）}

\begin{methodbox}{多框架对读}
\small
\begin{center}
\begin{tabular}{c|p{2.9cm}|p{5.0cm}|p{3.1cm}}
\toprule
视角 & 本章看点 & 别的视角怎么说 & 哪侧最省力 \\
\midrule
代数 & 解 $x=F(x)$、构造 $b_n=a_n-x^*$ & 把不动点当方程的根 & 求解析式 \\
几何 & 曲线与对角线的交点、变换的不动心 & 位置关系与轨迹形状 & 看收敛、判稳定 \\
集合 & 轨道 $\{a_n\}$ 与其聚点 & 值集、吸引域／排斥域 & 讨论存在与唯一 \\
映射 & 停留点 $F(p)=p$；Möbius 射影 & 换坐标系＝换表示 & 非线性 $\to$ 等比 \\
\bottomrule
\end{tabular}
\end{center}
卡的使用：四类不动点（数列／方程／变换／分式）共用同一句话——$F(p)=p$；算不下去就换表示（总图见第 0 章图~\ref{fig:four-view-map}）。
\end{methodbox}

\begin{warmupbox}{读法（先看这里）}
不动点只有一个定义：映射 $F$ 中满足 $F(p)=p$ 的点 $p$。
它有三副面孔：数列递推 $a_{n+1}=F(a_n)$（离散自迭代）、
方程 $f(x)=x$（曲线与对角线的交点）、几何变换（旋转中心、
位似中心、分式变换）。三者的共同动作是\key{换坐标系，把锚点
搬到原点}——这正是第 8 章齐次化、第 5 章平移／旋转的内核。
本章把这条线索单独展开，第 5 章与第 8 章的应用改为引用。
\end{warmupbox}

\begin{learnobjective}
\begin{itemize}
  \item 把数列看成 $\mathbb{N}\to\mathbb{R}$ 的（离散）函数，把递推看成映射的自迭代；
  \item 会用不动点方程 $x=F(x)$ 与稳定性判据 $|F'(x^*)|$ 处理一阶递推；
  \item 理解连续情形 $f(x)=x$ 的图解法与第 6 章隐零点、切线放缩的联系；
  \item 知道旋转中心、位似中心、仿射定点都是不动点的特例；
  \item 理解齐次坐标 $(u,v)$ 为什么能把分式迭代线性化（详见第 5 章）。
\end{itemize}
\end{learnobjective}

\subsection{第 1 节 数列是自然数集上的函数：递推即迭代}

数列 $a_n$ 本质是函数 $a:\mathbb{N}\to\mathbb{R}$，$n\mapsto a_n$。
一阶递推 $a_{n+1}=F(a_n)$ 就是映射 $F$ 的\key{自迭代}：
$a_1=F(a_0),\ a_2=F(F(a_0)),\dots$，数列是 $F$ 的一条轨道。
于是
\[
  \text{不动点}\ x^*=F(x^*)
  \iff \text{点 }(x^*,x^*)\text{ 同时在 }F\text{ 的图象与对角线 }y=x\text{ 上}.
\]
\textbf{稳定性判据}：在 $x^*$ 附近线性化 $F(x)\approx x^*+F'(x^*)(x-x^*)$，
故
\[
  |F'(x^*)|<1\Rightarrow\text{收缩（稳定）},\quad
  |F'(x^*)|>1\Rightarrow\text{扩张（不稳定）},\quad
  |F'(x^*)|=1\Rightarrow\text{看高阶项（临界）}.
\]

\begin{figure}[htbp]
  \centering
  \begin{minipage}{0.48\textwidth}
    \centering
    \begin{tikzpicture}[x=0.9cm,y=0.9cm,>=Stealth,semithick,font=\footnotesize]
      \draw[->,gray] (-1,0)--(3,0);
      \draw[->,gray] (0,-1)--(0,3);
      \draw[gray!60,dashed,domain=-1:3] plot(\x,\x);
      \draw[main,very thick,domain=-1:3] plot(\x,{0.5*\x+1});
      \fill[second] (2,2) circle (1.6pt) node[above right] {$x^*=2$};
      \draw[third,thick] (0,0)--(0,1)--(1,1)--(1,1.5)--(1.5,1.5)--(1.5,1.75)--(1.75,1.75);
      \node[gray!70!black] at (1.5,-0.65) {$F(x)=\tfrac12x+1$：收缩};
    \end{tikzpicture}
  \end{minipage}\hfill
  \begin{minipage}{0.48\textwidth}
    \centering
    \begin{tikzpicture}[x=0.9cm,y=0.9cm,>=Stealth,semithick,font=\footnotesize]
      \draw[->,gray] (-1,0)--(3,0);
      \draw[->,gray] (0,-1)--(0,3);
      \draw[gray!60,dashed,domain=-1:3] plot(\x,\x);
      \draw[main,very thick,domain=-0.5:2] plot(\x,{2*\x-1});
      \fill[second] (1,1) circle (1.6pt) node[below left] {$x^*=1$};
      \draw[third,thick] (1.2,1.2)--(1.2,1.4)--(1.4,1.4)--(1.4,1.8)--(1.8,1.8)--(1.8,2.6);
      \node[gray!70!black] at (1.5,-0.65) {$G(x)=2x-1$：扩张};
    \end{tikzpicture}
  \end{minipage}
  \caption{不动点的两种命运：同一张「$F$ 的图象＋对角线」图上，左图 $|F'|<1$，蛛网阶梯螺旋收拢到 $x^*$；右图 $|G'|>1$，阶梯被弹开。斜率与对角线 $y=x$ 的比较就是稳定性判据}
  \label{fig:cobweb-pair}
\end{figure}

\begin{methodbox}{不动点法三步}
\begin{enumerate}[label=(\arabic*), itemsep=2pt]
  \item \key{找}：解 $x^*=F(x^*)$（可能不止一个）；
  \item \key{判}：算 $|F'(x^*)|$ 定稳定／不稳定；
  \item \key{化}：令 $b_n=a_n-x^*$，则 $b_{n+1}\approx F'(x^*)\,b_n$；
        稳定时 $b_n\to0$，并可用等比模型写出通项。
\end{enumerate}
常见变体：分式递推用 $\dfrac{a_n-x_1}{a_n-x_2}$ 作商（第 5 章探究五），
双曲型递推用三角／指数代换。
\end{methodbox}

\subsection{第 2 节 连续情形：方程 $f(x)=x$ 与图解法}

把 $f(x)=x$ 改写成 $\varphi(x)=f(x)-x=0$：不动点就是 $\varphi$ 的零点。
几何上，它是曲线 $y=f(x)$ 与直线 $y=x$ 的交点。
\begin{example}[方程 $x=\cos x$ 在 $(0,\tfrac\pi2)$ 内有唯一解]
\vspace{-1mm}
\begin{solution}
$\varphi(x)=\cos x-x$，$\varphi(0)=1>0$，$\varphi(\tfrac\pi2)=-\tfrac\pi2<0$，
由零点定理至少有一解；又 $\varphi'(x)=-\sin x-1<0$ 严格递减，故唯一。
进一步 $|\varphi'|$ 与压缩性：$f(x)=\cos x$ 的导数 $|f'|=\sin x<1$ 在
$(0,\tfrac\pi2)$ 成立，故 $f$ 是压缩映射，迭代 $x_{n+1}=\cos x_n$
对任意初值都收敛到该唯一不动点——这就是第 6 章「隐零点＋单调性」
的几何版本。
\end{solution}
\end{example}
\textbf{与考纲的接口}：求超越方程的近似解（二分法）本质是不断把区间
压缩到不动点；证明 $f(x)\ge x$ 型不等式（如 $e^x\ge x+1$）则是在
比较曲线与对角线的位置（见第 6 章切线放缩）。

\subsection{第 3 节 几何变换的不动点：旋转中心、位似中心与仿射定点}

\begin{itemize}
  \item \key{旋转}：转角 $\theta\not\equiv0\pmod{2\pi}$ 的平面旋转只有
        \textbf{中心}一个不动点；
  \item \key{位似}：位似中心不动，其余点沿射线伸缩；
  \item \key{平移}：平面上没有不动点（它的「不动点」在无穷远处）；
  \item \key{仿射变换} $(x,y)\mapsto A\binom xy+\vec b$：不动点满足
        $(A-I)\binom xy=-\vec b$，当 $\det(A-I)\neq0$ 时唯一。
        椭圆的中心就是两方向伸缩变换的公共不动点（见第 4 章）。
\end{itemize}
直觉：任何「保结构」的变换都可以先找到它的不动点，再把坐标原点
搬过去——剩下的表达式往往立刻简化（第 8 章的平移手术刀、第 5 章
的旋转消交叉项都是这一招）。

\subsection{第 4 节 分式迭代与齐次坐标：射影视角}

一次分式变换 $x\mapsto\dfrac{ax+b}{cx+d}$ 本身不是线性的，但把 $x$
写成两个量的比 $x=u/v$ 之后，分子分母各自线性组合：
\[
  x_{n+1}=\frac{a\,x_n+b}{c\,x_n+d}
  \iff
  \binom{u_{n+1}}{v_{n+1}}
  =\begin{pmatrix}a&b\\ c&d\end{pmatrix}\binom{u_n}{v_n},
  \qquad x_n=\frac{u_n}{v_n}.
\]
关键在第二步：点 $(u,v)$ 与 $\lambda(u,v)$ 表示同一个 $x$（射影点），
所以「二维向量的线性作用」可以「投影」为一维分式作用。第 5 章
探究三已由此证明：\key{分式迭代的不动点 $=$ 矩阵的特征方向}
（推导与数值例子见第 5 章，此处不重复）。这里只看一件事：为什么
$(u,v)$ 与 $\lambda(u,v)$ 是同一点。

\begin{figure}[htbp]
  \centering
  \begin{tikzpicture}[>=Stealth,semithick,font=\footnotesize,x=1.1cm,y=1.1cm]
    \draw[->,gray] (-0.5,0)--(4.6,0) node[right] {$u$};
    \draw[->,gray] (0,-0.5)--(0,3.1) node[above] {$v$};
    \draw[main!60!black,dashed,very thick] (0,1)--(4.5,1) node[right] {$v=1$（投影线）};
    % 两条射线：同一条射线上取两个代表元
    \draw[second,very thick,->] (0,0)--(2.1,1.05);
    \draw[third,very thick,->] (0,0)--(1.05,2.1);
    \foreach \p/\lab/\pos in {(2,1)/{$x=2$}/above right,(1,2)/{$x=\tfrac12$}/above left}
      \fill[main] \p circle (1.6pt) node[\pos] {\lab};
    \fill (0,0) circle (1.2pt) node[below left] {$O$};
    \node[gray!70!black,align=center] at (2.2,-0.7) {同一条过原点的射线上的 $(u,v)$、$(2u,2v)$\\投影到 $v=1$ 上得到同一个 $x=u/v$};
  \end{tikzpicture}
  \caption{齐次坐标：用有序数对 $(u,v)$ 刻画比值 $x=u/v$。射线上的所有代表元落在投影线 $v=1$ 的同一点上，于是「一维分式迭代」升级为「二维线性迭代」}
  \label{fig:ray-projective}
\end{figure}

\textbf{双不动点的几何图像}：设 $M$ 的两个不动点为 $p,q$（射影直线上）。
取新坐标
\[
  X=\frac{x-p}{\,x-q\,},
\]
它把 $p$ 送到 $0$、$q$ 送到 $\infty$；在此坐标下 $M$ 退化为\emph{纯伸缩}
$X\mapsto\lambda X$，其中 $\lambda$ 正是转移矩阵两特征值之比（第 5 章探究五）。
于是
\[
  |\lambda|<1\Rightarrow X\to0\ \ (x\to p),\qquad
  |\lambda|>1\Rightarrow X\to\infty\ \ (x\to q).
\]
换言之：\key{两条特征方向在射影直线上就是两个不动点}；以 $p,q$ 为
$0,\infty$ 建立坐标后，分式迭代变成等比数列——这就是共轭变换的几何全貌。

\begin{figure}[htbp]
  \centering
  \begin{minipage}{0.48\textwidth}
    \centering
    \begin{tikzpicture}[>=Stealth,semithick,font=\footnotesize,x=1.2cm,y=1.2cm]
      \draw[->,gray] (-3.2,0)--(3.6,0);
      \node[gray!70!black,right] at (3.6,0) {$x$（射影直线）};
      \fill[main] (-2,0) circle (1.5pt) node[below] {$p$};
      \fill[second] (1.4,0) circle (1.5pt) node[below] {$q$};
      \draw[third,thick,->] (-1.2,0.28)--(-0.45,0.28);
      \draw[third,thick,->] (0.1,0.28)--(0.9,0.28);
      \node[gray!70!black,align=center] at (0.2,-1.0)
        {区间 $(p,q)$ 内 $X<0$，$|X|$ 越滚越大，\\$x$ 被拉向 $q$（$|\lambda|>1$）};
    \end{tikzpicture}
  \end{minipage}\hfill
  \begin{minipage}{0.48\textwidth}
    \centering
    \begin{tikzpicture}[>=Stealth,semithick,font=\footnotesize,x=1.0cm,y=1.0cm]
      \draw[->,gray] (-2.4,0)--(2.6,0) node[right] {$u$};
      \draw[->,gray] (0,-0.8)--(0,2.6) node[above] {$v$};
      \draw[main!60!black,dashed,thick] (-2.2,1)--(2.4,1) node[right] {$v=1$};
      \node[gray!70!black,left] at (-0.02,1) {$1$};
      \draw[main,very thick,->] (0,0)--(-1.05,0.55);
      \draw[main,very thick,dashed] (0,0)--(-2.1,1.05);
      \draw[second,very thick,->] (0,0)--(2.4,1.2);
      \fill[main] (-0.6,0.3) circle (1.4pt) node[below] {$(u,v)$};
      \node[main!70!black,above] at (-1.5,0.82) {$\times\lambda$};
      \fill[third] (-2,1) circle (1.4pt) node[below] {$p$};
      \fill[third] (2,1) circle (1.4pt) node[below] {$q$};
      \node[gray!70!black,align=center] at (0,-0.55)
        {两条特征直线 $\leftrightarrow$ 两个不动点};
    \end{tikzpicture}
  \end{minipage}
  \caption{双不动点的几何：左图在射影直线上，$X=(x-p)/(x-q)$ 使 $p\leftrightarrow0$、
  $q\leftrightarrow\infty$，迭代退化为 $X\mapsto\lambda X$（图示 $|\lambda|>1$，轨道被拉向 $q$）；
  右图回到 $(u,v)$ 平面，两条过原点的特征直线分别对应 $p,q$，作用即沿射线作伸缩}
  \label{fig:mobius-two-fixed}
\end{figure}

\subsection{第 5 节 七问连环母题：解答即笔记}

\refstepcounter{hvmasterq}
\begin{masterq}{不动点与收敛：$a_{n+1}=\dfrac12a_n+1$，$a_0=0$}
\begin{enumerate}[label=(\arabic*), itemsep=3pt]
\item \textbf{画图与数值感知}：算出 $a_1,a_2,a_3,a_4$，描点并猜极限。\logotag{核心公理}
\item \textbf{多视角结构翻译}：把该递推分别看成（i）$\mathbb{N}\to\mathbb{R}$
的函数迭代、（ii）平面直角坐标系中曲线 $y=\tfrac12x+1$ 与直线 $y=x$
的交点问题。\logotag{结构翻译}
\item \textbf{考纲规范通法}：用不动点法求通项，写出关键步骤
（令 $b_n=a_n-x^*$，\tk{} 填空）。\logotag{考纲通法}
\item \textbf{高观点降维秒杀}：用 $|F'(x^*)|$ 直接写出
$|a_n-x^*|=2\cdot2^{-n}$，并与第 6 章局部线性化对读。\logotag{大招公式}
\item \textbf{公式变形家族生成}：推导 $a_{n+1}=pa_n+q$ 的完整公式族
（$p\neq1$、$p=1$、$|p|>1$ 三种情形）。\logotag{公式族}
\item \textbf{参数变式与逆向追问}：由 $a_2=\tfrac32$、$a_3=\tfrac74$
反求 $p,q$；再回答：若 $\lim a_n=5$ 且 $p=\tfrac12$，求 $q$。\logotag{充要边界}
\item \textbf{边界陷阱与错题排雷}：某同学习惯「设 $\lim a_n=L$，代入得
$L=\tfrac12L+1$」，请指出缺少哪个前提，并举反例说明。\logotag{避坑铁律}
\end{enumerate}
\end{masterq}

\begin{solution}
(1) $a_1=1,\ a_2=\dfrac32,\ a_3=\dfrac74,\ a_4=\dfrac{15}8$，
趋向 $2$。

(2) （i）$F(x)=\tfrac12x+1$ 的自迭代；（ii）交点为
联立 $y=\tfrac12x+1$ 与 $y=x$ 的解 $(2,2)$。

(3) $x^*=\tfrac12x^*+1\Rightarrow x^*=2$。令 $b_n=a_n-2$，则
$b_{n+1}=a_{n+1}-2=\tfrac12(a_n-2)=\tfrac12 b_n$，$b_0=-2$，故
$b_n=-2\left(\tfrac12\right)^n$，
\[
  a_n=2-2^{\,1-n}.
\]
检验：$n=0$ 时 $a_0=0$，$n=1$ 时 $a_1=1$。

(4) $|F'(2)|=\tfrac12<1$，故 $|a_n-2|=|b_n|=2\cdot2^{-n}$ 指数衰减；
第 6 章的局部线性化 $F(x)\approx 2+\tfrac12(x-2)$ 正是这一步的连续版。

(5) $a_{n+1}=pa_n+q$：$x^*=px^*+q\Rightarrow x^*=\dfrac{q}{1-p}\ (p\neq1)$，
$a_n=x^*+p^n(a_0-x^*)$；$p=1$ 时退化为等差 $a_n=a_0+nq$；
$|p|>1$ 且 $a_0\neq x^*$ 时 $|a_n|\to\infty$。

(6) 由 $a_2=pa_1+q$、$a_3=pa_2+q$ 相减得
$p=\dfrac{a_3-a_2}{a_2-a_1}=\dfrac{1/4}{1/2}=\dfrac12$，
$q=a_2-pa_1=\dfrac32-\dfrac12=1$。
若 $\lim a_n=5$ 且 $p=\tfrac12$，则 $5=\dfrac q{1-1/2}=2q$，故 $q=\dfrac52$。

(7) 「取极限」需要先证明极限存在（单调有界或压缩映射）；反例：
$a_{n+1}=-a_n$、$a_0=1$ 振荡无极限，而方程 $x=-x$ 有唯一解 $0$——
\textbf{不动点存在不等于轨道收敛}。
\end{solution}

\refstepcounter{hvmasterq}
\begin{masterq}{分式迭代：$x_{n+1}=1+\dfrac1{x_n}$，$x_0=1$}
\begin{enumerate}[label=(\arabic*), itemsep=3pt]
\item \textbf{画图与数值感知}：算出 $x_1,x_2,x_3,x_4$，观察奇数项与
偶数项各自的走向。\logotag{核心公理}
\item \textbf{多视角结构翻译}：把它写成连分数 $x=1+\dfrac1{1+\dfrac1{1+\cdots}}$
与二次方程 $x^2-x-1=0$。\logotag{结构翻译}
\item \textbf{考纲规范通法}：证明奇数子列递增、偶数子列递减且有共同
极限 $\varphi=\dfrac{1+\sqrt5}{2}$（用 $\tk{}$ 搭骨架）。\logotag{考纲通法}
\item \textbf{高观点降维秒杀}：指出这是矩阵 $\begin{pmatrix}1&1\\1&0\end{pmatrix}$
的迭代，特征值之比给出收敛速率 $\dfrac{1}{\varphi^2}$（见第 5 章）。\logotag{大招公式}
\item \textbf{公式变形家族生成}：证明 $x_n=\dfrac{F_{n+2}}{F_{n+1}}$（$F_n$ 为
斐波那契数），并由此得到 $\varphi$ 的连分数与渐进分数。\logotag{公式族}
\item \textbf{参数变式与逆向追问}：若 $x_0=\varphi$ 会怎样？求使 $x_2=\varphi$
的 $x_0$；再把递推改成 $x_{n+1}=k+\dfrac1{x_n}$，求不动点与收敛范围。\logotag{充要边界}
\item \textbf{边界陷阱与错题排雷}：某同学由 $\lim x_n=L$ 直接写
$L=1+\dfrac1L$。请指出两处逻辑漏洞（存在性、正负号），并举出
$x_0<0$ 时的反例。\logotag{避坑铁律}
\end{enumerate}
\end{masterq}

\begin{solution}
(1) $x_1=2,\ x_2=\dfrac32,\ x_3=\dfrac53,\ x_4=\dfrac85$；奇数项递减、
偶数项递增，夹在 $1.6$ 附近。

(2) 迭代不变量：若 $x_n>0$，则 $x_{n+1}>1$；不动点满足
$x=1+\dfrac1x\Rightarrow x^2-x-1=0\Rightarrow \varphi=\dfrac{1+\sqrt5}{2}$
（负根 $-\dfrac1\varphi<0$ 舍去，因 $x_n>0$）。

(3) 设 $f(x)=1+\dfrac1x$，则 $x_{n+2}=f(f(x_n))=\dfrac{2x_n+1}{x_n+1}$。
直接证明：$x_{2k}$ 递增、$x_{2k+1}$ 递减，且
$|x_{n+1}-\varphi|=\dfrac{|x_n-\varphi|}{x_n\varphi}\le\dfrac1\varphi|x_n-\varphi|$
（因为 $x_n\ge1$），故两者夹逼到 $\varphi$。

(4) $\binom{u_{n+1}}{v_{n+1}}=\begin{pmatrix}1&1\\1&0\end{pmatrix}\binom{u_n}{v_n}$，
$x_n=u_n/v_n$；矩阵特征值为 $\varphi$ 与 $-1/\varphi$，故
$|x_n-\varphi|\sim C\varphi^{-2n}$（公比 $\left|\lambda_2/\lambda_1\right|
=\dfrac1{\varphi^2}$，见第 5 章探究五）。

(5) 归纳（取 $F_0=0,F_1=1$）：$x_0=\dfrac{F_2}{F_1}=\dfrac11=1$，且
$1+\dfrac1{x_n}=1+\dfrac{F_{n+1}}{F_{n+2}}
=\dfrac{F_{n+2}+F_{n+1}}{F_{n+2}}=\dfrac{F_{n+3}}{F_{n+2}}$，
故 $x_n=\dfrac{F_{n+2}}{F_{n+1}}$。

(6) $x_0=\varphi\Rightarrow$ 常数列 $\varphi$；反向递推
$x_n=\dfrac1{x_{n+1}-1}$，由 $x_2=\varphi$ 得 $x_1=\varphi$、$x_0=\varphi$。
一般地 $x_{n+1}=k+\dfrac1{x_n}$ 的不动点为
$\dfrac{k+\sqrt{k^2+4}}2$，$k>0$ 时正根稳定（$|f'|=1/x^{*2}<1$）。

(7) 漏洞一：未先证明极限存在；漏洞二：取极限时未说明保号。
方程 $L=1+\dfrac1L$ 的负根 $-\dfrac{1-\sqrt5}{2}$ 本身合法，
但它与 $x_n>0$ 的保号性矛盾，不能作为数列的极限。反例：
$x_0=-1$ 时 $x_1=0$，$x_2$ 无定义——定义域与符号必须先交代。
\end{solution}

\begin{conccard}{不动点：一把万能锚}{递推求通项与收敛性、超越方程近似解、几何变换的简化、分式迭代}{把「不动点存在」当成「轨道收敛」；未验证极限存在就取极限；忽略定义域与符号}
\vchoice{遇到「迭代／递推／方程 $f(x)=x$／变换简化」，先找不动点，
再判 $|F'|$，最后决定坐标系搬到哪里。}{记住三步：解 $x=F(x)$，
算 $|F'|$，搬原点。第 5、6、8 章都靠这一招联动。}
\end{conccard}

\begin{notezone}{章末归档：不动点二级结论 + 七问笔记功能}
\textbf{二级结论}：（1）一阶递推 $a_{n+1}=pa_n+q$ 的不动点
$x^*=\dfrac q{1-p}$，$a_n=x^*+p^n(a_0-x^*)$；
（2）稳定性 $|F'(x^*)|<1\Leftrightarrow$ 局部收缩；
（3）分式迭代的不动点＝矩阵特征方向（第 5 章）；
（4）齐次坐标 $(u,v)$ 把一维分式作用提升为二维线性作用；
（5）换坐标系（平移／旋转／齐次化）都是把不动点搬到原点。

\textbf{七问与笔记功能}：(1) 数值感知 $\to$ 概念引入；(2) 三视角
翻译 $\to$ 性质定义；(3) 考纲通法 $\to$ 规范书写；(4) 高观点秒杀
$\to$ 大招笔记；(5) 公式族 $\to$ 题型拓展；(6) 逆向追问 $\to$ 充要
边界；(7) 陷阱排雷 $\to$ 错题本。
\end{notezone}


% 示范章节（仅当 main.tex 中 \demomotrue 时挂载）
\ifdemomode
  % ==========================================================
% 示范章节：同一份源码，四种成品（约 40 行，只演示「怎么写」）
% 挂载：main.tex 中把 \demomodefalse 改为 \demomotrue，
%       或运行 python3 build_all.py --demo
% 覆盖写法：通用内容 + advancedonly + basiconly + \vchoice{}{}
%          + \tk{} + solution + warmupbox/trapbox/example
% ==========================================================
\section{示范：四合一多版本写法}

\begin{warmupbox}{读法（先看这里）}
本节不教新知识，只演示\textbf{版本分流写法}：修改 \verb|main.tex| 顶部的两行开关
（\verb|\teachertrue/\teacherfalse| 与 \verb|\basictrue/\basicfalse|），
即可从\emph{同一份源码}导出四种成品 PDF，不必复制、粘贴或维护任何副本。
\end{warmupbox}

设 $f(x)=x^{2}-2x$，考察它在区间 $[0,2]$ 上的结构。

\begin{learnobjective}
\begin{itemize}
  \item 会算 $f(0),f(1),f(2)$ 三个值并描点画草图；
  \item 会用集合数对、几何位置、代数方程三种语言描述同一条抛物线；
  \item 初步看懂「平方算子」$f\mapsto f^{2}$ 的零点行为。
\end{itemize}
\end{learnobjective}

\subsection*{第 (1) 问：具象感知与画图}
计算：$f(0)=\tk{0}$，$f(1)=\tk{-1}$，$f(2)=\tk{0}$；
描出 $(0,0),(1,-1),(2,0)$ 三点，并画出草图。

\vchoice{%
  三个函数值对应三个点；用集合数对写即
  $\{(x,y)\mid y=x^{2}-2x,\ x\in[0,2]\}$，几何上是开口向上、对称轴为
  $x=1$ 的一段抛物线。}{%
  三个点：$(0,0)$、$(1,-1)$、$(2,0)$；连起来是开口向上的小弧线，
  最低点在 $x=1$ 处。}

\subsection*{第 (2) 问：常规通法演算}
配方：$f(x)=x^{2}-2x=\tk{(x-1)^{2}-1}$，故顶点为 $\tk{(1,-1)}$，
对称轴为直线 $\tk{x=1}$。

\subsection*{第 (3) 问：三视角结构翻译}
同一条曲线可以写成三种语言：集合语言
$\{(x,f(x))\mid x\in\mathbb{R}\}$；几何语言「以 $x=1$ 为对称轴的抛物线」；
代数方程 $y=(x-1)^{2}-1$。三者说的是\emph{同一结构}。

\begin{advancedonly}
\subsection*{第 (4) 问：高观点降维秒杀（仅完整版可见）}
作平移 $u=x-1$，则 $y=u^{2}-1$：绝对尺度被消去，只剩一个比例变量 $u$，
这就是齐次化与平移归一的思想雏形。进一步看平方算子：
$f^{2}(x)=(x(x-2))^{2}$ 的零点是\textbf{水平相切的偶重零点}
（局部 $\approx 4(x-1)^{2}$ 的二重根平移），不是尖点——它解释了
「平方后符号不变、零点处压平」的现象。
\end{advancedonly}

\begin{basiconly}
\subsection*{第 (4)' 问：大白话口诀与半成品填空（仅薄弱生版可见）}
口诀：「平方不改号，零点压平地；倒数撕裂缝，无穷变成零。」
照样子填空：$f(1)=1-2=\tk{-1}$；顶点横坐标 $x=\tk{1}$。
\end{basiconly}

\begin{example}[用两种语言看同一结构]
\vchoice{双射同构映射}{一一对应的加工机器}的说法在本讲统一替换为
\vchoice{二次型结构}{抛物线的对称结构}。请用「顶点式」写出
$f(x)=x^{2}-4x+3$ 的对称轴：$\tk{x=2}$。
\end{example}
\begin{solution}
配方得 $f(x)=(x-2)^{2}-1$，故对称轴为 $x=2$，顶点 $(2,-1)$。
\end{solution}

\begin{trapbox}{顶点横坐标的符号}
配方时括号内是 $x-1$ 而不是 $x+1$：$y=(x-1)^{2}-1$ 的对称轴是 $x=1$。
若误写为 $x=-1$，是把一次项系数的符号弄反了。
\end{trapbox}

\fi

\appendix
% 附录 A：高观点结论的考纲合法化转译手册
\section{附录 A：高观点结论的考纲合法化转译手册}\label{app:legal}

\begin{commonerror}{铁律：草稿纸用高观点，答题卡写初等通法}
解答题中\textbf{严禁}出现「由仿射变换知」「由极点极线知」「由欧拉公式知」「由 Abel 变换知」「由洛必达法则知」等超纲字眼。这类词一旦出现，轻则扣步骤分，重则该问判零。高观点的正确用法是：在草稿纸上用它锁定答案、取等位置与参数边界，然后在卷面上用考纲认可的初等推导把结论\textbf{重新证一遍}——代数验算本身就是满分证明。
\end{commonerror}

\begin{conccard}[三级]{三类角色各就各位}{阅卷老师只看卷面}{把预判当成证明}
\textbf{草稿纸}：随便用，越快越准越好；\textbf{卷面}：只允许考纲工具（联立、韦达定理、作差构造、错位相减、分类讨论、切线不等式）；\textbf{采分点}：看关键等式有没有被显式写出。三者不冲突，缺的只是翻译。
\end{conccard}

下面六个模板覆盖高中压轴的高频场景，每个都按「预判 $\to$ 卷面 $\to$ 采分点」三段给出。

\subsection*{模板一：圆锥曲线中的仿射压缩}

\textbf{草稿纸预判}：椭圆沿短轴方向压缩 $b/a$ 倍即得圆，面积同比例缩小，故内接三角形最大面积为 $\dfrac{3\sqrt3}{4}ab$；斜率乘积定值 $-\dfrac{b^2}{a^2}$。

\textbf{卷面写法}：设过定点的直线为 $y=kx+m$，代入 $\dfrac{x^2}{a^2}+\dfrac{y^2}{b^2}=1$，写出韦达定理与判别式条件；由弦长公式 $|AB|=\sqrt{1+k^2}\,|x_1-x_2|$ 与点到直线距离写出面积
\[
  S=\frac12|AB|\cdot d ,
\]
整理成关于 $k,m$ 的函数，用基本不等式或求导在预判的取等点处取得最值，并验证此时满足判别式大于零。

\textbf{采分点}：韦达定理两式、弦长表达式、距离表达式、「当且仅当……时取等」以及取等条件的检验。仿射结论只用来决定往哪条直线上凑。

\subsection*{模板二：极点与极线的定点定值}

\textbf{草稿纸预判}：由极点 $P(x_0,y_0)$ 直接写出极线方程 $\dfrac{x_0x}{a^2}+\dfrac{y_0y}{b^2}=1$，目标直线即刻锁定，不必再猜。

\textbf{卷面写法}：设两交点 $A(x_1,y_1),B(x_2,y_2)$，分别代入曲线方程并作差：
\[
  \frac{x_1^2-x_2^2}{a^2}+\frac{y_1^2-y_2^2}{b^2}=0
  \ \Longrightarrow\
  \frac{(x_1+x_2)(x_1-x_2)}{a^2}+\frac{(y_1+y_2)(y_1-y_2)}{b^2}=0 .
\]
再结合 $A,B$ 都在已知直线上，整理出 $\dfrac{y_1-y_2}{x_1-x_2}$ 与 $\dfrac{y_1+y_2}{x_1+x_2}$ 的关系，最后代入要证的对称式化简即可。

\textbf{采分点}：两点式作差这一步必须显式写出；之后全是恒等变形。

\subsection*{模板三：不动点与特征根求通项}

\textbf{草稿纸预判}：解 $x^*=px^*+q$ 或 $\lambda^2-p\lambda-q=0$，读出公比与两根。

\textbf{卷面写法}：由原递推式恒等变形直接写出
\[
  a_{n+1}-x^*=p(a_n-x^*) ,
\]
或两根情形下
\[
  \frac{a_{n+1}-\lambda_1}{\,a_{n+1}-\lambda_2\,}=k\,\frac{a_n-\lambda_1}{\,a_n-\lambda_2\,} .
\]
\textbf{完全不需要解释} $x^*$、$\lambda_i$ 是怎么来的——作为「已知常数」使用，再用首项确定待定系数，最后验证 $n=1$ 成立。

\textbf{采分点}：构造式本身、首项代入、$n=1$ 的检验。

\subsection*{模板四：Abel 变换与错位相减}

\textbf{草稿纸预判}：用分部求和快速算出 $S_n$ 的闭式。

\textbf{卷面写法}：规范列出
\[
  S_n=a_1+a_2+\cdots+a_n,\qquad
  qS_n=\phantom{a_1+{}}a_1q+\cdots+a_nq ,
\]
错开一位对齐后写「两式相减得」，直接默写化简结果，并把 $n=1$ 与 $n\to\infty$（若涉及极限）单独验证。

\textbf{采分点}：两行原始展开、相减动作、最终合并式、首项检验。

\subsection*{模板五：三角恒等变换与欧拉公式}

\textbf{草稿纸预判}：忘了三倍角公式时，展开 $(\cos\theta+i\sin\theta)^3$ 取实部，秒得 $\cos3\theta=4\cos^3\theta-3\cos\theta$。

\textbf{卷面写法}：从两角和公式出发现场推导：
\[
  \cos3\theta=\cos(2\theta+\theta)=\cos2\theta\cos\theta-\sin2\theta\sin\theta=\cdots,
\]
再降幂为单角函数，最后写出所需形式。

\textbf{采分点}：起始公式明确、每一步降幂正确、最终形式与题目要求一致（平方关系、次数）。

\subsection*{模板六：导数压轴的隐零点与切线放缩}

\textbf{草稿纸预判}：用切线 $e^x\ge x+1$、$\ln x\le x-1$ 预判参数临界值 $a_0$；必要时用端点效应估出 $a_0$。

\textbf{卷面写法}：按 $a\le a_0$ 与 $a>a_0$ 分类讨论。对每一类，先证切线不等式（以 $e^x-x-1$ 求最小值的方式），再代入原函数完成放缩链；隐零点处写成「设 $x_0$ 满足 $f'(x_0)=0$，则……」，用整体代换避免解出具体数值。

\textbf{采分点}：分类标准清晰、切线不等式的完整证明、隐零点的设而不求与代换、端点值的检验。

\subsection*{自查清单}

\begin{itemize}
  \item 全篇是否出现超纲名词？若有，一律替换为等价代数操作。
  \item 关键等式是否都被显式写出（而不是「显然」「易得」）？
  \item 取等条件、判别式正负、$n=1$ 首项是否逐一检验？
  \item 参数分类是否有重叠或遗漏？
\end{itemize}

\begin{conccard}[一级]{一句话记住}{预判给方向，通法给分数}{两者都要}
\key{草稿纸负责快，卷面负责稳}。高观点决定你往哪条路上走，考纲通法决定你能拿多少分——只有前者，会丢步骤分；只有后者，会算不完。
\end{conccard}

% ==========================================================
% 附录 B：大学理工科荣誉课程（Honors Bridge）先修接口对照手册
% 说明：本文件为纯正文模块化片段（无 \documentclass），
%       由 main.tex 在 \appendix 之后 \input 挂载。
%       corebox / methodbox 现由 versions.tex 全局提供（四合一引擎），
%       配色与现有 main/third 主题一致，不再局部定义。
% ==========================================================

\section{附录 B：大学理工科荣誉课程先修接口对照手册}\label{app:honors}

\begin{corebox}{致准大学荣誉学院（Honors Program）学子}
普通高中生升入数学、物理、计算机、人工智能与电子信息等理工科专业时，普遍会经历一场剧烈的\textbf{“大一数学休克”}——例如看不懂《数学分析》严密的 $\varepsilon$-$\delta$ 极限语言，无法理解《线性代数》中抽象的基底变换与特征空间。

然而，如果你已经扎实地走完了本套讲义的第 0–9 章，\textbf{你实际上已经提前掌握了大学前两年硬核数学课的几何直觉与结构骨架}。本附录旨在为你完成最后一步\textbf{“学术暗号对接”}，将你熟悉的中学高观点，精准映射到大学专业课的正式定义与定理名称上。
\end{corebox}

\subsection*{一、 五大大学核心基础课“概念暗号”对接总表}

\begin{fastcard}{《数学分析》（Real Analysis / Honors Calculus）}
\textbf{你在本套讲义中学过的直觉与工具}：第 1 章有理数 $\mathbb{Q}$ 漏掉 $\sqrt{2}$ 孔洞，无理数填满实心数轴 $\mathbb{R}$；第 6 章局部切线近似（以直代曲）与切线支撑放缩。
\par\textbf{大学专业课正式学名与定理}：\textbf{实数完备性公理}（戴德金分割 / 确界原理）；\textbf{介值定理}与\textbf{微分中值定理}；\textbf{凸函数与 Jensen 不等式}。
\par\textbf{消除的“大一休克”痛点}：瞬间理解极限存在性证明，不再死背 $\varepsilon$-$\delta$ 符号，一眼看穿微分局部线性化本质。
\end{fastcard}

\begin{fastcard}{《线性代数》（Linear Algebra / 高等代数）}
\textbf{你在本套讲义中学过的直觉与工具}：第 5 章换基底对齐特征方向、旋转 $45^\circ$ 消除交叉项 $Bxy$；第 1 章函数变换中的主动（动图形）与被动（动坐标系）对偶；第 5 章转移矩阵特征向量锁定分式不动点。
\par\textbf{大学专业课正式学名与定理}：\textbf{线性空间与基底变换}（过渡矩阵 $P^{-1}AP$）；\textbf{实对称矩阵正交对角化}（谱定理）；\textbf{射影空间（Projective Space）}与齐次坐标。
\par\textbf{消除的“大一休克”痛点}：摆脱“矩阵只是数字表格”的机械计算误区，直接建立高维几何变换与特征空间直觉。
\end{fastcard}

\begin{fastcard}{《复变函数与信号系统》（Complex Analysis \& Signals）}
\textbf{你在本套讲义中学过的直觉与工具}：第 2 章虚数 $i$ 作为逆时针 $90^\circ$ 旋转算子；切线垂直速度演化出欧拉公式 $e^{i\theta}=\cos\theta+i\sin\theta$。
\par\textbf{大学专业课正式学名与定理}：\textbf{复数域的几何表示与保角映射}；\textbf{傅里叶级数与傅里叶变换}（正交旋转基展开）；工科电路\textbf{相量法（Phasor）}。
\par\textbf{消除的“大一休克”痛点}：在大学物理简谐振动、交流电路与信号处理中，直接把正弦量当成复平面旋转向量处理。
\end{fastcard}

\begin{fastcard}{《常微分方程与数值分析》（ODEs \& Numerical Methods）}
\textbf{你在本套讲义中学过的直觉与工具}：第 5 章差分 $\Delta a_n$ 与求和互逆、Abel 分部求和；蛛网图 $|f'(x^*)|<1$ 螺旋收敛于吸引不动点；生成函数 $G(x)=\sum a_n x^n$。
\par\textbf{大学专业课正式学名与定理}：\textbf{差分方程与常微分方程特征根法}；\textbf{巴拿赫压缩映射原理}（牛顿迭代法收敛性）；\textbf{Z 变换与拉普拉斯变换}。
\par\textbf{消除的“大一休克”痛点}：打通离散递归与连续微分方程，理解计算机迭代算法为何收敛以及混沌分岔的临界条件。
\end{fastcard}

\begin{fastcard}{《公理化概率论与机器学习》（Probability \& Machine Learning）}
\textbf{你在本套讲义中学过的直觉与工具}：第 7 章样本空间 $\Omega$ 归一化测度与条件子集缩放；二维散点云的竖直欧氏距离平方和最小投影线。
\par\textbf{大学专业课正式学名与定理}：\textbf{Kolmogorov 概率测度空间} $(\Omega, \mathcal{F}, P)$；\textbf{线性回归的最小二乘正交投影}（正规方程 $A^TA\hat{x}=A^Tb$）。
\par\textbf{消除的“大一休克”痛点}：跳过繁琐的古典计数，直接站在高维空间几何投影的高度理解统计学习与数据拟合。
\end{fastcard}

\subsection*{二、 三大核心概念的“中学直觉 $\to$ 大学严格化”跃迁示范}

为了让你在翻开大学教材的第一天就能写出严密的现代数学证明，以下展示三个最核心概念如何从本讲义的“几何直觉”平滑过渡到大学“公理化语言”：

\begin{methodbox}{跃迁 1：从“实心金条无漏洞”到《数学分析》的“确界原理”}
\begin{itemize}
    \item \textbf{中学高观点直觉}（第 1 章）：有理数轴 $\mathbb{Q}$ 像布满针眼的筛子，曲线 $y=x^2-2$ 会从 $\sqrt{2}$ 的孔洞漏过去；实数轴 $\mathbb{R}$ 填满了所有孔洞，因此有上界的递增序列必然会撞上一个确定的天花板。
    \item \textbf{大学《数学分析》严格表述（确界公理）}：设 $S \subset \mathbb{R}$ 为非空有上界实数集，则 $S$ 在 $\mathbb{R}$ 中必有\textbf{上确界（最小上界 $\sup S$）}。由此可直接推出“单调有界数列必有极限”与“闭区间套定理”。
\end{itemize}
\end{methodbox}

\begin{methodbox}{跃迁 2：从“旋转 $45^\circ$ 消交叉项 $xy$”到《线性代数》的“正交对角化”}
\begin{itemize}
    \item \textbf{中学高观点直觉}（第 5 章）：双曲线 $xy = 1$ 的主轴倾斜了 $45^\circ$，只要把坐标轴旋转 $45^\circ$ 对齐它的对称轴，交叉项就会消失，化为 $\frac{x'^2}{2} - \frac{y'^2}{2} = 1$。
    \item \textbf{大学《线性代数》严格表述（二次型主轴定理）}：将 $xy = 1$ 写成对称矩阵二次型：
    $$\begin{pmatrix} x & y \end{pmatrix} \begin{pmatrix} 0 & \frac{1}{2} \\ \frac{1}{2} & 0 \end{pmatrix} \begin{pmatrix} x \\ y \end{pmatrix} = 1$$
    求实对称矩阵 $A = \begin{pmatrix} 0 & 1/2 \\ 1/2 & 0 \end{pmatrix}$ 的特征值 $\lambda_1 = \frac{1}{2}, \lambda_2 = -\frac{1}{2}$ 及其单位正交特征向量（构成旋转 $45^\circ$ 的正交矩阵 $Q$），令 $\mathbf{x} = Q\mathbf{x}'$，立刻得到标准型 $\frac{1}{2}x'^2 - \frac{1}{2}y'^2 = 1$！
\end{itemize}
\end{methodbox}

\begin{methodbox}{跃迁 3：从“蛛网图台阶变小”到《泛函/数值分析》的“压缩映射原理”}
\begin{itemize}
    \item \textbf{中学高观点直觉}（第 5 章）：在蛛网图中，只要不动点处曲线斜率满足 $|f'(x^*)| < 1$，每次投影都会把误差缩短，序列最终螺旋收敛于 $x^*$。
    \item \textbf{大学严格表述（Banach Fixed-Point Theorem）}：设 $(X, d)$ 为完备度量空间，映射 $T: X \to X$ 满足存在常数 $0 \le L < 1$，使得对任意 $x, y \in X$ 均有 $d(T(x), T(y)) \le L \cdot d(x, y)$，则 $T$ 在 $X$ 中存在\textbf{唯一不动点 $x^*$}，且对任意初值 $x_0 \in X$，迭代序列 $x_{n+1} = T(x_n)$ 均以几何级数速度收敛于 $x^*$。
\end{itemize}
\end{methodbox}

\subsection*{三、 给荣誉学院新生的阅读进阶路线图}

当你顺利进入大学理工科荣誉班级后，推荐按以下书单无缝衔接本讲义的知识结构：
\begin{enumerate}
    \item \textbf{数学分析方向}：陶哲轩《陶哲轩实分析》（从集合与自然数公理重新严格建构整个实数与微积分体系，与本讲义第 0、1 章风格一脉相承）；
    \item \textbf{线性代数方向}：Sheldon Axler《Linear Algebra Done Right》（摒弃繁琐行列式，直接从向量空间与线性算子的不变子空间/特征向量切入，与本讲义换基底思想高度契合）；
    \item \textbf{复分析与几何直觉方向}：Tristan Needham《复分析：可视化方法》（用纯几何旋转与向量场讲透复变函数，完美延续本讲义第 2 章与 Manim 动画风格）。
\end{enumerate}


\end{document}
，避免导言区内容四处重复、彼此漂移。
% ============================================================

\title{数学导学案·高观点系列}
\author{Shawn Shang}
\date{2026 年 11 月}

\graphicspath{{../figures/}{../map/figures/}}

\begin{document}

\begin{titlepage}
  \thispagestyle{empty}
  \centering
  \makebox[\textwidth]{%
    \includegraphics[width=\paperwidth]{cover_cyan.png}}\par
  \vfill
  {\Huge\bfseries\color{ecolor!70!black} 数学导学案·高观点系列\par}
  \vskip 1.2em
  {\large 结构主义视角下的高中数学\par}
  \vskip 0.8em
  {\large 第 0 章 — 第 9 章\par}
  \vskip 0.8em
  {\large\versionlabel\par}
  \vskip 2.5em
  {\large Shawn Shang\par}
  \vfill
  {\normalsize 2026 年 11 月\par}
\end{titlepage}
\tableofcontents

% 让第一个 \section 编号为 0（第 0 章是总纲）
\setcounter{section}{-1}

% 高观点导学案：集合、逻辑演算与证明武器库（第 0 章·总纲）
% 主线：把“集合与逻辑”从概念背诵升级为贯穿高中解题的“题干同解变形与严格证明通用武器库”
% 四模块：集合语言与同解变形公理 → 命题演算三大变形引擎 → 两大间接证明支柱 → 武器库索引表与跨模块实战
% 全程用“集合包含 ↔ 逻辑等价 ↔ 量词否命题”三种语言互译，并用韦恩图与多米诺图把抽象逻辑可视化。

\section{集合、逻辑演算与证明武器库（高观点总纲）}

本章不引入新公式，而是把教材第 1 章的集合、常用逻辑用语重新锻造为一套“解题武器库”：
面对题干，先问它是哪一种\key{对象集合}（解集、值域、参数范围），再问它需要哪一种\key{变形引擎}
（同解变形、德摩根律、逆否等价、补集转化、反证法、数学归纳法）。四个模块环环相扣：
\key{集合语言与同解变形公理}（$A=B\iff A\subseteq B\land B\subseteq A$）$\to$
\key{命题演算三大变形引擎}（德摩根律、逆否等价、量词否定）$\to$
\key{两大间接证明支柱}（反证法、数学归纳法）$\to$
\key{武器库索引表与跨模块实战}。
核心思想：\key{题干中的每一个关键词，都对应一次集合运算或一次逻辑等价变形}；
解题失误往往不是计算错误，而是\key{没审清变形方向}（增根或漏根）或\key{没找对证明支柱}。

\begin{learnobjective}
\begin{itemize}
  \item 掌握集合相等判定公理 $A=B\iff(A\subseteq B)\land(B\subseteq A)$，能用它审查方程、不等式变形的同解性；
  \item 熟记同解变形防增漏根清单：平方（范围扩大易增根）、同除含参因式（范围缩小易漏根）、对数真数（定义域子集）；
  \item 掌握德摩根律 $\neg(p\land q)\iff\neg p\lor\neg q$ 与逆否等价 $(p\Rightarrow q)\iff(\neg q\Rightarrow\neg p)$，能用补集转化处理“至多、至少、不都、无解”；
  \item 掌握量词否定 $\neg\forall=\exists\neg$、$\neg\exists=\forall\neg$，会把“恒成立／能成立”翻译为值域的子集包含／交集非空；
  \item 理解反证法即证 $A\cap\complement_U B=\varnothing$，识别“唯一性、不存在、无理数、无限性”四类触发信号词；
  \item 理解数学归纳法的皮亚诺公理与良序原理本质，会用最小反例原理分析“归纳基础＋归纳递推”两步的缺一不可。
\end{itemize}
\end{learnobjective}

\begin{preview}
\subsection*{知识回顾}
\begin{itemize}
  \item $A\subseteq B$ 的定义是：对任意 $a\in A$，都有 \blank；
  \item 若 $A=\{x\mid x^2=1\}$，$B=\{-1,1\}$，则 $A$ 与 $B$ 的关系是 \blank（相等／互不包含）；
  \item 命题“若 $p$ 则 $q$”的逆否命题是 \blank；
  \item 德摩根律：$\neg(p\land q)\iff$ \blank。
\end{itemize}

\subsection*{教材速读}
阅读教材第 1 章“集合与常用逻辑用语”及本章导学案，完成下列问题：
\begin{enumerate}
  \item 补集 $\complement_U A$ 的定义是 $\{x\in U\mid \blank\}$；
  \item “$p$ 是 $q$ 的充分条件”用集合语言说是 \blank；
  \item “方程 $f(x)=0$ 无解”用集合语言说是 \blank；
  \item 数学归纳法证明的两步是 \blank 与 \blank。
\end{enumerate}
\end{preview}

\begin{explore}

\subsection{集合语言与同解变形公理}

\key{探究一：集合相等——互相包含才是相等}
集合相等不是“看起来一样”，而是\key{互相包含}：
\[
  A=B\iff(A\subseteq B)\land(B\subseteq A).
\]
证明 $A=B$ 必须证两个方向；只证 $A\subseteq B$ 只能得到包含，不能得到相等。
这条公理还是审查方程（不等式）变形的“合法性标准”：
把原方程的解集记为 $A$，变形后方程的解集记为 $B$，则
\begin{itemize}
  \item $A=B$：变形\key{同解}，解集不变，安全；
  \item $A\subsetneqq B$：变形后解集\key{变大}，多出的解是\key{增根}；
  \item $B\subsetneqq A$：变形后解集\key{变小}，丢掉的解是\key{漏根}。
\end{itemize}
因此解方程、解不等式的每一步都要问一句：\key{这一步把解集往哪个方向变了？}

\begin{center}
\footnotesize
\setlength{\tabcolsep}{3pt}
\renewcommand{\arraystretch}{1.4}
\begin{tabular}{p{2.6cm}p{2.8cm}p{3.2cm}p{3.2cm}}
\toprule
变形操作 & 解集关系 & 方向与风险 & 应对 \\
\midrule
两边加/减同一整式 & $A=B$ & 同解，安全 & 直接使用 \\
两边乘/除以非零常数 & $A=B$ & 同解，安全 & 先确认非零 \\
两边平方（或去绝对值） & $A\subsetneqq B$ & 范围扩大，易增根 & 必须验根 \\
两边同除含参/未知因式 & $B\subsetneqq A$ & 范围缩小，易漏根 & 移项分解，讨论零 \\
取对数/解对数方程 & $A\subsetneqq B$ & 真数 $>0$ 限制，易增根 & 先限定定义域 \\
\bottomrule
\end{tabular}
\end{center}

\key{风险一：平方变形范围扩大，易增根。}
由 $f(x)=g(x)$ 平方得 $f^2(x)=g^2(x)$，而
\[
  f^2=g^2\iff(f-g)(f+g)=0\iff f=g\ \text{或}\ f=-g,
\]
凭空多出 $f=-g$ 的解。必须验根。

\begin{example}
解方程 $\sqrt{x+1}=x-1$，并指出直接平方会多出哪个根。
\end{example}
\begin{solution}
两边平方（须先注意右边非负：$x-1\geq0$，即 $x\geq1$）：
\[
  x+1=(x-1)^2=x^2-2x+1\implies x^2-3x=0\implies x=0\ \text{或}\ x=3.
\]
检验：$x=0$ 时左边 $=1$，右边 $=-1$，不相等，是平方后产生的\key{增根}，舍去；
$x=3$ 时左边 $=2$，右边 $=2$，成立。故解集为 $\{3\}$。
\end{solution}

\key{风险二：同除含参因式范围缩小，易漏根。}
方程两边同除以含未知数（或参数）的因式，若该因式可能为 $0$，就会丢掉使因式为 $0$ 的解。
正确做法是移项因式分解，并单独讨论因式为 $0$ 的情形。

\begin{example}
解关于 $x$ 的方程 $x^2=ax$（$a$ 为参数），并说明“两边同除以 $x$”会漏掉什么。
\end{example}
\begin{solution}
移项分解：$x(x-a)=0$，故 $x=0$ 或 $x=a$。
当 $a\neq0$ 时解集为 $\{0,a\}$；当 $a=0$ 时方程为 $x^2=0$，解集为 $\{0\}$（两个相等实根）。
若直接两边同除以 $x$，得到 $x=a$，会漏掉 $x=0$（$a\neq0$ 时），
因为 $x$ 可能为 $0$，同除使解集由 $\{0,a\}$ 缩小为 $\{a\}$，属于“范围缩小易漏根”。
\end{solution}

\key{风险三：对数真数限制定义域子集，易增根。}
\[
  \log_a f(x)=\log_a g(x)\iff f(x)=g(x)>0,
\]
其中“$>0$”来自真数的定义域限制。定义域
$D=\{x\mid f(x)>0,\ g(x)>0\}$ 是一个\key{子集约束}，解必须落在 $D$ 内；忽略它就增根。
（同理，$a^{\log_a x}=x$ 要求 $x>0$，而 $\log_a a^x=x$ 对一切实数 $x$ 成立。）

\begin{example}
解方程 $\log_2(x^2-4)=\log_2(3x)$。
\end{example}
\begin{solution}
由对数定义，真数必须为正：$x^2-4>0$ 且 $3x>0$，即 $x>2$。
在此定义域内 $\log_2$ 是单射，故 $x^2-4=3x$，即 $x^2-3x-4=0$，解得 $x=4$ 或 $x=-1$。
$x=-1$ 不在定义域 $x>2$ 内，是增根，舍去。故解集为 $\{4\}$。
\end{solution}

\begin{conccard}[一级]{集合相等判定公理：$A=B\iff(A\subseteq B)\land(B\subseteq A)$}{证明集合相等、审查方程（不等式）同解变形}{只证 $A\subseteq B$ 一个方向就宣布相等；只算“怎么变形”，不算解集往哪变}
集合相等的本质是“互相包含”。证 $A=B$ 必须证两个方向：$A\subseteq B$ 且 $B\subseteq A$；
只证一个方向只能得到包含。把方程（不等式）的解集记为 $A$、变形后解集记为 $B$，
只有 $A=B$ 时变形才同解；$A\subsetneqq B$ 则增根，$B\subsetneqq A$ 则漏根。
这条公理是全部同解变形的“合法性审查标准”。
\end{conccard}

\begin{commonerror}{同解变形不审方向}
常见错误：平方后不验根、同除含参因式不讨论零、解对数方程不写定义域。
它们都源于同一个疏漏——只算“怎么变形”，不算“解集往哪个方向变”。
例如由 $\sqrt{x+1}=x-1$ 平方得 $x^2-3x=0$，解集由 $\{3\}$ 扩大为 $\{0,3\}$，多出的 $0$ 就是增根；
由 $x^2=ax$ 两边除以 $x$ 得 $x=a$，解集由 $\{0,a\}$ 缩小为 $\{a\}$，漏掉了 $0$。
\end{commonerror}

\subsection{命题演算三大变形引擎}

\key{探究二：德摩根律——“且”的否定是“或”}
\[
  \neg(p\land q)\iff\neg p\lor\neg q,\qquad
  \neg(p\lor q)\iff\neg p\land\neg q.
\]
对应的集合版本（补集的分配律）：
\[
  \complement_U(A\cap B)=\complement_U A\cup\complement_U B,\qquad
  \complement_U(A\cup B)=\complement_U A\cap\complement_U B.
\]
记忆：否定“且”要变“或”，否定“或”要变“且”，并且每个命题都要否定。

\key{引擎二：逆否等价——正难则反}
\[
  (p\Rightarrow q)\iff(\neg q\Rightarrow\neg p).
\]
原命题与逆否命题同真同假，这是“正难则反”的合法性来源。它还有集合版本：

\begin{derivation}
要证 $A\subseteq B\iff\complement_U B\subseteq\complement_U A$。
\begin{align*}
  A\subseteq B &\iff \forall x\,(x\in A\Rightarrow x\in B)\\
  &\iff \forall x\,(x\notin B\Rightarrow x\notin A) &&(\text{逆否等价})\\
  &\iff \forall x\,(x\in\complement_U B\Rightarrow x\in\complement_U A)\\
  &\iff \complement_U B\subseteq\complement_U A.
\end{align*}
\end{derivation}

可见：要证“小集合含于大集合”（$A\subseteq B$），可改证“大集合的补集含于小集合的补集”
（$\complement_U B\subseteq\complement_U A$），往往更容易。

\key{引擎三：量词否定与补集转化流程}
\[
  \neg(\forall x,P(x))\iff\exists x,\neg P(x),\qquad
  \neg(\exists x,P(x))\iff\forall x,\neg P(x).
\]
口诀：\key{量词互换，结论否定}。与德摩根律联用：
\[
  \neg\bigl(\forall x,\,p(x)\land q(x)\bigr)\iff\exists x,\,\neg p(x)\lor\neg q(x).
\]
题干关键词的补集转化流程：
\begin{itemize}
  \item “至少有一个”（$\exists$）的否定是“一个都没有”（$\forall,\neg$）；
  \item “至多有一个”= 否定“至少有两个”；
  \item “不都”= $\neg\forall$ = “至少有一个不”；
  \item “无解”= 解集 $\varnothing$ = 补集为全集 $U$。
\end{itemize}
流程：把“至多／至少／不都／无解”改写成“$\exists/\forall$ + 否定” $\to$ 用德摩根与量词否定化简
$\to$ 求补集（\key{补集法}）$\to$ 取补并检查端点。

\key{恒成立／能成立：值域的子集包含与交集非空。}
设 $f:D\to\mathbb{R}$ 的值域为 $R_f=\{f(x)\mid x\in D\}$，目标集合为 $M$，则
\[
  \forall x\in D,\ f(x)\in M \iff R_f\subseteq M,\qquad
  \exists x\in D,\ f(x)\in M \iff R_f\cap M\neq\varnothing.
\]
特例 $M=(a,+\infty)$：恒成立 $\iff\min R_f>a$（看最小值）；能成立 $\iff\max R_f>a$（看最大值）。
例如 $f(x)=x^2-2x$ 在 $[0,3]$ 上的值域为 $[-1,3]$，则 $f(x)>a$ 恒成立 $\iff a<-1$，
能成立 $\iff a<3$。

\begin{example}
写出命题“$\exists x\in\mathbb{R},\ x^2-2x+a\leq0$”的否定；
并说明“方程 $x^2-2x+a=0$ 无实根”用集合语言如何表述。
\end{example}
\begin{solution}
否定为“$\forall x\in\mathbb{R},\ x^2-2x+a>0$”（量词互换，结论否定）。
“无实根”即解集为空集：$\{x\mid x^2-2x+a=0\}=\varnothing$，
等价于“对任意 $x\in\mathbb{R}$，$x^2-2x+a\neq0$”，也等价于 $\Delta=4-4a<0$，即 $a>1$。
\end{solution}

\begin{example}
已知“存在 $x\in[1,2]$，使 $x^2-ax+1<0$”为真，求 $a$ 的取值范围（用补集法）。
\end{example}
\begin{solution}
因为 $x\in[1,2]$ 为正，原命题等价于
\[
  \exists x\in[1,2],\ a>x+\frac1x
  \iff a>\min_{x\in[1,2]}\Bigl(x+\frac1x\Bigr).
\]
由均值不等式 $x+\frac1x\geq2$，当 $x=1$ 时取等号，故最小值为 $2$，于是 $a>2$。

补集法验证：其否定为“$\forall x\in[1,2],\ x^2-ax+1\geq0$”，即 $a\leq x+\frac1x$ 恒成立，
得 $a\leq2$；取补集得 $a>2$，与上面一致。
\end{solution}

\begin{problem}
设 $f(x)=x^2-2x$ 在 $[0,3]$ 上的值域为 $R_f$。分别求“$f(x)>a$ 恒成立”与“$f(x)>a$ 能成立”时 $a$ 的范围。
\end{problem}
\begin{solution}
$R_f=[-1,3]$。恒成立 $\iff R_f\subseteq(a,+\infty)\iff a<-1$；
能成立 $\iff R_f\cap(a,+\infty)\neq\varnothing\iff a<3$。
\end{solution}

\begin{conccard}[二级]{逆否等价与补集反向包含：$p\Rightarrow q\iff\neg q\Rightarrow\neg p$；$A\subseteq B\iff\complement_U B\subseteq\complement_U A$}{证明充分条件、正难则反的题目}{把“逆命题”与“逆否命题”混同；把“否命题”当“命题的否定”}
原命题与逆否命题同真同假，这是“正难则反”的合法性来源。集合语言下，
要证 $A\subseteq B$ 可改证 $\complement_U B\subseteq\complement_U A$——
把“小推大”换成“大补推小补”，往往更容易。
\end{conccard}

\begin{conccard}[二级]{恒成立／能成立：值域的子集包含与交集非空}{含参恒成立、能成立求参数范围}{把“恒成立”与“能成立”方向用反；忽略值域端点是否取到}
设 $f$ 在 $D$ 上的值域为 $R_f$。$\forall x\in D,\ f(x)\in M\iff R_f\subseteq M$（看最小值）；
$\exists x\in D,\ f(x)\in M\iff R_f\cap M\neq\varnothing$（看最大值）。
特例 $M=(a,+\infty)$：恒成立 $\iff\min R_f>a$，能成立 $\iff\max R_f>a$。
\end{conccard}

\begin{commonerror}{把“否命题”当“命题的否定”}
命题“若 $p$ 则 $q$”的\key{否定}是“$p$ 且非 $q$”（$p\land\neg q$），它的\key{否命题}是“若非 $p$ 则非 $q$”（$\neg p\Rightarrow\neg q$）。
两者不同：否定只否定结论，否命题条件与结论都否定。
例如“若 $x>2$ 则 $x^2>4$”的否定是“存在 $x>2$ 使 $x^2\leq4$”；
其否命题“若 $x\leq2$ 则 $x^2\leq4$”是假命题（取 $x=-3$）。
\end{commonerror}

\begin{method}
补集转化法（正难则反）：遇到“至多、至少、不都、无解、并非”等关键词，先写出它的否定
（量词互换、联结词用德摩根），把“存在”转“任意”或把“任意”转“存在”，
求出否定命题成立的参数范围，再取补集。步骤：
\begin{enumerate}
  \item 定位量词（$\forall/\exists$）与联结词（$\land/\lor$）；
  \item 否定量词、否定结论，用德摩根律化简；
  \item 求否定命题对应的集合（补集法求参数）；
  \item 取补集，并检查端点。
\end{enumerate}
\end{method}

\subsection{两大间接证明支柱}

\key{探究三：反证法——证明反例集为空}
要证命题 $p\Rightarrow q$。设 $U$ 为论域，$A=\{x\mid p(x)\}$（条件成立），
$B=\{x\mid q(x)\}$（结论成立），则要证 $A\subseteq B$。
反证法不是直接证 $A\subseteq B$，而是证明\key{反例集为空}：
\[
  A\cap\complement_U B=\varnothing.
\]
做法：假设存在反例 $x_0\in A\cap\complement_U B$（满足条件却不满足结论），推出矛盾，
从而交集为空，即 $A\subseteq B$。

\begin{figure}[htbp]
  \centering
  \begin{tikzpicture}[>=Stealth, semithick, font=\footnotesize]
    \fill[gray!12] (0,0) rectangle (8,5);
    \draw[thick, main!70!black] (0,0) rectangle (8,5);
    \node[main!70!black] at (7.5,4.5) {$U$};
    \fill[main!10, main!70!black] (4.2,2.5) circle (1.8);
    \fill[second!20, second] (3.7,2.5) circle (0.9);
    \node[main!70!black] at (4.9,3.7) {$B$};
    \node[second] at (3.7,2.5) {$A$};
    \node[gray!60!black] at (1.1,2.6) {$\complement_U B$};
    \node[red!70!black] at (1.1,1.6) {$A\cap\complement_U B=\varnothing$};
  \end{tikzpicture}
  \caption{反证法的集合图景：条件集合 $A$ 完全落在结论集合 $B$ 内，反例集 $A\cap\complement_U B$ 为空}
  \label{fig:contradiction}
\end{figure}

反证法的第一步是明确写出“反设”：$A\not\subseteq B$，即存在 $x\in A$ 且 $x\notin B$。
图 \ref{fig:contradiction-assumption} 画出这一反设态；随后由条件与 $B$ 的定义推出矛盾，红色月牙收缩为空，回到图 \ref{fig:contradiction} 的结论态。

\begin{figure}[htbp]
  \centering
  \begin{tikzpicture}[>=Stealth, semithick, font=\footnotesize]
    \fill[gray!12] (0,0) rectangle (8,5);
    \draw[thick, main!70!black] (0,0) rectangle (8,5);
    \node[main!70!black] at (7.5,4.5) {$U$};

    \coordinate (B) at (2.7,2.6);
    \coordinate (A) at (4.9,2.6);

    % B（灰底）与 A；再把 A\B 涂成红色月牙
    \fill[main!14] (B) circle (1.6);
    \begin{scope}
      \clip (A) circle (1.4);
      \fill[red!50] (A) circle (1.4);
      \fill[main!14] (B) circle (1.6);
    \end{scope}

    \draw[main!60!black, thick] (B) circle (1.6);
    \draw[second!70!black, thick] (A) circle (1.4);

    \node[main!70!black] at (1.85,3.9) {$B$};
    \node[second!80!black] at (5.5,3.55) {$A$};
    \node[gray!60!black] at (0.75,4.4) {$\complement_U B$};

    \fill[red!80!black] (5.8,2.6) circle (2.2pt);
    \node[red!80!black, right] at (5.87,2.6) {$x$};
    \node[red!70!black, align=center] at (6.05,1.0) {$A\cap\complement_U B$\\ $\neq\varnothing$};
    \draw[red!60!black, ->, shorten >=1pt] (6.25,1.3) -- (5.9,2.1);

    \node[red!70!black] at (4.0,0.55) {反设：$\exists x\in A$ 且 $x\notin B$};
  \end{tikzpicture}
  \caption{反证法的反设态（正难则反）：假设 $A\not\subseteq B$，集合 $A$ 有一部分越出 $B$，红色月牙区域 $A\cap\complement_U B$ 非空，反例点 $x$ 落在其中}
  \label{fig:contradiction-assumption}
\end{figure}

\key{触发信号词：}
\begin{itemize}
  \item \key{唯一性}（“只有一个”“恰有一个”“唯一”）——假设存在两个不同的对象，推出矛盾（或证明任意两个相等）；
  \item \key{不存在}（“不存在”“没有”“无”）——假设存在，推出矛盾；
  \item \key{无理数}（“$\sqrt2$ 是无理数”）——假设可写成既约分数，推出分子分母同为偶数，与既约矛盾；
  \item \key{无限性}（“有无穷多个”）——假设只有有限个，推出还能构造出下一个，矛盾。
\end{itemize}

\begin{example}
用反证法证明：$\sqrt2$ 不是有理数。
\end{example}
\begin{solution}
假设 $\sqrt2$ 是有理数，则可写成既约分数 $\sqrt2=\dfrac pq$
（$p,q\in\mathbb{Z}^*$，$q>0$，$\gcd(p,q)=1$）。
两边平方得 $p^2=2q^2$，故 $p^2$ 为偶数，从而 $p$ 为偶数（奇数的平方仍为奇数）。
设 $p=2k$，代入得 $4k^2=2q^2$，即 $q^2=2k^2$，故 $q$ 也为偶数。
于是 $p,q$ 均为偶数，与 $\gcd(p,q)=1$ 矛盾。故假设不成立，$\sqrt2$ 是无理数。
\end{solution}

\key{探究四：数学归纳法——皮亚诺公理与最小反例原理}
数学归纳法不是“递推直觉”，而是自然数结构的公理推论。皮亚诺公理（$\mathbb{N}^*$）要点：
\begin{enumerate}
  \item $1\in\mathbb{N}^*$；
  \item 每个 $n\in\mathbb{N}^*$ 有唯一后继 $n+1$；
  \item $1$ 不是任何数的后继；
  \item 后继函数是单射：$m+1=n+1\Rightarrow m=n$；
  \item \key{归纳公理}：若 $S\subseteq\mathbb{N}^*$ 满足 $1\in S$ 且 $n\in S\Rightarrow n+1\in S$，则 $S=\mathbb{N}^*$。
\end{enumerate}
与归纳公理等价的是\key{良序原理}：$\mathbb{N}^*$ 的任一非空子集必有最小元。
由此得\key{最小反例原理}：若命题 $P(n)$ 不恒成立，则存在最小反例 $m$。
由 $m$ 最小知 $P(m-1)$ 成立，再用归纳递推得 $P(m)$ 成立，矛盾；故 $P(n)$ 恒成立。

因此数学归纳法两步缺一不可：
\begin{itemize}
  \item \key{归纳基础}：验证 $P(n_0)$ 成立（第一块骨牌倒下）；
  \item \key{归纳递推}：假设 $P(k)$ 成立，证 $P(k+1)$ 成立（前一块倒下必带倒后一块）。
\end{itemize}
只有基础是“只倒第一块”；只有递推是“没有启动”。归纳起点 $n_0$ 可以平移（不一定从 $1$ 开始）。

\begin{figure}[htbp]
  \centering
  \begin{tikzpicture}[>=Stealth, semithick, font=\footnotesize]
    \draw[->, gray] (0,0)--(10.8,0) node[right]{$n\in\mathbb{N}^*$};
    \foreach \i/\lab in {0/1,1/2,2/3,3/4,4/5,5/6} {
      \pgfmathsetmacro{\xx}{0.8+1.5*\i}
      \draw[main!70!black, fill=main!12, rounded corners=0.8pt] (\xx,0) rectangle ++(0.5,2.0);
      \node at (\xx+0.25,1.0) {$\lab$};
    }
    \foreach \i in {0,1,2,3,4} {
      \pgfmathsetmacro{\xx}{1.3+1.5*\i}
      \draw[->, second, thick] (\xx,1.0)--++(0.55,0);
    }
    \node[below, main] at (0.8,-0.15) {推倒第 $1$ 块};
    \node[above, second] at (5.0,2.3) {若第 $k$ 块倒下，则第 $k+1$ 块必倒下};
    \node[below, gray!70!black] at (8.2,-0.15) {$\cdots$};
    \node[below, gray!70!black] at (9.8,-0.15) {全部倒下};
  \end{tikzpicture}
  \caption{数学归纳法的多米诺模型：归纳基础推倒第 1 块，归纳递推保证“倒一块必带倒下一块”，于是全部倒下}
  \label{fig:induction}
\end{figure}

\begin{example}
用数学归纳法证明：$7^n-1$ 能被 $6$ 整除（$n\in\mathbb{N}^*$）。
\end{example}
\begin{solution}
(1) 基础：$n=1$ 时 $7^1-1=6$，能被 $6$ 整除。
(2) 递推：假设 $n=k$ 时 $7^k-1$ 能被 $6$ 整除，即 $7^k-1=6m$。则 $n=k+1$ 时
\[
  7^{k+1}-1=7\cdot7^k-1=7(7^k-1)+6=7\cdot6m+6=6(7m+1),
\]
能被 $6$ 整除。由(1)(2)，命题对一切 $n\in\mathbb{N}^*$ 成立。
\end{solution}

\subsection{武器库索引表与跨模块实战}

\key{《高中数学逻辑变形与证明触发信号对照表》}
把题干信号、集合／逻辑翻译、变形引擎与操作风险放在同一张表，解题时“对号入座”。

\begin{center}
\footnotesize
\setlength{\tabcolsep}{3pt}
\renewcommand{\arraystretch}{1.4}
\begin{tabular}{p{2.0cm}p{3.2cm}p{2.6cm}p{3.6cm}}
\toprule
题干信号 & 逻辑／集合翻译 & 变形引擎 & 操作与风险 \\
\midrule
集合相等 $A=B$ & 双向包含 & 集合相等公理 & 分证 $A\subseteq B$ 与 $B\subseteq A$ \\
\hline
平方、去绝对值 & $f^2=g^2\iff f=\pm g$ & 同解变形 & 解集扩大，必须验根 \\
\hline
同除含参因式 & 可能丢根 & 移项分解＋讨论 & 解集缩小，防漏根 \\
\hline
对数方程 & 真数 $>0$ & 同解变形＋定义域 & 忽略定义域则增根 \\
\hline
至少有一个 & $\exists x$ & 量词否定 & 否定为 $\forall x,\ \neg$ \\
\hline
至多有一个 & 否定“至少两个” & 量词否定＋德摩根 & 命题的否定改写 \\
\hline
不都 & $\neg\forall$ & 量词否定 & 否定为“至少有一个不” \\
\hline
无解、为空集 & $S=\varnothing$ & 补集 & 补集为全集 $U$ \\
\hline
恒成立 & $\forall x,\ f(x)\in M$ & 值域子集 & $R_f\subseteq M$，看最小值 \\
\hline
能成立 & $\exists x,\ f(x)\in M$ & 值域交集 & $R_f\cap M\neq\varnothing$，看最大值 \\
\hline
唯一性 & 存在且唯一 & 反证法 & 假设第二个，推出矛盾 \\
\hline
不存在、没有 & $\neg\exists$ & 反证法 & 假设存在，推出矛盾 \\
\hline
无理数 & 非有理数 & 反证法 & 设 $\sqrt2=\frac pq$ 既约 \\
\hline
无限性 & 非有限 & 反证法 & 假设有限，构造下一个 \\
\hline
对所有 $n\in\mathbb{N}^*$ & 自然数命题 & 数学归纳法 & 基础＋递推，缺一不可 \\
\hline
整除、离散不等式 & 自然数命题 & 数学归纳法／最小反例 & 归纳起点可平移至 $n_0$ \\
\bottomrule
\end{tabular}
\end{center}

\begin{example}
解不等式 $\dfrac{x^2-4}{x-2}\geq0$，并说明直接约去 $(x-2)$ 会丢掉什么。
\end{example}
\begin{solution}
定义域要求 $x-2\neq0$。在定义域内
\[
  \frac{x^2-4}{x-2}=\frac{(x-2)(x+2)}{x-2}=x+2,
\]
故原不等式等价于 $x+2\geq0$ 且 $x\neq2$，即 $x\geq-2$ 且 $x\neq2$，
解集为 $[-2,2)\cup(2,+\infty)$。
若直接约去 $(x-2)$ 得 $x+2\geq0$，解集为 $[-2,+\infty)$，把分母不能为 $0$ 的限制丢了（漏掉限制，不是漏根，而是扩大了解集），
所以必须保留 $x\neq2$。
\end{solution}

\begin{example}
若“存在 $x\in\mathbb{R}$，使 $x^2+2ax+4<0$”是假命题，求 $a$ 的取值范围。
\end{example}
\begin{solution}
该命题为假 $\iff$ 其否定“$\forall x\in\mathbb{R},\ x^2+2ax+4\geq0$”为真。
这是恒成立问题：$\Delta=(2a)^2-16\leq0\iff4a^2-16\leq0\iff-2\leq a\leq2$。
故 $a\in[-2,2]$。
\end{solution}

\begin{example}
证明：函数 $f(x)=x^3+x+1$ 有且只有一个零点。
\end{example}
\begin{solution}
\key{存在性}：$f(-1)=-1-1+1=-1<0$，$f(0)=1>0$，且 $f$ 是多项式（连续），
由零点存在定理，$(-1,0)$ 内至少有一个零点。
\key{唯一性（反证）}：假设有两个不相等的零点 $x_1,x_2$（设 $x_1<x_2$），则
\[
  x_1^3+x_1+1=0,\qquad x_2^3+x_2+1=0.
\]
两式相减得 $(x_1^3-x_2^3)+(x_1-x_2)=0$，即
\[
  (x_1-x_2)(x_1^2+x_1x_2+x_2^2+1)=0.
\]
因 $x_1\neq x_2$，需 $x_1^2+x_1x_2+x_2^2+1=0$。
但 $x_1^2+x_1x_2+x_2^2=\left(x_1+\frac{x_2}{2}\right)^2+\frac34x_2^2\geq0$，
故 $x_1^2+x_1x_2+x_2^2+1\geq1>0$，矛盾。故零点唯一。
\end{solution}

\begin{example}
用数学归纳法证明：$n^3+5n$ 能被 $6$ 整除（$n\in\mathbb{N}^*$）。
\end{example}
\begin{solution}
(1) 基础：$n=1$ 时 $1+5=6$，能被 $6$ 整除。
(2) 递推：假设 $n=k$ 时 $k^3+5k=6m$。则 $n=k+1$ 时
\begin{align*}
  (k+1)^3+5(k+1)&=k^3+3k^2+3k+1+5k+5\\
  &=(k^3+5k)+3k^2+3k+6\\
  &=6m+3k(k+1)+6.
\end{align*}
因为 $k(k+1)$ 是两个连续整数之积，必为偶数，故 $3k(k+1)$ 能被 $6$ 整除；又 $6$ 能被 $6$ 整除。
所以 $(k+1)^3+5(k+1)$ 能被 $6$ 整除。由(1)(2)，结论对一切 $n\in\mathbb{N}^*$ 成立。
\end{solution}

\begin{example}
用数学归纳法证明：对一切 $n\geq5$，$2^n>n^2$。
\end{example}
\begin{solution}
(1) 基础：$n=5$ 时 $2^5=32>25=5^2$，成立。
(2) 递推：假设 $n=k$（$k\geq5$）时 $2^k>k^2$。则 $n=k+1$ 时
\[
  2^{k+1}=2\cdot2^k>2k^2.
\]
要证 $2k^2>(k+1)^2=k^2+2k+1$，即 $k^2-2k-1>0$，即 $(k-1)^2>2$。
因 $k\geq5$，$(k-1)^2\geq16>2$，成立。故 $2^{k+1}>(k+1)^2$。
由数学归纳法，对一切 $n\geq5$，$2^n>n^2$（注意归纳起点平移到了 $5$）。
\end{solution}

\begin{example}
设命题 $p$：$\forall x\in\mathbb{R},\ x^2+ax+1>0$；命题 $q$：$\exists x\in\mathbb{R},\ x^2-2x+a=0$。
若 $p\lor q$ 为假，求 $a$ 的取值范围。
\end{example}
\begin{solution}
$p\lor q$ 为假 $\iff$ $p$ 假且 $q$ 假（德摩根律）。
$p$ 真 $\iff\Delta_1=a^2-4<0\iff-2<a<2$，故 $p$ 假 $\iff a\leq-2$ 或 $a\geq2$。
$q$ 真 $\iff\Delta_2=4-4a\geq0\iff a\leq1$，故 $q$ 假 $\iff a>1$。
取交：$(a\leq-2\ \text{或}\ a\geq2)\cap(a>1)$，得 $a\geq2$。
故 $a\in[2,+\infty)$。
\end{solution}

\begin{conccard}[三级]{反证法触发信号：唯一性、不存在、无理数、无限性}{结论含“只有一个／不存在／无理数／无穷多”的证明题}{把反证法当成“随便假设”；假设后没有推出矛盾就下结论}
反证法即证反例集 $A\cap\complement_U B=\varnothing$：假设反例存在，推出与公理、定理或临时假设矛盾的结果。
唯一性假设“至少两个”；不存在假设“存在”；无理数设既约分数 $\frac pq$；无限性设“只有有限个”。
矛盾必须明确写出，否则证明不完整。
\end{conccard}

\begin{conccard}[三级]{数学归纳法：归纳基础＋归纳递推}{证明对全体自然数（或 $n\geq n_0$）成立的命题}{只做递推不做基础；归纳假设没有真正用上；归纳起点误写为 $1$}
归纳公理：若 $1\in S$ 且 $n\in S\Rightarrow n+1\in S$，则 $S=\mathbb{N}^*$。
其根基是良序原理与最小反例原理：若存在反例，则存在最小反例，与递推矛盾。
使用时要检查两点：归纳基础是否验证；递推中是否用到了归纳假设。起点可按命题需要平移到 $n_0$。
\end{conccard}

\begin{summary}
本章把集合与逻辑锻造成一套“武器库”：集合相等公理 $A=B\iff(A\subseteq B)\land(B\subseteq A)$ 是审查一切同解变形的标准（$A\subsetneqq B$ 增根、$B\subsetneqq A$ 漏根）；德摩根律、逆否等价、量词否定与补集转化处理“至多、至少、不都、无解、恒成立、能成立”；反证法即证反例集 $A\cap\complement_U B=\varnothing$，由“唯一性、不存在、无理数、无限性”四类信号词触发；数学归纳法以皮亚诺公理与良序原理为根基，用最小反例原理保证“归纳基础＋归纳递推”两步缺一不可。解题时先翻译题干（集合／逻辑语言），再选引擎，最后审方向（防增漏根）。
\end{summary}

\end{explore}

\subsection{四视角等价地图：同一对象，多副面孔}

前四个模块其实都在做同一件事：给同一个数学对象换一副面孔。
代数、几何、集合、映射（对象本身）不是四门课，而是同一结构的四种\key{表示}；
只要两种表示之间存在\emph{保持结构的对应}（同构），命题就能互相翻译。
但「等价」不等于「一样难」：\textbf{同一问题在不同框架下的可解释性与计算量差别极大}，
所以解题的关键动作是\key{找一个同构映射，把问题搬到更顺的那一侧}。

\begin{center}
\footnotesize
\begin{tabular}{c|p{2.6cm}|p{3.2cm}|p{3.0cm}|p{3.4cm}}
\toprule
视角 & 看什么 & 擅长 & 短板 & 典型「搬运」 \\
\midrule
代数 & 符号、方程、恒等变形 & 机械化、可执行、可检验 & 缺直观，易迷失方向 & 几何 $\to$ 方程（解析几何）\\
几何 & 图形、变换、位置关系 & 直观、好找思路、便于猜答案 & 精度与严格性弱 & 方程 $\to$ 图形（数形结合）\\
集合 & 元素、包含、交并补、量词 & 严谨分类、盯住边界与反例 & 不给算法 & 命题 $\to$ 集合关系（充分必要）\\
映射／对象 & 对应关系、迭代、结构本身 & 统一语言、跨领域搬运结论 & 太抽象，需落地实例 & 矩阵 $=$ 线性映射；坐标 $=$ 同构 \\
\bottomrule
\end{tabular}
\end{center}

\begin{figure}[htbp]
\centering
\begin{tikzpicture}[>=Stealth,semithick,font=\footnotesize,
  box/.style={draw,rounded corners,inner sep=4pt,align=center,minimum width=2.1cm}]
  \node[box,draw=main] (A) at (-2.2,1.9) {代数\\方程};
  \node[box,draw=second] (G) at (2.2,1.9) {几何\\图形};
  \node[box,draw=third] (S) at (-2.2,-1.9) {集合\\逻辑};
  \node[box,draw=gray!60!black] (M) at (2.2,-1.9) {映射\\结构};
  \node[box,draw=main!70!black,minimum width=2.4cm] (C) at (0,0) {\textbf{同一个}\\数学对象};
  \draw[<->,main] (C)--(A) node[midway,above left,font=\scriptsize] {列方程};
  \draw[<->,second] (C)--(G) node[midway,above right,font=\scriptsize] {画出来};
  \draw[<->,third] (C)--(S) node[midway,below left,font=\scriptsize] {下定义};
  \draw[<->,gray!60!black] (C)--(M) node[midway,below right,font=\scriptsize] {建模型};
  \draw[->,gray!55,dashed] (A)--(G) node[midway,above,font=\scriptsize] {数形结合};
  \draw[->,gray!55,dashed] (S)--(M) node[midway,below,font=\scriptsize] {关系即映射};
\end{tikzpicture}
\caption{四视角等价地图：中心是对象本身，四个顶点是它的四种表示；
箭头都是「保持结构的映射」。解题时沿箭头换到顺手的一侧，算完再换回来}
\label{fig:four-view-map}
\end{figure}

\begin{figure}[htbp]
\centering
\begin{tikzpicture}[>=Stealth,semithick,font=\footnotesize,
  ch/.style={draw,rounded corners,inner sep=5pt,align=center,minimum width=2.0cm},
  lanename/.style={rotate=90,anchor=south,font=\scriptsize\bfseries,text=main!70!black}]
  % 三条主线 + 一条独立线
  \node[lanename] at (-1.1,6) {几何·结构线};
  \node[lanename] at (-1.1,3) {离散动力线};
  \node[lanename] at (-1.1,0.6) {连续分析线};
  % 第 0 层：语言底座
  \node[ch,draw=main,fill=main!8] (c0) at (1.6,7.6) {第 0 章\\集合·逻辑·证明};
  \node[ch,draw=gray!60!black,fill=gray!8] (c7) at (11.4,7.6) {第 7 章\\概率·统计};
  \draw[->,main] (c0)--(c7) node[midway,above,font=\scriptsize] {测度＝归一化计数};
  % 几何·结构线
  \node[ch,draw=second] (c1) at (1.6,6) {第 1 章\\函数};
  \node[ch,draw=second] (c2) at (5.0,6) {第 2 章\\三角·复数};
  \node[ch,draw=second] (c3) at (8.2,6) {第 3 章\\向量};
  \node[ch,draw=second] (c4) at (11.4,6) {第 4 章\\立几·圆锥曲线};
  \draw[->,second] (c1)--(c2);
  \draw[->,second] (c2)--(c3);
  \draw[->,second] (c3)--(c4);
  % 离散动力线
  \node[ch,draw=third] (c5) at (1.6,3) {第 5 章\\数列·递推};
  \node[ch,draw=third] (c9) at (11.4,3) {第 9 章\\不动点};
  \draw[->,third] (c5)--(c9) node[midway,above,font=\scriptsize] {递推求通项·收敛};
  % 连续分析线
  \node[ch,draw=main!60!black] (c6) at (5.0,0.6) {第 6 章\\导数·局部线性化};
  \node[ch,draw=main!60!black] (c8) at (8.2,0.6) {第 8 章\\齐次化};
  \draw[->,main!60!black] (c6)--(c9) node[midway,right,font=\scriptsize] {稳定性};
  \draw[->,main!60!black] (c8)--(c9) node[midway,right,font=\scriptsize] {射影};
  % 跨线接口
  \draw[->,gray!55,dashed] (c1)--(c5) node[midway,left,font=\scriptsize] {离散函数};
  \draw[->,gray!55,dashed] (c1)--(c6) node[midway,left,font=\scriptsize] {变化率};
  \draw[<->,gray!55,dashed] (c5)--(c6) node[midway,below,font=\scriptsize] {差分 $\leftrightarrow$ 导数};
  \draw[->,gray!55,dashed] (c4)--(c8) node[midway,right,font=\scriptsize] {条件最值};
  \draw[->,main!35] (c0)--(c1);
  \draw[->,gray!55,dashed] (c3)--(c9) node[midway,above right,font=\scriptsize] {特征方向};
\end{tikzpicture}
\caption{全书知识脉络图：实线是各自主线的推进，虚线是跨主线接口。
三条主线（几何·结构／离散动力／连续分析）最终在第 9 章不动点汇合；
第 8 章齐次化是射影观点的入口，第 0 章语言底座供全局调用}
\label{fig:book-map}
\end{figure}

\textbf{脉络怎么读}：几何·结构线从函数到圆锥曲线，解决「形的问题」；
离散动力线与连续分析线在 \key{差分 $\leftrightarrow$ 导数}处对接（同一个
「以直代曲／以比代差」思想），分别解决「列的问题」与「极限的问题」，
并共同汇入第 9 章的不动点理论；第 8 章齐次化把二元问题压到一元比例，
是通往射影／映射观点的接口。每一章章首的「多框架对读」卡，就是在这张
图上标明「你现在在哪条线、哪一侧」。

\textbf{怎么用它找思路}：拿到题先问「我现在站在哪个视角」——
卡在代数就画图（第 2、4 章的向量与圆锥曲线）；卡在几何就设坐标列方程
（第 4 章的仿射、第 8 章的齐次化）；被式子的次数与比值绕晕就升到射影／映射
观点（第 9 章的不动点与 Möbius）；讨论「有没有、有几个」就回到集合与量词
（本章的同解变形与反证法）。\textbf{同一结论往往能在四处各证一遍}，选最省力的那处书写。
此后各章章首统一附「多框架对读」卡，把这条暗线显性化。

\begin{inclasstest}
\begin{enumerate}
  \item 写出“$\exists x\in\mathbb{R},\ x^2-2x+1\leq0$”的否定，并判断其真假；\level{易}
  \item 解方程 $\sqrt{x+3}=x+1$，并说明验根的必要性；\level{中}
  \item 已知“存在 $x\in[1,3]$，使 $x^2-ax+2\leq0$”为真，求 $a$ 的取值范围；\level{中}
  \item 用反证法证明：若 $a,b\in\mathbb{Z}$ 且 $ab$ 能被 $3$ 整除，则 $a,b$ 中至少有一个能被 $3$ 整除。\level{中}
\end{enumerate}
\end{inclasstest}

\begin{afterclass}
\subsection*{基础巩固}

\begin{enumerate}
  \item 写出德摩根律与量词否定的四条公式：$\neg(p\land q)$、$\neg(p\lor q)$、$\neg(\forall x,P(x))$、$\neg(\exists x,P(x))$；
  \item 解方程 $x^2=2x$，并说明若两边同除以 $x$ 会漏掉哪个根。
\end{enumerate}

\subsection*{进阶提升}
\begin{enumerate}
  \item 已知“$\forall x\in\mathbb{R},\ x^2+2x+a>0$”为真，求 $a$；再写出它的否定命题为真时 $a$ 的范围；
  \item 用数学归纳法证明：$5^n-1$ 能被 $4$ 整除。
\end{enumerate}

\subsection*{拓展探究}
\begin{enumerate}
  \item 用反证法证明：$\sqrt3$ 是无理数。（提示：用到上面第 4 题的结论“若 $3\mid ab$，则 $3\mid a$ 或 $3\mid b$”。）
\end{enumerate}

\subsection*{教师版答案}
\begin{solution}
\textbf{基础 1：}$\neg(p\land q)\iff\neg p\lor\neg q$；$\neg(p\lor q)\iff\neg p\land\neg q$；
$\neg(\forall x,P(x))\iff\exists x,\neg P(x)$；$\neg(\exists x,P(x))\iff\forall x,\neg P(x)$。
\end{solution}
\begin{solution}
\textbf{基础 2：}移项分解 $x(x-2)=0$，得 $x=0$ 或 $x=2$，解集 $\{0,2\}$。
若两边同除以 $x$ 得 $x=2$，漏掉 $x=0$（同除使解集由 $\{0,2\}$ 缩小为 $\{2\}$）。
\end{solution}
\begin{solution}
\textbf{进阶 1：}$x^2+2x+a>0$ 恒成立 $\iff\Delta=4-4a<0\iff a>1$。
其否定“$\exists x\in\mathbb{R},\ x^2+2x+a\leq0$”为真 $\iff\Delta\geq0\iff a\leq1$。
\end{solution}
\begin{solution}
\textbf{进阶 2：}(1) $n=1$ 时 $5-1=4$，能被 $4$ 整除。
(2) 假设 $5^k-1=4m$，则 $5^{k+1}-1=5(5^k-1)+4=20m+4=4(5m+1)$，能被 $4$ 整除。
由数学归纳法，结论成立。
\end{solution}
\begin{solution}
\textbf{拓展：}假设 $\sqrt3$ 是有理数，写成既约分数 $\sqrt3=\frac pq$（$\gcd(p,q)=1$）。
平方得 $p^2=3q^2$，故 $3\mid p^2$；由提示（取 $a=b=p$）得 $3\mid p$，设 $p=3k$。
代入得 $9k^2=3q^2$，即 $q^2=3k^2$，故 $3\mid q$。于是 $p,q$ 均为 $3$ 的倍数，与既约矛盾。
故 $\sqrt3$ 是无理数。
\end{solution}
\reading{拓展阅读 00}{逻辑与证明的武器库}
\end{afterclass}

% 高观点导学案：函数的概念与基本初等函数
% 主线：关系 R⊆A×B（三档舞台）→ 方程/不等式＝关系的条件表达 → 两条限制轴（单值性→映射/函数；连续性与结构保持→基本初等函数）
% 全程三视角互译：代数有序数对集合 ↔ 几何坐标系点集 ↔ 映射箭头图

\section{函数的概念与基本初等函数（高观点结构版）}

\begin{methodbox}{多框架对读}
\small
\begin{center}
\begin{tabular}{c|p{2.9cm}|p{5.0cm}|p{3.1cm}}
\toprule
视角 & 本章看点 & 别的视角怎么说 & 哪侧最省力 \\
\midrule
代数 & 解析式、复合与解方程 & 有序数对集合的坐标运算 & 变形与恒等式 \\
几何 & 图象的平移／对称变换 & 坐标系中点集的移动 & 猜性质、定范围 \\
集合 & 定义域、值域、$f:A\to B$ & 取值范围即约束集合 & 审边界、找反例 \\
映射 & $y=f(x)$ 即对应法则 & 变换＝作用在图象上的映射 & 统一语言、整体搬迁 \\
\bottomrule
\end{tabular}
\end{center}
卡的使用：卡住时先判断自己站在哪一栏，再沿箭头换栏（总图见第 0 章图~\ref{fig:four-view-map}）。
\end{methodbox}

% ---------- 前置引入：舞台的诞生 ----------
\subsection{前置引入：舞台的诞生——集合自身的四级进化}

研究任何数学对象，现代数学靠两根支柱：\key{舞台}（集合自身的结构）与 \key{规则}（集合之间的映射）。本导学案主线一直在讲\key{规则}：$A\times B$ 中怎样挑点、怎样过滤出单值对应。但高中函数之所以拥有单调性、奇偶性、零点存在定理与导数，原因不在规则里，而在\key{舞台}上——定义域与值域选中的是哪种集合。

\begin{conccard}[一级]{集合装了多少「装备」，函数就有多少性质}{适用于所有「新定义」与概念辨析题}{只背性质条款，不问性质从何而来}
一般集合（例如 $\{\text{张三},\text{李四}\}$）只是一堆孤立元素：不能比大小，不能做加减，也没有「挨得很近」的概念。一旦把定义域与值域锁定在实数集 $\mathbb{R}$ 上，$\mathbb{R}$ 自带的三层结构便一次性注入函数。
\end{conccard}

\begin{itemize}
  \item \level{Lv.1} \key{原始集合}——一袋散装弹珠。只有「属于」与「元素个数」，无法比大小、不能运算：既谈不上单调性，也画不出坐标系中的图象。
  \item \level{Lv.2} \key{有序与代数集合}（$\mathbb{Z},\mathbb{Q}$）——排好队且能算账的刻度尺。\key{全序结构}（$x_1<x_2$ 可比大小）解锁\key{单调性与最值}；\key{代数运算与相反数}（$+,-,\times,\div$ 与 $-x$）解锁\key{奇偶性与周期性}。复数域 $\mathbb{C}$ 不能比大小，所以复变函数根本没有「单调递增」可言。
  \item \level{Lv.3} \key{稠密但有孔的集合}（$\mathbb{Q}$）——布满微观针眼的筛子。任意两点之间仍有无穷多个点，却漏掉了 $\sqrt{2},\pi,e$ 等无数无理数。致命后果见图~\ref{fig:qline}：$f(x)=x^2-2$ 从负值升到正值时，正好从 $\sqrt{2}$ 的孔洞中漏过去，\key{根本没有零点}。
  \item \level{Lv.4} \key{完备连续集合}（$\mathbb{R}$）——无缝焊接的实心金条。无理数填满全部孔洞，曲线无处可漏，\key{零点存在定理}（介值性）、极限以及「以直代曲」的导数才真正成为可能。
\end{itemize}

Lv.1 与 Lv.2 的差别不在元素本身，而在元素之间被装上了什么关系：前者只有「属于」与计数（图 \ref{fig:bag2axis} 左），后者多出了「谁大谁小」的全序与「怎么算账」的代数运算（图 \ref{fig:bag2axis} 右）。

\begin{figure}[htbp]
  \centering
  \begin{tikzpicture}[>=stealth,font=\small]
    % ---------- 左：Lv.1 一袋散装弹珠（无序点集） ----------
    \begin{scope}
      \draw[gray!60,very thick,fill=main!8] (0,0) ellipse (2.25cm and 1.6cm);
      \foreach \px/\py/\lab in {
        -1.45/0.62/-2, -0.95/-0.88/-1, -0.42/0.28/-\tfrac12,
        0.02/-0.55/0, 0.42/0.78/\tfrac12, 0.88/-0.18/1,
        1.35/0.5/\tfrac32, 1.12/-1.0/2, -1.78/-0.32/\tfrac52}
      {
        \fill[second] (\px,\py) circle (0.075);
        \node[gray!80,font=\scriptsize] at (\px+0.16,\py+0.16) {$\lab$};
      }
      \node at (0,-1.95) {\small 只有「属于」与元素个数};
      \node[second] at (0,-2.42) {\small 不能比大小，也不能加减};
    \end{scope}
    % ---------- 右：Lv.2 带刻度的数轴（有序集合 Q） ----------
    \begin{scope}[xshift=7.2cm]
      \draw[->,gray,very thick] (-2.7,0)--(3.1,0) node[right,gray] {$x$};
      \foreach \t in {-2,-1,0,1,2}
        \draw[gray,thick] (\t,0.09)--(\t,-0.09);
      \foreach \t in {-2,-1,0,1,2}
        \node[gray,font=\scriptsize] at (\t,-0.34) {$\t$};
      \foreach \t in {-2,-1,-0.5,0,0.5,1,1.5,2,2.5}
        \fill[main] (\t,0) circle (0.05);
      \node[third] at (0,1.15) {$(<)\quad +,\,-,\,\times,\,\div$};
      \node at (0,-1.35) {\small 排好队且能算账的刻度尺};
      \node[main] at (0,-1.82) {\small 全序解锁单调最值，运算解锁奇偶周期};
    \end{scope}
  \end{tikzpicture}
  \caption{从「属于」到「可比大小」：左为 Lv.1 原始集合——一袋散装弹珠，只有属于关系与元素个数，无法比大小、不能运算；右为 Lv.2 有序与代数集合——同一批元素排上带刻度的数轴 $\mathbb{Q}$，全序 $(<)$ 解锁单调性与最值，$+,-,\times,\div$ 与相反数解锁奇偶性与周期性}
  \label{fig:bag2axis}
\end{figure}

\begin{figure}[htbp]
  \centering
  \begin{tikzpicture}[x=0.92cm,y=0.92cm,>=stealth,font=\small]
    % ---------- 左：有理数轴（穿孔） ----------
    \begin{scope}
      \draw[->,white!55] (-0.7,0)--(2.9,0) node[right,white!70] {$x$};
      \foreach \x in {0.1,0.35,0.62,0.9,1.15,1.62,1.85,2.1,2.4,2.65}
        \draw[second,fill=white] (\x,0) circle (0.042);
      \draw[second,very thick,fill=white] (1.414,0) circle (0.075);
      \node[second] at (1.414,-0.42) {\small $\sqrt{2}$（洞）};
      \draw[main,thick,domain=0.55:2.0,samples=70,smooth] plot(\x,{\x*\x-2});
      \node[main] at (1.45,2.30) {$y=x^2-2$};
      \node[white!70] at (1.05,-2.02) {\small $\mathbb{Q}$：看似连续，实则满布针眼};
      \node[second] at (1.05,-2.56) {\small 穿过空洞：无交点};
    \end{scope}
    % ---------- 右：实数轴（实心） ----------
    \begin{scope}[xshift=7.6cm]
      \draw[->,third,very thick] (-0.7,0)--(2.9,0) node[right,third] {$x$};
      \foreach \x in {0.1,0.35,0.62,0.9,1.15,1.62,1.85,2.1,2.4,2.65}
        \fill[third!70] (\x,0) circle (0.032);
      \draw[main,thick,domain=0.55:2.0,samples=70,smooth] plot(\x,{\x*\x-2});
      \fill[second] (1.414,0) circle (0.075);
      \node[second] at (1.414,-0.42) {$(\sqrt{2},0)$};
      \node[third] at (1.05,-2.02) {\small $\mathbb{R}$：无缝连续};
      \node[second] at (1.05,-2.56) {\small 洞被填满：零点被迫现身};
    \end{scope}
  \end{tikzpicture}
  \caption{同一函数 $f(x)=x^2-2$ 在两种「舞台」上的命运：左，定义域取有理数集 $\mathbb{Q}$，曲线从 $\sqrt{2}$ 的空心孔洞中漏过，变号却无零点；右，定义域取实数集 $\mathbb{R}$，孔洞被无理数填满，零点 $(\sqrt{2},0)$ 必然出现}
  \label{fig:qline}
\end{figure}

\subsubsection*{开场三问}
\begin{enumerate}
  \item 复数域 $\mathbb{C}$ 上能否定义「单调递增」？为什么？（提示：$\mathbb{C}$ 上有没有「小于」关系）\blank
  \item 在有理数集 $\mathbb{Q}$ 上考察 $f(x)=x^2-3$：它满足 $f(1)<0<f(2)$，却没有零点。请用「孔洞」的语言解释。\blank
  \item 狄利克雷函数 $D(x)=\begin{cases}1,&x\in\mathbb{Q}\\0,&x\in\mathbb{R}\setminus\mathbb{Q}\end{cases}$ 以任一非零有理数为周期，却没有最小正周期。它击穿了「周期函数一定有最小正周期」这句直觉——你能再举一句被它击穿的说法吗？\blank
\end{enumerate}

\begin{method}
「看集合选武器」的反射：拿到一道涉及函数的题，先扫一眼定义域是哪类集合，再决定调用哪套工具——
$\mathbb{N}^{*}$（离散有序）$\Rightarrow$ 差分、数学归纳法、不动点；
$[a,b]\subseteq\mathbb{R}$（连续完备）$\Rightarrow$ 单调性、零点定理、导数局部线性化；
$\mathbb{Q}$ 与 $\mathbb{R}\setminus\mathbb{Q}$ 混用 $\Rightarrow$ 立刻警觉稠密性与孔洞性，准备构造反例。
\end{method}

\begin{commonerror}{把「有界」当成「能取到最值」}
$f(x)>a$ 在 $(0,1)$ 上恒成立，即使 $f$ 在该区间单调递增且 $\lim\limits_{x\to0^{+}}f(x)=2$，也只能推出 $a\le 2$ 而不是 $a<2$：端点取不到时，边界值扮演的是\key{下确界}而非最小值。这类「端点开合」丢分，考的正是实数的完备性。
\end{commonerror}

四级进化完成后，$\mathbb{R}$ 与自身作笛卡尔积，舞台张成连续平面（图 \ref{fig:r2plane}）——后续的二维区域与函数图象都落在这块舞台上。

\begin{figure}[htbp]
  \centering
  \begin{tikzpicture}[>=stealth,font=\small,scale=1.02]
    \draw[third!35,dashed,thin] (-2.5,-2) grid (2.5,2);
    \draw[->,gray,very thick] (-2.75,0)--(2.75,0) node[right,gray] {$x$};
    \draw[->,gray,very thick] (0,-2.25)--(0,2.25) node[above,gray] {$y$};
    \foreach \t in {-2,-1,1,2}{
      \node[gray,font=\scriptsize] at (\t,-0.18) {$\t$};
      \node[gray,font=\scriptsize] at (-0.2,\t) {$\t$};
    }
    \fill[main!12] plot[smooth cycle] coordinates {(0.1,-1.1) (1.5,-0.7) (2.05,0.35) (1.25,1.35) (-0.15,1.55) (-1.35,0.65) (-1.55,-0.6)};
    \draw[main!70!black,very thick] plot[smooth cycle] coordinates {(0.1,-1.1) (1.5,-0.7) (2.05,0.35) (1.25,1.35) (-0.15,1.55) (-1.35,0.65) (-1.55,-0.6)};
    \node[main!80!black] at (0.45,0.35) {二维区域（「面」）};
    \fill[second] (1.414,0) circle (0.06);
    \node[second,below right,font=\footnotesize] at (1.414,0) {$(\sqrt2,0)$};
    \node[third,font=\large] at (2.05,1.95) {$\mathbb{R}^2$};
  \end{tikzpicture}
  \caption{舞台就绪：完备的 $\mathbb{R}$ 与自身作笛卡尔积张成连续平面 $\mathbb{R}^2$（坐标轴、网格与一块二维区域的定场）。零维的点已被连成无缝的线，线再铺成二维的面——不等式解集、函数图象等后续对象都以这块平面为舞台}
  \label{fig:r2plane}
\end{figure}

\subsection{实数的三种等价面孔：分割、基本列与确界}

前面用「针眼筛子」与「实心金条」画出 $\mathbb{Q}$ 与 $\mathbb{R}$ 的差别。本节把它变成可以动手验证的数学：孔洞究竟在哪里？怎么补？补完的实数，为什么恰好就是数轴上的点？
（数系为何还要从 $\mathbb{R}$ 继续走向 $\mathbb{C}$，见第 2 章「数系扩张与虚数单位 $i$ 的几何诞生」。）

\key{探究一：稠密也会漏——有理数的上确界会失踪}
设
\[
  S=\{x\in\mathbb{Q}\mid x^2<2\}.
\]
$S\neq\varnothing$（$1\in S$），也有上界（$2$ 就是一个上界）。若在有理数范围内「有上界必有最小上界」，$S$ 应当存在最小上界。事实是：不存在。

\begin{derivation}
任取正有理数 $r$，构造它的「平分修正」
\[
  t=\frac{r+\dfrac{2}{r}}{2}\in\mathbb{Q},\qquad
  t^2-2=\frac{(r^2-2)^2}{4r^2}\geq0,\qquad
  t-r=-\frac{r^2-2}{2r}.
\]
\begin{itemize}
  \item 若 $r^2>2$（此时 $r$ 是 $S$ 的上界）：仍有 $t^2>2$，即 $t$ 也是上界；但 $t<r$。于是任何有理上界都能被更小的有理上界取代，「最小上界」无从谈起。
  \item 若 $r^2<2$（此时 $r\in S$）：$t>r$ 且 $t^2>2$，即 $t$ 是一个比 $r$ 更大的上界。
\end{itemize}
两种情形合起来：$\mathbb{Q}$ 中既没有 $S$ 的最小上界，也没有恰好落在 $\sqrt2$ 处的点。孔洞不是文学比喻，而是确界原理在 $\mathbb{Q}$ 中失效。
\end{derivation}

\key{探究二：补洞的三条等价路线}
\begin{itemize}
  \item \key{戴德金分割}：把 $\mathbb{Q}$ 切成上下两堆 $(A\mid B)$，要求 $A$ 无最大元；每个分割定义一个实数。上面由 $S$ 诱导的分割 $\bigl(\{x\in\mathbb{Q}:x^2<2\}\mid\mathbb{Q}\setminus\{x\in\mathbb{Q}:x^2<2\}\bigr)$ 恰好补上缺失的 $\sqrt2$。
  \item \key{柯西基本列}：把「彼此无限靠近」的有理数列（如 $3,\,3.1,\,3.14,\,\dots$）视为一个实数；两列若差趋于 $0$ 就视为同一个实数。$\pi$ 是这样被「逼」出来的。
  \item \key{公理刻画}（只描述、不构造）：确界原理、单调有界定理、闭区间套定理、柯西收敛准则四者等价，说的是同一句话——$\mathbb{R}$ 没有孔洞。高中阶段的极限、零点存在定理、介值性全部建立在这条公理上。
\end{itemize}

\begin{center}
\footnotesize
\setlength{\tabcolsep}{3pt}
\renewcommand{\arraystretch}{1.4}
\begin{tabular}{p{2.4cm}p{3.6cm}p{3.0cm}p{3.8cm}}
\toprule
视角 & 怎么补洞 & 典型产物 & 高中落点 \\
\midrule
戴德金分割 & 用「上下两堆」的卡口位置定义新数 & $\sqrt2$ 这类不可公度长度 & 数轴上的点与实数一一对应 \\
柯西基本列 & 用「越来越近的有理列」定义新数 & $\pi,\ e$ 这类极限 & 数列极限、级数求和的合法性 \\
公理刻画 & 直接声明「有上界就有上确界」 & 完备性本身 & 零点定理、最值定理的理论根基 \\
\bottomrule
\end{tabular}
\end{center}

\key{探究三：数轴上的点恰好是实数}
稠密 $\neq$ 连续：$\mathbb{Q}$ 已经稠密（任意两点之间还有有理数），却仍漏点；只有补上完备性，才得到「每一条几何长度都对应唯一实数」的连续数轴。
十进制表示还带来两个容易含糊的点：
\begin{itemize}
  \item $0.999\dots=1$——无限小数本身要靠完备性才能定义清楚；
  \item 「无理数 $=$ 无限不循环小数」只是无理数的表征，不能当作定义：要说明无限小数的全体恰好铺满完备数轴，绕不开分割或基本列。
\end{itemize}
这两条留作课后思考（见课后作业拓展探究）。

\begin{conccard}[一级]{实数的三种等价面孔：分割、基本列与确界}{判断一个集合能否承载函数的单调性、零点定理与极限}{把「实数 $=$ 无限小数」当定义；把「稠密」误当「连续」}
实数的三种定义（戴德金分割、柯西基本列、确界公理）彼此等价，共同特征是「没有孔洞」。
$\mathbb{Q}$ 的致命伤不在稀疏而在漏点：$S=\{x\in\mathbb{Q}:x^2<2\}$ 有上界却无最小上界，任何有理数都填不进 $\sqrt2$ 的卡口。
高中阶段不必背构造细节，但必须知道：极限、零点定理、最值定理的合法性来自完备性，而不只是来自「能比大小」。
\end{conccard}

\begin{commonerror}{把「稠密」当成「连续」}
常见错误：认为「任意两个有理数之间还有有理数，所以有理数轴已经连续」。稠密只保证「找不到相邻两点」，不保证「不漏点」：$\sqrt2,\pi,e$ 都不在 $\mathbb{Q}$ 里，曲线 $y=x^2-2$ 照样从孔洞中穿过。
判断标准：能否说「有上界的非空集合必有最小上界」。$\mathbb{Q}$ 不能，$\mathbb{R}$ 能。
\end{commonerror}

\begin{method}
「看集合选武器」再补一条反射：遇到「变号却无零点／有界却无最值／递增却撞不到天花板」，先怀疑定义域是 $\mathbb{Q}$ 而非 $\mathbb{R}$；遇到「长度、面积、开方、$\pi$」，先确认舞台是 $\mathbb{R}$。
\end{method}

% ---------- 以上为背景铺垫，以下进入主线 ----------

初中我们说"$y$ 随 $x$ 的变化而变化"，高中要说得更精确：函数是一种特殊的对应，它的图象是平面上的点集。本导学案用一条主线重新组织"函数"概念——
\key{关系}（$R\subseteq A\times B$ 的三档舞台：离散$\times$离散的孤立点、离散$\times$连续的竖直线段族、连续$\times$连续的二维区域）$\xrightarrow{\text{方程／不等式＝条件的表达法}}$ \key{两条限制轴}（轴一·单值性 $\to$ 映射与函数；轴二·连续性与结构保持 $\to$ 基本初等函数），并在每一步用"\key{代数有序数对集合 $\leftrightarrow$ 几何坐标系点集 $\leftrightarrow$ 映射箭头图}"三种语言互相翻译。

\begin{learnobjective}
\begin{itemize}
  \item 从\key{点集视角}理解二元关系与函数：把对应看成 $\mathbb{R}^2$ 平面上的点集，说清"哪些点被选中"；
  \item 理解\key{单值性}（垂线检验）如何把任意点集过滤成函数图象：每条竖直线 $x=a$ 与点集至多一个交点；
  \item 掌握\key{三视角互译}：代数有序数对集合 $\leftrightarrow$ 几何坐标系点集 $\leftrightarrow$ 映射箭头图；
  \item 理解\key{两条限制轴}：轴一「单值性」回答它是不是函数，轴二「连续性／结构保持」回答它好不好算，二者互不替代；
  \item 说明\key{方程与不等式是关系的条件表达法}，并在大一统框架下说出二元方程、函数、一元方程（不等式）的解分别是几维对象。
\end{itemize}
\end{learnobjective}

\begin{preview}
\subsection*{知识回顾}
\begin{itemize}
  \item 笛卡尔积：$A\times B=\{(a,b)\mid a\in A,\ b\in B\}$。若 $A=\{1,2\}$，$B=\{3\}$，则 $A\times B=$ \blank；
  \item 二元关系：$A\times B$ 的\key{任意子集} $R$ 都叫 $A$ 到 $B$ 的二元关系。若 $R=\{(1,3),(2,3)\}$，则 $R$ 中元素个数是 \blank；
  \item 不等式解集：$x^2-1\leq0$ 的解集是 \blank，它表示数轴上的 \blank；
  \item 平面直角坐标系中，点 $(x,y)$ 与有序数对 $(x,y)$ 是 \blank（一一对应／没有关系）。
\end{itemize}

\subsection*{课前速读}
阅读「函数的概念」一节，完成下列问题：
\begin{enumerate}
  \item 函数定义中的两个关键词是"\key{任意}"和"\key{唯一}"；
  \item 函数的三要素是 \blank、\blank、\blank；
  \item 下列对象中哪些是函数：(1) 圆的面积 $S$ 与半径 $r$；(2) 一个人的年龄与身高；(3) 平方对应 $x\mapsto x^2$。（答：\blank）
\end{enumerate}
\end{preview}

\begin{explore}

\subsection{有限离散关系：从 $A\times B$ 出发}

\key{探究一：把所有可能的配对都列出来}
设 $A=\{1,2,3\}$，$B=\{2,4,6\}$。
\begin{itemize}
  \item $A\times B$ 共有 $3\times3=9$ 个有序数对：
        $(1,2),(1,4),(1,6),(2,2),(2,4),(2,6),(3,2),(3,4),(3,6)$。
  \item 把它看成\key{网格点集}：在平面直角坐标系中描出这 $9$ 个点（图 \ref{fig:grid9}）。
  \item $A\times B$ 的任意子集 $R$ 都叫 $A$ 到 $B$ 的一个\key{二元关系}。例如 $R=\{(1,2),(2,4),(3,6)\}$ 只是"挑出其中几个配对"。
\end{itemize}

\begin{figure}[htbp]
  \centering
  \begin{tikzpicture}
    \begin{axis}[
      width=0.55\textwidth, height=6cm,
      xmin=0.5, xmax=3.5, ymin=1, ymax=7,
      xtick={1,2,3}, ytick={2,4,6},
      xlabel=$x$, ylabel=$y$,
      grid=major, grid style={dashed,gray!35},
      tick label style={font=\footnotesize},
      title={\small $A\times B$ 的 9 个有序数对（网格点）},
      title style={font=\small},
    ]
      \addplot[only marks, mark=*, mark size=2pt, main!70!black]
        coordinates {(1,2)(1,4)(1,6)(2,2)(2,4)(2,6)(3,2)(3,4)(3,6)};
    \end{axis}
  \end{tikzpicture}
  \caption{$A\times B$ 的点集视角：平面直角坐标系中的 $3\times3$ 网格点集。竖直方向看，每一列对应同一个 $x\in A$，各含 $3$ 个点（一对多，尚未通过垂线检验）；代数清单见正文，映射箭头图视角见图 \ref{fig:arrow}}
  \label{fig:grid9}
\end{figure}

\key{同一母体，三档舞台：}$R\subseteq A\times B$ 长什么样，首先取决于 $A,B$ 是哪种集合（即前文「舞台四级进化」装了多少装备）。三档并列如下：
\begin{itemize}
  \item \key{离散$\times$离散}（$A=\{1,2,3\},\ B=\{2,4,6\}$）：$R$ 只能是有限个孤立点——图 \ref{fig:grid9} 的 $9$ 点网格；
  \item \key{离散$\times$连续}（$A=\{1,2,3\},\ B=[0,4]$）：每个 $x$ 上方可以出现一整段竖直线段（点集是「三根竖条」，图 \ref{fig:three-tiers} 中），也可以只取有限个点——具体形态仍由所加的条件决定；
  \item \key{连续$\times$连续}（$A=B=[0,3]$）：$R$ 可以是一整块二维区域——图 \ref{fig:three-tiers} 右的矩形。
\end{itemize}
三档共用同一套三语言：代数写 $R=\{(x,y)\mid\text{条件}\}$，几何描点，箭头图看方向。\textbf{舞台决定点能密到什么程度，条件决定点被约束成什么形状}；二者缺一不可。特别地，$A\times B$（$9$ 点网格）本身也是一个二元关系——「全取」的关系；它代表\textbf{两个变元完全自由}的起点，而不是终点，后面的全部内容都在给这份自由逐级上锁。

\begin{figure}[htbp]
  \centering
  \begin{minipage}{0.31\textwidth}
    \centering
    \begin{tikzpicture}[x=0.55cm,y=0.55cm,>=stealth,font=\scriptsize]
      \draw[->,gray] (0.5,0)--(3.6,0) node[right] {$x$};
      \draw[->,gray] (0,-0.2)--(0,4.6) node[above] {$y$};
      \foreach \x in {1,2,3}{ \foreach \y in {2,4,6}{ \fill[main!70!black] (\x,\y) circle (1.6pt); } }
      \node at (2,-0.8) {离散$\times$离散：$9$ 个孤立点};
    \end{tikzpicture}
  \end{minipage}\hfill
  \begin{minipage}{0.31\textwidth}
    \centering
    \begin{tikzpicture}[x=0.55cm,y=0.55cm,>=stealth,font=\scriptsize]
      \draw[->,gray] (0.5,0)--(3.6,0) node[right] {$x$};
      \draw[->,gray] (0,-0.2)--(0,4.6) node[above] {$y$};
      \foreach \x/\yl/\yu in {1/0.2/1.6, 2/0.6/3.4, 3/0/2.2}{
        \draw[second,very thick] (\x,\yl)--(\x,\yu);
        \fill[second] (\x,\yl) circle (1.6pt);
        \fill[second] (\x,\yu) circle (1.6pt);
      }
      \node at (2,-0.8) {离散$\times$连续：三根竖条};
    \end{tikzpicture}
  \end{minipage}\hfill
  \begin{minipage}{0.31\textwidth}
    \centering
    \begin{tikzpicture}[x=0.55cm,y=0.55cm,>=stealth,font=\scriptsize]
      \draw[->,gray] (0.5,0)--(3.6,0) node[right] {$x$};
      \draw[->,gray] (0,-0.2)--(0,4.6) node[above] {$y$};
      \fill[main!12] (0,0) rectangle (3,4);
      \draw[main!70!black,thick] (0,0) rectangle (3,4);
      \node at (2,-0.8) {连续$\times$连续：一整块区域};
    \end{tikzpicture}
  \end{minipage}
  \caption{关系 $R\subseteq A\times B$ 的三档形态：同一母体，因 $A,B$ 的装备不同，分别长成孤立点云、竖直线段族、二维区域。三档都可以同时用「有序数对集合／坐标点集／箭头图」三种语言描述}
  \label{fig:three-tiers}
\end{figure}

\key{追问：}把这三档点集都当作「从 $x$ 到 $y$ 的对应」，它们各有什么缺陷，迫使我们继续加限制条件？
（提示：$9$ 点网格中取 $x=1$，$y$ 可以是 $2,4,6$ 中任意一个——对应不唯一；竖条与矩形中一个 $x$ 拖着无穷多个 $y$；而且离散点集没有「连续变化」的资格，圆滑曲线、连续性都无从谈起。）

\begin{example}
写出 $A=\{1,2,3\}$、$B=\{2,4,6\}$ 时 $A\times B$ 的全部 $9$ 个有序数对并在网格中描点；再任取两个不同的二元关系 $R_1,R_2\subseteq A\times B$（各含 $2$ 个元素）。
\end{example}
\begin{solution}
$9$ 个有序数对为 $(1,2),(1,4),(1,6),(2,2),(2,4),(2,6),(3,2),(3,4),(3,6)$，描点见图 \ref{fig:grid9}。
任取 $R_1=\{(1,2),(2,6)\}$，$R_2=\{(1,4),(3,2)\}$ 即可（答案不唯一）。它们都是 $A\times B$ 的子集，所以都是二元关系。
\end{solution}

\key{探究二：单值性——每列至多一个点}
对比两个关系：
\[
  R_1=\{(1,2),(1,4),(2,6)\},\qquad R_2=\{(1,2),(2,4),(3,6)\}.
\]
用三视角分别表示：
\begin{itemize}
  \item \key{代数}（有序数对集合）：如上，$R_1$ 中 $x=1$ 出现了两次；
  \item \key{几何}（点集）：在网格中描点（图 \ref{fig:r1r2}）。$R_1$ 在竖直线 $x=1$ 上有两个点 $(1,2),(1,4)$；$R_2$ 在每条竖直线 $x=1,2,3$ 上各有一个点；
  \item \key{映射}（箭头图）：图 \ref{fig:arrow}。$R_1$ 中 $1\to2$ 与 $1\to4$ 同时出现（\key{一对多}）；$R_2$ 中 $1\to2$、$2\to4$、$3\to6$（一对一）。
\end{itemize}

\begin{figure}[htbp]
  \centering
  \begin{minipage}{0.48\textwidth}
    \centering
    \begin{tikzpicture}
      \begin{axis}[
        width=\textwidth, height=5cm,
        xmin=0.5, xmax=3.5, ymin=1, ymax=7,
        xtick={1,2,3}, ytick={2,4,6},
        grid=major, grid style={dashed,gray!35},
        tick label style={font=\footnotesize},
        title={\small $R_1$：$x=1$ 处有两个点（一对多）},
        title style={font=\small},
      ]
        \addplot[only marks, mark=*, mark size=2pt, second]
          coordinates {(1,2)(1,4)(2,6)};
      \end{axis}
    \end{tikzpicture}
  \end{minipage}\hfill
  \begin{minipage}{0.48\textwidth}
    \centering
    \begin{tikzpicture}
      \begin{axis}[
        width=\textwidth, height=5cm,
        xmin=0.5, xmax=3.5, ymin=1, ymax=7,
        xtick={1,2,3}, ytick={2,4,6},
        grid=major, grid style={dashed,gray!35},
        tick label style={font=\footnotesize},
        title={\small $R_2$：每列恰有一个点（一对一）},
        title style={font=\small},
      ]
        \addplot[only marks, mark=*, mark size=2pt, main]
          coordinates {(1,2)(2,4)(3,6)};
      \end{axis}
    \end{tikzpicture}
  \end{minipage}
  \caption{$R_1$ 与 $R_2$ 的点集对比：竖直线 $x=1$ 与 $R_1$ 有两个交点，与 $R_2$ 只有一个交点}
  \label{fig:r1r2}
\end{figure}

\begin{figure}[htbp]
  \centering
  \begin{minipage}{0.48\textwidth}
    \centering
    \begin{tikzpicture}[>=Stealth, semithick]
      \node (a1) at (0,2.2) {$1$}; \node (a2) at (0,1.1) {$2$}; \node (a3) at (0,0) {$3$};
      \node (b2) at (3,2.2) {$2$}; \node (b4) at (3,1.1) {$4$}; \node (b6) at (3,0) {$6$};
      \draw[->] (a1)--(b2); \draw[->] (a1)--(b4); \draw[->] (a2)--(b6);
      \node[font=\footnotesize] at (0,2.8) {$A$}; \node[font=\footnotesize] at (3,2.8) {$B$};
    \end{tikzpicture}
    \\[2pt]{\footnotesize $R_1$：$1$ 同时指向 $2$ 和 $4$}
  \end{minipage}\hfill
  \begin{minipage}{0.48\textwidth}
    \centering
    \begin{tikzpicture}[>=Stealth, semithick]
      \node (a1) at (0,2.2) {$1$}; \node (a2) at (0,1.1) {$2$}; \node (a3) at (0,0) {$3$};
      \node (b2) at (3,2.2) {$2$}; \node (b4) at (3,1.1) {$4$}; \node (b6) at (3,0) {$6$};
      \draw[->] (a1)--(b2); \draw[->] (a2)--(b4); \draw[->] (a3)--(b6);
      \node[font=\footnotesize] at (0,2.8) {$A$}; \node[font=\footnotesize] at (3,2.8) {$B$};
    \end{tikzpicture}
    \\[2pt]{\footnotesize $R_2$：每个 $x$ 只指向一个 $y$}
  \end{minipage}
  \caption{映射箭头图：$R_1$ 一对多（不是函数），$R_2$ 一对一（是函数）}
  \label{fig:arrow}
\end{figure}

\key{判断规则（垂线检验雏形）：}点集 $P\subseteq\mathbb{R}^2$ 能作为某个函数的图象，当且仅当\key{每条竖直线 $x=a$ 与 $P$ 至多有一个交点}（"每列至多 $1$ 个点"）。据此 $R_2$ 是函数，$R_1$ 不是函数。

为什么 $R_2$ 是函数：对 $A$ 中每个 $x$，都有唯一 $y$ 与之对应；而 $R_1$ 中 $x=1$ 同时对应 $2$ 和 $4$，违反唯一性。

\key{把对应「翻译」成有序数对（图象的第三种身份）：}若对应关系为 $f$，把每个 $x$ 与它的像 $y=f(x)$ 并排写成 $(x,y)=(x,f(x))$。于是「哪些数对被选中」就是集合
\[
  \Gamma_f=\{(x,f(x))\mid x\in A\}.
\]
它读作映射，就是箭头图；读作有序数对集合，就是 $\Gamma_f$；坐标化到平面直角坐标系，就是函数的\key{图象}。三种说法在此合流——后文说「图象」，指的就是这个数对集合，不是额外多出来的对象。

\key{追问：}如果允许"一对多"，还能叫函数吗？"一对多"的对象在数学里是什么？
（提示：它仍是\key{关系}。例如二元方程 $x^2+y^2=4$ 的解集就是一对多的点集——这为探究三埋下伏笔。）

\begin{example}
用三视角（代数、几何、箭头图）表示关系 $R=\{(1,2),(2,4),(3,6)\}$，并判断它是否为函数。
\end{example}
\begin{solution}
代数：$R=\{(1,2),(2,4),(3,6)\}$。
几何：描点 $(1,2),(2,4),(3,6)$，三点在同一直线 $y=2x$ 上，且每条竖直线 $x=1,2,3$ 上各有一个点。
箭头图：$1\to2$，$2\to4$，$3\to6$。
是函数：每个 $x\in A$ 都有唯一 $y$ 与之对应（事实上 $y=2x$）。
\end{solution}

\begin{problem}
在 $A=\{1,2,3\}$、$B=\{2,4,6\}$ 的网格中，任取一个二元关系 $R$，使它不是函数，并说明理由；再任取一个函数关系，写出它的对应法则。
\end{problem}
\begin{solution}
例：$R=\{(1,2),(1,6),(2,4)\}$。$x=1$ 同时对应 $2$ 和 $6$，违反单值性，不是函数（竖直线 $x=1$ 与点集有两个交点）。
例：$R=\{(1,4),(2,2),(3,6)\}$ 是函数（每列至多一个点），对应法则可写为 $f(1)=4$，$f(2)=2$，$f(3)=6$。
\end{solution}

\key{从「完全自由」到「规则登场」：一条收紧链}

回到 $A=\{1,2,3\}$、$B=\{2,4,6\}$：$A\times B$ 的 $9$ 个配对是\emph{两个变元完全自由}时的全部候选。逐步加限制，自由度这样消失：
\begin{itemize}
  \item \key{二元关系}：对 $9$ 个配对逐个决定「取／不取」，共 $2^9=512$ 种——仍几乎无规律；
  \item \key{单值性}：每列至多取一个。每列有「不取，或取 $2,4,6$ 之一」共 $4$ 种选择，合计 $4^3=64$ 种；若要求 $A$ 中每个 $x$ 都有对应，则只剩 $3^3=27$ 种；
  \item \key{统一法则}：要求用同一条公式（结构保持等式）表达所有对应——$27$ 种里只剩极少数，例如 $(1,2),(2,4),(3,6)$ 可被 $y=2x$ 一网打尽。
\end{itemize}
函数不是外加的定义，而是这条收紧链上的一环：\textbf{自由消失多少，规则就登场多少}；基本初等函数位于链条顶端——规则收得最紧、结构保存得最好。

\subsection{连续二维区域：不等式解集的几何本质}

\key{先约定：方程与不等式是「从 $A\times B$ 里挑子集」的条件句}

含变元的等式或不等式，读法统一为「满足条件的有序数对全体」：
\begin{itemize}
  \item 一元方程 $x^2=4$：只有一个变元，解集是数轴上的点 $\{-2,2\}$；
  \item 二元方程 $x^2+y^2=4$：两个变元，解集是平面曲线（圆周）；
  \item 二元不等式 $x^2+y^2\leq4$：同一条件放松为「$\leq$」，解集从曲线「鼓成」闭圆盘。
\end{itemize}
所以方程、不等式不是与函数、映射并列的另一条线，而是\key{关系的具体表达法}：变元个数决定几何维度，等号／不等号决定边界还是区域。下面用这把尺子重新看连续点集。

\key{探究三：把 $9$ 个点换成无穷多个点}
把取值范围从离散集合推广到连续区间，考察两个平面点集：
\[
  R_3=\{(x,y)\mid x^2+y^2\leq4\},\qquad
  R_4=\{(x,y)\mid x+y\leq3,\ x\geq0,\ y\geq0\}.
\]
\begin{itemize}
  \item $R_3$ 是圆心在原点、半径为 $2$ 的\key{闭圆盘}（一块二维区域），其边界恰是二元方程 $x^2+y^2=4$，内部由不等式 $x^2+y^2\leq4$ 给出——「方程给边界、不等式给区域」的第一个实例；
  \item $R_4$ 是由 $x$ 轴、$y$ 轴与直线 $x+y=3$ 围成的\key{三角形区域}（图 \ref{fig:region}）。
  \item 它们都是 $\mathbb{R}^2$ 中"一块黑乎乎的区域"，而不是孤立的点。
\end{itemize}

关键观察：取 $x=1$，在 $R_3$ 中 $y\in[-\sqrt3,\sqrt3]$，有无穷多个取值；在 $R_4$ 中 $y\in[0,2]$，也有无穷多个取值。一个 $x$ 对应无穷多个 $y$——\key{失控}，无法由 $x$ 唯一预测 $y$。

\key{追问：}二元不等式（组）的解集在几何上是什么？（答：$\mathbb{R}^2$ 中的二维区域。）这正是未来\key{线性规划可行域}的伏笔：满足一组约束的 $(x,y)$ 全体就是平面区域，而最优解往往出现在区域的顶点或边界上。

\begin{figure}[htbp]
  \centering
  \begin{minipage}{0.48\textwidth}
    \centering
    \begin{tikzpicture}[scale=1.1]
      \draw[->] (-2.4,0)--(2.4,0) node[right]{$x$};
      \draw[->] (0,-2.4)--(0,2.4) node[above]{$y$};
      \fill[main!15] (0,0) circle (2cm);
      \draw[main!70!black, thick] (0,0) circle (2cm);
      \draw[dashed, second, thick] (1,-1.732)--(1,1.732)
        node[above right, font=\footnotesize]{$x=1:y\in[-\sqrt3,\sqrt3]$};
      \fill[second] (1,1.732) circle (1.5pt);
      \fill[second] (1,-1.732) circle (1.5pt);
    \end{tikzpicture}
    \\[2pt]{\footnotesize $R_3$：闭圆盘，竖直线 $x=1$ 与它交于无穷多点}
  \end{minipage}\hfill
  \begin{minipage}{0.48\textwidth}
    \centering
    \begin{tikzpicture}[scale=1.1]
      \draw[->] (-0.5,0)--(3.6,0) node[right]{$x$};
      \draw[->] (0,-0.5)--(0,3.6) node[above]{$y$};
      \fill[main!15] (0,0)--(3,0)--(0,3)--cycle;
      \draw[main!70!black, thick] (0,0)--(3,0)--(0,3)--cycle;
      \draw[dashed, second, thick] (1,0)--(1,2)
        node[right, font=\footnotesize]{$x=1:y\in[0,2]$};
    \end{tikzpicture}
    \\[2pt]{\footnotesize $R_4$：三角形区域，竖直线 $x=1$ 与它交于无穷多点}
  \end{minipage}
  \caption{连续二维区域：二元不等式（组）的解集是 $\mathbb{R}^2$ 中的"面"，不是曲线}
  \label{fig:region}
\end{figure}

\begin{example}
在坐标系中画出 $R_4=\{(x,y)\mid x+y\leq3,\ x\geq0,\ y\geq0\}$，并求当 $x=1$ 时 $y$ 的取值范围。
\end{example}
\begin{solution}
画出直线 $x+y=3$ 与两条坐标轴，取左下方的三角形内部（含边界），见图 \ref{fig:region} 右图。
当 $x=1$ 时，条件化为 $1+y\leq3$ 且 $y\geq0$，即 $0\leq y\leq2$。$y$ 有无穷多个取值，无法由 $x=1$ 唯一确定——所以 $R_4$ 不是函数。
\end{solution}

\begin{problem}
设 $R=\{(x,y)\mid x^2+y^2\leq4\}$。判断 $R$ 是否为函数，并用垂线检验说明理由。
\end{problem}
\begin{solution}
不是函数。取 $x=0$，则 $y\in[-2,2]$，一个 $x$ 对应无穷多个 $y$；几何上看，竖直线 $x=0$ 与圆盘相交于无穷多个点，违反"每列至多一个点"。
\end{solution}

\subsection{两条限制轴：从区域到基本初等函数}

\key{探究四：单值性与正则性，是两条互相独立的轴}

任意点集（撒点）常常同时「一对多」又「不连续」。要把它变成能研究的函数，需要在\textbf{两条不同的轴}上分别加限制，二者不可互相替代：

\begin{itemize}
  \item \key{轴一·代数／规则——单值性。}\key{回扣：}这正是前文那条收紧链——从 $A\times B$ 的完全自由（$2^9$ 种取法），到单值（$4^3$；若处处有定义则 $3^3=27$），再到统一法则。具体做法：规定每条竖直线 $x=a$ 上\textbf{至多}保留一个点，并且\textbf{明确指定}保留哪一个（不能随手挑）；未指定的 $x$ 允许没有对应点——这些被指定的 $x$ 全体就是定义域 $A$。满足单值性的关系叫\key{映射} $f:A\to B$；当 $A,B$ 取为数集时，就是\key{函数} $y=f(x)$。判别标准只有一条：垂线检验（图 \ref{fig:orderline} 第 (3) 幅）。再要求每个对应都能用同一条解析式表达，才得到「写得出来的函数」。
  \item \key{轴二·几何／分析——正则性与结构保持。}单值性\emph{不附带}任何光滑承诺：轴一走到尽头，图象照样可以处处间断（开场三问的 Dirichlet 函数），也可以连续却处处不可微。要让它「越来越好算」，必须继续加条件：连续性 $\to$ 可微性 $\to$ 某类「运算封闭／结构保持」的等式，最终落在\key{基本初等函数}上：
  \begin{itemize}
    \item 指数函数 $y=a^x$：把定义域的加法变成值域的乘法，$f(x_1+x_2)=f(x_1)f(x_2)$；
    \item 对数函数 $y=\log_a x$：把指数函数的数对 $(x,y)$ 调换为 $(y,x)$ 的逆关系，图象关于直线 $y=x$ 对称；
    \item 幂函数 $y=x^{\alpha}$：刻画幂次结构（面积、体积、增长阶）；
    \item 三角函数：刻画周期结构（旋转、振动、波动）。
  \end{itemize}
\end{itemize}

\key{两条轴的分工：}轴一回答「它是不是函数」，轴二回答「它好不好算」。而方程与不等式（见本节开头的约定）是描述点集条件的语言，在任何一档舞台、任何一条轴上都可用：用二元方程 $F(x,y)=0$ 描述轴二下的平面曲线（可以一对多，如圆 $x^2+y^2=4$）；用一元方程 $f(x)=g(x)$ 求两条函数图象交点的横坐标；用一元不等式 $f(x)>g(x)$ 表示「$f$ 的点集位于 $g$ 的点集上方」的 $x$ 的范围。

\begin{figure}[htbp]
  \centering
  \begin{minipage}{0.235\textwidth}
    \centering
    \begin{tikzpicture}[x=0.6cm,y=0.6cm]
      \draw[main!60!black,thick,fill=main!4] (-2.2,-2) rectangle (2.2,2);
      \pgfmathsetseed{11}
      \foreach \i in {1,...,140}{
        \pgfmathsetmacro{\px}{-2.05+4.1*rnd}
        \pgfmathsetmacro{\py}{-1.85+3.7*rnd}
        \fill[main!45] (\px,\py) circle (0.8pt);
      }
    \end{tikzpicture}
    \\[2pt]{\footnotesize (1) 无序：随手撒点——一个 $x$ 拖着一大段 $y$}
  \end{minipage}\hfill
  \begin{minipage}{0.235\textwidth}
    \centering
    \begin{tikzpicture}[x=0.6cm,y=0.6cm]
      \draw[main!60!black,thick,fill=main!4] (-2.2,-2) rectangle (2.2,2);
      \pgfmathsetseed{11}
      \foreach \i in {1,...,140}{
        \pgfmathsetmacro{\px}{-2.05+4.1*rnd}
        \pgfmathsetmacro{\py}{-1.85+3.7*rnd}
        \fill[main!22] (\px,\py) circle (0.8pt);
      }
      \draw[second,very thick] (-0.6,-2)--(-0.6,2);
      \foreach \yy in {-1.75,-1.3,-0.85,-0.4,0.05,0.5,0.95,1.4,1.8}
        \fill[second] (-0.6,\yy) circle (1.5pt);
      \node[second,font=\scriptsize,right] at (-0.52,1.82) {$x=a$};
    \end{tikzpicture}
    \\[2pt]{\footnotesize (2) 竖截：竖线 $x=a$ 只截一段，仍是一整列点}
  \end{minipage}\hfill
  \begin{minipage}{0.235\textwidth}
    \centering
    \begin{tikzpicture}[x=0.6cm,y=0.6cm]
      \draw[main!60!black,thick,fill=main!4] (-2.2,-2) rectangle (2.2,2);
      \draw[third!55,thin] plot coordinates {(-2.05,1.82)(-1.89,1.54)(-1.72,1.15)(-1.56,-0.49)(-1.39,0.13)(-1.23,-0.37)(-1.07,-0.38)(-0.9,-1.22)(-0.74,-1.25)(-0.57,-0.93)(-0.41,-1.39)(-0.25,-1.05)(-0.08,-1.8)(0.08,-1.7)(0.25,-0.61)(0.41,-1.05)(0.57,-0.81)(0.74,-0.33)(0.9,-0.83)(1.07,-1.18)(1.23,0.24)(1.39,0.12)(1.56,0.48)(1.72,0.69)(1.89,1.58)(2.05,0.87)};
      \foreach \p in {(-2.05,1.82),(-1.89,1.54),(-1.72,1.15),(-1.56,-0.49),(-1.39,0.13),(-1.23,-0.37),(-1.07,-0.38),(-0.9,-1.22),(-0.74,-1.25),(-0.57,-0.93),(-0.41,-1.39),(-0.25,-1.05),(-0.08,-1.8),(0.08,-1.7),(0.25,-0.61),(0.41,-1.05),(0.57,-0.81),(0.74,-0.33),(0.9,-0.83),(1.07,-1.18),(1.23,0.24),(1.39,0.12),(1.56,0.48),(1.72,0.69),(1.89,1.58),(2.05,0.87)}
        \fill[third] \p circle (1.5pt);
    \end{tikzpicture}
    \\[2pt]{\footnotesize (3) 每列一点：被钉住了，但点列仍在乱跳}
  \end{minipage}\hfill
  \begin{minipage}{0.235\textwidth}
    \centering
    \begin{tikzpicture}[x=0.6cm,y=0.6cm]
      \draw[main!60!black,thick,fill=main!4] (-2.2,-2) rectangle (2.2,2);
      \pgfmathsetseed{37}
      \foreach \i in {1,...,90}{
        \pgfmathsetmacro{\bx}{-2.05+4.1*(\i-1)/89}
        \pgfmathsetmacro{\by}{0.5*\bx*\bx-1+0.55*(2*rnd-1)}
        \fill[main!40] (\bx,\by) circle (0.7pt);
      }
      \draw[second,very thick,domain=-2.05:2.05,samples=60,smooth] plot(\x,{0.5*\x*\x-1});
      \node[second,font=\scriptsize] at (0,-1.78) {$y=\tfrac12x^2-1$};
    \end{tikzpicture}
    \\[2pt]{\footnotesize (4) 收束：粗带压扁成唯一的曲线}
  \end{minipage}
  \caption{图像线是怎样从点集里长出来的：同一块平面区域上的四步图像实验。(1) 无序——随手撒点，一个 $x$ 的正上方压着一大段 $y$，一对多；(2) 竖截——用竖直线 $x=a$ 去截点云，截出的仍是一整列（一段）点，一个 $x$ 对应多个 $y$；(3) 立规矩——每条竖线只留一个点，得到一般函数的图象，但点列仍无规律地上下跳动；(4) 收束——点越密，跳动的点列糊成一条有厚度的粗带，厚度趋于 $0$ 时只剩唯一的曲线 $y=\tfrac12x^2-1$。图象就是「按 $x$ 排好队、且每个 $x$ 只剩唯一 $y$」的点列极限}
  \label{fig:orderline}
\end{figure}

\begin{conccard}[一级]{单值性与垂线检验}{任意平面点集（或图象）}{把"每列至多一个点"误写成"每行至多一个点"；忽略定义域}
点集 $P\subseteq\mathbb{R}^2$ 是某函数 $y=f(x)$ 的图象，当且仅当每条竖直线 $x=a$ 与 $P$ 至多有一个交点。代数表述：对每个 $x$，满足 $(x,y)\in P$ 的 $y$ 至多一个。这条判据把"任意点集"过滤成"函数图象"，是探究二、探究三共用的核心工具。
\end{conccard}

\begin{conccard}[二级]{三视角互译}{二元关系、函数、方程的解集}{只在一个视角里思考；把"数对集合"与"点集"割裂}
同一对象有三种等价表示：代数（有序数对集合）$\leftrightarrow$ 几何（坐标系点集）$\leftrightarrow$ 映射（箭头图）。代数精确、几何直观、箭头图体现对应方向。解题时按需要在三种语言间切换：例如解 $f(x)>g(x)$，就是看 $f$ 的点集在 $g$ 的点集上方时 $x$ 的范围。
\end{conccard}

\begin{method}
三视角互译法：遇到函数、方程、不等式问题时，先写出集合（代数），再画图（几何），必要时画箭头图（映射），最后回到代数表达。例如判断 $f(x)>g(x)$ 的解：画两条曲线，看 $f$ 在 $g$ 上方的 $x$ 区间，再写出集合 $\{x\mid f(x)>g(x)\}$。
\end{method}

\begin{commonerror}{把二维区域误当函数}
误以为"式子中有 $x,y$ 就是函数"或"给出 $x$ 就能算出 $y$"。反例：$x^2+y^2=4$ 中 $x=0$ 对应 $y=\pm2$，是一对多，不是函数；$\{(x,y)\mid x+y\leq3\}$ 中 $x=1$ 对应无穷多个 $y$。判别时必须回到单值性：每个 $x$ 是否唯一确定 $y$。
\end{commonerror}

\begin{example}
说明方程 $x^2=4$ 的解可以看成两个函数图象交点的横坐标，并求出解。
\end{example}
\begin{solution}
$x^2=4$ 即 $x^2-4=0$。令 $f(x)=x^2-4$，$g(x)=0$，交点横坐标即 $f(x)=g(x)$ 的解：$x=\pm2$。
也可看作曲线 $y=x^2$ 与直线 $y=4$ 交点的横坐标，结果同样是 $x=\pm2$。
\end{solution}

\begin{problem}
设 $f(x)=x^2$，$g(x)=2x+3$。求不等式 $f(x)>g(x)$ 的解集，并指出它对应"哪段 $x$ 上 $f$ 的图象在 $g$ 上方"。
\end{problem}
\begin{solution}
$x^2>2x+3\iff x^2-2x-3>0\iff(x-3)(x+1)>0$，解集为 $(-\infty,-1)\cup(3,+\infty)$。
几何意义：在 $(-\infty,-1)$ 与 $(3,+\infty)$ 上，抛物线 $y=x^2$ 位于直线 $y=2x+3$ 的上方；在 $[-1,3]$ 上则位于下方（或相交）。
\end{solution}

\subsection{大一统结构对照}

把本节对象放进同一张表：它们都是「集合」，并且位于同一条包含链上——关系 $\supseteq$ 映射 $\supseteq$ 函数 $\supseteq$ 基本初等函数；方程与不等式则是按变元个数给关系分类的条件表达法。

\begin{center}
\begin{tabular}{|c|p{3.0cm}|p{5.2cm}|p{4.2cm}|}
\hline
\textbf{数学对象} & \textbf{集合与点集本质} & \textbf{几何维度直观理解} \\
\hline
二元关系 $R\subseteq A\times B$ & $A\times B$ 的任意子集（条件句挑出） & 点云：孤立点／竖直线段族／区域／曲线皆可 \\
\hline
映射 $f:A\to B$ & 每列至多一个点的单值关系 & 箭头图；数集情形为一维曲线 \\
\hline
二元不等式（组） & $\{(x,y)\mid \text{约束条件}\}$ & $\mathbb{R}^2$ 中的二维区域（"面"） \\
\hline
二元方程 $F(x,y)=0$ & $\{(x,y)\mid F(x,y)=0\}$ & 平面上的一维曲线（可一对多） \\
\hline
函数 $y=f(x)$ & $\{(x,f(x))\mid x\in A\}$ & 一维曲线，且每列至多一个点 \\
\hline
一元方程的解 & $\{x\mid f(x)=g(x)\}$ & 数轴上的零维点（交点横坐标） \\
\hline
一元不等式的解 & $\{x\mid f(x)>g(x)\}$ & 数轴上的一维区间（上下位置） \\
\hline
\end{tabular}
\end{center}

\begin{summary}
本节主线：关系 $R\subseteq A\times B$ 有三档舞台——离散$\times$离散的孤立点、离散$\times$连续的竖直线段族、连续$\times$连续的二维区域；方程与不等式是把关系写成条件句的表达法（变元个数定维度，等号／不等号定边界与区域）。在关系之上加限制要分两条独立的轴：轴一加单值性，得到映射与函数（判别标准只有垂线检验）；轴二加连续性、可微性与结构保持等式，得到基本初等函数（指数、对数、幂、三角）。三种语言（代数、几何、箭头图）全程互译。
\end{summary}

\subsection{函数与图形变换完全图谱}

\begin{conccard}[一级]{第一性原理：一切图像变换都是有序数对 $(x,y)$ 的映射}{抽象函数图像辨析、变换类压轴题}{只背“左加右减”口诀，不追踪每个点的去向}
函数图像是 $\mathbb{R}^2$ 中的点集 $\{(x,y)\mid y=f(x)\}$。所谓“图像变换”，本质是对集合中每个点的坐标施加映射。
掌握下面三级体系与对称周期生成律，就能把“看图识函数”变成可推导的代数操作。
\end{conccard}

\begin{center}
\renewcommand{\arraystretch}{1.4}
\begin{tabular}{|p{3.2cm}|p{4.0cm}|p{4.8cm}|}
\hline
\textbf{变换} & \textbf{点集规则} & \textbf{表达式与几何效应} \\
\hline
\multicolumn{3}{|l|}{\textbf{一级：刚体等距（形状大小不变）}} \\
\hline
平移 & $(x,y)\mapsto(x+h,\,y+k)$ & $y-k=f(x-h)$（逆向代换） \\
\cline{2-3}
轴对称（$x=a$） & $(x,y)\mapsto(2a-x,\,y)$ & $y=f(2a-x)$ \\
\cline{2-3}
中心对称（$(a,b)$） & $(x,y)\mapsto(2a-x,\,2b-y)$ & $y=2b-f(2a-x)$（$a=b=0$ 即奇函数） \\
\hline
\multicolumn{3}{|l|}{\textbf{二级：伸缩与翻折}} \\
\hline
轴向伸缩 & $(x,y)\mapsto(\frac1\omega x,\,Ay)$ & $y=Af(\omega x)$ \\
\cline{2-3}
绝对值翻折 & $(x,y)\mapsto(x,\,|y|)$；$x\mapsto|x|$ & $y=|f(x)|$（下方翻上）；$y=f(|x|)$（右侧复制到左侧，强制偶化） \\
\hline
\multicolumn{3}{|l|}{\textbf{三级：非线性算子（图像拓扑突变）}} \\
\hline
自变量平方／倒数 & $u=g(x)$ 作为内层 & $y=f(x^2)$ 必为偶函数且在 $x=0$ 处切线水平 \\
\cline{2-3}
函数值平方／倒数 & $(x,y)\mapsto(x,\,y^2)$；$(x,y)\mapsto(x,\,\frac1y)$ & 零点变极小值点／垂直渐近线（详见下文铁律） \\
\hline
\end{tabular}
\end{center}

\begin{method}
\textbf{双对称推周期：连续两次平行镜像反射 = 平移}。看到“两个对称条件”或“半周期反对称”，直接锁定最小周期，不必画图：
\begin{enumerate}
\item \textbf{双轴对称}：$f(a+x)=f(a-x)$ 且 $f(b+x)=f(b-x)$ $\Rightarrow$ $T=2|b-a|$；
\item \textbf{双中心对称}：$f(a+x)+f(a-x)=2c$ 且 $f(b+x)+f(b-x)=2c$ $\Rightarrow$ $T=2|b-a|$；
\item \textbf{一轴一中心}（最高频）：$f(a+x)=f(a-x)$ 且 $f(b+x)+f(b-x)=2c$ $\Rightarrow$ $T=4|b-a|$；
\item \textbf{滑移反对称}：$f(x+a)=-f(x)$、$f(x+a)=\pm\frac1{f(x)}$ 或 $f(x+a)=\frac{1-f(x)}{1+f(x)}$，迭代两次即复原 $\Rightarrow$ $T=2a$。
\end{enumerate}
几何直观：点 $P(x)$ 先关于 $x=a$ 反射到 $P_1(2a-x)$，再关于 $x=b$ 反射到 $P_2\bigl(2b-(2a-x)\bigr)=P(x+2(b-a))$——两次翻转方向复原，合成一次纯平移。
\end{method}

\begin{center}
\begin{tikzpicture}[scale=1.0]
\draw[->, gray, thick] (-1.2,0) -- (7.2,0) node[right] {$x$};
\draw[main, dashed, very thick] (1.5,-0.6) -- (1.5,2.4) node[above] {$x=a$};
\draw[second, dashed, very thick] (4.0,-0.6) -- (4.0,2.4) node[above] {$x=b$};
\fill[black] (0.5,1.5) circle (2pt) node[above] {$P(x)$};
\fill[main] (2.5,1.5) circle (2pt) node[below] {$P_1(2a-x)$};
\fill[second] (5.5,1.5) circle (2pt) node[below] {$P_2\bigl(x+2(b-a)\bigr)$};
\draw[<->, gray, thick] (0.5,-0.35) -- (5.5,-0.35) node[midway, below, font=\small\bfseries] {合成平移 $T=2(b-a)$};
\end{tikzpicture}
\end{center}

\begin{conccard}[二级]{导数对称遗传律：周期全遗传，对称奇偶互换}{含参导数压轴、切线分析}{把 $f$ 与 $f'$ 的对称性混为一谈}
对 $f$ 的恒等式两边求导（链式法则）即得：
\begin{itemize}
\item $f$ 周期为 $T$ $\iff$ $f'$ 周期仍为 $T$；
\item $f$ 关于直线 $x=a$ 轴对称 $\Rightarrow$ $f'$ 关于点 $(a,0)$ 中心对称，且 $f'(a)=0$；
\item $f$ 关于点 $(a,b)$ 中心对称 $\Rightarrow$ $f'$ 关于直线 $x=a$ 轴对称。
\end{itemize}
\end{conccard}

\begin{method}
\textbf{非线性算子的五条铁律}（以 $g(x)=\dfrac1{f(x)}$ 为例，$y=[f(x)]^2$ 见表中）：
\begin{enumerate}
\item \key{符号不变}：$f>0$ 处 $g>0$，$f<0$ 处 $g<0$，上下位置不动；
\item \key{零点撕裂}：$f(x_0)=0$ 变成垂直渐近线 $x=x_0$；
\item \key{无穷变零}：$f\to\infty$ 处 $g\to0$（趋近 $x$ 轴）；
\item \key{单调反转}：同一同号区间内 $g$ 与 $f$ 单调性相反（不能跨渐近线谈单调）；
\item \key{极值互换}：正的极大值 $M$ 变成极小值 $\frac1M$；而 $y=\pm1$ 的点是\key{不动点}，位置纹丝不动。
\end{enumerate}
$y=f(x^2)$ 与 $y=f\!\left(\frac1x\right)$ 则作用在自变量上：$f(x^2)$ 是偶函数，
左半支与右半支互为镜像（同一 $|x|$ 处函数值相同）；在 $x\ge0$ 上它等于把
$y=f(u)\,(u\ge0)$ 的横坐标按 $u=x^2$ 重标得到的曲线——靠近原点被压缩、
远离原点被拉伸；$f\!\left(\frac1x\right)$ 则把 $x\to0$ 的行为与
$x\to\pm\infty$ 的行为\key{对调}。

\textbf{平方算子 $y=[f(x)]^2$ 的五条铁律}（与倒数算子对照记忆）：
\begin{itemize}
  \item \key{非负化}：整个图像被压进 $y\ge0$，负部沿 $x$ 轴\key{翻折}上来；
  \item \key{零点不变}：$f(x_0)=0\iff f^2(x_0)=0$，只是零点的「重数」加倍；
  \item \key{无穷仍无穷}：$f\to\infty$ 时 $f^2\to+\infty$ 增长更快；
  \item \key{单调分区}：由 $(f^2)'=2ff'$，$f>0$ 处 $f^2$ 与 $f'$ 同号、
        $f<0$ 处反号，故在 $f$ 的正区间保持单调方向、负区间反转；
  \item \key{极值缩放}：$|f|$ 的极大值 $M$ 变成 $f^2$ 的极大值 $M^2$
        （$M\ge1$ 时变大、$0<M<1$ 时变小），$|f|$ 的小极小值则被压得更小；
        $y=1$ 上的点是\key{不动点}。
\end{itemize}
\end{method}

\begin{center}
\begin{tikzpicture}[scale=1.15]
\draw[->, gray, thick] (-1.8,0) -- (3.8,0) node[right] {$x$};
\draw[->, gray, thick] (0,-2.2) -- (0,3.0) node[above] {$y$};
\draw[gray, dashed, thick] (2,-2.2) -- (2,3.0) node[above] {$x=2$};
\draw[main, thick, domain=-1.4:3.4, samples=60] plot(\x, {(\x)*(\x-2)});
\node[main] at (3.1,1.6) {$f(x)=x^2-2x$};
\draw[second, thick, domain=-1.5:-0.12, samples=60] plot(\x, {1/((\x)*(\x-2))});
\draw[second, thick, domain=0.12:1.9, samples=80] plot(\x, {1/((\x)*(\x-2))});
\draw[second, thick, domain=2.12:3.5, samples=60] plot(\x, {1/((\x)*(\x-2))});
\fill[second] (1,-1) circle (2pt) node[below left] {$(1,-1)$};
\end{tikzpicture}
{\small\itshape 抛物线与其倒数：零点 $0,2$ 变为竖直渐近线，极小值 $-1$ 变为极大值 $-1$。}
\end{center}

\begin{example}[双对称推周期与导数遗传]
已知 $f(x)$ 在 $\mathbb{R}$ 上可导，$f(x+1)$ 为偶函数，$f(x+2)$ 为奇函数，且 $f(1)=2$，求 $f(2027)+f'(2025)$。
\begin{solution}
$f(x+1)$ 偶 $\Rightarrow f(1-x)=f(1+x)$，故 $f$ 关于 $x=1$ 轴对称；
$f(x+2)$ 奇 $\Rightarrow f(2-x)=-f(2+x)$，故 $f$ 关于 $(2,0)$ 中心对称。
一轴一中心 $\Rightarrow T=4|2-1|=4$。
由中心对称取 $t=1$：$f(3)=-f(1)=-2$，于是 $f(2027)=f(4\times506+3)=f(3)=-2$。
$f$ 关于 $x=1$ 轴对称 $\Rightarrow f'$ 关于 $(1,0)$ 中心对称，故 $f'(1)=0$；又 $f'$ 周期仍为 $4$，
$f'(2025)=f'(4\times506+1)=f'(1)=0$。合计 $\key{-2}$。
\end{solution}
\end{example}

\begin{example}[用奇偶性速判函数图像]
函数 $g(x)=\dfrac{x}{e^x-1}$ 的大致图像如何？
\begin{solution}
先看定义域 $x\neq0$。由切线不等式 $e^x>x+1\ (x\neq0)$ 知 $e^x-1$ 与 $x$ 同号，故 $g(x)>0$ 恒成立（全在 $x$ 轴上方）；且 $\lim\limits_{x\to0}g(x)=1$。
当 $x\to+\infty$ 时指数增长远快于一次项，$g(x)\to0^+$；当 $x\to-\infty$ 时 $e^x\to0$，$g(x)\approx\dfrac{x}{-1}=-x\to+\infty$，即左端贴着斜率 $-1$ 的直线向上。无需二阶求导即可锁定图像。
\end{solution}
\end{example}

\begin{commonerror}{把“对称”当“周期”直接用}
$f(a+x)=f(a-x)$ 只说明关于 $x=a$ 对称，一个对称条件\key{产生不了周期}；必须有两个对称条件，
或用滑移反对称结构。此外 $y=f(x^2)$ 是偶函数，而 $y=[f(x)]^2$ 只保证非负——两者都“偶化”方向不同，不可互推。
\end{commonerror}

图形变换还有一种对偶记法，能把「左加右减」的口诀风险降到最低：动曲线与动坐标系是同一事实的两种记账方式（图 \ref{fig:activepassive}）。

\begin{methodbox}{主被动对偶律：曲线平移，等价于坐标系反向平移}
设曲线 $C:y=f(x)$ 主动向右平移 $h$ 个单位，得到新曲线 $y=f(x-h)$；若改为让坐标轴向左平移 $h$ 个单位、曲线在空间中保持不动，则新读数满足 $x'=x+h$、$y'=y$，同一物理状态在新系下的方程仍写作 $y'=f(x'-h)$。
两式描述的是同一幅终态图：顶点的读数相同，只是「谁在动」的归属不同。若条件按曲线平移给出（如「图象右移 $2$」），用主动记法；若条件按新原点、新读数给出，用被动记法。两者方向相反是必然，切勿把坐标轴的移动方向也套用「左加右减」。
\end{methodbox}

\begin{figure}[htbp]
  \centering
  \begin{tikzpicture}[>=stealth,font=\small]
    % ---------- 左：主动变换（坐标系不动，曲线右移） ----------
    \begin{scope}
      \draw[->,gray,thick] (-1.4,0)--(3.6,0) node[right,gray,font=\footnotesize] {$x$};
      \draw[->,gray,thick] (0,-0.7)--(0,3.0) node[above,gray,font=\footnotesize] {$y$};
      \draw[gray!70,dashed,thick,domain=-1.15:1.15,samples=50,smooth] plot(\x,{\x*\x});
      \draw[main,very thick,domain=0.85:3.15,samples=60,smooth] plot(\x,{(\x-2)*(\x-2)});
      \fill[gray!70] (0,0) circle (0.05);
      \node[gray!70,below left,font=\footnotesize] at (0,0) {$(0,0)$};
      \fill[main] (2,0) circle (0.05);
      \node[main,below right,font=\footnotesize] at (2,0) {$(2,0)$};
      % 平移箭头画在 x 轴下方空白处，避免压曲线
      \draw[<->,second,thick] (0.15,-0.35)--(1.85,-0.35)
        node[midway,above,font=\footnotesize] {右移 $2$};
      \node[main,font=\footnotesize] at (2.55,2.5) {$y=(x-2)^2$};
      \node[font=\footnotesize] at (1.1,-1.45) {主动：坐标系不动，曲线右移 $2$};
    \end{scope}
    % ---------- 右：被动变换（曲线不动，坐标轴左移） ----------
    \begin{scope}[xshift=7.7cm]
      \draw[gray!55,dashed,thick] (-1.4,0)--(2.2,0);
      \draw[gray!55,dashed,thick] (0,-0.7)--(0,3.0);
      \node[gray!55,font=\scriptsize,right] at (-1.35,2.55) {旧系（虚线）};
      \draw[->,second,thick] (-3.5,0)--(1.7,0) node[right,second,font=\footnotesize] {$x'$};
      \draw[->,second,thick] (-2,-0.7)--(-2,3.0) node[above,second,font=\footnotesize] {$y'$};
      \draw[main,very thick,domain=-1.15:1.15,samples=50,smooth] plot(\x,{\x*\x});
      \fill[main] (0,0) circle (0.05);
      \node[gray!60,font=\scriptsize,left] at (-0.06,0.28) {旧 $O$};
      \node[second,below left,font=\footnotesize] at (-2,0) {$O'$};
      \node[main,font=\footnotesize,right] at (0.12,-0.32) {顶点 $(2,0)$};
      \draw[<->,third,thick] (-2,-0.85)--(0,-0.85)
        node[midway,below,font=\footnotesize] {轴左移 $2$};
      \draw[third!70!black,->,shorten <=1pt,shorten >=1pt] (0.2,-0.68)--(0.06,-0.12);
      \node[third!70!black,font=\scriptsize,right] at (0.25,-0.72) {读数 $x'=2$};
      \node[main,font=\footnotesize] at (0.95,2.25) {$y'=(x'-2)^2$};
      \node[font=\footnotesize] at (-0.85,-1.45) {被动：曲线不动，坐标轴左移 $2$};
    \end{scope}
  \end{tikzpicture}
  \caption{主被动对偶：左图坐标系不动，抛物线 $y=x^2$ 主动右移 $2$ 得 $y=(x-2)^2$；右图抛物线原地不动，坐标轴被动左移 $2$（虚线为旧系），新读数 $x'=x+2$，顶点在新系下仍读作 $(2,0)$。两幅终态图完全相同——「动曲线」与「动坐标系」是同一平移的两种归属，不要混用方向}
  \label{fig:activepassive}
\end{figure}

% ---------- 收官微专题：$45^\circ$ 镜面反射定律 ----------

\begin{conccard}[二级]{唯有 $k=\pm1$ 的镜面，能保住「函数」的资格}{反函数对称、指对同构压轴选填、解析几何对称速算}{把斜线对称当成左右平移；漏掉 $b$ 的补偿项}
把对称轴从竖直线 $x=a$ 推广到一般直线 $y=kx+b$ 时，反射会让横、纵坐标\textbf{互相耦合}，反射后的曲线往往不再通过垂线检验，\key{跌出函数集合}，退回一般二元方程 $F(x,y)=0$。只有 $k=\pm1$（倾角 $45^\circ$ 或 $135^\circ$）时，耦合系数全部约尽，退化为干净的坐标互换——这正是反函数与指对同构压轴题的共同底座。
\end{conccard}

先把最典型的镜面 $l:y=x$ 看清楚（图 \ref{fig:expinvlog}）：反射把指数函数的点集 $\{(x,y)\}$ 整体翻成 $\{(y,x)\}$，对数曲线由此生成——这就是「反函数＝逆关系点集」的几何来源。

\begin{figure}[htbp]
  \centering
  \begin{tikzpicture}[x=0.9cm,y=0.9cm,>=stealth,font=\small]
    \draw[->,gray,thick] (-1.7,-1.7)--(3.6,-1.7) node[right] {$x$};
    \draw[->,gray,thick] (-1.7,-1.7)--(-1.7,3.8) node[above] {$y$};
    \draw[gray,dashed,thick,domain=-1.4:3.4,samples=40] plot(\x,\x) node[above right] {$y=x$};
    \draw[main,very thick,domain=-1.3:1.35,samples=80,smooth] plot(\x,{exp(\x)}) node[right] {$y=e^x$};
    \draw[second,very thick,domain=0.28:3.4,samples=90,smooth] plot(\x,{ln(\x)}) node[right] {$y=\ln x$};
    \fill[main] (0,1) circle (0.07);       \node[main,left] at (0,1) {$(0,1)$};
    \fill[main] (1,2.718) circle (0.07);   \node[main,above right] at (1,2.718) {$(1,e)$};
    \fill[main] (-1,0.368) circle (0.07);  \node[main,below left] at (-1,0.368) {$(-1,e^{-1})$};
    \fill[second] (1,0) circle (0.07);      \node[second,below] at (1,0) {$(1,0)$};
    \fill[second] (2.718,1) circle (0.07);  \node[second,right] at (2.718,1) {$(e,1)$};
    \fill[second] (0.368,-1) circle (0.07); \node[second,below left] at (0.368,-1) {$(e^{-1},-1)$};
    \draw[->,third,thick] (0,1) to[bend left=12] (1,0);
    \draw[->,third,thick] (1,2.718) to[bend left=18] (2.718,1);
    \draw[->,third,thick] (-1,0.368) to[bend right=12] (0.368,-1);
  \end{tikzpicture}
  \caption{同一坐标系中的互逆反射：$y=e^x$ 与 $y=\ln x$ 关于 $y=x$ 对称，每个点 $(a,b)$ 与它的像 $(b,a)$ 用跨越镜面的箭头一一相连，例如 $(0,1)\leftrightarrow(1,0)$、$(1,e)\leftrightarrow(e,1)$、$(-1,e^{-1})\leftrightarrow(e^{-1},-1)$。对数曲线就是指数曲线点集 $\{(x,y)\}$ 经 $y=x$ 翻转后的 $\{(y,x)\}$}
  \label{fig:expinvlog}
\end{figure}

\begin{method}
\textbf{$45^\circ$ 镜面反射替换速查表}（镜面记为 $l$，点 $P(x,y)$ 的像为 $P'(x',y')$）：
\begin{center}
\renewcommand{\arraystretch}{1.4}
\small
\begin{tabular}{|p{2.0cm}|p{3.6cm}|p{3.9cm}|p{4.2cm}|}
\hline
\textbf{镜面} & \textbf{点集映射} & \textbf{曲线替换口径} & \textbf{典型用途} \\
\hline
$l:y=x$ & $(x,y)\mapsto(y,\,x)$ & $F(y,\,x)=0$ & 反函数：$y=a^x$ 与 $y=\log_a x$ \\
\hline
$l:y=-x$ & $(x,y)\mapsto(-y,\,-x)$ & $F(-y,\,-x)=0$ & 双轴对称联动：$y=\dfrac1x$ \\
\hline
$l:y=x+b$ & $(x,y)\mapsto(y-b,\,x+b)$ & $F(y-b,\,x+b)=0$ & 平移型反函数、指对同构压轴 \\
\hline
$l:y=-x+b$ & $(x,y)\mapsto(b-y,\,b-x)$ & $F(b-y,\,b-x)=0$ & 斜率 $-1$ 直线族（$y=x$ 的垂线） \\
\hline
\end{tabular}
\end{center}
\vspace{0.3em}
\key{秒杀推论}：若两条曲线关于 $y=x$ 对称（如互为反函数），且被斜率为 $-1$ 的直线 $x+y=m$ 所截，则两交点 $A,B$ 必关于 $y=x$ 对称，中点 $\bigl(\tfrac m2,\tfrac m2\bigr)$ 落在 $y=x$ 上，于是
$$x_A+x_B=m,\qquad y_A+y_B=m .$$
\end{method}

\begin{figure}[htbp]
  \centering
  \begin{tikzpicture}[x=0.95cm,y=0.95cm,>=stealth,font=\small]
    \draw[->,gray,thick] (-1.2,-1.2)--(5.2,-1.2) node[right] {$x$};
    \draw[->,gray,thick] (-1.2,-1.2)--(-1.2,5.2) node[above] {$y$};
    \draw[main,dashed,very thick,domain=-1.1:3.9,samples=40] plot(\x,\x) node[above right] {$y=x$};
    \draw[second,very thick,domain=0.09:4.8,samples=90] plot(\x,{ln(\x)}) node[right] {$y=\ln x$};
    \draw[third,very thick,domain=-1:1.2,samples=70] plot(\x,{exp(\x)}) node[above] {$y=e^x$};
    \draw[red,very thick] (-1,5)--(5,-1) node[below right] {$x+y=4$};
    \fill[main] (2,2) circle (2.4pt) node[above left] {$M(2,2)$};
    \fill[red] (1.07,2.93) circle (2.4pt) node[above left] {$A(x_1,y_1)$};
    \fill[red] (2.93,1.07) circle (2.4pt) node[below right] {$B(y_1,x_1)$};
  \end{tikzpicture}
  \caption{反函数 $y=e^x$ 与 $y=\ln x$ 关于 $y=x$ 对称；斜率为 $-1$ 的直线 $x+y=4$ 与二者各交于 $A,B$，两点关于 $y=x$ 对称，中点 $M(2,2)$ 落在 $y=x$ 上，故 $x_1+x_2=4$}
  \label{fig:mirror45}
\end{figure}

\begin{example}[指对同构：不解方程，三步出答案]
设 $x_1$ 满足 $e^{x_1}+x_1=4$，$x_2$ 满足 $\ln x_2+x_2=4$，求 $x_1+x_2$。
\begin{solution}
两式分别改写为 $e^{x_1}=4-x_1$、$\ln x_2=4-x_2$。
于是 $A(x_1,e^{x_1})$ 是曲线 $y=e^x$ 与直线 $y=4-x$ 的交点，$B(x_2,\ln x_2)$ 是曲线 $y=\ln x$ 与同一直线的交点（图 \ref{fig:mirror45}）。
$y=e^x$ 与 $y=\ln x$ 互为反函数，关于 $y=x$ 对称；直线 $y=4-x$ 的斜率为 $-1$，与 $y=x$ 垂直，故 $A,B$ 关于 $y=x$ 对称，即 $x_2=e^{x_1}$。
所以 $x_1+x_2=x_1+e^{x_1}=\key{4}$。全程无需解超越方程。
\end{solution}
\end{example}

\begin{example}[平移型镜面：对称轴不是 $y=x$ 时]
求曲线 $C_1:y=e^x$ 关于直线 $l:y=x+1$ 对称的曲线 $C_2$ 的方程，并验证点 $P(1,e)$ 与其对称点 $P'$ 分别落在 $C_1,C_2$ 上。
\begin{solution}
用 $l:y=x+b$（$b=1$）的替换口径 $F(y-b,\,x+b)=0$，把 $C_1$ 的方程 $y=e^x$ 中的 $x$ 换成 $y-1$、$y$ 换成 $x+1$：
$$x+1=e^{\,y-1}\implies y=\ln(x+1)+1,$$
即 $C_2:y=\ln(x+1)+1$。
点 $P(1,e)$ 的对称点为 $P'=(y-b,\,x+b)=(e-1,\,2)$；代入 $C_2$：$\ln\bigl((e-1)+1\bigr)+1=\ln e+1=2$，恰好落在 $C_2$ 上。
（顺带一提：$x=0$ 时 $e^0=1=0+1$，两条曲线在镜面上的点 $(0,1)$ 处相遇。）
\end{solution}
\end{example}

\begin{commonerror}{把「$y=x+b$ 对称」误当成「左移 $b$」}
关于 $y=x+b$ 的对称\emph{不是}平移：它同时交换了横纵坐标的角色。若把 $C_1:y=e^x$ 关于 $y=x+1$ 的对称曲线误写成 $y=e^{x-1}$（左移思路），代入点 $P(1,e)$ 的对称点 $(e-1,2)$ 立即矛盾——正确的替换永远是 $x\leftrightarrow y-b$ 的\emph{互换}，而不是同向加减。
\end{commonerror}

\end{explore}

\begin{inclasstest}
\begin{enumerate}
  \item 用三视角（代数、几何、箭头图）表示关系 $R=\{(1,1),(2,4),(3,9)\}$，判断它是否为函数，并写出对应法则；\level{易}
  \item 判断点集 $P=\{(x,y)\mid x^2+y^2=1\}$ 是否为函数，用垂线检验说明理由；\level{易}
  \item 点集 $\{(x,y)\mid x^2+y^2\leq4\}$ 与 $\{(x,y)\mid x+y\leq3\}$ 的交集表示什么几何图形？它是否可能为某个函数的图象？说明理由；\level{中}
  \item 方程 $x^3=x$ 的解可看成哪两个函数图象交点的横坐标？求出所有解。\level{中}
\end{enumerate}
\end{inclasstest}

\begin{afterclass}
\subsection*{基础巩固}
\begin{enumerate}
  \item 写出 $A=\{1,2\}$、$B=\{3,4\}$ 的笛卡尔积 $A\times B$，并任取一个二元关系；
  \item 判断 $R=\{(1,2),(2,2),(3,2)\}$ 是否为函数，并说明理由。
\end{enumerate}

\subsection*{进阶提升}
\begin{enumerate}
  \item 画出 $R=\{(x,y)\mid y\geq x^2,\ x^2+y^2\leq4\}$ 表示的平面区域，并判断它是否为函数；
  \item 求不等式 $x^2<2x+3$ 的解集，并说明它对应两曲线怎样的位置关系。
\end{enumerate}

\subsection*{拓展探究}
\begin{enumerate}
  \item 查资料说明：为什么"对数能把乘除变成加减"？用 $y=a^x$ 与 $y=\log_a x$ 的互逆关系（$(x,y)\leftrightarrow(y,x)$）加以解释。
  \item （对应「实数的三种等价面孔」）设 $S=\{x\in\mathbb{Q}:x^2<2\}$。证明：$S$ 在 $\mathbb{Q}$ 中没有最小上界——要求对任意有理上界 $r$，构造出更小的有理上界；并由此说明曲线 $y=x^2-2$ 为什么能穿过有理数轴的孔洞。
  \item （表征与定义之别）某同学说：「$0.999\dots$ 和 $1$ 之间总差一点点，所以 $0.999\dots\neq1$。」请指出错在哪里，并用等比数列求和或方程法给出严格说明。
\end{enumerate}

\subsection*{教师版答案}
\begin{solution}
\textbf{基础 1：}$A\times B=\{(1,3),(1,4),(2,3),(2,4)\}$；任取一个子集即可，如 $R=\{(1,4),(2,3)\}$。
\end{solution}
\begin{solution}
\textbf{基础 2：}是函数。虽然三个 $x$ 都对应同一个 $y=2$，但每个 $x$ 对应的 $y$ 唯一（多对一是允许的），竖直线 $x=1,2,3$ 与点集各只有一个交点。
\end{solution}
\begin{solution}
\textbf{进阶 1：}区域是抛物线 $y=x^2$ 上方、圆盘 $x^2+y^2\leq4$ 内部（含边界）的公共部分，即圆盘被抛物线截出的"弓形"部分。它不是函数：取 $x=0$，$y\in[0,2]$，一个 $x$ 对应无穷多个 $y$（竖直线 $x=0$ 与区域有无穷多交点）。
\end{solution}
\begin{solution}
\textbf{进阶 2：}$x^2<2x+3\iff(x-3)(x+1)<0$，解集 $(-1,3)$。几何上，在 $(-1,3)$ 上抛物线 $y=x^2$ 位于直线 $y=2x+3$ 的下方。
\end{solution}
\begin{solution}
\textbf{拓展：}设 $x=a^{m}$，$y=a^{n}$，则 $m=\log_a x$，$n=\log_a y$。乘法 $xy=a^{m+n}$ 取对数得 $\log_a(xy)=m+n=\log_a x+\log_a y$，即乘法化为加法；除法 $x/y=a^{m-n}$ 同理化为减法。反函数关系体现在：指数函数把指数 $m$ 映为幂 $a^m$，对数函数把幂还原为指数，数对 $(x,y)$ 与 $(y,x)$ 关于 $y=x$ 对称。
\end{solution}
\begin{solution}
\textbf{拓展 2：}设 $r\in\mathbb{Q}$、$r>0$ 且 $r^2>2$（此时 $r$ 是上界），令 $t=\dfrac{r+\dfrac{2}{r}}{2}\in\mathbb{Q}$。由 $t-r=-\dfrac{r^2-2}{2r}<0$ 知 $t<r$，由 $t^2-2=\dfrac{(r^2-2)^2}{4r^2}>0$ 知 $t$ 仍是上界。故任何有理上界都不是最小，$S$ 在 $\mathbb{Q}$ 中无最小上界。几何上，$y=x^2-2$ 的零点恰是 $\sqrt2$；$\sqrt2\notin\mathbb{Q}$ 意味着这条连续曲线穿过数轴时，落点位置在有理数里没有任何对应——孔洞不仅存在，还被曲线「踩」了出来。
\end{solution}
\begin{solution}
\textbf{拓展 3：}设 $x=0.999\dots$，则 $10x=9.999\dots$，两式相减得 $9x=9$，即 $x=1$（方程法）；也可写成等比级数 $0.999\dots=0.9\bigl(1+0.1+0.01+\dots\bigr)=\dfrac{0.9}{1-0.1}=1$。「总差一点点」的错误在于：若 $1-x=\varepsilon>0$，则 $x\leq1-\varepsilon<1$，与 $x$ 的十进制表示各位都逼近 $1$ 矛盾——所谓「一点点」是一个比任何正实数都小的量，不存在。这也说明无限小数的含义要靠完备性才能定义清楚。
\end{solution}
\end{afterclass}
\reading{拓展阅读 07}{函数的结构特化}

% 高观点导学案：三角函数、解三角形与复平面旋转（第 2 章·高观点统一版）
% 主线：以「乘以 e^{iβ} = 逆时针旋转 β」为唯一母公理，让整套三角恒等变换从指数运算律自然生长
% 四模块：数系扩张与 i 的几何诞生 → 结构同构（旋转必为指数函数）→ 切线速度锁定底数 → 母公式引擎
% 另配一节：解三角形与边角互化。全程用「复平面旋转 ↔ 实虚部投影」互译，原生 TikZ 图可视化。

\section{三角函数、解三角形与复平面旋转（高观点统一版）}

\begin{methodbox}{多框架对读}
\small
\begin{center}
\begin{tabular}{c|p{2.9cm}|p{5.0cm}|p{3.1cm}}
\toprule
视角 & 本章看点 & 别的视角怎么说 & 哪侧最省力 \\
\midrule
代数 & 和差倍半公式、恒等变形 & 符号系统的机械推导 & 硬算与验算 \\
几何 & 单位圆上的旋转与投影 & 角度相加＝复数相乘 & 猜公式、看符号 \\
集合 & $\theta\mapsto(\cos\theta,\sin\theta)$ & 参数方程即一个对应 & 界定义域与值域 \\
映射 & 乘 $e^{i\beta}$＝逆时针转 $\beta$ & 旋转群与指数函数同构 & 统一推导所有公式 \\
\bottomrule
\end{tabular}
\end{center}
卡的使用：公式卡壳时回到单位圆「转一遍」；证明题优先用旋转映射（总图见第 0 章图~\ref{fig:four-view-map}）。
\end{methodbox}

本章不新增公式，而是给三角函数一章装上一个\key{母引擎}：
在复平面上，\key{乘以 $e^{i\beta}$ 就等于把向量逆时针旋转 $\beta$ 角}（模长不变）。
由这一条母公理出发，两角和差、二倍角三倍角、诱导公式、和差化积与积化和差，
全部只是指数运算律 $e^{i(\alpha+\beta)}=e^{i\alpha}e^{i\beta}$ 的代数推论。
四个模块环环相扣：
\key{数系扩张与 $i$ 的几何诞生}（为什么需要虚数）$\to$
\key{结构同构}（单位圆旋转必为指数函数）$\to$
\key{切线速度锁定底数}（欧拉公式 $e^{i\theta}=\cos\theta+i\sin\theta$）$\to$
\key{母公式引擎}（一条指数律长出全套公式），
最后用 \key{解三角形与边角互化} 把复数直觉落回高考主战场。
核心思想：\key{三角恒等变换不是几十个孤立公式，而是同一个旋转群的三种语言}（指数、坐标、三角形）。

\begin{learnobjective}
\begin{itemize}
  \item 理解虚数单位 $i$ 的几何意义（逆时针旋转 $90^\circ$），知道 $i^2=-1$ 是「转两次直角 $=$ 掉头」的代数翻译；
  \item 理解数系扩张 $\mathbb{N}\to\mathbb{Z}\to\mathbb{Q}\to\mathbb{R}\to\mathbb{C}$ 每一步都在为运算封闭与几何变换解锁新维度；
  \item 会用单位复数 $z(\theta)=\cos\theta+i\sin\theta$ 验证 $z(\alpha)z(\beta)=z(\alpha+\beta)$，理解「加法变乘法」是指数函数的基因；
  \item 会用切线速度 $(e^{k\theta})'\big|_{\theta=0}=k$ 锁定底数，写出欧拉公式 $e^{i\theta}=\cos\theta+i\sin\theta$，并得出 $e^{i\pi}+1=0$；
  \item 会用 $e^{i(\alpha+\beta)}$ 展开、棣莫弗定理与指数分解，现场推导两角和差、二倍角三倍角与和差化积／积化和差；
  \item 掌握正弦定理（外接圆半径 $R$）、余弦定理与面积公式，会按「一次齐次用正弦、二次齐次用余弦」做边角互化，并识别钝角三角形中的隐藏定义域。
\end{itemize}
\end{learnobjective}

\begin{preview}
\subsection*{知识回顾}
\begin{itemize}
  \item 实数轴上，乘以正数 $2$ 的几何意义是 \blank（拉伸／旋转）；乘以 $-1$ 的几何意义是 \blank；
  \item $(-1)\times(-1)=$ \blank；
  \item 指数函数 $a^x$ 满足的运算律是 $a^{\alpha+\beta}=$ \blank；
  \item 单位圆上辐角 $\theta$ 对应的点是 $(\cos\theta,\sin\theta)$，它到原点的距离是 \blank。
\end{itemize}

\subsection*{课前速读}
阅读三角函数一章及本章导学案，完成下列问题：
\begin{enumerate}
  \item $i^2=$ \blank；
  \item 欧拉公式 $e^{i\theta}=$ \blank；
  \item 正弦定理 $\dfrac{a}{\sin A}=$ \blank（用外接圆半径 $R$ 表示）；
  \item 余弦定理 $a^2=$ \blank。
\end{enumerate}
\end{preview}

\begin{explore}

\subsection{数系扩张与虚数单位 $i$ 的几何诞生}

\key{探究一：乘法即变换——从一维缩放、掉头到二维旋转}
在实数轴上盯住数字 $1$，看对它做乘法时的几何运动：
\begin{itemize}
  \item 乘以正数（如 $\times2$）：线段被\key{拉伸}为原来的 $2$ 倍，方向不变——纯\key{缩放}；
  \item 乘以 $-1$：把 $1$ 变成 $-1$，几何意义是绕原点逆时针旋转 \key{$180^\circ$}（半圈）；
  \item 连续乘两次 $-1$：$(-1)\times(-1)=1$，即转 $180^\circ+180^\circ=360^\circ$，恰好回到原位。
\end{itemize}
于是提出关键问题：\key{有没有一种操作，连续做两次恰好等于乘 $-1$（即一次只逆时针转 $90^\circ$）？}
设这个操作是 $i$，则
\[
  i\times i=-1\iff i^2=-1.
\]
一维数轴上根本容不下「转 $90^\circ$」后的落点，数轴被迫向垂直方向张开，出现\key{虚轴}，
一维直线跃升为二维\key{复平面}。

\begin{figure}[htbp]
  \centering
  \begin{tikzpicture}[>=Stealth, semithick, font=\footnotesize]
    % 左图：乘 -1 转 180°
    \draw[->, gray] (-2.5,0)--(2.5,0) node[below right] {实轴};
    \fill (0,0) circle (1pt);
    \draw[->, main!80!black, thick] (0,0)--(1,0) node[midway,above] {$1$};
    \draw[second, thick] (1,0) arc (0:180:1);
    \draw[->, second, thick] (0,0)--(-1,0) node[midway,above] {$-1$};
    \node[second] at (0,1.35) {$180^\circ$};
    \node[main!70!black] at (0,-2.05) {乘 $-1$};
    % 右图：乘 i 转 90°，张开虚轴
    \begin{scope}[xshift=6.2cm]
      \draw[->, gray] (-1.4,0)--(2.5,0) node[below right] {实轴};
      \draw[->, dashed, gray!70!black] (0,-1.7)--(0,2.3) node[above] {虚轴};
      \fill (0,0) circle (1pt);
      \draw[->, main!80!black, thick] (0,0)--(1,0) node[midway,below] {$1$};
      \draw[third, thick] (1,0) arc (0:90:1);
      \draw[->, third, thick] (0,0)--(0,1) node[midway,right] {$i$};
      \draw[dashed, gray] (0,1) arc (90:180:1);
      \draw[->, dashed, gray] (0,0)--(-1,0) node[midway,below] {$i^2=-1$};
      \node[third] at (0.62,0.30) {$90^\circ$};
      % 说明文字放在虚轴下端之下（虚轴只画到 $-1.7$），避免被竖直虚线穿过
      \node[main!70!black] at (0,-2.05) {乘 $i$（连乘两次得 $-1$）};
    \end{scope}
  \end{tikzpicture}
  \caption{乘法的几何：乘 $-1$ 逆时针转 $180^\circ$；乘 $i$ 只转 $90^\circ$，一维数轴被迫张开虚轴，成为复平面}
  \label{fig:rotate-i}
\end{figure}

点睛：实数 $1$ 负责向前走，$-1$ 负责向后转，$i$ 负责向左转；
连喊两次「向左转」正好就是「向后转」（$i^2=-1$）。

乘 $i$ 的作用不止于「转两次得 $-1$」：对任意辐角 $\theta$ 的单位向量，乘 $i$ 都把它的两个分量整体倒置（图~\ref{fig:i-component-flip}）。
\begin{figure}[htbp]
  \centering
  \begin{tikzpicture}[>=Stealth, semithick, font=\footnotesize]
    % ===== 左面板：转前 z=cosθ+isinθ =====
    \draw[->, gray] (-1.9,0)--(1.9,0) node[below right] {实轴};
    \draw[->, gray] (0,-0.9)--(0,1.95) node[above] {虚轴};
    \draw[gray!55, dashed] (0,0) circle (1.6);
    \draw[black!70, thick] (0,0)--(40:1.6) node[above right, font=\scriptsize] {$z$};
    \draw[blue!70!black, very thick] (0,0)--({1.6*cos(40)},0);
    \node[blue!70!black, below] at ({0.8*cos(40)},0) {$\cos\theta$};
    \draw[green!55!black, very thick] ({1.6*cos(40)},0)--(40:1.6);
    \node[green!55!black, left] at ({1.6*cos(40)+0.02},{0.8*sin(40)}) {$\sin\theta$};
    % 直角符号：两个投影分量互相垂直
    \draw[black!60, thin] ({1.6*cos(40)-0.22},0)--({1.6*cos(40)-0.22},0.22)--({1.6*cos(40)},0.22);
    \draw[black!70, thick] (0.42,0) arc (0:40:0.42) node[midway, above, font=\scriptsize] {$\theta$};
    \node[main!70!black] at (0,-0.7) {$z=\cos\theta+i\sin\theta$};
    % ===== 右面板：乘 i 后 iz=-sinθ+icosθ =====
    \begin{scope}[xshift=4.4cm]
      \draw[->, gray] (-1.9,0)--(1.9,0) node[below right] {实轴};
      \draw[->, gray] (0,-0.9)--(0,1.95) node[above] {虚轴};
      \draw[gray!55, dashed] (0,0) circle (1.6);
      \draw[dashed, gray!70!black, thick] (0,0)--(40:1.6);
      % 直角符号：z 与 iz 垂直（随 θ 取向）
      \draw[third, thin] ({0.24*cos(40)},{0.24*sin(40)})--({0.24*cos(40)-0.24*sin(40)},{0.24*sin(40)+0.24*cos(40)})--({-0.24*sin(40)},{0.24*cos(40)});
      % 90° 弧：从 θ 扫到 θ+90°
      \draw[->, third, thick] (40:0.5) arc (40:130:0.5);
      \node[third, left] at ({0.72*cos(95)},{0.72*sin(95)}) {$90^\circ$};
      \draw[->, second, very thick] (0,0)--({1.6*cos(130)},{1.6*sin(130)}) node[above left, font=\scriptsize] {$iz$};
      % 倒置后的两个分量：蓝竖（cosθ）、绿横（-sinθ）
      \draw[blue!70!black, very thick] (0,0)--(0,{1.6*cos(40)});
      \node[blue!70!black, right] at (0,{0.8*cos(40)}) {$\cos\theta$};
      \draw[green!55!black, very thick] (0,{1.6*cos(40)})--({-1.6*sin(40)},{1.6*cos(40)});
      \node[green!55!black, below] at ({-0.8*sin(40)},{1.6*cos(40)}) {$-\sin\theta$};
      \draw[dashed, gray] ({1.6*cos(130)},{1.6*sin(130)})--({1.6*cos(130)},0);
      \node[main!70!black] at (0,-0.7) {$iz=-\sin\theta+i\cos\theta$};
    \end{scope}
  \end{tikzpicture}
  \caption{乘 $i$ 的分量倒置：$z=\cos\theta+i\sin\theta$（横蓝 $\cos\theta$、纵绿 $\sin\theta$）乘 $i=e^{i\pi/2}$ 逆时针转 $90^\circ$ 后，蓝分量竖到虚轴成 $\cos\theta$、绿分量倒向负实轴成 $-\sin\theta$，即 $iz=-\sin\theta+i\cos\theta$}
  \label{fig:i-component-flip}
\end{figure}

\begin{example}
三次方程 $x^3=15x+4$ 显然有实根 $x=4$。用卡丹公式却会得到
\[
  x=\sqrt[3]{2+\sqrt{-121}}+\sqrt[3]{2-\sqrt{-121}}
  =\sqrt[3]{2+11i}+\sqrt[3]{2-11i}.
\]
请验证 $(2+i)^3=2+11i$，并说明虚部如何抵消得到 $x=4$。
\end{example}
\begin{solution}
\[
  (2+i)^3=8+12i+6i^2+i^3=8+12i-6-i=2+11i,
\]
同理 $(2-i)^3=2-11i$。于是
\[
  x=(2+i)+(2-i)=4,
\]
两个虚部一正一负抵消，还原出真实的实数解。
可见虚数不是凭空捏造的幻想，而是连接实数问题与实数解的「地下通道」：
引入 $i$ 后，$n$ 次代数方程必有 $n$ 个根（代数基本定理），数系运算彻底封闭。
\end{solution}

\begin{center}
\footnotesize
\setlength{\tabcolsep}{3pt}
\renewcommand{\arraystretch}{1.4}
\begin{tabular}{p{2.8cm}p{3.4cm}p{5.4cm}}
\toprule
扩张阶段 & 解决的方程障碍 & 解锁的几何空间与变换 \\
\midrule
$\mathbb{N}\to\mathbb{Z}$（引入负数） & $x+3=1$ 无解 & 数轴向负半轴延伸，解锁反向平移与 $180^\circ$ 掉头 \\
$\mathbb{Z}\to\mathbb{Q}$（引入分数） & $2x=1$ 无解 & 填满离散整数点之间的空隙，解锁任意比例缩放 \\
$\mathbb{Q}\to\mathbb{R}$（引入无理数） & $x^2=2$ 无解 & 消除数轴上的「针眼漏洞」，得到连续完备实数线 \\
$\mathbb{R}\to\mathbb{C}$（引入虚数 $i$） & $x^2+1=0$ 无解 & 一维直线跃升为二维平面，解锁任意角度旋转与代数封闭 \\
\bottomrule
\end{tabular}
\end{center}

\begin{conccard}[一级]{虚数单位 $i$ 的几何意义：$i^2=-1$}{理解复数的几何来源、判断复数乘法的旋转角与伸缩倍数}{把 $i$ 只当成「$\sqrt{-1}$」的代数符号，丢掉它的旋转几何；把 $i^2=-1$ 与「转 $180^\circ$」脱钩}
$1\times i$ 表示逆时针转 $90^\circ$，$1\times(-1)$ 表示转 $180^\circ$，而 $i\times i=-1$ 正是「转两次直角等于掉头」。
任何复数 $z$ 乘以 $i$，都是把 $z$ 的辐角加 $90^\circ$、模长不变。
\end{conccard}

\key{再换一个角度：每次数系扩充，保留了什么、失去了什么}
数系扩充不是「只增不减」地加数：每一步都在尽量保留旧运算律的前提下换来解题能力，同时不可避免地牺牲掉某些结构性性质。下表是四次扩充的准确清单（表中「失去」指该性质在新数系中不再整体成立）：

\begin{center}
\footnotesize
\setlength{\tabcolsep}{4pt}
\renewcommand{\arraystretch}{1.5}
\begin{tabular}{p{2.6cm}p{4.6cm}p{4.4cm}}
\toprule
扩充过程 & 保留的性质 & 不能继续保留的性质 \\
\midrule
$\mathbb{N}\to\mathbb{Z}$ & 加法、乘法及其运算律；相反数的出现使减法封闭 & 良序性（$\mathbb{N}$ 的任意非空子集有最小元；$\mathbb{Z}$ 中负整数全体就没有最小元） \\
$\mathbb{Z}\to\mathbb{Q}$ & 四则运算律与大小顺序；除法封闭、非零元都有逆元 & 离散性（不存在相邻的两个有理数，整数之间不再有空隙） \\
$\mathbb{Q}\to\mathbb{R}$ & 域结构与全序；开方、长度、极限有了合法舞台 & 可数性（$\mathbb{Q}$ 可逐个列出，$\mathbb{R}$ 不能——对角线法） \\
$\mathbb{R}\to\mathbb{C}$ & 加、减、乘、除四则运算律；代数封闭 & 与四则运算相容的大小顺序（不能像实数那样比大小） \\
\bottomrule
\end{tabular}
\end{center}

值得强调的是：前三次扩充都保留甚至加强了「能比大小」的能力，只有最后一步必须放弃它。这不是缺陷，而是换取代数封闭性的代价。

\key{为什么复数不能比较大小}
假设在 $\mathbb{C}$ 上存在一种全序 $<$，满足两条与运算相容的常规要求：($i$) $a<b\Rightarrow a+c<b+c$；($ii$) $a>0,\ b>0\Rightarrow ab>0$（于是非零数的平方必为正）。
\begin{derivation}
$i\neq0$，所以要么 $i>0$，要么 $i<0$。
\begin{itemize}
  \item 若 $i>0$：由 ($ii$) 得 $i^2>0$，但 $i^2=-1<0$，矛盾。
  \item 若 $i<0$：由 ($i$) 两边加 $-i$ 得 $0<-i$，即 $-i>0$；再由 ($ii$) 得 $(-i)^2>0$，而 $(-i)^2=i^2=-1<0$，同样矛盾。
\end{itemize}
两种情形都矛盾，故这样的全序不存在。
\end{derivation}
因此只能说「$z_1$ 与 $z_2$ 不相等」，不能说「$z_1$ 比 $z_2$ 大」；比较两个复数的唯一合法方式是比\key{模}（$|z|$ 是实数，可以比大小）或比\key{辐角}（方向）。这与实数数轴上的「谁左谁右」是完全不同的两回事。

\begin{commonerror}{把历史的顺序当成逻辑的唯一顺序}
常见讲法：「人类为了描述旋转 $90^\circ$ 才发明了 $i$。」这在数学史上不准确：虚数最早来自三次方程求根时出现的负数开平方（卡尔达诺《大术》1545 年已有记载，邦贝利随后系统研究），复平面的几何解释（把复数看成平面上的点、把乘法看成旋转伸缩）是更晚才发展出来的认识。
更严谨的教学表述是：人类先为突破实数范围内的代数求解限制而引入虚数，后来才发现它还赋予复数乘法漂亮的几何意义。
本节用「乘法即变换」作为引入路径，是\key{逻辑在先、直观优先}的教学设计，不是历史事实的顺序；讲课时应明确区分，避免学生形成「复数是为了画图方便而被发明」的错误印象。
\end{commonerror}

\begin{method}
判断一个数系够不够用的三个反射：
\begin{itemize}
  \item 方程层面：目标方程在当前数系内无解吗？（$x+3=1,\ 2x=1,\ x^2=2,\ x^2+1=0$ 各逼出一次扩充。）
  \item 运算层面：当前数系对四则运算封闭吗？对极限运算封闭吗？
  \item 结构层面：为了求解新问题，愿意放弃哪条旧性质？（序、离散性、可数性……）
\end{itemize}
\end{method}

\key{课堂讨论：数系扩充有没有终点}
\begin{itemize}
  \item \key{问题 1}：实数已经能表示数轴上的每一个点，为什么还不够用？
        （引导：完备性解决的是「连续」；方程 $x^2+1=0$ 需要的是「代数封闭」，两者是不同的要求。）
  \item \key{问题 2}：为什么复数不能比较大小？请用平方运算自己推出矛盾。
  \item \key{问题 3}：不能比大小，是不是说明复数比实数「更差」？
        （引导：失去序，换来代数封闭与二维几何；交流电路、振动、信号处理都靠复数把运算变简洁。）
\end{itemize}
给后续章节留一个伏笔：从 $\mathbb{R}$ 到 $\mathbb{C}$ 也改变了对函数的理解——实函数依赖大小顺序与数轴结构，复函数则可以借助复平面与几何变换建立新的理论。

\subsection{结构同构：单位圆旋转必为指数函数}

\key{探究二：角度相加，复数相乘}
把单位圆上辐角为 $\theta$ 的点写成复数
\[
  z(\theta)=\cos\theta+i\sin\theta.
\]
计算两个单位复数的乘积：
\begin{derivation}
\begin{align*}
  z(\alpha)z(\beta)&=(\cos\alpha+i\sin\alpha)(\cos\beta+i\sin\beta)\\
  &=(\cos\alpha\cos\beta-\sin\alpha\sin\beta)+i(\sin\alpha\cos\beta+\cos\alpha\sin\beta)\\
  &=\cos(\alpha+\beta)+i\sin(\alpha+\beta)=z(\alpha+\beta).
\end{align*}
\end{derivation}
（最后一步用了两角和公式——这正是我们稍后要用旋转重新解释的对象。）
惊人之处：\key{输入的角度相加，输出的复数相乘}。

唤醒旧知：哪类函数的核心基因是 $f(\alpha+\beta)=f(\alpha)f(\beta)$？
学生应脱口而出\key{指数函数} $f(x)=b^x$。
而 $z(0)=\cos0+i\sin0=1$，所以 $z(\theta)$ 在代数结构上就是一个底数为某复常数 $b$ 的指数函数 $z(\theta)=b^\theta$——
\key{单位圆旋转群与指数函数同构}。

\begin{conccard}[二级]{结构同构：$z(\alpha)z(\beta)=z(\alpha+\beta)$}{判断两个复数相乘的几何效果、理解三角公式的共同来源}{只记「模相乘、辐角相加」的结论，不从结构公理出发；把 $z(\theta)$ 与指数函数割裂}
单位复数 $z(\theta)=\cos\theta+i\sin\theta$ 满足「角度相加 $\to$ 复数相乘」，且 $z(0)=1$。
凡是满足 $f(\alpha+\beta)=f(\alpha)f(\beta)$ 的连续函数必为指数函数，故 $z(\theta)$ 一定可以写成 $b^\theta$ 的形式。
\end{conccard}

几何上看，$z(\alpha)\cdot e^{i\beta}$ 就是把 $z_1=z(\alpha)$ 绕原点再逆时针转 $\beta$，终点辐角为 $\alpha+\beta$（图~\ref{fig:rotate-sum}）。
\begin{figure}[htbp]
  \centering
  \begin{tikzpicture}[>=Stealth, semithick, font=\footnotesize]
    \draw[->, gray] (-1.8,0)--(3.5,0) node[below right] {$\mathrm{Re}$};
    \draw[->, gray] (0,-1.7)--(0,2.9) node[above] {$\mathrm{Im}$};
    \draw[gray!55, dashed] (0,0) circle (2.2);
    % 三段同心弧：α、β、α+β（半径与标注角错层，避免相叠）
    \draw[main, thick] (0.5,0) arc (0:20:0.5);
    \node[main] at (10:0.68) {$\alpha$};
    \draw[->, third, thick] (20:0.85) arc (20:75:0.85);
    \node[third] at (47:1.06) {$\beta$};
    \draw[->, gray!60!black, thick] (1.2,0) arc (0:75:1.2);
    \node[gray!60!black] at (12:1.45) {$\alpha+\beta$};
    % 旋转轨迹：z1 沿单位圆转到 z2
    \draw[->, third, thin] (20:2.2) arc (20:75:2.2);
    % z1=e^{iα} 与 z2=e^{i(α+β)}
    \draw[->, main!80!black, thick] (0,0)--(20:2.2) node[above right] {$z_1=e^{i\alpha}$};
    \draw[->, second, very thick] (0,0)--(75:2.2) node[above left] {$z_2=e^{i(\alpha+\beta)}$};
    % z2 的两轴投影：实部红、虚部青
    \draw[red!80!black, very thick] (0,0)--({2.2*cos(75)},0);
    \node[red!80!black, below] at ({1.1*cos(75)},0) {$\cos(\alpha+\beta)$};
    \draw[dashed, gray] (75:2.2)--(0,{2.2*sin(75)});
    \draw[cyan!80!black, very thick] ({2.2*cos(75)},0)--(75:2.2);
    \node[cyan!80!black, right] at ({2.2*cos(75)},1.2) {$\sin(\alpha+\beta)$};
    \node[main!70!black] at (0.9,-1.25) {$z_1\cdot e^{i\beta}=e^{i(\alpha+\beta)}=z_2$};
  \end{tikzpicture}
  \caption{旋转叠加：$z_1=e^{i\alpha}$ 再乘 $e^{i\beta}$ 即绕原点逆时针转 $\beta$，终点 $z_2=e^{i(\alpha+\beta)}$；$z_2$ 在实轴、虚轴上的投影 $\cos(\alpha+\beta)$、$\sin(\alpha+\beta)$ 与 $(\cos\alpha+i\sin\alpha)(\cos\beta+i\sin\beta)$ 的实部、虚部一一对应——两角和公式就是这么「转」出来的}
  \label{fig:rotate-sum}
\end{figure}

\subsection{切线速度锁定底数}

\key{探究三：切线速度——底数只能是 $e$}
对比两个「指数运动」的动力学：
\begin{itemize}
  \item 实指数 $e^{x}$：在 $x=0$ 处位置为 $1$，速度方向与位置\key{同向}（沿实轴向右），越跑越快，轨迹是直线爆炸；
  \item 复指数 $e^{i\theta}$：在 $\theta=0$ 处位置为 $1$，速度方向必须与位置\key{垂直}（沿虚轴向上，即 $+i$），只转弯不伸长，轨迹是单位圆。
\end{itemize}
把 $z(\theta)$ 看作质点位置，求导得切线速度：
\[
  z'(\theta)=-\sin\theta+i\cos\theta\implies z'(0)=i.
\]
而实指数 $f(\theta)=e^{k\theta}$ 满足
\[
  (e^{k\theta})'\Big|_{\theta=0}=k\cdot e^{0}=k.
\]
要让圆周运动在出发点的瞬时速度恰好是 $+i$，唯一可能就是 \key{$k=i$}。于是
\[
  e^{i\theta}=\cos\theta+i\sin\theta.
\]
代入 $\theta=\pi$：
\[
  e^{i\pi}=\cos\pi+i\sin\pi=-1\implies e^{i\pi}+1=0,
\]
即被誉为「上帝公式」的欧拉恒等式，把 $0,1,e,i,\pi$ 五个常数用最纯粹的圆周旋转焊在一起。

\begin{center}
\footnotesize
\setlength{\tabcolsep}{3pt}
\renewcommand{\arraystretch}{1.4}
\begin{tabular}{p{2.2cm}p{2.0cm}p{2.6cm}p{3.2cm}p{3.0cm}}
\toprule
函数 & 起点位置 & 初始速度方向 & 速度与位置关系 & 轨迹 \\
\midrule
$e^{x}$ & $1$ & $+1$（实轴向右） & 同向 & 沿实轴指数爆炸 \\
$e^{i\theta}$ & $1$ & $+i$（虚轴向上） & 垂直 & 绕单位圆匀速旋转 \\
\bottomrule
\end{tabular}
\end{center}

\begin{figure}[htbp]
  \centering
  \begin{tikzpicture}[>=Stealth, semithick, font=\footnotesize]
    \draw[->, gray] (-2.6,0)--(2.9,0) node[below right] {$\mathrm{Re}$};
    \draw[->, gray] (0,-2.3)--(0,2.5) node[above] {$\mathrm{Im}$};
    \draw[gray!55, dashed] (0,0) circle (2);
    % 一般位置 z(θ)：θ≈50°
    \draw[->, dashed, gray!80!black] (0,0)--(1.286,1.532) node[above] {$z(\theta)$};
    \draw[->, dashed, gray!80!black] (1.286,1.532)--(0.826,1.918) node[above left, font=\scriptsize] {$i\,z(\theta)$};
    % 起点 (1,0) 与切线速度 +i
    \fill (2,0) circle (1.2pt);
    \draw[->, main!80!black, thick] (0,0)--(2,0) node[midway,below] {$z(0)=1$};
    \draw[->, second, thick] (2,0)--(2,0.95) node[right] {$z'(0)=i$};
    \draw[gray, thin] (1.52,0)--(1.52,0.48)--(2,0.48);
    \draw[gray] (0.72,0) arc (0:50:0.72);
    \node[gray!80!black] at (0.98,0.22) {$\theta$};
    \node[main!70!black] at (0,-2.8) {位置向量与切线速度垂直：只转弯，不伸长};
  \end{tikzpicture}
  \caption{复平面单位圆：起点 $(1,0)$ 处位置向量为 $1$，切线速度 $+i$ 与其垂直；任意位置的速度 $i\,z(\theta)$ 始终垂直于位置}
  \label{fig:euler-tangent}
\end{figure}

\begin{conccard}[一级]{欧拉公式：$e^{i\theta}=\cos\theta+i\sin\theta$}{推导三角公式、判断复数旋转角、快速验算}{只把它当结论背诵；忘记「乘以 $e^{i\beta}$ 就是逆时针转 $\beta$」这一几何含义}
由 $z(\alpha)z(\beta)=z(\alpha+\beta)$ 与 $z(0)=1$ 知 $z(\theta)$ 必为指数函数 $b^\theta$；
再由切线速度 $(e^{k\theta})'\big|_{\theta=0}=k$ 与 $z'(0)=i$ 锁定 $k=i$。
于是乘以 $e^{i\beta}$ 等价于逆时针旋转 $\beta$，这就是本章全部三角公式的母公理。
\end{conccard}

\subsection{母公式引擎：一条指数律长出全套公式}

\key{探究四：$e^{i(\alpha+\beta)}=e^{i\alpha}e^{i\beta}$ 的展开与投影}

\key{（1）诱导公式——乘 $i=e^{i\pi/2}$ 即转 $90^\circ$。}
\[
  i(\cos\theta+i\sin\theta)=-\sin\theta+i\cos\theta,
\]
对比实部虚部，一眼读出
\[
  \cos\left(\theta+\frac{\pi}{2}\right)=-\sin\theta,\qquad
  \sin\left(\theta+\frac{\pi}{2}\right)=\cos\theta.
\]
同理乘 $-1=e^{i\pi}$ 得 $\cos(\theta+\pi)=-\cos\theta$、$\sin(\theta+\pi)=-\sin\theta$；
取共轭 $\overline{e^{i\theta}}=e^{-i\theta}$ 即关于实轴对称，这正是奇偶性
$\cos(-\theta)=\cos\theta$、$\sin(-\theta)=-\sin\theta$ 的几何本质。
「奇变偶不变，符号看象限」不再是口诀，而是「转了几个直角」的图形事实。

\key{（2）两角和与差——多项式展开一步到位。}
\begin{derivation}
\begin{align*}
  e^{i(\alpha+\beta)}&=e^{i\alpha}e^{i\beta}
  =(\cos\alpha+i\sin\alpha)(\cos\beta+i\sin\beta)\\
  &=(\cos\alpha\cos\beta-\sin\alpha\sin\beta)
  +i(\sin\alpha\cos\beta+\cos\alpha\sin\beta).
\end{align*}
\end{derivation}
左边又是 $\cos(\alpha+\beta)+i\sin(\alpha+\beta)$，对比实部与虚部：
\[
  \cos(\alpha+\beta)=\cos\alpha\cos\beta-\sin\alpha\sin\beta,\qquad
  \sin(\alpha+\beta)=\sin\alpha\cos\beta+\cos\alpha\sin\beta.
\]
把 $\beta$ 换成 $-\beta$（用 $\cos(-\beta)=\cos\beta$、$\sin(-\beta)=-\sin\beta$），立刻得到差角公式：
\[
  \cos(\alpha-\beta)=\cos\alpha\cos\beta+\sin\alpha\sin\beta,\qquad
  \sin(\alpha-\beta)=\sin\alpha\cos\beta-\cos\alpha\sin\beta.
\]

\key{（3）二倍角与三倍角——棣莫弗定理。}
由 $e^{in\theta}=(e^{i\theta})^n$ 得
\[
  (\cos\theta+i\sin\theta)^n=\cos n\theta+i\sin n\theta.
\]
$n=2$ 展开即得二倍角；$n=3$ 用二项式定理：
\[
  (\cos\theta+i\sin\theta)^3
  =(\cos^3\theta-3\cos\theta\sin^2\theta
  +i(3\cos^2\theta\sin\theta-\sin^3\theta),
\]
于是
\[
  \cos3\theta=4\cos^3\theta-3\cos\theta,\qquad
  \sin3\theta=3\sin\theta-4\sin^3\theta.
\]

\key{（4）和差化积与积化和差——同一指数分解。}
由欧拉公式解出
\[
  \cos\theta=\frac{e^{i\theta}+e^{-i\theta}}{2},\qquad
  \sin\theta=\frac{e^{i\theta}-e^{-i\theta}}{2i}.
\]
例如
\begin{align*}
  \cos\alpha\cos\beta
  &=\frac{(e^{i\alpha}+e^{-i\alpha})(e^{i\beta}+e^{-i\beta})}{4}\\
  &=\frac{e^{i(\alpha+\beta)}+e^{-i(\alpha+\beta)}+e^{i(\alpha-\beta)}+e^{-i(\alpha-\beta)}}{4}\\
  &=\frac{\cos(\alpha+\beta)+\cos(\alpha-\beta)}{2},
\end{align*}
这就是积化和差；把 $\alpha+\beta=A$、$\alpha-\beta=B$ 反解即得和差化积。
结论：和差、乘积、倍角只是同一组指数幂的不同代数外形。

\begin{example}
用棣莫弗定理把 $\cos3\theta$ 用 $\cos\theta$ 表示。
\end{example}
\begin{solution}
\[
  \cos3\theta=\operatorname{Re}(\cos\theta+i\sin\theta)^3
  =\cos^3\theta-3\cos\theta\sin^2\theta
  =\cos^3\theta-3\cos\theta(1-\cos^2\theta)
  =4\cos^3\theta-3\cos\theta.
\]
\end{solution}

\begin{example}
求 $\cos\theta+\cos2\theta+\cdots+\cos n\theta$（$\theta\neq2k\pi$）。
\end{example}
\begin{solution}
令 $q=e^{i\theta}$，则 $\cos k\theta=\operatorname{Re}q^k$，于是
\[
  S=\sum_{k=1}^{n}q^k=\frac{q(1-q^n)}{1-q}
  =\frac{e^{i\theta}(1-e^{in\theta})}{1-e^{i\theta}}
  =e^{i\frac{(n+1)\theta}{2}}\cdot\frac{\sin\frac{n\theta}{2}}{\sin\frac{\theta}{2}}.
\]
取实部即得
\[
  \cos\theta+\cos2\theta+\cdots+\cos n\theta
  =\frac{\sin\frac{n\theta}{2}}{\sin\frac{\theta}{2}}\cos\frac{(n+1)\theta}{2}.
\]
实数域里要凑 $\sin\frac{\theta}{2}$ 裂项，技巧性强；复数域里只是公比为 $q=e^{i\theta}$ 的等比数列求和。
\end{solution}

\begin{conccard}[一级]{考纲安全边界：复数法用于选填秒杀，解答题须还原为实数格式}{选择填空快速求值、验算；忘记公式时现场推导；快速判断符号}{在解答题中直接写 $e^{i\theta}$、$i$ 的推导过程而跳过实数三角变换步骤，造成阅卷扣分}
高考解答题默认在实数域 $\mathbb{R}$ 内评分。复数法与欧拉公式的正确用法是：
选填小题秒杀、快速验算、忘记公式时现场推导。
解答题书写必须还原为实数三角变换格式（和差倍半公式、正弦定理、余弦定理的标准步骤），
不能把 $e^{i\theta}$ 与 $i$ 直接写进解答题过程。
\end{conccard}

\begin{commonerror}{指数展开时丢符号：忘了 $i^2=-1$}
用 $e^{i(\alpha+\beta)}$ 展开时，交叉项来自 $i\sin\alpha\cdot i\sin\beta=i^2\sin\alpha\sin\beta=-\sin\alpha\sin\beta$。
若忘记 $i^2=-1$，就会把实部写成 $\cos\alpha\cos\beta+\sin\alpha\sin\beta$，得到错误的「和角公式」。
记住：\key{每出现一个 $i^2$ 就要换成 $-1$}，符号自然正确。
\end{commonerror}

\begin{method}
母公式引擎使用流程（复数法三步）：
\begin{enumerate}
  \item \key{翻译}：把三角式写成 $e^{i\theta}$ 的指数形式（$\cos\theta=\frac{e^{i\theta}+e^{-i\theta}}2$，$\sin\theta=\frac{e^{i\theta}-e^{-i\theta}}{2i}$）；
  \item \key{指数运算}：用 $e^{i(\alpha+\beta)}=e^{i\alpha}e^{i\beta}$、棣莫弗定理、等比数列求和等纯代数工具计算；
  \item \key{取实虚部还原}：需要余弦取实部，需要正弦取虚部，最后写回实数三角形式。
\end{enumerate}
\end{method}

\subsection{解三角形与边角互化}

\key{探究五：正余弦定理——把三角形放进外接圆}
正弦定理：
\[
  \frac{a}{\sin A}=\frac{b}{\sin B}=\frac{c}{\sin C}=2R,
\]
其中 $R$ 是 $\triangle ABC$ 外接圆半径。几何意义：作外接圆，由圆周角定理得 $a=2R\sin A$，其余同理。
余弦定理：
\[
  a^2=b^2+c^2-2bc\cos A
  \iff \cos A=\frac{b^2+c^2-a^2}{2bc}.
\]
面积公式：
\[
  S=\frac12bc\sin A=\frac12ca\sin B=\frac12ab\sin C=\frac{abc}{4R}.
\]

\key{边角互化的选择标准。}
\begin{itemize}
  \item 关于 $a,b,c$（或 $\sin A,\sin B,\sin C$）的\key{一次齐次式} $\to$ 用正弦定理，把边化角或把角化边；
  \item 关于 $a^2,b^2,c^2$（或 $\sin^2$）的\key{二次齐次式} $\to$ 用余弦定理。
\end{itemize}

\begin{example}
在 $\triangle ABC$ 中，已知 $a=2b\cos C$，判断三角形的形状。
\end{example}
\begin{solution}
由正弦定理把边化角：$\sin A=2\sin B\cos C$。
又 $A=\pi-(B+C)$，故 $\sin A=\sin(B+C)=\sin B\cos C+\cos B\sin C$，于是
\[
  \sin B\cos C+\cos B\sin C=2\sin B\cos C
  \implies \cos B\sin C-\sin B\cos C=0
  \implies \sin(C-B)=0.
\]
因为 $B,C\in(0,\pi)$，所以 $C-B=0$，即 $B=C$。故三角形为等腰三角形。
（一次齐次式，选正弦定理。）
\end{solution}

\begin{example}
在 $\triangle ABC$ 中，已知 $a^2=b^2+c^2-bc$，求角 $A$。
\end{example}
\begin{solution}
与余弦定理 $a^2=b^2+c^2-2bc\cos A$ 对比，得 $2\cos A=1$，即 $\cos A=\frac12$。
又 $A\in(0,\pi)$，故 $A=\dfrac{\pi}{3}$（$60^\circ$）。
（二次齐次式，选余弦定理。）
\end{solution}

\key{钝角三角形的隐藏定义域。}
当题目给出「钝角三角形」或隐含某角为钝角时，角的范围被额外压缩：
钝角 $\gamma\in\left(\dfrac{\pi}{2},\pi\right)$，两个锐角满足 $\alpha+\beta<\dfrac{\pi}{2}$。
用余弦定理处理时，$\cos\gamma<0$ 给出边的平方约束，而三角形存在条件给出另一侧约束，两者共同构成隐藏定义域。

\begin{example}
在 $\triangle ABC$ 中，已知 $a=2$，$b=3$，且 $C$ 为钝角，求边 $c$ 的取值范围。
\end{example}
\begin{solution}
由余弦定理
\[
  \cos C=\frac{a^2+b^2-c^2}{2ab}=\frac{13-c^2}{12}.
\]
$C$ 为钝角 $\iff \cos C<0\iff c^2>13\iff c>\sqrt{13}$。
同时三角形存在要求 $c<a+b=5$（且 $c>|b-a|=1$，已被 $c>\sqrt{13}$ 包含）。
故 $c\in(\sqrt{13},5)$。
若只写 $c^2>13$，就漏掉了三角形存在给出的上界 $c<5$——这正是「钝角三角形」隐藏的定义域。
\end{solution}

\begin{conccard}[二级]{边角互化：一次齐次用正弦，二次齐次用余弦}{解三角形中的求值、求角、判断形状、求范围}{齐次次数判断错误；用正弦定理求角时忽略两解；忽略钝角三角形对边角范围的额外约束}
正弦定理 $\frac{a}{\sin A}=2R$ 把边与角的正弦互相转化，适合一次齐次式；
余弦定理 $a^2=b^2+c^2-2bc\cos A$ 把三边与一角联系起来，适合二次齐次式。
涉及「钝角三角形」时，必须同时写出 $\cos\gamma<0$ 与三角形存在条件，得到完整的定义域。
\end{conccard}

\begin{commonerror}{解三角形忽略隐藏定义域}
两类常见失误：
\begin{itemize}
  \item 用正弦定理由 $\sin B=\frac{b\sin A}{a}$ 求角时，只取锐角解，漏掉钝角解；
  \item 题目说「钝角三角形」却只列 $\cos\gamma<0$，漏掉三角形存在条件（两边之和大于第三边）。
\end{itemize}
例如由 $a=2,b=3$ 且 $C$ 为钝角求 $c$ 时，只写 $c>\sqrt{13}$ 会漏掉 $c<5$。
正确做法是把「角的范围」与「三角形存在条件」同时列出，求交集。
\end{commonerror}

\begin{summary}
本章以「乘以 $e^{i\beta}=$ 逆时针旋转 $\beta$」为唯一母公理，把三角函数从零散公式还原为同一个旋转群的三种语言：
数系扩张 $\mathbb{N}\to\mathbb{Z}\to\mathbb{Q}\to\mathbb{R}\to\mathbb{C}$ 让 $i^2=-1$（转 $90^\circ$ 两次等于掉头）成为二维复平面的入场券；
单位复数 $z(\theta)=\cos\theta+i\sin\theta$ 满足 $z(\alpha)z(\beta)=z(\alpha+\beta)$ 且 $z(0)=1$，故必为指数函数；
由切线速度 $(e^{k\theta})'\big|_{\theta=0}=k$ 与 $z'(0)=i$ 锁定 $k=i$，得到欧拉公式 $e^{i\theta}=\cos\theta+i\sin\theta$ 与 $e^{i\pi}+1=0$；
再用 $e^{i(\alpha+\beta)}$ 展开、棣莫弗定理与 $\cos\theta=\frac{e^{i\theta}+e^{-i\theta}}2$、$\sin\theta=\frac{e^{i\theta}-e^{-i\theta}}{2i}$，
一次性长出两角和差、二倍角三倍角、诱导公式与和差化积／积化和差。
解三角形部分用正弦定理（外接圆半径 $R$）、余弦定理与面积公式，按「一次齐次用正弦、二次齐次用余弦」做边角互化，并警惕钝角三角形的隐藏定义域。
考纲边界：复数法用于选填秒杀与验算，解答题必须还原为实数三角变换格式。
\end{summary}

\end{explore}

\begin{inclasstest}
\begin{enumerate}
  \item 计算 $i\cdot(1+i)$，并说明它在复平面上表示怎样的旋转与伸缩；\level{易}
  \item 用 $e^{i(\alpha+\beta)}=e^{i\alpha}e^{i\beta}$ 展开并对比实部与虚部，写出 $\cos(\alpha-\beta)$ 与 $\sin(\alpha-\beta)$；\level{中}
  \item 用棣莫弗定理把 $\cos3\theta$ 用 $\cos\theta$ 表示；\level{中}
  \item 在 $\triangle ABC$ 中，$a=2$，$b=3$，$C=120^\circ$，求边 $c$ 与面积 $S$。\level{中}
\end{enumerate}
\end{inclasstest}

\begin{afterclass}
\subsection*{基础巩固}
\begin{enumerate}
  \item 计算 $i^3$、$i^4$ 与 $(1+i)^2$，并说出各自的几何意义；
  \item 由 $e^{i\theta}=\cos\theta+i\sin\theta$ 解出 $\cos\theta,\sin\theta$ 的指数表达式，并验证 $\cos^2\theta+\sin^2\theta=1$。
\end{enumerate}

\subsection*{进阶提升}
\begin{enumerate}
  \item 用 $i=e^{i\pi/2}$ 推导 $\cos\left(\theta-\dfrac{\pi}{2}\right)$ 与 $\sin\left(\theta-\dfrac{\pi}{2}\right)$；
  \item 在 $\triangle ABC$ 中，$a=2$，$b=3$，$c=4$，判断三角形的形状并求面积。
\end{enumerate}

\subsection*{拓展探究}
\begin{enumerate}
  \item 用复数等比数列求 $\cos\theta+\cos2\theta+\cdots+\cos n\theta$（$\theta\neq2k\pi$），写出结果并说明它为什么只是「取实部」。
\end{enumerate}

\subsection*{教师版答案}
\begin{solution}
\textbf{基础 1：}$i^3=i^2\cdot i=-i$（逆时针转 $270^\circ$）；$i^4=(i^2)^2=1$（转 $360^\circ$ 回到原位）；
$(1+i)^2=1+2i+i^2=2i$（模由 $\sqrt2$ 变为 $2$，辐角由 $45^\circ$ 变为 $90^\circ$）。
\end{solution}
\begin{solution}
\textbf{基础 2：}由 $e^{i\theta}=\cos\theta+i\sin\theta$、$e^{-i\theta}=\cos\theta-i\sin\theta$ 解得
\[
  \cos\theta=\frac{e^{i\theta}+e^{-i\theta}}{2},\qquad
  \sin\theta=\frac{e^{i\theta}-e^{-i\theta}}{2i}.
\]
于是
\[
  \cos^2\theta+\sin^2\theta
  =\frac{(e^{i\theta}+e^{-i\theta})^2}{4}+\frac{(e^{i\theta}-e^{-i\theta})^2}{(2i)^2}
  =\frac{e^{2i\theta}+2+e^{-2i\theta}-e^{2i\theta}+2-e^{-2i\theta}}{4}=1.
\]
\end{solution}
\begin{solution}
\textbf{进阶 1：}乘以 $-i=e^{-i\pi/2}$（顺时针转 $90^\circ$）：
\[
  (\cos\theta+i\sin\theta)(-i)=\sin\theta-i\cos\theta,
\]
对比实虚部得
\[
  \cos\left(\theta-\frac{\pi}{2}\right)=\sin\theta,\qquad
  \sin\left(\theta-\frac{\pi}{2}\right)=-\cos\theta.
\]
\end{solution}
\begin{solution}
\textbf{进阶 2：}最大边为 $c=4$，
\[
  \cos C=\frac{a^2+b^2-c^2}{2ab}=\frac{4+9-16}{12}=-\frac14<0,
\]
故 $C$ 为钝角，三角形是钝角三角形。
\[
  \sin C=\sqrt{1-\frac1{16}}=\frac{\sqrt{15}}4,\qquad
  S=\frac12ab\sin C=\frac12\cdot2\cdot3\cdot\frac{\sqrt{15}}4=\frac{3\sqrt{15}}4.
\]
\end{solution}
\begin{solution}
\textbf{拓展：}令 $q=e^{i\theta}$，则
\[
  \sum_{k=1}^{n}q^k=\frac{q(1-q^n)}{1-q}
  =\frac{e^{i\theta}(1-e^{in\theta})}{1-e^{i\theta}}
  =e^{i\frac{(n+1)\theta}{2}}\cdot\frac{\sin\frac{n\theta}{2}}{\sin\frac{\theta}{2}}.
\]
取实部得
\[
  \cos\theta+\cos2\theta+\cdots+\cos n\theta
  =\frac{\sin\frac{n\theta}{2}}{\sin\frac{\theta}{2}}\cos\frac{(n+1)\theta}{2}.
\]
因为 $\cos k\theta=\operatorname{Re}(e^{ik\theta})$，所以它只是等比数列求和后「取实部」。
\end{solution}
\reading{拓展阅读 08}{复数与三角的大一统}
\end{afterclass}

% 高观点导学案：平面向量与空间向量（第 3 章·高观点统一版）
% 主线：既有大小又有方向的量 → 点集的位移 → 基底下的坐标 → 内积投影算子 → 空间中的方向与法向
% 四模块：基底分解与坐标表示 → 数量积与投影算子 → 极化恒等式与面积 → 空间向量（方向与法向）
% 全程突出向量体系的独立地位：向量不是三角／几何的附庸，而是「几何空间代数化」的通用线性语言。

\section{平面向量与空间向量（高观点统一版）}

\begin{methodbox}{多框架对读}
\small
\begin{center}
\begin{tabular}{c|p{2.9cm}|p{5.0cm}|p{3.1cm}}
\toprule
视角 & 本章看点 & 别的视角怎么说 & 哪侧最省力 \\
\midrule
代数 & 坐标运算、数量积公式 & 有序数对的算术 & 硬算定坐标 \\
几何 & 分解、投影、法向量 & 图形中看方向与垂直 & 找关系、省计算 \\
集合 & 向量＝有向线段构成的集合 & 基底＝给集合装坐标系 & 界线性相关无关 \\
映射 & 点 $\to$ 坐标是同构；投影是算子 & 换基底＝换表示 & 整体搬迁结论 \\
\bottomrule
\end{tabular}
\end{center}
卡的使用：几何题先问「有没有现成的算子（投影／法向量）」，再决定建系还是综合（总图见第 0 章图~\ref{fig:four-view-map}）。
\end{methodbox}

第 1 章研究「对应」（函数），第 2 章研究「旋转」（三角函数）；本章登场的向量是第三条独立主线：
它研究\key{位移的合成与分解}。向量体系相对三角与函数具有显性独立地位——
三角函数是旋转群的语言，函数是对应的语言，而向量是\key{几何空间代数化}的通用线性语言：
不依赖角度公式、不依赖函数图象，只用「大小$+$方向」「基底分解」「数量积」三件工具，
就能把平面与空间中的距离、夹角、平行、垂直、面积全部翻译成代数。
四个模块层层展开：
\key{基底分解与坐标表示}（任一向量 $=$ 基底线性组合，唯一）$\to$
\key{数量积与投影算子}（数量积 $=$ 余弦定理的代数化）$\to$
\key{极化恒等式与面积}（共起点三向量乘积的秒杀通道）$\to$
\key{空间向量：方向与法向}（用法向量统一处理空间位置关系）。
核心思想：\key{选定一组基底，几何就变成了线性代数}。

\begin{learnobjective}
\begin{itemize}
  \item 理解向量「既有大小又有方向」的双重属性，区分向量与数量，理解自由向量（可平移）与零向量；
  \item 理解平面向量基本定理：不共线向量张成整个平面，任一向量有唯一分解；坐标即有序实数对，与第 1 章「一一对应」呼应；
  \item 掌握数量积的定义、坐标运算与投影公式，能把余弦定理与数量积互译（数量积就是余弦定理的代数化）；
  \item 掌握极化恒等式 $\overrightarrow{AB}\cdot\overrightarrow{AC}=|\overrightarrow{AM}|^2-|\overrightarrow{MB}|^2$（$M$ 为 $BC$ 中点），会用它秒杀共起点三向量乘积的最值问题；知道面积形式 $S=\frac12|\vec a\times\vec b|$；
  \item 会把空间向量坐标化，理解平面法向量、直线方向向量，用法向量统一处理线面／面面的平行、垂直与夹角（线面角、二面角）。
\end{itemize}
\end{learnobjective}

\begin{preview}
\subsection*{知识回顾}
\begin{itemize}
  \item 实数加法满足交换律：$a+b=$ \blank；类比思考：向量加法满足交换律吗？\blank；
  \item 平面直角坐标系中，点 $P$ 的坐标是 \blank（有序实数对／无序数组），它体现了点与 \blank 的一一对应；
  \item 两直线（斜率存在时）垂直的判定：斜率之积等于 \blank；
  \item 三角形面积公式 $S=\dfrac12\times$底$\times$ \blank；平行四边形面积 $=$ 底 $\times$ \blank。
\end{itemize}

\subsection*{课前速读}
阅读「平面向量」与「空间向量」两章内容，完成下列问题：
\begin{enumerate}
  \item 平面向量基本定理要求两个基底向量 \blank（不共线／相等）；
  \item 设 $\vec a=(x_1,y_1)$，$\vec b=(x_2,y_2)$，则 $\vec a\cdot\vec b=$ \blank；
  \item 与平面 \blank 的向量叫做该平面的法向量；
  \item 设直线方向向量为 $\vec s$、平面法向量为 $\vec n$，线面角 $\varphi$ 满足 $\sin\varphi=$ \blank。
\end{enumerate}
\end{preview}

\begin{explore}

\subsection{基底分解：平面向量基本定理与坐标表示}

\key{探究一：任一向量都能被基底「唯一拼装」出来}

\key{（1）从「既有大小又有方向的量」到「点集的位移」。}
向量是既有大小又有方向的量，用有向线段表示；
两个向量只要大小相等、方向相同就叫\key{相等向量}，与起点位置无关（\key{自由向量}，可平移）。
位移视角：$\overrightarrow{AB}$ 表示点 $A$ 到点 $B$ 的位移。把所有起点平移到原点 $O$，
向量 $\overrightarrow{OP}$ 与平面点 $P$ 一一对应——这就是「向量 $=$ 点集的位移」的准确含义。
规定：零向量 $\vec 0$（长度为 $0$，方向任意）；$\vec a$ 的相反向量记 $-\vec a$。

\key{（2）平面向量基本定理：不共线的两个向量张成整个平面。}
\begin{derivation}
定理内容：如果 $\vec e_1,\vec e_2$ 是同一平面内的两个\key{不共线}向量，
那么对这一平面内的任一向量 $\vec a$，存在\key{唯一}一对实数 $\lambda_1,\lambda_2$，使得
\[
  \vec a=\lambda_1\vec e_1+\lambda_2\vec e_2.
\]
存在性：过 $\vec a$ 的终点分别作 $\vec e_2,\vec e_1$ 方向的平行线，围成平行四边形，
两边即 $\lambda_1\vec e_1$ 与 $\lambda_2\vec e_2$（平行四边形法则）。
唯一性：设 $\vec a=\lambda_1\vec e_1+\lambda_2\vec e_2=\mu_1\vec e_1+\mu_2\vec e_2$，则
\[
  (\lambda_1-\mu_1)\vec e_1+(\lambda_2-\mu_2)\vec e_2=\vec 0.
\]
若 $\lambda_1\neq\mu_1$，则 $\vec e_1=\dfrac{\mu_2-\lambda_2}{\lambda_1-\mu_1}\vec e_2$，
与 $\vec e_1,\vec e_2$ 不共线矛盾；故 $\lambda_1=\mu_1$，同理 $\lambda_2=\mu_2$。唯一性得证。
\end{derivation}
$\{\vec e_1,\vec e_2\}$ 称为一组\key{基底}。注意基底不唯一、也不要求垂直——
这是与「坐标轴」的本质区别，也是「选得巧，计算就少」的前提。

\key{（3）坐标：基底下的有序实数对。}
取\key{正交单位基底} $\vec i,\vec j$（互相垂直、长度都为 $1$），则
\[
  \vec a=x\vec i+y\vec j \iff \vec a=(x,y),
\]
$(x,y)$ 就是 $\vec a$ 的\key{坐标}。当起点在原点时，向量 $(x,y)$ 恰好指向点 $(x,y)$——
几何对象（向量／点）、代数对象（有序实数对）、对应关系三者同构，
这正是第 1 章「笛卡尔积坐标化」思想的延续：\key{坐标 $=$ 向量在基底下的读数}。
坐标运算规则：
\[
  \vec a\pm\vec b=(x_1\pm x_2,\,y_1\pm y_2),\qquad
  \lambda\vec a=(\lambda x_1,\,\lambda y_1),\qquad
  |\vec a|=\sqrt{x_1^2+y_1^2}.
\]

\begin{figure}[htbp]
  \centering
  \begin{tikzpicture}[>=Stealth, semithick, font=\footnotesize]
    % ===== 左：斜基底的一般分解 =====
    \coordinate (O) at (0,0);
    \coordinate (E1) at (2,0.5);
    \coordinate (E2) at (0.6,1.6);
    \coordinate (P1) at (4,1);
    \coordinate (P2) at (4.6,2.6);
    \draw[->, gray] (-0.5,0)--(5,0);
    \draw[->, gray] (0,-0.5)--(0,3);
    \fill (O) circle (1pt);
    \draw[->, third, thick] (O)--(E1) node[right] {$\vec{e}_1$};
    \draw[->, third, thick] (O)--(E2) node[above left] {$\vec{e}_2$};
    \draw[->, main!80!black, very thick] (O)--(P2) node[above right] {$\vec{a}=x\vec{e}_1+y\vec{e}_2$};
    \draw[->, second, thick] (O)--(P1) node[midway,below] {$x\vec{e}_1$};
    \draw[->, second, thick] (P1)--(P2) node[midway,right] {$y\vec{e}_2$};
    \draw[dashed, second!70!black] (E2)--(P2);
    \node[main!70!black] at (0,-0.85) {一般基底：平行四边形法则唯一分解};
    % ===== 右：正交单位基底与坐标 =====
    \begin{scope}[xshift=6.6cm]
      \draw[->, gray] (-0.5,0)--(4.2,0) node[below right] {$x$};
      \draw[->, gray] (0,-0.5)--(0,3.2) node[above left] {$y$};
      \foreach \x in {1,2,3} \draw[gray!25, very thin] (\x,-0.4)--(\x,3);
      \foreach \y in {1,2} \draw[gray!25, very thin] (-0.4,\y)--(4,\y);
      \fill (0,0) circle (1pt);
      \draw[->, third, thick] (0,0)--(1,0) node[below] {$\vec{i}=\vec{e}_1$};
      \draw[->, third, thick] (0,0)--(0,1) node[left] {$\vec{j}=\vec{e}_2$};
      \draw[->, main!80!black, very thick] (0,0)--(3,2) node[above right] {$\vec{a}=(3,2)$};
      \draw[->, second, thick] (0,0)--(3,0) node[midway,below] {$3\vec{i}$};
      \draw[->, second, thick] (3,0)--(3,2) node[midway,right] {$2\vec{j}$};
      \draw[dashed, second!70!black] (0,2)--(3,2);
      \node[main!70!black] at (2,-0.85) {正交单位基底：读数即坐标};
    \end{scope}
  \end{tikzpicture}
  \caption{平面向量的基底分解与坐标：左图为任意不共线基底 $\{\vec e_1,\vec e_2\}$ 下 $\vec a=x\vec e_1+y\vec e_2$ 的唯一分解；右图取正交单位基底 $\{\vec i,\vec j\}$ 时，读数 $(x,y)$ 就是平面直角坐标}
  \label{fig:basis-coord}
\end{figure}

\key{（4）「基底选得巧，计算就少」。}
分解的目的不是形式，而是让未知关系落到已知方向上。看一个反复出现的模型：

\begin{example}
在 $\triangle ABC$ 中，$D$ 在 $BC$ 上且 $BD:DC=1:2$，用基底 $\{\overrightarrow{AB},\overrightarrow{AC}\}$ 表示 $\overrightarrow{AD}$。
\end{example}
\begin{solution}
取 $\overrightarrow{AB},\overrightarrow{AC}$ 为基底。由 $BD:DC=1:2$ 知 $\overrightarrow{BD}=\dfrac13\overrightarrow{BC}$，于是
\[
  \overrightarrow{AD}=\overrightarrow{AB}+\overrightarrow{BD}
  =\overrightarrow{AB}+\frac13\overrightarrow{BC}
  =\overrightarrow{AB}+\frac13(\overrightarrow{AC}-\overrightarrow{AB})
  =\frac23\overrightarrow{AB}+\frac13\overrightarrow{AC}.
\]
取 $D$ 为中点时即得常用结论 $\overrightarrow{AD}=\dfrac12(\overrightarrow{AB}+\overrightarrow{AC})$。
\key{一旦基底选定，所有向量都变成有序数对，几何题就变成解方程组。}
\end{solution}

\begin{conccard}[一级]{平面向量基本定理与坐标表示}{把任意向量用两个不共线向量线性表示；判断三点共线；坐标运算}{认为「基底必须互相垂直」；分解不唯一时怀疑定理；写出坐标却忘记它依赖所选基底}
$\{\vec e_1,\vec e_2\}$ 不共线 $\implies$ 平面内任一向量 $\vec a$ 都可唯一写成 $\lambda_1\vec e_1+\lambda_2\vec e_2$。
特别地，$\overrightarrow{OP}=x\overrightarrow{OA}+y\overrightarrow{OB}$ 且 $x+y=1\iff P,A,B$ 三点共线（等和线）。
\end{conccard}

\subsection{数量积与投影算子}

\key{探究二：数量积是「投影乘模长」，更是余弦定理的代数化}

\key{（1）定义与几何意义。}
\[
  \vec a\cdot\vec b=|\vec a|\,|\vec b|\cos\theta,\qquad \theta=\langle\vec a,\vec b\rangle\in[0,\pi].
\]
几何读法：先取 $\vec a$ 在 $\vec b$ 方向上的投影 $|\vec a|\cos\theta$（可正可负可为零），再乘以 $|\vec b|$。
特例：同向 $\vec a\cdot\vec b=|\vec a||\vec b|$；反向 $=-|\vec a||\vec b|$；垂直 $=0$。
显然 $\vec a\cdot\vec a=|\vec a|^2$——\key{数量积自带求模公式}。

\key{（2）坐标运算：把定义交给分配律。}
\begin{derivation}
\begin{align*}
  \vec a\cdot\vec b
  &=(x_1\vec i+y_1\vec j)\cdot(x_2\vec i+y_2\vec j)\\
  &=x_1x_2\,\underbrace{\vec i\cdot\vec i}_{=1}
    +(x_1y_2+x_2y_1)\,\underbrace{\vec i\cdot\vec j}_{=0}
    +y_1y_2\,\underbrace{\vec j\cdot\vec j}_{=1}
  =x_1x_2+y_1y_2.
\end{align*}
\end{derivation}
即「对应坐标相乘再相加」。由此还可直接验证交换律 $\vec a\cdot\vec b=\vec b\cdot\vec a$ 与分配律。
夹角公式：$\cos\theta=\dfrac{\vec a\cdot\vec b}{|\vec a||\vec b|}=\dfrac{x_1x_2+y_1y_2}{\sqrt{x_1^2+y_1^2}\sqrt{x_2^2+y_2^2}}$。

\key{（3）投影：数量投影与投影向量。}
\[
  \text{数量投影}=\frac{\vec a\cdot\vec b}{|\vec b|}=|\vec a|\cos\theta,\qquad
  \mathrm{proj}_{\vec b}\vec a=\frac{\vec a\cdot\vec b}{|\vec b|^2}\,\vec b.
\]
记忆钩子：\key{投影向量 $=$ 数量投影 $\times$ 单位向量}，即 $\mathrm{proj}_{\vec b}\vec a=\left(\vec a\cdot\dfrac{\vec b}{|\vec b|}\right)\dfrac{\vec b}{|\vec b|}$。
注意分母是 $|\vec b|^2$（保证量纲是长度），且投影方向可以逆着 $\vec b$（钝角时为负）。

\begin{figure}[htbp]
  \centering
  \begin{tikzpicture}[>=Stealth, semithick, font=\footnotesize]
    \coordinate (O) at (0,0);
    \coordinate (B) at (4,0);
    \coordinate (A) at (2.6,2.2);
    \coordinate (H) at (2.6,0);
    \draw[->, gray] (-0.4,0)--(4.6,0);
    \draw[->, gray] (0,-0.4)--(0,2.6);
    \fill (O) circle (1pt) node[below left] {$O$};
    \draw[->, main!80!black, very thick] (O)--(B) node[below right] {$\vec{b}$};
    \draw[->, main!80!black, very thick] (O)--(A) node[above] {$\vec{a}$};
    \draw[dashed, second, thick] (A)--(H);
    \draw[->, second, thick] (O)--(H);
    \draw[gray, thin] (2.35,0)--(2.35,0.25)--(2.6,0.25);
    \draw[gray] (0.9,0) arc (0:40:0.9);
    \node[gray!80!black] at (1.2,0.26) {$\theta$};
    % 两条公式统一放到底部居中说明，避免压到 x 轴与图形外撑宽图幅
    \node[second, align=center] at (2.1,-0.95)
      {向量投影 $\mathrm{proj}_{\vec b}\vec a=\dfrac{\vec a\cdot\vec b}{|\vec b|^2}\vec b$\\[3pt]
       数量投影 $=\dfrac{\vec a\cdot\vec b}{|\vec b|}=|\vec a|\cos\theta$};
  \end{tikzpicture}
  \caption{数量积的投影几何：$\vec a$ 在 $\vec b$ 方向上的垂足为 $H$，$OH$ 即投影向量；数量积 $=|\vec b|\times|\vec a|\cos\theta$，即「投影乘模长」}
  \label{fig:dot-proj}
\end{figure}

\key{（4）与余弦定理互译：数量积就是余弦定理的代数化。}
\begin{derivation}
在 $\triangle ABC$ 中记 $\overrightarrow{CA}=\vec u$，$\overrightarrow{CB}=\vec v$，则 $|\vec u|=b$，$|\vec v|=a$，$|\vec u-\vec v|=|\overrightarrow{BA}|=c$。于是
\[
  c^2=|\vec u-\vec v|^2=|\vec u|^2+|\vec v|^2-2\,\vec u\cdot\vec v
  =a^2+b^2-2ab\cos C,
\]
反解即 $\cos C=\dfrac{a^2+b^2-c^2}{2ab}$。
\end{derivation}
同一个事实的两种语言：\key{余弦定理是数量积的模长读法，数量积是余弦定理的代数化}。
第 2 章用「复平面投影」理解 $\cos$，本章用「向量投影」理解 $\cos$——两种语言指向同一个对象。

\begin{example}
已知 $|\vec a|=4$，$|\vec b|=5$，夹角 $\theta=120^\circ$，求 $\vec a\cdot\vec b$、数量投影与投影向量。
\end{example}
\begin{solution}
\[
  \vec a\cdot\vec b=4\times5\times\cos120^\circ=20\times\left(-\frac12\right)=-10.
\]
数量投影 $=\dfrac{\vec a\cdot\vec b}{|\vec b|}=\dfrac{-10}{5}=-2$（负号说明投影逆着 $\vec b$）；
投影向量 $=\dfrac{\vec a\cdot\vec b}{|\vec b|^2}\vec b=\dfrac{-10}{25}\vec b=-\dfrac25\vec b$。
\end{solution}

\begin{conccard}[一级]{数量积与投影}{求模长、判垂直、求夹角、做投影；坐标运算；与余弦定理互译}{把投影向量分母写成 $|\vec b|$；认为 $\vec a\cdot\vec b$ 是向量；默认投影非负}
$\vec a\cdot\vec b=|\vec a||\vec b|\cos\theta=x_1x_2+y_1y_2$；
$\mathrm{proj}_{\vec b}\vec a=\dfrac{\vec a\cdot\vec b}{|\vec b|^2}\vec b$，数量投影 $=\dfrac{\vec a\cdot\vec b}{|\vec b|}$。
$\vec a\perp\vec b\iff\vec a\cdot\vec b=0$（非零向量）。
\end{conccard}

\begin{commonerror}{投影运算的三个高频失误}
\paragraph{失误一：分母多除一次模长。}
投影向量是 $\dfrac{\vec a\cdot\vec b}{|\vec b|^2}\vec b$。若误写 $\dfrac{\vec a\cdot\vec b}{|\vec b|}\vec b$，
左边是「数量投影（长度）再乘向量」，量纲变成长度的平方，必然出错。
正确链条：数量投影 $\dfrac{\vec a\cdot\vec b}{|\vec b|}$（一个数）$\to$ 乘单位向量 $\dfrac{\vec b}{|\vec b|}$ $\to$ 投影向量。
\paragraph{失误二：把数量积当成向量。}
$\vec a\cdot\vec b$ 是一个\key{实数}（可正可负可为零），不是向量；
更不能写成 $\vec a\cdot\vec b\,\vec b$ 这种不伦不类的式子。
\paragraph{失误三：默认投影非负、夹角非钝。}
$\theta\in[0,\pi]$，钝角时 $\cos\theta<0$，数量投影与数量积都是负数；
例如上例 $\theta=120^\circ$ 时数量投影为 $-2$。符号携带方向信息，不能丢。
\end{commonerror}

\subsection{极化恒等式与面积}

\key{探究三：极化恒等式——共起点三向量乘积的秒杀通道}

（恒等式的核心是两向量之积 $\overrightarrow{AB}\cdot\overrightarrow{AC}$；
下面的中点形式由它配方推出。）

\key{（1）一般的极化恒等式。}
对任意向量 $\vec u,\vec v$：
\[
  \vec u\cdot\vec v=\frac14\left(|\vec u+\vec v|^2-|\vec u-\vec v|^2\right).
\]
推导只需展开右边：$|\vec u+\vec v|^2-|\vec u-\vec v|^2=4\,\vec u\cdot\vec v$。

\key{（2）三角形中的中点版本。}
设 $M$ 是 $BC$ 的中点，则 $\overrightarrow{AB}+\overrightarrow{AC}=2\overrightarrow{AM}$，$\overrightarrow{AB}-\overrightarrow{AC}=\overrightarrow{CB}$，代入得
\begin{derivation}
\[
  \overrightarrow{AB}\cdot\overrightarrow{AC}
  =\frac14\left(|2\overrightarrow{AM}|^2-|\overrightarrow{CB}|^2\right)
  =|\overrightarrow{AM}|^2-\frac14|\overrightarrow{CB}|^2
  =|\overrightarrow{AM}|^2-|\overrightarrow{MB}|^2.
\]
\end{derivation}
这个恒等式把「\key{共起点的两个向量之积}」改写成「\key{动点到中点的距离平方} $-$ 定值」，
是处理动态最值问题的标准武器。

\begin{example}
已知圆 $O$ 的半径为 $2$，弦 $BC=2$，点 $A$ 在圆上运动，求 $\overrightarrow{AB}\cdot\overrightarrow{AC}$ 的最大值。
\end{example}
\begin{solution}
取 $M$ 为 $BC$ 中点，则 $|\overrightarrow{MB}|=1$ 为定值，由极化恒等式
\[
  \overrightarrow{AB}\cdot\overrightarrow{AC}=|\overrightarrow{AM}|^2-1.
\]
问题转化为求 $|\overrightarrow{AM}|$ 的最大值。由垂径定理 $|OM|=\sqrt{2^2-1^2}=\sqrt3$，
$A$ 在圆上运动时 $|AM|$ 的最大值为 $R+|OM|=2+\sqrt3$（$A$ 位于射线 $MO$ 的延长线与圆的交点）。
因此
\[
  \left(\overrightarrow{AB}\cdot\overrightarrow{AC}\right)_{\max}=(2+\sqrt3)^2-1=6+4\sqrt3.
\]
若直接设坐标硬算，将陷入三角换元；\key{极化恒等式一步就把「双动向量乘积」化为「单动点距离」。}
\end{solution}

\key{（3）面积形式：叉积的模。}
与数量积配对的是\key{叉积（外积）}：$\vec a\times\vec b$ 是垂直于 $\vec a,\vec b$ 所在平面的向量，
其模
\[
  |\vec a\times\vec b|=|\vec a|\,|\vec b|\,|\sin\theta|
\]
恰等于以 $\vec a,\vec b$ 为邻边的平行四边形面积；坐标计算公式为
\[
  |\vec a\times\vec b|=|x_1y_2-x_2y_1|,\qquad
  S_{\triangle ABC}=\frac12|\overrightarrow{AB}\times\overrightarrow{AC}|
  =\frac12\,|x_1y_2-x_2y_1|.
\]
对照记忆：\key{点积管「角与投影」（余弦），叉积管「面积与正弦」}，二者互补。

\begin{example}
已知 $\vec a=(3,1)$，$\vec b=(2,4)$，求以 $\vec a,\vec b$ 为邻边的三角形面积。
\end{example}
\begin{solution}
\[
  S=\frac12|x_1y_2-x_2y_1|=\frac12|3\times4-1\times2|=\frac12\times10=5.
\]
\end{solution}

注记（高观点）：把三个极化恒等式按三角形面积加权求和，就得到著名的奔驰定理
$S_{\triangle BOC}\cdot\overrightarrow{OA}+S_{\triangle COA}\cdot\overrightarrow{OB}+S_{\triangle AOB}\cdot\overrightarrow{OC}=\vec 0$
（$O$ 为 $\triangle ABC$ 内一点）——它是「重心坐标」的面积版；第 7 章统计回归的样本中心点与之同源。

\begin{conccard}[二级]{极化恒等式：共起点三向量乘积}{共起点两向量数量积的最值／范围问题，尤其是「一定弦、一动顶点」模型}{认不出中点结构，仍去建系硬算；把 $|\overrightarrow{MB}|$ 误当 $|BC|$；忽略动点到中点距离的取值范围需按圆／线段讨论}
$\overrightarrow{AB}\cdot\overrightarrow{AC}=|\overrightarrow{AM}|^2-|\overrightarrow{MB}|^2$（$M$ 为 $BC$ 中点）。
更一般地 $\vec u\cdot\vec v=\frac14(|\vec u+\vec v|^2-|\vec u-\vec v|^2)$。
配套面积公式：$S_{\triangle ABC}=\frac12|\overrightarrow{AB}\times\overrightarrow{AC}|=\frac12|x_1y_2-x_2y_1|$。
\end{conccard}

\subsection{空间向量：方向向量与法向量}

\key{探究四：用法向量统一处理空间中的平行、垂直与夹角}

\key{（1）空间向量坐标化。}
以 $O$ 为原点、$\vec i,\vec j,\vec k$ 为三条两两垂直的单位向量建立空间直角坐标系，则
\[
  \vec a=x\vec i+y\vec j+z\vec k\iff \vec a=(x,y,z),\qquad
  |\vec a|=\sqrt{x^2+y^2+z^2},
\]
\[
  \vec a\cdot\vec b=x_1x_2+y_1y_2+z_1z_2,\qquad
  \cos\langle\vec a,\vec b\rangle=\frac{\vec a\cdot\vec b}{|\vec a||\vec b|}.
\]
几何空间代数化的关键一步：\key{建系 $=$ 给每个点分配坐标，之后所有几何量都是坐标运算}。

\key{（2）平面的法向量。}
与平面 $\alpha$ 垂直的向量 $\vec n$ 叫做 $\alpha$ 的\key{法向量}。
一个平面有无穷多个法向量，但它们彼此平行（共线），所以\key{法向量的方向就是平面的「朝向」}。
若 $\vec n=(A,B,C)$，$P_0(x_0,y_0,z_0)\in\alpha$，则平面方程为
\[
  A(x-x_0)+B(y-y_0)+C(z-z_0)=0\quad\text{即}\quad Ax+By+Cz+D=0.
\]
求法向量的通用方法：找平面内两个不共线向量 $\vec a,\vec b$，令 $\vec n\cdot\vec a=0$、$\vec n\cdot\vec b=0$ 解方程组。

\key{（3）线面角与二面角。}
设直线 $l$ 的方向向量为 $\vec s$，平面 $\alpha$ 的法向量为 $\vec n$：
\[
  \textbf{线面角 }\varphi:\quad \sin\varphi=\frac{|\vec s\cdot\vec n|}{|\vec s|\,|\vec n|}
  \quad(\varphi\in[0,\tfrac{\pi}{2}]).
\]
注意是\key{正弦}：直线与法向量的夹角与线面角互余。
设两平面法向量为 $\vec n_1,\vec n_2$：
\[
  \textbf{二面角 }\theta:\quad |\cos\theta|=\frac{|\vec n_1\cdot\vec n_2|}{|\vec n_1|\,|\vec n_2|},
\]
锐钝要结合图形判断（或用法向量方向调整符号）。

\begin{figure}[htbp]
  \centering
  \begin{tikzpicture}[>=Stealth, semithick, font=\small]
    % 两个半平面（平行四边形示意），公共棱为中间那条水平线
    \fill[main!8] (-2.6,0)--(2.6,0)--(1.9,2.4)--(-3.3,2.4)--cycle;
    \fill[second!10] (-2.6,0)--(2.6,0)--(1.9,-2.0)--(-3.3,-2.0)--cycle;
    \draw[thick, main!60!black] (-2.6,0)--(2.6,0)--(1.9,2.4)--(-3.3,2.4)--cycle;
    \draw[thick, second!50!black] (-2.6,0)--(2.6,0)--(1.9,-2.0)--(-3.3,-2.0)--cycle;
    % 棱与棱上一点
    \draw[very thick, gray] (-2.6,0)--(2.6,0);
    \fill (0,0) circle (1.5pt) node[below left] {$P$};
    % 二面角的平面角：在两个半平面内分别作垂直于棱的射线
    \draw[->, third, very thick] (0,0)--(-1.1,1.9) node[above left] {$\vec u\subset\alpha$};
    \draw[->, third, very thick] (0,0)--(-1.1,-1.6) node[below left] {$\vec v\subset\beta$};
    % 直角符号：u\perp l，v\perp l（放大并贴近射线）
    \draw[third!70!black, thin] (0.4,0)--(0.2,0.3464)--(-0.2,0.3464);
    \draw[third!70!black, thin] (0.4,0)--(0.1736,-0.3296)--(-0.2264,-0.3296);
    % 平面角 \theta：弧线从 \vec u 跨过 \vec v
    \draw[main, thick] (0,0) ++(120:0.95) arc (120:235:0.95);
    \node[main] at (-1.15,0) {$\theta$};
    % 法向量
    \draw[->, second, thick] (1.2,1.2)--(2.2,2.6) node[above right] {$\vec n_1$};
    \draw[->, main, thick] (1.2,-1.0)--(2.2,-2.4) node[below right] {$\vec n_2$};
    \node[main!70!black] at (-2.6,1.7) {$\alpha$};
    \node[second!60!black] at (-2.6,-1.4) {$\beta$};
    \node[gray!70!black] at (2.5,0.28) {$l$};
    \node[gray!70!black, align=center] at (0,-3.05) {沿棱 $l$ 取点 $P$，在 $\alpha,\beta$ 内分别作垂直于 $l$ 的射线，\\两条射线的夹角即二面角的平面角};
  \end{tikzpicture}
  \caption{二面角的几何图：两平面 $\alpha,\beta$ 交于棱 $l$；$\vec u,\vec v$ 分别躺在两个半平面内且都垂直于 $l$（直角符号），夹角就是二面角的平面角 $\theta$。法向量 $\vec n_1,\vec n_2$ 分别垂直于两平面，其夹角与 $\theta$ 相等或互补，故 $|\cos\theta|=\dfrac{|\vec n_1\cdot\vec n_2|}{|\vec n_1||\vec n_2|}$，锐钝需结合图形判断}
  \label{fig:dihedral-angle}
\end{figure}

\begin{center}
\footnotesize
\setlength{\tabcolsep}{3pt}
\renewcommand{\arraystretch}{1.4}
\begin{tabular}{p{2.4cm}p{6.8cm}p{3.6cm}}
\toprule
位置关系 & 向量判据 & 说明 \\
\midrule
线线平行 & $\vec s_1\parallel\vec s_2$（即 $\vec s_1\times\vec s_2=\vec 0$） & 方向向量共线 \\
线线垂直 & $\vec s_1\cdot\vec s_2=0$ & 方向向量垂直 \\
线面平行 & $\vec s\cdot\vec n=0$ 且 $l\not\subset\alpha$ & 方向向量与法向量垂直 \\
线面垂直 & $\vec s\parallel\vec n$ & 方向向量与法向量共线 \\
面面平行 & $\vec n_1\parallel\vec n_2$ & 法向量共线 \\
面面垂直 & $\vec n_1\cdot\vec n_2=0$ & 法向量垂直 \\
\bottomrule
\end{tabular}
\end{center}

\begin{figure}[htbp]
  \centering
  \begin{tikzpicture}[>=Stealth, semithick, font=\footnotesize]
    \coordinate (O) at (0,0);
    \coordinate (X) at (-2.3,-1.05);
    \coordinate (Y) at (2.9,-0.75);
    \coordinate (XY) at (0.6,-1.8);
    % 平面 α（xOy 平面，z=0）
    \fill[third!8] (O)--(X)--(XY)--(Y)--cycle;
    \draw[third!50!black, thin] (O)--(X)--(XY)--(Y)--cycle;
    \node[third!70!black] at (2.9,-1.05) {$\alpha$};
    % 坐标轴
    \draw[->, gray] (O)--(X) node[below left] {$x$};
    \draw[->, gray] (O)--(Y) node[below right] {$y$};
    \draw[->, gray] (O)--(0,3.1) node[above] {$z$};
    % 法向量 n（竖直向上，垂直于平面）
    \draw[->, third, very thick] (-0.35,-0.5)--(-0.35,1.3) node[pos=0.85,right] {$\vec{n}$};
    % 直线方向向量 s 穿过平面
    \coordinate (R) at (1.15,-0.85);
    \coordinate (S) at (1.7,1.5);
    \draw[->, main!80!black, very thick] (R)--(S) node[midway,right] {$\vec{s}$};
    \draw[dashed, gray] (S)--(1.7,0);
    % 线面角 φ
    \coordinate (I) at (1.349,0);
    \draw (1.799,0) arc (0:77:0.45);
    \node at (1.98,0.4) {$\varphi$};
    \node[main!70!black] at (0.6,-2.35) {方向向量 $\vec s$、法向量 $\vec n$ 与线面角 $\varphi$};
  \end{tikzpicture}
  \caption{空间向量与法向量：平面 $\alpha$ 的法向量 $\vec n$ 指示平面朝向；直线方向向量 $\vec s$ 穿过平面，$\varphi$ 为 $\vec s$ 与其在平面内投影的夹角（线面角）}
  \label{fig:space-normal}
\end{figure}

\begin{example}
正方体 $ABCD-A_1B_1C_1D_1$ 的棱长为 $1$，求二面角 $A_1-BD-C$ 的余弦值。
\end{example}
\begin{solution}
以 $D$ 为原点建系：$D(0,0,0)$，$A(1,0,0)$，$C(0,1,0)$，$B(1,1,0)$，$A_1(1,0,1)$。
平面 $BDC$ 即底面 $z=0$，取其法向量 $\vec n_1=(0,0,1)$。
平面 $A_1BD$ 内有 $\overrightarrow{DB}=(1,1,0)$，$\overrightarrow{DA_1}=(1,0,1)$，
\[
  \vec n_2=\overrightarrow{DB}\times\overrightarrow{DA_1}
  =(1\times1-0\times0,\;0\times1-1\times1,\;1\times0-1\times1)=(1,-1,-1).
\]
于是
\[
  |\cos\theta|=\frac{|\vec n_1\cdot\vec n_2|}{|\vec n_1|\,|\vec n_2|}
  =\frac{|-1|}{1\times\sqrt3}=\frac{\sqrt3}{3}.
\]
取 $BD$ 中点 $M$，$\overrightarrow{MA_1}=(-0.5,-0.5,1)$、$\overrightarrow{MC}=(-0.5,-0.5,0)$ 均与 $BD$ 垂直，
两者点积为正，故二面角 $A_1-BD-C$ 为锐角，其余弦值为 $\dfrac{\sqrt3}{3}$。
传统几何法需要作平面角、证垂直、再解三角形；向量法只需建系、算两个法向量、代公式。
\end{solution}

\begin{example}
平面 $\alpha,\beta$ 的法向量分别为 $\vec n_1=(2,-1,1)$，$\vec n_2=(1,1,-1)$，直线 $l$ 的方向向量为 $\vec s=(1,2,2)$。
（1）判断 $l$ 与 $\alpha$ 的位置关系并求线面角；（2）判断 $\alpha$ 与 $\beta$ 的位置关系。
\end{example}
\begin{solution}
（1）$\vec s\cdot\vec n_1=1\times2+2\times(-1)+2\times1=2\neq0$，
所以 $l$ 与 $\alpha$ \key{相交}（不平行也不在面内）。线面角 $\varphi$ 满足
\[
  \sin\varphi=\frac{|\vec s\cdot\vec n_1|}{|\vec s|\,|\vec n_1|}
  =\frac{2}{\sqrt{1+4+4}\cdot\sqrt{4+1+1}}=\frac{2}{3\sqrt6}=\frac{\sqrt6}{9}.
\]
（2）$\vec n_1\cdot\vec n_2=2-1-1=0$，所以两平面的法向量垂直，即 \key{$\alpha\perp\beta$}。
\end{solution}

\begin{conccard}[二级]{法向量统一处理空间位置关系}{线面／面面的平行、垂直、夹角问题；点到平面距离；存在性问题}{把线面角公式错写成余弦；二面角不判断锐钝；建系时没说明三条轴两两垂直或坐标写错}
直线方向向量 $\vec s$、法向量 $\vec n$：
线面平行 $\iff\vec s\cdot\vec n=0$；线面垂直 $\iff\vec s\parallel\vec n$；
线面角 $\sin\varphi=\dfrac{|\vec s\cdot\vec n|}{|\vec s||\vec n|}$；
二面角 $|\cos\theta|=\dfrac{|\vec n_1\cdot\vec n_2|}{|\vec n_1||\vec n_2|}$。
（第 4 章立体几何主干直接调用这套判据与下面的易错清单，不再重复。）
\end{conccard}

\begin{commonerror}{空间向量建系的三个高频失误}
\paragraph{失误一：建系不交代垂直性。}
使用空间直角坐标系前，必须说明所选三条轴（或从同一点出发的三条棱）两两垂直；
若题目只给出斜角或没有垂直条件，应先证明垂直关系再建系。
\paragraph{失误二：线面角套错公式。}
线面角 $\varphi$ 与「方向向量和法向量的夹角」互余，所以
$\sin\varphi=\dfrac{|\vec s\cdot\vec n|}{|\vec s||\vec n|}$，而不是余弦。
把 $\vec s$ 换成它在面内的投影，才是余弦关系，切勿混用。
\paragraph{失误三：二面角不判断锐钝。}
$|\cos\theta|$ 只给大小；必须结合图形（或取棱上一点，用两个半平面内的向量验号）确定二面角是锐角还是钝角，否则一半情形会丢。
\end{commonerror}

\begin{advancedonly}
\key{超纲延伸一：换基底与相似律——同一线性变换，在不同基下「读数不同，变换不动」。}

设线性变换 $T$ 在标准基 $\{\vec e_1,\vec e_2\}$ 下的矩阵为
$A=\begin{pmatrix}2&1\\0&1\end{pmatrix}$（即 $T(\vec e_1)=(2,0)$，$T(\vec e_2)=(1,1)$）。
取 $\vec v=(2,1)$，则
\[
  T(\vec v)=A\vec v=\begin{pmatrix}2&1\\0&1\end{pmatrix}\begin{pmatrix}2\\1\end{pmatrix}
  =\tk{\begin{pmatrix}5\\1\end{pmatrix}}.
\]
换一组新基 $\mathcal B'=\{\vec u_1,\vec u_2\}$：$\vec u_1=(1,1)$，$\vec u_2=(-1,1)$，
过渡矩阵 $P=(\vec u_1\ \vec u_2)=\begin{pmatrix}1&-1\\1&1\end{pmatrix}$。
同一个向量换基后读数要左乘 $P^{-1}$；由 $\det P=2$ 得
$P^{-1}=\dfrac12\begin{pmatrix}1&1\\-1&1\end{pmatrix}$，于是
\[
  [\vec v]_{\mathcal B'}=P^{-1}\begin{pmatrix}2\\1\end{pmatrix}
  =\tk{\begin{pmatrix}\tfrac32\\[2pt]-\tfrac12\end{pmatrix}},\qquad
  [T(\vec v)]_{\mathcal B'}=P^{-1}\begin{pmatrix}5\\1\end{pmatrix}
  =\tk{\begin{pmatrix}3\\-2\end{pmatrix}}.
\]
而 $T$ 在新基下的矩阵不再是 $A$，它由相似律给出：
\[
  A'=P^{-1}AP
  =\frac12\begin{pmatrix}1&1\\-1&1\end{pmatrix}
   \begin{pmatrix}2&1\\0&1\end{pmatrix}
   \begin{pmatrix}1&-1\\1&1\end{pmatrix}
  =\tk{\begin{pmatrix}2&0\\-1&1\end{pmatrix}}.
\]
验证（一步）：用新矩阵乘新坐标，必须还原像向量的读数，
\[
  A'\,[\vec v]_{\mathcal B'}
  =\begin{pmatrix}2&0\\-1&1\end{pmatrix}
   \begin{pmatrix}\tfrac32\\[2pt]-\tfrac12\end{pmatrix}
  =\tk{\begin{pmatrix}3\\-2\end{pmatrix}}=[T(\vec v)]_{\mathcal B'}.
\]
\key{$A'=P^{-1}AP$：换基底只是换一把尺子，坐标读数变了，变换本身纹丝不动。}
顺带一提：$\det A=\det A'=2$——相似矩阵的不变量，正是「换个角度，本质仍在」。
\end{advancedonly}

\begin{advancedonly}
\key{超纲延伸二：镜面反射公式——投影走多远，就回头走多远。}

以 $\vec b$ 所在的整条直线为镜面（不只是向量 $\vec b$ 那一截），
$\vec a$ 的反射像由投影直接写出：
\[
  \vec a'=2\,\mathrm{proj}_{\vec b}\vec a-\vec a.
\]
取 $\vec a=(2.5,1.5)$，$\vec b=(3,1)$。先投影：
\[
  \vec a\cdot\vec b=2.5\times3+1.5\times1=\tk{9},\qquad
  |\vec b|^2=3^2+1^2=\tk{10},
\]
\[
  \mathrm{proj}_{\vec b}\vec a=\frac{\vec a\cdot\vec b}{|\vec b|^2}\,\vec b
  =\frac{9}{10}(3,1)=\tk{(2.7,0.9)}.
\]
再越过镜面走同样远：
\[
  \vec a'=2(2.7,0.9)-(2.5,1.5)=\tk{(2.9,0.3)}.
\]
几何解释：垂足 $\mathrm{proj}_{\vec b}\vec a$ 恰好是 $\vec a$ 与 $\vec a'$ 连线的中点，
也就是说镜面是这条连线的垂直平分线——从 $\vec a$ 到镜面有多远，像就退到另一侧多远。

\begin{figure}[htbp]
  \centering
  \begin{tikzpicture}[>=Stealth, semithick, font=\footnotesize]
    \coordinate (O) at (0,0);
    \coordinate (A) at (2.5,1.5);
    \coordinate (B) at (3,1);
    \coordinate (F) at (2.7,0.9);
    \coordinate (Ap) at (2.9,0.3);
    % 镜面：\vec b 所在的整条直线（双向延长）
    \draw[dashed, second, thick] (-1.5,-0.5)--(4.2,1.4);
    \node[second, below left] at (-1.5,-0.5) {镜面};
    % 坐标轴
    \draw[->, gray] (-1.0,0)--(4.6,0);
    \draw[->, gray] (0,-0.9)--(0,2.6);
    \fill (O) circle (1pt);
    % 投影向量（恰好躺在镜面上）与垂足
    \draw[->, third, very thick] (O)--(F);
    \fill[third] (F) circle (1.5pt);
    \node[third] at (1.55,0.30) {$\mathrm{proj}_{\vec b}\vec a$};
    % 原向量、像向量与镜面方向向量
    \draw[->, main!80!black, very thick] (O)--(A) node[above left] {$\vec a=(2.5,1.5)$};
    \draw[->, second, very thick] (O)--(B) node[above right] {$\vec b=(3,1)$};
    \draw[->, main, very thick] (O)--(Ap) node[right] {$\vec a'=(2.9,0.3)$};
    % 垂线段与垂足处的直角符号
    \draw[dashed, gray] (A)--(F)--(Ap);
    \draw[third!70!black, thin] (2.85,0.95)--(2.80,1.10)--(2.65,1.05);
  \end{tikzpicture}
  \caption{镜面反射：以 $\vec b$ 所在直线（虚线）为镜面，$\mathrm{proj}_{\vec b}\vec a$ 是 $\vec a$ 的垂足，$\vec a'=2\,\mathrm{proj}_{\vec b}\vec a-\vec a$；垂足同时也是 $\vec a$ 与 $\vec a'$ 连线的中点}
  \label{fig:mirror-reflection}
\end{figure}
\end{advancedonly}

\begin{method}
向量解题三步法：
\begin{enumerate}
  \item \key{翻译}：把几何条件（中点、平行、垂直、夹角、面积）译成向量语言；
        不确定从哪下手时，先选一组基底或建立坐标系。有垂直关系优先建系。
  \item \key{运算}：坐标化或基底化之后，用数量积处理长度与夹角，
        用「中点公式 $+$ 极化恒等式」处理共起点向量乘积的最值，
        用方向向量与法向量处理空间位置关系；追求「选得巧，计算就少」。
  \item \key{回译}：把代数结果翻回几何结论，并检查退化情形
        （零向量、两向量共线、点在平面内、分母为零）。
\end{enumerate}
\end{method}

\begin{summary}
本章把向量当作「既有大小又有方向的量」$\to$「点集的位移」$\to$「基底下的坐标」$\to$「内积投影算子」$\to$「空间中的方向与法向」层层展开，凸显向量体系相对三角与函数的独立地位：它是几何空间代数化的通用线性语言。
平面向量基本定理保证任一向量被不共线基底唯一分解，坐标就是有序实数对；数量积 $\vec a\cdot\vec b=|\vec a||\vec b|\cos\theta$ 的几何意义是「投影乘模长」，投影向量 $\mathrm{proj}_{\vec b}\vec a=\frac{\vec a\cdot\vec b}{|\vec b|^2}\vec b$，它与余弦定理互为代数化与几何读法；
极化恒等式 $\overrightarrow{AB}\cdot\overrightarrow{AC}=|\overrightarrow{AM}|^2-|\overrightarrow{MB}|^2$ 把共起点三向量乘积化为动点到中点的距离平方差，是动态最值的秒杀通道，面积则由叉积的模 $S=\frac12|\vec a\times\vec b|$ 给出；
空间向量坐标化之后，平面的法向量成为统一语言：线面平行／垂直、线面角、二面角全部化为方向向量与法向量的数量积运算。
\end{summary}

\end{explore}

\begin{inclasstest}
\begin{enumerate}
  \item 已知 $\vec a=(1,2)$，$\vec b=(3,-1)$，求 $\vec a\cdot\vec b$、$|\vec a|$ 以及 $\vec a$ 在 $\vec b$ 上的投影向量；\level{易}
  \item 已知 $\vec e_1,\vec e_2$ 不共线，$\vec m=\vec e_1+2\vec e_2$，$\vec n=3\vec e_1+k\vec e_2$，若 $\vec m\parallel\vec n$，求实数 $k$；\level{易}
  \item 半径为 $2$ 的圆中，弦 $BC=2$，点 $A$ 在圆上运动，求 $\overrightarrow{AB}\cdot\overrightarrow{AC}$ 的取值范围；\level{中}
  \item 平面 $\alpha$ 的法向量 $\vec n_1=(1,2,-1)$，平面 $\beta$ 的法向量 $\vec n_2=(2,1,1)$，判断 $\alpha$ 与 $\beta$ 的位置关系，并求所成锐二面角的余弦值。\level{中}
\end{enumerate}
\end{inclasstest}

\begin{afterclass}
\subsection*{基础巩固}
\begin{enumerate}
  \item 已知 $\vec a=(2,-1)$，$\vec b=(-3,4)$，求 $\vec a+\vec b$、$\vec a\cdot\vec b$ 与两向量夹角的余弦值；
  \item 设 $\vec e_1,\vec e_2$ 不共线，$\vec a=\vec e_1+2\vec e_2$。若换用新基底 $\vec e_1'=\vec e_1+\vec e_2$，$\vec e_2'=\vec e_1-\vec e_2$，求 $\vec a$ 在新基底下的坐标，并说明坐标为什么随基底改变而改变。
\end{enumerate}

\subsection*{进阶提升}
\begin{enumerate}
  \item 在 $\triangle ABC$ 中，$AB=3$，$AC=5$，$BC=4$，用极化恒等式求 $\overrightarrow{AB}\cdot\overrightarrow{AC}$；
  \item 三棱锥 $D-ABC$ 中，$A(1,0,0)$，$B(0,1,0)$，$C(0,0,1)$，$D(0,0,0)$，求二面角 $A-BC-D$ 的余弦值。
\end{enumerate}

\subsection*{拓展探究}
\begin{enumerate}
  \item 用向量方法证明：三角形三条中线交于一点，且该点分每条中线为 $2:1$（提示：先用平面向量基本定理比较系数求交点分比）。
\end{enumerate}

\subsection*{教师版答案}
\begin{solution}
\textbf{基础 1：}$\vec a+\vec b=(2-3,\,-1+4)=(-1,3)$；
$\vec a\cdot\vec b=2\times(-3)+(-1)\times4=-6-4=-10$；
\[
  |\vec a|=\sqrt{2^2+(-1)^2}=\sqrt5,\quad |\vec b|=\sqrt{(-3)^2+4^2}=5,
\]
\[
  \cos\theta=\frac{\vec a\cdot\vec b}{|\vec a||\vec b|}=\frac{-10}{5\sqrt5}=-\frac{2\sqrt5}{5}.
\]
\end{solution}
\begin{solution}
\textbf{基础 2：}设 $\vec a=\lambda\vec e_1'+\mu\vec e_2'$，即
\[
  \vec a=\lambda(\vec e_1+\vec e_2)+\mu(\vec e_1-\vec e_2)
  =(\lambda+\mu)\vec e_1+(\lambda-\mu)\vec e_2.
\]
与 $\vec a=\vec e_1+2\vec e_2$ 比较系数（基本定理保证唯一）：
\[
  \lambda+\mu=1,\quad \lambda-\mu=2\implies\lambda=\frac32,\ \mu=-\frac12.
\]
故新坐标为 $\left(\dfrac32,-\dfrac12\right)$。
坐标是「向量在基底下的读数」，基底变了读数自然改变；变的只是坐标，向量本身不变。
\end{solution}
\begin{solution}
\textbf{进阶 1：}取 $BC$ 中点 $M$，$|MB|=2$，中线长
\[
  |AM|^2=\frac{2AB^2+2AC^2-BC^2}{4}
  =\frac{2\times9+2\times25-16}{4}=\frac{52}{4}=13.
\]
由极化恒等式
\[
  \overrightarrow{AB}\cdot\overrightarrow{AC}=|\overrightarrow{AM}|^2-|\overrightarrow{MB}|^2=13-4=9.
\]
（检验：$\cos A=\frac{9+25-16}{2\times3\times5}=\frac{3}{5}$，故 $\overrightarrow{AB}\cdot\overrightarrow{AC}=3\times5\times\frac35=9$。）
\end{solution}
\begin{solution}
\textbf{进阶 2：}平面 $ABC$：由 $A,B,C$ 的坐标知方程为 $x+y+z=1$，法向量 $\vec n_1=(1,1,1)$；
平面 $DBC$ 过原点且含 $B,C$，方程为 $x=0$，法向量 $\vec n_2=(1,0,0)$。
\[
  |\cos\theta|=\frac{|\vec n_1\cdot\vec n_2|}{|\vec n_1||\vec n_2|}
  =\frac{1}{\sqrt3\times1}=\frac{\sqrt3}{3}.
\]
取 $BC$ 中点 $M(0,\tfrac12,\tfrac12)$，$\overrightarrow{MA}=(1,-\tfrac12,-\tfrac12)$、$\overrightarrow{MD}=(0,-\tfrac12,-\tfrac12)$ 均与 $\overrightarrow{BC}=(0,1,-1)$ 垂直，且点积为正，故二面角 $A-BC-D$ 为锐角，余弦值为 $\dfrac{\sqrt3}{3}$。
\end{solution}
\begin{solution}
\textbf{拓展：}取 $A$ 为原点，记 $\overrightarrow{AB}=\vec b$，$\overrightarrow{AC}=\vec c$。
设 $D,E$ 分别为 $BC,CA$ 的中点，则
\[
  \overrightarrow{AD}=\frac{\vec b+\vec c}{2},\qquad
  \overrightarrow{BE}=\overrightarrow{AE}-\overrightarrow{AB}=\frac{\vec c}{2}-\vec b.
\]
设中线 $AD$ 与 $BE$ 交于点 $P$，令 $\overrightarrow{AP}=t\,\overrightarrow{AD}=\dfrac{t}{2}\vec b+\dfrac{t}{2}\vec c$；
另一方面 $\overrightarrow{AP}=\overrightarrow{AB}+s\,\overrightarrow{BE}=(1-s)\vec b+\dfrac{s}{2}\vec c$。
由平面向量基本定理比较系数：
\[
  \frac t2=1-s,\qquad \frac t2=\frac s2
  \implies s=t,\ \frac t2=1-t\implies t=\frac23.
\]
故 $AP:PD=BP:PE=2:1$。由对称性，$P$ 也在中线 $CF$（$F$ 为 $AB$ 中点）上，即三条中线共点，且该点（重心）分每条中线为 $2:1$。
这正是解析几何与统计中「均值点」的向量面孔（第 7 章样本中心与之同源）。
\end{solution}
\reading{拓展阅读 09}{向量：几何空间的通用线性语言}
\end{afterclass}

% 高观点导学案：立体几何与圆锥曲线大一统（第 4 章·高观点统一版）
% 主线：仿射变换（圆压缩为椭圆）→ 三维圆锥截面（Dandelin）→ 考纲安全落地双层法 → 解析几何主干 → 立体几何主干（向量法统一）
% 核心思想：圆锥曲线不是三类孤立的二次曲线，而是同一个三维对象（圆锥面）与不同倾角平面的相遇方式；
%           高观点负责「预判答案」，考纲通法负责「写出满分」。

\section{立体几何与圆锥曲线大一统（高观点版）}

\begin{methodbox}{多框架对读}
\small
\begin{center}
\begin{tabular}{c|p{2.9cm}|p{5.0cm}|p{3.1cm}}
\toprule
视角 & 本章看点 & 别的视角怎么说 & 哪侧最省力 \\
\midrule
代数 & 标准方程、判别式、弦长 & 坐标满足的条件 & 硬算与判别 \\
几何 & 轨迹形状、仿射变换、二面角 & 图形的位置与度量 & 猜结论、估范围 \\
集合 & 曲线＝点集；平行垂直＝关系 & 条件即约束的集合 & 界存在与唯一 \\
映射 & 仿射把圆映成椭圆；投影降维 & 换表示＝换坐标系 & 整体保持性质 \\
\bottomrule
\end{tabular}
\end{center}
卡的使用：解析几何先问「这是哪个变换下的不变量」，再决定代数硬算还是几何巧解（总图见第 0 章图~\ref{fig:four-view-map}）。
\end{methodbox}

第 3 章把向量确立为「几何空间代数化」的通用语言；本章再向前一步，打通立体几何与解析几何的维度壁垒：
椭圆不过是圆在\key{仿射变换}下的像，圆锥曲线不过是圆锥面在不同倾角平面上的\key{截线}——
维度不是障碍，变换才是本质。五个模块层层展开：
\key{仿射变换：从圆到椭圆的降维打击}（圆沿一个方向均匀压缩即得椭圆，共线、平行、中点、相切、面积比全部保持）$\to$
\key{Dandelin 视角}（椭圆／双曲线／抛物线是同一个圆锥面被不同倾角平面所截的结果）$\to$
\key{考纲安全落地的双层法}（高观点 10 秒预设定点定值，考纲通法写满分步骤）$\to$
\key{解析几何主干}（定义、标准方程、直线与圆锥曲线的位置关系、定点定值通法）$\to$
\key{立体几何主干}（线面关系、空间角与距离，向量法全面落地）。
核心思想：\key{圆锥曲线不是三类孤立的二次曲线，而是同一个三维对象与不同平面的相遇方式}。

\begin{learnobjective}
\begin{itemize}
  \item 理解仿射变换 $T:(x,y)\mapsto\left(x,\frac{b}{a}y\right)$，会把椭圆 $\frac{x^2}{a^2}+\frac{y^2}{b^2}=1$ 看作圆 $x^2+y^2=a^2$ 的压缩像；
  \item 掌握仿射不变量清单（共线、平行、中点、相切、面积比 $S_{\text{椭圆}}=\frac{b}{a}S_{\text{圆}}$），会用它秒算椭圆内接三角形面积最值；
  \item 知道圆锥曲线的 Dandelin 球证法：焦点—准线定义的统一，$e=\frac{\cos\beta}{\cos\alpha}$；
  \item 掌握「高观点预判 $+$ 考纲通法」双层解题策略：极点极线预设定点定值，设线—联立—韦达—判别式写满分步骤；
  \item 掌握椭圆／双曲线／抛物线的定义与标准方程，直线与圆锥曲线位置关系的判别式、弦长与中点弦公式，以及定点定值问题的通法步骤；
  \item 掌握线面平行／垂直的判定与性质、空间角（线面角、二面角）与点到平面距离，会用向量法统一处理（与第 3 章衔接）。
\end{itemize}
\end{learnobjective}

\begin{preview}
\subsection*{知识回顾}
\begin{itemize}
  \item 圆的标准方程是 \blank；把圆沿 $y$ 轴按比例 $\frac{b}{a}$ 均匀压缩后得到的图形是 \blank；
  \item 椭圆上任意一点到两焦点的距离之和等于 \blank；双曲线则对应距离之差的绝对值等于 \blank；
  \item 直线与一元二次方程联立后，判别式 $\Delta>0$ 表示有 \blank 个不同实根（交点）；
  \item 直线与平面平行的判定定理：平面外一条直线与此平面内一条直线 \blank，则线面平行。
\end{itemize}

\subsection*{教材速读}
阅读「圆锥曲线」与「立体几何」两章教材，完成下列问题：
\begin{enumerate}
  \item 椭圆定义要求 $2a$ \blank $2c$（填 $>$ 或 $<$）；双曲线要求 $2a$ \blank $2c$；
  \item 抛物线 $y^2=2px\,(p>0)$ 的焦点坐标是 \blank，准线方程是 \blank；
  \item 弦长公式 $|AB|=\sqrt{1+k^2}\cdot$ \blank；
  \item 点到平面的距离可用等体积法或 \blank 法求解。
\end{enumerate}
\end{preview}

\begin{explore}

\subsection{仿射变换：从圆到椭圆的降维打击}

\key{探究一：把椭圆「压」回圆，把难题「降」回熟题}

\key{（1）一个变换看出椭圆的身世。}
设圆 $C:x^2+y^2=a^2$。作变换
\[
  T:\begin{pmatrix}x\\y\end{pmatrix}\longmapsto\begin{pmatrix}x\\ \frac{b}{a}y\end{pmatrix},\qquad 0<b\le a.
\]
它把每一点沿 $y$ 轴方向按比例 $\frac{b}{a}$ \key{均匀压缩}：水平方向不动，竖直方向整体变扁。
设 $(x,y)\in C$，记 $(X,Y)=T(x,y)$，则 $x=X,\ y=\frac{a}{b}Y$，代入圆方程得
\[
  X^2+\left(\frac{a}{b}Y\right)^2=a^2
  \iff \frac{X^2}{a^2}+\frac{Y^2}{b^2}=1.
\]
于是 \key{椭圆就是圆在仿射变换 $T$ 下的像}。$T$ 是线性变换（矩阵 $\mathrm{diag}\left(1,\frac{b}{a}\right)$），所以它仍然把直线变成直线——这是下面所有「保持」的根源。

\key{（2）仿射不变量清单。}
因为 $T$ 把直线变成直线、且保持线段上的分比，以下事实全部保持：
\begin{itemize}
  \item \key{共线性}：三点共线 $\xrightarrow{T}$ 三点共线；
  \item \key{平行性}：$m\parallel m'\iff T(m)\parallel T(m')$；
  \item \key{中点}：线段的中点变为像线段的中点（$T$ 保持定比分点）；
  \item \key{相切}：直线与圆锥曲线相切（联立后 $\Delta=0$）的关系保持，因为 $T$ 可逆且把直线变直线；
  \item \key{面积}：任意平面图形的面积都乘以同一常数 $|\det T|=\frac{b}{a}$，故
  \[
    S_{\text{椭圆}}=\frac{b}{a}\,S_{\text{圆}}.
  \]
\end{itemize}
\key{面积比与图形的具体形状无关}，这是仿射变换最锋利的性质。

\begin{figure}[htbp]
  \centering
  \begin{tikzpicture}[>=Stealth, semithick, font=\footnotesize]
    % ---------- 左：圆与内接正三角形 ----------
    \draw[->, gray] (-2.5,0)--(2.7,0) node[below right] {$x$};
    \draw[->, gray] (0,-2.3)--(0,2.6) node[above left] {$y$};
    \fill[second!10] (0,0) circle (2);
    \draw[main!70!black, very thick] (0,0) circle (2);
    \fill[third!18] (0,2)--(-1.732,-1)--(1.732,-1)--cycle;
    \draw[third!70!black, very thick] (0,2)--(-1.732,-1)--(1.732,-1)--cycle;
    \fill (0,0) circle (1pt) node[below left] {$O$};
    \node[gray!70!black] at (1.15,-0.32) {$a$};
    \node[second!70!black] at (0,-2.02) {圆 $x^2+y^2=a^2$：$S_{\max}=\dfrac{3\sqrt3}{4}a^2$};
    % ---------- 右：压缩后的椭圆 ----------
    \begin{scope}[xshift=6.9cm]
      \draw[->, gray] (-2.5,0)--(2.7,0) node[below right] {$x$};
      \draw[->, gray] (0,-2.3)--(0,2.6) node[above left] {$y$};
      \fill[second!10] (0,0) ellipse (2 and 1.2);
      \draw[main!70!black, very thick] (0,0) ellipse (2 and 1.2);
      \fill[third!18] (0,1.2)--(-1.732,-0.6)--(1.732,-0.6)--cycle;
      \draw[third!70!black, very thick] (0,1.2)--(-1.732,-0.6)--(1.732,-0.6)--cycle;
      \fill (0,0) circle (1pt) node[below left] {$O$};
      \node[gray!70!black] at (1.15,-0.32) {$a$};
      \node[gray!70!black] at (-0.45,1.2) {$b$};
      \node[second!70!black] at (0,-2.02) {椭圆：$S_{\max}=\dfrac{b}{a}S_{\text{圆}}=\dfrac{3\sqrt3}{4}ab$};
    \end{scope}
    % ---------- 中间：压缩映射 ----------
    \draw[->, main!80!black, very thick] (3.0,2.0)--(3.7,2.0)
      node[midway,above] {$\displaystyle\binom{x}{y}\mapsto\binom{x}{\frac{b}{a}y}$};
  \end{tikzpicture}
  \caption{仿射压缩：左图为圆 $x^2+y^2=a^2$（取 $a=2$）及其内接正三角形；右图为沿 $y$ 轴按比例 $\frac{b}{a}=0.6$ 压缩后的椭圆 $\frac{x^2}{a^2}+\frac{y^2}{b^2}=1$（$b=1.2$），内接正三角形同步变形为内接三角形，面积同乘 $\frac{b}{a}$}
  \label{fig:affine-compress}
\end{figure}

\key{（3）降维打击：面积最值一步出答案。}
因为 $T$ 把所有三角形的面积同乘一个正常数，\key{「面积最大」在变换下保持}：
椭圆 $\frac{x^2}{a^2}+\frac{y^2}{b^2}=1$ 的内接三角形中，面积最大者正是圆 $x^2+y^2=a^2$ 的内接正三角形在 $T$ 下的像。
圆的内接正三角形面积为 $\frac{3\sqrt3}{4}a^2$，于是
\[
  S_{\max}=\frac{b}{a}\cdot\frac{3\sqrt3}{4}a^2=\frac{3\sqrt3}{4}ab.
\]
一个只含 $a,b$ 的干净答案——选填题可直接秒答。

\begin{example}
求椭圆 $E:\dfrac{x^2}{4}+\dfrac{y^2}{3}=1$ 的内接三角形面积的最大值。
\end{example}
\begin{solution}
这里 $a=2,\ b=\sqrt3$，压缩比为 $\frac{b}{a}=\frac{\sqrt3}{2}$。对应的圆半径为 $2$，其内接正三角形面积为
\[
  S_{\text{圆}}=\frac{3\sqrt3}{4}\cdot 2^2=3\sqrt3.
\]
故
\[
  S_{\max}=\frac{b}{a}S_{\text{圆}}=\frac{\sqrt3}{2}\cdot 3\sqrt3=\frac92
  \quad\left(\text{或直接用 }\frac{3\sqrt3}{4}ab=\frac{3\sqrt3}{4}\cdot2\cdot\sqrt3=\frac92\right).
\]
若按常规方法设三个动点坐标硬算，将陷入多元三角换元；\key{仿射变换一步就把「陌生题」变回「课本题」}。
\end{solution}

\begin{conccard}[一级]{仿射变换：椭圆是圆的压缩像}{椭圆内接三角形面积最值、共线／平行／中点类的比例问题、面积比的快速换算}{认为仿射变换保持角度与垂直；把「面积比恒定」误用于周长比；压缩后忘记新图形是椭圆而非相似形}
$T:(x,y)\mapsto\left(x,\frac{b}{a}y\right)$ 把圆 $x^2+y^2=a^2$ 变为椭圆 $\frac{x^2}{a^2}+\frac{y^2}{b^2}=1$；
共线、平行、中点、相切保持，面积同乘 $|\det T|=\frac{b}{a}$，故 $S_{\text{椭圆}}=\frac{b}{a}S_{\text{圆}}$，$S_{\triangle\max}=\frac{3\sqrt3}{4}ab$。
\end{conccard}

\begin{commonerror}{仿射变换「能保持什么、不能保持什么」}
\paragraph{失误一：以为角度也保持。}
仿射变换 $T$ 只保证「直线 $\to$ 直线、分比不变」，\key{不保持角度、不保持垂直、不保持长度}。
例如圆内接正三角形的一个内角是 $60^\circ$，压缩成椭圆的内接三角形后三个角一般不再是 $60^\circ$。
所以「椭圆 $\xrightarrow{\text{仿射}}$ 圆」可以搬用共线、平行、中点、面积比，却\textbf{不能}搬用直角与角度关系。
\paragraph{失误二：把面积比套到周长比上。}
周长在 $y$ 方向被压缩后会改变（弧长的积分并不同乘一个常数），只有\key{面积}这类二维测度才同乘 $|\det T|$。
\paragraph{失误三：忽略 $b>a$ 的情形。}
若沿 $y$ 轴拉伸（$\frac{b}{a}>1$），同样的推导给出 $\frac{x^2}{a^2}+\frac{y^2}{b^2}=1$（$b>a$），焦点在 $y$ 轴上。
公式 $S_{\max}=\frac{3\sqrt3}{4}ab$ 对两种情形都成立，但焦点位置要看 $a,b$ 谁大。
\end{commonerror}

\subsection{Dandelin 视角：圆锥曲线是三维截线}

\key{探究二：把解析几何的二次曲线放回立体几何的三维空间}

\key{（1）圆锥面与截面倾角。}
以 $z$ 轴为中心轴、半顶角为 $\alpha$ 的圆锥面方程为
\[
  x^2+y^2=z^2\tan^2\alpha .
\]
用一个平面去截它，截线的类型只取决于平面与中心轴的夹角 $\beta$（取锐角）：
\begin{itemize}
  \item $\beta>\alpha$：平面较「平」，只切断一侧锥面，得\key{椭圆}（$\beta=90^\circ$ 时退化为圆）；
  \item $\beta=\alpha$：平面恰好与一条母线平行，得\key{抛物线}；
  \item $\beta<\alpha$：平面较「陡」，同时切断两侧锥面，得\key{双曲线}。
\end{itemize}
并且离心率由这两个角完全决定：
\[
  e=\frac{\cos\beta}{\cos\alpha}.
\]
检验：$\beta=90^\circ\Rightarrow e=0$（圆）；$\beta=\alpha\Rightarrow e=1$（抛物线）；$\beta<\alpha\Rightarrow\cos\beta>\cos\alpha\Rightarrow e>1$（双曲线）。
\key{三类圆锥曲线从此不是三个定义，而是同一个对象的三种截法。}

\key{（2）Dandelin 球：把「和／差为定值」做成肉眼可见。}
在截面与锥面围成的空间内放入与圆锥面相切、又与截面相切的球（Dandelin 球），它们与截面的切点就是焦点（椭圆、双曲线用两个球，抛物线用一个球）。
设 $P$ 为截线上任一点，过 $P$ 的母线先后穿过两个切圆所在的平面。由「从一点向球引两条切线，切线长相等」：
\[
  |PF_1|=P\text{ 到第一个切圆的母线段长},\qquad |PF_2|=P\text{ 到第二个切圆的母线段长}.
\]
两段之和（椭圆）或之差（双曲线）恰为两切圆之间的母线长——这是一个与 $P$ 无关的定值。于是
\[
  \text{椭圆：}\ |PF_1|+|PF_2|=\text{定值};\qquad
  \text{双曲线：}\ \bigl||PF_1|-|PF_2|\bigr|=\text{定值}.
\]
\key{（3）统一的焦点—准线关系。}
抛物线情形只有一个球，但它已经足以给出统一定义。设 $\sigma$ 为球与圆锥的切圆所在平面，$l=\pi\cap\sigma$ 即准线（$\pi$ 为截面），$\delta$ 为 $\pi$ 与 $\sigma$ 的夹角。
设 $P$ 为截线上任一点，$h$ 为 $P$ 到 $\sigma$ 的距离，$T$ 为过 $P$ 的母线与切圆的交点。由「从一点向球引两条切线，切线长相等」：
\[
  |PF|=|PT|=\frac{h}{\cos\alpha},
  \qquad
  d(P,l)=\frac{h}{\sin\delta}.
\]
两式相除得
\[
  \frac{|PF|}{d(P,l)}=\frac{\sin\delta}{\cos\alpha}.
\]
而 $\sigma\perp$ 轴线，故 $\sin\delta=\cos\beta$，代入即得
\[
  \frac{|PF|}{d(P,l)}=\frac{\cos\beta}{\cos\alpha}=e.
\]
对椭圆、双曲线，把同一计算分别用到两个 Dandelin 球上，每一对「焦点—准线」都满足同一比值 $e$。
一般地，三种曲线统一为 \key{焦点—准线定义}：到定点（焦点）与定直线（准线）的距离之比为常数 $e$ 的点的轨迹，$e<1$ 为椭圆，$e=1$ 为抛物线，$e>1$ 为双曲线。
这正是「立体几何中的球与平面」给「解析几何中的二次曲线」发的身份证。

\begin{figure}[htbp]
  \centering
  \begin{tikzpicture}[>=Stealth, semithick, font=\footnotesize]
    % ================= 面板一：椭圆 =================
    \begin{scope}[xshift=0cm]
      \draw[dashed, gray!60] (0,-2.55)--(0,2.55);
      \draw[gray!70, thin] (0,0)--(1.162,2.492) (0,0)--(-1.162,2.492)
                         (0,0)--(1.162,-2.492) (0,0)--(-1.162,-2.492);
      \draw[second, very thick] (-2.068,0.685)--(2.068,1.415);
      \node[second!70!black] at (1.65,0.95) {截面};
      \draw[main!45!black, fill=main!8] (0,0.7347) circle (0.3106);
      \fill (0,0) circle (1pt);
      \fill[main] (-0.054,1.041) circle (1.2pt);
      \fill[main] (0.135,1.074) circle (1.2pt);
      % 两个焦点过近，合并为一个标注，并用引导线指出，避免标签互相叠字或压轴
      \node[main] at (1.05,0.45) {$\scriptstyle F_1,F_2$};
      \draw[main!55, thin] (0.78,0.52)--(0.18,1.0);
      \draw[main!80!black, thick] (1.62,2.42) ellipse (0.42 and 0.20);
      \node[main!70!black] at (0,-2.95) {$\beta>\alpha\Rightarrow$ 椭圆$\ (e<1)$};
    \end{scope}
    % ================= 面板二：抛物线 =================
    \begin{scope}[xshift=4.9cm]
      \draw[dashed, gray!60] (0,-2.55)--(0,2.55);
      \draw[gray!70, thin] (0,0)--(1.162,2.492) (0,0)--(-1.162,2.492)
                         (0,0)--(1.162,-2.492) (0,0)--(-1.162,-2.492);
      \draw[second, very thick] (-0.465,0.503)--(0.740,3.086);
      \node[second!70!black] at (0.72,0.62) {截面};
      \draw[main!45!black, fill=main!8] (0,0.75) circle (0.317);
      \fill (0,0) circle (1pt);
      \fill[main] (-0.287,0.884) circle (1.2pt) node[above] {$\scriptstyle F$};
      \draw[main!80!black, thick] (1.28,2.62) parabola bend (1.68,2.34) (2.08,2.62);
      \node[main!70!black] at (0,-2.95) {$\beta=\alpha\Rightarrow$ 抛物线$\ (e=1)$};
    \end{scope}
    % ================= 面板三：双曲线 =================
    \begin{scope}[xshift=9.8cm]
      \draw[dashed, gray!60] (0,-2.55)--(0,2.55);
      \draw[gray!70, thin] (0,0)--(1.162,2.492) (0,0)--(-1.162,2.492)
                         (0,0)--(1.162,-2.492) (0,0)--(-1.162,-2.492);
      \draw[second, very thick] (-0.330,-0.671)--(0.330,2.871);
      \node[second!70!black] at (0.42,2.55) {截面};
      \draw[main!45!black, fill=main!8] (0,0.3495) circle (0.1477);
      \draw[main!45!black, fill=main!8] (0,-0.8370) circle (0.3537);
      \fill (0,0) circle (1pt);
      \fill[main] (-0.026,0.495) circle (1.2pt) node[left] {$\scriptstyle F_1$};
      \fill[main] (-0.061,-0.489) circle (1.2pt) node[left] {$\scriptstyle F_2$};
      \draw[main!80!black, thick] (1.30,2.10) .. controls (1.78,2.32) and (1.78,2.52) .. (1.30,2.74);
      \draw[main!80!black, thick] (2.02,2.10) .. controls (1.54,2.32) and (1.54,2.52) .. (2.02,2.74);
      \node[main!70!black] at (0,-2.95) {$\beta<\alpha\Rightarrow$ 双曲线$\ (e>1)$};
    \end{scope}
  \end{tikzpicture}
  \caption{Dandelin 示意（轴截面）：同一个圆锥面被不同倾角的平面所截。$\beta>\alpha$ 得椭圆（两个 Dandelin 球各切出一个焦点），$\beta=\alpha$ 得抛物线（一个球切出焦点），$\beta<\alpha$ 得双曲线（两个球分置两侧锥面）；右上方小图标为截线实形}
  \label{fig:dandelin}
\end{figure}

\begin{example}
圆锥的半顶角 $\alpha=30^\circ$，分别用与轴线夹角 $\beta=60^\circ,30^\circ,15^\circ$ 的平面去截，判断截线类型并求离心率。
\end{example}
\begin{solution}
由 $e=\dfrac{\cos\beta}{\cos\alpha}$，其中 $\cos30^\circ=\frac{\sqrt3}{2}$：
\[
  \beta=60^\circ:\ e=\frac{\cos60^\circ}{\cos30^\circ}=\frac{1/2}{\sqrt3/2}=\frac{\sqrt3}{3}<1\Rightarrow\text{椭圆};
\]
\[
  \beta=30^\circ:\ e=\frac{\cos30^\circ}{\cos30^\circ}=1\Rightarrow\text{抛物线};
\]
\[
  \beta=15^\circ:\ e=\frac{\cos15^\circ}{\cos30^\circ}\approx\frac{0.966}{0.866}\approx1.12>1\Rightarrow\text{双曲线}.
\]
只需要比较两个角的大小，就能立刻判定类型——这是「维度打通」带来的预判力。
\end{solution}

\begin{conccard}[二级]{圆锥曲线的统一：截面与焦点—准线}{判断截线类型、求离心率、理解三类曲线定义的统一性}{把 $\beta$ 与 $\alpha$ 的大小关系记反；忘记 $\beta=90^\circ$ 是圆；以为双曲线需要「两个分开的圆锥」}
圆锥面 $x^2+y^2=z^2\tan^2\alpha$ 被与轴线夹角为 $\beta$ 的平面所截：
$\beta>\alpha$ 椭圆（$\beta=90^\circ$ 为圆），$\beta=\alpha$ 抛物线，$\beta<\alpha$ 双曲线；
统一离心率 $e=\dfrac{\cos\beta}{\cos\alpha}$，统一定义为 $\dfrac{|PF|}{d(P,l)}=e$。
\end{conccard}

\subsection{考纲安全落地：高观点预判与通法满分的双层法}

\key{探究三：上层用高观点「看透」，下层用考纲通法「写全」}

高考解析几何解答题的特点是：答案往往可以被高观点瞬间看穿，但阅卷只看考纲允许的推理链。
因此最稳的策略是\key{双层法}：先用高观点（极点极线、仿射变换）确定目标与答案，再用考纲通法（设线、联立、韦达、判别式）把过程写得滴水不漏。

\key{（1）高观点的两个预判工具。}
\begin{itemize}
  \item \key{极点极线预设定点定值}：椭圆 $\frac{x^2}{a^2}+\frac{y^2}{b^2}=1$ 外一点 $P(x_0,y_0)$ 的两条切线的切点连线（切点弦）方程为
  \[
    \frac{x_0x}{a^2}+\frac{y_0y}{b^2}=1,
  \]
  它被称为 $P$ 的\key{极线}。若 $P$ 在定直线 $x=m$ 上运动，则极线 $\frac{mx}{a^2}+\frac{y_0y}{b^2}=1$ 恒过定点 $\left(\frac{a^2}{m},0\right)$——不必动笔算，定点已出。
  \item \key{仿射变换秒算面积最值}：涉及椭圆内接图形面积的最值，先仿射回圆，用圆的结果乘 $\frac{b}{a}$ 即可。
\end{itemize}

\key{（2）考纲通法的标准四步。}
设线 $\to$ 联立 $\to$ 韦达 $\to$ 回代（并用判别式检验）。每一步都要写在卷面上，因为分步给分。

\begin{example}[两层对照示范]
椭圆 $E:\dfrac{x^2}{4}+\dfrac{y^2}{2}=1$，动点 $P$ 在定直线 $x=2$ 上，过 $P$ 作 $E$ 的两条切线，切点分别为 $A,B$。求证：直线 $AB$ 恒过定点，并求出该定点。
\end{example}
\begin{solution}
\key{第一步（高观点预判）。}设 $P(2,t)$。由极线公式，$AB$ 即 $P$ 的极线：
\[
  \frac{2x}{4}+\frac{ty}{2}=1\iff x+ty=2,
\]
它显然恒过 $(2,0)$。十秒钟，答案已经到手。

\key{第二步（考纲通法落地）。}设 $A(x_1,y_1),B(x_2,y_2)$。
由隐函数求导（第 6 章）：椭圆两边对 $x$ 求导得 $\frac{2x}{4}+\frac{2yy'}{2}=0$，即 $y'=-\frac{b^2x}{a^2y}=-\frac{x}{2y}$。
所以椭圆在 $A$ 处的切线方程为
\[
  y-y_1=-\frac{x_1}{2y_1}(x-x_1).
\]
利用 $\frac{x_1^2}{4}+\frac{y_1^2}{2}=1$ 化简，得切线的简洁形式
\[
  \frac{x_1x}{4}+\frac{y_1y}{2}=1.
\]
（若尚未学求导，也可设切线斜率 $k$，联立椭圆后令 $\Delta=0$，推出同一结果。）
因为 $P(2,t)$ 同时在两条切线上，代入得
\[
  \frac{2x_1}{4}+\frac{ty_1}{2}=1,\qquad \frac{2x_2}{4}+\frac{ty_2}{2}=1,
  \iff x_1+ty_1=2,\quad x_2+ty_2=2.
\]
这说明 $A,B$ 都在直线 $x+ty=2$ 上；两点确定一条直线，故
\[
  AB:\ x+ty=2.
\]
令 $y=0$ 得 $x=2$，与 $t$ 无关，所以 $AB$ \key{恒过定点 $(2,0)$}。

\key{对照结论}：高观点负责「知道往哪走」（预判定点 $(2,0)$），通法负责「写得干净」（求导得切线、两点定直线）。二者分工明确，考场最稳。
\end{solution}

\begin{figure}[htbp]
  \centering
  \begin{tikzpicture}[>=Stealth, semithick, font=\footnotesize]
    \draw[->, gray] (-2.9,0)--(4.4,0) node[below right] {$x$};
    \draw[->, gray] (0,-2.2)--(0,2.9) node[above left] {$y$};
    % 椭圆
    \fill[second!8] (0,0) ellipse (2 and 1.4142);
    \draw[main!70!black, very thick] (0,0) ellipse (2 and 1.4142);
    % 极线（切点弦）
    \draw[third!70!black, thick, dashed] (-0.5,2.75)--(2.0,-1.0);
    % 两条切线
    \draw[main!80!black, thick] (3,1)--($(1.771,-0.657)!-0.4!(3,1)$);
    \draw[main!80!black, thick] (3,1)--($(0.411,1.384)!-0.4!(3,1)$);
    % 过 P 的动弦
    \draw[second, thick] (-1.809,-0.603)--(3.75,1.25);
    % 点与标注
    \fill (0,0) circle (1pt);
    \fill[main] (3,1) circle (1.6pt) node[right] {$P(x_0,y_0)$};
    \fill[third] (1.771,-0.657) circle (1.6pt) node[below right] {$A$};
    \fill[third] (0.411,1.384) circle (1.6pt) node[above left] {$B$};
    \fill[second] (1.809,0.603) circle (1.6pt) node[above right] {$M$};
    \fill[second] (-1.809,-0.603) circle (1.6pt) node[below left] {$N$};
    \node[third!70!black, anchor=west] at (0.15,2.05) {$AB$：$P$ 的极线（切点弦）};
    \node[second!70!black, anchor=west] at (0.15,-1.7) {过 $P$ 的动弦 $MN$};
  \end{tikzpicture}
  \caption{极点极线与切点弦：椭圆 $\frac{x^2}{4}+\frac{y^2}{2}=1$ 外一点 $P(3,1)$，$PA,PB$ 为切线、$A,B$ 为切点，则 $AB$ 即 $P$ 的极线；过 $P$ 的任意弦 $MN$ 与极线的关系给出定点定值问题的预判。当 $P$ 在定直线 $x=m$ 上运动时，极线 $\frac{mx}{a^2}+\frac{y_0y}{b^2}=1$ 恒过定点 $\left(\frac{a^2}{m},0\right)$}
  \label{fig:pole-polar}
\end{figure}

注记（高观点对偶）：把上面的推理倒过来即得逆命题——过定点 $(2,0)$ 的直线与椭圆交于 $A,B$，则 $A,B$ 处切线的交点必在直线 $x=2$ 上。
这在极点极线理论中称为 \key{La Hire 对偶}：点 $P$ 在点 $Q$ 的极线上 $\iff$ 点 $Q$ 在点 $P$ 的极线上。

\begin{method}[双层法：高观点预判 $+$ 通法满分]
\begin{enumerate}
  \item \key{预判}：先问「答案大概长什么样」。定点定值类问题用极线公式扫一眼；面积最值类问题先仿射回圆算一遍；
  \item \key{设线}：设动直线（含参数）或设动点坐标，注意斜率不存在的情形要单独讨论；
  \item \key{联立}：与圆锥曲线方程联立，消元得到关于 $x$（或 $y$）的一元二次方程；
  \item \key{韦达}：先写 $\Delta>0$ 保证两交点，再写 $x_1+x_2,\ x_1x_2$；
  \item \key{回代}：把目标量（数量积、斜率、面积、中点）用韦达结果表示，整理成关于参数的式子；
  \item \key{定值}：令参数系数为零（或「对任意参数成立」列方程组）求出定点／定值，最后回代检验。
\end{enumerate}
\end{method}

\subsection{解析几何主干：定义、方程与直线的位置关系}

\key{探究四：把「二次曲线」拆成三件通用工具}

\key{（1）三大定义与标准方程。}
\begin{center}
\footnotesize
\setlength{\tabcolsep}{4pt}
\renewcommand{\arraystretch}{1.35}
\begin{tabular}{p{1.5cm}p{4.6cm}p{3.5cm}p{3.0cm}p{1.9cm}}
\toprule
曲线 & 定义 & 标准方程 & 焦点 & 离心率 \\
\midrule
椭圆 & $|PF_1|+|PF_2|=2a\ (2a>2c)$ & $\dfrac{x^2}{a^2}+\dfrac{y^2}{b^2}=1\ (a>b>0)$ & $(\pm c,0),\ c^2=a^2-b^2$ & $e=\dfrac ca\in(0,1)$ \\
双曲线 & $\bigl||PF_1|-|PF_2|\bigr|=2a\ (2a<2c)$ & $\dfrac{x^2}{a^2}-\dfrac{y^2}{b^2}=1$ & $(\pm c,0),\ c^2=a^2+b^2$ & $e=\dfrac ca>1$ \\
抛物线 & $|PF|=d(P,l)$ & $y^2=2px\ (p>0)$ & $\left(\dfrac p2,0\right)$，准线 $x=-\dfrac p2$ & $e=1$ \\
\bottomrule
\end{tabular}
\end{center}
三条定义其实是同一件事（到焦点与到准线的距离比为 $e$）的特例；$a,b,c$ 的关系只需记「椭圆 $c^2=a^2-b^2$、双曲线 $c^2=a^2+b^2$、抛物线只有一个 $p$」。

\key{（2）直线与圆锥曲线的位置关系：联立 $+$ 判别式。}
设直线 $l:y=kx+m$，与圆锥曲线方程联立、消去 $y$，得到关于 $x$ 的一元二次方程 $Ax^2+Bx+C=0$（$A$ 可能为 $0$），则
\[
  \Delta>0\iff \text{两个不同公共点};\quad
  \Delta=0\iff \text{相切（一个公共点）};\quad
  \Delta<0\iff \text{无公共点}.
\]
对抛物线等特殊情形，还要检查二次项系数 $A=0$（直线平行于对称轴）时的退化情况——此时方程退化为一次方程，最多一个交点。

\key{（3）弦长公式。}
设两交点为 $A(x_1,y_1),B(x_2,y_2)$，则
\[
  |AB|=\sqrt{1+k^2}\,|x_1-x_2|
  =\sqrt{1+k^2}\cdot\frac{\sqrt{\Delta}}{|A|}.
\]
若直线写成 $x=my+n$，则改用 $|AB|=\sqrt{1+m^2}\,|y_1-y_2|$。记忆钩子：\key{弦长 $=$ 水平（竖直）投影 $\times$ 放大因子}。

\key{（4）中点弦：点差法。}
设 $A(x_1,y_1),B(x_2,y_2)$ 在椭圆 $\frac{x^2}{a^2}+\frac{y^2}{b^2}=1$ 上，$M(x_0,y_0)$ 为 $AB$ 的中点。两式相减：
\[
  \frac{(x_1+x_2)(x_1-x_2)}{a^2}+\frac{(y_1+y_2)(y_1-y_2)}{b^2}=0.
\]
代入 $x_1+x_2=2x_0,\ y_1+y_2=2y_0$，得
\[
  k_{AB}=\frac{y_1-y_2}{x_1-x_2}=-\frac{b^2x_0}{a^2y_0},
  \qquad\text{等价地}\quad k_{AB}\cdot k_{OM}=-\frac{b^2}{a^2}.
\]
\key{点差法的实质是「中点 $\Rightarrow$ 斜率」}，但必须回代检验直线与椭圆确有两交点。

\begin{example}
椭圆 $E:\dfrac{x^2}{4}+\dfrac{y^2}{2}=1$ 的一条弦 $AB$ 以 $M(1,1)$ 为中点，求直线 $AB$ 的方程与弦长。
\end{example}
\begin{solution}
由点差法（$a^2=4,b^2=2,x_0=y_0=1$）：
\[
  k_{AB}=-\frac{b^2x_0}{a^2y_0}=-\frac{2}{4}=-\frac12.
\]
故直线 $AB:y-1=-\frac12(x-1)$，即
\[
  x+2y-3=0.
\]
检验（韦达）：代入椭圆得
\[
  \frac{x^2}{4}+\frac{\left(-\frac x2+\frac32\right)^2}{2}=1
  \iff 3x^2-6x+1=0,
\]
$\Delta=36-12=24>0$，确有两个交点。$x_1+x_2=2$（中点横坐标为 $1$，一致），
\[
  |x_1-x_2|=\frac{\sqrt{24}}{3}=\frac{2\sqrt6}{3},
  \qquad
  |AB|=\sqrt{1+\left(-\frac12\right)^2}\cdot\frac{2\sqrt6}{3}
  =\frac{\sqrt5}{2}\cdot\frac{2\sqrt6}{3}=\frac{\sqrt{30}}{3}.
\]
\end{solution}

\begin{conccard}[二级]{直线与圆锥曲线：判别式、弦长、中点弦}{位置关系判定、弦长计算、中点弦方程、弦中点轨迹}{忽略二次项系数为 $0$ 的退化；点差法不检验 $\Delta$；弦长公式漏乘 $\sqrt{1+k^2}$ 或把分母写成 $\sqrt\Delta$}
联立得 $Ax^2+Bx+C=0$：$\Delta>0$ 相交、$\Delta=0$ 相切、$\Delta<0$ 相离；
$|AB|=\sqrt{1+k^2}\cdot\frac{\sqrt\Delta}{|A|}$；
中点弦斜率 $k=-\dfrac{b^2x_0}{a^2y_0}$，且 $k\cdot k_{OM}=-\dfrac{b^2}{a^2}$。
\end{conccard}

\begin{commonerror}{直线与圆锥曲线联立的三个高频失误}
\paragraph{失误一：忘记讨论二次项系数为 $0$。}
消元后若 $x^2$ 项系数含参数（如直线与抛物线联立），它可能为 $0$，此时方程退化为一元一次方程，只有一个交点。
必须先声明「当 $A=0$ 时……；当 $A\ne0$ 时……」。
\paragraph{失误二：点差法求出斜率就直接写答案。}
点差法只保证「若 $M$ 是中点，则斜率必为 $k$」，并不保证这样的弦真的存在。
必须回代验证 $\Delta>0$（上面例题中 $\Delta=24>0$ 才可写）。
\paragraph{失误三：弦长公式张冠李戴。}
$|x_1-x_2|=\frac{\sqrt\Delta}{|A|}$ 只是水平投影；弦长还要乘 $\sqrt{1+k^2}$。
若直线设为 $x=my+n$，则应乘 $\sqrt{1+m^2}$ 并对 $y$ 用韦达，切勿混用。
\end{commonerror}

\begin{conccard}[三级]{定点定值问题的通法骨架}{含参数的动直线与圆锥曲线相交，证明某点／某量为定值}{只算特殊位置就下结论；忽略判别式；参数没有全部消去就说「定值」}
设线（含参数）$\to$ 联立 $\to$ 韦达（先 $\Delta>0$）$\to$ 把目标量表示成参数的式子 $\to$ 令各参数系数为零（或对任意参数成立）$\to$ 解出定点／定值并回代检验。
高观点（极线、仿射）用于提前预判答案，通法用于写满步骤。
\end{conccard}

\subsection{立体几何主干：位置关系、空间角与向量法}

\key{探究五：判定靠逻辑，计算靠向量}

\key{（1）线面平行与垂直：判定和性质是一体两面。}
\begin{itemize}
  \item \key{线面平行}：平面外一条直线与此平面内一条直线平行 $\implies$ 线面平行（判定）；
        线面平行 $\implies$ 过该直线的任一平面与原平面的交线与该直线平行（性质）。
  \item \key{线面垂直}：一条直线与平面内\key{两条相交}直线都垂直 $\implies$ 线面垂直（判定）；
        线面垂直 $\implies$ 垂直于平面内任意直线（性质）；垂直于同一平面的两直线互相平行。
  \item \key{面面关系}：线面平行 $+$ 线线平行可推面面平行；线面垂直可推面面垂直。它们是上述两个判定的「升级版」。
\end{itemize}
解题时的固定动作：\key{要证平行，先找「面内一条线与它平行」；要证垂直，先找「面内两条相交线与它垂直」}。

\key{（2）空间角与距离的定义。}
\begin{itemize}
  \item \key{线面角}：斜线与它在平面内的射影所成的角，范围 $\left[0,\frac\pi2\right]$；线面垂直时为 $\frac\pi2$，线面平行或在面内时为 $0$；
  \item \key{二面角}：在棱上取一点，在两个半平面内分别作垂直于棱的射线，两射线的夹角即二面角的平面角；
  \item \key{点到平面的距离}：点到平面的垂线段长。几何法常用\key{等体积法}（换底面，体积不变）；向量法用公式 $d=\frac{|\overrightarrow{AP}\cdot\vec n|}{|\vec n|}$。
\end{itemize}

\key{（3）向量法：把第 3 章的工具全部调用过来。}
以方向向量 $\vec s$ 与法向量 $\vec n$ 为「两个抓手」，空间位置关系与度量可全部坐标化。
线面平行／垂直、线面角、二面角四条判据与三个高频失误已在第 3 章
完整给出（点面距公式也出自同一套向量法），此处不重复；
本章只补一条应用——点到平面的距离
\[
  d=\frac{|\overrightarrow{AP}\cdot\vec n|}{|\vec n|}\quad(A\in\alpha),
\]
即 $\overrightarrow{AP}$ 在法向量方向上投影的绝对值；几何法常用等体积法对照。

\begin{example}
四棱锥 $P-ABCD$ 中，底面 $ABCD$ 是边长为 $2$ 的正方形，$PA\perp$ 底面 $ABCD$，$PA=2$。
（1）求证：$BD\perp$ 平面 $PAC$；（2）求点 $A$ 到平面 $PCD$ 的距离；（3）求二面角 $A-PC-D$ 的余弦值。
\end{example}
\begin{solution}
以 $A$ 为原点，$\overrightarrow{AB},\overrightarrow{AD},\overrightarrow{AP}$ 分别为 $x,y,z$ 轴建立空间直角坐标系（三条轴两两垂直），则
\[
  A(0,0,0),\ B(2,0,0),\ C(2,2,0),\ D(0,2,0),\ P(0,0,2).
\]
（1）$\overrightarrow{BD}=(-2,2,0)$，$\overrightarrow{AC}=(2,2,0)$，$\overrightarrow{AP}=(0,0,2)$。
\[
  \overrightarrow{BD}\cdot\overrightarrow{AC}=-4+4=0,\qquad
  \overrightarrow{BD}\cdot\overrightarrow{AP}=0,
\]
故 $BD\perp AC$ 且 $BD\perp AP$。又 $AC\cap AP=A$，所以 $BD\perp$ 平面 $PAC$。

（2）$\overrightarrow{PC}=(2,2,-2)$，$\overrightarrow{PD}=(0,2,-2)$，平面 $PCD$ 的法向量
\[
  \vec n=\overrightarrow{PC}\times\overrightarrow{PD}=(0,4,4)\parallel(0,1,1).
\]
点 $A$ 到平面 $PCD$ 的距离
\[
  d=\frac{|\overrightarrow{AD}\cdot\vec n|}{|\vec n|}
  =\frac{|(0,2,0)\cdot(0,1,1)|}{\sqrt2}=\frac{2}{\sqrt2}=\sqrt2.
\]
几何法对照（等体积法）：
$V_{P-ACD}=\frac13 S_{\triangle ACD}\cdot PA=\frac13\cdot2\cdot2=\frac43$；
又 $PD^2+CD^2=8+4=12=PC^2$，故 $\triangle PCD$ 在 $D$ 处为直角，$S_{\triangle PCD}=\frac12\cdot2\sqrt2\cdot2=2\sqrt2$。
由 $\frac13 S_{\triangle PCD}\cdot h=\frac43$ 得 $h=\sqrt2$，与向量法一致。

（3）平面 $APC$ 的法向量 $\vec n_1=\overrightarrow{AP}\times\overrightarrow{AC}=(-4,4,0)\parallel(-1,1,0)$；平面 $PCD$ 的法向量 $\vec n_2=(0,1,1)$。
\[
  |\cos\theta|=\frac{|\vec n_1\cdot\vec n_2|}{|\vec n_1|\,|\vec n_2|}
  =\frac{|0+1+0|}{\sqrt2\cdot\sqrt2}=\frac12.
\]
取 $PC$ 的中点 $M(1,1,1)$，把 $\overrightarrow{MA}=(-1,-1,-1)$、$\overrightarrow{MD}=(-1,1,-1)$ 去掉沿 $\overrightarrow{PC}$ 方向的分量，得到二面角平面角的两条边，计算得其夹角余弦为 $\frac12>0$，故二面角 $A-PC-D$ 为锐角，其余弦值为
\[
  \cos\theta=\frac12.
\]
传统几何法需要作平面角、证垂直、再解三角形；向量法只需建系、算两个法向量、代公式。
\end{solution}

\begin{conccard}[一级]{立体几何的「判定—性质」闭环}{证平行／垂直时的辅助线设计}{只会背定理不会用；找错「面内两条相交直线」；用缺少相交条件的向量}
证线面平行：面外线 $\parallel$ 面内线；证线面垂直：直线 $\perp$ 面内两条相交直线；
线面垂直 $\Rightarrow$ 线线垂直；线面平行 $+$ 相交平面 $\Rightarrow$ 线线平行。
几何法与向量法各有胜场：垂直关系明显时优先几何法（等体积、三垂线），坐标好写时优先向量法。
\end{conccard}

\begin{summary}
本章用「仿射变换 $+$ 三维截面 $+$ 向量法」三件武器打通了立体几何与解析几何的维度壁垒。
椭圆是圆沿 $y$ 轴按比例 $\frac{b}{a}$ 的均匀压缩像，仿射不变量（共线、平行、中点、相切、面积比）全部保持，故椭圆内接三角形最大面积为 $\frac{3\sqrt3}{4}ab$；
Dandelin 球说明椭圆／双曲线／抛物线是同一个圆锥面被不同倾角平面所截的结果，统一为焦点—准线定义，离心率 $e=\frac{\cos\beta}{\cos\alpha}$；
考场上层的「高观点预判」（极线公式预设定点、仿射变换秒算面积）与下层的「考纲通法」（设线—联立—韦达—判别式）合成双层法，既看得透又写得全；
解析几何主干落实为定义与标准方程表、判别式判位置、弦长公式与中点弦点差法，以及定点定值问题的六步通法；
立体几何主干落实为线面平行／垂直的判定与性质、空间角与距离，并由第 3 章的方向向量与法向量统一处理。
\key{维度不是障碍，变换才是本质；答案是预判出来的，步骤是通法写出来的。}
\end{summary}

\end{explore}

\begin{inclasstest}
\begin{enumerate}
  \item 把圆 $x^2+y^2=9$ 沿 $y$ 轴按比例 $\frac23$ 均匀压缩，写出所得椭圆方程，并求其内接三角形面积的最大值；\level{易}
  \item 求抛物线 $y^2=8x$ 的焦点坐标与准线方程；\level{易}
  \item 若直线 $y=x+m$ 与椭圆 $\dfrac{x^2}{4}+\dfrac{y^2}{2}=1$ 有两个不同的交点，求实数 $m$ 的取值范围；当两者相切时 $m$ 取何值？\level{中}
  \item 三棱锥 $O-ABC$ 中，$OA,OB,OC$ 两两垂直，$OA=1,OB=2,OC=3$，求直线 $OC$ 与平面 $ABC$ 所成角的正弦值。\level{中}
\end{enumerate}
\end{inclasstest}

\begin{afterclass}
\subsection*{基础巩固}
\begin{enumerate}
  \item 求椭圆 $\dfrac{x^2}{16}+\dfrac{y^2}{12}=1$ 的半长轴、半短轴、半焦距、焦点坐标与离心率；椭圆上任意一点到两焦点距离之和是多少？
  \item 求双曲线 $\dfrac{x^2}{9}-\dfrac{y^2}{16}=1$ 的实轴长、虚轴长、离心率与渐近线方程。
\end{enumerate}

\subsection*{进阶提升}
\begin{enumerate}
  \item 椭圆 $E:\dfrac{x^2}{4}+y^2=1$，动点 $P$ 在定直线 $x=4$ 上，过 $P$ 作 $E$ 的两条切线，切点分别为 $A,B$，求证：直线 $AB$ 恒过定点，并求出该定点。
  \item 椭圆 $E:\dfrac{x^2}{9}+\dfrac{y^2}{4}=1$ 中，以 $M\left(\dfrac32,1\right)$ 为中点的弦所在直线方程为何？该弦的长度是多少？
\end{enumerate}

\subsection*{拓展探究}
\begin{enumerate}
  \item （极点极线的一般化）已知椭圆 $E:\dfrac{x^2}{a^2}+\dfrac{y^2}{b^2}=1$，$P(x_0,y_0)$ 为 $E$ 外一点，过 $P$ 作两条切线，切点分别为 $A,B$。
  \begin{enumerate}
    \item 求证：切点弦（极线）$AB$ 的方程为 $\dfrac{x_0x}{a^2}+\dfrac{y_0y}{b^2}=1$；
    \item 若 $P$ 在定直线 $x=m\ (m\neq0)$ 上运动，求 $AB$ 恒过的定点，并用第 6 章的隐函数求导检验 $a=2,b=\sqrt2,x_0=2$ 的特例。
  \end{enumerate}
  （提示：设切点 $(x_1,y_1)$，用隐函数求导求切线斜率，化简时利用 $\frac{x_1^2}{a^2}+\frac{y_1^2}{b^2}=1$。）
\end{enumerate}

\subsection*{教师版答案}
\begin{solution}
\textbf{基础 1：}$a^2=16,b^2=12\Rightarrow a=4,b=2\sqrt3$；$c^2=a^2-b^2=4\Rightarrow c=2$。
焦点 $(\pm2,0)$，离心率 $e=\dfrac ca=\dfrac12$；椭圆上任意一点到两焦点距离之和为 $2a=8$。
\end{solution}
\begin{solution}
\textbf{基础 2：}$a^2=9,b^2=16\Rightarrow a=3,b=4$；$c^2=a^2+b^2=25\Rightarrow c=5$。
实轴长 $2a=6$，虚轴长 $2b=8$，离心率 $e=\dfrac53$，渐近线 $y=\pm\dfrac{b}{a}x=\pm\dfrac43x$。
\end{solution}
\begin{solution}
\textbf{进阶 1：}设 $P(4,t)$。由极线公式（或设切点用求导推导），切点弦 $AB$ 的方程为
\[
  \frac{4x}{4}+\frac{ty}{1}=1\iff x+ty=1.
\]
令 $y=0$ 得 $x=1$，与 $t$ 无关，故 $AB$ 恒过定点 $(1,0)$。
\end{solution}
\begin{solution}
\textbf{进阶 2：}由点差法，$k_{AB}=-\dfrac{b^2x_0}{a^2y_0}=-\dfrac{4\cdot\frac32}{9\cdot1}=-\dfrac23$。
直线 $AB:y-1=-\dfrac23\left(x-\dfrac32\right)$，即
\[
  2x+3y-6=0.
\]
代入椭圆：$\dfrac{x^2}{9}+\dfrac{\left(2-\frac23x\right)^2}{4}=1
\iff \dfrac29x^2-\dfrac23x=0$，得 $x=0$ 或 $x=3$。
对应交点为 $(0,2)$ 与 $(3,0)$（恰为椭圆的两个顶点），中点确为 $\left(\dfrac32,1\right)$，弦长
\[
  |AB|=\sqrt{(3-0)^2+(0-2)^2}=\sqrt{13}.
\]
\end{solution}
\begin{solution}
\textbf{拓展：（1）}设 $A(x_1,y_1)$。椭圆两边对 $x$ 求导：$\dfrac{2x}{a^2}+\dfrac{2yy'}{b^2}=0\Rightarrow y'=-\dfrac{b^2x}{a^2y}$。
故 $A$ 处切线为 $y-y_1=-\dfrac{b^2x_1}{a^2y_1}(x-x_1)$，整理：
\[
  a^2y_1y-a^2y_1^2=-b^2x_1x+b^2x_1^2
  \iff b^2x_1x+a^2y_1y=b^2x_1^2+a^2y_1^2
  =a^2b^2\left(\frac{x_1^2}{a^2}+\frac{y_1^2}{b^2}\right)=a^2b^2,
\]
两边除以 $a^2b^2$ 得切线方程 $\dfrac{x_1x}{a^2}+\dfrac{y_1y}{b^2}=1$。
因为 $P(x_0,y_0)$ 在切线上，所以 $\dfrac{x_1x_0}{a^2}+\dfrac{y_1y_0}{b^2}=1$；同理 $\dfrac{x_2x_0}{a^2}+\dfrac{y_2y_0}{b^2}=1$。
这表明 $A,B$ 都在直线 $\dfrac{x_0x}{a^2}+\dfrac{y_0y}{b^2}=1$ 上，两点定一直线，即为 $AB$ 的方程。

\textbf{（2）}令 $x_0=m$，极线为 $\dfrac{mx}{a^2}+\dfrac{y_0y}{b^2}=1$。令 $y=0$ 得 $x=\dfrac{a^2}{m}$，与 $y_0$ 无关，故 $AB$ 恒过定点 $\left(\dfrac{a^2}{m},0\right)$。
特例检验：$a=2,b=\sqrt2,m=2$，定点为 $\left(\dfrac{4}{2},0\right)=(2,0)$；极线方程 $\dfrac{2x}{4}+\dfrac{y_0y}{2}=1$ 即 $x+y_0y=2$，当 $y_0=1$ 时 $x+y=2$ 过 $(2,0)$，与课堂例题一致。
\end{solution}
\reading{拓展阅读 10}{圆锥曲线：同一个三维对象的三种截面}
\end{afterclass}

% 高观点导学案：数列与离散动态系统（第 5 章·高观点版）
% 主线：差分与求和互逆（离散微积分）→ 递推即离散动力系统（不动点与蛛网图）→ 线性递推与特征值（矩阵同构）
%       → Abel 变换（错位相减的机械化）→ 共轭变换（不动点为何生效）→ 生成函数（强基拓展，超出考纲）
% 核心思想：数列不是公式的仓库，而是离散动力系统在时间轴上的轨迹；解题的本质是换坐标系，
%           把陌生递推化为等比／等差／代数方程。

\section{数列与离散动态系统（高观点版）}

\begin{methodbox}{多框架对读}
\small
\begin{center}
\begin{tabular}{c|p{2.9cm}|p{5.0cm}|p{3.1cm}}
\toprule
视角 & 本章看点 & 别的视角怎么说 & 哪侧最省力 \\
\midrule
代数 & 通项公式、求和与差分 & 符号对象的恒等变形 & 硬算、裂项、Abel \\
几何 & 蛛网图与特征螺线 & 轨迹在坐标平面上的形状 & 看收敛、猜极限 \\
集合 & $\mathbb{N}\to\mathbb{R}$ 的离散函数 & 轨道是值的集合 & 审定义域与边界 \\
映射 & $x_{n+1}=F(x_n)$、共轭变换 & 换坐标系＝换表示 & 非线性 $\to$ 等比／等差 \\
\bottomrule
\end{tabular}
\end{center}
卡的使用：递推题先问「这是哪个映射的自迭代」，再决定搬去几何还是代数（总图见第 0 章图~\ref{fig:four-view-map}）。
\end{methodbox}

第 3 章把向量确立为「几何空间代数化」的通用语言，第 4 章用仿射变换打通了圆锥曲线的维度壁垒；
本章回到最朴素的对象——数列，把它从「套公式」升级为「离散动力系统 $+$ 降维变换」：
数列是定义域为 $\mathbb{N}^*$ 的离散函数，递推是离散时间的状态演化，求和与差分是离散的微积分，
而不动点、特征值、共轭、Abel、生成函数则是把复杂递推「降维」成等比／等差／代数方程的变换。
六个模块层层展开：
\key{差分与求和互逆}（离散版微积分基本定理，由 $S_n$ 反求通项、差分降次）$\to$
\key{递推即离散动力系统}（不动点、蛛网图、$|f'(x^*)|$ 稳定性判据）$\to$
\key{线性递推与特征值}（转移矩阵 $\begin{pmatrix}0&1\\q&p\end{pmatrix}$、特征方向即不动点方向）$\to$
\key{Abel 变换}（错位相减的机械化，与分部积分互译）$\to$
\key{共轭变换}（不动点为何生效：换坐标系、莫比乌斯分式）$\to$
\key{生成函数}（强基拓展，\key{超出高考考纲}）。
核心思想：\key{数列不是公式的仓库，而是离散动力系统在时间轴上的轨迹；解题的本质是换坐标系，把陌生递推化为等比／等差。}

\begin{learnobjective}
\begin{itemize}
  \item 理解前向差分 $\Delta a_n=a_{n+1}-a_n$ 与求和 $\sum$ 互为逆运算，掌握离散版微积分基本定理 $\sum_{k=1}^n\Delta a_k=a_{n+1}-a_1$，会由和式 $S_n$ 反求通项；
  \item 理解「差分降次」与导数降次的类比，会看差分表判断通项的次数；
  \item 会把一阶递推 $a_{n+1}=f(a_n)$ 看作离散动力系统，会解不动点 $x^*=f(x^*)$、画蛛网图，并用 $|f'(x^*)|$ 判稳定性（$<1$ 吸引、$>1$ 排斥）；
  \item 掌握二阶线性递推的特征方程 $\lambda^2-p\lambda-q=0$，理解它与转移矩阵 $\begin{pmatrix}0&1\\q&p\end{pmatrix}$ 的特征值同构，特征方向即分式递推的不动点方向；
  \item 会用 Abel 变换 $\sum_{k=1}^na_kb_k=a_nB_n-\sum_{k=1}^{n-1}(a_{k+1}-a_k)B_k$ 机械化求「等差 $\times$ 等比」型求和；
  \item 理解共轭变换（换坐标系）是不动点法生效的底层原因，会作 $b_n=a_n-x^*$ 的等比化与 $\dfrac{a_n-p}{a_n-q}$ 的分式化；
  \item 了解生成函数 $G(x)=\sum a_nx^n$（\key{超出高考考纲，仅竞赛／强基}）。
\end{itemize}
\end{learnobjective}

\begin{preview}
\subsection*{知识回顾}
\begin{itemize}
  \item 等差数列的通项公式是 \blank，前 $n$ 项和公式是 \blank；
  \item 等比数列的通项公式是 \blank，前 $n$ 项和公式是 \blank；
  \item 设 $S_n$ 是数列 $\{a_n\}$ 的前 $n$ 项和，则 $n\ge2$ 时 $a_n=$ \blank，且必须单独验证 \blank；
  \item 函数 $y=f(x)$ 的图像与直线 $y=x$ 的交点横坐标满足方程 \blank。
\end{itemize}

\subsection*{教材速读}
阅读「数列」一章教材，完成下列问题：
\begin{enumerate}
  \item 由 $S_n$ 求 $a_n$ 的「分解式」$a_n=\begin{cases}S_1,&n=1\\ S_n-S_{n-1},&n\ge2\end{cases}$ 中，$n=1$ 时 \blank（填「需要」或「不需要」）单独检验；
  \item 一阶线性递推 $a_{n+1}=pa_n+q\ (p\ne1)$ 的不动点 $x^*$ 满足 $x^*=$ \blank；
  \item 斐波那契数列 $F_1=F_2=1,\ F_{n+2}=F_{n+1}+F_n$ 的特征方程是 \blank；
  \item 「错位相减法」主要适用于 \blank 数列与 \blank 数列相乘型求和。
\end{enumerate}
\end{preview}

\begin{explore}

\subsection{差分与求和互逆：离散版微积分基本定理}

\key{探究一：$\Delta$ 与 $\sum$ 是同一个东西的两面}

\key{（1）前向差分：数列的「离散导数」。}
定义
\[
  \Delta a_n=a_{n+1}-a_n .
\]
$\Delta$ 把数列 $\{a_n\}$ 映成新数列 $\{\Delta a_n\}$；二阶差分 $\Delta^2a_n=\Delta(\Delta a_n)=a_{n+2}-2a_{n+1}+a_n$，以此类推。
以 $a_n=n^2$ 为例，列表如下：
\[
\begin{array}{c|ccccc}
  n & 1 & 2 & 3 & 4 & 5 \\ \hline
  a_n=n^2 & 1 & 4 & 9 & 16 & 25 \\
  \Delta a_n & 3 & 5 & 7 & 9 & \\
  \Delta^2a_n & 2 & 2 & 2 & &
\end{array}
\]
一阶差分 $\Delta a_n=2n+1$ 是一次等差数列，二阶差分 $\Delta^2a_n=2$ 是常数——每差一次，次数降一次。

\key{（2）差分与求和互逆：离散版微积分基本定理。}
因为求和时中间项全部抵消（telescoping）：
\[
  \sum_{k=1}^{n}\Delta a_k=\sum_{k=1}^{n}(a_{k+1}-a_k)=a_{n+1}-a_1 .
\]
这就是离散版的「牛顿—莱布尼茨公式」：$\sum$ 与 $\Delta$ 互为逆运算。
反过来，若已知前 $n$ 项和 $S_n=\sum_{k=1}^na_k$，则
\[
  a_n=S_n-S_{n-1}\quad(n\ge2),\qquad a_1=S_1 .
\]
所以「由和式反求通项」不是技巧，而是差分的逆运算；注意 $n=1$ 要单独检验。

\key{（3）差分降次：类比导数。}
若 $a_n$ 是 $k$ 次多项式，则 $\Delta a_n$ 是 $k-1$ 次多项式（类比 $(x^k)'=kx^{k-1}$）。由此得到判据：
\key{若 $\Delta^2a_n$ 恒为常数，则 $a_n$ 是二次式；若 $\Delta a_n$ 是等差数列，则 $a_n$ 是二次数列}。
差分表因此可以「降次」，把高次结构一层层剥成常数。

\begin{example}
已知 $S_n=n^2+2n$，求 $\{a_n\}$ 的通项公式。
\end{example}
\begin{solution}
$a_1=S_1=1+2=3$；当 $n\ge2$ 时
\[
  a_n=S_n-S_{n-1}=(n^2+2n)-\bigl((n-1)^2+2(n-1)\bigr)=2n+1 .
\]
$n=1$ 时 $2n+1=3=S_1$，可并入，故
\[
  a_n=2n+1 .
\]
差分表视角：$\Delta a_n=2$ 为常数，故 $a_n$ 是一次等差数列，首项 $3$、公差 $2$。
\end{solution}

\begin{example}
用「差分求和」求 $\sum_{k=1}^nk^2$。
\end{example}
\begin{solution}
利用恒等式
\[
  (k+1)^3-k^3=3k^2+3k+1 ,
\]
左边正是 $k^3$ 的差分结构。对 $k=1,2,\dots,n$ 求和，左边 telescoping：
\[
  \sum_{k=1}^n\bigl((k+1)^3-k^3\bigr)=(n+1)^3-1 ,
\]
右边为
\[
  \sum_{k=1}^n(3k^2+3k+1)=3\sum_{k=1}^nk^2+\frac{3n(n+1)}2+n .
\]
联立解得
\[
  \sum_{k=1}^nk^2=\frac{n(n+1)(2n+1)}6 .
\]
关键：把 $k^2$ 写成「差分的展开」，求和就自动相消。
\end{solution}

\begin{conccard}[一级]{差分与求和互逆：$\Delta\leftrightarrow\sum$}{由 $S_n$ 反求通项、裂项相消、判断通项次数}{忘记 $n=1$ 单独检验；把 $\Delta a_n$ 的符号与 $a_n$ 的符号混同；误以为 $\Delta$ 保持多项式次数}
$\Delta a_n=a_{n+1}-a_n$；$\sum_{k=1}^n\Delta a_k=a_{n+1}-a_1$（离散牛顿—莱布尼茨）；
$a_n=S_n-S_{n-1}\ (n\ge2)$，$a_1=S_1$；
$k$ 次多项式差分一次降一次，$\Delta^2a_n$ 为常数 $\iff a_n$ 为二次式。
\end{conccard}

\subsection{递推即离散动力系统：不动点与蛛网图}

（不动点的统一理论——连续情形、几何变换与射影看法——见第 9 章；
本节只取数列递推这一路应用。）

\key{探究二：把 $a_{n+1}=f(a_n)$ 看成状态演化}

\key{（1）离散动力系统。}
数列 $\{a_n\}$ 是定义域为 $\mathbb{N}^*$ 的函数；一阶递推
\[
  a_{n+1}=f(a_n)
\]
表示：已知当前状态 $a_n$，下一步状态由同一个规则 $f$ 给出。重复作用 $f$ 就得到轨迹 $a_1,a_2,a_3,\dots$。
这正是「离散时间、离散状态」的动力系统。

\key{（2）不动点：系统的平衡态。}
若
\[
  x^*=f(x^*),
\]
则从 $x^*$ 出发永远停在 $x^*$：$f(x^*)=x^*,f(f(x^*))=x^*,\dots$。
几何上，$x^*$ 是曲线 $y=f(x)$ 与直线 $y=x$ 的交点横坐标。
对 $a_{n+1}=pa_n+q\ (p\ne1)$，不动点为 $x^*=\dfrac{q}{1-p}$。

\key{（3）蛛网图：把迭代画成阶梯。}
画法：从 $x$ 轴上的 $a_1$ 出发，竖直画到曲线 $y=f(x)$ 得 $(a_1,a_2)$；水平画到直线 $y=x$ 得 $(a_2,a_2)$；
再竖直画到曲线得 $(a_2,a_3)$；如此反复。折线像「蜘蛛网」一样阶梯式逼近或远离不动点。

\key{（4）稳定性判据：$|f'(x^*)|$ 决定吸引还是排斥。}
在不动点附近写 $a_n=x^*+\varepsilon_n$，则
\[
  \varepsilon_{n+1}=f(x^*+\varepsilon_n)-x^*\approx f'(x^*)\,\varepsilon_n .
\]
于是偏差按因子 $\lambda=f'(x^*)$ 缩放：
\begin{itemize}
  \item $|f'(x^*)|<1$：\key{吸引（稳定）}，偏差逐次缩小，轨迹收敛到 $x^*$；
  \item $|f'(x^*)|>1$：\key{排斥（不稳定）}，偏差逐次放大，轨迹逃离 $x^*$；
  \item $f'(x^*)<0$ 时偏差每步变号，轨迹\key{摆动}收敛；$0<f'(x^*)<1$ 时偏差不变号，轨迹\key{单调}收敛。
\end{itemize}
对线性 $f(x)=px+q$，近似就是精确：$b_{n+1}=pb_n$，即等比数列。

\begin{figure}[htbp]
  \centering
  \begin{tikzpicture}[>=Stealth, semithick, font=\footnotesize]
    % ---------- 左：吸引不动点 ----------
    \begin{scope}[xshift=0cm]
      \draw[->, gray] (-0.2,0)--(5.6,0) node[below right] {$x$};
      \draw[->, gray] (0,-0.2)--(0,5.2) node[above left] {$y$};
      \draw[gray!60, dashed] (-0.2,-0.2)--(5.2,5.2) node[right] {$y=x$};
      \draw[main!70!black, very thick] (-0.1,1.95)--(5.4,4.7) node[right] {$y=\frac12x+2$};
      \fill[second] (4,4) circle (1.8pt) node[above left] {$x^*=4$};
      \draw[third!70!black, thick] (0.5,0)--(0.5,2.25)--(2.25,2.25)--(2.25,3.125)
        --(3.125,3.125)--(3.125,3.5625)--(3.5625,3.5625);
      \fill[third] (0.5,0) circle (1.2pt) node[below] {$a_1$};
      \fill[third] (2.25,2.25) circle (1.2pt) node[below right] {$a_2$};
      \fill[third] (3.125,3.125) circle (1.2pt) node[below right] {$a_3$};
      \node[main!70!black] at (2.6,-1.0) {$|f'(4)|=\frac12<1$：阶梯收缩，收敛};
    \end{scope}
    % ---------- 右：排斥不动点 ----------
    \begin{scope}[xshift=7.4cm]
      \draw[->, gray] (-0.2,0)--(5.6,0) node[below right] {$x$};
      \draw[->, gray] (0,-0.2)--(0,5.2) node[above left] {$y$};
      \draw[gray!60, dashed] (-0.2,-0.2)--(5.2,5.2) node[right] {$y=x$};
      \draw[main!70!black, very thick] (0.75,0)--(3.6,4.56) node[right] {$y=1.6x-1.2$};
      \fill[second] (2,2) circle (1.8pt) node[above left] {$x^*=2$};
      \draw[third!70!black, thick] (2.25,0)--(2.25,2.4)--(2.4,2.4)--(2.4,2.64)
        --(2.64,2.64)--(2.64,3.024)--(3.024,3.024);
      \fill[third] (2.25,0) circle (1.2pt) node[below] {$a_1$};
      \fill[third] (2.64,2.64) circle (1.2pt) node[below right] {$a_3$};
      \fill[third] (3.024,3.024) circle (1.2pt) node[below right] {$a_4$};
      \node[main!70!black] at (2.6,-1.0) {$|g'(2)|=1.6>1$：阶梯放大，发散};
    \end{scope}
  \end{tikzpicture}
  \caption{蛛网图：从 $a_1$ 出发，竖直打到 $y=f(x)$、水平反射到 $y=x$，如此阶梯迭代。左图 $|f'(x^*)|<1$，阶梯收缩到不动点；右图 $|g'(x^*)|>1$，阶梯放大逃离不动点}
  \label{fig:cobweb}
\end{figure}

\begin{example}
设 $a_1=0,\ a_{n+1}=\dfrac12a_n+2$，求通项并判断数列的收敛性。
\end{example}
\begin{solution}
不动点满足 $x^*=\frac12x^*+2$，解得 $x^*=4$。两式相减：
\[
  a_{n+1}-4=\frac12(a_n-4) .
\]
故 $\{a_n-4\}$ 是首项 $a_1-4=-4$、公比 $\frac12$ 的等比数列：
\[
  a_n-4=-4\left(\frac12\right)^{n-1}=-2^{3-n},
  \qquad
  a_n=4-2^{3-n} .
\]
因 $|f'(4)|=\frac12<1$，不动点 $4$ 吸引，$a_n\to4$（单调递增）。
\end{solution}

\begin{conccard}[二级]{不动点与稳定性：$|f'(x^*)|$ 判据}{一阶递推 $a_{n+1}=f(a_n)$ 的收敛性预判、蛛网图定性、线性化降维}{把「令 $a_{n+1}=a_n=x^*$」当成假设相邻项相等；解出 $x^*$ 不检验；对 $f'(x^*)$ 不存在或 $|f'(x^*)|=1$ 的临界情形乱下结论}
不动点 $x^*=f(x^*)$ 是曲线 $y=f(x)$ 与 $y=x$ 的交点；
$b_n=a_n-x^*$ 满足 $b_{n+1}\approx f'(x^*)b_n$；
$|f'(x^*)|<1$ 吸引（稳定），$|f'(x^*)|>1$ 排斥（不稳定）；$f'(x^*)<0$ 摆动收敛。
\end{conccard}

\begin{commonerror}{不动点与稳定性判据的三个高频失误}
\paragraph{失误一：以为「令 $a_{n+1}=a_n=x^*$」是假设相邻两项相等。}
不动点 $x^*$ 是\key{解方程 $x^*=f(x^*)$ 得到的平衡点}，不是把相邻两项强行看成相等。
它的真实作用是「作差消去常数项」（见探究五）。
\paragraph{失误二：不验证 $x^*$ 是否真是不动点。}
解出 $x^*$ 后必须回代 $f(x^*)=x^*$；尤其是分式递推，使分母为零的「解」要舍去。
\paragraph{失误三：乱用 $|f'(x^*)|$ 判据。}
该判据要求 $f$ 在 $x^*$ 附近可导且 $x^*$ 是不动点；对 $f'(x^*)$ 不存在（如 $f(x)=\sqrt x$ 在 $0$ 处）
或 $|f'(x^*)|=1$ 的临界情形，结论要另行分析。
\end{commonerror}

\subsection{线性递推与特征值：状态向量与矩阵同构}

\key{探究三：把递推升维成矩阵乘法}

\key{（1）状态向量：二阶递推 $=$ 二维线性系统。}
二阶线性递推
\[
  a_{n+2}=pa_{n+1}+qa_n
\]
等价于状态向量 $\mathbf v_n=\dbinom{a_n}{a_{n+1}}$ 的演化：
\[
  \binom{a_{n+1}}{a_{n+2}}
  =\begin{pmatrix}0&1\\ q&p\end{pmatrix}\binom{a_n}{a_{n+1}}
  =A\,\mathbf v_n .
\]
于是 $\mathbf v_n=A^{n-1}\mathbf v_1$：\key{递推的每一小步，都是矩阵 $A$ 的一次乘法}。
矩阵 $A$ 的特征方程为
\[
  \det(\lambda I-A)=\begin{vmatrix}\lambda&-1\\-q&\lambda-p\end{vmatrix}
  =\lambda^2-p\lambda-q=0 ,
\]
与递推的特征方程完全一致——这正是「递推特征方程」的矩阵来源。

\key{（2）特征值决定通项结构。}
设特征根为 $\lambda_1,\lambda_2$。
\begin{itemize}
  \item $\lambda_1\ne\lambda_2$：$a_n=c_1\lambda_1^{\,n-1}+c_2\lambda_2^{\,n-1}$；
  \item $\lambda_1=\lambda_2=\lambda$（重根）：$a_n=(c_1+c_2n)\lambda^{\,n-1}$。
\end{itemize}
原理：特征向量方向在 $A$ 作用下只伸缩不变，故沿特征方向的「坐标」按 $\lambda^n$ 增长；
两个特征方向张成状态空间，通项就是它们的线性组合。重根时出现 $n\lambda^{n-1}$ 项（类比 Jordan 块）。

\key{（3）特征方向即不动点方向（射影视角）。}
对分式递推（莫比乌斯变换）
\[
  x_{n+1}=\frac{ax_n+b}{cx_n+d}\quad(ad-bc\ne0),
\]
引入齐次坐标 $x_n=\dfrac{u_n}{v_n}$，则
\[
  \binom{u_{n+1}}{v_{n+1}}=\begin{pmatrix}a&b\\c&d\end{pmatrix}\binom{u_n}{v_n} .
\]
$x_n$ 不变（$x_{n+1}=x_n=x^*$）当且仅当 $\dbinom{x^*}{1}$ 是矩阵的特征向量：
\[
  \begin{pmatrix}a&b\\c&d\end{pmatrix}\binom{x^*}{1}=\lambda\binom{x^*}{1}
  \iff ax^*+b=\lambda x^*,\ \ cx^*+d=\lambda .
\]
消去 $\lambda$ 即得分式不动点方程 $x^*=\dfrac{ax^*+b}{cx^*+d}$，且 $\lambda=cx^*+d$。
若有两个相异特征值 $\lambda_1=cx_1+d,\lambda_2=cx_2+d$，则双不动点构造的公比正是 $\dfrac{\lambda_2}{\lambda_1}$（见探究五）。
\key{结论：矩阵的特征方向（射影点）就是分式递推的不动点；特征值之比就是等比构造的公比。}

\begin{figure}[htbp]
  \centering
  \begin{tikzpicture}[>=Stealth, semithick, font=\footnotesize]
    \draw[->, gray] (-4.6,0)--(4.6,0) node[below right] {$u$};
    \draw[->, gray] (0,-1.5)--(0,3.3) node[above left] {$v$};
    \draw[third, very thick] (-4.4,1)--(4.4,1) node[right] {$v=1$（投影数轴）};
    % 两条特征方向直线都收在 $|y|\le 2.3$ 内，避免穿进底部说明
    \draw[main!70!black, thick, dashed] (-2.0,-2.0)--(3.6,3.6) node[above right] {$v=u\ (\lambda_1=5)$};
    \draw[second!70!black, thick, dashed] (3.9,-1.3)--(-3.9,1.3) node[above left] {$v=-\frac u3\ (\lambda_2=1)$};
    \fill[main] (1,1) circle (1.8pt) node[above right] {$x_1=1$};
    \fill[second] (-3,1) circle (1.8pt) node[below left] {$x_2=-3$};
    \fill (0,0) circle (1pt) node[below left] {$O$};
    \node[gray!70!black] at (0,-2.75) {$M=\begin{pmatrix}2&3\\1&4\end{pmatrix}$ 的两条特征方向直线与投影线 $v=1$ 的交点};
  \end{tikzpicture}
  \caption{射影同构：矩阵 $M=\begin{pmatrix}2&3\\1&4\end{pmatrix}$ 的特征方向 $v=u$、$v=-\frac u3$ 与投影线 $v=1$ 的交点，正是分式递推 $x_{n+1}=\frac{2x_n+3}{x_n+4}$ 的两个不动点 $x_1=1,\ x_2=-3$；公比 $=\lambda_2/\lambda_1=1/5$}
  \label{fig:eigen-fixed}
\end{figure}

\begin{example}
求递推 $a_{n+2}=3a_{n+1}-2a_n$ 在 $a_1=1,\ a_2=3$ 下的通项。
\end{example}
\begin{solution}
这里 $p=3,\ q=-2$，转移矩阵
\[
  A=\begin{pmatrix}0&1\\-2&3\end{pmatrix},\qquad
  \det(\lambda I-A)=\lambda^2-3\lambda+2=0
  \Rightarrow \lambda_1=1,\ \lambda_2=2 .
\]
故设 $a_n=c_1\cdot1^{n-1}+c_2\cdot2^{n-1}$。由初值：
\[
  c_1+c_2=1,\qquad c_1+2c_2=3
  \ \Rightarrow\ c_2=2,\ c_1=-1 .
\]
所以
\[
  a_n=2^n-1 .
\]
检验：$a_3=3a_2-2a_1=9-2=7=2^3-1$。$\surd$
\end{solution}

\begin{conccard}[二级]{特征值法：通项的终极统一模型}{二阶（及高阶）线性常系数递推、斐波那契类递推、重根情形}{特征方程写错符号（应为 $\lambda^2-p\lambda-q=0$）；重根时仍套用两个相异根的公式；初值代入时指数起点 $n-1$ 与 $n$ 混用}
状态向量 $\mathbf v_n=\binom{a_n}{a_{n+1}}$，$\mathbf v_{n+1}=A\mathbf v_n$，$A=\begin{pmatrix}0&1\\q&p\end{pmatrix}$；
特征方程 $\lambda^2-p\lambda-q=0$；
$\lambda_1\ne\lambda_2$ 时 $a_n=c_1\lambda_1^{n-1}+c_2\lambda_2^{n-1}$；
$\lambda_1=\lambda_2=\lambda$ 时 $a_n=(c_1+c_2n)\lambda^{n-1}$。
\end{conccard}

\subsection{Abel 变换：错位相减的机械化}

\key{探究四：离散版分部积分}

\key{（1）Abel 恒等式。}
设 $B_k=\sum_{i=1}^kb_i$，$B_0=0$。则
\[
  \sum_{k=1}^na_kb_k=a_nB_n-\sum_{k=1}^{n-1}(a_{k+1}-a_k)B_k .
\]
推导如下：
\[
\begin{aligned}
  \sum_{k=1}^na_kb_k
  &=\sum_{k=1}^na_k(B_k-B_{k-1})
  =\sum_{k=1}^na_kB_k-\sum_{k=1}^na_kB_{k-1}\\
  &=a_nB_n+\sum_{k=1}^{n-1}a_kB_k-\sum_{k=2}^{n}a_kB_{k-1}
  =a_nB_n-\sum_{k=1}^{n-1}(a_{k+1}-a_k)B_k .
\end{aligned}
\]

\key{（2）与分部积分互译。}
连续版 $\int u\,dv=uv-\int v\,du$；离散版把 $u\to a_k,\ dv\to b_k,\ v\to B_k,\ du\to \Delta a_k$：
\[
\begin{array}{c|c}
  \text{分部积分} & \text{Abel 变换} \\ \hline
  uv & a_nB_n \\
  \displaystyle\int v\,du & \displaystyle\sum_{k=1}^{n-1}(a_{k+1}-a_k)B_k
\end{array}
\]
记忆：\key{「乘积求和 $=$ 末端乘积 $-$ 差分乘部分和的求和」}。

\key{（3）一次秒求「等差 $\times$ 等比」。}
当 $a_k$ 等差时，$\Delta a_k=a_{k+1}-a_k=d$ 是常数；当 $b_k$ 等比时，$B_k$ 有现成公式。代入即得，无需「错位」。

\begin{example}
求 $S_n=\sum_{k=1}^n(2k-1)\cdot2^k$。
\end{example}
\begin{solution}
取 $a_k=2k-1$（等差，$\Delta a_k=2$），$b_k=2^k$（等比，$B_k=2^{k+1}-2$）。由 Abel：
\[
\begin{aligned}
  S_n&=a_nB_n-\sum_{k=1}^{n-1}2\cdot B_k\\
  &=(2n-1)(2^{n+1}-2)-2\sum_{k=1}^{n-1}(2^{k+1}-2)\\
  &=(2n-1)(2^{n+1}-2)-2\bigl(2^{n+1}-4-2(n-1)\bigr)\\
  &=(2n-3)2^{n+1}+6 .
\end{aligned}
\]
检验：$n=1$ 时 $S_1=1\cdot2=2$，公式 $(2-3)\cdot4+6=2$ $\surd$；
$n=2$ 时 $S_2=2+12=14$，公式 $1\cdot8+6=14$ $\surd$。
\end{solution}

\begin{figure}[htbp]
  \centering
  \begin{tikzpicture}[>=Stealth, semithick, font=\footnotesize, xscale=0.55, yscale=0.75]
    % 坐标轴
    \draw[->, gray] (-0.3,0)--(16.2,0) node[below right] {$B_k$};
    \draw[->, gray] (0,-0.3)--(0,5.8) node[above left] {$a_k$};
    % 阶梯柱：$\sum a_k b_k$
    \fill[main!10] (0,0) rectangle (1,5);
    \fill[main!10] (1,0) rectangle (3,4);
    \fill[main!10] (3,0) rectangle (6,3);
    \fill[main!10] (6,0) rectangle (10,2);
    \fill[main!10] (10,0) rectangle (15,1);
    \draw[main!60!black, thick] (0,0) rectangle (1,5);
    \draw[main!60!black, thick] (1,0) rectangle (3,4);
    \draw[main!60!black, thick] (3,0) rectangle (6,3);
    \draw[main!60!black, thick] (6,0) rectangle (10,2);
    \draw[main!60!black, thick] (10,0) rectangle (15,1);
    % 大矩形 $a_nB_n$
    \draw[second!70!black, thick, dashed] (0,0) rectangle (15,1);
    \node[second!70!black] at (7.5,0.5) {$a_nB_n$};
    % 橙色差值区域
    \fill[second!14] (0,1) rectangle (1,5);
    \fill[second!14] (1,1) rectangle (3,4);
    \fill[second!14] (3,1) rectangle (6,3);
    \fill[second!14] (6,1) rectangle (10,2);
    % 标注
    \foreach \x/\lab in {0/$B_0$,1/$B_1$,3/$B_2$,6/$B_3$,10/$B_4$,15/$B_5$}
      \draw (\x,0.08)--(\x,-0.08) node[below] {\lab};
    \node[main!70!black] at (0.5,2.5) {$a_1b_1$};
    \node[main!70!black] at (2,2.5) {$a_2b_2$};
    \node[main!70!black] at (4.5,2.5) {$a_3b_3$};
    \node[main!70!black] at (8,2.5) {$a_4b_4$};
    \node[main!70!black] at (12.5,2.5) {$a_5b_5$};
    % 差值区域标注放在全图上方的空白带（阶梯最高点为 $y=5$）
    \node[second!70!black] at (7.5,5.45) {差值区域：面积 $=\sum\limits_{k=1}^{4}(a_k-a_5)b_k=\sum\limits_{k=1}^{4}B_k=20$};
  \end{tikzpicture}
  \caption{Abel 变换面积示意（取 $a_k=6-k$，$b_k=k$，则 $B_k=\frac{k(k+1)}2$）：蓝色阶梯面积为 $\sum_{k=1}^5a_kb_k=35$；虚线大矩形面积为 $a_5B_5=15$；橙色差值区域面积为 $\sum_{k=1}^{4}B_k=1+3+6+10=20$。三者关系即 $\sum a_kb_k=a_nB_n+\sum_{k=1}^{n-1}(a_k-a_{k+1})B_k$}
  \label{fig:abel-area}
\end{figure}

\begin{conccard}[一级]{Abel 变换：离散分部求和}{等差 $\times$ 等比求和、乘积型数列求和、数列不等式放缩}{把 $a_k$ 与 $b_k$ 角色用反；忘记 $B_0=0$；$n=1$ 时不检验公式退化}
设 $B_k=\sum_{i=1}^kb_i$，则
$\sum_{k=1}^na_kb_k=a_nB_n-\sum_{k=1}^{n-1}(a_{k+1}-a_k)B_k$；
它是 $\int u\,dv=uv-\int v\,du$ 的离散版本；
取 $a_k$ 为等差、$b_k$ 为等比，可三行写出求和。
\end{conccard}

\begin{commonerror}{错位相减与 Abel 变换的高频失误}
\paragraph{失误一：错位相减时项数对不齐。}
「$S_n$ 与 $qS_n$」相减时，错位后首尾会各多出一项，必须明确写出剩余的首项与末项再合并。
\paragraph{失误二：忘记 $B_0=0$ 与 $n=1$ 的检验。}
Abel 公式中 $B_0=0$；当 $n=1$ 时 $\sum_{k=1}^{n-1}$ 为空，公式退化为 $a_1B_1$，要单独验证。
\paragraph{失误三：把 $a_k$ 与 $b_k$ 的角色用反。}
公式中 $a_k$ 是被差分的那个（应选等差／简单数列），$b_k$ 是被求和的那个（应选等比／可求和的数列）。
若反过来，$\Delta a_k$ 不再简单，公式失效。
\end{commonerror}

\subsection{共轭变换：不动点为何生效}

\key{探究五：换坐标系，把复杂运动还原为缩放}

\key{（1）不是「假设相邻两项相等」，而是「换原点」。}
「令 $a_{n+1}=a_n=x^*$」只是记忆口诀。真实动作是：先解出平衡点 $x^*$，再令
\[
  b_n=a_n-x^* .
\]
则
\[
  b_{n+1}=a_{n+1}-x^*=f(a_n)-f(x^*)\approx f'(x^*)(a_n-x^*)=f'(x^*)b_n .
\]
对线性 $f(x)=px+q$，这是精确的：$b_{n+1}=pb_n$。
几何上相当于把数轴原点从 $0$ 平移到 $x^*$：原来的「拉伸 $p$ 倍 $+$ 平移 $q$」复合运动，
在新坐标系下变成纯粹的「以 $x^*$ 为中心的缩放」。这就是\key{共轭变换} $g=h^{-1}\circ f\circ h$（其中 $h(x)=x-x^*$）。

\key{（2）一般情形：$\lambda=f'(x^*)$ 是线性化乘子。}
即使 $f$ 非线性，在不动点附近也有
\[
  b_{n+1}=f(x^*+b_n)-x^*=f'(x^*)b_n+O(b_n^2),
\]
所以 $\lambda=f'(x^*)$ 控制偏差的衰减／放大——这正是探究二稳定性判据的来源。
\key{不动点法生效的底层原因，就是「减去不动点」把非线性递推线性化为等比。}

\key{（3）莫比乌斯分式：配 $\dfrac{a_n-p}{a_n-q}$。}
对 $a_{n+1}=\dfrac{aa_n+b}{ca_n+d}$，设两个相异不动点 $x_1,x_2$。计算偏差：
\[
  a_{n+1}-x_1=\frac{(ad-bc)(a_n-x_1)}{(ca_n+d)(cx_1+d)},\qquad
  a_{n+1}-x_2=\frac{(ad-bc)(a_n-x_2)}{(ca_n+d)(cx_2+d)} .
\]
两式相除，随 $n$ 变化的分母 $(ca_n+d)$ 恰好约去：
\[
  \frac{a_{n+1}-x_1}{a_{n+1}-x_2}
  =\frac{cx_2+d}{cx_1+d}\cdot\frac{a_n-x_1}{a_n-x_2} .
\]
于是 $\left\{\dfrac{a_n-x_1}{a_n-x_2}\right\}$ 是等比数列，公比为 $\dfrac{cx_2+d}{cx_1+d}$。
若只有一个重不动点 $x_0$，则取倒数：
\[
  \frac1{a_{n+1}-x_0}-\frac1{a_n-x_0}=\text{常数},
\]
化为等差数列。

\begin{example}
设 $a_1=2,\ a_{n+1}=\dfrac{2a_n+3}{a_n+4}$，求通项公式。
\end{example}
\begin{solution}
不动点方程 $x=\dfrac{2x+3}{x+4}\Rightarrow x^2+2x-3=0\Rightarrow x_1=1,\ x_2=-3$。
\[
  \frac{a_{n+1}-1}{a_{n+1}+3}
  =\frac{\dfrac{2a_n+3}{a_n+4}-1}{\dfrac{2a_n+3}{a_n+4}+3}
  =\frac{a_n-1}{5a_n+15}
  =\frac15\cdot\frac{a_n-1}{a_n+3} .
\]
故 $\left\{\dfrac{a_n-1}{a_n+3}\right\}$ 是首项 $\dfrac{2-1}{2+3}=\dfrac15$、公比 $\dfrac15$ 的等比数列：
\[
  \frac{a_n-1}{a_n+3}=\left(\frac15\right)^n
  \ \Rightarrow\
  a_n=\frac{5^n+3}{5^n-1} .
\]
检验：$n=1$ 时 $\dfrac{5+3}{5-1}=2$ $\surd$。
\end{solution}

\begin{conccard}[二级]{共轭变换：不动点为何生效}{一阶线性递推、分式（莫比乌斯）递推、嵌套迭代 $f^{(n)}(x^*)$}{把口诀「令 $a_{n+1}=a_n=x^*$」当成真的假设相邻项相等；双不动点时忘记检验 $x_1\ne x_2$；重不动点时仍强行作比}
一阶线性：$b_n=a_n-x^*$，$b_{n+1}=f'(x^*)b_n$（精确）；
一般情形：$b_{n+1}=f'(x^*)b_n+O(b_n^2)$，$\lambda=f'(x^*)$ 是线性化乘子；
莫比乌斯双不动点：$\dfrac{a_n-x_1}{a_n-x_2}$ 等比，公比 $\dfrac{cx_2+d}{cx_1+d}$；
重不动点：$\dfrac1{a_n-x_0}$ 等差。
\end{conccard}

\subsection{强基拓展：生成函数（超出高考考纲）}

\key{探究六：把整个数列打包成一个函数}

\key{温馨提示：本模块涉及形式幂级数，超出高考考纲，仅面向竞赛／强基与学有余力者。}

\key{（1）生成函数：序列 $\to$ 幂级数。}
定义
\[
  G(x)=\sum_{n\ge0}a_nx^n .
\]
它把无穷数列 $\{a_n\}$ 打包成一个形式幂级数：
\begin{itemize}
  \item 下标平移 $a_{n+1}$ 对应「除以 $x$ 并减去首项」：$\sum_{n\ge0}a_{n+1}x^n=\dfrac{G(x)-a_0}{x}$；
  \item 递推关系对应 $G(x)$ 的代数方程；
  \item 解出 $G(x)$ 再展开，$x^n$ 的系数就是 $a_n$。
\end{itemize}

\key{（2）例：斐波那契数列。}
设 $F_0=0,\ F_1=1,\ F_{n+2}=F_{n+1}+F_n$，$G(x)=\sum_{n\ge0}F_nx^n$。则
\[
\begin{aligned}
  G(x)&=x+\sum_{n\ge2}F_nx^n
  =x+\sum_{n\ge2}(F_{n-1}+F_{n-2})x^n\\
  &=x+x\sum_{n\ge2}F_{n-1}x^{n-1}+x^2\sum_{n\ge2}F_{n-2}x^{n-2}
  =x+xG(x)+x^2G(x) .
\end{aligned}
\]
故
\[
  G(x)=\frac{x}{1-x-x^2}
  =\frac{x}{(1-\varphi x)(1-\psi x)},\qquad
  \varphi=\frac{1+\sqrt5}{2},\ \psi=\frac{1-\sqrt5}{2} .
\]
部分分式分解：
\[
  G(x)=\frac1{\sqrt5}\left(\frac1{1-\varphi x}-\frac1{1-\psi x}\right)
  =\frac1{\sqrt5}\sum_{n\ge0}(\varphi^n-\psi^n)x^n ,
\]
比较系数即得比内公式
\[
  F_n=\frac{\varphi^n-\psi^n}{\sqrt5} .
\]
同一个 $G(x)$ 还能给出组合意义（第 7 章）：$\dfrac1{1-x-x^2}$ 的系数是「用 $1\times1$ 与 $1\times2$ 骨牌铺满 $1\times n$ 棋盘」的方法数。

\key{（3）适用边界。}
生成函数擅长：线性递推（齐次／非齐次）、卷积型求和、组合计数（卡特兰数等）。
它把「解递推」变成「解代数方程 $+$ 展开」，是强基与竞赛的核武器；
但解答题若用此法，需补写合法的代数推导。

\begin{conccard}[三级]{生成函数：递推 $\to$ 代数方程}{斐波那契等线性递推、非齐次递推、卷积与组合计数}{超出高考考纲，解答题不能直接套用；忘记首项 $a_0$ 导致平移错误；展开时忽略收敛域（形式幂级数只需形式运算）}
$G(x)=\sum_{n\ge0}a_nx^n$；$\sum_{n\ge0}a_{n+1}x^n=\dfrac{G(x)-a_0}{x}$；
由递推导出 $G(x)$ 的封闭形式，再展开读系数；
斐波那契：$G(x)=\dfrac{x}{1-x-x^2}$，$F_n=\dfrac{\varphi^n-\psi^n}{\sqrt5}$。
\end{conccard}

\begin{method}[变换—求解—逆变换：本章五件降维武器]
\begin{enumerate}
  \item \key{换坐标系}：看到 $a_{n+1}=f(a_n)$，先解不动点 $x^*=f(x^*)$，作 $b_n=a_n-x^*$；
  \item \key{升维}：看到高阶线性递推，写成状态向量 $\mathbf v_{n+1}=A\mathbf v_n$，求特征值；
  \item \key{互逆}：看到求和，想差分 $\Delta$ 与裂项；看到乘积求和，想 Abel 变换；
  \item \key{打包}：递推复杂时，用生成函数 $G(x)$ 把它变成代数方程（强基）；
  \item \key{还原}：在新坐标系算完后，用逆变换 $h^{-1}$ 把答案翻回原数列，并检验首项与边界。
\end{enumerate}
\end{method}

\begin{summary}
本章把数列从「套公式」升级为「离散动力系统 $+$ 降维变换」：
差分 $\Delta$ 与求和 $\sum$ 互为逆运算（离散版微积分基本定理），由 $S_n$ 反求 $a_n$、差分降次判断次数；
一阶递推 $a_{n+1}=f(a_n)$ 是不动点迭代，$x^*=f(x^*)$ 是平衡态，蛛网图把迭代画成阶梯，$|f'(x^*)|<1$ 吸引、$>1$ 排斥；
二阶线性递推对应转移矩阵 $\begin{pmatrix}0&1\\q&p\end{pmatrix}$，特征方程 $\lambda^2-p\lambda-q=0$ 决定通项结构，特征方向即分式递推的不动点方向；
Abel 变换是离散分部积分，把「等差 $\times$ 等比」求和机械化；
共轭变换（换坐标系）解释了不动点法为何生效：$b_n=a_n-x^*$ 把非线性递推线性化为等比；
生成函数（强基拓展）把递推打包成代数方程。

\begin{center}
\footnotesize
\setlength{\tabcolsep}{4pt}
\renewcommand{\arraystretch}{1.35}
\begin{tabular}{p{2.1cm}p{4.9cm}p{4.2cm}p{1.9cm}}
\toprule
降维变换 & 核心操作 & 把什么变成什么 & 对应探究 \\
\midrule
差分 $\Delta$ & $\Delta a_n=a_{n+1}-a_n$ & 求和 $\to$ 相消；高次 $\to$ 低次 & 探究一 \\
共轭（平移／分式） & $b_n=a_n-x^*$；$\dfrac{a_n-x_1}{a_n-x_2}$ & 拉伸$+$平移／非线性分式 $\to$ 等比／等差 & 探究二、五 \\
特征值对角化 & $\mathbf v_{n+1}=A\mathbf v_n$ & 递推 $\to$ 特征值幂 & 探究三 \\
Abel 变换 & $\sum a_kb_k=a_nB_n-\sum\Delta a_kB_k$ & 乘积求和 $\to$ 差分 $\times$ 部分和 & 探究四 \\
生成函数 & $G(x)=\sum a_nx^n$ & 递推 $\to$ 代数方程 & 探究六 \\
\bottomrule
\end{tabular}
\end{center}

\key{数列不是公式的仓库，而是离散动力系统的轨迹；解题的本质是换坐标系，把陌生递推化为等比／等差／代数方程。}
\end{summary}

\end{explore}

% ---------- 跨章联动：降维双引擎（正文收尾；原在答案区之后，已前移） ----------
\subsection{跨章方法论大一统：平移消常数，旋转消交叉项}

前面第 4 章用\key{仿射}压缩椭圆回圆、第 8 章用\key{齐次化}处理条件最值，
本章用\key{不动点}处理递推。它们看似分属三个板块，底层却是同一个动作：
\key{当前的坐标原点没选好，方程里才拖着讨厌的常数尾巴}。把原点精准搬到
锚点上，非齐次项会瞬间蒸发。

\begin{conccard}[二级]{换坐标系的两步走}{适用于所有「带常数项」「带一次项」「带交叉项」的结构}{只背结论，不问原点该搬去哪}
第一步\key{平移}（$x'=x-x_0,\ y'=y-y_0$）：把原点搬到不动点、对称中心或定点 $P$，利用「锚点自身满足方程」这一性质抵消常数项与一次项（第 8 章「手术刀之二」）；第二步\key{旋转}（乘以 $e^{i\theta}$）：把坐标轴对齐到主轴／特征方向，消灭二次交叉项 $Bxy$，并把复特征根的螺旋迭代还原为实数周期震荡。
\end{conccard}

\begin{center}
\renewcommand{\arraystretch}{1.4}
\begin{tabular}{|p{2.8cm}|p{5.0cm}|p{5.4cm}|}
\hline
\textbf{对比维度} & \textbf{数列：不动点法} & \textbf{条件最值／圆锥曲线：齐次化（第 8 章）} \\
\hline
非齐次障碍 & 递推拖着常数尾巴：$a_{n+1}=pa_n+q$ & 平移到定点后仍剩一次项与常数项 \\
\hline
锚点满足的方程 & $x^*=px^*+q$ & $\dfrac{x_0^2}{a^2}+\dfrac{y_0^2}{b^2}=1$ \\
\hline
平移后发生了什么 & 令 $b_n=a_n-x^*$，得 $b_{n+1}=pb_n$，$q$ 归零 & 常数项化为 $\dfrac{x_0^2}{a^2}+\dfrac{y_0^2}{b^2}-1=0$，自动抵消 \\
\hline
剩下的纯粹结构 & 只含公比 $p$ 的等比缩放 & 只含斜率 $k=\dfrac{y'}{x'}$ 的一元二次方程 \\
\hline
\end{tabular}
\end{center}

\begin{problem}[旋转 $45^\circ$：把 $xy=1$ 还原成标准双曲线]
证明反比例函数 $xy=1$ 的图象是等轴双曲线，并写出其焦点位置。
\end{problem}
\begin{solution}
将坐标轴逆时针旋转 $45^\circ$，即
\[
  x=\frac{\sqrt2}{2}(x'-y'),\qquad y=\frac{\sqrt2}{2}(x'+y') .
\]
代入 $xy=1$：
\[
  \frac{\sqrt2}{2}(x'-y')\cdot\frac{\sqrt2}{2}(x'+y')=1
  \ \Longrightarrow\
  \frac{x'^2}{2}-\frac{y'^2}{2}=1 .
\]
交叉项 $xy$ 变成了平方差：这正是以 $x'$ 轴为实轴的等轴双曲线，半实轴 $a=b=\sqrt2$，故 $c=\sqrt{a^2+b^2}=2$，焦点在新坐标系的 $(\pm2,0)$，也就是原坐标系的 $(\sqrt2,\sqrt2)$ 与 $(-\sqrt2,-\sqrt2)$。
\end{solution}

\begin{center}
  \begin{tikzpicture}[x=1.05cm,y=1.05cm,>=stealth,font=\small]
    \draw[->,white!55] (-3,0)--(3,0) node[right,white!70] {$x$};
    \draw[->,white!55] (0,-3)--(0,3) node[above,white!70] {$y$};
    % 旋转后的主轴
    \draw[->,second,dashed,very thick] (-2.3,-2.3)--(2.5,2.5) node[above right,second] {\small $x'$（新实轴）};
    \draw[->,second,dashed,very thick] (2.3,-2.3)--(-2.5,2.5) node[above left,second] {\small $y'$（新虚轴）};
    \draw[->,second,thick] (1.1,0) arc (0:45:1.1);
    \node[second] at (1.15,0.38) {\small $45^\circ$};
    % 双曲线 xy=1 的两支
    \draw[main,very thick,domain=0.42:2.6,samples=80,smooth] plot(\x,{1/\x});
    \draw[main,very thick,domain=-2.6:-0.42,samples=80,smooth] plot(\x,{1/\x});
    \node[main] at (1.9,1.15) {\small $xy=1$};
    \node[white!70] at (0,-3.4) {\small 交叉项 $xy$ 与旋转 $45^\circ$ 后的新主轴 $x',y'$};
  \end{tikzpicture}\\[0.3em]
  {\small\itshape 同一曲线的两种读法：原坐标系下是交叉项方程 $xy=1$；把轴转 $45^\circ$ 后成为标准等轴双曲线 $\frac{x'^2}{2}-\frac{y'^2}{2}=1$}
\end{center}

\paragraph{第二步：旋转也管数列。} 当二阶递推 $a_{n+2}=pa_{n+1}+qa_n$ 的特征方程判别式 $\Delta<0$ 时，两根为共轭复数 $\lambda_{1,2}=r(\cos\theta\pm i\sin\theta)$。此时几何图像不再是沿直线趋近不动点，而是每步\key{缩放 $r$ 倍并旋转 $\theta$ 角}的对数螺线；通项写成
\[
  a_n=r^{\,n-1}\bigl(c_1\cos n\theta+c_2\sin n\theta\bigr) .
\]
$r=1$ 且 $\theta=\frac{2\pi}{T}$ 时数列退化为纯周期，例如 $a_{n+2}=a_{n+1}-a_n$ 的 $\lambda=e^{\pm i\frac{\pi}{3}}$，周期为 $6$。这与「把轴旋转到主轴方向」是同一件事：都是把耦合项转没。

\begin{figure}[htbp]
  \centering
  \begin{tikzpicture}[>=Stealth, semithick, font=\footnotesize, x=1.15cm, y=1.15cm]
    \draw[->, gray] (-1.9,0)--(1.95,0);
    \draw[->, gray] (0,-1.35)--(0,1.55);
    % 迭代点 v0..v5（M = r R(40°)，r=0.8），相邻两点长度比 0.8、辐角差 40°
    \coordinate (v0) at (1.5,0.2);
    \coordinate (v1) at (0.8163,0.8942);
    \coordinate (v2) at (0.0402,0.9679);
    \coordinate (v3) at (-0.4733,0.6138);
    \coordinate (v4) at (-0.6058,0.1327);
    \coordinate (v5) at (-0.4394,-0.2302);
    \fill (0,0) circle (1.2pt) node[below right] {$O$};
    \draw[main, very thick] (v0)--(v1)--(v2)--(v3)--(v4)--(v5);
    \foreach \i/\p/\pos in {0/v0/above right,1/v1/right,2/v2/above left,3/v3/left,4/v4/below left,5/v5/below}
      \fill[second] (\p) circle (1.5pt) node[\pos] {$\mathbf v_{\i}$};
    \node[gray!70!black, align=center] at (0,-1.15) {$M=rR(\theta)$，$r=0.8,\ \theta=40^\circ$：每步长度乘 $r$、方向加 $\theta$；\\$r<1$ 螺线收拢到唯一不动点 $O$，$r>1$ 远离，$r=1$ 且 $\theta=2\pi/T$ 时为周期 $T$};
  \end{tikzpicture}
  \caption{复特征值的迭代图像：$v_{n+1}=Mv_n$ 且 $M=rR(\theta)$ 时，每一步都把向量缩放 $r$ 倍并旋转 $\theta$ 角，点列沿螺线运动——这就是「实部控伸缩、虚部控旋转」的离散版本}
  \label{fig:spiral-complex}
\end{figure}

\begin{commonerror}{卷面上禁止出现「由仿射变换知」「由极点极线知」}
高观点只用于草稿纸预判答案与取等位置，答题卡上必须写考纲认可的初等推导（联立、韦达定理、作差构造等比）。详见附录~\ref{app:legal}。
\end{commonerror}

\begin{inclasstest}
\begin{enumerate}
  \item 设 $S_n=2n^2-n$，求数列 $\{a_n\}$ 的通项公式；\level{易}
  \item 设 $a_1=1,\ a_{n+1}=2a_n+3$，求通项并判断 $\{a_n\}$ 是否收敛；\level{中}
  \item 用 Abel 变换求 $\sum_{k=1}^nk\cdot2^k$；\level{中}
  \item 求 $a_{n+2}=4a_{n+1}-4a_n$ 在 $a_1=1,\ a_2=4$ 下的通项（重根情形）。\level{中}
\end{enumerate}
\end{inclasstest}

\begin{afterclass}
\subsection*{基础巩固}
\begin{enumerate}
  \item 已知 $\Delta a_n=2n+1$，且 $a_1=1$，求 $a_n$；
  \item 证明 $\sum_{k=1}^n\Delta a_k=a_{n+1}-a_1$，并用它求 $\sum_{k=1}^n\dfrac1{k(k+1)}$。
\end{enumerate}

\subsection*{进阶提升}
\begin{enumerate}
  \item 用特征值法求 $a_{n+2}=2a_{n+1}+3a_n$ 在 $a_1=1,\ a_2=5$ 下的通项；
  \item 设 $a_1=0,\ a_{n+1}=\dfrac{3a_n+1}{a_n+3}$，求通项公式。
\end{enumerate}

\subsection*{拓展探究}
\begin{enumerate}
  \item （强基／竞赛，超出高考考纲）设 $a_0=1,\ a_{n+1}=2a_n+1$，用生成函数求通项公式。
\end{enumerate}

\subsection*{教师版答案}
\begin{solution}
\textbf{基础 1：}由离散版微积分基本定理，
\[
  a_n=a_1+\sum_{k=1}^{n-1}\Delta a_k
  =1+\sum_{k=1}^{n-1}(2k+1)
  =1+2\cdot\frac{(n-1)n}2+(n-1)=n^2 .
\]
（检验：$n=1$ 时 $a_1=1$。）
\end{solution}
\begin{solution}
\textbf{基础 2：}求和时中间项全部抵消：
\[
  \sum_{k=1}^n\Delta a_k=\sum_{k=1}^n(a_{k+1}-a_k)=a_{n+1}-a_1 .
\]
又 $\dfrac1{k(k+1)}=\dfrac1k-\dfrac1{k+1}=\Delta\!\left(-\dfrac1k\right)$，故
\[
  \sum_{k=1}^n\frac1{k(k+1)}
  =\sum_{k=1}^n\Delta\!\left(-\frac1k\right)
  =\left(-\frac1{n+1}\right)-(-1)=\frac n{n+1} .
\]
\end{solution}
\begin{solution}
\textbf{进阶 1：}$A=\begin{pmatrix}0&1\\3&2\end{pmatrix}$，特征方程
\[
  \lambda^2-2\lambda-3=0\Rightarrow\lambda_1=3,\ \lambda_2=-1 .
\]
设 $a_n=c_1\cdot3^{n-1}+c_2\cdot(-1)^{n-1}$，由 $a_1=1,\ a_2=5$：
\[
  c_1+c_2=1,\quad 3c_1-c_2=5
  \ \Rightarrow\ c_1=\frac32,\ c_2=-\frac12 .
\]
故
\[
  a_n=\frac{3^n+(-1)^n}2 .
\]
（检验：$a_3=2\cdot5+3\cdot1=13=\frac{27-1}2$。）
\end{solution}
\begin{solution}
\textbf{进阶 2：}不动点方程 $x=\dfrac{3x+1}{x+3}\Rightarrow x^2=1\Rightarrow x_1=1,\ x_2=-1$。
\[
  \frac{a_{n+1}-1}{a_{n+1}+1}
  =\frac{(3a_n+1)-(a_n+3)}{(3a_n+1)+(a_n+3)}
  =\frac{2a_n-2}{4a_n+4}
  =\frac12\cdot\frac{a_n-1}{a_n+1} .
\]
首项 $\dfrac{0-1}{0+1}=-1$，故 $\dfrac{a_n-1}{a_n+1}=-\left(\dfrac12\right)^{n-1}$，解得
\[
  a_n=\frac{1-\left(\frac12\right)^{n-1}}{1+\left(\frac12\right)^{n-1}}
  =\frac{2^{n-1}-1}{2^{n-1}+1} .
\]
（检验：$n=1$ 时 $a_1=0$；$n=2$ 时 $a_2=\frac13=\frac{3\cdot0+1}{0+3}$。）
\end{solution}
\begin{solution}
\textbf{拓展：}设 $G(x)=\sum_{n\ge0}a_nx^n$。由 $a_{n+1}=2a_n+1$：
\[
  \frac{G(x)-a_0}{x}=2G(x)+\sum_{n\ge0}x^n=2G(x)+\frac1{1-x} .
\]
代入 $a_0=1$：
\[
  G(x)-1=2xG(x)+\frac{x}{1-x}
  \ \Rightarrow\
  G(x)(1-2x)=\frac1{1-x}
  \ \Rightarrow\
  G(x)=\frac1{(1-x)(1-2x)} .
\]
部分分式：
\[
  G(x)=\frac2{1-2x}-\frac1{1-x}
  =\sum_{n\ge0}(2\cdot2^n-1)x^n
  =\sum_{n\ge0}(2^{n+1}-1)x^n ,
\]
故 $a_n=2^{n+1}-1$。（检验：$a_0=1,\ a_1=3$。）
\end{solution}

\reading{拓展阅读 11}{数列：离散动力系统的轨迹}
\end{afterclass}

% 高观点导学案：导数与局部线性化引擎（第 6 章·高观点版）
% 主线：局部线性化（以直代曲）→ 导数即变化率（差分的连续版）→ 五大不等式的结构本质（支撑边界）
%       → 切线放缩与隐零点 → 条件最值四大标准动作 → 含参恒成立/能成立的集合边界
% 核心思想：把「导数」与「不等式」统一在「以直代曲」这一个引擎下——曲线在一点附近用切线代替，
%           不等式是为点集划定不可逾越的边界；解题四步：识别结构 → 局部线性化 → 求最值 → 验取等。

\section{导数与局部线性化引擎（高观点版）}

\begin{methodbox}{多框架对读}
\small
\begin{center}
\begin{tabular}{c|p{2.9cm}|p{5.0cm}|p{3.1cm}}
\toprule
视角 & 本章看点 & 别的视角怎么说 & 哪侧最省力 \\
\midrule
代数 & 求导法则、恒等变形 & 符号对象的机械运算 & 硬算单调区间 \\
几何 & 切线、增减、凹凸的形状 & 曲线在放大镜下变直线 & 看图猜极值 \\
集合 & 零点集、符号区间 & 定义域值域的边界 & 界参数范围 \\
映射 & 局部线性化：曲线 $\to$ 直线 & 换到切线坐标系再比大小 & 五不等式一源 \\
\bottomrule
\end{tabular}
\end{center}
卡的使用：证不等式先想「是不是某个函数的单调性」，再把切线当替代函数（总图见第 0 章图~\ref{fig:four-view-map}）。
\end{methodbox}

第 1 章确立了函数与单调性，第 5 章把差分 $\Delta$ 确立为离散意义下的变化率；
本章把「导数」与「不等式」统一在同一个引擎下——\key{局部线性化（以直代曲）}：
任何光滑曲线在一点的无穷小邻域内都可以用它的切线近似代替，而这条切线的斜率就是导数。
于是「求导数、判符号」成为判定单调性／极值／最值的统一操作；
五大核心不等式分别是射影定理、内积公理、齐次化（完整展开见第 8 章）、度量三角公理与凸凹函数切线支撑这五种结构下的「支撑边界」；
条件最值与含参成立性问题，则归结为「边界值能否取到」。
六个模块层层展开：
\key{局部线性化}（放大镜下曲成直）$\to$
\key{导数即变化率}（差分的连续版，统一判定单调性／极值／最值）$\to$
\key{五大不等式的结构本质}（一张对比表：射影、内积、齐次、度量、凸支撑）$\to$
\key{切线放缩与隐零点}（$e^x\ge x+1$、$\ln x\le x-1$ 的严格证明与多步取等检验）$\to$
\key{条件最值四大标准动作}（1 的代换、消元主元、拆项分离、切线放缩）$\to$
\key{含参恒成立／能成立}（点集位置、集合包含与量词的值域边界）。
核心思想：\key{导数与不等式是同一个引擎的两面——以直代曲给出局部逼近，支撑边界给出整体控制；解题的本质是「识别结构 $\to$ 局部线性化 $\to$ 求最值 $\to$ 验取等」。}

\begin{learnobjective}
\begin{itemize}
  \item 理解局部线性化公式 $f(x)\approx f(x_0)+f'(x_0)(x-x_0)$ 与「以直代曲」的几何含义，会用切线作近似计算；同一个一阶逼近在第 5 章的离散递推中就是线性化乘子 $f'(x^*)$；
  \item 掌握导数的极限定义与导函数，理解它是差分的连续版本（$\Delta a_n$ 是 $a_{n+1}-a_n$，$f'(x)$ 是 $\frac{\Delta y}{\Delta x}$ 的极限）；
  \item 会用导数符号统一判定单调性、极值与最值，理解驻点与极值点的区别；
  \item 从射影定理、内积公理、齐次化、度量公理、凸凹性五个结构本质理解五大不等式，并记住各自的取等条件（齐次化专论见第 8 章）；
  \item 会用作差求导严格证明 $e^x\ge x+1$、$\ln x\le x-1$，能在多步放缩中检验各步取等条件是否同时成立；
  \item 掌握隐零点「设而不求、整体代换」的基本操作；
  \item 掌握条件最值的四大标准动作（1 的代换、消元主元、拆项分离常数、切线放缩），会检验取等的可行性；
  \item 会用点集上下位置与集合包含关系统一解释恒成立／能成立问题，掌握量词到最值的四类转化与含参二次不等式的分类讨论树。
\end{itemize}
\end{learnobjective}

\begin{preview}
\subsection*{知识回顾}
\begin{itemize}
  \item 函数 $y=f(x)$ 在区间 $[x_0,x_0+\Delta x]$ 上的平均变化率是 \blank；
  \item 导数的定义：$f'(x_0)=$ \blank（写成 $\Delta x\to0$ 时的极限）；
  \item 若 $f'(x)>0$ 在区间 $I$ 上恒成立，则 $f$ 在 $I$ 上 \blank；
  \item 基本不等式：$a,b>0$ 时 \blank $\ge$ \blank，当且仅当 \blank 时取等；
  \item 一元二次方程 $ax^2+bx+c=0$ 的判别式 $\Delta=$ \blank，$\Delta>0$ 时有两个 \blank 实根。
\end{itemize}

\subsection*{教材速读}
阅读「导数」一章教材，完成下列问题：
\begin{enumerate}
  \item 可导与连续的关系是：可导 \blank 连续，连续 \blank 可导（填「一定」或「不一定」）；
  \item 写出三个基本求导公式：$(x^n)'=$ \blank，$(e^x)'=$ \blank，$(\ln x)'=$ \blank；
  \item 函数极值的第一充分条件：在 $x_0$ 两侧 $f'(x)$ 异号，则 $x_0$ 是 \blank；
  \item 闭区间上连续函数的最值只可能在 \blank 与 \blank 处取得。
\end{enumerate}
\end{preview}

\begin{explore}

\subsection{局部线性化：放大镜下曲成直}

\key{探究一：切线是曲线在一点处的「最佳直线近似」}

\key{（1）从平均变化率到瞬时变化率。}
函数在 $x_0$ 附近的平均变化率是 $\dfrac{\Delta y}{\Delta x}=\dfrac{f(x_0+\Delta x)-f(x_0)}{\Delta x}$。
让 $\Delta x\to0$，若极限存在，就得到导数
\[
  f'(x_0)=\lim_{\Delta x\to0}\frac{f(x_0+\Delta x)-f(x_0)}{\Delta x}.
\]
这正是第 5 章差分 $\Delta a_n=a_{n+1}-a_n$ 的连续版：把「步长为 1 的差分」换成「步长趋于 0 的差商极限」。

\key{（2）局部线性化公式：以直代曲。}
过点 $P(x_0,f(x_0))$ 的切线方程为
\[
  y=f(x_0)+f'(x_0)(x-x_0),
\]
而在 $x_0$ 附近，曲线与切线的误差是关于 $(x-x_0)$ 的高阶无穷小：
\[
  f(x)=f(x_0)+f'(x_0)(x-x_0)+o(x-x_0).
\]
也就是说，\key{在 $x_0$ 的无穷小邻域内，曲线可以用切线代替——这就是「以直代曲」}。
几何显微镜放得越大，曲线看起来越像一条直线（图 \ref{fig:local-linear}）。

\begin{figure}[htbp]
  \centering
  \begin{tikzpicture}[>=Stealth, semithick, font=\footnotesize]
    % ---------- 左：全局视角 ----------
    \begin{scope}
      \draw[->, gray] (-1.3,0)--(3.0,0) node[below right] {$x$};
      \draw[->, gray] (0,-2.3)--(0,4.4) node[above left] {$y$};
      \draw[main!70!black, very thick, domain=-1:2, smooth, samples=60] plot (\x, {\x*\x});
      \node[main!70!black] at (1.72,3.15) {$y=x^2$};
      \draw[second, thick, domain=-0.4:2.5] plot (\x, {2*\x-1});
      \node[second] at (2.64,3.9) {$y=2x-1$};
      \fill (1,1) circle (1.6pt) node[above right] {$P(1,1)$};
      \draw[dashed, gray] (0.85,0.7) rectangle (1.15,1.3);
      % 说明文字整体右移，避开左图的 $y$ 轴
      \node[gray!70!black, align=left] at (0.2,-2.15) {虚线框为右图的放大区域};
    \end{scope}
    % ---------- 右：放大视角 ----------
    \begin{scope}[xshift=5.8cm]
      \draw[->, gray] (0.78,0.55)--(1.26,0.55) node[above right] {$x$};
      \draw[->, gray] (0.75,0.58)--(0.75,1.46) node[above left] {$y$};
      \draw[second, thick, dashed, domain=0.84:1.16] plot (\x, {2*\x-1});
      \draw[main!70!black, very thick, domain=0.84:1.16, smooth, samples=60] plot (\x, {\x*\x});
      \fill (1,1) circle (1.6pt) node[above right] {$P(1,1)$};
      \node[gray!70!black, align=center] at (1.01,-0.1) {放大约 $25$ 倍后，曲线（实线）\\与切线（虚线）几乎重合};
    \end{scope}
  \end{tikzpicture}
  \caption{以直代曲：抛物线 $y=x^2$ 在 $P(1,1)$ 处的切线为 $y=2x-1$。左图全局视角下曲线明显弯曲；右图把视野放大到 $P$ 附近，曲线与切线在视觉上无法区分——$f(x)\approx f(x_0)+f'(x_0)(x-x_0)$}
  \label{fig:local-linear}
\end{figure}

\begin{example}
用切线近似求 $\sqrt{4.1}$ 与 $\sin31^\circ$（精确到 $0.001$）。
\end{example}
\begin{solution}
取 $f(x)=\sqrt x$，在 $x_0=4$ 处：$f(4)=2$，$f'(x)=\frac1{2\sqrt x}$，$f'(4)=\frac14$，故
\[
  \sqrt{4.1}\approx 2+\frac14\cdot0.1=2.025 .
\]
（真值 $\sqrt{4.1}\approx2.0248$，误差约 $1.5\times10^{-4}$。）

取 $g(x)=\sin x$，在 $x_0=\frac\pi6$ 处：$g'(x)=\cos x$，$g'\!\left(\frac\pi6\right)=\frac{\sqrt3}{2}$，自变量增量
$\Delta x=\frac\pi{180}$，故
\[
  \sin31^\circ\approx \frac12+\frac{\sqrt3}{2}\cdot\frac{\pi}{180}
  \approx0.5+0.01511=0.515 .
\]
（真值 $\sin31^\circ\approx0.51504$。）
\end{solution}

\begin{conccard}[一级]{局部线性化：以直代曲}{光滑曲线的近似计算、误差估计、理解导数的几何意义}{把「近似」当成「相等」；放大倍数不够时强行替代；忘记公式只在 $x_0$ 附近有效}
$f'(x_0)=\lim\limits_{\Delta x\to0}\dfrac{f(x_0+\Delta x)-f(x_0)}{\Delta x}$；
切线 $y=f(x_0)+f'(x_0)(x-x_0)$；
$f(x)=f(x_0)+f'(x_0)(x-x_0)+o(x-x_0)$，即以直代曲；
导数是差分 $\Delta a_n$ 的连续版本。
\end{conccard}

\subsection{导数即变化率：统一判定单调性、极值、最值}

\key{探究二：$f'(x)$ 的正负号就是曲线的「上下坡」}

\key{（1）导函数与符号判定。}
把每一点的导数收集起来得到导函数 $f'(x)$。在区间 $I$ 上：
\begin{itemize}
  \item $f'(x)>0$ 恒成立 $\iff$ $f$ 在 $I$ 上单调递增；
  \item $f'(x)<0$ 恒成立 $\iff$ $f$ 在 $I$ 上单调递减。
\end{itemize}
这与第 5 章用差分判数列单调性完全同构：$\Delta a_n>0\iff a_{n+1}>a_n$（递增数列）。

\key{（2）极值：$f'$ 变号的点。}
费马引理：可导函数在极值点处必有 $f'(x_0)=0$（驻点）。
第一充分条件：$f'$ 在 $x_0$ 两侧异号 $\iff$ $x_0$ 是极值点。汇总如下：

\begin{center}\footnotesize\renewcommand{\arraystretch}{1.3}
\begin{tabular}{c|c|c}
$f'(x)$ 在 $x_0$ 附近的符号 & $f(x)$ 的形态 & 极值 \\\hline
$+$，然后 $-$ & 先增后减 & 极大值 \\
$-$，然后 $+$ & 先减后增 & 极小值 \\
不变号（如 $y=x^3$ 在 $0$ 处） & 单调 & 无极值 \\
\end{tabular}
\end{center}

\key{（3）最值：端点 $+$ 极值点。}
闭区间 $[a,b]$ 上的连续函数，最值只可能在端点 $a,b$ 与极值点处取得，比较这些点的函数值即可。

\begin{example}
求 $f(x)=x^3-3x$ 的单调区间与极值；若在 $[-2,2]$ 上求最值。
\end{example}
\begin{solution}
$f'(x)=3x^2-3=3(x+1)(x-1)$。
\begin{itemize}
  \item $x<-1$ 或 $x>1$ 时 $f'(x)>0$；$-1<x<1$ 时 $f'(x)<0$。
  \item 故 $f$ 在 $(-\infty,-1)$、$(1,+\infty)$ 上递增，在 $(-1,1)$ 上递减。
  \item 极大值 $f(-1)=-1+3=2$；极小值 $f(1)=1-3=-2$。
\end{itemize}
在 $[-2,2]$ 上额外比较端点：$f(-2)=-8+6=-2$，$f(2)=8-6=2$。
故最小值为 $-2$（在 $x=-2$ 与 $x=1$ 处），最大值为 $2$（在 $x=-1$ 与 $x=2$ 处）。
$\surd$
\end{solution}

\begin{commonerror}{驻点、极值点与最值的三个高频失误}
\paragraph{失误一：以为 $f'(x_0)=0$ 就一定是极值点。}
$f'(x_0)=0$ 只是必要条件。反例：$f(x)=x^3$ 在 $x=0$ 处 $f'(0)=0$，但两侧 $f'>0$ 不变号，不是极值点。
必须检验两侧符号是否改变。
\paragraph{失误二：忽略定义域。}
例如 $f(x)=x\ln x$ 的定义域是 $(0,+\infty)$，讨论单调性时不能把 $x\le0$ 也算进去；
求闭区间最值时也不能漏掉端点。
\paragraph{失误三：把「单点处 $f'(x_0)>0$」当成「区间上递增」。}
导数符号必须在整个区间上保持，才能推出单调性；个别点的符号说明不了整体趋势。
\end{commonerror}

\begin{conccard}[一级]{导数符号法：单调性、极值、最值一条龙}{含参数的函数性质讨论、恒成立问题中求最值}{驻点当成极值点；漏定义域；漏端点；把 $f'(x)\ge0$ 中的孤立零点误判为单调性破坏}
$f'(x)>0\Rightarrow$ 递增，$f'(x)<0\Rightarrow$ 递减；
极值点处 $f'(x_0)=0$ 且两侧异号；
闭区间最值 $=$ 端点值与极值值中的最大／最小者。
\end{conccard}

\subsection{五大不等式的结构本质}

\key{探究三：五条不等式，五个「支撑边界」}

同一个「$\ge$」在不同结构下有不同来历。看懂它们的结构本质，取等条件就会自己跳出来。

\begin{center}\footnotesize\setlength{\tabcolsep}{4pt}\renewcommand{\arraystretch}{1.4}
\begin{tabular}{p{1.9cm}p{3.5cm}p{5.2cm}p{2.6cm}}
\toprule
不等式 & 核心形式 & 结构本质 & 取等条件 \\
\midrule
基本不等式 & $\dfrac{a+b}{2}\ge\sqrt{ab}\ (a,b>0)$ & 半圆射影定理：直径 $a+b$ 的半径 $\ge$ 斜边上的高 $\sqrt{ab}$ & $a=b$ \\[2pt]
柯西不等式 & $(a_1^2+a_2^2)(b_1^2+b_2^2)\ge(a_1b_1+a_2b_2)^2$ & 向量内积公理 $|\vec u|^2|\vec v|^2\ge(\vec u\cdot\vec v)^2$（即 $|\cos\theta|\le1$） & $\vec u\parallel\vec v$（共线） \\[2pt]
权方和不等式 & $\dfrac{a_1^2}{b_1}+\dfrac{a_2^2}{b_2}\ge\dfrac{(a_1+a_2)^2}{b_1+b_2}\ (b_i>0)$ & 柯西在 $\left(\dfrac{a_i}{\sqrt{b_i}}\right)$ 与 $(\sqrt{b_i})$ 下的齐次化特例 & $\dfrac{a_1}{b_1}=\dfrac{a_2}{b_2}$ \\[2pt]
绝对值三角不等式 & $\bigl||a|-|b|\bigr|\le|a\pm b|\le|a|+|b|$ & 度量空间的距离三角不等式（三角形两边之和 $\ge$ 第三边） & 共线同向（右端）／反向（左端） \\[2pt]
切线放缩 & $e^x\ge x+1$，$\ln x\le x-1$ & 凸／凹函数的切线支撑边界（Jensen 不等式的一阶情形） & $x=0$／$x=1$ \\
\bottomrule
\end{tabular}
\end{center}

注：齐次化作为第一性原理的完整方法链见第 8 章；本章只把它用作
条件最值的一个入口，权方和则是柯西（内积公理）的齐次化特例。

\key{（1）基本不等式 $=$ 半圆射影定理。}
作直径 $AB=a+b$ 的半圆，圆心 $O$，则半径 $R=\dfrac{a+b}{2}$。
设分界点 $C$ 满足 $AC=a$、$CB=b$，过 $C$ 作 $CP\perp AB$ 交半圆于 $P$。
因为直径所对圆周角为直角，$\triangle APB$ 是直角三角形，由射影定理
\[
  PC^2=AC\cdot CB=ab\ \Rightarrow\ PC=\sqrt{ab}.
\]
又在直角三角形 $OPC$ 中 $OP$ 是斜边，故 $PC\le OP=R$，即
\[
  \sqrt{ab}\le\frac{a+b}{2},
\]
当且仅当 $C$ 与 $O$ 重合（即 $a=b$）时取等（图 \ref{fig:amgm-semicircle}）。

\begin{figure}[htbp]
  \centering
  \begin{tikzpicture}[>=Stealth, semithick, font=\footnotesize, scale=1.05]
    % 直径 AB=a+b，取 a=2.4,b=1.2 => R=1.8，C 位于 x=0.6，PC=sqrt(2.88)≈1.697
    \draw[->, gray] (-2.35,0)--(2.35,0);
    \draw[thick] (-1.8,0)--(1.8,0);
    \draw[thick] (1.8,0) arc (0:180:1.8);
    \fill (-1.8,0) circle (1.5pt) node[below left] {$A$};
    \fill (1.8,0) circle (1.5pt) node[below right] {$B$};
    \fill (0,0) circle (1.5pt) node[below] {$O$};
    \fill (0.6,0) circle (1.5pt) node[below] {$C$};
    \fill (0.6,1.697) circle (1.5pt) node[above right] {$P$};
    \draw[main!70!black, thick] (0.6,0)--(0.6,1.697);
    % 标签放在 $PC$ 右侧、斜边 $PB$ 下方的空档处
    \node[main!70!black] at (0.95,0.55) {$\sqrt{ab}$};
    \draw[second, thick, dashed] (0,0)--(0,1.8);
    % 半径标签锚定在虚线右侧起排，避免压线
    \node[second, anchor=west] at (0.72,1.25) {$\dfrac{a+b}{2}$};
    \draw[third!70!black, thick] (-1.8,0)--(0.6,1.697)--(1.8,0);
    \node at (-0.6,-0.4) {$a$};
    \node at (1.2,-0.4) {$b$};
    \node[gray!70!black, align=center] at (0,-2.15) {$PC^2=AC\cdot CB=ab\Rightarrow PC=\sqrt{ab}\le OP=\dfrac{a+b}{2}$\\（当且仅当 $C$ 与 $O$ 重合即 $a=b$ 时取等）};
  \end{tikzpicture}
  \caption{半圆射影定理：直径 $AB=a+b$，半径 $R=\frac{a+b}{2}$ 永远是斜边上的高 $PC=\sqrt{ab}$ 的上界——这就是基本不等式的几何模型}
  \label{fig:amgm-semicircle}
\end{figure}

\key{（2）柯西 $=$ 向量内积公理；权方和 $=$ 柯西的齐次化。}
$|\vec u\cdot\vec v|=|\vec u|\,|\vec v|\,|\cos\theta|\le|\vec u|\,|\vec v|$，平方即得二维柯西，
取等当且仅当两向量共线。代数证法（判别式法）：
\begin{derivation}
令 $\varphi(x)=(a_1x+b_1)^2+(a_2x+b_2)^2\ge0$ 对一切实数 $x$ 成立，展开得
\[
  (a_1^2+a_2^2)x^2+2(a_1b_1+a_2b_2)x+(b_1^2+b_2^2)\ge0 .
\]
二次函数恒非负，故判别式 $\Delta\le0$：
\[
  4(a_1b_1+a_2b_2)^2-4(a_1^2+a_2^2)(b_1^2+b_2^2)\le0,
\]
即 $(a_1^2+a_2^2)(b_1^2+b_2^2)\ge(a_1b_1+a_2b_2)^2$。
取等当且仅当存在 $x$ 使 $a_1x+b_1=a_2x+b_2=0$，即 $(a_1,a_2)\parallel(b_1,b_2)$。
\end{derivation}
权方和只是把 $\vec u=\left(\dfrac{a_1}{\sqrt{b_1}},\dfrac{a_2}{\sqrt{b_2}}\right)$、
$\vec v=(\sqrt{b_1},\sqrt{b_2})$ 代入柯西：
\[
  \left(\frac{a_1^2}{b_1}+\frac{a_2^2}{b_2}\right)(b_1+b_2)
  \ge\left(\frac{a_1}{\sqrt{b_1}}\cdot\sqrt{b_1}+\frac{a_2}{\sqrt{b_2}}\cdot\sqrt{b_2}\right)^2
  =(a_1+a_2)^2 .
\]

\begin{conccard}[二级]{五大不等式的结构本质与取等条件}{证明不等式、选择放缩工具、检验取等}{只背公式不看结构；张冠李戴（该用柯西时硬凑基本不等式）；忽视取等条件导致「最值取不到」}
基本不等式＝半圆射影定理（$a=b$ 取等）；
柯西＝内积公理（共线取等）；
权方和＝柯西的齐次化（$\frac{a_1}{b_1}=\frac{a_2}{b_2}$ 取等）；
绝对值三角＝度量三角公理（共线取等）；
切线放缩＝凸凹函数的切线支撑（切点取等）。
\end{conccard}

\subsection{切线放缩与隐零点}

\key{探究四：切线是全局支撑边界；隐零点靠「设而不求」}

\key{（1）$e^x\ge x+1$ 与 $\ln x\le x-1$ 的严格证明。}
\begin{derivation}
设 $g(x)=e^x-x-1$，则 $g'(x)=e^x-1$。
$x<0$ 时 $g'<0$，$x>0$ 时 $g'>0$，故 $g$ 在 $x=0$ 处取最小值 $g(0)=0$。
所以
\[
  e^x\ge x+1\quad(\text{当且仅当 }x=0\text{ 取等}).
\]
设 $h(x)=x-1-\ln x\ (x>0)$，则 $h'(x)=1-\frac1x$。
$0<x<1$ 时 $h'<0$，$x>1$ 时 $h'>0$，故 $h$ 在 $x=1$ 处取最小值 $h(1)=0$。
所以
\[
  \ln x\le x-1\quad(\text{当且仅当 }x=1\text{ 取等}).
\]
几何含义：$y=e^x$ 是凸函数，恒在切线 $y=x+1$ 上方；$y=\ln x$ 是凹函数，恒在切线 $y=x-1$ 下方（图 \ref{fig:tangent-exp}）。
\end{derivation}

\begin{figure}[htbp]
  \centering
  \begin{tikzpicture}[>=Stealth, semithick, font=\footnotesize]
    \draw[->, gray] (-2.2,0)--(1.95,0) node[below right] {$x$};
    \draw[->, gray] (0,-1.2)--(0,5.1) node[above left] {$y$};
    \fill[main!10] plot[domain=-2:1.6, samples=60] (\x, {exp(\x)})
      -- plot[domain=1.6:-2, samples=60] (\x, {\x+1}) -- cycle;
    \draw[main!70!black, very thick, domain=-2:1.6, smooth, samples=80] plot (\x, {exp(\x)});
    \draw[second, thick, domain=-2:1.6] plot (\x, {\x+1});
    \fill (0,1) circle (1.6pt) node[above right] {$(0,1)$};
    \node[main!70!black] at (0.45,3.25) {$y=e^x$};
    \node[second] at (-1.85,-0.55) {$y=x+1$};
    % 结论文字移到左上空白区，避免被 $e^x$ 曲线或切线穿过
    \node[gray!70!black] at (-0.35,4.35) {$e^x\ge x+1$};
    \node[gray!70!black] at (-0.35,3.95) {仅在 $x=0$ 取等};
  \end{tikzpicture}
  \caption{$y=e^x$ 恒在切线 $y=x+1$ 上方（凸函数的切线支撑）：阴影为差值区域 $e^x-(x+1)\ge0$，两者仅在切点 $(0,1)$ 接触}
  \label{fig:tangent-exp}
\end{figure}

\key{（2）多步放缩：各步取等条件必须同时成立。}
\begin{example}
证明：对 $x>0$，$e^x-\ln x>2$。
\end{example}
\begin{solution}
由 $e^x\ge x+1$（取等 $x=0$）与 $-\ln x\ge1-x$（取等 $x=1$）相加得
\[
  e^x-\ln x\ge (x+1)+(1-x)=2 .
\]
但两式取等条件分别是 $x=0$ 与 $x=1$，不可能同时成立，故等号永远取不到：
\[
  e^x-\ln x>2 .
\]
\key{多步放缩时，每一步的取等条件必须相容，否则最后的「$\ge$」其实是严格的「$>$」。}
\end{solution}

\key{（3）隐零点：设而不求，整体代换。}
当 $f'(x)=0$ 的解 $x_0$ 无法用初等形式写出时，不必硬解：设 $x_0$ 为零点，
利用 $f'(x_0)=0$ 提供的等量关系做整体代换，把超越结构化为代数结构。

\begin{example}
证明：对一切 $x>-2$，$e^x>\ln(x+2)$。
\end{example}
\begin{solution}
设 $f(x)=e^x-\ln(x+2)$，$f'(x)=e^x-\dfrac1{x+2}$，$f''(x)=e^x+\dfrac1{(x+2)^2}>0$，
故 $f'$ 在 $(-2,+\infty)$ 上严格递增。
又
\[
  f'\!\left(-\tfrac12\right)=e^{-1/2}-\frac23<0,\qquad f'(0)=1-\frac12>0,
\]
所以存在唯一的 $x_0\in\left(-\dfrac12,0\right)$ 使 $f'(x_0)=0$，即
\[
  e^{x_0}=\frac1{x_0+2}\quad\Longrightarrow\quad \ln(x_0+2)=-x_0 .
\]
$f$ 在 $(-2,x_0)$ 上递减、在 $(x_0,+\infty)$ 上递增，最小值为
\[
  f(x_0)=e^{x_0}-\ln(x_0+2)=\frac1{x_0+2}+x_0 .
\]
由 $x_0\in\left(-\dfrac12,0\right)$ 得 $\dfrac1{x_0+2}>\dfrac12$ 且 $x_0>-\dfrac12$，故
\[
  f(x_0)>\frac12-\frac12=0 .
\]
因此 $f(x)\ge f(x_0)>0$，即 $e^x>\ln(x+2)$ 对一切 $x>-2$ 成立。$\surd$
\end{solution}

\begin{conccard}[一级]{切线放缩与隐零点}{证明含 $e^x$、$\ln x$ 的不等式；求不可解方程的极值}{忘记检验各步取等条件；不敢把「$\ge$」收紧为「$>$」；找到 $x_0$ 后不会用 $f'(x_0)=0$ 代换；忽略 $x_0$ 的取值范围}
$e^x\ge x+1$（取等 $x=0$），$\ln x\le x-1$（取等 $x=1$）；
多步放缩取等条件不相容 $\Rightarrow$ 严格不等；
隐零点：设 $f'(x_0)=0$，用该等式整体代换，并确定 $x_0$ 的范围。
\end{conccard}

\subsection{条件最值：四大标准动作}

\key{探究五：把多元、超越的结构，逐个「降维」成对勾函数}

\key{动作一：齐次化与「1」的代换（入口；完整方法链见第 8 章）。}
已知 $ma+nb=k$ 求 $\dfrac pa+\dfrac qb$ 的最值：乘以 $\dfrac{ma+nb}k=1$，展开后用基本不等式。

\begin{example}
$a,b>0$，$a+b=1$，求 $\dfrac1a+\dfrac4b$ 的最小值。
\end{example}
\begin{solution}
「1 的代换」：
\[
  \frac1a+\frac4b=\left(\frac1a+\frac4b\right)(a+b)
  =5+\frac ba+\frac{4a}b\ge5+2\sqrt4=9,
\]
当且仅当 $\dfrac ba=\dfrac{4a}b$ 即 $b=2a$ 时取等，结合 $a+b=1$ 得 $a=\dfrac13,b=\dfrac23$。
「权方和秒杀」：
\[
  \frac1a+\frac4b=\frac{1^2}{a}+\frac{2^2}{b}\ge\frac{(1+2)^2}{a+b}=9,
\]
取等条件同为 $a:b=1:2$。
\end{solution}

\key{动作二：消元与主元转化（双变量 $\to$ 对勾函数 $t+\dfrac kt$）。}

\begin{example}
$x,y>0$ 满足 $x+2y+xy=30$，求 $x+2y$ 的最小值。
\end{example}
\begin{solution}
由等式解出 $y=\dfrac{30-x}{x+2}$。令 $t=x+2$，则 $t>2$，且 $y>0\Rightarrow t<32$，故 $t\in(2,32)$。
\[
  x+2y=t-2+\frac{2(32-t)}t=t+\frac{64}t-4\ge2\sqrt{64}-4=12,
\]
当且仅当 $t=\dfrac{64}t$ 即 $t=8\in(2,32)$ 时取等，此时 $x=6,y=3$（检验：$6+6+18=30$）。
\end{solution}

\key{动作三：拆项分离常数。}

\begin{example}
求 $f(x)=\dfrac{x^2+3x+5}{x+1}$ 在 $x>-1$ 时的最小值。
\end{example}
\begin{solution}
令 $t=x+1>0$，则 $x=t-1$，分子
\[
  x^2+3x+5=(t-1)^2+3(t-1)+5=t^2+t+3,
\]
故
\[
  f(x)=t+1+\frac3t\ge1+2\sqrt3,
\]
当且仅当 $t=\dfrac3t$ 即 $t=\sqrt3$（$x=\sqrt3-1$）时取等。
\end{solution}

\key{动作四：切线放缩放去超越函数。}

\begin{example}
证明：$x>0$ 时 $\ln(x+1)+\dfrac1{x+1}>1$。
\end{example}
\begin{solution}
在 $\ln u\le u-1$ 中取 $u=\dfrac1t\ (t>0)$，得 $\ln t\ge1-\dfrac1t$（取等 $t=1$）。
令 $t=x+1>1$，则
\[
  \ln(x+1)\ge1-\frac1{x+1}\quad\Longrightarrow\quad
  \ln(x+1)+\frac1{x+1}\ge1 .
\]
取等需要 $t=1$ 即 $x=0$，但 $x>0$，故严格成立：
\[
  \ln(x+1)+\frac1{x+1}>1 .
\]
\key{用切线（一阶 Taylor）把 $\ln$、$e^x$ 替换为整式，是消去超越函数的标准动作。}
\end{solution}

\begin{commonerror}{易错诊断室：连续放缩的取等冲突与换元定义域}
\paragraph{反例一：连续两次用基本不等式，取等条件互相打架。}
已知 $x,y>0$，$x+2y=1$，某同学求 $\dfrac1{x^2}+\dfrac1{y^2}$ 的「最小值」：
\[
  \frac1{x^2}+\frac1{y^2}\ge\frac{2}{xy},\qquad
  xy=\frac{x\cdot2y}{2}\le\frac12\left(\frac{x+2y}{2}\right)^2=\frac18,
\]
于是 $\dfrac1{x^2}+\dfrac1{y^2}\ge\dfrac{2}{1/8}=16$。
问题：第一步取等要求 $x=y$，第二步取等要求 $x=2y$，二者不能同时成立，
所以 $16$ 根本取不到（真实最小值大于 $16$）。
\key{每写一次放缩，就把取等条件记下来；所有条件必须相容，最值才可能取到。}

\paragraph{反例二：换元后忽略新元的取值范围。}
求 $f(x)=\dfrac{x^2+5}{\sqrt{x^2+4}}$ 的最小值。令 $t=\sqrt{x^2+4}$，则 $t\ge2$，且
\[
  f(x)=t+\frac1t .
\]
若直接写 $t+\dfrac1t\ge2$，取等要求 $t=1$，但 $t\ge2$，取不到！
正确做法：$\varphi(t)=t+\dfrac1t$ 在 $[2,+\infty)$ 上递增，故最小值为
\[
  \varphi(2)=2+\frac12=\frac52 .
\]
\key{换元必先写出新元的定义域；取等点若不在定义域内，必须改用单调性求最值。}
\end{commonerror}

\begin{conccard}[一级]{条件最值四大标准动作}{二元条件最值、分式函数最值、含超越函数的不等式}{取等条件冲突未检验；换元后忘定义域；该用权方和时硬展开；切线放缩用于非切点附近却当作等式}
（1）1 的代换：乘 $\frac{ma+nb}k=1$ 展开（或直接用权方和）；齐次化的完整展开见第 8 章，本节只作条件最值的入口；
（2）消元主元：解出一元，换 $t$ 凑 $t+\frac kt$；
（3）拆项分离常数：$t=x+1$ 类换元降幂；
（4）切线放缩：用 $e^x\ge x+1$ 等消去超越项。
每步都要检验：取等条件相容吗？在定义域内吗？
\end{conccard}

\subsection{含参恒成立与能成立：点集位置与集合包含}

\key{探究六：把「成立性」翻译成「边界值的大小关系」}

\key{（1）点集上下位置 $\to$ 恒成立。}
$f(x)\ge a$ 在 $D$ 上恒成立，几何上是直线 $y=a$ 始终位于曲线 $y=f(x)$ 的下方（允许相切），
等价于
\[
  a\le f(x)_{\min}\quad(x\in D).
\]
类比地，$f(x)>a$ 恒成立需 $a\le f_{\min}$ 且验证最小值点处是否严格（最小值能取到时需 $a<f_{\min}$）。

\key{（2）一元二次含参不等式：分类讨论树。}
以 $ax^2+px+q>0$ 为例，固定顺序：先看二次项系数 $a$（零／正／负），
再看判别式 $\Delta$（正／零／负），最后比较两根大小（图 \ref{fig:quad-tree}）。

\begin{figure}[htbp]
  \centering
  \begin{tikzpicture}[>=Stealth, semithick, font=\footnotesize,
    box/.style={draw, rounded corners, align=center, inner sep=3pt, text width=2.2cm}]
    \node[box, text width=3.8cm] (root) at (0,0) {解关于 $x$ 的含参不等式\\$ax^{2}+px+q>0$\\($a,p,q$ 为参数)};
    \node[box] (l) at (-4.6,-1.6) {$a=0$\\退化为一次};
    \node[box] (m) at (0,-1.6) {$a>0$\\开口向上};
    \node[box] (r) at (4.6,-1.6) {$a<0$\\开口向下};
    \node[box] (ll) at (-5.2,-3.3) {$px+q>0$\\按 $p$ 符号定向};
    \node[box] (m1) at (-2.6,-3.3) {$\Delta>0$\\两根之外};
    \node[box] (m2) at (0,-3.3) {$\Delta=0$\\$x\ne -\frac{p}{2a}$};
    \node[box] (m3) at (2.6,-3.3) {$\Delta<0$\\解集为 $\mathbb R$};
    \node[box] (rr) at (5.2,-3.3) {乘 $-1$ 变号\\化为 $a>0$};
    \draw (root)--(l) (root)--(m) (root)--(r);
    \draw (l)--(ll) (m)--(m1) (m)--(m2) (m)--(m3) (r)--(rr);
    \node[gray!70!black, align=center, text width=13cm] at (0,-4.9)
      {注：$\Delta>0$ 时两根为 $x_{1,2}=\dfrac{-p\pm\sqrt{\Delta}}{2a}$，\\还必须比较两根大小，然后才能写下解集区间。};
  \end{tikzpicture}
  \caption{含参一元二次不等式的分类讨论树：二次项系数（零／正／负）$\to$ 判别式 $\Delta$（正／零／负）$\to$ 两根大小}
  \label{fig:quad-tree}
\end{figure}

\begin{example}
解关于 $x$ 的不等式 $ax^2-(a+1)x+1>0$。
\end{example}
\begin{solution}
因式分解：$ax^2-(a+1)x+1=(x-1)(ax-1)>0$。
\begin{itemize}
  \item $a=0$：化为 $-(x-1)>0$，解集 $(-\infty,1)$；
  \item $a>0$：两根为 $1$ 与 $\dfrac1a$。
    $a>1$ 时 $\dfrac1a<1$，解集 $\left(-\infty,\dfrac1a\right)\cup(1,+\infty)$；
    $a=1$ 时为 $(x-1)^2>0$，解集 $\{x\mid x\ne1\}$；
    $0<a<1$ 时 $\dfrac1a>1$，解集 $(-\infty,1)\cup\left(\dfrac1a,+\infty\right)$；
  \item $a<0$：开口向下，两根 $\dfrac1a<0<1$，解集 $\left(\dfrac1a,1\right)$。
\end{itemize}
\end{solution}

\key{（3）量词问题：值域边界的对偶。}
「任意」「存在」混合时，统一翻译为两个函数值域的边界关系（设最值都能取到）：

\begin{center}\footnotesize\renewcommand{\arraystretch}{1.3}
\begin{tabular}{c|c}
量词结构 & 等价条件 \\\hline
$\forall x_1\in D_1,\ \forall x_2\in D_2:\ f(x_1)\ge g(x_2)$ & $f_{\min}\ge g_{\max}$ \\[2pt]
$\forall x_1\in D_1,\ \exists x_2\in D_2:\ f(x_1)\ge g(x_2)$ & $f_{\min}\ge g_{\min}$ \\[2pt]
$\exists x_1\in D_1,\ \forall x_2\in D_2:\ f(x_1)\ge g(x_2)$ & $f_{\max}\ge g_{\max}$ \\[2pt]
$\exists x_1\in D_1,\ \exists x_2\in D_2:\ f(x_1)\ge g(x_2)$ & $f_{\max}\ge g_{\min}$ \\
\end{tabular}
\end{center}

\begin{example}
设 $f(x)=e^x-x$，$g(x)=x^2-2x+a$，$x\in[0,2]$。若对任意 $x_1\in[0,2]$，都存在 $x_2\in[0,2]$ 使 $f(x_1)\ge g(x_2)$，求 $a$ 的取值范围。
\end{example}
\begin{solution}
条件等价于 $f_{\min}\ge g_{\min}$。
$f'(x)=e^x-1\ge0$ 在 $[0,2]$ 上成立，$f$ 递增，$f_{\min}=f(0)=1$。
$g(x)=(x-1)^2+a-1$，$g_{\min}=g(1)=a-1$。
故
\[
  1\ge a-1\ \Rightarrow\ a\le2 .
\]
\end{solution}

\begin{conccard}[二级]{恒成立／能成立：点集位置与值域边界}{含参不等式恒成立、双变量量词问题、参数范围求解}{把 $\forall$ 与 $\exists$ 的最值方向搞反；忽略最值是否能取到；严格不等式端点处理错误；二次不等式分类讨论顺序混乱}
第 0 章的集合版（值域子集包含）在这里换成最值语言：$f(x)\ge a$ 恒成立 $\iff a\le f_{\min}$；
$\forall x_1,\exists x_2:\ f(x_1)\ge g(x_2)\iff f_{\min}\ge g_{\min}$（其余三种对偶见表）；
含参二次不等式按「$a$ 的符号 $\to\Delta\to$ 两根大小」三步分类。
\end{conccard}

\begin{method}[局部线性化引擎的使用流程]
\begin{enumerate}
  \item \key{识别结构}：见 $e^x$、$\ln x$ 想切线放缩；见双变量等式想 1 的代换／消元；见含参成立性想最值与值域边界；
  \item \key{局部线性化}：在关键点求导、写切线，用 $f(x)\approx f(x_0)+f'(x_0)(x-x_0)$ 把非线性问题线性化；
  \item \key{求最值}：导数符号法找单调区间与极值，再比较端点；
  \item \key{验取等}：记录每一步的取等条件，检查是否相容、是否在定义域内——不相容则「$\ge$」改「$>$」，不在定义域内则改用单调性；
  \item \key{回代检验}：用特殊值（如 $x=0,1$）或量纲、符号快速验证结论。
\end{enumerate}
\end{method}

\begin{summary}
本章把「导数」与「不等式」统一在\key{局部线性化（以直代曲）}这一个引擎之下：
曲线在一点附近由切线 $f(x_0)+f'(x_0)(x-x_0)$ 近似代替，导数是差分的连续版本；
由 $f'(x)$ 的符号统一判定单调性、极值与最值；
五大核心不等式分别起源于射影定理、内积公理、齐次化、度量三角公理与凸凹函数的切线支撑，取等条件各不相同；
$e^x\ge x+1$、$\ln x\le x-1$ 的严格证明与多步取等的相容性检验，构成切线放缩的完整操作规程；
隐零点靠「设而不求、整体代换」；条件最值有四大标准动作（1 的代换、消元主元、拆项分离、切线放缩）；
含参恒成立／能成立问题则归结为点集的上下位置与值域的边界大小。

\begin{center}
\footnotesize
\setlength{\tabcolsep}{4pt}
\renewcommand{\arraystretch}{1.35}
\begin{tabular}{p{2.3cm}p{4.6cm}p{4.6cm}p{1.9cm}}
\toprule
引擎部件 & 核心操作 & 把什么变成什么 & 对应探究 \\
\midrule
局部线性化 & $f(x)\approx f(x_0)+f'(x_0)(x-x_0)$ & 非线性 $\to$ 线性；变化率 $\to$ 切线斜率 & 探究一 \\
导数符号法 & $f'(x)>0$ 增、$f'(x)<0$ 减 & 单调性／极值／最值 $\to$ 符号判定 & 探究二 \\
结构本质表 & 射影／内积／齐次／度量／凸支撑 & 不等式 $\to$ 几何与代数边界 & 探究三 \\
切线放缩 & $e^x\ge x+1$，$\ln x\le x-1$ & 超越 $\to$ 整式（须验取等） & 探究四 \\
隐零点代换 & 设 $x_0$，用 $f'(x_0)=0$ 整体代换 & 不可解方程 $\to$ 代数关系 & 探究四 \\
条件最值四动作 & 1 的代换／消元主元／拆项／放缩 & 多变量、超越 $\to$ 一元低次 & 探究五 \\
成立性量词化 & $\forall\exists$ 四类转化 & 参数成立性 $\to$ 值域边界 & 探究六 \\
\bottomrule
\end{tabular}
\end{center}

\key{导数与不等式是同一个引擎的两面：以直代曲给出局部逼近，支撑边界给出整体控制；解题的本质是「识别结构 $\to$ 局部线性化 $\to$ 求最值 $\to$ 验取等」。}
\end{summary}

\end{explore}

\begin{inclasstest}
\begin{enumerate}
  \item 用切线近似求 $\sqrt{4.02}$ 的近似值（精确到 $0.001$）；\level{易}
  \item 求 $f(x)=x^3-3x^2-9x+5$ 的单调区间与极值；\level{中}
  \item 证明：对任意 $x>0$，$e^x-\ln x>2$；\level{中}
  \item 解关于 $x$ 的不等式 $ax^2-(a+1)x+1>0$。\level{难}
\end{enumerate}
\end{inclasstest}

\begin{afterclass}
\subsection*{基础巩固}
\begin{enumerate}
  \item 求函数 $f(x)=x\ln x\ (x>0)$ 的最小值；
  \item 证明：对任意实数 $x$，$e^x\ge ex$（提示：求 $f(x)=e^x-ex$ 的最小值）。
\end{enumerate}

\subsection*{进阶提升}
\begin{enumerate}
  \item 已知 $a,b>0$ 且 $a+2b=1$，求 $\dfrac1a+\dfrac1b$ 的最小值；
  \item 设 $f(x)=e^x-ax\ (a>0)$，讨论其单调性；若 $f(x)\ge0$ 对一切 $x\in\mathbb R$ 恒成立，求 $a$ 的取值范围。
\end{enumerate}

\subsection*{拓展探究}
\begin{enumerate}
  \item （强基／竞赛）把权方和不等式推广到三元：证明
  $\dfrac{a_1^2}{b_1}+\dfrac{a_2^2}{b_2}+\dfrac{a_3^2}{b_3}\ge\dfrac{(a_1+a_2+a_3)^2}{b_1+b_2+b_3}\ (b_i>0)$，
  并用它求 $a,b,c>0$、$a+b+c=3$ 时 $\dfrac1a+\dfrac4b+\dfrac9c$ 的最小值。
\end{enumerate}

\subsection*{教师版答案}
\begin{solution}
\textbf{基础 1：}$f'(x)=\ln x+1$。令 $f'(x)=0$ 得 $x=\dfrac1e$。
$0<x<\dfrac1e$ 时 $f'<0$，$x>\dfrac1e$ 时 $f'>0$，故最小值为
\[
  f\!\left(\frac1e\right)=\frac1e\ln\frac1e=-\frac1e .
\]
\end{solution}
\begin{solution}
\textbf{基础 2：}设 $f(x)=e^x-ex$，$f'(x)=e^x-e$。
$x<1$ 时 $f'<0$，$x>1$ 时 $f'>0$，故 $f$ 在 $x=1$ 处取最小值 $f(1)=0$，
即 $e^x\ge ex$，当且仅当 $x=1$ 取等（这正是 $y=e^x$ 在 $(1,e)$ 处的切线 $y=ex$）。
\end{solution}
\begin{solution}
\textbf{进阶 1：}「1 的代换」：
\[
  \frac1a+\frac1b=\left(\frac1a+\frac1b\right)(a+2b)
  =3+\frac{2b}{a}+\frac ab\ge3+2\sqrt2,
\]
当且仅当 $\dfrac{2b}{a}=\dfrac ab$ 即 $a=\sqrt2\,b$ 时取等，结合 $a+2b=1$ 得
$b=\dfrac1{2+\sqrt2}$，$a=\dfrac{\sqrt2}{2+\sqrt2}$。最小值 $3+2\sqrt2$。
\end{solution}
\begin{solution}
\textbf{进阶 2：}$f'(x)=e^x-a$。$x<\ln a$ 时 $f'<0$，$x>\ln a$ 时 $f'>0$，
故 $f$ 在 $(-\infty,\ln a)$ 递减、在 $(\ln a,+\infty)$ 递增，最小值为
\[
  f(\ln a)=a-a\ln a=a(1-\ln a).
\]
$f(x)\ge0$ 恒成立 $\iff a(1-\ln a)\ge0\iff\ln a\le1\iff a\le e$。
故 $a$ 的取值范围是 $(0,e]$。
\end{solution}
\begin{solution}
\textbf{拓展：}由柯西不等式，取
$\vec u=\left(\dfrac{a_1}{\sqrt{b_1}},\dfrac{a_2}{\sqrt{b_2}},\dfrac{a_3}{\sqrt{b_3}}\right)$，
$\vec v=(\sqrt{b_1},\sqrt{b_2},\sqrt{b_3})$，得
\[
  \left(\frac{a_1^2}{b_1}+\frac{a_2^2}{b_2}+\frac{a_3^2}{b_3}\right)(b_1+b_2+b_3)
  \ge(a_1+a_2+a_3)^2,
\]
即所要的不等式，取等当且仅当 $\dfrac{a_1}{b_1}=\dfrac{a_2}{b_2}=\dfrac{a_3}{b_3}$。
应用：
\[
  \frac1a+\frac4b+\frac9c=\frac{1^2}{a}+\frac{2^2}{b}+\frac{3^2}{c}
  \ge\frac{(1+2+3)^2}{a+b+c}=\frac{36}{3}=12,
\]
当且仅当 $a:b:c=1:2:3$ 且 $a+b+c=3$，即 $a=\dfrac12,b=1,c=\dfrac32$ 时取等。最小值为 $12$。
\end{solution}
\reading{拓展阅读 12}{导数：以直代曲的局部线性化}
\end{afterclass}
% 高观点导学案：计数原理、概率测度空间与统计线性逼近（第 7 章·高观点版）
% 主线：把「不确定性」纳入集合框架——样本空间是集合，概率是其上的归一化计数测度；
%       条件概率 = 样本全集缩小后的空间缩放；全概率/贝叶斯 = 完备互斥划分；
%       回归直线 = 散点云的竖直最小二乘投影；独立性检验 = 小概率事件原理的统计反证法；
%       压轴收官：新定义情境题的四步破题法（翻译 → 特例 → 归纳 → 严格证明）。

\section{计数原理、概率测度空间与统计线性逼近（高观点版）}

\begin{methodbox}{多框架对读}
\small
\begin{center}
\begin{tabular}{c|p{2.9cm}|p{5.0cm}|p{3.1cm}}
\toprule
视角 & 本章看点 & 别的视角怎么说 & 哪侧最省力 \\
\midrule
代数 & 计数公式、期望方差 & 符号对象的机械运算 & 硬算分布列 \\
几何 & 散点、回归直线、面积测度 & 图形的位置与拟合 & 看趋势、估相关 \\
集合 & 样本空间、事件、完备划分 & 事件＝子集；概率＝测度 & 界互斥与独立 \\
映射 & $X:\Omega\to\mathbb R$；最小二乘＝投影 & 换到残差正交的坐标系 & 一图读懂相关 \\
\bottomrule
\end{tabular}
\end{center}
卡的使用：概率题先画维恩图与树图（集合），统计题先想「投影到哪个子空间」（总图见第 0 章图~\ref{fig:four-view-map}）。
\end{methodbox}

第 0 章确立了集合与逻辑这一全书语言，第 4 章把向量确立为空间的通用线性语言，第 6 章把曲线局部化为直线；
本章把「不确定性」也纳入同一个框架——随机试验的样本空间 $\Omega$ 是一个集合，概率是定义在其上的\key{归一化计数测度}；
条件概率 $P(B\mid A)$ 的本质是把样本全集 $\Omega$ 缩小为子集 $A$ 之后的\key{空间缩放}；
线性回归方程是二维散点云在竖直方向上的\key{最小二乘投影直线}，独立性检验则是「小概率事件原理」支配下的\key{统计学反证法}。
四个模块层层展开：
\key{计数原理的集合本质}（并集基数与笛卡尔积基数，排列＝有序单射、组合＝无序子集族，捆绑／插空／隔板三技法）$\to$
\key{概率空间的归一化测度}（古典与几何概型的统一除法，条件概率＝空间缩放，全概率与贝叶斯＝完备互斥划分）$\to$
\key{统计线性逼近}（百分位数规则，回归直线必过样本中心，独立性检验的统计反证）$\to$
\key{新定义情境题破题模板}（翻译符号集合 $\to$ 构造小规模特例 $\to$ 归纳猜想通项 $\to$ 反证法／数学归纳法严格证明）。
核心思想：\key{随机性不是结构的对立面，而是有限集合上的测度；统计不是死记公式，而是「投影逼近 $+$ 小概率反证」；压轴新定义题的本质是「把陌生符号翻译回熟悉的集合与结构」。}

\begin{learnobjective}
\begin{itemize}
  \item 理解分类加法原理即互斥子集并集的基数公式 $|A\cup B|=|A|+|B|$，分步乘法原理即笛卡尔积的基数公式 $|A\times B|=|A|\cdot|B|$；
  \item 理解排列是有序单射序列、组合是无序 $k$ 元子集族，掌握 $A_n^k$ 与 $C_n^k$ 的关系 $A_n^k=C_n^k\cdot k!$；
  \item 掌握捆绑法、插空法、隔板法的适用结构，会把「相邻／不相邻／相同元素分堆」问题翻译为标准计数模型；
  \item 会把古典概型与几何概型统一为归一化测度 $P(A)=\mu(A)/\mu(\Omega)$，说出概率三条公理（非负性、规范性、可加性）；
  \item 理解条件概率 $P(B\mid A)=\dfrac{P(A\cap B)}{P(A)}$ 是样本空间缩小为 $A$ 后的重新归一化，会用全概率公式与贝叶斯公式；
  \item 掌握百分位数的计算规则，理解回归直线 $\hat y=\hat bx+\hat a$ 是竖直残差平方和最小的投影直线且必过 $(\bar x,\bar y)$；
  \item 理解独立性检验 $K^2$ 统计量背后的「小概率事件原理」，能正确陈述检验结论；
  \item 掌握新定义情境题的四步破题法（读懂符号集合 $\to$ 小规模特例 $\to$ 归纳猜想 $\to$ 反证法／数学归纳法严格证明）。
\end{itemize}
\end{learnobjective}

\begin{preview}
\subsection*{知识回顾}
\begin{itemize}
  \item 若 $A\cap B=\varnothing$，则 $|A\cup B|=$ \blank；
  \item 完成一件事有两类互斥的途径，第一类有 $n_1$ 种方法、第二类有 $n_2$ 种方法，则共有 \blank 种方法；
  \item 古典概型：$P(A)=$ \blank；几何概型：$P(A)=$ \blank；
  \item 样本相关系数 $r$ 满足 $|r|$ \blank $1$，$r>0$ 表示两个变量 \blank 相关；
  \item 反证法的逻辑结构：假设结论不成立，推出 \blank，从而肯定原结论。
\end{itemize}

\subsection*{教材速读}
阅读「计数原理」「概率」「成对数据的统计」相关教材，完成下列问题：
\begin{enumerate}
  \item 排列数公式 $A_n^m=$ \blank；组合数公式 $C_n^m=$ \blank；
  \item 条件概率公式 $P(B\mid A)=$ \blank（要求 \blank $\ne0$）；
  \item 全概率公式：若 $A_1,\dots,A_n$ 是 $\Omega$ 的完备互斥划分，则 $P(B)=$ \blank；
  \item 回归直线 $\hat y=\hat bx+\hat a$ 必过样本中心 \blank，且 $\hat a=$ \blank（用 $\bar x,\bar y,\hat b$ 表示）；
  \item 独立性检验统计量 $K^2=$ \blank（设 $2\times2$ 列联表中四个数据为 $a,b,c,d$）。
\end{enumerate}
\end{preview}

\begin{explore}

\subsection{计数原理：并集、笛卡尔积与三类经典技法}

\key{探究一：计数不是数数，而是集合的度量}

\key{（1）两个原理的集合本质。}
设完成一件事的所有可能结果构成集合 $\Omega$。若 $\Omega$ 被分成两个互斥的子集 $\Omega=A\sqcup B$（即 $A\cap B=\varnothing$），则
\[
  |\Omega|=|A\sqcup B|=|A|+|B|,
\]
这就是\key{分类加法原理}：各类之间「不重不漏」，求总数就是求互斥子集并集的基数。
若完成一件事要分两步：第一步的结果集为 $A$、第二步的结果集为 $B$，则所有可能结果组成笛卡尔积
\[
  \Omega=A\times B=\{(a,b)\mid a\in A,\ b\in B\},\qquad |\Omega|=|A\times B|=|A|\cdot|B|,
\]
这就是\key{分步乘法原理}：每一步的选择共同确定一个有序对，求总数就是求积集的基数。
一句话对照：加法管「或」（不交并），乘法管「且」（有序积）。

\begin{example}
从 5 名男同学和 4 名女同学中选出 3 人组成学习小组（不分职务），有多少种选法？若从中选出 3 人分别担任语文、数学、英语三科课代表，又有多少种选法？
\end{example}
\begin{solution}
第一问只关心「选出的集合」，是无序的 $3$ 元子集族，用组合：
\[
  C_9^3=\frac{9!}{3!\,6!}=84 .
\]
第二问关心「谁对应哪一科」，是有序安排：先选后排，由分步乘法原理
\[
  C_9^3\cdot A_3^3=84\times6=504\quad(\text{或直接 }A_9^3=9\times7\times8=504).
\]
两问相差 $3!=6$ 倍——正是「无序」与「有序」的差别。
\end{solution}

\key{（2）排列与组合：有序单射与无序子集。}
从 $n$ 个不同元素中取出 $k$ 个排成一列，相当于构造一个单射 $\{1,2,\dots,k\}\hookrightarrow\{1,2,\dots,n\}$，其个数为排列数
\[
  A_n^k=n(n-1)\cdots(n-k+1)=\frac{n!}{(n-k)!}.
\]
从 $n$ 个不同元素中取出 $k$ 个组成子集（不计顺序），其个数为组合数
\[
  C_n^k=\frac{A_n^k}{k!}=\frac{n!}{k!\,(n-k)!},
\]
因为「先选后排」＝「选子集」乘「子集全排列」：$A_n^k=C_n^k\cdot k!$——这正是分步乘法原理的一次应用。
特别地，$C_n^k=C_n^{n-k}$（选 $k$ 个留下等价于选 $n-k$ 个淘汰）。

\key{（3）三类经典技法的结构对号。}
\begin{itemize}
  \item \key{捆绑法（必相邻）}：把必须相邻的 $m$ 个元素捆成一个「大元素」，先与其他元素一起排，再解开内部排列：
    $n$ 个元素中指定 $m$ 个相邻的排法 $=(n-m+1)!\cdot A_m^m$。
  \item \key{插空法（不相邻）}：先排其余 $n$ 个元素（产生 $n+1$ 个空位，含两端），再把 $m$ 个元素插入 $m$ 个不同的空位：$A_n^n\cdot A_{n+1}^{\,m}$。（此处的 $n$ 指「其余元素个数」，与上一行捆绑法中 $n$ 为总元素数的口径不同。）
  \item \key{隔板法（相同元素分堆）}：不定方程 $x_1+x_2+\cdots+x_m=n\ (x_i\in\mathbb N^*)$ 的正整数解，对应把 $n$ 个相同小球排成一排后在 $n-1$ 个空隙中插入 $m-1$ 块板，解数为
    \[
      C_{n-1}^{\,m-1}.
    \]
    若允许 $x_i=0$，先令 $x_i'=x_i+1$ 化为正整数解，得 $C_{n+m-1}^{\,m-1}$。
\end{itemize}

\begin{example}
（1）6 人站成一排，甲、乙必须相邻，有多少种站法？
（2）6 人站成一排，甲、乙互不相邻，有多少种站法？
（3）求 $x+y+z=10$ 的正整数解的组数。
\end{example}
\begin{solution}
（1）捆绑：甲乙捆成一个大元素，与其余 4 人共 5 个元素全排列，甲乙内部再交换：
\[
  A_5^5\cdot A_2^2=120\times2=240 .
\]
（2）插空：先排其余 4 人，产生 5 个空位，甲乙选 2 个空位排列：
\[
  A_4^4\cdot A_5^2=24\times20=480 .
\]
（3）隔板：10 个相同小球排一排，9 个空隙插 2 块板分成 3 份（每份至少 1）：
\[
  C_9^2=36 .
\]
\end{solution}

\begin{figure}[htbp]
  \centering
  \begin{tikzpicture}[>=Stealth, semithick, font=\footnotesize]
    % ---------- 左：文氏图 ----------
    \draw[gray] (-1.7,-1.6) rectangle (3.9,1.6);
    \begin{scope}
      \clip (0,0) circle (1.15);
      \fill[second!20] (-1.8,-1.7) rectangle (4.0,1.7);
    \end{scope}
    \begin{scope}
      \clip (1.4,0) circle (1.15);
      \fill[third!20] (-1.8,-1.7) rectangle (4.0,1.7);
    \end{scope}
    \begin{scope}
      \clip (0,0) circle (1.15);
      \fill[main!32] (1.4,0) circle (1.15);
    \end{scope}
    \draw (0,0) circle (1.15);
    \draw (1.4,0) circle (1.15);
    \node at (-0.62,0) {$A$};
    \node at (2.02,0) {$B$};
    \node[main!70!black] at (0.7,-0.05) {\small \color{white}$A\cap B$};
    \node[gray!70!black] at (0.7,1.4) {$A\cup B$};
    \node[gray!70!black] at (-1.5,1.4) {$U$};
    % ---------- 右：卡诺图式网格 ----------
    \begin{scope}[xshift=5.4cm, yshift=0.5cm]
      \draw (0,0) rectangle (3.0,2.4);
      \draw (1.5,0)--(1.5,2.4);
      \draw (0,1.2)--(3.0,1.2);
      \node at (0.75,1.8) {$A\cap B$};
      \node at (2.25,1.8) {$A\cap\bar B$};
      \node at (0.75,0.6) {$\bar A\cap B$};
      \node at (2.25,0.6) {$\bar A\cap\bar B$};
      \node[above] at (0.75,2.4) {$B$};
      \node[above] at (2.25,2.4) {$\bar B$};
      \node[left] at (0,1.8) {$A$};
      \node[left] at (0,0.6) {$\bar A$};
    \end{scope}
  \end{tikzpicture}
  \caption{集合运算的卡诺图式表示：左图文氏图中 $A\cup B$ 是「或」、$A\cap B$ 是「且」；右图 $2\times2$ 网格穷举两个事件的四种组合，对应概率论中样本空间的完备互斥划分}
  \label{fig:ch07-venn}
\end{figure}

\begin{commonerror}{易错诊断室：可重复与不可重复、有序与无序}
\paragraph{失误一：有序、无序混淆。}
从 5 人中选 3 人「组成代表团」用 $C_5^3=10$；若「分别担任三个不同职务」就要用 $A_5^3=60$。
题干中出现「分配、担任、排成一列、编号」等词，结果自带顺序，必须按有序处理。
\paragraph{失误二：平均分组忘记除序。}
把 6 人平均分成 3 组（每组 2 人）：若直接写 $C_6^2C_4^2C_2^2=90$，等于把三个无区别的组又按顺序数了一遍，
正确的组数为
\[
  \frac{C_6^2C_4^2C_2^2}{3!}=\frac{90}{6}=15 .
\]
\key{分组若无编号（与「分到三个不同地方」区别开），组间要除以组数的阶乘。}
\paragraph{失误三：可重复与不可重复混淆。}
3 个人从 5 种口味冰淇淋中各选一种：口味可被多人重复选择时，由分步乘法原理有 $5^3=125$ 种；
若规定三人口味互不相同，则相当于从 5 种中取 3 种的排列 $A_5^3=60$。
一字之差，原理全变——务必先确认「选过的还能不能再选」。
\end{commonerror}

\begin{conccard}[一级]{计数原理的集合本质与三技法}{一切计数问题，尤其是含「相邻／不相邻／平均分组」的排列组合题}{有序无序混淆；平均分组忘除序；可重复当成不可重复；隔板法忽略「每份至少一个」}
分类加法＝互斥子集并集基数：$|A\sqcup B|=|A|+|B|$；
分步乘法＝笛卡尔积基数：$|A\times B|=|A|\cdot|B|$；
排列＝有序单射序列 $A_n^k=\frac{n!}{(n-k)!}$，组合＝无序子集族 $C_n^k=\frac{A_n^k}{k!}$；
捆绑：$(n-m+1)!\,A_m^m$；插空：$A_n^nA_{n+1}^{\,m}$；
隔板：$x_1+\cdots+x_m=n$ 的正整数解数 $=C_{n-1}^{\,m-1}$。
\end{conccard}

\subsection{概率测度：归一化、空间缩放与完备划分}

\key{探究二：概率是定义在样本空间上的「归一化集合测度」}

\key{（1）古典概型与几何概型的统一。}
任何随机试验的样本空间 $\Omega$ 都是一个集合。在其上定义一个「大小」函数 $\mu$（测度）：
古典概型中 $\mu$ 是计数测度（$\mu(A)=|A|$），几何概型中 $\mu$ 是长度／面积／体积。
只要每个基本结果等可能，就有统一的归一化公式
\[
  P(A)=\frac{\mu(A)}{\mu(\Omega)}.
\]
由此可直接验证概率的三条公理：非负性 $P(A)\ge0$、规范性 $P(\Omega)=1$、互斥可加性 $P(A\sqcup B)=P(A)+P(B)$——它们不过是测度的归一化性质。

\begin{example}
在区间 $[0,1]$ 上任取两点 $X,Y$，求 $|X-Y|\le\dfrac12$ 的概率。
\end{example}
\begin{solution}
样本空间是单位正方形 $\Omega=[0,1]^2$，$\mu(\Omega)=1$。
满足 $|X-Y|>\dfrac12$ 的区域是直线 $Y=X\pm\dfrac12$ 之外的两个三角形，每个面积为 $\dfrac12\left(\dfrac12\right)^2=\dfrac18$。
故
\[
  P(|X-Y|\le\tfrac12)=1-2\times\tfrac18=\tfrac34 .
\]
古典（数点）与几何（量面积）在此只是测度 $\mu$ 的不同，算法同为「除以全集」。
\end{solution}

\key{（2）条件概率 $P(B\mid A)$：样本空间的缩放。}
事件 $A$ 发生后，样本全集由 $\Omega$ 缩小为子集 $A$，每个结果仍在 $A$ 内等可能。
新样本空间上的归一化给出
\[
  P(B\mid A)=\frac{P(A\cap B)}{P(A)}\qquad(P(A)>0).
\]
几何读法：\key{把 $A$ 当作新的「全宇宙」，概率质量重新按 $P(A)$ 缩放}。它仍满足三条公理（规范性 $P(\Omega\mid A)=1$、可加性），因此是以 $A$ 为底的合法概率。

\begin{example}
10 张考签中有 1 张中奖，5 人依次不放回抽取。证明：无论先抽后抽，每人中奖概率都等于 $\dfrac1{10}$（抽签公平性）。
\end{example}
\begin{solution}
第 1 人显然为 $\frac1{10}$。
第 2 人中奖 $=$「第 1 人未中且第 2 人中」，由乘法公式与缩小的样本空间：
\[
  P=P(\text{第 1 人未中})\cdot P(\text{第 2 人中}\mid\text{第 1 人未中})
  =\frac{9}{10}\times\frac{1}{9}=\frac{1}{10}.
\]
第 3 人同理：$\frac{9}{10}\times\frac{8}{9}\times\frac{1}{8}=\frac1{10}$。
第 4、5 人同理（依次缩小样本空间后概率仍为 $\frac1{10}$）。
可见「先抽优势」是错觉——条件概率把空间缩小后自动抹平了顺序差异。
\end{solution}

\key{（3）全概率公式与贝叶斯公式：完备互斥划分。}
若 $\Omega=A_1\sqcup A_2\sqcup\cdots\sqcup A_n$（$A_i$ 两两不交、并为全集，称为\key{完备互斥划分}），则
\[
  P(B)=\sum_{i=1}^{n}P(A_i)\,P(B\mid A_i)\quad(\text{全概率公式}),
\]
它是面积可加性 $\mu(B)=\sum\mu(A_i\cap B)$ 的直接翻译：把 $B$ 沿划分切成互不重叠的碎片再求和。
反过来，已知 $B$ 发生了，反推它是哪个碎片贡献的：
\[
  P(A_j\mid B)=\frac{P(A_j)\,P(B\mid A_j)}{\sum\limits_{i=1}^{n}P(A_i)\,P(B\mid A_i)}\quad(\text{贝叶斯公式}).
\]
全概率是「由因及果」，贝叶斯是「由果溯因」：$P(A_j)$ 是先验概率，$P(A_j\mid B)$ 是纳入证据 $B$ 后的后验概率。

\begin{example}
某厂三条流水线产量占比为 $50\%,30\%,20\%$，次品率分别为 $2\%,3\%,5\%$。任取一件产品，（1）求是次品的概率；（2）若它是次品，求它来自第一条流水线的概率。
\end{example}
\begin{solution}
记 $A_i$ 为「产品来自第 $i$ 条流水线」，$B$ 为「取到次品」。
（1）由全概率公式
\[
  P(B)=0.5\times0.02+0.3\times0.03+0.2\times0.05=0.010+0.009+0.010=0.029 .
\]
（2）由贝叶斯公式
\[
  P(A_1\mid B)=\frac{0.5\times0.02}{0.029}=\frac{0.010}{0.029}=\frac{10}{29}\approx0.345 .
\]
\end{solution}

\begin{figure}[htbp]
  \centering
  \begin{tikzpicture}[>=Stealth, semithick, font=\footnotesize]
    % ---------- 左：完备划分的面积分解 ----------
    \draw[thick] (0,0) rectangle (8,3);
    \draw (4,0)--(4,3);
    \draw (6.4,0)--(6.4,3);
    \fill[main!30] (0,2.4) rectangle (4,3);
    \fill[main!30] (4,2.1) rectangle (6.4,3);
    \fill[main!30] (6.4,1.5) rectangle (8,3);
    \draw[thick] (0,0) rectangle (8,3);
    \draw (4,0)--(4,3);
    \draw (6.4,0)--(6.4,3);
    \node[main!80!black] at (2,2.7) {$0.010$};
    \node[main!80!black] at (5.2,2.55) {$0.009$};
    \node[main!80!black] at (7.2,2.25) {$0.010$};
    \node[align=center] at (2,-0.42) {$A_1$\\ $P(A_1)=0.5$};
    \node[align=center] at (5.2,-0.42) {$A_2$\\ $P(A_2)=0.3$};
    \node[align=center] at (7.2,-0.42) {$A_3$\\ $P(A_3)=0.2$};
    \node[gray!70!black] at (-0.35,3.3) {$\Omega$};
    \node[align=center] at (4,-1.15) {$P(B)=0.010+0.009+0.010=0.029$\\(列宽 $\propto P(A_i)$，块高 $\propto P(B\mid A_i)$ 统一放大 $30$ 倍)\\(暗色块面积 $\propto P(A_i\cap B)$，数字为 $P(A_i\cap B)$)};
    % ---------- 右：概率树 ----------
    \begin{scope}[xshift=9.2cm]
      \node[circle,fill=black,inner sep=1.6pt] (r) at (0,1.85) {};
      \node[above left] at (r) {$\Omega$};
      \node[circle,fill=black,inner sep=1.6pt] (a1) at (1.8,3.25) {};
      \node[circle,fill=black,inner sep=1.6pt] (a2) at (1.8,1.85) {};
      \node[circle,fill=black,inner sep=1.6pt] (a3) at (1.8,0.45) {};
      \node[right] at (a1) {$A_1$};
      \node[right] at (a2) {$A_2$};
      \node[right] at (a3) {$A_3$};
      \node[circle,fill=main,inner sep=1.6pt] (b1) at (4.2,3.75) {};
      \node[circle,fill=gray,inner sep=1.6pt] (nb1) at (4.2,2.75) {};
      \node[circle,fill=main,inner sep=1.6pt] (b2) at (4.2,2.35) {};
      \node[circle,fill=gray,inner sep=1.6pt] (nb2) at (4.2,1.35) {};
      \node[circle,fill=main,inner sep=1.6pt] (b3) at (4.2,0.95) {};
      \node[circle,fill=gray,inner sep=1.6pt] (nb3) at (4.2,-0.05) {};
      \draw (r)--(a1) (r)--(a2) (r)--(a3);
      \draw[main] (a1)--(b1); \draw[gray] (a1)--(nb1);
      \draw[main] (a2)--(b2); \draw[gray] (a2)--(nb2);
      \draw[main] (a3)--(b3); \draw[gray] (a3)--(nb3);
      \node at (0.55,2.75) {$0.5$};
      \node at (0.9,2.15) {$0.3$};
      \node at (0.55,1.0) {$0.2$};
      \node[main,left] at (2.95,3.65) {$0.02$}; \node[gray,left] at (2.95,2.9) {$0.98$};
      \node[main,left] at (2.95,2.25) {$0.03$}; \node[gray,left] at (2.95,1.5) {$0.97$};
      \node[main,left] at (2.95,0.85) {$0.05$}; \node[gray,left] at (2.95,0.1) {$0.95$};
      \node[main,right=2mm] at (b1) {$B$: $0.010$};
      \node[gray,right=2mm] at (nb1) {$B^c$: $0.490$};
      \node[main,right=2mm] at (b2) {$B$: $0.009$};
      \node[gray,right=2mm] at (nb2) {$B^c$: $0.291$};
      \node[main,right=2mm] at (b3) {$B$: $0.010$};
      \node[gray,right=2mm] at (nb3) {$B^c$: $0.190$};
      \node[align=center] at (2.1,-0.8) {$P(B)=0.010+0.009+0.010=0.029$\\[2pt]$P(A_1\mid B)=\frac{0.010}{0.029}=\frac{10}{29}\approx0.345$};
    \end{scope}
  \end{tikzpicture}
  \caption{样本空间的完备互斥划分（全概率·贝叶斯）：左图中 $\Omega=A_1\sqcup A_2\sqcup A_3$，列宽正比于 $P(A_i)$、每列内暗色块高度正比于 $P(B\mid A_i)$（统一放大 $30$ 倍的示意图），故暗色块面积正比于 $P(A_i\cap B)$；右图沿树枝把每个 $A_i$ 再按条件概率分叉，路径概率求和得 $P(B)$，贝叶斯公式则由结果 $B$ 沿树反向回溯到原因 $A_1$}
  \label{fig:ch07-partition}
\end{figure}

\begin{conccard}[一级]{归一化测度：条件概率、全概率与贝叶斯}{几何概型、条件概率、由果溯因类应用题、医学筛查类情境}{把 $P(B\mid A)$ 与 $P(A\mid B)$ 混为一谈；划分不满足「不重不漏」就套全概率；只记住公式不理解「空间缩放」的几何意义}
$P(A)=\mu(A)/\mu(\Omega)$（古典＝计数测度，几何＝长度／面积／体积）；
条件概率＝空间缩放：$P(B\mid A)=P(A\cap B)/P(A)$；
全概率：$P(B)=\sum_i P(A_i)P(B\mid A_i)$；
贝叶斯（由果溯因）：$P(A_j\mid B)=\dfrac{P(A_j)P(B\mid A_j)}{\sum_i P(A_i)P(B\mid A_i)}$。
\end{conccard}

\subsection{统计线性逼近：最小二乘投影与统计反证}

\key{探究三：回归直线是散点云的「最小二乘投影」；独立性检验是统计反证法}

\key{（1）百分位数：位置的度量。}
把 $n$ 个数据从小到大排列，计算 $i=n\times p\%$：
\begin{itemize}
  \item 若 $i$ 不是整数，第 $p$ 百分位数是第 $\lceil i\rceil$ 项（向上取整的下一项）；
  \item 若 $i$ 是整数，第 $p$ 百分位数是第 $i$ 项与第 $i+1$ 项的平均数。
\end{itemize}

\begin{example}
求数据 $12,7,9,15,10,8,11,13,6,14$ 的第 $25$ 百分位数与第 $50$ 百分位数。
\end{example}
\begin{solution}
排序：$6,7,8,9,10,11,12,13,14,15$，共 $n=10$ 个。
\[
  i_1=10\times25\%=2.5\ (\text{非整数})\Rightarrow \text{第 }25\text{ 百分位数是第 }3\text{ 项 }=8;
\]
\[
  i_2=10\times50\%=5\ (\text{整数})\Rightarrow \text{第 }50\text{ 百分位数}=\frac{\text{第 }5\text{ 项}+\text{第 }6\text{ 项}}{2}=\frac{10+11}{2}=10.5 .
\]
\end{solution}

\key{（2）回归直线：竖直残差平方和最小的投影线。}
给定散点 $(x_1,y_1),\dots,(x_n,y_n)$，在直线族 $\hat y=bx+a$ 中挑选一条，使各点到直线的\key{竖直距离}（残差 $e_i=y_i-\hat y_i$）的平方和最小：
\[
  Q(a,b)=\sum_{i=1}^{n}(y_i-bx_i-a)^2\ \longrightarrow\ \min .
\]

\begin{derivation}
对 $a,b$ 分别求偏导并令其为零：
\[
  \frac{\partial Q}{\partial a}=-2\sum_{i=1}^n(y_i-bx_i-a)=0
  \ \Longrightarrow\ na+b\sum x_i=\sum y_i\ \Longrightarrow\ a=\bar y-b\bar x .
\]
\[
  \frac{\partial Q}{\partial b}=-2\sum_{i=1}^nx_i(y_i-bx_i-a)=0 .
\]
把 $a=\bar y-b\bar x$ 代入，利用 $\sum(x_i-\bar x)=0$，得
\[
  \sum_{i=1}^n(x_i-\bar x)(y_i-\bar y)-b\sum_{i=1}^n(x_i-\bar x)^2=0,
\]
故
\[
  \hat b=\frac{\sum\limits_{i=1}^n(x_i-\bar x)(y_i-\bar y)}{\sum\limits_{i=1}^n(x_i-\bar x)^2},
  \qquad \hat a=\bar y-\hat b\,\bar x .
\]
几何性质：$(\hat a,\hat b)$ 满足 $\hat a=\bar y-\hat b\,\bar x$，即点 $(\bar x,\bar y)$ 恒在直线上——\key{回归直线必过样本中心 $(\bar x,\bar y)$}。
这与第 6 章「以直代曲」、第 4 章「正交投影」同源：最小二乘就是把观测向量 $\vec y$ 向由 $\vec 1$ 与 $\vec x$ 张成的平面作投影，残差与平面正交。
\end{derivation}

\begin{example}
已知 5 组成对数据 $(1,2),(2,3),(3,5),(4,6),(5,8)$，求回归方程并计算样本相关系数 $r$。
\end{example}
\begin{solution}
\[
  \bar x=3,\quad \bar y=\frac{2+3+5+6+8}{5}=4.8 .
\]
\[
  \sum(x_i-\bar x)(y_i-\bar y)=(-2)(-2.8)+(-1)(-1.8)+0\times0.2+1\times1.2+2\times3.2=15,
\]
\[
  \sum(x_i-\bar x)^2=4+1+0+1+4=10,\quad
  \sum(y_i-\bar y)^2=7.84+3.24+0.04+1.44+10.24=22.8 .
\]
所以
\[
  \hat b=\frac{15}{10}=1.5,\qquad \hat a=4.8-1.5\times3=0.3,
  \qquad \hat y=1.5x+0.3 .
\]
\[
  r=\frac{15}{\sqrt{10\times22.8}}=\frac{15}{\sqrt{228}}\approx0.993 ,
\]
高度正相关（$|r|\le1$ 正是由第 6 章柯西不等式保证的）。
\end{solution}

\begin{figure}[htbp]
  \centering
  \begin{tikzpicture}[>=Stealth, semithick, font=\footnotesize]
    \draw[->,gray] (-0.4,0)--(5.9,0) node[below right] {$x$};
    \draw[->,gray] (0,-0.6)--(0,8.9) node[above left] {$y$};
    \foreach \x in {1,...,5} \draw (\x,0.06)--(\x,-0.06) node[below] {\x};
    \foreach \y in {2,4,6,8} \draw (0.06,\y)--(-0.06,\y) node[left] {\y};
    \draw[second,very thick,domain=-0.2:5.5,samples=2] plot (\x,{1.5*\x+0.3});
    \node[second] at (0.55,4.6) {$\hat y=1.5x+0.3$};
    \foreach \x/\y in {1/2,2/3,3/5,4/6,5/8}{
      \draw[second!80,dashed,very thick] (\x,\y)--(\x,{1.5*\x+0.3});
      \fill[main] (\x,\y) circle (1.7pt);
    }
    \node[second,right] at (1.04,1.9) {$e_1$};
    \node[second,right] at (2.04,3.15) {$e_2$};
    \node[second,left] at (2.96,4.98) {$e_3$};
    \node[second,right] at (4.04,6.15) {$e_4$};
    \node[second,right] at (5.04,7.9) {$e_5$};
    \draw[third,dashed,thick] (2.68,4.8)--(3.32,4.8);
    \draw[third,dashed,thick] (3,4.48)--(3,5.12);
    \fill[third] (3,4.8) circle (1.7pt);
    \node[third,right] at (3.35,4.86) {$(\bar x,\bar y)=(3,4.8)$};
    \fill[main] (2.9,8.55) circle (1.7pt); \node[right] at (2.98,8.55) {观测点};
    \draw[second,very thick] (4.1,8.55)--(4.5,8.55); \node[right] at (4.58,8.55) {回归直线};
  \end{tikzpicture}
  \caption{最小二乘回归：回归直线 $\hat y=1.5x+0.3$ 使各观测点到直线的竖直残差 $e_1,\dots,e_5$（橙色虚线）平方和最小；它必然穿过样本中心 $(\bar x,\bar y)=(3,4.8)$}
  \label{fig:ch07-leastsq}
\end{figure}

\key{（3）独立性检验：小概率事件原理的统计反证法。}
$2\times2$ 列联表
\[
  \begin{array}{c|cc|c}
    & y_1 & y_2 & \text{合计}\\\hline
    x_1 & a & b & a+b\\
    x_2 & c & d & c+d\\\hline
    \text{合计} & a+c & b+d & n
  \end{array}
\]
若两变量无关（零假设 $H_0$），应有 $ad\approx bc$。统计量
\[
  K^2=\frac{n(ad-bc)^2}{(a+b)(c+d)(a+c)(b+d)}
\]
度量实际数据偏离「无关」的程度。逻辑结构与第 0 章的反证法完全同构：先承认 $H_0$ 成立，则 $K^2$ 很大是\key{小概率事件}；它竟然发生了，就有理由拒绝 $H_0$。常用临界值：$K^2\ge3.841$ 对应犯错概率不超过 $5\%$，$K^2\ge6.635$ 对应不超过 $1\%$。

\begin{example}
调查 $200$ 人（男、女各 $100$ 人）对某项政策的支持情况：男性中 $70$ 人支持、$30$ 人反对；女性中 $50$ 人支持、$50$ 人反对。是否有 $99\%$ 以上的把握认为性别与支持态度有关？
\end{example}
\begin{solution}
\[
  K^2=\frac{200\times(70\times50-30\times50)^2}{100\times100\times120\times80}
  =\frac{200\times2000^2}{9.6\times10^7}
  =\frac{8\times10^8}{9.6\times10^7}\approx8.33 .
\]
因为 $8.33>6.635$，所以在犯错误的概率不超过 $0.01$ 的前提下，认为性别与支持态度有关——即有 $99\%$ 以上的把握。
\end{solution}

\begin{commonerror}{独立性检验的三个高频失误}
\paragraph{失误一：把结论说过头。}
正确表述是「在犯错误的概率不超过 $0.01$ 的前提下认为有关」或「有 $99\%$ 的把握认为有关」，
而不是「两变量有 $99\%$ 的关系」——$K^2$ 只衡量证据强度，不量化关联程度。
\paragraph{失误二：忽略样本量 $n$。}
公式分子含 $n$，同样比例的列联表，样本量扩大 $10$ 倍，$K^2$ 也扩大 $10$ 倍，结论可能完全不同；小样本下比例再悬殊也不宜下强结论。
\paragraph{失误三：把检验结论当成因果关系。}
「有关」不等于「导致」：独立性检验只能发现关联，因果方向须由实际背景另行论证。
\end{commonerror}

\begin{conccard}[一级]{最小二乘回归与独立性检验}{成对数据的回归计算、相关性判断、列联表分析}{把回归直线当成必过原点的比例函数；以为 $K^2$ 大就证明因果关系；计算 $\hat b$ 时忘除以 $\sum(x_i-\bar x)^2$；百分位数忽略 $i$ 是否为整数}
$\hat b=\frac{\sum(x_i-\bar x)(y_i-\bar y)}{\sum(x_i-\bar x)^2}$，$\hat a=\bar y-\hat b\bar x$，必过 $(\bar x,\bar y)$；
百分位数看 $i=n\cdot p\%$ 是否整数；
$K^2=\frac{n(ad-bc)^2}{(a+b)(c+d)(a+c)(b+d)}$，对照 $3.841$（$0.05$）、$6.635$（$0.01$）下结论。
\end{conccard}

\subsection{压轴收官：新定义情境题的四步破题法}

\key{探究四：新定义题不是考知识，是考「翻译」}

新高考压轴常以全新符号、新运算、新集合开场，陌生感是纸老虎。通用四步法：
\key{第一步，读懂符号集合}——把每个新名词逐字翻译为已有的集合语言（元素是什么？满足什么条件？运算封闭吗？）；
\key{第二步，构造小规模特例试探}——取 $n=1,2,3$ 等最小情形，逐一列举，列出数表；
\key{第三步，归纳猜想通项}——从特例中找规律、猜公式（归纳只指方向，不能当证明）；
\key{第四步，严格证明}——用双射（一一对应）、数学归纳法或反证法把猜想钉死。

\begin{method}[新定义情境题的四步破题流程]
\begin{enumerate}
  \item \key{翻译}：逐句把新定义的符号「转译」成集合、映射、计数或概率语言，画出韦恩图或数表；
  \item \key{特例}：取最小的参数（$n=1,2,3$ 或最小非平凡情形）把结果全部列出，观察结构；
  \item \key{归纳}：由特例数据猜想通项／公式／计数法则，并明确写出猜想；
  \item \key{证明}：优先寻找双射（把新对象一一对应到已知模型），否则用数学归纳法（与自然数相关）或反证法（唯一性、不存在性）。
\end{enumerate}
\end{method}

\begin{problem}
\key{压轴题一（新定义 $+$ 计数）}设 $n\ge k\ge1$。称 $\{1,2,\dots,n\}$ 的 $k$ 元子集 $A=\{a_1,a_2,\dots,a_k\}$（$a_1<a_2<\dots<a_k$）为「稀疏集」，若其中任意两个不同元素之差的绝对值不小于 $2$（即不含相邻整数）。记稀疏集的个数为 $f(n,k)$。
（1）列举 $n=5,\ k=2$ 的全部稀疏集，并求 $f(5,2)$；再求 $f(5,3)$；
（2）依（1）猜想 $f(n,k)$ 的表达式，并给出严格证明。
\end{problem}
\begin{solution}
\key{【第一步：读懂符号集合】}「稀疏集」＝$\{1,\dots,n\}$ 中不含相邻数的 $k$ 元子集，$f(n,k)$ 是这类子集的个数。
\key{【第二步：小规模特例】}（1）$n=5,k=2$：枚举得
\[
  \{1,3\},\{1,4\},\{1,5\},\{2,4\},\{2,5\},\{3,5\},
\]
共 $6$ 个，即 $f(5,2)=6=C_4^2$。
$n=5,k=3$：只能取 $\{1,3,5\}$，故 $f(5,3)=1=C_3^3$。
再试 $f(4,2)=3=C_3^2$、$f(6,2)=10=C_5^2$。
\key{【第三步：归纳猜想】}各值都等于 $C_{\,n-k+1}^{\,k}$，故猜想
\[
  f(n,k)=C_{\,n-k+1}^{\,k}.
\]
\key{【第四步：严格证明（双射法）】}设稀疏集 $A=\{a_1<\dots<a_k\}$，令
\[
  b_i=a_i-(i-1)\quad(i=1,\dots,k).
\]
由 $a_{i+1}-a_i\ge2$ 得 $b_{i+1}-b_i=(a_{i+1}-a_i)-1\ge1$，即 $b_1<b_2<\dots<b_k$；
又 $b_1=a_1\ge1$，$b_k=a_k-(k-1)\le n-k+1$，故 $\{b_1,\dots,b_k\}$ 是 $\{1,\dots,n-k+1\}$ 的任意 $k$ 元子集。
反之，任取 $1\le b_1<\dots<b_k\le n-k+1$，令 $a_i=b_i+(i-1)$，则 $a_{i+1}-a_i=b_{i+1}-b_i+1\ge2$，回到一个稀疏集。
两个方向互为逆映射，构成\key{一一对应}，故
\[
  f(n,k)=C_{\,n-k+1}^{\,k}.\ \surd
\]
\key{这一步同时解释了「为什么恰好是隔板模型」：减去的 $(i-1)$ 相当于预先在每个选中元素后面放好一个「隔离空位」，隔板数从 $n-1$ 缩为 $n-k+1$。}
\end{solution}

\begin{problem}
\key{压轴题二（新定义 $+$ 概率统计）}一批共 $6$ 件产品中有 $2$ 件次品，现逐个检查（不放回）。定义「转正等待次数」$X$ 为：首次抽到正品时已经检查过的产品件数。
（1）求 $X$ 的分布列与数学期望 $E(X)$；
（2）若改为有放回抽取，$X$ 服从什么分布？比较两种情形的期望并解释差异原因。
\end{problem}
\begin{solution}
\key{【第一步：翻译】}$X$ 是取值于 $\{1,2,3\}$ 的离散随机变量（最多查 3 件必出正品，因为只有 2 件次品）。
\key{【第二步：小特例手算】}
\[
  P(X=1)=\frac{4}{6}=\frac{2}{3};
\]
\[
  P(X=2)=P(\text{第 1 件次品})\cdot P(\text{第 2 件正品}\mid\text{第 1 件次品})
  =\frac{2}{6}\times\frac{4}{5}=\frac{4}{15};
\]
\[
  P(X=3)=P(\text{前两件均为次品})=\frac{2}{6}\times\frac{1}{5}=\frac{1}{15}.
\]
\key{【第三步：归纳校验】}$\frac23+\frac{4}{15}+\frac{1}{15}=1$，分布列为
\[
  \begin{array}{c|ccc}
    X & 1 & 2 & 3\\\hline
    P & \dfrac{2}{3} & \dfrac{4}{15} & \dfrac{1}{15}
  \end{array}
\]
\key{【第四步：期望与对比】}
\[
  E(X)=1\times\frac{2}{3}+2\times\frac{4}{15}+3\times\frac{1}{15}
  =\frac{10}{15}+\frac{8}{15}+\frac{3}{15}=\frac{21}{15}=\frac{7}{5}.
\]
（2）有放回时每次命中正品概率恒为 $\frac{4}{6}=\frac23$，$X$ 服从参数 $p=\frac23$ 的几何分布：
\[
  P(X=k)=\left(\frac{1}{3}\right)^{k-1}\cdot\frac{2}{3}\ (k\ge1),\qquad E(X)=\frac{1}{p}=\frac{3}{2}.
\]
$\frac{3}{2}>\frac{7}{5}$：不放回时次品一旦被抽出，剩下产品中正品浓度上升，等待时间自然缩短——这正是条件概率「空间缩放」在期望上的体现。
\end{solution}

\begin{conccard}[二级]{新定义情境题四步破题法}{强基、新高考压轴中「新运算、新集合、新统计量」类题目}{把归纳猜想的结论直接当证明使用；跳过特例直接硬算；翻译时漏掉隐含约束（如 $n\ge k$）}
第一步 翻译：逐句转译为集合／映射语言，画图表；
第二步 特例：$n=1,2,3$ 全部列举；
第三步 归纳：写出猜想公式；
第四步 证明：双射（计数）、数学归纳法（自然数命题）、反证法（唯一性与不存在）。
\end{conccard}

\begin{summary}
本章把「不确定性」彻底纳入第 0 章的集合框架：样本空间 $\Omega$ 是集合，概率是归一化计数测度 $P(A)=\mu(A)/\mu(\Omega)$；
计数问题的本质是并集基数（分类加法）与笛卡尔积基数（分步乘法），排列与组合分别是「有序单射」与「无序子集族」；
条件概率 $P(B\mid A)$ 是把样本全集缩小为 $A$ 后的空间缩放，全概率与贝叶斯公式源于样本空间的完备互斥划分（由因及果与由果溯因）；
线性回归直线是散点云的竖直最小二乘投影，必过样本中心 $(\bar x,\bar y)$；独立性检验则是「小概率事件原理」下的统计反证法；
压轴新定义题按「翻译符号 $\to$ 特例试探 $\to$ 归纳猜想 $\to$ 严格证明」四步破题。

\begin{center}
\footnotesize
\setlength{\tabcolsep}{4pt}
\renewcommand{\arraystretch}{1.35}
\begin{tabular}{p{2.5cm}p{4.7cm}p{4.5cm}p{1.4cm}}
\toprule
引擎部件 & 核心操作 & 把什么变成什么 & 对应探究 \\
\midrule
计数两原理 & $|A\sqcup B|=|A|+|B|$；$|A\times B|=|A|\,|B|$ & 分类（或）$\to$ 并；分步（且）$\to$ 积 & 探究一 \\
排列与组合 & $A_n^k=\frac{n!}{(n-k)!}$；$C_n^k=\frac{A_n^k}{k!}$ & 有序单射序列 $\to$ 无序 $k$ 元子集族 & 探究一 \\
三技法 & 捆绑／插空／隔板 $C_{n-1}^{m-1}$ & 相邻、不相邻、相同分堆 $\to$ 标准模型 & 探究一 \\
归一化测度 & $P(A)=\mu(A)/\mu(\Omega)$ & 古典与几何 $\to$ 同一除法 & 探究二 \\
条件概率 & $P(B\mid A)=\frac{P(A\cap B)}{P(A)}$ & 样本全集 $\Omega\to$ 子集 $A$：缩放归一 & 探究二 \\
全概率·贝叶斯 & $P(B)=\sum_iP(A_i)P(B\mid A_i)$ & 由因及果 $\leftrightarrow$ 由果溯因 & 探究二 \\
最小二乘回归 & $\min\sum(y_i-\hat bx_i-\hat a)^2$，过 $(\bar x,\bar y)$ & 散点云 $\to$ 竖直投影直线 & 探究三 \\
独立性检验 & $K^2$ 对临界值 & 关联判断 $\to$ 小概率反证 & 探究三 \\
四步破题 & 翻译 $\to$ 特例 $\to$ 归纳 $\to$ 严格证 & 陌生符号 $\to$ 熟悉结构 & 探究四 \\
\bottomrule
\end{tabular}
\end{center}

\key{随机性不是结构的对立面，而是有限集合上的测度；统计不是死记公式，而是「最小二乘投影 $+$ 小概率反证」的合体；新定义压轴的本质是把陌生符号翻译回熟悉的集合与结构。}
\end{summary}

\end{explore}

\begin{inclasstest}
\begin{enumerate}
  \item 从 6 名同学中选出 3 人，分别担任语文、数学、英语三科课代表，共有多少种不同的选法？\level{易}
  \item 把 10 个相同的小球放入 3 个不同的盒子，每个盒子至少放 1 个，共有多少种放法？\level{易}
  \item 口袋中装有 5 个红球、3 个白球，不放回地取 3 个球。已知取出的第一个球是红球，求取出的第三个球也是红球的概率。\level{中}
  \item 已知某组成对数据的样本中心为 $(\bar x,\bar y)=(3,5)$，且 $\sum(x_i-\bar x)(y_i-\bar y)=12$，$\sum(x_i-\bar x)^2=4$。求回归直线方程，并指出它必过的定点。\level{难}
\end{enumerate}
\end{inclasstest}

\begin{afterclass}
\subsection*{基础巩固}
\begin{enumerate}
  \item 6 人站成一排，甲、乙、丙三人必须站在一起，有多少种站法？
  \item 袋中有 3 个红球、2 个白球，不放回地依次取出 2 个球，求第二次取到红球的概率。
\end{enumerate}

\subsection*{进阶提升}
\begin{enumerate}
  \item 甲、乙两人射击的命中率分别为 $0.8$ 与 $0.7$，各射击一次。已知恰有一人命中，求命中者是甲的概率。
  \item 将 5 个不同的小球放入 3 个不同的盒子，每个盒子至少放 1 个，有多少种放法？
\end{enumerate}

\subsection*{拓展探究}
\begin{enumerate}
  \item （强基／竞赛）设 $n\ge k\ge1$。称 $\{1,2,\dots,n\}$ 的 $k$ 元子集 $A$ 为「强稀疏集」，若其中任意两个不同元素之差至少为 $3$。记其个数为 $g(n,k)$。
  （1）列出 $n=6,k=2$ 的全部强稀疏集，求 $g(6,2)$；
  （2）猜想 $g(n,k)$ 的表达式并用双射法证明；
  （3）比较 $g(n,k)$ 与本章 $f(n,k)$ 的大小，说明理由。
\end{enumerate}

\subsection*{教师版答案}
\begin{solution}
\textbf{基础 1：}把甲、乙、丙三人捆成一个大元素，与其余 3 人共 4 个元素全排列，三人内部再全排列：
\[
  A_4^4\cdot A_3^3=24\times6=144 .
\]
\end{solution}
\begin{solution}
\textbf{基础 2：}方法一（全概率）：
\[
  P=\frac35\times\frac24+\frac25\times\frac34=\frac{6}{20}+\frac{6}{20}=\frac{3}{5}.
\]
方法二（抽签公平性）：先后手对称，第二次取到红球的概率等于第一次取到红球的概率 $\frac35$。
\end{solution}
\begin{solution}
\textbf{进阶 1：}记 $A$ 为「甲命中」，$B$ 为「乙命中」，$A,B$ 独立。
\[
  P(\text{恰有一人命中})=P(A\bar B)+P(\bar A B)=0.8\times0.3+0.2\times0.7=0.38 .
\]
由条件概率（把样本空间缩小为「恰有一人命中」）：
\[
  P(A\mid \text{恰有一人命中})=\frac{0.24}{0.38}=\frac{12}{19}\approx0.632 .
\]
\end{solution}
\begin{solution}
\textbf{进阶 2：}由容斥原理（全集减去有空盒的放法）：
\[
  3^5-C_3^1\cdot2^5+C_3^2\cdot1^5=243-96+3=150 .
\]
\end{solution}
\begin{solution}
\textbf{拓展：}（1）$n=6,k=2$：枚举得
\[
  \{1,4\},\{1,5\},\{1,6\},\{2,5\},\{2,6\},\{3,6\},
\]
共 $6$ 个，即 $g(6,2)=6$。
（2）猜想：$g(n,k)=C_{\,n-2k+2}^{\,k}$（约定 $m<k$ 时 $C_m^k=0$，对应不可能情形）。
双射证明：设强稀疏集 $A=\{a_1<\dots<a_k\}$，$a_{i+1}-a_i\ge3$。令
\[
  b_i=a_i-2(i-1)\quad(i=1,\dots,k),
\]
则 $b_{i+1}-b_i=(a_{i+1}-a_i)-2\ge1$，即 $b_1<\dots<b_k$；且 $b_1\ge1$，$b_k=a_k-2(k-1)\le n-2k+2$，
故 $\{b_1,\dots,b_k\}$ 是 $\{1,\dots,n-2k+2\}$ 的任意 $k$ 元子集；反之 $a_i=b_i+2(i-1)$ 给出逆映射。一一对应，故 $g(n,k)=C_{\,n-2k+2}^{\,k}$。
（3）「强稀疏集」一定是「稀疏集」，反之不然（如 $\{1,3\}$ 稀疏但不强稀疏），前者集合族是后者的真子族，故除 $k=1$ 等平凡情形外恒有
\[
  g(n,k)<f(n,k).
\]
\end{solution}

\reading{拓展阅读 13}{概率：从计数测度到样本空间}
\end{afterclass}

% ==========================================================
% 第 8 章 齐次化：把二元问题压成一元比例
% 说明：本文件为纯正文片段，由 main_body.tex 挂载；
%       使用 versions.tex 的 \tk / \vchoice / advancedonly /
%       basiconly / solution / 四个盒子。
% ==========================================================

\section{齐次化：把二元问题压成一元比例（高观点结构版）}

\begin{methodbox}{多框架对读}
\small
\begin{center}
\begin{tabular}{c|p{2.9cm}|p{5.0cm}|p{3.1cm}}
\toprule
视角 & 本章看点 & 别的视角怎么说 & 哪侧最省力 \\
\midrule
代数 & 同除配凑、$1$ 的代换、换元 & 符号的恒等变形 & 硬算最值 \\
几何 & 比值 $t=\frac yx$＝射线斜率 & 点 $(x,y)$ 落在一条过原点射线上 & 看比例、判范围 \\
集合 & 齐次式按次数分层 & 同次＝缩放不变的性质 & 界次数与符号 \\
映射 & 齐次坐标 $\to$ 射影直线 & $(u,v)\sim\lambda(u,v)$，$x=u/v$ & 升到第 9 章接口 \\
\bottomrule
\end{tabular}
\end{center}
卡的使用：二元分式最值先数次数、再定「$1$ 的代换还是 $t=\frac yx$」，最后回代数收尾（总图见第 0 章图~\ref{fig:four-view-map}）。
\end{methodbox}

\begin{warmupbox}{读法（先看这里）}
齐次化不是「大招」，而是一句结构主义的第一性原理：\textbf{把
$n$ 个自由变量的绝对尺度约掉，只留下 $n-1$ 个比例。}
遇到二元（$x,y$）的对称式，先问一句：它们能不能合并成一个
分数 $t=\dfrac{y}{x}$？能，就降维；不能，就先配平次数。

（本章与第 5、6 章互为表里：第 6 章「条件最值四大标准动作」给出
「1 的代换」这一入口，第 5 章用矩阵特征方向处理分式递推；本章把
两者背后的原理——齐次化——上升为第一性原理并展开成完整方法链。）
\end{warmupbox}

\begin{corebox}{第一性原理}
设 $P(x,y)=0$ 是一条曲线方程，$Q(x,y)$ 是待求的目标式。若存在
正数 $\lambda$ 使得 $Q(\lambda x,\lambda y)=\lambda^{d}\,Q(x,y)$，
则 $Q$ 是 $d$ 次齐次的，它的取值只依赖于比例 $t=y/x$，与点的
绝对位置无关。非齐次的问题之所以难，是因为混进了不同次数的
项——\textbf{齐次化就是把这些项强行配到同一个次数上}。
\end{corebox}

\begin{methodbox}{次数扫描雷达（动手前先扫一遍）}
给每一项标一个「次数」，像雷达一样先扫描整个式子：
\begin{center}
\begin{tikzpicture}[scale=0.9, font=\small]
\draw[->] (-0.4,0)--(9.6,0) node[right] {$x$ 轴：项};
\draw[->] (0,-0.5)--(0,3.4);
\foreach \x/\h/\lab/\deg [count=\i] in {
  1.4/0.5/常数项/0,
  3.2/1.0/一次项/1,
  5.0/1.5/二次项/2,
  6.8/2.0/分式项/-1,
  8.6/2.5/根式项/{1/2}} {
  \draw[main, thick] (\x,0)--(\x,\h);
  \node[above] at (\x,\h) {\lab};
  \node[below] at (\x,-0.18) {$\deg$ 次};
}
\node[below] at (5.0,-0.66) {（柱高仅为扫描顺序，次数见下方标签）};
\end{tikzpicture}
\end{center}
（常数 $0$ 次、$x$ 与 $y$ 各 $1$ 次、$x^2/xy/y^2$ 为 $2$ 次、
$\dfrac1x$ 为 $-1$ 次、$\sqrt{x}$ 为 $\tfrac12$ 次。）
\textbf{判据}：一个式子是齐次的，当且仅当所有项次数相同。
\end{methodbox}

\begin{methodbox}{方法归纳：同除主元最高次幂（降维的钥匙）}\label{met:homo-divide}
若已经确认各项同次，就同除「主元」的最高次幂，把 $n$ 个变量
压成比例。例如二元齐次式同除 $x^2$（同除 $y^2$ 亦可，此时
比例换成 $x/y=1/t$），只剩
$t=\dfrac{y}{x}$：
\[
\frac{x^2+2xy+y^2}{x^2}=1+2t+t^2 .
\]
一元化之后，剩下的就是单变量函数的最值问题——\textbf{这就是
齐次化的全部收益：把二元压成一元}。
\end{methodbox}

\subsection{第 1 节 常规通法与齐次化的第一次对照}

\begin{example}[已知 $x>0,\ y>0,\ 2x+y=1$，求 $xy$ 的最大值]
\vspace{-1mm}
\begin{solution}
\textbf{考纲通法}：$y=1-2x$，故 $xy=x(1-2x)=-2x^2+x$，
开口向下，顶点在 $x=\dfrac14$，最大值为 $\tk{\dfrac18}$。

\textbf{齐次化视角}：目标与约束同时齐次化，把分母补成
\[
\frac{xy}{(2x+y)^2}=\frac{xy}{1}.
\]
令 $t=\dfrac{y}{x}>0$，则上式右端化为
\[
\frac{t}{(2+t)^2}.
\]
由 AM–GM（$(t+2)^2\ge 8t$），有 $\dfrac{t}{(2+t)^2}\le\dfrac18$，
等号在 $t=2$ 即 $y=2x$、$x=\dfrac14,\ y=\dfrac12$ 时成立。
两种算法给出同一个答案 $\tk{\dfrac18}$。
\end{solution}
\end{example}

\begin{conccard}{齐次化入口：把约束等式当 1 乘进目标式}{二元对称式的最值、范围题}{只会代入消元，不会「补次数」}
约束条件「$2x+y=1$」提供了 \textbf{一个} 齐次化入口：任何
目标式都可以乘上 $(2x+y)^k$ 而不改变取值，从而把混次的式子
全部抬到同一次数。\vchoice{乘方指数 $k$ 的选取原则是「让所有项次数相等」。}{看到等式条件，就想想能不能拿它当 1 来乘。}
\end{conccard}

\subsection{第 2 节 五大配平手术刀}

\begin{methodbox}{手术刀之一「1 的代换」升次}
条件等式本身等于 $1$（或可化成 $1$），把它乘进目标式，
低次项立刻被抬成齐次式。
\[
\left(\frac1x+\frac1y\right)\cdot 1
=\left(\frac1x+\frac1y\right)(x+y)
=2+\frac{y}{x}+\frac{x}{y}\ge 4 .
\]
\end{methodbox}

\begin{methodbox}{手术刀之二 平移原点，把方程归一}
把原点移到图形的中心（或对称中心），一次项/常数项会被
消掉，方程立刻变成标准型。
\[
x^2+y^2-2x=0\ \xrightarrow{\ X=x-1,\ Y=y\ }\ X^2+Y^2=1,
\]
于是 $x+y=X+Y+1\in[\tk{1-\sqrt2},\ \tk{1+\sqrt2}]$。
\end{methodbox}

\begin{methodbox}{手术刀之三 平方升次，剥掉根号}
$\sqrt{x}$ 的次数是 $\tfrac12$，平方一次就升到整次。
已知 $\sqrt{x}+\sqrt{y}=1$，则
\[
x+y=(\sqrt{x}+\sqrt{y})^2-2\sqrt{xy}=1-2\sqrt{xy},
\]
再用 $\sqrt{x}\cdot\sqrt{y}\le\left(\dfrac{\sqrt{x}+\sqrt{y}}{2}\right)^2=\dfrac14$
得 $x+y\ge\tk{\dfrac12}$。
\end{methodbox}

\begin{methodbox}{手术刀之四 分解与组合降次}
高次式先分解，往往能拆出可以约分的因式。
\[
\frac{1}{x(x+1)}=\frac1x-\frac1{x+1},\qquad
x^3+y^3=(x+y)(x^2-xy+y^2).
\]
\end{methodbox}

\begin{methodbox}{手术刀之五 旋转消交叉项}
二次曲线
\[
Ax^2+Bxy+Cy^2+\cdots=0
\]
的交叉项 $Bxy$ 用平移消不掉，换一招：把坐标轴旋转角度 $\theta$，
\[
x=x'\cos\theta-y'\sin\theta,\qquad y=x'\sin\theta+y'\cos\theta .
\]
代回后令 $x'y'$ 的系数为 $0$，即得消交叉项的条件
\[
\cot 2\theta=\frac{A-C}{B}
\qquad\text{（等价地 }\tan 2\theta=\frac{B}{A-C}\text{；当 }A=C\text{ 时用前者）}.
\]
特例 $xy=1$：这里 $A=C=0,\ B=1$，故 $\cot 2\theta=0$，取 $\theta=45^\circ$，
\[
xy=1\ \xrightarrow{\ \theta=45^\circ\ }\ \frac{x'^2}{2}-\frac{y'^2}{2}=1,
\]
等轴双曲线还原为标准型（$a=b=\sqrt2$）。
\textbf{完整推导、图与数列中的复特征值应用见第 5 章。}
\end{methodbox}

\begin{center}
\begin{tikzpicture}[scale=0.85, font=\small]
% 椭圆平移齐次化示意：从原点出发的斜率射线
\draw[gray!60, thick] (0,0) ellipse (3.2 and 2.0);
\draw[->, main, very thick] (0,0)--(3.6,2.4);
\draw[dashed, third] (0,0)--(3.6,-1.2);
\node[right] at (3.65,2.4) {$k=\dfrac{y}{x}>0$};
\node[right] at (3.65,-1.2) {$k<0$};
\fill[second] (0,0) circle (2pt) node[below left] {新原点};
\node at (0,-2.5) {$k=\dfrac{y}{x}$ 的符号决定直线落在哪个象限};
\end{tikzpicture}
\end{center}

\begin{conccard}{五把配平手术刀：升次、平移、平方、分解、旋转}{混次式、根式、分式迭代、含交叉项的二次曲线}{只背大招名称，不动手配平次数}
五把手术刀的共同目标是把方程「摆正」：前四把在\textbf{次数}上
配平，第五把在\textbf{坐标方向}上配平，消掉交叉项 $Bxy$。
判断该用哪一把，只需回看次数扫描的结果——缺几次就补几次，
多几次就降几次；遇到交叉项 $Bxy$ 就转坐标轴（第 5 章有完整推导），
直到方程露出标准型。
\vchoice{严格地说，配平后问题就化为关于 $t=y/x$ 的一元函数，导数或基本不等式都可直接上手。}{口诀：次数不齐先配平，同除主元变一元；方向不正先旋转。}
\end{conccard}

\begin{advancedonly}
\subsection{第 3 节 射影桥梁：为什么齐次化天然有效}

第 5 章探究三已经证明：分式迭代（莫比乌斯变换）
\[
x_{n+1}=\frac{a\,x_n+b}{c\,x_n+d}\qquad(ad-bc\neq0)
\]
的不动点恰是矩阵 $\begin{pmatrix}a&b\\c&d\end{pmatrix}$ 的特征方向，
特征值之比就是双不动点构造的公比（推导与图见第 5 章
图~\ref{fig:eigen-fixed}；不动点的一般理论与双不动点的几何图见第 9 章
图~\ref{fig:mobius-two-fixed}）。这里只补上齐次化的动机：

\textbf{齐次式天生属于射影。} 把 $(u,v)$ 与 $\lambda(u,v)$ 视为同一个
射影点，比值 $x=u/v$ 就是它的坐标；分子、分母同次的式子在这条射线
上取确定值。于是「找不动点」在射影语言里就是「求特征方向」——
齐次化把非线性的分式迭代重新线性化了。

\begin{example}[设 $x_{n+1}=\dfrac{2x_n+1}{x_n+2}$，$x_0=1$，求 $\displaystyle\lim_{n\to\infty}x_n$]
\vspace{-1mm}
\begin{solution}
矩阵 $M=\begin{pmatrix}2&1\\1&2\end{pmatrix}$ 的特征值为
$\lambda_1=3,\lambda_2=1$，特征向量分别沿
$\binom{1}{1}$ 与 $\binom{1}{-1}$，对应不动点 $x=1$ 与 $x=-1$。
$x_0=1$ 恰在 $x=1$ 处，故极限为 $\tk{1}$。
\end{solution}
\end{example}
\end{advancedonly}

\begin{basiconly}
\subsection{第 3 节 大白话版：比例不变，大小随意}

把 $(x,y)$ 想成一条过原点的射线上的点。把 $x,y$ 同时放大
$k$ 倍，点在射线上滑动，但「谁是谁的几倍」没变。齐次式
只认这条射线，不认点在射线上的位置——所以它可以只用一个
字母 $t$ 说话。

\textbf{照样子填空}：$\dfrac{xy}{(2x+y)^2}$ 的分子分母同除
$x^2$ 后得到 $\dfrac{\underline{\quad t\quad}}{(2+\underline{\quad t\quad})^2}$。
\end{basiconly}

\subsection{第 4 节 跨章实战：五题对照（含三道双解）}

\begin{problem}[闪电战 1] 已知 $x>0,\ y>0,\ x+2y=1$，求 $x^2+y^2$ 的最小值。\end{problem}
\begin{solution}
齐次化：$(x^2+y^2)(x+2y)^2=x^2+y^2$。令 $t=\dfrac{y}{x}$，
左边化为 $\dfrac{1+t^2}{(1+2t)^2}$。由柯西不等式
$(1+2t)^2\le(1^2+2^2)(1+t^2)$，得
$\dfrac{1+t^2}{(1+2t)^2}\ge\dfrac15$，等号在 $t=2$
即 $y=2x$、$x=\dfrac15,\ y=\dfrac25$ 处取得。最小值为 $\tk{\dfrac15}$。
\end{solution}

\begin{problem}[闪电战 2] 已知 $x^2+y^2-2x-4y=0$，求 $x+y$ 的取值范围。\end{problem}
\begin{solution}
配方：$(x-1)^2+(y-2)^2=5$，令
$X=x-1,\ Y=y-2$，则 $X^2+Y^2=5$，
$x+y=X+Y+3\in[3-\sqrt{10},\ 3+\sqrt{10}]$。
\end{solution}

\begin{problem}[双解对照 1] 已知 $x>0,\ y>0,\ \dfrac1x+\dfrac2y=1$，求 $x+y$ 的最小值。\end{problem}
\begin{solution}
\textbf{通法}：$x+y=(x+y)\left(\dfrac1x+\dfrac2y\right)
=3+\dfrac{y}{x}+\dfrac{2x}{y}\ge3+2\sqrt2$。

\textbf{齐次化}：令 $t=\dfrac{y}{x}$，约束化为
$\dfrac1x+\dfrac{2}{tx}=1\Rightarrow x=1+\dfrac2t$，
于是 $x+y=x(1+t)=\left(1+\dfrac2t\right)(1+t)
=1+t+\dfrac2t+2\ge3+2\sqrt2$。两种写法同解。
\end{solution}

\begin{problem}[双解对照 2] 曲线 $xy=1$ 与直线 $y=kx+2$ 有交点，求 $k$ 的范围。\end{problem}
\begin{solution}
代入得 $kx^2+2x-1=0$。有实根要求判别式
$\Delta=4+4k\ge0\Rightarrow k\ge-1$；$k=0$ 时方程为
$2x-1=0$，对应交点 $\left(\dfrac12,2\right)$，仍在曲线上。
故 $k\in[-1,+\infty)$，最小值在 $k=-1$ 处取到（相切）。
\end{solution}

\begin{problem}[双解对照 3] 已知 $\sqrt{x}+\sqrt{y}=2$，求 $x+y$ 的最小值与最大值。\end{problem}
\begin{solution}
平方：$x+y=4-2\sqrt{xy}$。由
$\sqrt{xy}\le\dfrac{\sqrt{x}+\sqrt{y}}{2}=1$ 得 $x+y\ge2$；
又 $x+y=(\sqrt{x}+\sqrt{y})^2-2\sqrt{xy}\le4$（当一项为 $0$ 时取等）。
故最小值 $\tk{2}$，最大值 $\tk{4}$。
\end{solution}

\begin{inclasstest}{当堂检测（4 题）}
\begin{enumerate}[label=\arabic*., itemsep=2pt]
\item 次数扫描：$\dfrac{x^2}{y}$ 是\underline{\hspace{2em}}次齐次式。
\item 已知 $x>0,\ y>0,\ x+y=1$，则 $xy$ 的最大值是\underline{\hspace{2em}}。
\item 曲线 $x^2+y^2-6x=0$ 平移后的标准方程是\underline{\hspace{2em}}。
\item 已知 $x>0,\ y>0$ 且 $\dfrac1x+\dfrac1y=1$，则 $x+y$ 的最小值是\underline{\hspace{2em}}。
\end{enumerate}
\begin{solution}
\tk{\text{一}}；\tk{\dfrac14}；\tk{(x-3)^2+y^2=9}；\tk{4}。
\end{solution}
\end{inclasstest}

\begin{homework}{分层作业}
\textbf{基础层（步步填空）}：已知 $x>0,\ y>0,\ x+3y=1$，则
$\dfrac1x+\dfrac3y=1+\underline{\quad}+\dfrac{3y}{x}+\underline{\quad}
\ge\underline{\quad}$。

\textbf{提高层（双解对照）}：已知 $x>0,\ y>0,\ 2x+3y=1$，
求 $\dfrac{1}{x}+\dfrac{1}{y}$ 的最小值（写出通法与齐次化两种）。

\begin{solution}
基础层：$\tk{\dfrac{3x}{y}}$，$\tk{9}$，$\tk{16}$。\\
提高层：$\left(\dfrac1x+\dfrac1y\right)(2x+3y)
=5+\dfrac{3y}{x}+\dfrac{2x}{y}\ge5+2\sqrt6$，最小值
$\tk{5+2\sqrt6}$。
\end{solution}
\end{homework}

\begin{commonerror}{易错点：把「同除」用反了方向}
常见的错误是把 $\dfrac{xy}{(2x+y)^2}$ 齐次化时只除一次，
错写成 $\dfrac{y/x}{2+y/x}$——\textbf{分母少除了一个 $x$}。
正确写法是分子分母\textbf{同除同一个量} $x^2$：
\[
  \dfrac{xy}{(2x+y)^2}=\dfrac{y/x}{(2+y/x)^2}.
\]
次数差才保持不变；除完应当立刻检查「分子还剩几次、分母还剩几次」。
\end{commonerror}

\begin{conccard}[二级]{什么时候必须齐次化，什么时候不必}{条件最值、曲线与直线的位置关系、根式最值}{不管题目长什么样都硬凑 $t=y/x$；或只会代入消元，从不配次数}
先看次数结构：目标式与约束能配成同次（或经平方、平移后可配平）时，
齐次化一步降为一元 $t$；消元后是一元二次、只问有无实根时，判别式
更直接；含 $\sqrt{\phantom{x}}$ 的半次项先平方升次，再决定走哪条路。
\end{conccard}

% ==========================================================
%  第 5 节 七问连环母题（解答即笔记 · Scaffolded Master Problem）
% ==========================================================

\subsection{第 5 节 七问连环母题：解答即笔记}

\begin{warmupbox}{母题怎么用（先看这里）}
本节两道大题，每道挂载七个递进小问。\vchoice{它们不是练习题，
而是「会说话的笔记」：第 (1)–(3) 问替代概念与规范书写笔记，
第 (4)–(6) 问替代大招、公式族与充要边界笔记，第 (7) 问替代
错题本。做完一道，这一章的框架就长在你的解题纸上。}{每道题顺着
七步走即可：先算具体数值，再翻译结构，然后规范推导、降维秒杀、
推广公式、逆向追问，最后排雷。}
\end{warmupbox}

\refstepcounter{hvmasterq}
\begin{masterq}{线性约束 $x+y=1$：从 $\dfrac1x+\dfrac1y$ 长出一族不等式}
设 $x>0,\ y>0,\ x+y=1$。
\begin{enumerate}[label=(\arabic*), itemsep=3pt]
\item \textbf{画图与数值感知}：计算 $(x,y)=\left(\dfrac14,\dfrac34\right)$、
$\left(\dfrac12,\dfrac12\right)$、$\left(\dfrac34,\dfrac14\right)$ 处
$\dfrac1x+\dfrac1y$ 的值，描点并补出草图。\logotag{核心公理}
\item \textbf{多视角结构翻译}：用集合语言写出点集；指出它在坐标系中的
几何形状；再用恒等式把 $\dfrac1x+\dfrac1y$ 改写成只含 $xy$ 的式子。\logotag{结构翻译}
\item \textbf{考纲规范通法}：用基本不等式求 $\dfrac1x+\dfrac1y$ 的最小值，
写出取等条件与关键步骤（用 \tk{} 填空）。\logotag{考纲通法}
\item \textbf{高观点降维秒杀}：令 $t=\dfrac{y}{x}$，把 $\dfrac1x+\dfrac1y$
写成 $t$ 的函数并用一行不等式收束，与第 (3) 问核对。\logotag{大招公式}
\item \textbf{公式变形家族生成}：推广到 $x+y=s$、$ax+by=1$，并推导
$\dfrac1{x^2}+\dfrac1{y^2}$ 的下界。\logotag{公式族}
\item \textbf{参数变式与逆向追问}：已知 $\dfrac1x+\dfrac1y=5$ 且 $x+y=1$，
反求 $x,y$；再回答：使 $\dfrac1x+\dfrac1y\ge m$ 恒成立的 $m$ 最大是多少？\logotag{充要边界}
\item \textbf{边界陷阱与错题排雷}：某同学删去「$x,y>0$」后仍写
$\dfrac1x+\dfrac1y\ge4$。请举反例推翻它，并说明原题有没有最大值。\logotag{避坑铁律}
\end{enumerate}
\end{masterq}

\begin{solution}
(1) $\dfrac{16}{3}$，$\tk{4}$，$\dfrac{16}{3}$；图象关于 $x=\dfrac12$ 对称，
两端趋于 $+\infty$。

(2) $\{(x,y)\mid x>0,\ y>0,\ x+y=1\}$；第一象限内连接 $(1,0)$ 与
$(0,1)$ 的线段；$\dfrac1x+\dfrac1y=\dfrac{x+y}{xy}=\tk{\dfrac{1}{xy}}$。

(3) $xy\le\left(\dfrac{x+y}{2}\right)^2=\dfrac14$，故
$\dfrac1x+\dfrac1y=\dfrac1{xy}\ge\tk{4}$，当且仅当 $x=y=\dfrac12$。

(4) $x=\dfrac{1}{1+t},\ y=\dfrac{t}{1+t}$，于是
$\dfrac1x+\dfrac1y=(1+t)+\dfrac{1+t}{t}=2+t+\dfrac1t\ge\tk{4}$
（$t=1$ 取等），与 (3) 一致。

(5) (a) $x+y=s$：$\dfrac1x+\dfrac1y=\dfrac{s}{xy}\ge\dfrac{s}{s^2/4}
=\tk{\dfrac4s}$；
(b) $ax+by=1$：$\dfrac1x+\dfrac1y=\left(\dfrac1x+\dfrac1y\right)(ax+by)
=a+b+\dfrac{ay}{x}+\dfrac{bx}{y}\ge a+b+2\sqrt{ab}=\tk{(\sqrt a+\sqrt b)^2}$；
(c) $\dfrac1{x^2}+\dfrac1{y^2}=\left(\dfrac1x+\dfrac1y\right)^2-\dfrac{2}{xy}
\ge16-8=\tk{8}$。

(6) $xy=\dfrac15$，故 $x,y$ 是 $z^2-z+\dfrac15=0$ 的两根，即
$\tk{\dfrac{5\pm\sqrt5}{10}}$；$m_{\max}=\tk{4}$（恰为下确界）。

(7) 反例：$x=2,\ y=-1$ 满足 $x+y=1$，但 $\dfrac1x+\dfrac1y=-\dfrac12<4$。
原题中 $x\to0^+$ 时 $\dfrac1x+\dfrac1y\to+\infty$，故只有最小值、无最大值。
\end{solution}

\refstepcounter{hvmasterq}
\begin{masterq}{二次约束 $x^2+y^2=1$：齐次化与「1 的代换」双剑合璧}
设 $x^2+y^2=1\ (x>0,\ y>0)$。
\begin{enumerate}[label=(\arabic*), itemsep=3pt]
\item \textbf{画图与数值感知}：计算 $\left(\dfrac{\sqrt2}{2},\dfrac{\sqrt2}{2}\right)$、
$\left(\dfrac35,\dfrac45\right)$、$\left(\dfrac12,\dfrac{\sqrt3}{2}\right)$ 处
$\dfrac1{x^2}+\dfrac1{y^2}$ 与 $\dfrac{x}{y}+\dfrac{y}{x}$ 的值，并在单位圆
第一象限弧上描点。\logotag{核心公理}
\item \textbf{多视角结构翻译}：用集合 $\{(\cos\theta,\sin\theta)\mid
0<\theta<\tfrac\pi2\}$ 重写点集；指出几何形状；说明 $x^2+y^2=1$ 为什么
可以当作「1」乘进任何式子。\logotag{结构翻译}
\item \textbf{考纲规范通法}：求 $\dfrac1{x^2}+\dfrac1{y^2}$ 的最小值，
写出取等条件（用 \tk{} 填空）。\logotag{考纲通法}
\item \textbf{高观点降维秒杀}：把 $\dfrac1{x^2}+\dfrac1{y^2}$ 配成分子分母
同次（乘「1」再除「1」），令 $t=\dfrac{y}{x}$ 一行出结果，并顺手得到
$\dfrac{x}{y}+\dfrac{y}{x}$ 的下界。\logotag{大招公式}
\item \textbf{公式变形家族生成}：推导 $x^2+y^2=r^2$、$x^2+2y^2=1$ 情形下
$\dfrac1{x^2}+\dfrac1{y^2}$、$\dfrac1{x^2}+\dfrac2{y^2}$ 与
$\dfrac1{x^4}+\dfrac1{y^4}$ 的下界。\logotag{公式族}
\item \textbf{参数变式与逆向追问}：已知 $\dfrac1{x^2}+\dfrac1{y^2}=5$ 且
$x^2+y^2=1$，求 $x^2,y^2$；若去掉「$x,y>0$」只保留 $x^2+y^2=1$，
最小值是否改变？\logotag{充要边界}
\item \textbf{边界陷阱与错题排雷}：错解「$xy\le\dfrac12\Rightarrow
\dfrac1{xy}\ge2\Rightarrow\dfrac1{x^2y^2}\ge4$」的哪一步依赖了符号条件？
写出防错铁律。\logotag{避坑铁律}
\end{enumerate}
\end{masterq}

\begin{solution}
(1) $\left(\dfrac{\sqrt2}{2},\dfrac{\sqrt2}{2}\right)$：$\dfrac1{x^2}+\dfrac1{y^2}=4$，
$\dfrac{x}{y}+\dfrac{y}{x}=2$；
$\left(\dfrac35,\dfrac45\right)$：$\dfrac{25}{9}+\dfrac{25}{16}
=\tk{\dfrac{625}{144}}$，$\dfrac{x}{y}+\dfrac{y}{x}=\dfrac34+\dfrac43
=\tk{\dfrac{25}{12}}$；
$\left(\dfrac12,\dfrac{\sqrt3}{2}\right)$：$4+\dfrac43=\tk{\dfrac{16}{3}}$。

(2) 几何上是单位圆在第一象限的四分之一圆弧；代数上 $x^2+y^2=1$ 是
「1 的代换」的合法来源（乘以它不改变任何式子的取值）。

(3) $\dfrac1{x^2}+\dfrac1{y^2}=\dfrac{x^2+y^2}{x^2y^2}=\dfrac1{x^2y^2}$，
而 $xy\le\dfrac{x^2+y^2}{2}=\dfrac12\Rightarrow x^2y^2\le\dfrac14$，
故原式 $\ge\tk{4}$，等号当且仅当 $x=y=\dfrac{\sqrt2}{2}$。

(4) 乘「1」：$\dfrac1{x^2}+\dfrac1{y^2}
=\dfrac{(x^2+y^2)^2}{x^2y^2}$；令 $t=\dfrac{y}{x}$ 得
$\dfrac{(1+t^2)^2}{t^2}=\left(t+\dfrac1t\right)^2\ge4$；同一行顺手得到
$\dfrac{x}{y}+\dfrac{y}{x}=t+\dfrac1t\ge\tk{2}$。

(5) (a) $x^2+y^2=r^2\Rightarrow\dfrac1{x^2}+\dfrac1{y^2}
=\dfrac{r^4}{x^2y^2}\ge\tk{\dfrac{4}{r^2}}$；
(b) $\left(\dfrac1{x^2}+\dfrac2{y^2}\right)(x^2+2y^2)
=5+2\left(\dfrac{y^2}{x^2}+\dfrac{x^2}{y^2}\right)\ge5+4=9$，
故 $\dfrac1{x^2}+\dfrac2{y^2}\ge\tk{9}$（等号 $x^2=y^2=\dfrac13$）；
(c) $\dfrac1{x^4}+\dfrac1{y^4}\ge\dfrac{2}{x^2y^2}\ge\tk{8}$。

(6) 令 $u=x^2,\ v=y^2$，则 $u+v=1,\ \dfrac{u+v}{uv}=5\Rightarrow uv=\dfrac15$，
故 $u,v=\tk{\dfrac{5\pm\sqrt5}{10}}$；去掉正数限制后结论不变（只含平方项），
但仍需 $x,y\neq0$。

(7) 「$\dfrac1{xy}\ge2$」这一步要求 $xy>0$；若 $xy<0$，不等号方向相反。
本题最终正确只是因为又平方了一次。防错铁律：取倒数或用基本不等式前，
先确认各项同号。
\end{solution}

\begin{basiconly}
\subsection{第 6 节 大白话版：两道母题的半成品脚手架}
母题一：两人分一个蛋糕，$x,y$ 是各拿多少。怎么分，倒数的和最小？
先试四等分点：$x=\dfrac14$ 时 $\dfrac1x+\dfrac1y=\dfrac{16}{3}$；
正中间 $x=\dfrac12$ 时等于 \underline{\hspace{2em}}，看起来最小。
于是猜想：\underline{\hspace{2em}} 时最小，最小值 \underline{\hspace{2em}}。

母题二：单位圆上找一点。试 $x=y=\dfrac{\sqrt2}{2}$，得
$\dfrac1{x^2}+\dfrac1{y^2}=4$；试 $x=\dfrac12,\ y=\dfrac{\sqrt3}{2}$，
得 $\dfrac{16}{3}$。故最小值是 \underline{\hspace{2em}}，
取到在 \underline{\hspace{2em}}。
\end{basiconly}

\begin{advancedonly}
\textbf{射影眼光看这两道母题}：第 (4) 问的比值 $t=y/x$ 是射影直线上的
坐标，齐次式只认射影点、不认代表元（同第 3 节的射影视角）。把母题一
的约束换成 $ax+by=1$ 后，(5)(b) 的取等条件 $\dfrac{ay}{x}=\dfrac{bx}{y}$
给出最优比例 $t^*=\sqrt{b/a}$，即最优射线的方向比为 $\sqrt b:\sqrt a$——
配平问题变成求一条射线与一条直线的交点。
\end{advancedonly}

\begin{conccard}{连环母题：做一题 = 学一节}{概念引入、性质翻译、规范书写、大招结论、公式族、充要边界、错题本七类笔记一次生成}{把母题当普通大题，只做最后一问}
\vchoice{做完这两道题，留白处就是一份《齐次化》随身笔记：第 (1)(2) 问
是概念与性质，第 (3) 问是考场书写，第 (4)(5) 问是大招与公式表，
第 (6) 问是逆向思维，第 (7) 问是错题本。}{基础较弱的同学先完成
第 (1)–(3) 问即可达标；第 (4)–(7) 问留给学有余力者。}
\end{conccard}

\begin{notezone}{章末归档：齐次化二级结论 + 七问笔记功能}
\textbf{二级结论}（方法链压缩版）：（1）齐次式的值只依赖比例 $t=y/x$；
（2）约束等式可以当 $1$ 乘进目标式完成升次（第 6 章入口、本章展开）；
（3）半次项平方可升为整次；（4）平移可以消去一次项，一般理论见
第 5 章；（5）分式迭代的不动点对应矩阵特征方向（第 5 章）。

\textbf{七问与笔记功能}：(1) 画图与数值感知 $\to$ 概念引入笔记；
(2) 三视角结构翻译 $\to$ 性质定义笔记；(3) 考纲规范通法 $\to$ 规范
书写笔记；(4) 高观点降维秒杀 $\to$ 二级结论与大招笔记；(5) 公式变形
家族 $\to$ 题型拓展笔记；(6) 参数变式与逆向追问 $\to$ 充要条件与逻辑
笔记；(7) 边界陷阱排雷 $\to$ 错题本易错点笔记。
\end{notezone}

\reading{拓展阅读 14}{齐次化：比例结构与射影桥梁}

% ==========================================================
% 第 9 章 不动点：迭代、方程与几何变换的交汇
% 说明：本文件为纯正文片段，由 main_body.tex 挂载；
%       使用 versions.tex 的 \tk / \vchoice / advancedonly /
%       basiconly / solution / 盒子与七问母题环境。
% ==========================================================

\section{不动点：迭代、方程与几何变换的交汇（高观点版）}

\begin{methodbox}{多框架对读}
\small
\begin{center}
\begin{tabular}{c|p{2.9cm}|p{5.0cm}|p{3.1cm}}
\toprule
视角 & 本章看点 & 别的视角怎么说 & 哪侧最省力 \\
\midrule
代数 & 解 $x=F(x)$、构造 $b_n=a_n-x^*$ & 把不动点当方程的根 & 求解析式 \\
几何 & 曲线与对角线的交点、变换的不动心 & 位置关系与轨迹形状 & 看收敛、判稳定 \\
集合 & 轨道 $\{a_n\}$ 与其聚点 & 值集、吸引域／排斥域 & 讨论存在与唯一 \\
映射 & 停留点 $F(p)=p$；Möbius 射影 & 换坐标系＝换表示 & 非线性 $\to$ 等比 \\
\bottomrule
\end{tabular}
\end{center}
卡的使用：四类不动点（数列／方程／变换／分式）共用同一句话——$F(p)=p$；算不下去就换表示（总图见第 0 章图~\ref{fig:four-view-map}）。
\end{methodbox}

\begin{warmupbox}{读法（先看这里）}
不动点只有一个定义：映射 $F$ 中满足 $F(p)=p$ 的点 $p$。
它有三副面孔：数列递推 $a_{n+1}=F(a_n)$（离散自迭代）、
方程 $f(x)=x$（曲线与对角线的交点）、几何变换（旋转中心、
位似中心、分式变换）。三者的共同动作是\key{换坐标系，把锚点
搬到原点}——这正是第 8 章齐次化、第 5 章平移／旋转的内核。
本章把这条线索单独展开，第 5 章与第 8 章的应用改为引用。
\end{warmupbox}

\begin{learnobjective}
\begin{itemize}
  \item 把数列看成 $\mathbb{N}\to\mathbb{R}$ 的（离散）函数，把递推看成映射的自迭代；
  \item 会用不动点方程 $x=F(x)$ 与稳定性判据 $|F'(x^*)|$ 处理一阶递推；
  \item 理解连续情形 $f(x)=x$ 的图解法与第 6 章隐零点、切线放缩的联系；
  \item 知道旋转中心、位似中心、仿射定点都是不动点的特例；
  \item 理解齐次坐标 $(u,v)$ 为什么能把分式迭代线性化（详见第 5 章）。
\end{itemize}
\end{learnobjective}

\subsection{第 1 节 数列是自然数集上的函数：递推即迭代}

数列 $a_n$ 本质是函数 $a:\mathbb{N}\to\mathbb{R}$，$n\mapsto a_n$。
一阶递推 $a_{n+1}=F(a_n)$ 就是映射 $F$ 的\key{自迭代}：
$a_1=F(a_0),\ a_2=F(F(a_0)),\dots$，数列是 $F$ 的一条轨道。
于是
\[
  \text{不动点}\ x^*=F(x^*)
  \iff \text{点 }(x^*,x^*)\text{ 同时在 }F\text{ 的图象与对角线 }y=x\text{ 上}.
\]
\textbf{稳定性判据}：在 $x^*$ 附近线性化 $F(x)\approx x^*+F'(x^*)(x-x^*)$，
故
\[
  |F'(x^*)|<1\Rightarrow\text{收缩（稳定）},\quad
  |F'(x^*)|>1\Rightarrow\text{扩张（不稳定）},\quad
  |F'(x^*)|=1\Rightarrow\text{看高阶项（临界）}.
\]

\begin{figure}[htbp]
  \centering
  \begin{minipage}{0.48\textwidth}
    \centering
    \begin{tikzpicture}[x=0.9cm,y=0.9cm,>=Stealth,semithick,font=\footnotesize]
      \draw[->,gray] (-1,0)--(3,0);
      \draw[->,gray] (0,-1)--(0,3);
      \draw[gray!60,dashed,domain=-1:3] plot(\x,\x);
      \draw[main,very thick,domain=-1:3] plot(\x,{0.5*\x+1});
      \fill[second] (2,2) circle (1.6pt) node[above right] {$x^*=2$};
      \draw[third,thick] (0,0)--(0,1)--(1,1)--(1,1.5)--(1.5,1.5)--(1.5,1.75)--(1.75,1.75);
      \node[gray!70!black] at (1.5,-0.65) {$F(x)=\tfrac12x+1$：收缩};
    \end{tikzpicture}
  \end{minipage}\hfill
  \begin{minipage}{0.48\textwidth}
    \centering
    \begin{tikzpicture}[x=0.9cm,y=0.9cm,>=Stealth,semithick,font=\footnotesize]
      \draw[->,gray] (-1,0)--(3,0);
      \draw[->,gray] (0,-1)--(0,3);
      \draw[gray!60,dashed,domain=-1:3] plot(\x,\x);
      \draw[main,very thick,domain=-0.5:2] plot(\x,{2*\x-1});
      \fill[second] (1,1) circle (1.6pt) node[below left] {$x^*=1$};
      \draw[third,thick] (1.2,1.2)--(1.2,1.4)--(1.4,1.4)--(1.4,1.8)--(1.8,1.8)--(1.8,2.6);
      \node[gray!70!black] at (1.5,-0.65) {$G(x)=2x-1$：扩张};
    \end{tikzpicture}
  \end{minipage}
  \caption{不动点的两种命运：同一张「$F$ 的图象＋对角线」图上，左图 $|F'|<1$，蛛网阶梯螺旋收拢到 $x^*$；右图 $|G'|>1$，阶梯被弹开。斜率与对角线 $y=x$ 的比较就是稳定性判据}
  \label{fig:cobweb-pair}
\end{figure}

\begin{methodbox}{不动点法三步}
\begin{enumerate}[label=(\arabic*), itemsep=2pt]
  \item \key{找}：解 $x^*=F(x^*)$（可能不止一个）；
  \item \key{判}：算 $|F'(x^*)|$ 定稳定／不稳定；
  \item \key{化}：令 $b_n=a_n-x^*$，则 $b_{n+1}\approx F'(x^*)\,b_n$；
        稳定时 $b_n\to0$，并可用等比模型写出通项。
\end{enumerate}
常见变体：分式递推用 $\dfrac{a_n-x_1}{a_n-x_2}$ 作商（第 5 章探究五），
双曲型递推用三角／指数代换。
\end{methodbox}

\subsection{第 2 节 连续情形：方程 $f(x)=x$ 与图解法}

把 $f(x)=x$ 改写成 $\varphi(x)=f(x)-x=0$：不动点就是 $\varphi$ 的零点。
几何上，它是曲线 $y=f(x)$ 与直线 $y=x$ 的交点。
\begin{example}[方程 $x=\cos x$ 在 $(0,\tfrac\pi2)$ 内有唯一解]
\vspace{-1mm}
\begin{solution}
$\varphi(x)=\cos x-x$，$\varphi(0)=1>0$，$\varphi(\tfrac\pi2)=-\tfrac\pi2<0$，
由零点定理至少有一解；又 $\varphi'(x)=-\sin x-1<0$ 严格递减，故唯一。
进一步 $|\varphi'|$ 与压缩性：$f(x)=\cos x$ 的导数 $|f'|=\sin x<1$ 在
$(0,\tfrac\pi2)$ 成立，故 $f$ 是压缩映射，迭代 $x_{n+1}=\cos x_n$
对任意初值都收敛到该唯一不动点——这就是第 6 章「隐零点＋单调性」
的几何版本。
\end{solution}
\end{example}
\textbf{与考纲的接口}：求超越方程的近似解（二分法）本质是不断把区间
压缩到不动点；证明 $f(x)\ge x$ 型不等式（如 $e^x\ge x+1$）则是在
比较曲线与对角线的位置（见第 6 章切线放缩）。

\subsection{第 3 节 几何变换的不动点：旋转中心、位似中心与仿射定点}

\begin{itemize}
  \item \key{旋转}：转角 $\theta\not\equiv0\pmod{2\pi}$ 的平面旋转只有
        \textbf{中心}一个不动点；
  \item \key{位似}：位似中心不动，其余点沿射线伸缩；
  \item \key{平移}：平面上没有不动点（它的「不动点」在无穷远处）；
  \item \key{仿射变换} $(x,y)\mapsto A\binom xy+\vec b$：不动点满足
        $(A-I)\binom xy=-\vec b$，当 $\det(A-I)\neq0$ 时唯一。
        椭圆的中心就是两方向伸缩变换的公共不动点（见第 4 章）。
\end{itemize}
直觉：任何「保结构」的变换都可以先找到它的不动点，再把坐标原点
搬过去——剩下的表达式往往立刻简化（第 8 章的平移手术刀、第 5 章
的旋转消交叉项都是这一招）。

\subsection{第 4 节 分式迭代与齐次坐标：射影视角}

一次分式变换 $x\mapsto\dfrac{ax+b}{cx+d}$ 本身不是线性的，但把 $x$
写成两个量的比 $x=u/v$ 之后，分子分母各自线性组合：
\[
  x_{n+1}=\frac{a\,x_n+b}{c\,x_n+d}
  \iff
  \binom{u_{n+1}}{v_{n+1}}
  =\begin{pmatrix}a&b\\ c&d\end{pmatrix}\binom{u_n}{v_n},
  \qquad x_n=\frac{u_n}{v_n}.
\]
关键在第二步：点 $(u,v)$ 与 $\lambda(u,v)$ 表示同一个 $x$（射影点），
所以「二维向量的线性作用」可以「投影」为一维分式作用。第 5 章
探究三已由此证明：\key{分式迭代的不动点 $=$ 矩阵的特征方向}
（推导与数值例子见第 5 章，此处不重复）。这里只看一件事：为什么
$(u,v)$ 与 $\lambda(u,v)$ 是同一点。

\begin{figure}[htbp]
  \centering
  \begin{tikzpicture}[>=Stealth,semithick,font=\footnotesize,x=1.1cm,y=1.1cm]
    \draw[->,gray] (-0.5,0)--(4.6,0) node[right] {$u$};
    \draw[->,gray] (0,-0.5)--(0,3.1) node[above] {$v$};
    \draw[main!60!black,dashed,very thick] (0,1)--(4.5,1) node[right] {$v=1$（投影线）};
    % 两条射线：同一条射线上取两个代表元
    \draw[second,very thick,->] (0,0)--(2.1,1.05);
    \draw[third,very thick,->] (0,0)--(1.05,2.1);
    \foreach \p/\lab/\pos in {(2,1)/{$x=2$}/above right,(1,2)/{$x=\tfrac12$}/above left}
      \fill[main] \p circle (1.6pt) node[\pos] {\lab};
    \fill (0,0) circle (1.2pt) node[below left] {$O$};
    \node[gray!70!black,align=center] at (2.2,-0.7) {同一条过原点的射线上的 $(u,v)$、$(2u,2v)$\\投影到 $v=1$ 上得到同一个 $x=u/v$};
  \end{tikzpicture}
  \caption{齐次坐标：用有序数对 $(u,v)$ 刻画比值 $x=u/v$。射线上的所有代表元落在投影线 $v=1$ 的同一点上，于是「一维分式迭代」升级为「二维线性迭代」}
  \label{fig:ray-projective}
\end{figure}

\textbf{双不动点的几何图像}：设 $M$ 的两个不动点为 $p,q$（射影直线上）。
取新坐标
\[
  X=\frac{x-p}{\,x-q\,},
\]
它把 $p$ 送到 $0$、$q$ 送到 $\infty$；在此坐标下 $M$ 退化为\emph{纯伸缩}
$X\mapsto\lambda X$，其中 $\lambda$ 正是转移矩阵两特征值之比（第 5 章探究五）。
于是
\[
  |\lambda|<1\Rightarrow X\to0\ \ (x\to p),\qquad
  |\lambda|>1\Rightarrow X\to\infty\ \ (x\to q).
\]
换言之：\key{两条特征方向在射影直线上就是两个不动点}；以 $p,q$ 为
$0,\infty$ 建立坐标后，分式迭代变成等比数列——这就是共轭变换的几何全貌。

\begin{figure}[htbp]
  \centering
  \begin{minipage}{0.48\textwidth}
    \centering
    \begin{tikzpicture}[>=Stealth,semithick,font=\footnotesize,x=1.2cm,y=1.2cm]
      \draw[->,gray] (-3.2,0)--(3.6,0);
      \node[gray!70!black,right] at (3.6,0) {$x$（射影直线）};
      \fill[main] (-2,0) circle (1.5pt) node[below] {$p$};
      \fill[second] (1.4,0) circle (1.5pt) node[below] {$q$};
      \draw[third,thick,->] (-1.2,0.28)--(-0.45,0.28);
      \draw[third,thick,->] (0.1,0.28)--(0.9,0.28);
      \node[gray!70!black,align=center] at (0.2,-1.0)
        {区间 $(p,q)$ 内 $X<0$，$|X|$ 越滚越大，\\$x$ 被拉向 $q$（$|\lambda|>1$）};
    \end{tikzpicture}
  \end{minipage}\hfill
  \begin{minipage}{0.48\textwidth}
    \centering
    \begin{tikzpicture}[>=Stealth,semithick,font=\footnotesize,x=1.0cm,y=1.0cm]
      \draw[->,gray] (-2.4,0)--(2.6,0) node[right] {$u$};
      \draw[->,gray] (0,-0.8)--(0,2.6) node[above] {$v$};
      \draw[main!60!black,dashed,thick] (-2.2,1)--(2.4,1) node[right] {$v=1$};
      \node[gray!70!black,left] at (-0.02,1) {$1$};
      \draw[main,very thick,->] (0,0)--(-1.05,0.55);
      \draw[main,very thick,dashed] (0,0)--(-2.1,1.05);
      \draw[second,very thick,->] (0,0)--(2.4,1.2);
      \fill[main] (-0.6,0.3) circle (1.4pt) node[below] {$(u,v)$};
      \node[main!70!black,above] at (-1.5,0.82) {$\times\lambda$};
      \fill[third] (-2,1) circle (1.4pt) node[below] {$p$};
      \fill[third] (2,1) circle (1.4pt) node[below] {$q$};
      \node[gray!70!black,align=center] at (0,-0.55)
        {两条特征直线 $\leftrightarrow$ 两个不动点};
    \end{tikzpicture}
  \end{minipage}
  \caption{双不动点的几何：左图在射影直线上，$X=(x-p)/(x-q)$ 使 $p\leftrightarrow0$、
  $q\leftrightarrow\infty$，迭代退化为 $X\mapsto\lambda X$（图示 $|\lambda|>1$，轨道被拉向 $q$）；
  右图回到 $(u,v)$ 平面，两条过原点的特征直线分别对应 $p,q$，作用即沿射线作伸缩}
  \label{fig:mobius-two-fixed}
\end{figure}

\subsection{第 5 节 七问连环母题：解答即笔记}

\refstepcounter{hvmasterq}
\begin{masterq}{不动点与收敛：$a_{n+1}=\dfrac12a_n+1$，$a_0=0$}
\begin{enumerate}[label=(\arabic*), itemsep=3pt]
\item \textbf{画图与数值感知}：算出 $a_1,a_2,a_3,a_4$，描点并猜极限。\logotag{核心公理}
\item \textbf{多视角结构翻译}：把该递推分别看成（i）$\mathbb{N}\to\mathbb{R}$
的函数迭代、（ii）平面直角坐标系中曲线 $y=\tfrac12x+1$ 与直线 $y=x$
的交点问题。\logotag{结构翻译}
\item \textbf{考纲规范通法}：用不动点法求通项，写出关键步骤
（令 $b_n=a_n-x^*$，\tk{} 填空）。\logotag{考纲通法}
\item \textbf{高观点降维秒杀}：用 $|F'(x^*)|$ 直接写出
$|a_n-x^*|=2\cdot2^{-n}$，并与第 6 章局部线性化对读。\logotag{大招公式}
\item \textbf{公式变形家族生成}：推导 $a_{n+1}=pa_n+q$ 的完整公式族
（$p\neq1$、$p=1$、$|p|>1$ 三种情形）。\logotag{公式族}
\item \textbf{参数变式与逆向追问}：由 $a_2=\tfrac32$、$a_3=\tfrac74$
反求 $p,q$；再回答：若 $\lim a_n=5$ 且 $p=\tfrac12$，求 $q$。\logotag{充要边界}
\item \textbf{边界陷阱与错题排雷}：某同学习惯「设 $\lim a_n=L$，代入得
$L=\tfrac12L+1$」，请指出缺少哪个前提，并举反例说明。\logotag{避坑铁律}
\end{enumerate}
\end{masterq}

\begin{solution}
(1) $a_1=1,\ a_2=\dfrac32,\ a_3=\dfrac74,\ a_4=\dfrac{15}8$，
趋向 $2$。

(2) （i）$F(x)=\tfrac12x+1$ 的自迭代；（ii）交点为
联立 $y=\tfrac12x+1$ 与 $y=x$ 的解 $(2,2)$。

(3) $x^*=\tfrac12x^*+1\Rightarrow x^*=2$。令 $b_n=a_n-2$，则
$b_{n+1}=a_{n+1}-2=\tfrac12(a_n-2)=\tfrac12 b_n$，$b_0=-2$，故
$b_n=-2\left(\tfrac12\right)^n$，
\[
  a_n=2-2^{\,1-n}.
\]
检验：$n=0$ 时 $a_0=0$，$n=1$ 时 $a_1=1$。

(4) $|F'(2)|=\tfrac12<1$，故 $|a_n-2|=|b_n|=2\cdot2^{-n}$ 指数衰减；
第 6 章的局部线性化 $F(x)\approx 2+\tfrac12(x-2)$ 正是这一步的连续版。

(5) $a_{n+1}=pa_n+q$：$x^*=px^*+q\Rightarrow x^*=\dfrac{q}{1-p}\ (p\neq1)$，
$a_n=x^*+p^n(a_0-x^*)$；$p=1$ 时退化为等差 $a_n=a_0+nq$；
$|p|>1$ 且 $a_0\neq x^*$ 时 $|a_n|\to\infty$。

(6) 由 $a_2=pa_1+q$、$a_3=pa_2+q$ 相减得
$p=\dfrac{a_3-a_2}{a_2-a_1}=\dfrac{1/4}{1/2}=\dfrac12$，
$q=a_2-pa_1=\dfrac32-\dfrac12=1$。
若 $\lim a_n=5$ 且 $p=\tfrac12$，则 $5=\dfrac q{1-1/2}=2q$，故 $q=\dfrac52$。

(7) 「取极限」需要先证明极限存在（单调有界或压缩映射）；反例：
$a_{n+1}=-a_n$、$a_0=1$ 振荡无极限，而方程 $x=-x$ 有唯一解 $0$——
\textbf{不动点存在不等于轨道收敛}。
\end{solution}

\refstepcounter{hvmasterq}
\begin{masterq}{分式迭代：$x_{n+1}=1+\dfrac1{x_n}$，$x_0=1$}
\begin{enumerate}[label=(\arabic*), itemsep=3pt]
\item \textbf{画图与数值感知}：算出 $x_1,x_2,x_3,x_4$，观察奇数项与
偶数项各自的走向。\logotag{核心公理}
\item \textbf{多视角结构翻译}：把它写成连分数 $x=1+\dfrac1{1+\dfrac1{1+\cdots}}$
与二次方程 $x^2-x-1=0$。\logotag{结构翻译}
\item \textbf{考纲规范通法}：证明奇数子列递增、偶数子列递减且有共同
极限 $\varphi=\dfrac{1+\sqrt5}{2}$（用 $\tk{}$ 搭骨架）。\logotag{考纲通法}
\item \textbf{高观点降维秒杀}：指出这是矩阵 $\begin{pmatrix}1&1\\1&0\end{pmatrix}$
的迭代，特征值之比给出收敛速率 $\dfrac{1}{\varphi^2}$（见第 5 章）。\logotag{大招公式}
\item \textbf{公式变形家族生成}：证明 $x_n=\dfrac{F_{n+2}}{F_{n+1}}$（$F_n$ 为
斐波那契数），并由此得到 $\varphi$ 的连分数与渐进分数。\logotag{公式族}
\item \textbf{参数变式与逆向追问}：若 $x_0=\varphi$ 会怎样？求使 $x_2=\varphi$
的 $x_0$；再把递推改成 $x_{n+1}=k+\dfrac1{x_n}$，求不动点与收敛范围。\logotag{充要边界}
\item \textbf{边界陷阱与错题排雷}：某同学由 $\lim x_n=L$ 直接写
$L=1+\dfrac1L$。请指出两处逻辑漏洞（存在性、正负号），并举出
$x_0<0$ 时的反例。\logotag{避坑铁律}
\end{enumerate}
\end{masterq}

\begin{solution}
(1) $x_1=2,\ x_2=\dfrac32,\ x_3=\dfrac53,\ x_4=\dfrac85$；奇数项递减、
偶数项递增，夹在 $1.6$ 附近。

(2) 迭代不变量：若 $x_n>0$，则 $x_{n+1}>1$；不动点满足
$x=1+\dfrac1x\Rightarrow x^2-x-1=0\Rightarrow \varphi=\dfrac{1+\sqrt5}{2}$
（负根 $-\dfrac1\varphi<0$ 舍去，因 $x_n>0$）。

(3) 设 $f(x)=1+\dfrac1x$，则 $x_{n+2}=f(f(x_n))=\dfrac{2x_n+1}{x_n+1}$。
直接证明：$x_{2k}$ 递增、$x_{2k+1}$ 递减，且
$|x_{n+1}-\varphi|=\dfrac{|x_n-\varphi|}{x_n\varphi}\le\dfrac1\varphi|x_n-\varphi|$
（因为 $x_n\ge1$），故两者夹逼到 $\varphi$。

(4) $\binom{u_{n+1}}{v_{n+1}}=\begin{pmatrix}1&1\\1&0\end{pmatrix}\binom{u_n}{v_n}$，
$x_n=u_n/v_n$；矩阵特征值为 $\varphi$ 与 $-1/\varphi$，故
$|x_n-\varphi|\sim C\varphi^{-2n}$（公比 $\left|\lambda_2/\lambda_1\right|
=\dfrac1{\varphi^2}$，见第 5 章探究五）。

(5) 归纳（取 $F_0=0,F_1=1$）：$x_0=\dfrac{F_2}{F_1}=\dfrac11=1$，且
$1+\dfrac1{x_n}=1+\dfrac{F_{n+1}}{F_{n+2}}
=\dfrac{F_{n+2}+F_{n+1}}{F_{n+2}}=\dfrac{F_{n+3}}{F_{n+2}}$，
故 $x_n=\dfrac{F_{n+2}}{F_{n+1}}$。

(6) $x_0=\varphi\Rightarrow$ 常数列 $\varphi$；反向递推
$x_n=\dfrac1{x_{n+1}-1}$，由 $x_2=\varphi$ 得 $x_1=\varphi$、$x_0=\varphi$。
一般地 $x_{n+1}=k+\dfrac1{x_n}$ 的不动点为
$\dfrac{k+\sqrt{k^2+4}}2$，$k>0$ 时正根稳定（$|f'|=1/x^{*2}<1$）。

(7) 漏洞一：未先证明极限存在；漏洞二：取极限时未说明保号。
方程 $L=1+\dfrac1L$ 的负根 $-\dfrac{1-\sqrt5}{2}$ 本身合法，
但它与 $x_n>0$ 的保号性矛盾，不能作为数列的极限。反例：
$x_0=-1$ 时 $x_1=0$，$x_2$ 无定义——定义域与符号必须先交代。
\end{solution}

\begin{conccard}{不动点：一把万能锚}{递推求通项与收敛性、超越方程近似解、几何变换的简化、分式迭代}{把「不动点存在」当成「轨道收敛」；未验证极限存在就取极限；忽略定义域与符号}
\vchoice{遇到「迭代／递推／方程 $f(x)=x$／变换简化」，先找不动点，
再判 $|F'|$，最后决定坐标系搬到哪里。}{记住三步：解 $x=F(x)$，
算 $|F'|$，搬原点。第 5、6、8 章都靠这一招联动。}
\end{conccard}

\begin{notezone}{章末归档：不动点二级结论 + 七问笔记功能}
\textbf{二级结论}：（1）一阶递推 $a_{n+1}=pa_n+q$ 的不动点
$x^*=\dfrac q{1-p}$，$a_n=x^*+p^n(a_0-x^*)$；
（2）稳定性 $|F'(x^*)|<1\Leftrightarrow$ 局部收缩；
（3）分式迭代的不动点＝矩阵特征方向（第 5 章）；
（4）齐次坐标 $(u,v)$ 把一维分式作用提升为二维线性作用；
（5）换坐标系（平移／旋转／齐次化）都是把不动点搬到原点。

\textbf{七问与笔记功能}：(1) 数值感知 $\to$ 概念引入；(2) 三视角
翻译 $\to$ 性质定义；(3) 考纲通法 $\to$ 规范书写；(4) 高观点秒杀
$\to$ 大招笔记；(5) 公式族 $\to$ 题型拓展；(6) 逆向追问 $\to$ 充要
边界；(7) 陷阱排雷 $\to$ 错题本。
\end{notezone}


% 示范章节（仅当 main.tex 中 \demomotrue 时挂载）
\ifdemomode
  % ==========================================================
% 示范章节：同一份源码，四种成品（约 40 行，只演示「怎么写」）
% 挂载：main.tex 中把 \demomodefalse 改为 \demomotrue，
%       或运行 python3 build_all.py --demo
% 覆盖写法：通用内容 + advancedonly + basiconly + \vchoice{}{}
%          + \tk{} + solution + warmupbox/trapbox/example
% ==========================================================
\section{示范：四合一多版本写法}

\begin{warmupbox}{读法（先看这里）}
本节不教新知识，只演示\textbf{版本分流写法}：修改 \verb|main.tex| 顶部的两行开关
（\verb|\teachertrue/\teacherfalse| 与 \verb|\basictrue/\basicfalse|），
即可从\emph{同一份源码}导出四种成品 PDF，不必复制、粘贴或维护任何副本。
\end{warmupbox}

设 $f(x)=x^{2}-2x$，考察它在区间 $[0,2]$ 上的结构。

\begin{learnobjective}
\begin{itemize}
  \item 会算 $f(0),f(1),f(2)$ 三个值并描点画草图；
  \item 会用集合数对、几何位置、代数方程三种语言描述同一条抛物线；
  \item 初步看懂「平方算子」$f\mapsto f^{2}$ 的零点行为。
\end{itemize}
\end{learnobjective}

\subsection*{第 (1) 问：具象感知与画图}
计算：$f(0)=\tk{0}$，$f(1)=\tk{-1}$，$f(2)=\tk{0}$；
描出 $(0,0),(1,-1),(2,0)$ 三点，并画出草图。

\vchoice{%
  三个函数值对应三个点；用集合数对写即
  $\{(x,y)\mid y=x^{2}-2x,\ x\in[0,2]\}$，几何上是开口向上、对称轴为
  $x=1$ 的一段抛物线。}{%
  三个点：$(0,0)$、$(1,-1)$、$(2,0)$；连起来是开口向上的小弧线，
  最低点在 $x=1$ 处。}

\subsection*{第 (2) 问：常规通法演算}
配方：$f(x)=x^{2}-2x=\tk{(x-1)^{2}-1}$，故顶点为 $\tk{(1,-1)}$，
对称轴为直线 $\tk{x=1}$。

\subsection*{第 (3) 问：三视角结构翻译}
同一条曲线可以写成三种语言：集合语言
$\{(x,f(x))\mid x\in\mathbb{R}\}$；几何语言「以 $x=1$ 为对称轴的抛物线」；
代数方程 $y=(x-1)^{2}-1$。三者说的是\emph{同一结构}。

\begin{advancedonly}
\subsection*{第 (4) 问：高观点降维秒杀（仅完整版可见）}
作平移 $u=x-1$，则 $y=u^{2}-1$：绝对尺度被消去，只剩一个比例变量 $u$，
这就是齐次化与平移归一的思想雏形。进一步看平方算子：
$f^{2}(x)=(x(x-2))^{2}$ 的零点是\textbf{水平相切的偶重零点}
（局部 $\approx 4(x-1)^{2}$ 的二重根平移），不是尖点——它解释了
「平方后符号不变、零点处压平」的现象。
\end{advancedonly}

\begin{basiconly}
\subsection*{第 (4)' 问：大白话口诀与半成品填空（仅薄弱生版可见）}
口诀：「平方不改号，零点压平地；倒数撕裂缝，无穷变成零。」
照样子填空：$f(1)=1-2=\tk{-1}$；顶点横坐标 $x=\tk{1}$。
\end{basiconly}

\begin{example}[用两种语言看同一结构]
\vchoice{双射同构映射}{一一对应的加工机器}的说法在本讲统一替换为
\vchoice{二次型结构}{抛物线的对称结构}。请用「顶点式」写出
$f(x)=x^{2}-4x+3$ 的对称轴：$\tk{x=2}$。
\end{example}
\begin{solution}
配方得 $f(x)=(x-2)^{2}-1$，故对称轴为 $x=2$，顶点 $(2,-1)$。
\end{solution}

\begin{trapbox}{顶点横坐标的符号}
配方时括号内是 $x-1$ 而不是 $x+1$：$y=(x-1)^{2}-1$ 的对称轴是 $x=1$。
若误写为 $x=-1$，是把一次项系数的符号弄反了。
\end{trapbox}

\fi

\appendix
% 附录 A：高观点结论的考纲合法化转译手册
\section{附录 A：高观点结论的考纲合法化转译手册}\label{app:legal}

\begin{commonerror}{铁律：草稿纸用高观点，答题卡写初等通法}
解答题中\textbf{严禁}出现「由仿射变换知」「由极点极线知」「由欧拉公式知」「由 Abel 变换知」「由洛必达法则知」等超纲字眼。这类词一旦出现，轻则扣步骤分，重则该问判零。高观点的正确用法是：在草稿纸上用它锁定答案、取等位置与参数边界，然后在卷面上用考纲认可的初等推导把结论\textbf{重新证一遍}——代数验算本身就是满分证明。
\end{commonerror}

\begin{conccard}[三级]{三类角色各就各位}{阅卷老师只看卷面}{把预判当成证明}
\textbf{草稿纸}：随便用，越快越准越好；\textbf{卷面}：只允许考纲工具（联立、韦达定理、作差构造、错位相减、分类讨论、切线不等式）；\textbf{采分点}：看关键等式有没有被显式写出。三者不冲突，缺的只是翻译。
\end{conccard}

下面六个模板覆盖高中压轴的高频场景，每个都按「预判 $\to$ 卷面 $\to$ 采分点」三段给出。

\subsection*{模板一：圆锥曲线中的仿射压缩}

\textbf{草稿纸预判}：椭圆沿短轴方向压缩 $b/a$ 倍即得圆，面积同比例缩小，故内接三角形最大面积为 $\dfrac{3\sqrt3}{4}ab$；斜率乘积定值 $-\dfrac{b^2}{a^2}$。

\textbf{卷面写法}：设过定点的直线为 $y=kx+m$，代入 $\dfrac{x^2}{a^2}+\dfrac{y^2}{b^2}=1$，写出韦达定理与判别式条件；由弦长公式 $|AB|=\sqrt{1+k^2}\,|x_1-x_2|$ 与点到直线距离写出面积
\[
  S=\frac12|AB|\cdot d ,
\]
整理成关于 $k,m$ 的函数，用基本不等式或求导在预判的取等点处取得最值，并验证此时满足判别式大于零。

\textbf{采分点}：韦达定理两式、弦长表达式、距离表达式、「当且仅当……时取等」以及取等条件的检验。仿射结论只用来决定往哪条直线上凑。

\subsection*{模板二：极点与极线的定点定值}

\textbf{草稿纸预判}：由极点 $P(x_0,y_0)$ 直接写出极线方程 $\dfrac{x_0x}{a^2}+\dfrac{y_0y}{b^2}=1$，目标直线即刻锁定，不必再猜。

\textbf{卷面写法}：设两交点 $A(x_1,y_1),B(x_2,y_2)$，分别代入曲线方程并作差：
\[
  \frac{x_1^2-x_2^2}{a^2}+\frac{y_1^2-y_2^2}{b^2}=0
  \ \Longrightarrow\
  \frac{(x_1+x_2)(x_1-x_2)}{a^2}+\frac{(y_1+y_2)(y_1-y_2)}{b^2}=0 .
\]
再结合 $A,B$ 都在已知直线上，整理出 $\dfrac{y_1-y_2}{x_1-x_2}$ 与 $\dfrac{y_1+y_2}{x_1+x_2}$ 的关系，最后代入要证的对称式化简即可。

\textbf{采分点}：两点式作差这一步必须显式写出；之后全是恒等变形。

\subsection*{模板三：不动点与特征根求通项}

\textbf{草稿纸预判}：解 $x^*=px^*+q$ 或 $\lambda^2-p\lambda-q=0$，读出公比与两根。

\textbf{卷面写法}：由原递推式恒等变形直接写出
\[
  a_{n+1}-x^*=p(a_n-x^*) ,
\]
或两根情形下
\[
  \frac{a_{n+1}-\lambda_1}{\,a_{n+1}-\lambda_2\,}=k\,\frac{a_n-\lambda_1}{\,a_n-\lambda_2\,} .
\]
\textbf{完全不需要解释} $x^*$、$\lambda_i$ 是怎么来的——作为「已知常数」使用，再用首项确定待定系数，最后验证 $n=1$ 成立。

\textbf{采分点}：构造式本身、首项代入、$n=1$ 的检验。

\subsection*{模板四：Abel 变换与错位相减}

\textbf{草稿纸预判}：用分部求和快速算出 $S_n$ 的闭式。

\textbf{卷面写法}：规范列出
\[
  S_n=a_1+a_2+\cdots+a_n,\qquad
  qS_n=\phantom{a_1+{}}a_1q+\cdots+a_nq ,
\]
错开一位对齐后写「两式相减得」，直接默写化简结果，并把 $n=1$ 与 $n\to\infty$（若涉及极限）单独验证。

\textbf{采分点}：两行原始展开、相减动作、最终合并式、首项检验。

\subsection*{模板五：三角恒等变换与欧拉公式}

\textbf{草稿纸预判}：忘了三倍角公式时，展开 $(\cos\theta+i\sin\theta)^3$ 取实部，秒得 $\cos3\theta=4\cos^3\theta-3\cos\theta$。

\textbf{卷面写法}：从两角和公式出发现场推导：
\[
  \cos3\theta=\cos(2\theta+\theta)=\cos2\theta\cos\theta-\sin2\theta\sin\theta=\cdots,
\]
再降幂为单角函数，最后写出所需形式。

\textbf{采分点}：起始公式明确、每一步降幂正确、最终形式与题目要求一致（平方关系、次数）。

\subsection*{模板六：导数压轴的隐零点与切线放缩}

\textbf{草稿纸预判}：用切线 $e^x\ge x+1$、$\ln x\le x-1$ 预判参数临界值 $a_0$；必要时用端点效应估出 $a_0$。

\textbf{卷面写法}：按 $a\le a_0$ 与 $a>a_0$ 分类讨论。对每一类，先证切线不等式（以 $e^x-x-1$ 求最小值的方式），再代入原函数完成放缩链；隐零点处写成「设 $x_0$ 满足 $f'(x_0)=0$，则……」，用整体代换避免解出具体数值。

\textbf{采分点}：分类标准清晰、切线不等式的完整证明、隐零点的设而不求与代换、端点值的检验。

\subsection*{自查清单}

\begin{itemize}
  \item 全篇是否出现超纲名词？若有，一律替换为等价代数操作。
  \item 关键等式是否都被显式写出（而不是「显然」「易得」）？
  \item 取等条件、判别式正负、$n=1$ 首项是否逐一检验？
  \item 参数分类是否有重叠或遗漏？
\end{itemize}

\begin{conccard}[一级]{一句话记住}{预判给方向，通法给分数}{两者都要}
\key{草稿纸负责快，卷面负责稳}。高观点决定你往哪条路上走，考纲通法决定你能拿多少分——只有前者，会丢步骤分；只有后者，会算不完。
\end{conccard}

% ==========================================================
% 附录 B：大学理工科荣誉课程（Honors Bridge）先修接口对照手册
% 说明：本文件为纯正文模块化片段（无 \documentclass），
%       由 main.tex 在 \appendix 之后 \input 挂载。
%       corebox / methodbox 现由 versions.tex 全局提供（四合一引擎），
%       配色与现有 main/third 主题一致，不再局部定义。
% ==========================================================

\section{附录 B：大学理工科荣誉课程先修接口对照手册}\label{app:honors}

\begin{corebox}{致准大学荣誉学院（Honors Program）学子}
普通高中生升入数学、物理、计算机、人工智能与电子信息等理工科专业时，普遍会经历一场剧烈的\textbf{“大一数学休克”}——例如看不懂《数学分析》严密的 $\varepsilon$-$\delta$ 极限语言，无法理解《线性代数》中抽象的基底变换与特征空间。

然而，如果你已经扎实地走完了本套讲义的第 0–9 章，\textbf{你实际上已经提前掌握了大学前两年硬核数学课的几何直觉与结构骨架}。本附录旨在为你完成最后一步\textbf{“学术暗号对接”}，将你熟悉的中学高观点，精准映射到大学专业课的正式定义与定理名称上。
\end{corebox}

\subsection*{一、 五大大学核心基础课“概念暗号”对接总表}

\begin{fastcard}{《数学分析》（Real Analysis / Honors Calculus）}
\textbf{你在本套讲义中学过的直觉与工具}：第 1 章有理数 $\mathbb{Q}$ 漏掉 $\sqrt{2}$ 孔洞，无理数填满实心数轴 $\mathbb{R}$；第 6 章局部切线近似（以直代曲）与切线支撑放缩。
\par\textbf{大学专业课正式学名与定理}：\textbf{实数完备性公理}（戴德金分割 / 确界原理）；\textbf{介值定理}与\textbf{微分中值定理}；\textbf{凸函数与 Jensen 不等式}。
\par\textbf{消除的“大一休克”痛点}：瞬间理解极限存在性证明，不再死背 $\varepsilon$-$\delta$ 符号，一眼看穿微分局部线性化本质。
\end{fastcard}

\begin{fastcard}{《线性代数》（Linear Algebra / 高等代数）}
\textbf{你在本套讲义中学过的直觉与工具}：第 5 章换基底对齐特征方向、旋转 $45^\circ$ 消除交叉项 $Bxy$；第 1 章函数变换中的主动（动图形）与被动（动坐标系）对偶；第 5 章转移矩阵特征向量锁定分式不动点。
\par\textbf{大学专业课正式学名与定理}：\textbf{线性空间与基底变换}（过渡矩阵 $P^{-1}AP$）；\textbf{实对称矩阵正交对角化}（谱定理）；\textbf{射影空间（Projective Space）}与齐次坐标。
\par\textbf{消除的“大一休克”痛点}：摆脱“矩阵只是数字表格”的机械计算误区，直接建立高维几何变换与特征空间直觉。
\end{fastcard}

\begin{fastcard}{《复变函数与信号系统》（Complex Analysis \& Signals）}
\textbf{你在本套讲义中学过的直觉与工具}：第 2 章虚数 $i$ 作为逆时针 $90^\circ$ 旋转算子；切线垂直速度演化出欧拉公式 $e^{i\theta}=\cos\theta+i\sin\theta$。
\par\textbf{大学专业课正式学名与定理}：\textbf{复数域的几何表示与保角映射}；\textbf{傅里叶级数与傅里叶变换}（正交旋转基展开）；工科电路\textbf{相量法（Phasor）}。
\par\textbf{消除的“大一休克”痛点}：在大学物理简谐振动、交流电路与信号处理中，直接把正弦量当成复平面旋转向量处理。
\end{fastcard}

\begin{fastcard}{《常微分方程与数值分析》（ODEs \& Numerical Methods）}
\textbf{你在本套讲义中学过的直觉与工具}：第 5 章差分 $\Delta a_n$ 与求和互逆、Abel 分部求和；蛛网图 $|f'(x^*)|<1$ 螺旋收敛于吸引不动点；生成函数 $G(x)=\sum a_n x^n$。
\par\textbf{大学专业课正式学名与定理}：\textbf{差分方程与常微分方程特征根法}；\textbf{巴拿赫压缩映射原理}（牛顿迭代法收敛性）；\textbf{Z 变换与拉普拉斯变换}。
\par\textbf{消除的“大一休克”痛点}：打通离散递归与连续微分方程，理解计算机迭代算法为何收敛以及混沌分岔的临界条件。
\end{fastcard}

\begin{fastcard}{《公理化概率论与机器学习》（Probability \& Machine Learning）}
\textbf{你在本套讲义中学过的直觉与工具}：第 7 章样本空间 $\Omega$ 归一化测度与条件子集缩放；二维散点云的竖直欧氏距离平方和最小投影线。
\par\textbf{大学专业课正式学名与定理}：\textbf{Kolmogorov 概率测度空间} $(\Omega, \mathcal{F}, P)$；\textbf{线性回归的最小二乘正交投影}（正规方程 $A^TA\hat{x}=A^Tb$）。
\par\textbf{消除的“大一休克”痛点}：跳过繁琐的古典计数，直接站在高维空间几何投影的高度理解统计学习与数据拟合。
\end{fastcard}

\subsection*{二、 三大核心概念的“中学直觉 $\to$ 大学严格化”跃迁示范}

为了让你在翻开大学教材的第一天就能写出严密的现代数学证明，以下展示三个最核心概念如何从本讲义的“几何直觉”平滑过渡到大学“公理化语言”：

\begin{methodbox}{跃迁 1：从“实心金条无漏洞”到《数学分析》的“确界原理”}
\begin{itemize}
    \item \textbf{中学高观点直觉}（第 1 章）：有理数轴 $\mathbb{Q}$ 像布满针眼的筛子，曲线 $y=x^2-2$ 会从 $\sqrt{2}$ 的孔洞漏过去；实数轴 $\mathbb{R}$ 填满了所有孔洞，因此有上界的递增序列必然会撞上一个确定的天花板。
    \item \textbf{大学《数学分析》严格表述（确界公理）}：设 $S \subset \mathbb{R}$ 为非空有上界实数集，则 $S$ 在 $\mathbb{R}$ 中必有\textbf{上确界（最小上界 $\sup S$）}。由此可直接推出“单调有界数列必有极限”与“闭区间套定理”。
\end{itemize}
\end{methodbox}

\begin{methodbox}{跃迁 2：从“旋转 $45^\circ$ 消交叉项 $xy$”到《线性代数》的“正交对角化”}
\begin{itemize}
    \item \textbf{中学高观点直觉}（第 5 章）：双曲线 $xy = 1$ 的主轴倾斜了 $45^\circ$，只要把坐标轴旋转 $45^\circ$ 对齐它的对称轴，交叉项就会消失，化为 $\frac{x'^2}{2} - \frac{y'^2}{2} = 1$。
    \item \textbf{大学《线性代数》严格表述（二次型主轴定理）}：将 $xy = 1$ 写成对称矩阵二次型：
    $$\begin{pmatrix} x & y \end{pmatrix} \begin{pmatrix} 0 & \frac{1}{2} \\ \frac{1}{2} & 0 \end{pmatrix} \begin{pmatrix} x \\ y \end{pmatrix} = 1$$
    求实对称矩阵 $A = \begin{pmatrix} 0 & 1/2 \\ 1/2 & 0 \end{pmatrix}$ 的特征值 $\lambda_1 = \frac{1}{2}, \lambda_2 = -\frac{1}{2}$ 及其单位正交特征向量（构成旋转 $45^\circ$ 的正交矩阵 $Q$），令 $\mathbf{x} = Q\mathbf{x}'$，立刻得到标准型 $\frac{1}{2}x'^2 - \frac{1}{2}y'^2 = 1$！
\end{itemize}
\end{methodbox}

\begin{methodbox}{跃迁 3：从“蛛网图台阶变小”到《泛函/数值分析》的“压缩映射原理”}
\begin{itemize}
    \item \textbf{中学高观点直觉}（第 5 章）：在蛛网图中，只要不动点处曲线斜率满足 $|f'(x^*)| < 1$，每次投影都会把误差缩短，序列最终螺旋收敛于 $x^*$。
    \item \textbf{大学严格表述（Banach Fixed-Point Theorem）}：设 $(X, d)$ 为完备度量空间，映射 $T: X \to X$ 满足存在常数 $0 \le L < 1$，使得对任意 $x, y \in X$ 均有 $d(T(x), T(y)) \le L \cdot d(x, y)$，则 $T$ 在 $X$ 中存在\textbf{唯一不动点 $x^*$}，且对任意初值 $x_0 \in X$，迭代序列 $x_{n+1} = T(x_n)$ 均以几何级数速度收敛于 $x^*$。
\end{itemize}
\end{methodbox}

\subsection*{三、 给荣誉学院新生的阅读进阶路线图}

当你顺利进入大学理工科荣誉班级后，推荐按以下书单无缝衔接本讲义的知识结构：
\begin{enumerate}
    \item \textbf{数学分析方向}：陶哲轩《陶哲轩实分析》（从集合与自然数公理重新严格建构整个实数与微积分体系，与本讲义第 0、1 章风格一脉相承）；
    \item \textbf{线性代数方向}：Sheldon Axler《Linear Algebra Done Right》（摒弃繁琐行列式，直接从向量空间与线性算子的不变子空间/特征向量切入，与本讲义换基底思想高度契合）；
    \item \textbf{复分析与几何直觉方向}：Tristan Needham《复分析：可视化方法》（用纯几何旋转与向量场讲透复变函数，完美延续本讲义第 2 章与 Manim 动画风格）。
\end{enumerate}


\end{document}
。
%
%  四种组合（改两行即可，无需复制任何文件）：
%    \teachertrue  \basicfalse   完整教师版（默认母版）
%    \teacherfalse \basicfalse   完整学生版（隐藏解析与填空答案）
%    \teachertrue  \basictrue    薄弱生教师版（显化脚手架 + 完整答案）
%    \teacherfalse \basictrue    薄弱生学生版（保留 \tk 留白横线）
%
%  本文件只依赖已加载的宏包（xcolor / tcolorbox / environ），
%  不修改 common/，故对 ../hv-bank 等其他分册零影响。
% ============================================================

% ---------- 布尔开关 ----------
\newif\ifteacher          % 教师解析版 / 学生留白版
\newif\ifbasic            % 基础通关版 / 完整高观点版
\newif\ifdemomode         % 是否挂载示范章节 chapters/demo_multiversion

% 安全默认（main.tex 的两行开关覆盖此处）
\teachertrue
\basicfalse
\demomodefalse

% ---------- 零残留隐藏所需 ----------
\usepackage{environ}      % \NewEnviron（隐藏时不留空白）

% ---------- 版本标签（封面 / 页眉自动显示） ----------
\newcommand{\versionlabel}{%
  \ifteacher 教师解析版\else 学生留白版\fi
  \,\textperiodcentered\,
  \ifbasic 基础通关版\else 完整高观点版\fi}

% ============================================================
%  一、版本内容分流环境
% ============================================================

% 仅完整版渲染；薄弱生版彻底忽略（不占位）
% 用途：莫比乌斯矩阵、射影同构、压轴第 3–4 阶梯难题等
\NewEnviron{advancedonly}{%
  \ifbasic\else
    \par\medskip
    \BODY
    \par\medskip
  \fi
}

% 仅薄弱生版渲染；完整版彻底忽略
% 用途：课前 15 分钟计算过关、大白话口诀、半成品步步填空
\NewEnviron{basiconly}{%
  \ifbasic
    \par\medskip
    \BODY
    \par\medskip
  \fi
}

% ============================================================
%  二、双模填空宏 \tk{...} 与内联版本选择 \vchoice{}{}
% ============================================================

% 教师版：红色下划线答案；学生版：等宽下划线留白（版面稳定）
% 内容一律写成裸数学式（不带 $）。\underline 继承当前模式，
% 而 \sbox 的盒子内容始终是文本模式、必须自带数学定界符。
\newcommand{\tk}[1]{%
  \ifteacher
    \ifmmode
      {\color{second}\underline{#1}}%
    \else
      {\color{second}\underline{$\displaystyle#1$}}%
    \fi
  \else
    \sbox0{$\displaystyle#1$}%
    \underline{\hspace{\wd0}}%
  \fi
}

% 同一句子里的措辞切换：完整版文本 / 薄弱生大白话文本
\newcommand{\vchoice}[2]{%
  \ifbasic #2\else #1\fi}

% ============================================================
%  三、解答环境重定义（教师版带边框【剖析与标准解答】）
% ============================================================
\RenewDocumentEnvironment{solution}{+b}{%
  \ifteacher
    \par\medskip
    \begin{tcolorbox}[
      breakable, enhanced, sharp corners,
      colback=main!3, colframe=main!60!black, boxrule=0.5pt,
      title={\bfseries 剖析与标准解答}, coltitle=black, fonttitle=\small,
      top=5pt,
    ]
    #1
    \end{tcolorbox}%
  \else
    % 学生版：不留内容，只保留少量书写留白
    \par\noteblank[6ex]%
  \fi
}{}

% ============================================================
%  四、语义化盒子（corebox / methodbox / trapbox / warmupbox）
% ============================================================

% 核心概念盒（绿）
\newtcolorbox{corebox}[1]{
  breakable, enhanced, sharp corners,
  colback=main!4, colframe=main!60!black, boxrule=0.6pt,
  title={\bfseries #1}, coltitle=black, fonttitle=\small,
  attach boxed title to top left={xshift=2mm, yshift=-2mm},
  boxed title style={sharp corners, colback=white,
    colframe=main!60!black, boxrule=0.4pt},
  top=7pt,
}

% 方法归纳盒（青）
\newtcolorbox{methodbox}[1]{
  breakable, enhanced, sharp corners,
  colback=third!5, colframe=third!60!black, boxrule=0.5pt,
  title={\bfseries #1}, coltitle=black, fonttitle=\small,
  attach boxed title to top left={xshift=2mm, yshift=-2mm},
  boxed title style={sharp corners, colback=white,
    colframe=third!60!black, boxrule=0.4pt},
  top=7pt,
}

% 易错陷阱盒（红）
\newtcolorbox{trapbox}[1]{
  breakable, enhanced, sharp corners,
  colback=second!4, colframe=second!65!black, boxrule=0.5pt,
  title={\bfseries 陷阱：#1}, coltitle=black, fonttitle=\small,
  attach boxed title to top left={xshift=2mm, yshift=-2mm},
  boxed title style={sharp corners, colback=white,
    colframe=second!65!black, boxrule=0.4pt},
  top=7pt,
}

% 课前热身盒（蓝）
\newtcolorbox{warmupbox}[1]{
  breakable, enhanced, sharp corners,
  colback=structurecolor!4, colframe=structurecolor!70!black, boxrule=0.5pt,
  title={\bfseries #1}, coltitle=black, fonttitle=\small,
  top=6pt,
}

% ============================================================
%  五、按章（节）自动编号的 example / exercise 环境
%  注：elegantnote 自带 example 计数器（\c@example），故另立
%      hvexample / hvexercise 计数，避免与模板冲突；编号仍按节累计。
% ============================================================
\newcounter{hvexample}[section]
\renewcommand{\thehvexample}{\thesection.\arabic{hvexample}}

\RenewDocumentEnvironment{example}{O{}}{%
  \refstepcounter{hvexample}%
  \par\medskip\noindent
  {\bfseries\color{main}例题\thehvexample\IfValueT{#1}{：#1}\quad}%
}{\par\medskip}

\newcounter{hvexercise}[section]
\renewcommand{\thehvexercise}{\thesection.\arabic{hvexercise}}

\NewDocumentEnvironment{exercise}{O{}}{%
  \refstepcounter{hvexercise}%
  \par\medskip\noindent
  {\bfseries\color{third}练习\thehvexercise\IfValueT{#1}{：#1}\quad}%
}{\par\medskip}

% ============================================================
%  六、七问连环母题（Scaffolded Master Problem）载体环境
% ------------------------------------------------------------
%  定位：一道母题挂载 (1)–(7) 个递进小问，学生做题留白处写下的
%        推理本身就是这一章的学习笔记。写作规范见 ARCHITECTURE.md。
% ============================================================

\newcounter{hvmasterq}[section]
\renewcommand{\thehvmasterq}{\thesection.\arabic{hvmasterq}}

% 母题盒：标题自动编号；盒内用 enumerate[label=(\arabic*)] 写七问
% 注意：计数器由调用方在 \begin{masterq} 前手动 \refstepcounter{hvmasterq}，
%       不要把 \refstepcounter 写进本定义的选项里（tcolorbox 会反复展开选项，
%       导致编译卡在分页计算）。
\newtcolorbox{masterq}[1]{
  breakable, enhanced, sharp corners,
  colback=structurecolor!3, colframe=structurecolor!70!black, boxrule=0.6pt,
  title={\bfseries 七问连环母题\thehvmasterq：#1},
  coltitle=black, fonttitle=\small,
  attach boxed title to top left={xshift=2mm, yshift=-2mm},
  boxed title style={sharp corners, colback=white,
    colframe=structurecolor!70!black, boxrule=0.4pt},
  top=7pt,
}

% 每问末尾的笔记归档标签：\logotag{存档类别}（如「核心公理」「大招公式」）
\newcommand{\logotag}[1]{%
  \hfill\textcolor{third}{\small\ttfamily【笔记归档·#1】}}
